% Standalone source. Compile with pdflatex (three passes recommended).
% All mathematical text is embedded; no external figures, inputs, or .bib file are required.
\documentclass[11pt,a4paper]{article}
\usepackage[T1]{fontenc}
\usepackage[utf8]{inputenc}
\usepackage{lmodern,amsmath,amssymb,amsthm,mathtools,booktabs,enumitem,longtable,array}
\usepackage[margin=26mm,headheight=15pt,headsep=8mm]{geometry}
\usepackage{microtype,xcolor,fancyhdr,titlesec}
\usepackage{etoolbox}
\usepackage[colorlinks=true,linkcolor=blue!40!black,urlcolor=blue!45!black,citecolor=blue!40!black,bookmarksnumbered=true]{hyperref}
\usepackage{bookmark}
\hypersetup{pdftitle={Sharp Bounds and Exact Classifications: An Integrated Research Manuscript},pdfauthor={KEIO},pdfsubject={Sixty contributions, corrected proofs, and further deductions}}
\newtheorem{theorem}{Theorem}[section]
\newtheorem{lemma}[theorem]{Lemma}
\newtheorem{proposition}[theorem]{Proposition}
\newtheorem{corollary}[theorem]{Corollary}
\theoremstyle{definition}\newtheorem{definition}[theorem]{Definition}
\theoremstyle{remark}\newtheorem{remark}[theorem]{Remark}
\newcommand{\R}{\mathbb R}\newcommand{\C}{\mathbb C}\newcommand{\N}{\mathbb N}
\newcommand{\Z}{\mathbb Z}\newcommand{\Q}{\mathbb Q}\newcommand{\F}{\mathbb F}
\newcommand{\Pcal}{\mathcal P}\newcommand{\PP}{\mathcal P}\newcommand{\Pp}{\mathcal P}
\newcommand{\E}{\mathbb E}\newcommand{\one}{\mathbf 1}\newcommand{\1}{\mathbf 1}
\newcommand{\Id}{\mathrm I}\newcommand{\D}{\operatorname{D}}
\newcommand{\dd}{\frac{d}{dx}}\newcommand{\Span}{\operatorname{span}}\newcommand{\Res}{\operatorname{Res}}
\DeclareMathOperator{\coker}{coker}\DeclareMathOperator{\diag}{diag}\DeclareMathOperator{\rank}{rank}
\DeclareMathOperator{\ord}{ord}\DeclareMathOperator{\lcm}{lcm}\DeclareMathOperator{\im}{im}
\DeclareMathOperator{\tr}{tr}\DeclareMathOperator{\supp}{supp}\DeclareMathOperator{\TV}{TV}
\DeclareMathOperator{\Lip}{Lip}\DeclareMathOperator{\Var}{Var}
\DeclareMathOperator*{\argmin}{argmin}\DeclareMathOperator{\conv}{conv}
\newcommand{\cF}{\mathcal F}\newcommand{\cH}{\mathcal H}
\newcommand{\ip}[2]{\langle #1,#2\rangle}\newcommand{\norm}[1]{\lVert #1\rVert}\newcommand{\abs}[1]{\lvert #1\rvert}
\numberwithin{equation}{section}
\setlength{\emergencystretch}{3em}
\setlength{\parskip}{3pt plus 1pt}
\setlength{\parindent}{1.2em}
\setcounter{tocdepth}{1}
\makeatletter
\renewcommand*\l@section{\@dottedtocline{1}{0em}{2.5em}}
\makeatother
\setcounter{secnumdepth}{3}
\titleformat{\section}{\Large\bfseries\raggedright}{\thesection}{0.7em}{}
\titleformat{\subsection}{\large\bfseries}{\thesubsection}{0.65em}{}
\titlespacing*{\section}{0pt}{2.5ex}{1.7ex}
\titlespacing*{\subsection}{0pt}{2.2ex}{1.2ex}
\pagestyle{fancy}\fancyhf{}
\renewcommand{\sectionmark}[1]{}
\renewcommand{\subsectionmark}[1]{}
\fancyhead[L]{\small\textsc{Sharp Bounds and Exact Classifications}}
\fancyhead[R]{\small\nouppercase{\rightmark}}
\fancyfoot[C]{\thepage}
\renewcommand{\headrulewidth}{0.3pt}
\fancypagestyle{plain}{\fancyhf{}\fancyfoot[C]{\thepage}\renewcommand{\headrulewidth}{0pt}}
\raggedbottom
\begin{document}
\hypersetup{pageanchor=false}
\begin{titlepage}
\centering
\vspace*{20mm}
{\small\scshape Integrated research manuscript\par}
\vspace{12mm}
{\Huge\bfseries Sharp Bounds and\par Exact Classifications\par}
\vspace{9mm}
{\Large Algebra, Analysis, Information,\par and Electrical Networks\par}
\vspace{12mm}
{\large Sixty contributions with corrected proofs\par and further deductions\par}
\vspace{15mm}
{\Large KEIO\par}
\vspace{4mm}
{\normalsize Consolidated edition: 24 October 2026\par}
\vfill
\begin{minipage}{0.86\textwidth}\small
This standalone manuscript consolidates the complete mathematical content of six ten-stage research programmes: Showa, Meiji, Keio, Taisho, Reiwa, and Heisei. These names identify source programmes. The supplied corrections are incorporated into the relevant statements and examples. Further deductions are separated from the sixty source contributions and include complete proofs. Historical priority and external peer review are not asserted by this consolidation.
\end{minipage}
\vspace{10mm}
\end{titlepage}
\pagenumbering{roman}
\hypersetup{pageanchor=true}
\markboth{Abstract}{Abstract}
\section*{Abstract}
\phantomsection\addcontentsline{toc}{section}{Abstract}
This manuscript develops exact classifications and sharp quantitative bounds in six related methodological settings: finite-order positivity of stable polynomials and Lorentz forms; identifiability and minimax recovery of latent mixing laws; integral presentations and graph critical groups; resistance optimization and random spanning trees; entropy inequalities and marginal compatibility; and polynomial or rational differential spectra. Its sixty source contributions are organized by subject and retained with their proofs, including auxiliary classifications, extremizers, algorithms, and counterexamples. The main themes are finite certificates, exact thresholds, matching upper and lower bounds, and the identification of information that an invariant can and cannot retain. The compilation integrates four corrections and distinguishes established external ingredients from the statements derived here. A final part records additional results obtained during integration, with their assumptions, proofs, and unresolved boundaries stated explicitly.

\medskip
\noindent\textbf{Keywords.} Real stability; alpha-permanent; Lorentz cone; binomial mixture; Wasserstein minimax risk; Smith normal form; critical group; effective resistance; entropy inequality; quasi-exact solvability.
\clearpage\markboth{Contents}{Contents}
\begingroup\small\setlength{\parskip}{0pt}\tableofcontents\endgroup
\clearpage\pagenumbering{arabic}
\section*{Integrated introduction and reading guide}
\phantomsection\addcontentsline{toc}{section}{Integrated introduction and reading guide}
\markboth{Integrated introduction}{Reading guide}
The six parts below study different mathematical objects rather than instances of one unproved universal principle. Their common working method is to expose a finite certificate or a sharp obstruction, and then prove both the positive and negative sides of the resulting classification. The original ten-stage development of each programme is preserved so that no auxiliary result or proof is lost. Stronger later results are explicitly linked to earlier partial ones. The final part contains further deductions obtained in preparing this integrated edition.

\subsection*{The six programmes}
\begin{description}[style=nextline,leftmargin=0pt,labelsep=0pt]
\item[Part I: stable polynomials and Lorentz positivity.]
Boundary reconstruction and corrupted principal minors lead to finite-order inverse-power questions. The central conclusions are the sharp universal third-order exponent for nonnegative-coefficient real-stable polynomials, all finite-order thresholds on the three-dimensional Lorentz cone, and quantitative rank-two permanent rigidity. Higher-dimensional fifth-order thresholds connect the analysis to spherical energy.
\item[Part II: statistical inverse problems.]
Finite and infinite stopping experiments are separated by identifiability, while recovery rates depend on how information is distributed across observation lengths. The programme develops minimax curves, information eigenvalues, a critical logarithmic lower bound, and an exact root-$n$ criterion on a bounded positive density class.
\item[Part III: integral graph invariants.]
Unimodular compression replaces spectral splitting when arithmetic primes matter. This yields complete blowup critical-group formulas, general coupled-presentation identities, and inverse classifications showing which tree information survives every uniform blowup.
\item[Part IV: electrical and spectral design.]
The results range from sharp network designs and uncertainty radii to topology-sensitive random-tree inflation and candidate-constrained greedy counterexamples. Claims about constrained additions are kept separate from the problem allowing every absent edge.
\item[Part V: entropy and compatibility.]
Local overlap obstructions are compared with global marginal feasibility. Fractional-partition inequalities receive exact coefficients and sharp witnesses; singular equality laws and quantitative rounding are treated with explicit conditioning assumptions.
\item[Part VI: differential spectra.]
Degree filtrations and resonance chains provide finite spectral certificates. Polynomial modules of specified forms are classified, with exact energies and nodal indices in selected families. The scope of polynomial Jordan chains is corrected before any extension to rational spaces is made.
\end{description}

\subsection*{Additional conclusions obtained in this edition}
Part VII proves six groups of further deductions. Their exact hypotheses and complete proofs appear in Sections~\ref{ext-showa-section}--\ref{ext-heisei-section}.
\begin{enumerate}[leftmargin=*]
\item \textbf{A second rank-two transition.} The regular projective polygon minimizes the alpha-permanent precisely for $\alpha\ge\mu_m$, with an explicit scalar equation for $\mu_m$, a nonregular equality family at the boundary, and $\tfrac12-\mu_m\sim m^2/2^{m+1}$. This minimization threshold differs from the positivity threshold (Theorems~\ref{ext-showa-minimum} and~\ref{ext-showa-asymptotic}).
\item \textbf{Two statistical resources.} With $n$ individuals and expected additional-observation budget $b$, the minimax Wasserstein risk is $\asymp\min\{1,\max(n^{-1/2},b^{-1/3})\}$. A complementary strict lower bound explains why a single fixed finite-mean cap law cannot attain the boundary rate (Section~\ref{ext-meiji-section}).
\item \textbf{What two cluster scales reveal.} An exact inverse classification, a single prime-power decoding rule, and $\exp(\Omega(\sqrt N))$ degree-distinct collision families determine what all two-scale critical groups can recover. Independently labeled cluster sizes recover the whole tree from group orders (Section~\ref{ext-keio-section}).
\item \textbf{Sparse greedy obstructions.} Candidate-constrained greedy augmentation can have vanishing relative gain and unbounded final-objective ratio even on unit-edge graphs of maximum degree five and treewidth at most two (Section~\ref{ext-taisho-section}).
\item \textbf{Exact recombination constants.} For any cut law $\nu$, the universal entropy-production constant in the Caputo--Sinclair normalization is $\delta_\nu=\min_{|S|\ge2}q_\nu(S)/(|S|-1)$, where $q_\nu(S)$ is the probability of splitting $S$. A strictly positive rare-event witness proves sharpness; exact continuous- and discrete-time contraction profiles follow. The optimization allows the product reference measure to vary (Theorems~\ref{ext-reiwa-production-exact} and~\ref{ext-reiwa-profiles}).
\item \textbf{The complete rational spectral space.} Positive growth at infinity and explicit local pole-raising conditions give a uniform denominator, a finite maximal invariant space, and all rational eigenvalues and Jordan chains. A counterexample shows why the local assumption matters (Theorem~\ref{ext-heisei-global}).
\end{enumerate}

\subsection*{Scope, provenance, and mathematical status}
All sixty source contributions occur below, identified by a programme letter and a two-digit round: S01--S10, M01--M10, K01--K10, T01--T10, R01--R10, and H01--H10. The global section and theorem numbering is new; explicitly tagged equations carry their source code. Their proofs are embedded in this single source file, not supplied by external inputs. A source concordance and an incorporated-corrections record close the manuscript.

The ordinary meaning of a theorem in this manuscript is a statement proved under its displayed hypotheses. It does not certify that the theorem is historically new. Real-stability theory, the complete-monotonicity framework, Bernstein--Durrmeyer diagonalization, Smith normal forms, matrix--forest identities, polymatroid minorants, and differential indicial methods have established antecedents, cited in the relevant contributions. Numerical checks are supplementary and are not substituted for proofs.

Lee and Kennedy's \emph{Minimax Rates for Learning Smooth Populations of Parameters} \cite{integration-lee-kennedy} provides a related heterogeneous-trial Legendre/Hahn information framework. Its stated target is pointwise smooth-density estimation, whereas the random-cap results here concern Wasserstein recovery. In particular the information diagonalization is not described as an invention of this manuscript. The Lorentz results address specified cones and derivative orders within the broader finite-order question of Scott and Sokal. The dimension-four fifth-order argument uses Tumanov's published five-point energy theorem. The candidate-restricted greedy construction does not settle the unrestricted all-nonedges formulation of Achterberg and Kooij.

\subsection*{How to read cumulative results}
For the strongest endpoints of the original programmes, start with S05, S08--S10; M06--M10; K04--K10; T08--T10; R05 and R09--R10; and H08--H10. Earlier statements and proofs are retained; integration notes identify subsequent resolutions of questions left open at an earlier stage.


\clearpage
\part{Stable Polynomials, Matrix Reconstruction, and Lorentz Positivity}


\clearpage
\section[S01 --- Sharp finite-moment reconstruction on zero-forcing graphs]{Sharp finite-moment reconstruction on zero-forcing graphs}\label{contrib-S01}
\markboth{Stable Polynomials, Matrix Reconstruction, and Lorentz Positivity}{S01 --- Showa}
\noindent\textsc{S01 \quad Showa programme, source round 1}\par\medskip



\begin{quote}\small\textbf{Overview.} For a real symmetric matrix with known graph and known signs of its nonzero off-diagonal entries, we give a finite, constructive recovery theorem from boundary moments. If a prescribed zero-forcing schedule colors the graph in $t$ parallel rounds, moments through degree $2t+1$ determine every entry, without simple-spectrum, genericity, or diagonal assumptions. The bound is sharp uniformly in $t$: a diagonal entry at the remote endpoint of a path is invisible through degree $2t$. We also describe exactly the gauge ambiguity when edge signs are not supplied. The underlying zero-forcing identifiability principle is established literature; the contribution offered for assessment is the explicit sharp horizon and its polynomial-vector certificate, not a claim to have solved a previously open problem.\end{quote}

\subsection{Problem and relation to prior work}\label{S01-sec-1}
Let $G=(V,E)$ be a finite simple graph, with vertices labeled $1,\ldots,n$. Write $\mathcal S(G)$ for the real symmetric matrices $A$ with
\[
 A_{uv}\ne0\quad\Longleftrightarrow\quad\{u,v\}\in E\qquad(u\ne v).
\]
Diagonal entries are arbitrary. For $B\subseteq V$, let $E_B$ be the matrix whose columns are the coordinate vectors $e_b$, in a specified order. The measured data are
\begin{equation}\label{S01-eq:moments}
 M_k=E_B^{T}A^kE_B,\qquad 0\le k\le h.
\end{equation}
These are exact algebraic data; no claim is made that differentiating experimental trajectories to high order is numerically benign.

A forcing round starts with a colored set $S$. A vertex $u\in S$ with exactly one neighbor $v\notin S$ may force $v$. Simultaneous forces are permitted; if more than one source forces the same target, retain one. A schedule is a sequence
\[
 B=S_0\subset S_1\subset\cdots\subset S_t=V
\]
in which every newly colored vertex is the target of a permitted force from $S_{s-1}$. Such a $B$ is a zero-forcing set.

Hamiltonian identification through an infecting gateway was developed by Burgarth and Maruyama \cite{ref001}. Zero forcing also yields identifiability for combinatorially symmetric hidden Markov models \cite{ref002}; it is therefore essential not to advertise the bare uniqueness implication as new. The argument below isolates a finite measurement horizon, makes its dependence on a specified parallel schedule explicit, and supplies a matching worst-case obstruction. No exhaustive novelty determination has been made.

\subsection{A sharp finite-horizon theorem}\label{S01-sec-2}
\begin{theorem}\label{S01-thm:main}
Fix $G,B$, a $t$-round zero-forcing schedule, and the sign of each nonzero off-diagonal entry. Every $A\in\mathcal S(G)$ with those signs is uniquely determined by $M_0,\ldots,M_{2t+1}$. Recovery uses only arithmetic and positive square roots. No eigenvalue separation or simplicity is assumed.

For every integer $t\ge0$, some graph with propagation time $t$ has two matrices with the same prescribed edge signs and identical moments through degree $2t$, but different degree-$2t+1$ moments. Thus $2t+1$ cannot be reduced in a theorem depending only on $t$.
\end{theorem}

To prove the theorem we work with formal vector polynomials
\[
 p(x)=\sum_{j=0}^d x^j c_j,\qquad c_j\in\R^{|B|},
\]
and their realizations $p[A]=\sum_j A^jE_Bc_j\in\R^n$. The moments determine two pairings whenever the indicated indices are available:
\begin{align}
 \langle p,q\rangle_M&=\sum_{i,j}c_i^TM_{i+j}d_j=p[A]^Tq[A],\label{S01-eq:gram}\\
 \langle p,xq\rangle_M&=\sum_{i,j}c_i^TM_{i+j+1}d_j=p[A]^TAq[A].\label{S01-eq:shift}
\end{align}
These identities follow immediately by symmetry of $A$. Formal polynomials with zero realization need not be removed or identified; their pairings with all relevant vectors vanish automatically.

\begin{lemma}[One forcing step]\label{S01-lem:step}
Suppose for every $w\in S_s$ a known polynomial $p_w$ of degree at most $s$ satisfies $p_w[A]=e_w$. If $u\in S_s$ forces $v\notin S_s$, then moments through $2s+2$ construct a polynomial $p_v$ of degree at most $s+1$ with $p_v[A]=e_v$.
\end{lemma}
\begin{proof}
For $w\in S_s$, compute
\[
 a_{wu}=\langle p_w,xp_u\rangle_M.
\]
This needs moments only through $2s+1$. Because $v$ is the unique neighbor of $u$ outside $S_s$, the polynomial
\[
 q_v=xp_u-\sum_{w\in S_s}a_{wu}p_w
\]
satisfies $q_v[A]=a_{vu}e_v$. Its degree is at most $s+1$, so its squared norm $\langle q_v,q_v\rangle_M=a_{vu}^2$ uses moments only through $2s+2$. The norm is strictly positive. Using the prescribed sign $\sigma_{vu}$, set
\[
 a_{vu}=\sigma_{vu}\sqrt{\langle q_v,q_v\rangle_M},\qquad
 p_v=q_v/a_{vu}.
\]
These expressions have the required realization and degree. All old polynomials are retained unchanged.
\end{proof}

\begin{proof}[Proof of Theorem \ref{S01-thm:main}]
For $b\in B$, take $p_b$ to be its constant coordinate vector in $\R^{|B|}$. Apply Lemma \ref{S01-lem:step} to every force in a round, using the set at the start of that round. After $t$ rounds, every $e_v$ has a polynomial representation of degree at most $t$. The reconstruction is then
\[
 A_{uv}=\langle p_u,xp_v\rangle_M,
\]
which requires moments through $2t+1$. The formulas are deterministic after a source is selected for each target and give the same result for any matrix consistent with the data.

For sharpness, take the path $0,1,\ldots,t$ and $B=\{0\}$, with arbitrary fixed positive edge weights $b_1,\ldots,b_t$. Fix every diagonal entry except $A_{tt}=d$. A term in $(A^k)_{00}$ is a weighted closed walk of length $k$ starting at $0$, allowing diagonal stays. To use $d$, it must travel to $t$, stay there, and return, which requires at least $2t+1$ steps. Consequently the moments for $k\le2t$ do not depend on $d$. At degree $2t+1$ there is exactly one walk involving $d$, giving the term
\[
 d\prod_{j=1}^{t}b_j^2.
\]
All remaining terms are independent of $d$. For $t=0$, the one-vertex matrix itself gives the same conclusion.
\end{proof}

\subsection{What happens without sign information?}\label{S01-sec-3}
The boundary moments cannot identify arbitrary edge signs absolutely. Let $D$ be diagonal with entries $\pm1$ and $D_{bb}=1$ for every $b\in B$. Then
\[
 E_B^T(DAD)^kE_B=E_B^TA^kE_B.
\]
The following shows that under zero forcing this is the only ambiguity.

\begin{theorem}[Exact boundary-fixed gauge fiber]\label{S01-thm:gauge}
Suppose $B$ has a $t$-round forcing schedule. For $A,\widetilde A\in\mathcal S(G)$, the following are equivalent:
\begin{enumerate}[label=(\roman*)]
 \item $E_B^TA^kE_B=E_B^T\widetilde A^kE_B$ for $0\le k\le2t+1$;
 \item $\widetilde A=DAD$ for a diagonal sign matrix $D$ satisfying $D_{bb}=1$ on $B$;
 \item the boundary resolvents agree: $E_B^T(zI-A)^{-1}E_B=E_B^T(zI-\widetilde A)^{-1}E_B$ wherever defined.
\end{enumerate}
\end{theorem}
\begin{proof}
Choose one forcing edge into each vertex outside $B$. These edges form a rooted forest: every edge points from an earlier to a later round, each nonroot has one parent, and following parents reaches $B$. There is a unique diagonal sign matrix fixing $B$ that makes all selected forcing-edge entries of $DAD$ positive. The one-step construction in Lemma \ref{S01-lem:step} needs signs only on selected forcing edges, not on all edges. Therefore it reconstructs this canonical representative uniquely from the finite moments. Equal moments imply equal canonical representatives, hence (i)$\Rightarrow$(ii). Conjugation proves (ii)$\Rightarrow$(iii). Expanding the resolvent at infinity proves (iii)$\Rightarrow$(i).
\end{proof}

\begin{remark}
If $G$ has several connected components, zero forcing requires at least one boundary vertex in every component. The theorem remains valid with a rooted forest, not just a connected graph. Although the canonical gauge makes forcing edges positive, other edges may be negative; their signs are recovered by the final pairing, so no cycle-sign information is discarded.
\end{remark}

\subsection{An exact certificate and a worked path}\label{S01-sec-4}
A useful output is not merely a reconstructed matrix but the list of polynomial vectors $p_v$. A checker can verify the moment identities
\[
 \langle p_u,p_v\rangle_M=\delta_{uv},\qquad
 \langle p_u,xp_v\rangle_M=A_{uv}
\]
and the recorded forcing equations. These are finite identities in the input moments. This is a transparent algebraic certificate, not formal proof-assistant verification.

For a three-vertex path with $B=\{1\}$ and positive edges, write $m_k=(A^k)_{11}$. The first two recovery stages are
\[
 a=m_1,\qquad b=\sqrt{m_2-m_1^2},\qquad
 p_1=1,\quad p_2=(x-a)/b.
\]
Then
\[
 d=\langle p_2,xp_2\rangle_M,
 \quad q_3=xp_2-bp_1-dp_2,
 \quad c=\sqrt{\langle q_3,q_3\rangle_M},
 \quad p_3=q_3/c.
\]
Finally $f=\langle p_3,xp_3\rangle_M$ gives
\[
 A=\begin{pmatrix}a&b&0\\b&d&c\\0&c&f\end{pmatrix}.
\]
Here $d$ uses $m_3$, $c$ uses $m_4$, and $f$ uses $m_5$. This displays precisely the horizon $2t+1=5$, including the otherwise easily missed last diagonal measurement.

\subsection{Scope and unresolved questions}\label{S01-sec-5}
The theorem is unconditional under the stated exact-data assumptions. It does not say that zero forcing is necessary for identification on a particular matrix family. Nor does it claim that $2t+1$ is minimal for every graph or every boundary set. The graph-dependent minimal horizon may be substantially smaller. Uniform numerical conditioning also cannot follow from the theorem: as a forcing-edge weight approaches zero, divisions in the reconstruction become ill-conditioned. Quantitative stability and optimal graph-dependent horizons are different research questions.





\clearpage
\section[S02 --- Optimal adversarial correction of triangle-minor tomography]{Optimal adversarial correction of triangle-minor tomography}\label{contrib-S02}
\markboth{Stable Polynomials, Matrix Reconstruction, and Lorentz Positivity}{S02 --- Showa}
\noindent\textsc{S02 \quad Showa programme, source round 2}\par\medskip



\begin{quote}\small\textbf{Overview.} We study an adversarial version of symmetric principal-minor reconstruction. For an $n\times n$ real symmetric matrix with all off-diagonal entries nonzero, suppose its order-one and order-two principal minors are exact, while the order-three minors contain at most $f$ arbitrary corruptions and $e$ announced erasures. We prove that $2f+e<n-2$ is sufficient for exact recovery up to diagonal sign conjugation, by an explicit $O(n^3)$ majority decoder, and is necessary for a uniform guarantee. The impossibility examples can be chosen with all eigenvalues strictly between zero and one. The result combines classical switching invariants with a disjoint-witness decoding argument; no claim that the underlying principal-minor identifiability or two-graph algebra is new is made.\end{quote}

\subsection{Observation model and exact statement}\label{S02-sec-1}
Let $n\ge3$ and $A=A^T\in\R^{n\times n}$ satisfy $a_{ij}\ne0$ for $i\ne j$. For $S\subseteq[n]$, let $\Delta_S=\det A[S,S]$. We observe all $\Delta_{\{i\}}$ and $\Delta_{\{i,j\}}$ exactly. Among the $\binom n3$ order-three observations, a set of at most $e$ is declared unavailable and at most $f$ additional entries can be changed arbitrarily. The corrupted locations are unknown. All budgets are total counts, not fractions per vertex.

Two matrices are switching equivalent if $\widetilde A=DAD$ for a diagonal matrix $D$ with entries in $\{\pm1\}$. This ambiguity is unavoidable because every principal minor is preserved.

\begin{theorem}[Sharp mixed-error threshold]\label{S02-thm:main}
Under the observation model above, if
\begin{equation}\label{S02-eq:threshold}
 2f+e<n-2,
\end{equation}
then $A$ can be recovered up to switching in $O(n^3)$ arithmetic operations and sign comparisons. A canonical output has all entries in its first off-diagonal row positive.

Conversely, if $2f+e\ge n-2$, there are two nonswitching-equivalent matrices and a common admissible observation record, each within the stated budgets. This remains true when both matrices have spectrum contained in $(0,1)$.
\end{theorem}

The exact-minor problem and its relation to short cycles are established \cite{ref003}. Signed-graph switching has long been connected with binary codes \cite{ref004}. Our focus differs: measurements rather than edges are corrupted, and the theorem gives an explicit optimal mixed-error radius for the complete support graph. This is a finite theorem with a complete proof; novelty of this precise formulation has not been exhaustively verified.

\subsection{Reducing determinants to triangle signs}\label{S02-sec-2}
The trusted data determine
\[
 d_i=a_{ii}=\Delta_{\{i\}},\qquad
 w_{ij}=|a_{ij}|=\sqrt{d_id_j-\Delta_{\{i,j\}}}>0.
\]
Write $s_{ij}=\operatorname{sgn}(a_{ij})=s_{ji}$. For a triple $T=\{i,j,k\}$ define
\[
 \tau_{ijk}=s_{ij}s_{ik}s_{jk}\in\{\pm1\},
 \qquad
 b_{ijk}=d_id_jd_k-d_iw_{jk}^2-d_jw_{ik}^2-d_kw_{ij}^2.
\]
The determinant expansion gives
\begin{equation}\label{S02-eq:triangle}
 \Delta_{\{i,j,k\}}=b_{ijk}+2w_{ij}w_{ik}w_{jk}\tau_{ijk}.
\end{equation}
Convert a nonerased reported determinant $y_{ijk}$ to
\[
 \widehat\tau_{ijk}=\operatorname{sgn}(y_{ijk}-b_{ijk}),
\]
using either sign when the difference is zero. Since the weights are positive, every uncorrupted measurement produces the correct sign. Thus there are at most $f$ wrong signs and $e$ erased signs. A malicious numerical value need not equal either feasible determinant; it still causes at most one sign error.

\begin{lemma}[Tetrahedron identity]\label{S02-lem:tetra}
For any four distinct vertices $i,j,k,\ell$,
\[
 \tau_{ijk}\tau_{ij\ell}\tau_{ik\ell}\tau_{jk\ell}=1.
\]
Conversely, any assignment of triangle signs satisfying these identities comes from edge signs, uniquely up to switching.
\end{lemma}
\begin{proof}
In the displayed product each of the six edge signs appears exactly twice. For the converse, set $s_{1i}=1$ for $i>1$, and set $s_{ij}=\tau_{1ij}$ for $i,j>1$. Triangles through vertex $1$ are correct by construction. For $i,j,k>1$, the tetrahedron identity on $\{1,i,j,k\}$ gives
\[
 s_{ij}s_{ik}s_{jk}=\tau_{1ij}\tau_{1ik}\tau_{1jk}=\tau_{ijk}.
\]
Any edge signing can be switched uniquely to make every $s_{1i}$ positive, by taking the switching sign at vertex $i$ to be its old $s_{1i}$ and fixing the sign at vertex $1$. Hence the canonical representative is unique.
\end{proof}

\subsection{The decoder and its proof}\label{S02-sec-3}
Fix the reference vertex $1$. For each $2\le i<j\le n$, construct the following $n-2$ candidate votes for $\tau_{1ij}$:
\begin{enumerate}[label=(\alph*)]
 \item the direct value $\widehat\tau_{1ij}$;
 \item for every $k\notin\{1,i,j\}$, the product
 \[
  \widehat\tau_{1ik}\widehat\tau_{1jk}\widehat\tau_{ijk}.
 \]
\end{enumerate}
Discard a vote if any sign used in it is erased. Take the majority of the remaining votes. Define the output by $\widehat a_{ii}=d_i$, $\widehat a_{1i}=w_{1i}$, and
\[
 \widehat a_{ij}=w_{ij}\,\operatorname{Majority}(\text{votes for }\tau_{1ij})
 \quad(i,j>1).
\]
Symmetrize by $\widehat a_{ji}=\widehat a_{ij}$.

\begin{lemma}[Disjoint witness sets]\label{S02-lem:disjoint}
For fixed $(i,j)$, the triangle coordinates used by distinct votes are disjoint. Therefore at most $e$ votes are discarded and at most $f$ retained votes are wrong.
\end{lemma}
\begin{proof}
The direct coordinate $\{1,i,j\}$ occurs in no product vote. For a product vote indexed by $k$, its coordinates are $\{1,i,k\}$, $\{1,j,k\}$, and $\{i,j,k\}$. None equals one of the three coordinates belonging to a different index $\ell$: any possible equality would force $k=\ell$ or put an excluded vertex in the index. Hence a single erased or corrupted triangle influences at most one vote for this fixed pair. An unaltered, unerased product vote is correct by Lemma \ref{S02-lem:tetra}. Products containing two errors can happen to be correct; that only improves the bound.
\end{proof}

\begin{proof}[Sufficiency in Theorem \ref{S02-thm:main}]
At least $n-2-e$ votes survive. At most $f$ are wrong. By \eqref{S02-eq:threshold}, $n-2-e>2f$, so the correct sign is a strict majority. Every canonical edge sign is therefore recovered. Lemma \ref{S02-lem:tetra} proves switching equivalence. There are $O(n^2)$ edge signs, each using $O(n)$ votes; all preprocessing also takes at most $O(n^3)$ operations.
\end{proof}

\subsection{Exact distance and impossibility}\label{S02-sec-4}
Let $\mathcal T_n$ be the set of all triangle-sign arrays induced by edge signings of $K_n$. Regard $-1$ as the binary symbol $1$. Then $\mathcal T_n$ is a binary linear code: multiplication of sign arrays corresponds to addition of edge-sign vectors modulo two. Its dimension is
\[
 \binom n2-(n-1)=\binom{n-1}{2},
\]
since vertex switching has dimension $n-1$ and the canonical signs on pairs outside vertex $1$ are free.

\begin{proposition}[Exact minimum distance]\label{S02-prop:distance}
Distinct words of $\mathcal T_n$ differ on at least $n-2$ triangles, and this lower bound is attained.
\end{proposition}
\begin{proof}
Divide the two sign arrays coordinatewise and choose a triangle $\{i,j,k\}$ with sign $-1$. For each $\ell\notin\{i,j,k\}$, the tetrahedron identity forces at least one of
\[
 \{i,j,\ell\},\quad\{i,k,\ell\},\quad\{j,k,\ell\}
\]
to have sign $-1$. Those triples are distinct for different $\ell$, and none is the initial triangle. This gives $1+(n-3)=n-2$ discrepancies. Flipping a single edge changes precisely the $n-2$ triangles containing that edge, attaining equality.
\end{proof}

\begin{proof}[Necessity in Theorem \ref{S02-thm:main}]
Choose two edge signings differing on exactly one edge. Their triangle arrays differ on a set $D$ of size $n-2$. Erase $e'=\min(e,n-2)$ coordinates of $D$. The remaining $n-2-e'$ coordinates can be partitioned into two sets of size at most $f$ exactly when $n-2-e'\le2f$, which follows from the hypothesis. On one part report the first array's value, and on the other report the second array's value. Outside $D$, report their common value. The record is within $f$ errors and $e$ erasures of each array.

To realize this ambiguity with well-behaved matrices, let all diagonal entries equal $1/2$ and all off-diagonal magnitudes equal $\varepsilon$, where $0<\varepsilon<1/(2(n-1))$. Use the two signings for the off-diagonal signs. The sum of absolute off-diagonal entries in every row is $(n-1)\varepsilon<1/2$. For a real unit vector $x$, the inequality $2|x_ix_j|\le x_i^2+x_j^2$ gives
\[
 \left|x^T(A-\tfrac12I)x\right|\le(n-1)\varepsilon.
\]
Thus both spectra lie in $(0,1)$. Their order-one and order-two principal minors agree, while their triangle determinants differ precisely on $D$ by \eqref{S02-eq:triangle}. The same mixed observation construction applies to the real determinant values. The differing triangle invariant proves that the matrices are not switching equivalent.
\end{proof}

\subsection{Small deterministic noise and limitations}\label{S02-sec-5}
The exact trusted-data model admits a useful noise extension without a new probabilistic argument. Suppose the diagonals and magnitudes remain exact, and every nonexceptional triangle determinant obeys
\[
 |y_{ijk}-\Delta_{\{i,j,k\}}|<2w_{ij}w_{ik}w_{jk}.
\]
Then \eqref{S02-eq:triangle} shows that its thresholded sign is correct. Consequently the same decoder succeeds under \eqref{S02-eq:threshold} even if every nonexceptional measurement is noisy. The strict inequality matters: at equality the observation can reach the midpoint $b_{ijk}$.

For example, with $n=7$, at most two arbitrary corrupted triangle measurements can be corrected in the absence of erasures because $4<5$. With one arbitrary corruption and two erasures, $2+2<5$ also guarantees recovery. Three corruptions cannot be uniformly corrected, even for kernels bounded away from both $0$ and $I$.

This theorem assumes a complete nonzero support graph and exact low-order magnitudes. Missing edges remove triangle witnesses. Errors in low-order minors affect many thresholds simultaneously and cannot be charged as isolated triangle errors. Extending the argument to sparse support graphs requires a separate measurement design and cannot be inferred from the present result.





\clearpage
\section[S03 --- A third-order logarithmic test for partition matroids and cactus graphs]{A third-order logarithmic test for partition matroids and cactus graphs}\label{contrib-S03}
\markboth{Stable Polynomials, Matrix Reconstruction, and Lorentz Positivity}{S03 --- Showa}
\noindent\textsc{S03 \quad Showa programme, source round 3}\par\medskip



\begin{quote}\small\textbf{Overview.} We give an elementary third-order test for a very restrictive positivity property of matroid basis polynomials. For a finite matroid $M$, every mixed third derivative in three distinct variables of $\log B_M$ is nonnegative throughout the positive orthant if and only if every circuit has size at most two. Equivalently, after loops are removed, $M$ is a direct sum of rank-one matroids, and $B_M$ is a product of linear forms in disjoint variable sets. Thus third-order logarithmic positivity already forces alternating signs at every derivative order and complete monotonicity of $B_M^{-\beta}$ for every $\beta>0$. For cotree polynomials this characterizes cactus graphs. The inverse-power classification has established antecedents in Scott--Sokal; the finite third-order obstruction and its elementary certificate are the focus here, not a claimed solution of their broader open classification problems.\end{quote}

\subsection{Definitions and the structural theorem}\label{S03-sec-1}
A finite matroid $M$ on ground set $E$ consists of a family of independent subsets satisfying heredity and augmentation. Its maximal independent sets are its bases, and its minimal dependent sets are its circuits. A loop is an element in no basis; a coloop is an element in every basis. The basis generating polynomial is
\[
 B_M(x)=\sum_{B\text{ a basis of }M}\prod_{e\in B}x_e.
\]
It is positive when all $x_e>0$, including for the rank-zero matroid, whose polynomial is $1$.

For a smooth positive function $P$, say that its logarithm has alternating positive derivatives if, for every $k\ge1$ and every sequence of variables $i_1,\ldots,i_k$ (repetition allowed),
\begin{equation}\label{S03-eq:alternate}
 (-1)^{k-1}\partial_{i_1}\cdots\partial_{i_k}\log P\ge0.
\end{equation}
This definition imposes no sign on $\log P$ itself. A smooth function $F$ is completely monotone if $(-1)^k\partial_{i_1}\cdots\partial_{i_k}F\ge0$ for all $k\ge0$.

\begin{theorem}\label{S03-thm:matroid}
For a finite matroid $M$, the following conditions are equivalent.
\begin{enumerate}[label=(\roman*)]
 \item For every three distinct elements $e,f,g$, $\partial_e\partial_f\partial_g\log B_M(x)\ge0$ for every $x\in(0,\infty)^E$.
 \item Every circuit of $M$ has size one or two.
 \item The nonloops partition into nonempty classes $P_1,\ldots,P_r$ such that the bases contain exactly one element of each class. Consequently
 \[
 B_M(x)=\prod_{j=1}^r\left(\sum_{e\in P_j}x_e\right).
 \]
 \item $\log B_M$ satisfies \eqref{S03-eq:alternate} at every order.
 \item $B_M^{-\beta}$ is completely monotone for every real $\beta>0$.
\end{enumerate}
If (i) fails, it fails strictly at a point with positive rational coordinates.
\end{theorem}

Scott and Sokal \cite{ref005} study complete monotonicity of inverse basis-polynomial powers and prove substantially more detailed exponent-dependent classifications. In particular their small-exponent results imply the structural class in (v). Our proof is self-contained and emphasizes that the structural class can already be detected by one order of logarithmic derivatives. Complete log-concavity is a different, much broader property; it must not be confused with \eqref{S03-eq:alternate}.

\subsection{Minor closure of the third-order condition}\label{S03-sec-2}
For completeness, deletion $M\setminus e$ restricts the independent sets to $E\setminus\{e\}$. For a nonloop $e$, the independent sets of the contraction $M/e$ are those $I\subseteq E\setminus\{e\}$ for which $I\cup\{e\}$ is independent in $M$. Contracting a loop agrees with deleting it. The basis polynomial has the usual elementary decomposition
\begin{equation}\label{S03-eq:delcon}
 B_M=B_{M\setminus e}+x_eB_{M/e}
\end{equation}
when $e$ is neither a loop nor a coloop: bases either avoid $e$ or contain it. For a loop $B_M=B_{M\setminus e}$; for a coloop $B_M=x_eB_{M/e}$ and $M/e=M\setminus e$ as matroids on the remaining elements.

\begin{lemma}\label{S03-lem:minor}
Condition (i) of Theorem \ref{S03-thm:matroid} is inherited by every matroid minor.
\end{lemma}
\begin{proof}
Let the three differentiated variables be different from $e$. When $e$ is neither a loop nor a coloop, \eqref{S03-eq:delcon} gives
\[
 B_M(x',t)\longrightarrow B_{M\setminus e}(x')\quad(t\downarrow0)
\]
and
\[
 t^{-1}B_M(x',t)\longrightarrow B_{M/e}(x')\quad(t\to\infty).
\]
On every compact subset of the positive orthant of $x'$, these convergences hold with every derivative in $x'$. Both limiting polynomials are strictly positive there, so logarithms and their derivatives also converge. The term $\log t$ in the second limit has zero derivatives in $x'$. Nonnegative third derivatives therefore remain nonnegative in both limits. Loop and coloop cases follow from the displayed factorizations. Iterating proves the claim.
\end{proof}

\begin{lemma}[The three-element obstruction]\label{S03-lem:obstruction}
Let $U_{2,3}$ be the matroid whose bases are all two-element subsets of $\{1,2,3\}$. For $Q=x_1x_2+x_1x_3+x_2x_3$,
\[
 \partial_1\partial_2\partial_3\log Q\big|_{(1,1,1)}=-\frac{2}{27}<0.
\]
\end{lemma}
\begin{proof}
For distinct $i,j,k$,
\[
 \partial_{ijk}\log Q
 =\frac{Q_{ijk}}{Q}
 -\frac{Q_{ij}Q_k+Q_{ik}Q_j+Q_{jk}Q_i}{Q^2}
 +\frac{2Q_iQ_jQ_k}{Q^3}.
\]
At $(1,1,1)$ the quantities are $Q=3$, $Q_i=2$, $Q_{ij}=1$, and $Q_{123}=0$. Hence the value is $-6/9+16/27=-2/27$.
\end{proof}

\subsection{Proof of the classification}\label{S03-sec-3}
\begin{proof}[Proof of Theorem \ref{S03-thm:matroid}]
Suppose (i) holds. If $C$ is a circuit of size $s\ge3$, delete all elements outside $C$. The restricted matroid has every proper subset of $C$ independent, hence is $U_{s-1,s}$. Contract any $s-3$ elements of $C$. By the definition of contraction, the resulting three-element matroid has every two-element set independent and its three-element set dependent; it is $U_{2,3}$. This contradicts Lemmas \ref{S03-lem:minor} and \ref{S03-lem:obstruction}. Thus (i) implies (ii).

Assume (ii), remove loops, and declare two distinct elements equivalent if they form a two-element circuit. This is an equivalence relation after adjoining equality. To see transitivity, if $\{a,b\}$ and $\{b,c\}$ are circuits, circuit elimination gives a circuit contained in $\{a,c\}$; there are no loops, so this circuit is $\{a,c\}$. Circuit elimination itself follows from the standard independent-set matroid axioms and is used here only in this two-element case: if $\{a,c\}$ were independent, augmentation between it and $\{b\}$ would make either $\{a,b\}$ or $\{b,c\}$ independent, a contradiction.

An independent set contains at most one member of each equivalence class. Conversely, a set with at most one member from each class contains no circuit, and is therefore independent. Maximality says that a basis contains exactly one member of each class, proving (iii).

Under (iii), $\log B_M=\sum_j\log L_j$ with $L_j=\sum_{e\in P_j}x_e$. A derivative involving variables from different classes is zero. If all $k$ variables lie in one class, the derivative is
\[
 (-1)^{k-1}(k-1)!L_j^{-k}.
\]
Variables belonging to loops give zero. Thus (iv) holds, and (iv) immediately implies (i).

Still under (iii), $B_M^{-\beta}=\prod_jL_j^{-\beta}$. For a total of $k$ derivatives in variables from a single class,
\[
 (-1)^k\partial_{i_1}\cdots\partial_{i_k}L_j^{-\beta}
 =\beta(\beta+1)\cdots(\beta+k-1)L_j^{-\beta-k}\ge0.
\]
Disjointness of the variable classes proves complete monotonicity of the product, hence (v).

Finally assume (v). Write $h=\log B_M$ and fix three distinct variables. Direct differentiation gives
\[
 -\partial_{efg}e^{-\beta h}
 =e^{-\beta h}\left[
 \beta h_{efg}-\beta^2(h_eh_{fg}+h_fh_{eg}+h_gh_{ef})
 +\beta^3h_eh_fh_g\right].
\]
The left side is nonnegative. Divide by $\beta e^{-\beta h}$ and let $\beta\downarrow0$ to obtain $h_{efg}\ge0$, giving (i).

If (i) fails, the preceding implication shows there is a point and a triple at which the third derivative is strictly negative. Since the derivative is a continuous rational function with positive denominator, the strict inequality persists in an open neighborhood. Positive rational points are dense, proving the last assertion.
\end{proof}

\subsection{A graph-theoretic formulation}\label{S03-sec-4}
For a connected finite loopless multigraph $G=(V,E)$, let
\[
 \Psi_G(r)=\sum_{T\text{ spanning tree}}\prod_{e\notin T}r_e
\]
be its cotree polynomial. A cactus is a connected graph in which every edge belongs to at most one simple cycle; a pair of parallel edges counts as a cycle of length two.

\begin{corollary}[Cactus test]\label{S03-cor:cactus}
For a connected loopless multigraph $G$, the following are equivalent:
\begin{enumerate}[label=(\alph*)]
 \item $G$ is a cactus;
 \item $\partial_e\partial_f\partial_g\log\Psi_G\ge0$ throughout the positive orthant for every three distinct edges;
 \item $\log\Psi_G$ has alternating positive derivatives at every order;
 \item $\Psi_G^{-\beta}$ is completely monotone for every $\beta>0$.
\end{enumerate}
\end{corollary}
\begin{proof}
For a cactus, every spanning tree includes all bridges and omits exactly one edge from each cycle. The choices are independent, so
\[
 \Psi_G(r)=\prod_{C\text{ a cycle}}\left(\sum_{e\in C}r_e\right).
\]
The derivative calculations in the main theorem prove (b)--(d).

If $G$ is not a cactus, it contains a theta subgraph: two distinct vertices joined by three internally disjoint paths. Indeed, take two distinct cycles with a common edge; a maximal segment of one outside the other, together with the two arcs of the other between its endpoints, gives such a theta. Extend this theta to a connected spanning subgraph by adding tree branches. Delete all other edges, contract the branches, and contract all but one edge in each of the three paths. The result is the graph with two vertices and three parallel edges.

The cotree polynomial of the last graph is $r_1r_2+r_1r_3+r_2r_3$. To pass to the minor in cotree variables, contraction corresponds to setting $r_e\downarrow0$; deletion of a nonbridge corresponds to taking $r_e\to\infty$ and dividing by $r_e$. These facts follow by splitting spanning trees according to whether they contain $e$. They preserve the third-derivative nonnegativity in retained variables by the same proof as Lemma \ref{S03-lem:minor}. The strict negative value $-2/27$ contradicts (b). Finally (c) implies (b), and (d) implies (b) by the small-$\beta$ calculation already proved.
\end{proof}

\subsection{Certificates and what is not proved}\label{S03-sec-5}
The negative certificate can be constructed from any circuit of size at least three, or from a theta in the graph case. Record a deletion/contraction sequence to the obstruction. Lift the target point $(1,1,1)$ by assigning successively sufficiently small or large positive rational values to eliminated variables. Convergence guarantees that the retained derivative stays negative after all lifts. A finite exact checker evaluates
\[
 P^2P_{efg}-P(P_{ef}P_g+P_{eg}P_f+P_{fg}P_e)+2P_eP_fP_g<0,
\]
which is the numerator of $P^3\partial_{efg}\log P$.
No polynomial bit-complexity bound is asserted for this naive lifting method.

The theorem concerns positivity at \emph{all} positive weight vectors. Testing only the equal-weight vector is not equivalent. Nor does failure here contradict pairwise negative correlation, ordinary log-concavity, or complete monotonicity for selected larger exponents. The exponent-dependent problem studied by Scott--Sokal is not solved by a third-order test.





\clearpage
\section[S04 --- The exact third-order inverse-power threshold for rank-two basis polynomials]{The exact third-order inverse-power threshold for rank-two basis polynomials}\label{contrib-S04}
\markboth{Stable Polynomials, Matrix Reconstruction, and Lorentz Positivity}{S04 --- Showa}
\noindent\textsc{S04 \quad Showa programme, source round 4}\par\medskip
\noindent\textit{Integration note.} The order-three conclusion is extended to arbitrary real-stable polynomials in Contribution~\ref{contrib-S05}.\par\medskip



\begin{quote}\small\textbf{Overview.} Let $E_{2,n}(x)=\sum_{i<j}x_ix_j$ on the positive orthant. We determine exactly when $E_{2,n}^{-\beta}$ has the alternating derivative signs of a completely monotone function through total order three, with repeated variables included. For every $n\ge3$, the answer is $\beta=0$ or $\beta\ge1/4$; the threshold does not depend on $n$. The proof rests on the sharp inequality $4E_{2,n}(E_i+E_j+E_k)\le9E_iE_jE_k$, with an explicit nonnegative remainder. We extend the result to every rank-two matroid, record a deterministic sharpness sequence, and exhibit a fifth derivative that fails at the boundary exponent. This is a finite-order classification, not a claim of complete monotonicity or a resolution of the full inverse-power problem.\end{quote}

\subsection{The finite-order question}\label{S04-sec-1}
For $s\ge0$, call a smooth function $F$ \emph{alternating through order $s$} if
\[
 (-1)^{|\alpha|}\partial^\alpha F(x)\ge0
 \qquad\text{for every multi-index $\alpha$ with }|\alpha|\le s.
\]
Repeated coordinate derivatives are part of this definition. Complete monotonicity requires this for every $s$.

Set $Q=E_{2,n}=\sum_{i<j}x_ix_j$ and $Q_i=\partial_iQ$. The full complete-monotonicity problem for inverse elementary symmetric polynomials is part of the work of Scott and Sokal \cite{ref005}. In particular $(xy+xz+yz)^{-\beta}$ is completely monotone precisely for $\beta=0$ or $\beta\ge1/2$. The finite order-three property below has a strictly smaller critical exponent. The distinction prevents an elementary low-order calculation from being mistaken for a solution of an all-orders problem.

\begin{theorem}\label{S04-thm:main}
Let $n\ge3$ and $\beta\in\R$. The function $E_{2,n}^{-\beta}$ is alternating through order three on $(0,\infty)^n$ if and only if
\[
 \beta=0\quad\text{or}\quad\beta\ge\frac14.
\]
For $n=2$, the corresponding condition is simply $\beta\ge0$.
\end{theorem}

This paper supplies a complete algebraic proof and an explicit sharp inequality. The precise finite-order result is offered as a focused refinement; an exhaustive claim of first publication is not made.

\subsection{A sharp three-gradient inequality}\label{S04-sec-2}
\begin{lemma}\label{S04-lem:pair}
For distinct $i,j$, $Q_iQ_j\ge Q$.
\end{lemma}
\begin{proof}
Write $x=x_i$, $y=x_j$, $S=\sum_{\ell\ne i,j}x_\ell$, and $R=\sum_{a<b,\ a,b\ne i,j}x_ax_b$. Then
\[
 Q=xy+S(x+y)+R,\qquad Q_iQ_j=(y+S)(x+S).
\]
Their difference is $S^2-R=\sum_{\ell\ne i,j}x_\ell^2+R\ge0$.
\end{proof}

\begin{lemma}[Sharp inequality]\label{S04-lem:triple}
For three distinct indices $i,j,k$,
\begin{equation}\label{S04-eq:ineq}
 4Q(Q_i+Q_j+Q_k)\le9Q_iQ_jQ_k.
\end{equation}
The constant $9/4$ in $Q(Q_i+Q_j+Q_k)/(Q_iQ_jQ_k)\le9/4$ is sharp for every $n\ge3$.
\end{lemma}
\begin{proof}
Rename the selected variables $x,y,z$ and put
\[
 U=x+y+z,\quad V=xy+xz+yz,\quad W=xyz,
\]
\[
 S=\sum_{\ell\notin\{i,j,k\}}x_\ell,\qquad
 R=\sum_{a<b,\ a,b\notin\{i,j,k\}}x_ax_b,
 \qquad D=\frac{S^2}{2}-R=\frac12\sum_{\ell\notin\{i,j,k\}}x_\ell^2.
\]
Then
\[
 Q=V+SU+\frac{S^2}{2}-D,
 \qquad Q_i+Q_j+Q_k=2U+3S,
\]
and
\[
 Q_iQ_jQ_k=(y+z+S)(x+z+S)(x+y+S)
 =S(U+S)^2+V(U+S)-W.
\]
A direct expansion yields the exact identity
\begin{align}\label{S04-eq:certificate}
 &9Q_iQ_jQ_k-4Q(Q_i+Q_j+Q_k)\notag\\
 &\qquad=UV-9W+S(U^2-3V)+2S^2U+3S^3+4D(2U+3S).
\end{align}
Every term on the right is nonnegative. Indeed
\[
 U^2-3V=\tfrac12\big((x-y)^2+(y-z)^2+(z-x)^2\big)\ge0,
\]
and $UV\ge9W$ follows by multiplying $U\ge3W^{1/3}$ and $V\ge3W^{2/3}$, the arithmetic--geometric mean inequalities. The remaining terms have nonnegative coefficients and variables.

For $n=3$, equality holds at $x=y=z>0$, since $S=D=0$ and $UV=9W$. For $n>3$, fix the selected three variables at $1$ and let all other coordinates decrease to $0$. The ratio tends to $9/4$, proving sharpness within the positive orthant as a supremum. Equality need not be attained there when $n>3$.
\end{proof}

\subsection{Every derivative pattern through degree three}\label{S04-sec-3}
We use the rising factorial $(\beta)_r=\beta(\beta+1)\cdots(\beta+r-1)$. Assume first $\beta>0$. The first derivatives obey
\[
 -\partial_iQ^{-\beta}=\beta Q^{-\beta-1}Q_i>0.
\]
Since $Q_{ii}=0$, diagonal second derivatives are positive. For $i\ne j$,
\begin{equation}\label{S04-eq:second}
 \partial_{ij}Q^{-\beta}
 =\beta Q^{-\beta-2}\big[(\beta+1)Q_iQ_j-Q\big]\ge0
\end{equation}
by Lemma \ref{S04-lem:pair}.

There are exactly three multiplicity patterns at order three.
\begin{enumerate}[label=(\roman*)]
 \item All indices equal:
 \[
 -\partial_{iii}Q^{-\beta}=(\beta)_3Q^{-\beta-3}Q_i^3>0.
 \]
 \item Exactly two indices equal, say $i,i,j$ with $i\ne j$:
 \[
 -\partial_{iij}Q^{-\beta}
 =\beta(\beta+1)Q^{-\beta-3}Q_i
 \big[(\beta+2)Q_iQ_j-2Q\big]\ge0.
 \]
 The bracket is at least $\beta Q_iQ_j$ by Lemma \ref{S04-lem:pair}.
 \item Three distinct indices:
 \begin{align}\label{S04-eq:third}
 -\partial_{ijk}Q^{-\beta}
 &=\beta(\beta+1)Q^{-\beta-3}\notag\\
 &\quad\times\big[(\beta+2)Q_iQ_jQ_k-Q(Q_i+Q_j+Q_k)\big].
 \end{align}
 Lemma \ref{S04-lem:triple} makes the bracket nonnegative whenever $\beta+2\ge9/4$.
\end{enumerate}
Thus $\beta\ge1/4$ is sufficient, and $\beta=0$ gives the constant function.

\begin{proof}[Necessity in Theorem \ref{S04-thm:main}]
For $\beta<0$, $Q^{-\beta}$ increases strictly in every coordinate, violating the first-derivative sign. Suppose $0<\beta<1/4$. If $n=3$, evaluate \eqref{S04-eq:third} at $(1,1,1)$. The bracket is
\[
 8(\beta+2)-18=8\beta-2<0.
\]
If $n>3$, take $x_i=x_j=x_k=1$ and all remaining coordinates equal to $\varepsilon>0$. As $\varepsilon\downarrow0$, the same bracket tends to $8\beta-2$, so it is negative for every sufficiently small positive $\varepsilon$.

For $n=2$, $Q=x_1x_2$ and $Q^{-\beta}=x_1^{-\beta}x_2^{-\beta}$ is completely monotone for $\beta\ge0$, while $\beta<0$ again fails at the first derivative.
\end{proof}

\subsection{Extension to all rank-two matroids}\label{S04-sec-4}
A rank-two matroid's nonloops partition into parallel classes $P_1,\ldots,P_m$, where a pair is dependent precisely when it belongs to one class. Bases are pairs chosen from two different classes. In particular $m\ge2$. Define $L_a=\sum_{e\in P_a}x_e$. Then
\begin{equation}\label{S04-eq:matroid}
 B_M(x)=\sum_{a<b}L_aL_b=E_{2,m}(L_1,\ldots,L_m).
\end{equation}

\begin{corollary}\label{S04-cor:ranktwo}
If a rank-two matroid has exactly two nonloop parallel classes, $B_M^{-\beta}$ is alternating through order three exactly for $\beta\ge0$. If it has at least three such classes, the exact exponent set is $\{0\}\cup[1/4,\infty)$.
\end{corollary}
\begin{proof}
A derivative in a variable $x_e$ is exactly the derivative in its class variable $L_a$; mixed derivatives from one class become repeated derivatives in $L_a$. Thus sufficiency follows from Theorem \ref{S04-thm:main}, including every multiplicity pattern. For necessity, any positive value of a class sum can be realized by distributing it equally among the elements of that class. Choose one differentiated element in each of three distinct classes and use the sharpness configuration. Loops do not occur in the polynomial and have zero derivatives.
\end{proof}

\subsection{A fifth-order obstruction at the boundary}\label{S04-sec-5}
The order-three theorem must not be extrapolated. Set $Q=xy+xz+yz$. A direct differentiation gives
\begin{equation}\label{S04-eq:fifth}
 (-1)^5\partial_x^2\partial_y^2\partial_z Q^{-\beta}\big|_{(1,1,1)}
 =\frac{4\beta(\beta+1)(\beta+2)(8\beta^2+8\beta-3)}{3^{\beta+5}}.
\end{equation}
Here is a coefficient check that avoids relying on symbolic software. Put $x=1-u,y=1-v,z=1-w$; then
\[
 Q=3-2(u+v+w)+(uv+uw+vw).
\]
Expand $Q^{-\beta}=3^{-\beta}\sum_{r\ge0}(\beta)_r[2(u+v+w)-(uv+uw+vw)]^r/(r!3^r)$ and collect the coefficient of $u^2v^2w$. Multiplying it by $2!2!1!$ gives \eqref{S04-eq:fifth}. Only $r=3,4,5$ contribute, so this is a finite algebraic verification.

In particular, at $\beta=1/4$ the last polynomial factor is $-1/2$. Therefore $Q^{-1/4}$ satisfies every alternating derivative inequality through order three and fails one of order five. More generally the displayed obstruction applies for
\[
 0<\beta<\frac{\sqrt{10}-2}{4}.
\]
The true complete-monotonicity threshold for this three-variable polynomial is $1/2$ by \cite[Proposition 3.8]{ref005}. We do not determine every intermediate finite-order threshold.

\subsection{Research status}\label{S04-sec-6}
The principal calculation is the dimension-free sharp inequality \eqref{S04-eq:ineq}, with the full remainder \eqref{S04-eq:certificate}. It gives a uniform theorem at every number of variables and for every rank-two matroid, rather than a list of low-dimensional numerical checks. All numerical experiments used while developing the identity are dispensable: every claim is verified by the displayed formulas. This paper neither upgrades the previous round's structural theorem by assertion nor treats a finite derivative check as an all-orders proof. Exact order-four and higher thresholds are left open here.





\clearpage
\section[S05 --- A sharp universal cubic inequality for real-stable polynomials]{A sharp universal cubic inequality for real-stable polynomials}\label{contrib-S05}
\markboth{Stable Polynomials, Matrix Reconstruction, and Lorentz Positivity}{S05 --- Showa}
\noindent\textsc{S05 \quad Showa programme, source round 5}\par\medskip



\begin{quote}\small\textbf{Overview.} We prove a sharp universal finite-order inverse-power theorem. If $P$ is any nonzero real-stable polynomial with nonnegative coefficients, with no restriction on its degree or number of variables, then $P^{-\beta}$ has every alternating derivative sign through total order three on the positive orthant whenever $\beta\ge1/4$. The constant $1/4$ is optimal already for $xy+xz+yz$. The proof gives a cubic differential inequality, an explicit normalized three-derivative constraint, and the equality case. A stability-preserving differential lift handles arbitrary degree, repeated indices, and absent variables. All ingredients needed for the theorem are proved here. This is a finite-order result: complete monotonicity at all orders is neither asserted nor implied.\end{quote}

\subsection{Statement and mathematical context}\label{S05-sec-1}
A nonzero polynomial $P\in\R[x_1,\ldots,x_n]$ is \emph{real stable} if
\[
 P(z_1,\ldots,z_n)\ne0\qquad\text{whenever }\operatorname{Im}z_i>0\text{ for all }i.
\]
Assume throughout that its coefficients are nonnegative. Then $P(x)>0$ for every $x\in(0,\infty)^n$. For a smooth function $F$, alternating signs through order three means
\[
 (-1)^{|\alpha|}\partial^\alpha F(x)\ge0\qquad(0\le|\alpha|\le3).
\]
Here $\alpha$ is an arbitrary multi-index, so repeated derivatives are included.

\begin{theorem}[Universal sharp threshold]\label{S05-thm:universal}
For every nonzero real-stable polynomial $P$ with nonnegative coefficients and every $\beta\ge1/4$, the function $P^{-\beta}$ has alternating signs through order three on the positive orthant. The assertion also holds trivially for $\beta=0$.

These are exactly the exponents for which this conclusion holds for \emph{every} polynomial in the stated class: every $0<\beta<1/4$ fails for $P(x,y,z)=xy+xz+yz$, and every $\beta<0$ fails for $P(x)=1+x$.
\end{theorem}

\begin{theorem}[Sharp cubic differential inequality]\label{S05-thm:cubic}
Under the same assumptions, for any three indices $i,j,k$, not necessarily distinct,
\begin{equation}\label{S05-eq:cubic}
 P^2P_{ijk}-\frac54P(P_iP_{jk}+P_jP_{ik}+P_kP_{ij})
 +\frac{45}{16}P_iP_jP_k\ge0.
\end{equation}
If $P_iP_jP_k>0$ at the evaluation point, equality holds if and only if
\begin{equation}\label{S05-eq:equality}
 \frac{PP_{ij}}{P_iP_j}=\frac{PP_{ik}}{P_iP_k}=\frac{PP_{jk}}{P_jP_k}=\frac34,
 \qquad P_{ijk}=0.
\end{equation}
The derivatives in these formulas are ordinary derivatives, with their usual multiplicities when indices repeat.
\end{theorem}

Real stability, Rayleigh inequalities, and negative dependence form an established theory \cite{ref006,ref007}. Scott and Sokal \cite{ref005} study the much stronger all-orders inverse-power problem. The present theorem concerns a sharp universal order-three truncation. We found no identical statement in the primary sources inspected, but this is not an exhaustive novelty certificate or a claim to resolve a named all-orders conjecture.

\subsection{Two elementary stability lemmas}\label{S05-sec-2}
The main technical issue is that a general stable polynomial is not multiaffine. We avoid assuming that it is.

\begin{lemma}[A differential lift]\label{S05-lem:lift}
If $P$ is stable, then
\[
 \mathcal L_iP(x,y)=P(x)+y\partial_iP(x)
\]
is stable as a polynomial in all the old variables and the new variable $y$. The polynomial is not identically zero. Real coefficients and nonnegative coefficients are preserved whenever present.
\end{lemma}
\begin{proof}
Fix the old variables other than $x_i$ in the upper half-plane. The resulting univariate polynomial $g(z)$ is nonzero and has no roots in that half-plane. If it is constant, the conclusion is immediate. Otherwise write $g(z)=c\prod_{a=1}^d(z-\rho_a)$, where every $\operatorname{Im}\rho_a\le0$. At a point $z$ in the upper half-plane,
\[
 \frac{g'(z)}{g(z)}=\sum_{a=1}^d\frac1{z-\rho_a}
\]
has strictly negative imaginary part. If $g(z)+yg'(z)=0$, then $y=-1/(g'(z)/g(z))$ also has strictly negative imaginary part, which is impossible for an upper-half-plane $y$. The nonzero claim follows from the coefficient of $y^0$, which is $P$. The assertions about coefficients follow directly from the formula.
\end{proof}

\begin{lemma}[Real Rayleigh inequality in two multiaffine coordinates]\label{S05-lem:rayleigh}
Let $H$ be a real-stable polynomial, multiaffine in two particular variables $u,v$, with no degree restriction in its other variables. After fixing the other variables to arbitrary real values, write
\[
 H(u,v)=auv+bu+cv+d.
\]
Then $bc-ad\ge0$. Consequently $H_uH_v-HH_{uv}\ge0$ at every real point.
\end{lemma}
\begin{proof}
Suppose $bc-ad<0$. The denominator $au+c$ is not identically zero. Choose $u$ in the upper half-plane with $au+c\ne0$ and put
\[
 v=-\frac{bu+d}{au+c}.
\]
The real-coefficient fractional-linear formula gives
\[
 \operatorname{Im}v=\frac{(ad-bc)\operatorname{Im}u}{|au+c|^2}>0.
\]
If there are no other variables, this already contradicts stability. Otherwise perturb each fixed real variable by a sufficiently small positive imaginary part. The four coefficients and the displayed root $v$ vary continuously, and the denominator stays nonzero. Thus $v$ remains in the upper half-plane and gives a zero with every variable in that half-plane, again a contradiction. Finally $H_uH_v-HH_{uv}=bc-ad$ identically in $u,v$.
\end{proof}

Lemma \ref{S05-lem:rayleigh} is the necessary direction of the familiar multiaffine Rayleigh criterion. The elementary proof is included so that no converse characterization or polarization theorem is required for our main result.

\subsection{The exact local three-variable constraint}\label{S05-sec-3}
Fix $x\in(0,\infty)^n$ and indices $i,j,k$. Apply three lifts, with three \emph{different new variables} even if some old indices coincide:
\[
 H(x;y_1,y_2,y_3)
 =(1+y_1\partial_i)(1+y_2\partial_j)(1+y_3\partial_k)P(x).
\]
The full polynomial is stable by Lemma \ref{S05-lem:lift}, and is multiaffine in the $y$ variables. At the fixed point $x$, its coefficients are precisely the derivatives of $P$ with the indicated multiplicities.

If one of $P_i,P_j,P_k$ vanishes at a positive point, nonnegative coefficients imply that $P$ is independent of that variable. Every derivative of $P^{-\beta}$ involving that variable vanishes. We may therefore suppose all three derivatives are positive and put
\[
 a_1=P_i/P,\qquad a_2=P_j/P,\qquad a_3=P_k/P.
\]
Substitute $y_\ell=z_\ell/a_\ell$ and divide by $P$. The normalized auxiliary polynomial at the fixed old point is
\begin{equation}\label{S05-eq:normal}
 f(z)=1+z_1+z_2+z_3+r_{12}z_1z_2+r_{13}z_1z_3+r_{23}z_2z_3+qz_1z_2z_3,
\end{equation}
where
\[
 r_{12}=\frac{PP_{ij}}{P_iP_j},\quad
 r_{13}=\frac{PP_{ik}}{P_iP_k},\quad
 r_{23}=\frac{PP_{jk}}{P_jP_k},\quad
 q=\frac{P^2P_{ijk}}{P_iP_jP_k}.
\]
We only need the real Rayleigh inequalities of $f$, which follow directly from Lemma \ref{S05-lem:rayleigh} applied to $H$ before fixing the old variables. Thus no boundary-specialization assertion about stability is hidden in the argument.

\begin{lemma}[Normalized constraint]\label{S05-lem:constraint}
The numbers above satisfy $0\le r_{ab}\le1$, $q\ge0$, and, with $R=r_{12}+r_{13}+r_{23}$,
\begin{equation}\label{S05-eq:disc}
 (q-(R-2))^2\le4(1-r_{12})(1-r_{13})(1-r_{23}).
\end{equation}
In particular, writing $s=3-R\in[0,3]$,
\begin{equation}\label{S05-eq:envelope}
 q\ge\max\left\{0,\,1-s-2(s/3)^{3/2}\right\}.
\end{equation}
\end{lemma}
\begin{proof}
Nonnegative coefficients give $r_{ab},q\ge0$, and the Rayleigh inequality at the origin gives $r_{ab}\le1$. Keeping $z_3=t$ real and taking the Rayleigh difference in $z_1,z_2$ gives the nonnegative quadratic
\[
 (1-r_{12})+(r_{13}+r_{23}-r_{12}-q)t+(r_{13}r_{23}-q)t^2\ge0
 \qquad(t\in\R).
\]
A quadratic nonnegative for every real $t$ has nonpositive discriminant, including the degenerate constant and zero cases. Expanding that discriminant gives exactly \eqref{S05-eq:disc}. Therefore
\[
 q\ge R-2-2\sqrt{(1-r_{12})(1-r_{13})(1-r_{23})}.
\]
The three nonnegative factors under the square root have sum $s$, so their product is at most $(s/3)^3$ by arithmetic--geometric mean. Combining this with $q\ge0$ gives \eqref{S05-eq:envelope}.
\end{proof}

\subsection{The universal exponent and equality}\label{S05-sec-4}
The first two derivative signs follow from the same lift argument. For any pair of active indices, including a repeated index, $PP_{ij}\le P_iP_j$. Hence for $\beta>0$,
\[
 -\partial_iP^{-\beta}=\beta P^{-\beta-1}P_i\ge0,
\]
\[
 \partial_{ij}P^{-\beta}
 =\beta P^{-\beta-2}\big[(\beta+1)P_iP_j-PP_{ij}\big]\ge0.
\]
The third derivative formula is
\begin{align}\label{S05-eq:third}
 -\partial_{ijk}P^{-\beta}
 =\beta P^{-\beta-3}\big[& (\beta+1)(\beta+2)P_iP_jP_k\notag\\
 &-(\beta+1)P(P_iP_{jk}+P_jP_{ik}+P_kP_{ij})+P^2P_{ijk}\big].
\end{align}
This formula is valid with repeated indices as well.

Divide the bracket by $P_iP_jP_k$. With $s=3-R$, the result is
\begin{equation}\label{S05-eq:betaform}
 \Phi_\beta=\beta^2+s\beta+(q+s-1).
\end{equation}
It increases in $\beta\ge0$, since $2\beta+s\ge0$. Thus it suffices to prove $\Phi_{1/4}\ge0$.

If $s\ge3/4$, then $q\ge0$ gives
\[
 \Phi_{1/4}=q+\frac54s-\frac{15}{16}\ge0.
\]
If $0\le s\le3/4$, use \eqref{S05-eq:envelope} and set $t=\sqrt{s/3}\in[0,1/2]$. Then
\begin{align*}
 \Phi_{1/4}
 &\ge\frac1{16}+\frac{s}{4}-2(s/3)^{3/2}\\
 &=\frac{1+12t^2-32t^3}{16}
 =\frac{(1-2t)(16t^2+2t+1)}{16}\ge0.
\end{align*}
This proves Theorem \ref{S05-thm:universal}'s sufficiency and the cubic inequality \eqref{S05-eq:cubic}.

For the equality statement, the first case can have equality only if $s=3/4$ and $q=0$. The second case is strictly positive for $s<3/4$, so equality again forces those two values. Now \eqref{S05-eq:disc} gives
\[
 \frac1{16}\le4(1-r_{12})(1-r_{13})(1-r_{23}).
\]
But the three factors have sum $3/4$, so their product is at most $1/64$. Equality must hold in arithmetic--geometric mean, yielding $r_{12}=r_{13}=r_{23}=3/4$. Conversely those values and $q=0$ give equality directly. This proves Theorem \ref{S05-thm:cubic} in full.

\subsection{Sharpness, local thresholds, and consequences}\label{S05-sec-5}
The polynomial $Q(x,y,z)=xy+xz+yz$ is real stable: when all three variables are in the upper half-plane, each reciprocal has negative imaginary part, and
\[
 Q=xyz(1/x+1/y+1/z)\ne0.
\]
At $(1,1,1)$, $r_{12}=r_{13}=r_{23}=3/4$ and $q=0$. Formula \eqref{S05-eq:betaform} becomes
\[
 \Phi_\beta=\beta^2+\tfrac34\beta-\tfrac14
 =(\beta-\tfrac14)(\beta+1),
\]
which is negative for $0<\beta<1/4$. For $\beta<0$, the first derivative of $(1+x)^{-\beta}$ is positive. Thus the universal exponent set is exactly the one stated.

There is also an exact \emph{local} answer for a specified active triple at a specified point. Put $c=q+s-1$. The third-derivative sign holds for all positive exponents if $c\ge0$. If $c<0$, its exact positive threshold is
\[
 \beta_* =\frac{\sqrt{s^2-4c}-s}{2}\le\frac14.
\]
The maximal local threshold $1/4$ occurs precisely at \eqref{S05-eq:equality}. This statement follows from the increasing quadratic \eqref{S05-eq:betaform}; it is not a claim that one triple determines all higher derivatives.

The theorem applies, in particular, to every real-stable basis generating polynomial and to generating polynomials of strongly Rayleigh distributions. It also applies to nonmultiaffine stable polynomials and to products and positive coordinate rescalings of them. By expanding nonnegative directional derivatives as nonnegative linear combinations of coordinate derivatives, the order-three inequalities hold for any three directions in $[0,\infty)^n$.

No all-orders conclusion is possible at the boundary exponent. For the same $Q$, an exact fifth-derivative calculation gives
\[
 (-1)^5\partial_x^2\partial_y^2\partial_zQ^{-\beta}\big|_{(1,1,1)}
 =\frac{4\beta(\beta+1)(\beta+2)(8\beta^2+8\beta-3)}{3^{\beta+5}},
\]
which is negative at $\beta=1/4$. It can be checked by expanding
$Q(1-u,1-v,1-w)=3-2(u+v+w)+uv+uw+vw$ and extracting the $u^2v^2w$ coefficient. The full complete-monotonicity threshold of this example is $1/2$ \cite{ref005}.

\subsection{Scope of the advance}\label{S05-sec-6}
The preceding round proved an order-three result for elementary symmetric quadratics and rank-two matroids by a special inequality. Here the assumptions are broadened to every degree and every real-stable polynomial with nonnegative coefficients. The differential lift is essential: merely applying a multiaffine formula to repeated derivatives of a general polynomial would be invalid. The proof also identifies the exact normalized extremal jet and gives a local exponent formula. What remains unproved is a sharp universal description at derivative orders four and above, or an exhaustive literature determination of the present theorem's novelty.





\clearpage
\section[S06 --- Sharp fourth-order inverse-power thresholds on Lorentz cones]{Sharp fourth-order inverse-power thresholds on Lorentz cones}\label{contrib-S06}
\markboth{Stable Polynomials, Matrix Reconstruction, and Lorentz Positivity}{S06 --- Showa}
\noindent\textsc{S06 \quad Showa programme, source round 6}\par\medskip


\begin{quote}\small\textbf{Overview.} We determine exactly the exponents for which an inverse Lorentz quadratic has alternating mixed directional derivatives through order four. In dimension three the positive critical exponent is $(\sqrt3-1)/2$; in every dimension at least four it is $(\sqrt{33}-3)/6$. In dimensions one and two every nonnegative exponent works. The corresponding third-order threshold is $1/4$ in all dimensions at least three. The sharp fourth-order configurations are a square and a regular tetrahedron of null directions. The proof is elementary: a Lorentz normalization converts the problem into an optimization over four unit vectors, and two dimension-sensitive Gram inequalities give the exact answer. These are finite-order results, not assertions of complete monotonicity of all orders.\end{quote}
\subsection{The finite-order question and exact answer}\label{S06-sec-1}
For $n\ge2$ let
\[
 L_n=\{(x_0,x')\in\R\times\R^{n-1}:x_0>\|x'\|\},\qquad
 Q_n(x)=x_0^2-\|x'\|^2.
\]
For $n=1$ set $L_1=(0,\infty)$ and $Q_1(x)=x^2$. A smooth function $f$ on $L_n$ is \emph{alternating through order $K$} if
\begin{equation}\label{S06-alt}
 (-1)^kD_{v_1}\cdots D_{v_k}f(x)\ge0
 \quad(0\le k\le K,\ x,v_1,\ldots,v_k\in L_n).
\end{equation}
Directions may repeat. By continuity the same inequalities hold for directions in $\overline L_n$.

Scott and Sokal explicitly ask how inverse-power conditions depend on a finite derivative order $K$ \cite[Section 1.5]{ref005}. Their all-orders theorem for a Lorentz quadratic gives the exponent set $\{0\}\cup[(n-2)/2,\infty)$ when $n\ge3$ \cite[Theorem 1.9]{ref005}. We answer the distinct finite-order problem completely for $K\le4$.
\begin{theorem}\label{S06-main}
Let $\beta\in\R$. The function $Q_n^{-\beta}$ is alternating through order four if and only if $\beta$ belongs to
\[
\begin{cases}
[0,\infty),&n\le2,\\
\{0\}\cup[(\sqrt3-1)/2,\infty),&n=3,\\
\{0\}\cup[(\sqrt{33}-3)/6,\infty),&n\ge4.
\end{cases}
\]
For orders one and two the answer is $[0,\infty)$ in every dimension. For order three it is $[0,\infty)$ if $n\le2$, and $\{0\}\cup[1/4,\infty)$ if $n\ge3$.
\end{theorem}
All constants concern the entire cone and all mixed directions. They must not be read as an exact answer on every smaller prescribed orthant. The proof supplies explicit negative derivatives below every asserted threshold. We found no identical finite-order classification in the primary sources inspected; this is not an exhaustive novelty certificate.

\subsection{Normalization and derivative polynomials}\label{S06-sec-2}
Write $e=(1,0,\ldots,0)$. A positive scaling and a time-orientation-preserving Lorentz transformation reduce any base point of $L_n$ to $e$, without changing signs of mixed derivatives. For completeness, first rotate $x'$ to $(r,0,\ldots,0)$; on the first two coordinates the matrix
\[
 \frac1{\sqrt{x_0^2-r^2}}
 \begin{pmatrix}x_0&-r\\-r&x_0\end{pmatrix}
\]
preserves $Q_n$, preserves $L_n$, and sends $(x_0,r)$ to $(\sqrt{Q_n(x)},0)$. A scaling completes the reduction.

At $e$, each nonzero closed-cone direction can be positively rescaled to
\[
 v_i=\tfrac12(1,u_i),\qquad \|u_i\|\le1.
\]
Then $Q_n(e)=1$, $D_{v_i}Q_n(e)=1$, and
\begin{equation}\label{S06-r}
 r_{ij}:=D_{v_i}D_{v_j}Q_n(e)=\frac{1-u_i\cdot u_j}{2}.
\end{equation}
Since all higher derivatives of $Q_n$ vanish, the product rule gives
\begin{align}
 -D_{v_1}Q_n^{-\beta}(e)&=\beta,\label{S06-first}\\
 D_{v_1}D_{v_2}Q_n^{-\beta}(e)&=\beta[(\beta+1)-r_{12}],\label{S06-second}\\
 -D_{v_1}D_{v_2}D_{v_3}Q_n^{-\beta}(e)
 &=\beta(\beta+1)[\beta+2-r_{12}-r_{13}-r_{23}],\label{S06-third}\\
 D_{v_1}\cdots D_{v_4}Q_n^{-\beta}(e)
 &=\beta(\beta+1)F_\beta,\label{S06-fourth}\\
 F_\beta&=(\beta+2)(\beta+3)-(\beta+2)R+S,\nonumber
\end{align}
where $R=\sum_{i<j}r_{ij}$ and $S=r_{12}r_{34}+r_{13}r_{24}+r_{14}r_{23}$.
One can verify these formulas without a recurrence: partitions of the derivative labels into singletons and pairs contribute the chain-rule terms; there are six one-pair partitions and three two-pair partitions at order four.

The expression in each line is affine in each vector $u_i$ separately. Every vector in the closed unit ball is a convex combination of unit vectors. Therefore positivity for unit $u_i$ implies positivity throughout the ball. This reduction does not assume that derivative labels are distinct in the original problem.

For three unit vectors,
\[
 r_{12}+r_{13}+r_{23}
 =\frac94-\frac14\|u_1+u_2+u_3\|^2\le\frac94.
\]
Three vectors with pairwise inner product $-1/2$ exist in dimension at least two. Together with $0\le r_{12}\le1$, this proves the order-two and order-three assertions for $n\ge3$. Negative $\beta$ fails \eqref{S06-first}.

\subsection{Two four-vector inequalities}\label{S06-sec-3}
For four unit vectors in a Euclidean space define
\[
 T=\sum_{i<j}u_i\cdot u_j,\quad
 U=(u_1\cdot u_2)(u_3\cdot u_4)
 +(u_1\cdot u_3)(u_2\cdot u_4)
 +(u_1\cdot u_4)(u_2\cdot u_3),
 \quad t=\Big\|\sum_i u_i\Big\|^2.
\]
Thus $T=(t-4)/2$. Expanding \eqref{S06-r} gives
\begin{equation}\label{S06-F}
 F_\beta=\beta^2+\beta-\frac34+\frac{2\beta+3}{8}t+\frac14U.
\end{equation}
\begin{lemma}[Dimension-free bound]\label{S06-gram}
For four unit vectors in any Euclidean space,
\begin{equation}\label{S06-Ub}
 U\ge\frac13-\frac53t+\frac{11}{96}t^2.
\end{equation}
At $t=0$ equality holds precisely when all six pairwise inner products equal $-1/3$.
\end{lemma}
\begin{proof}
Let $G=(u_i\cdot u_j)_{i,j=1}^4$, and decompose its matrix relative to $a=\tfrac12(1,1,1,1)$ and $a^\perp$:
\[
 G=\begin{pmatrix}t/4&b^T\\b&C\end{pmatrix}.
\]
Here $G\succeq0$, $C\succeq0$, and $\tr C=4-t/4$. Positivity of each $2\times2$ compression gives
\[
 \|b\|^2\le\frac t4\tr C=t-\frac{t^2}{16},\qquad
 \tr(C^2)\ge\frac{(4-t/4)^2}{3}.
\]
The second inequality is Cauchy--Schwarz for the three eigenvalues of $C$, including zeros. Direct expansion of the six Gram entries yields
\begin{equation}\label{S06-identity}
 U=\frac{t^2}{64}-\frac32\|b\|^2+\frac14\tr(C^2)-1.
\end{equation}
For an independent derivation, put $s=\sum_i u_i$, $M=\sum_i u_i u_i^T$, and use
\[
 U=\frac{t^2}{8}+\frac{\tr(M^2)-4}{4}-\frac{s^TMs}{2}.
\]
The identities $\tr M^2=\tr G^2$ and $s^TMs=4a^TG^2a$ imply \eqref{S06-identity}. Substitution of the two bounds gives \eqref{S06-Ub}. If $t=0$, positivity forces $b=0$; equality requires $C=(4/3)I_3$, or $G=(4/3)(I-aa^T)$. This has all off-diagonal entries $-1/3$, and conversely that Gram matrix attains equality.
\end{proof}
\begin{lemma}[Planar bound]\label{S06-plane}
If the four unit vectors lie in a Euclidean space of dimension at most two, then
\begin{equation}\label{S06-Up}
 U\ge1-t.
\end{equation}
Equality at $t=0$ is attained by $(u_1,u_2,u_3,u_4)=(e_1,e_2,-e_1,-e_2)$.
\end{lemma}
\begin{proof}
Pad the ambient space to dimension two and again set $M=\sum_i u_i u_i^T$. Its two eigenvalues are $2+h,2-h$ for $0\le h\le2$. Consequently $s^TMs\le(2+h)t$. The identity used in the preceding proof implies
\[
 U\ge\frac{t^2}{8}+1+\frac{h^2}{2}-t-\frac{ht}{2}
 =1-t+\frac12\left(h-\frac t2\right)^2\ge1-t.
\]
The displayed square configuration has $s=0$ and $M=2I$.
\end{proof}

\subsection{Proof of the sharp thresholds}\label{S06-sec-4}
In spatial dimension at least three, \eqref{S06-F} and Lemma~\ref{S06-gram} imply
\begin{equation}\label{S06-lower}
 F_\beta\ge\beta^2+\beta-\frac23
 +\frac{6\beta-1}{24}t+\frac{11}{384}t^2.
\end{equation}
For $\beta\ge b_4:=(\sqrt{33}-3)/6$, all three coefficients on the right are nonnegative. Notice that $b_4>1/6$ and $b_4>1/4$, so the lower-order signs already proved apply as well. For a regular tetrahedron of unit vectors, $t=0$, $U=1/3$, and
\[
 F_\beta=\beta^2+\beta-\frac23.
\]
It is negative for $0<\beta<b_4$. The four associated null directions belong to $\overline L_n$. Replacing every $u_i$ by $(1-\epsilon)u_i$ gives directions in $L_n$ and preserves strict negativity for sufficiently small $\epsilon>0$. This proves necessity under exactly the open-cone definition \eqref{S06-alt}.

In spatial dimension two, Lemma~\ref{S06-plane} instead gives
\[
 F_\beta\ge\beta^2+\beta-\frac12+\frac{2\beta+1}{8}t.
\]
Thus $\beta\ge b_3:=(\sqrt3-1)/2$ suffices; here also $b_3>1/4$. The square configuration has $t=0$, $U=1$, hence $F_\beta=\beta^2+\beta-1/2$, proving the matching necessity by the same interior perturbation.

Finally, $Q_2(x)=(x_0-x_1)(x_0+x_1)$ is a product of two positive linear forms on $L_2$. Every nonnegative inverse power of each factor has alternating derivatives of all orders along cone directions, and the Leibniz rule preserves these signs in their product. The one-dimensional case is immediate. The case $\beta=0$ is the constant function. This completes the proof of Theorem~\ref{S06-main}.

\subsection{Consequences, extremizers, and limitations}\label{S06-sec-5}
\begin{corollary}\label{S06-quad}
Let $q$ be a nondegenerate real quadratic form of signature $(1,n-1)$, positive on one component $C$ of $\{q>0\}$. The exponent sets of Theorem~\ref{S06-main} apply unchanged to $q^{-\beta}$ on $C$.
\end{corollary}
\begin{proof}
An invertible real linear change of variables reduces $q$ to $Q_n$ and maps the chosen component to $L_n$ or $-L_n$. It transports every mixed directional inequality, in both directions.
\end{proof}
For $n\ge4$ and $\beta=b_4$, equality of a fourth derivative in four normalized null directions forces $t=0$ in \eqref{S06-lower}, and then forces the regular tetrahedral Gram matrix. Thus the extremizer is rigid up to an orthogonal transformation and individual positive direction rescaling. In dimension three the planar proof forces $s=0$ and $M=2I$ at equality; the square realizes these conditions.

If a quadratic polynomial on an orthant is the pullback of $Q_n$ by a linear map taking the orthant into $L_n$, the corresponding sufficient bounds follow immediately by the chain rule. For example $E_{2,4}(x)=\sum_{i<j}x_ix_j$ has signature $(1,3)$, and at $x=(1,1,1,1)$ its four coordinate directions give the tetrahedral normalized Gram matrix. Therefore its inverse power is alternating through order four on the positive orthant exactly for $\beta=0$ or $\beta\ge b_4$. This recovers the same sharp constant on a familiar stable polynomial, not just on an abstract cone.

The planar cone result does not assert that the square is realizable by four coordinate directions of $E_{2,3}$. Such a restriction changes the optimization problem. Nor do the finite-order thresholds approach the all-orders threshold by any argument given here. Determining the exact order-five threshold, and the full $K$-dependence, remains outside this paper. The recent general-existence preprint \cite{ref008} concerns large exponents and all orders; no result from it is used here.




\clearpage
\section[S07 --- Explicit rank-two alpha-permanent witnesses near the half-integer threshold]{Explicit rank-two alpha-permanent witnesses near the half-integer threshold}\label{contrib-S07}
\markboth{Stable Polynomials, Matrix Reconstruction, and Lorentz Positivity}{S07 --- Showa}
\noindent\textsc{S07 \quad Showa programme, source round 7}\par\medskip
\noindent\textit{Integration note.} The witness family studied here is proved universally extremal in Contribution~\ref{contrib-S08}; quantitative rigidity follows in Contribution~\ref{contrib-S10}.\par\medskip


\begin{quote}\small\textbf{Overview.} For the rank-two correlation matrix $G_m=(\cos(\pi(i-j)/m))_{i,j=0}^{m-1}$ we evaluate the $\alpha$-permanent by a finite positive-binomial sum. Its unique positive sign-change parameter $a_m$ increases strictly to $1/2$, with $1/2-a_m\sim m/2^{m+1}$. Consequently every $0<\alpha<1/2$ has an explicit rank-two positive-semidefinite counterexample whose size is $\log_2(1/(1-2\alpha))+\log_2\log_2(1/(1-2\alpha))+O(1)$ as $\alpha\uparrow1/2$. A finite inequality certifies negativity without a numerical root computation. The construction does not claim optimal matrix size or settle positivity in any exponent interval above $1/2$. The proof uses an elementary disk beta integral and the product formula for a regular polygon.\end{quote}
\subsection{An explicit witness family}\label{S07-sec-1}
For a real $m\times m$ matrix $A$ define
\[
 \operatorname{per}_{\alpha}(A)
 =\sum_{\sigma\in S_m}\alpha^{c(\sigma)}\prod_{i=0}^{m-1}A_{i,\sigma(i)},
\]
where $c(\sigma)$ is the number of cycles, including fixed points. Put
\[
 w_j=(\cos(\pi j/m),\sin(\pi j/m)),\qquad
 G_m=(w_i\cdot w_j)_{i,j=0}^{m-1}.
\]
For $m\ge3$, $G_m$ is a real correlation matrix of rank two.
\begin{theorem}\label{S07-exact}
Let $m\ge3$ and $\alpha>0$. With $(a)_m=a(a+1)\cdots(a+m-1)$,
\begin{equation}\label{S07-formula}
 \operatorname{per}_\alpha(G_m)
 =\frac{2(2\alpha)_m}{4^m}
 \left[1+(2\alpha-1)\sum_{k=1}^{\lfloor m/2\rfloor}
 \frac{\binom m{2k}}{2\alpha+2k-1}\right].
\end{equation}
There is a unique $a_m\in(0,1/2)$ such that this permanent is negative for $0<\alpha<a_m$, zero at $a_m$, and positive for $\alpha>a_m$. The sequence $a_m$ is strictly increasing and
\begin{equation}\label{S07-asym}
 \frac12-a_m\sim\frac{m}{2^{m+1}}.
\end{equation}
\end{theorem}
The all-size positivity problem for $\alpha$-permanents was resolved by Br\"and\'en \cite{ref009}, using Scott--Sokal inverse-power theory \cite{ref005}. In particular, existence of negative real positive-semidefinite examples below $1/2$ is not new. The aim here is an explicit rank-two family, an exact finite formula, and a quantitative size estimate. No identical evaluation or threshold asymptotic was found in the primary sources inspected, but the search is not an exhaustive novelty certificate.

\subsection{The determinant coefficient identity}\label{S07-sec-2}
We record the squarefree case of the generalized MacMahon identity, with its short proof so the argument does not depend on an external generating-function theorem.
\begin{lemma}\label{S07-mac}
For any square matrix $A$ and $Z=\diag(z_0,\ldots,z_{m-1})$,
\[
 [z_0\cdots z_{m-1}]\det(I-ZA)^{-\alpha}
 =\operatorname{per}_\alpha(A).
\]
\end{lemma}
\begin{proof}
As formal series,
\[
 \det(I-ZA)^{-\alpha}
 =\exp\left(\alpha\sum_{\ell\ge1}\frac{\tr((ZA)^\ell)}{\ell}\right).
\]
In a squarefree coefficient, each trace term uses distinct labels. Division by $\ell$ identifies the $\ell$ cyclic starting points of an oriented cycle. The exponential selects an unordered collection of disjoint cycles covering all labels. Such a collection is a permutation, and each cycle contributes one factor of $\alpha$ and its matrix-entry product.
\end{proof}
Apply Sylvester's determinant identity to the $2\times m$ matrix with columns $w_j$. The polynomial in Lemma~\ref{S07-mac} becomes
\[
 \det\left(I_2-\sum_j z_jw_jw_j^T\right).
\]
Identify a real symmetric $2\times2$ matrix with
\[
 \begin{pmatrix}t+x&y\\y&t-x\end{pmatrix},\qquad \det=t^2-x^2-y^2=:Q(t,x,y).
\]
The matrix $I_2$ corresponds to $e=(1,0,0)$, and $w_jw_j^T$ corresponds to
\[
 v_j=\tfrac12(1,u_j),\qquad
 u_j=(\cos(2\pi j/m),\sin(2\pi j/m)).
\]
Therefore
\begin{equation}\label{S07-der}
 \operatorname{per}_\alpha(G_m)=(-1)^mD_{v_0}\cdots D_{v_{m-1}}Q^{-\alpha}(e).
\end{equation}

\subsection{A disk beta identity}\label{S07-sec-3}
The following calculation uses a positive integral only for $\alpha>1/2$. We later extend the resulting polynomial identity, not the positive measure, to smaller exponents.
\begin{lemma}\label{S07-disk}
For $\alpha>1/2$, let $V$ have density
\[
 \frac{\alpha-1/2}{\pi}(1-\|v\|^2)^{\alpha-3/2}\,\mathbf1_{\{\|v\|<1\}}
\]
on the unit disk. Let $T$, independently, have gamma density $t^{2\alpha-1}e^{-t}/\Gamma(2\alpha)$ on $(0,\infty)$. Then for $a>\|b\|$,
\begin{equation}\label{S07-diskid}
 (a^2-\|b\|^2)^{-\alpha}
 =\mathbb E\exp\{-T(a-1-b\cdot V)\}.
\end{equation}
\end{lemma}
\begin{proof}
The density integrates to one by polar coordinates. Rotational symmetry and a beta integral give
\[
 \mathbb E V_1^{2k}=\frac{(1/2)_k}{(\alpha+1/2)_k},\qquad
 \mathbb E V_1^{2k+1}=0.
\]
One way to check the even moment is to integrate first over $v_2$; the marginal density of $v_1$ is proportional to $(1-v_1^2)^{\alpha-1}$, after which $s=v_1^2$ gives the beta ratio. Integrating over $T$ reduces the right side of \eqref{S07-diskid} to $\mathbb E(a-b\cdot V)^{-2\alpha}$. Rotate $b$ to the first axis and expand the binomial series, absolutely convergent because $\|b\|/a<1$. Duplication of rising factorials gives
\[
 \frac{(2\alpha)_{2k}}{(2k)!}
 \frac{(1/2)_k}{(\alpha+1/2)_k}
 =\frac{(\alpha)_k}{k!}.
\]
The resulting series is $a^{-2\alpha}(1-\|b\|^2/a^2)^{-\alpha}$.
\end{proof}
Differentiating \eqref{S07-diskid} at $e$ in the null directions $v_j$ gives
\begin{equation}\label{S07-moment}
 (-1)^mD_{v_0}\cdots D_{v_{m-1}}Q^{-\alpha}(e)
 =\frac{(2\alpha)_m}{2^m}\mathbb E\prod_j(1-u_j\cdot V).
\end{equation}
Differentiation is justified by gamma integrability in a neighborhood of $e$.

\subsection{The regular-polygon product and exact formula}\label{S07-sec-4}
For $0\le r\le1$, angular averaging gives
\begin{equation}\label{S07-angular}
 \frac1{2\pi}\int_0^{2\pi}\prod_{j=0}^{m-1}
 \bigl(1-r\cos(\theta-2\pi j/m)\bigr)\,d\theta
 =\frac{(1+\sqrt{1-r^2})^m+(1-\sqrt{1-r^2})^m}{2^m}.
\end{equation}
For a direct verification, use the polynomial identity
\[
 \prod_{j=0}^{m-1}(z-\cos(\theta-2\pi j/m))
 =2^{1-m}(T_m(z)-\cos(m\theta)),
\]
where $T_m(\cos\theta)=\cos(m\theta)$. The identity follows because the two degree-$m$ polynomials have the same roots and leading coefficient; repeated roots follow by continuity. The formula
$T_m(z)=\tfrac12[(z+\sqrt{z^2-1})^m+(z-\sqrt{z^2-1})^m]$
then proves \eqref{S07-angular}, including the endpoints by continuity.

Under the disk density in Lemma~\ref{S07-disk}, the variable $W=\sqrt{1-\|V\|^2}$ has density $(2\alpha-1)w^{2\alpha-2}$ on $(0,1)$. Equations \eqref{S07-der}, \eqref{S07-moment}, and \eqref{S07-angular} yield
\[
 \operatorname{per}_\alpha(G_m)
 =\frac{(2\alpha)_m}{4^m}(2\alpha-1)
 \int_0^1w^{2\alpha-2}\bigl[(1+w)^m+(1-w)^m\bigr],dw.
\]
Expand the bracket in even powers. The constant term contributes $2$ after multiplication by $2\alpha-1$, and the other terms give \eqref{S07-formula} for $\alpha>1/2$.

Both sides of \eqref{S07-formula} are polynomials in $\alpha$: every denominator $2\alpha+2k-1$ divides $(2\alpha)_m$, since $2k-1\le m-1$. Agreement on an interval therefore proves the identity for every real $\alpha$. In particular there is no illicit continuation of a positive integral to $\alpha<1/2$.

\subsection{Sign changes and explicit certificates}\label{S07-sec-5}
Set $\delta=1-2\alpha$ and
\[
 S_m(\delta)=\sum_{k=1}^{\lfloor m/2\rfloor}\frac{\binom m{2k}}{2k-\delta}.
\]
For $0<\alpha<1/2$, the sign in \eqref{S07-formula} is the sign of $1-\delta S_m(\delta)$. The function $\delta S_m(\delta)$ is continuous and strictly increasing on $[0,1]$, is zero at zero, and at one is at least $\binom m2>1$. Thus there is a unique $\delta_m\in(0,1)$ with
\begin{equation}\label{S07-root}
 \delta_m S_m(\delta_m)=1,
 \qquad a_m=\frac{1-\delta_m}{2}.
\end{equation}
For $\alpha\ge1/2$ every term in \eqref{S07-formula} is nonnegative and the constant term is positive. Also $S_{m+1}(\delta)>S_m(\delta)$ for $0<\delta<1$, by the strict increase of $\binom m2$ and nondecrease of every other coefficient. Hence $\delta_m$ decreases strictly and $a_m$ increases strictly.

An explicit sufficient condition for negativity is
\begin{equation}\label{S07-certificate}
 (1-2\alpha)(2^{m-1}-1)^2>m2^{m-2}.
\end{equation}
Indeed, Cauchy--Schwarz gives
\[
 S_m(0)\ge
 \frac{\left(\sum_{k\ge1}\binom m{2k}\right)^2}
 {\sum_{k\ge1}2k\binom m{2k}}
 =\frac{(2^{m-1}-1)^2}{m2^{m-2}},
\]
and $S_m(\delta)>S_m(0)$. For rational $\alpha$, condition \eqref{S07-certificate} is a finite integer-arithmetic certificate. The matrix entries are explicit real algebraic numbers, as cosines of rational multiples of $\pi$.

Some exact initial thresholds are
\[
 a_3=\frac14,\qquad a_4=\frac{\sqrt3-1}{2},\qquad
 a_5=\frac{\sqrt{41}-3}{8},\qquad a_6=\frac{\sqrt{34}-4}{4}.
\]
For example, $\operatorname{per}_\alpha(G_5)=\alpha(\alpha+1)(\alpha+2)(4\alpha^2+3\alpha-2)/4$. At the endpoint, \eqref{S07-formula} gives the exact positive value
\[
 \operatorname{per}_{1/2}(G_m)=\frac{2m!}{4^m}.
\]

\subsection{Threshold asymptotics and witness size}\label{S07-sec-6}
\begin{lemma}\label{S07-bin}
As $m\to\infty$, $S_m(0)\sim2^m/m$.
\end{lemma}
\begin{proof}
Let $B$ be binomial with parameters $m,1/2$. Then
\[
 2^{-m}mS_m(0)=\mathbb E\left[\frac mB\,
 \mathbf1_{\{B>0,\ B\ {\rm even}\}}\right].
\]
The probability of even $B$ is exactly $1/2$. On $|B-m/2|\le\epsilon m$, the ratio $m/B$ lies between $(1/2+\epsilon)^{-1}$ and $(1/2-\epsilon)^{-1}$. On the complement and $B\ge1$, it is at most $m$. The fourth central moment is $(3m^2-2m)/16$, so Markov's inequality makes this complementary contribution $O_\epsilon(m^{-1})$. The omitted event $B=0$ has probability $2^{-m}$. First let $m\to\infty$ and then $\epsilon\downarrow0$; the expectation converges to one.
\end{proof}
Since $S_m(\delta_m)\ge S_m(0)$, equation \eqref{S07-root} gives $\delta_m\le1/S_m(0)\to0$. For $0\le\delta<1$ and $2k\ge2$,
\[
 \frac1{2k}\le\frac1{2k-\delta}\le\frac1{1-\delta/2}\frac1{2k}.
\]
Consequently $S_m(\delta_m)/S_m(0)\to1$, so \eqref{S07-root} and Lemma~\ref{S07-bin} give $\delta_m\sim m/2^m$, proving \eqref{S07-asym}.

Let $m_*(\alpha)$ be the smallest $m\ge3$ for which this particular family has negative $\alpha$-permanent. Writing $\delta=1-2\alpha\downarrow0$, strict monotonicity and the asymptotic just proved imply
\[
 m_*(\alpha)=\log_2(1/\delta)+\log_2\log_2(1/\delta)+O(1).
\]
Indeed $\delta_m\sim m2^{-m}$ reduces the crossing to $m-\log_2m=\log_2(1/\delta)+O(1)$; evaluating the left side at the displayed expression plus or minus a fixed sufficiently large constant brackets the crossing.

This logarithmic law is optimal only \emph{within the stated regular-polygon family}. We do not prove that another rank-two configuration cannot produce smaller counterexamples, nor that the roots $a_m$ are the universal order-$m$ inverse-Lorentz thresholds when $m\ge5$. The exact family nonetheless gives explicit tests at every exponent below the first positive half-integer, without invoking a nonconstructive failure of complete monotonicity.




\clearpage
\section[S08 --- All finite-order thresholds on the three-dimensional Lorentz cone]{All finite-order thresholds on the three-dimensional Lorentz cone}\label{contrib-S08}
\markboth{Stable Polynomials, Matrix Reconstruction, and Lorentz Positivity}{S08 --- Showa}
\noindent\textsc{S08 \quad Showa programme, source round 8}\par\medskip


\begin{quote}\small\textbf{Overview.} We determine the exact exponent threshold for every finite order of complete monotonicity of $(t^2-x^2-y^2)^{-\alpha}$ on the three-dimensional Lorentz cone. For order $m\ge3$, the threshold $a_m$ is the unique number in $(0,1/2)$ satisfying an explicit binomial-sum equation. Equivalently, every real positive-semidefinite $m\times m$ matrix of rank at most two has nonnegative $\alpha$-permanent exactly when $\alpha=0$ or $\alpha\ge a_m$. Regular projective polygons are the unique nondegenerate equality configurations at the threshold. The proof expands a circle-rooted polynomial in an orthogonal Fourier basis and compares paired Bernstein terms under an order-reversing beta functional. In particular $1/2-a_m\sim m/2^{m+1}$. This solves the entire finite-order problem for this one cone, not for general hyperbolic polynomials or higher-dimensional Lorentz cones.\end{quote}
\subsection{Exact finite-order and matrix statements}\label{S08-sec-1}
Write
\[
 L=\{(t,x,y):t>\sqrt{x^2+y^2}\},\quad Q(t,x,y)=t^2-x^2-y^2.
\]
A smooth $f$ is \emph{alternating through order $m$ on $L$} if
\[
 (-1)^kD_{v_1}\cdots D_{v_k}f(z)\ge0
 \quad(0\le k\le m,\ z,v_1,\ldots,v_k\in L).
\]
Repeated directions are allowed. Set $a_1=a_2=0$. For $m\ge3$, define $\delta_m\in(0,1)$ by
\begin{equation}\label{S08-root}
 \delta_m\sum_{k=1}^{\lfloor m/2\rfloor}
 \frac{\binom m{2k}}{2k-\delta_m}=1,
 \qquad a_m=(1-\delta_m)/2.
\end{equation}
Existence and uniqueness are proved below.
\begin{theorem}\label{S08-main}
For every integer $m\ge1$ and real $\alpha$, $Q^{-\alpha}$ is alternating through order $m$ on $L$ if and only if
\[
 \alpha\in\{0\}\cup[a_m,\infty).
\]
The thresholds $a_m$ are strictly increasing for $m\ge2$ and satisfy
\[
 \frac12-a_m\sim\frac{m}{2^{m+1}}.
\]
\end{theorem}
For a matrix $A$, put $\operatorname{per}_\alpha(A)=\sum_{\sigma\in S_m}\alpha^{c(\sigma)}\prod_iA_{i,\sigma(i)}$, where $c(\sigma)$ counts cycles.
\begin{theorem}\label{S08-matrix}
Fix $m\ge3$ and $\alpha\ge0$. The inequality $\operatorname{per}_\alpha(A)\ge0$ holds for every real positive-semidefinite $m\times m$ matrix of rank at most two if and only if $\alpha=0$ or $\alpha\ge a_m$.
At $\alpha=a_m$, a correlation matrix of rank at most two has zero $\alpha$-permanent if and only if, up to a simultaneous permutation and switching matching row and column signs, it equals
\[
 G_m=(\cos(\pi(i-j)/m))_{i,j=0}^{m-1}.
\]
\end{theorem}
Scott and Sokal ask how inverse-power conditions depend on a finite derivative order \cite[Section 1.5]{ref005}. Theorem~\ref{S08-main} answers that question for the fixed cone $L$. The all-orders threshold $1/2$ and the all-size matrix-positivity context are established theory \cite{ref005,ref009}. The new step here is universal optimality and rigidity of the finite-order regular-polygon witnesses. The author found no identical theorem in the inspected primary literature; no exhaustive novelty certification is claimed.

\subsection{Reduction to unit circles}\label{S08-sec-2}
A time-oriented Lorentz transformation and a positive scaling take any base point of $L$ to $e=(1,0,0)$, preserving derivative signs. Explicitly, after rotating $(x,y)$ to $(r,0)$, the $2\times2$ block
\[
 \frac1{\sqrt{t^2-r^2}}\begin{pmatrix}t&-r\\-r&t\end{pmatrix}
\]
sends $(t,r)$ to $(\sqrt Q,0)$ and preserves the quadratic form. Each nonzero closed-cone direction can then be positively rescaled to $v_j=\tfrac12(1,u_j)$ with $\|u_j\|\le1$.
Mixed derivatives are affine in each $u_j$ separately. Every point of the unit disk is a convex combination of unit-circle points. It therefore suffices for positivity to take
\[
 u_j=(\cos\theta_j,\sin\theta_j),\qquad
 v_j=\tfrac12(1,u_j).
\]
For necessity, any strictly negative null-direction derivative persists after replacing $u_j$ by $(1-\epsilon)u_j$, so it also gives an open-cone counterexample.

For the matrix formulation, if a diagonal entry of $A\succeq0$ is zero, its row and column vanish and its $\alpha$-permanent is zero. Otherwise diagonal normalization multiplies the permanent by $\prod_i A_{ii}>0$, so reduce to a correlation matrix $A=(w_i\cdot w_j)$ with unit $w_j\in\R^2$. Writing $w_j=(\cos(\theta_j/2),\sin(\theta_j/2))$ identifies $w_jw_j^T$ with $v_j$ under
\[
 \begin{pmatrix}t+x&y\\y&t-x\end{pmatrix}\longleftrightarrow(t,x,y).
\]
The formal identity
\begin{equation}\label{S08-mac}
 [z_1\cdots z_m]\det(I-ZA)^{-\alpha}=\operatorname{per}_\alpha(A)
\end{equation}
follows by expanding $\exp(\alpha\sum_{\ell\ge1}\tr((ZA)^\ell)/\ell)$: trace terms encode cycles and the exponential encodes unordered disjoint cycle collections. Sylvester's identity gives
\begin{equation}\label{S08-der}
 \operatorname{per}_\alpha(A)=(-1)^mD_{v_1}\cdots D_{v_m}Q^{-\alpha}(e).
\end{equation}
Thus it remains to control the same circle-direction expression in both problems.

\subsection{The beta functional and Fourier decomposition}\label{S08-sec-3}
For a polynomial $f$ and $0<\delta<1$, define
\begin{equation}\label{S08-functional}
 \mathcal L_\delta(f)=f(0)-\delta\int_0^1
 \frac{f(w)-f(0)}{w^{1+\delta}}\,dw.
\end{equation}
The integral converges, and for $f(w)=\sum_{j\ge0}f_jw^j$ it equals
\[
 \mathcal L_\delta(f)=f_0-\delta\sum_{j\ge1}\frac{f_j}{j-\delta}.
\]
Crucially, if $f(0)=h(0)$ and $f(w)\le h(w)$ on $[0,1]$, then
\begin{equation}\label{S08-reverse}
 \mathcal L_\delta(f)\ge\mathcal L_\delta(h).
\end{equation}
This is an order reversal, not a positive-measure representation.

Let
\[
 p(z)=\prod_{j=1}^m(z-e^{i\theta_j})=\sum_{k=0}^mc_kz^k,
 \qquad b_k=|c_k|^2,\quad B=\sum_{k=0}^mb_k.
\]
Since all roots lie on the unit circle, $p$ is self-inversive, and hence
\begin{equation}\label{S08-self}
 b_k=b_{m-k},\qquad b_0=b_m=1,\qquad B\ge2.
\end{equation}
For example self-inversivity follows by comparing $p(z)$ with $z^m\overline{p(1/\overline z)}$: the two polynomials have the same roots and proportional coefficients with a unit-modulus factor.
Set
\[
 f_k(w)=(1+w)^{m-k}(1-w)^k,\qquad
 g(w)=\sum_{k=0}^mb_k f_k(w),\qquad h(w)=f_0(w)+f_m(w).
\]
\begin{lemma}\label{S08-rep}
For $0<\alpha<1/2$ and $\delta=1-2\alpha$,
\begin{equation}\label{S08-representation}
 (-1)^mD_{v_1}\cdots D_{v_m}Q^{-\alpha}(e)
 =\frac{(2\alpha)_m}{4^m}\mathcal L_\delta(g).
\end{equation}
For $\alpha=1/2$ the same expression is $m!B/4^m$. For $\alpha>1/2$ it is nonnegative.
\end{lemma}
\begin{proof}
Initially assume $\alpha>1/2$. On the unit disk take the probability density
$(\alpha-1/2)\pi^{-1}(1-\|v\|^2)^{\alpha-3/2}$, and independently take a gamma variable $T$ of shape $2\alpha$ and rate one. A beta integral gives
$\mathbb E V_1^{2k}=(1/2)_k/(\alpha+1/2)_k$, while odd moments vanish. Thus the binomial series and
$(2\alpha)_{2k}=4^k(\alpha)_k(\alpha+1/2)_k$ imply
\[
 Q(t,x,y)^{-\alpha}=\mathbb E e^{-T(t-1-xV_1-yV_2)}.
\]
Differentiating yields $(2\alpha)_m2^{-m}\mathbb E\prod_j(1-u_j\cdot V)$, proving nonnegativity in this range.

For fixed radius $r$, put $w=\sqrt{1-r^2}$ and $\rho=r/(1+w)$. The identity
$1-r\cos\phi=\tfrac12(1+w)|1-\rho e^{i\phi}|^2$
and Fourier orthogonality show
\[
 \frac1{2\pi}\int_0^{2\pi}\prod_j(1-r\cos(\phi-\theta_j))\,d\phi
 =\frac1{2^m}\sum_{k=0}^mb_k(1+w)^{m-k}(1-w)^k
 =2^{-m}g(w).
\]
The radius substitution makes $w$ have density $(2\alpha-1)w^{2\alpha-2}$. Consequently the derivative equals
\[
 \frac{(2\alpha)_m}{4^m}(2\alpha-1)\int_0^1w^{2\alpha-2}g(w)\,dw.
\]
Expand $g$ in powers of $w$. The constant term contributes $g(0)$ and each other coefficient is multiplied by $(2\alpha-1)/(2\alpha+j-1)$. Every denominator with $1\le j\le m$ divides $(2\alpha)_m$. Thus the resulting identity is polynomial in $\alpha$ and remains valid for every real $\alpha$. For $0<\alpha<1/2$ it is exactly \eqref{S08-representation}. At $\alpha=1/2$ only the constant term remains, yielding $m!B/4^m$.
\end{proof}

\subsection{The sharp comparison}\label{S08-sec-4}
\begin{lemma}\label{S08-comparison}
For $0\le w\le1$,
\begin{equation}\label{S08-bernstein}
 g(w)\le\frac B2h(w),\qquad g(0)=\frac B2h(0)=B.
\end{equation}
If some $b_k>0$ for $0<k<m$, the inequality is strict for every $0<w<1$.
\end{lemma}
\begin{proof}
For $x=1+w$, $y=1-w$, convexity of the function
$s\mapsto x^{m-s}y^s+x^sy^{m-s}$ on $[0,m]$ gives
\[
 f_k(w)+f_{m-k}(w)\le x^m+y^m=h(w).
\]
For $0<w<1$ this convexity is strict, so the inequality is strict for $0<k<m$. Multiply by $b_k/2$, sum, and use \eqref{S08-self}.
\end{proof}
Now
\[
 \mathcal L_\delta(h)
 =2\left[1-\delta\sum_{k=1}^{\lfloor m/2\rfloor}
 \frac{\binom m{2k}}{2k-\delta}\right].
\]
Equations \eqref{S08-reverse} and \eqref{S08-bernstein} yield the central inequality
\begin{equation}\label{S08-central}
 \mathcal L_\delta(g)\ge\frac B2\mathcal L_\delta(h).
\end{equation}
The quantity $\delta\sum_{k\ge1}\binom m{2k}/(2k-\delta)$ is strictly increasing in $\delta\in[0,1]$, goes from zero to a value greater than one when $m\ge3$, and strictly increases with $m$. This proves existence and uniqueness in \eqref{S08-root}, and strict increase of $a_m$.

If $a_m\le\alpha<1/2$, then $\delta\le\delta_m$, so the right side of \eqref{S08-central} is nonnegative. Lemma~\ref{S08-rep} handles $\alpha\ge1/2$. For $0<\alpha<a_m$, choose the $m$ angles to be a regular polygon. Then $p(z)=z^m-\zeta$ for some $|\zeta|=1$, all interior $b_k$ vanish, $g=h$, and the derivative is strictly negative. At the critical exponent, equality in \eqref{S08-central} is possible only when all interior $b_k$ vanish, which is precisely a regular polygon of the doubled angles $\theta_j$. This proves Theorem~\ref{S08-matrix}, including its sign-switching interpretation.

For $m=1,2$, direct differentiation gives $-D_vQ^{-\alpha}(e)=\alpha$ and
$D_vD_wQ^{-\alpha}(e)=\alpha[(\alpha+1)-(1-u\cdot u')/2]\ge0$ for $\alpha\ge0$. Negative exponents fail at order one. Since the thresholds increase, the order-$m$ positivity just proved implies all lower-order inequalities. The disk reduction and the null-to-open perturbation from Section~\ref{S08-sec-2} prove Theorem~\ref{S08-main}.

\subsection{Extremal values and quantitative convergence}\label{S08-sec-5}
In the range $a_m\le\alpha\le1/2$, equations \eqref{S08-representation} and \eqref{S08-central} imply, for every real rank-at-most-two correlation matrix $A$,
\[
 \operatorname{per}_\alpha(A)\ge\frac B2\operatorname{per}_\alpha(G_m)
 \ge\operatorname{per}_\alpha(G_m).
\]
The regular-projective-polygon matrix thus minimizes the permanent in this parameter range. For $a_m<\alpha\le1/2$ equality also forces $B=2$, hence the same polygon configuration. No analogous minimum assertion for $\alpha>1/2$ is made.

For completeness we prove the asymptotic threshold formula. Let
$S_m(0)=\sum_{k\ge1}\binom m{2k}/(2k)$. If $N$ is binomial with parameters $m,1/2$, then
\[
 m2^{-m}S_m(0)=\mathbb E\left[\frac mN\mathbf1_{\{N>0,\ N\ {\rm even}\}}\right]\longrightarrow1.
\]
Indeed $\mathbb P(N\text{ even})=1/2$; on $|N-m/2|\le\epsilon m$, $m/N$ is uniformly close to two as $\epsilon\downarrow0$. The complementary expectation is $O_\epsilon(1/m)$ because $m/N\le m$ for $N\ge1$ and $\mathbb E(N-m/2)^4=(3m^2-2m)/16$. The event $N=0$ has probability $2^{-m}$. Therefore $S_m(0)\sim2^m/m$.
Equation \eqref{S08-root} gives $\delta_m\le1/S_m(0)\to0$, while
\[
 S_m(0)\le S_m(\delta_m)\le\frac{S_m(0)}{1-\delta_m/2}.
\]
Hence $\delta_m\sim1/S_m(0)\sim m/2^m$, as asserted.

The first values are $a_3=1/4$, $a_4=(\sqrt3-1)/2$, and $a_5=(\sqrt{41}-3)/8$. Unlike a merely sufficient family of witnesses, Theorem~\ref{S08-matrix} now proves that the first negative member $G_m$ has the \emph{smallest possible size among all real rank-at-most-two positive-semidefinite matrices}. As $\alpha\uparrow1/2$ this size is
\[
 \log_2\frac1{1-2\alpha}+\log_2\log_2\frac1{1-2\alpha}+O(1).
\]
This optimal-size claim is confined to rank at most two. Arbitrary matrix rank, higher-dimensional Lorentz cones, and general stable polynomials are different problems.




\clearpage
\section[S09 --- Fifth-order Lorentz thresholds and five-point spherical energy]{Fifth-order Lorentz thresholds and five-point spherical energy}\label{contrib-S09}
\markboth{Stable Polynomials, Matrix Reconstruction, and Lorentz Positivity}{S09 --- Showa}
\noindent\textsc{S09 \quad Showa programme, source round 9}\par\medskip


\begin{quote}\small\textbf{Overview.} We determine the exact fifth-order inverse-power threshold for a Lorentz quadratic in every dimension. Besides the trivial dimensions one and two, the thresholds are $(\sqrt{41}-3)/8$ in dimension three, $(\sqrt{57}-3)/8$ in dimension four, and $(\sqrt{66}-3)/8$ in every dimension at least five. The sharp null-direction configurations are respectively a regular pentagon, a triangular bipyramid, and a regular four-simplex. The higher-dimensional proof is elementary Gram optimization. The dimension-four proof explicitly uses Tumanov's established five-point biquadratic-energy theorem, followed by a spectral estimate; that theorem is a cited input, not reproved here. This identifies a precise connection between finite-order complete monotonicity and spherical energy minimization.\end{quote}
\subsection{Exact statement and external input}\label{S09-sec-1}
Let $Q_n(t,x)=t^2-\|x\|^2$ on $L_n=\{(t,x):t>\|x\|\}\subset\R^n$, with the usual interpretation $Q_1(t)=t^2$. We say $f$ is alternating through order five if
\[
 (-1)^kD_{v_1}\cdots D_{v_k}f(z)\ge0
 \quad(0\le k\le5,\ z,v_1,\ldots,v_k\in L_n).
\]
\begin{theorem}\label{S09-main}
For real $\beta$, $Q_n^{-\beta}$ is alternating through order five if and only if $\beta\in\{0\}\cup[b_n,\infty)$, where
\[
 b_n=\begin{cases}
 0,&n\le2,\\
 (\sqrt{41}-3)/8,&n=3,\\
 (\sqrt{57}-3)/8,&n=4,\\
 (\sqrt{66}-3)/8,&n\ge5.
 \end{cases}
\]
At the positive critical exponents, normalized null-direction extremizers are a regular pentagon for $n=3$, a triangular bipyramid for $n=4$, and a regular four-simplex for $n\ge5$, up to orthogonal transformations and permutation.
\end{theorem}
The finite-order inverse-power question is explicitly raised by Scott and Sokal \cite[Section 1.5]{ref005}; their all-orders Lorentz theorem is a different result. The present classification is only for order five and quadratic cones. The dimension-four case uses the following existing theorem in an equivalent form.
\begin{theorem}[Tumanov, cited input]\label{S09-tumanov}
If $u_1,\ldots,u_5$ are unit vectors in $\R^3$, $s=\sum_i u_i$, and $M=\sum_i u_i u_i^T$, then
\begin{equation}\label{S09-tum}
 \tr(M^2)+2\|s\|^2\ge\frac{17}{2}.
\end{equation}
Equality holds precisely for a triangular bipyramid: an antipodal pair and an equilateral triangle in the perpendicular plane.
\end{theorem}
This is Tumanov's minimization of $\sum_{i<j}(1+u_i\cdot u_j)^2$ \cite[Sections 2--6]{ref011}, rewritten using
\[
 \sum_{i<j}(1+u_i\cdot u_j)^2=\tfrac12\tr(M^2)+\|s\|^2+\tfrac52.
\]
Its minimum is $27/4$. Tumanov's theorem is unconditional published mathematics; its proof is not reproduced in this paper. All remaining ingredients are proved below. No claim is made that this spherical-energy result is new.

\subsection{Fifth derivatives and Euler descent}\label{S09-sec-2}
Lorentz transformations and positive scaling reduce the base point to $e=(1,0)$. Explicitly, rotate $x$ to $(r,0,\ldots,0)$ and apply the block
$Q_n(t,x)^{-1/2}\left(\begin{smallmatrix}t&-r\\-r&t\end{smallmatrix}\right)$, then scale. Each nonzero closed-cone direction can be positively rescaled to $v_i=\tfrac12(1,u_i)$ with $\|u_i\|\le1$. Multilinearity permits reducing each $u_i$ to the unit sphere by convex combinations. Strict negative values on that sphere persist on replacing $u_i$ by $(1-\epsilon)u_i$, so necessity is valid for open-cone directions too.

Write
\[
 r_{ij}=\frac{1-u_i\cdot u_j}{2},\qquad
 R=\sum_{i<j}r_{ij},\qquad
 S=\sum_{\{\{i,j\},\{k,l\}\}}r_{ij}r_{kl},
\]
where $S$ sums over the fifteen unordered pairs of disjoint edges on five labels. Partitions of the five derivative labels into singletons and pairs give
\begin{equation}\label{S09-der}
 (-1)^5D_{v_1}\cdots D_{v_5}Q_n^{-\beta}(e)
 =\beta(\beta+1)(\beta+2)F_\beta,
 \quad F_\beta=(\beta+3)(\beta+4)-(\beta+3)R+S.
\end{equation}
Indeed the zero-pair, one-pair, and two-pair partitions have $1,10,15$ possibilities respectively, and all higher derivatives of $Q_n$ vanish.

For $\beta>0$, positivity at order five implies all lower-order signs. If $g(z)=D_{v_1}\cdots D_{v_k}Q_n^{-\beta}(z)$ with fixed directions, then $g$ is homogeneous of degree $-2\beta-k$, hence
\[
 D_zg(z)=-(2\beta+k)g(z).
\]
Thus an alternating sign at order $k+1$, with the additional direction equal to the base point, implies the sign at order $k$. Iterating proves the assertion. This avoids any assumption that the lower-order thresholds are already known. Negative $\beta$ fails the first derivative; $\beta=0$ gives a constant.

\subsection{Five-vector identities and a dimension-free bound}\label{S09-sec-3}
For unit $u_i$, put
\[
 s=\sum_i u_i,\quad t=\|s\|^2,\quad M=\sum_i u_i u_i^T,\quad
 U=\sum_{\{\{i,j\},\{k,l\}\}}(u_i\cdot u_j)(u_k\cdot u_l).
\]
Expanding the definitions gives
\begin{align}
 R&=(25-t)/4,\label{S09-R}\\
 F_\beta&=\beta^2+\tfrac34\beta-\tfrac98
       +\tfrac{2\beta+3}{8}t+\tfrac14U,\label{S09-F}\\
 U&=\frac{t^2}{8}-\frac t4+\frac{\tr(M^2)}4
       -\frac{s^TMs}{2}-\frac58.\label{S09-U}
\end{align}
To check the last identity, write $g_{ij}=u_i\cdot u_j$ and $T=\sum_{i<j}g_{ij}=(t-5)/2$. The square of $T$ consists of squared edges, twice adjacent edge products, and twice the disjoint edge products $U$. Also
\[
 \sum_i\left(\sum_{j\ne i}g_{ij}\right)^2
 =s^TMs-2t+5,\qquad
 \sum_{i<j}g_{ij}^2=(\tr M^2-5)/2.
\]
Eliminating the adjacent edge products gives \eqref{S09-U}.
\begin{lemma}\label{S09-fivegram}
Five unit vectors in any Euclidean space satisfy
\begin{equation}\label{S09-gram}
 U\ge\frac{15}{16}-\frac78t-\frac1{16}t^2.
\end{equation}
At $t=0$, equality holds precisely when their Gram matrix has all off-diagonal entries $-1/4$.
\end{lemma}
\begin{proof}
Let $G=(g_{ij})$ and use the decomposition relative to $a=5^{-1/2}(1,\ldots,1)$ and $a^\perp$:
\[
 G=\begin{pmatrix}t/5&b^T\\b&C\end{pmatrix},\qquad
 c:=\tr C=5-t/5.
\]
If $\lambda$ is the largest eigenvalue of $C$, positivity of the $2\times2$ compressions, or the Schur complement, gives $\|b\|^2\le(t/5)\lambda$. The other three eigenvalues give
$\tr C^2\ge\lambda^2+(c-\lambda)^2/3$. Substitution in \eqref{S09-U} yields
\[
 U\ge\frac{7t^2}{200}-\frac t4-\frac{2t\lambda}{5}
       +\frac14\left(\lambda^2+\frac{(c-\lambda)^2}{3}\right)-\frac58.
\]
Minimizing this quadratic over all real $\lambda$ is a valid lower bound; its minimum, at $\lambda=c/4+3t/5$, is exactly \eqref{S09-gram}. At $t=0$, $b=0$ and equality requires $C=(5/4)I_4$, equivalent to $G=(5/4)(I-aa^T)$.
\end{proof}
Combine \eqref{S09-F} and \eqref{S09-gram}:
\begin{equation}\label{S09-large}
 F_\beta\ge\beta^2+\tfrac34\beta-\frac{57}{64}
       +\frac{8\beta+5}{32}t-\frac{t^2}{64}.
\end{equation}
If $0\le t\le9$, the last two terms are nonnegative for $\beta\ge0$. If $t\ge9$, then $R\le4$ by \eqref{S09-R}, and $S\ge0$, so \eqref{S09-der} gives $F_\beta\ge\beta(\beta+3)$. Hence $\beta\ge(\sqrt{66}-3)/8$ suffices in every dimension. For a regular four-simplex, $t=0$ and $U=15/16$, so \eqref{S09-large} is an equality. Such a simplex exists in spatial dimension at least four. This proves sharpness for $n\ge5$, and the equality characterization follows from Lemma~\ref{S09-fivegram}.

\subsection{Spatial dimension three: the bipyramid reduction}\label{S09-sec-4}
Suppose $u_i\in\R^3$, and put $\lambda=\lambda_{\max}(M)\in[5/3,5]$. Equations \eqref{S09-F}--\eqref{S09-U} can be rearranged as
\begin{equation}\label{S09-H}
 F_\beta=\beta^2+\tfrac34\beta-\tfrac34+\frac{H_\beta}{16},
\quad H_\beta=\tr M^2-\tfrac{17}{2}+(4\beta+5)t-2s^TMs+\tfrac12t^2.
\end{equation}
We claim $H_\beta\ge0$ for all $\beta\ge1/2$. If $\lambda\le5/2$, Theorem~\ref{S09-tumanov} and $s^TMs\le\lambda t$ give
\[
 H_\beta\ge(4\beta+3-2\lambda)t+\tfrac12t^2
 \ge(4\beta-2)t+\tfrac12t^2\ge0.
\]
If $\lambda\ge5/2$, the other two eigenvalues give
\[
 \tr M^2\ge\lambda^2+\frac{(5-\lambda)^2}{2}.
\]
Writing $A(\lambda)=\tfrac32\lambda^2-5\lambda+4$, we obtain
\[
 H_\beta\ge A(\lambda)+(7-2\lambda)t+\tfrac12t^2.
\]
For $5/2\le\lambda\le7/2$ this is at least $A(5/2)=7/8$. For $7/2\le\lambda\le5$, minimizing in $t\in\R$ gives
\[
 H_\beta\ge A(\lambda)-\tfrac12(2\lambda-7)^2
 =-\tfrac12\lambda^2+9\lambda-\tfrac{41}{2}\ge\tfrac{39}{8}>0.
\]
The claim follows. Since $b_4=(\sqrt{57}-3)/8>1/2$, equation \eqref{S09-H} proves sufficiency at and above $b_4$.

For the triangular bipyramid, $s=0$, $\tr M^2=17/2$, and $U=3/2$, so
$F_\beta=\beta^2+3\beta/4-3/4$. This is negative for $0<\beta<b_4$, proving sharpness. At $\beta=b_4$, equality in the small-$\lambda$ estimate forces $t=0$; Theorem~\ref{S09-tumanov} then forces the bipyramid. This completes the dimension-four case using exactly the stated spherical-energy input.

\subsection{Spatial dimension two: a self-contained pentagon argument}\label{S09-sec-5}
We include the planar argument to make the dimension-three entry independent of any earlier submission. Write $u_j=(\cos\theta_j,\sin\theta_j)$ and
\[
 p(z)=\prod_{j=1}^5(z-e^{i\theta_j})=\sum_{k=0}^5c_kz^k,
 \quad b_k=|c_k|^2,\quad B=\sum b_k.
\]
Unit-circle roots imply $b_k=b_{5-k}$ and $b_0=b_5=1$. Define
$f_k(w)=(1+w)^{5-k}(1-w)^k$, $g=\sum b_kf_k$, and $h=f_0+f_5$.
For $0<\beta<1/2$, put $\delta=1-2\beta$ and
\[
 \mathcal L_\delta(f)=f(0)-\delta\int_0^1
 (f(w)-f(0))w^{-1-\delta}\,dw.
\]
The integral converges for every polynomial. Fourier orthogonality and a beta integral give
\begin{equation}\label{S09-planar}
 (-1)^5D_{v_1}\cdots D_{v_5}Q_3^{-\beta}(e)
 =\frac{(2\beta)_5}{4^5}\mathcal L_\delta(g).
\end{equation}
Here are the details of this identity. For $\beta>1/2$, take a disk variable $V$ with density $(\beta-1/2)\pi^{-1}(1-\|v\|^2)^{\beta-3/2}$ and an independent gamma variable $T$ of shape $2\beta$. The moment formula
$\mathbb EV_1^{2k}=(1/2)_k/(\beta+1/2)_k$
and the binomial expansion prove
$Q_3(t,x,y)^{-\beta}=\mathbb E e^{-T(t-1-xV_1-yV_2)}$.
Differentiation gives $(2\beta)_5 2^{-5}\mathbb E\prod_j(1-u_j\cdot V)$.
For $w=\sqrt{1-r^2}$, factor
$1-r\cos\phi=(1+w)|1-r e^{i\phi}/(1+w)|^2/2$.
Angular averaging removes unequal Fourier powers and yields $g(w)/2^5$; the radial density becomes $(2\beta-1)w^{2\beta-2}$. Expanding $g$ now proves \eqref{S09-planar} by polynomial continuation in $\beta$, since each denominator $2\beta+j-1$, $1\le j\le5$, divides $(2\beta)_5$. The integral measure is used only above $1/2$; its continuation below $1/2$ is the displayed signed functional, not a positive measure.

Convexity in $k$ gives $f_k(w)+f_{5-k}(w)\le h(w)$ on $[0,1]$, strictly on $(0,1)$ for interior $k$. Thus $g\le Bh/2$ and both sides agree at zero. The functional reverses this comparison:
\[
 \mathcal L_\delta(g)\ge\tfrac B2\mathcal L_\delta(h),\qquad
 \mathcal L_\delta(h)=2\left[1-\delta\left(\frac{10}{2-\delta}+\frac5{4-\delta}\right)\right].
\]
The latter is nonnegative exactly when $\beta\ge(\sqrt{41}-3)/8$ in $(0,1/2)$. Above $1/2$ the positive gamma representation gives all-order signs; at $1/2$ its limiting derivative is positive. Equality at the positive critical exponent forces every interior $c_k$ to vanish, hence a regular pentagon. This also gives a strictly negative witness below the threshold.
Finally $Q_2=(t-x)(t+x)$ is a product of positive linear forms on its cone, so every nonnegative inverse power has all alternating signs. The one-dimensional case is immediate. Theorem~\ref{S09-main} is proved.

\subsection{Scope}\label{S09-sec-6}
The dimension-four threshold is not a numerical conjecture: its proof is reduced above to a published exact five-point energy theorem and explicit spectral inequalities. Conversely the proof is not wholly independent of that theorem. The dimension-three and dimension-at-least-five proofs need no such energy input. All statements concern the full Lorentz cone; restriction to a smaller coordinate orthant can change the sharp threshold. Neither general real-stable order-five positivity nor the all-orders elementary-symmetric-polynomial problem is settled here.




\clearpage
\section[S10 --- Quantitative rigidity and exact certificates for critical rank-two permanents]{Quantitative rigidity and exact certificates for critical rank-two permanents}\label{contrib-S10}
\markboth{Stable Polynomials, Matrix Reconstruction, and Lorentz Positivity}{S10 --- Showa}
\noindent\textsc{S10 \quad Showa programme, source round 10}\par\medskip


\begin{quote}\small\textbf{Overview.} At the sharp order-$m$ positivity threshold for real rank-two $\alpha$-permanents, we give an exact weighted sum-of-squares certificate in the coefficients of a polynomial with roots on the unit circle. The least weight is explicitly $\delta_m2^m/m$, and all weights converge uniformly to one. This yields a uniform coefficient-energy asymptotic, an explicit nondegenerate angular Hessian modulo rotations, and a quantitative inverse theorem: a small permanent deficit forces the Gram matrix close to a regular projective polygon. The distance bound has exponent $1/2$, and that exponent is optimal. All statements and their proofs are included; the result is a quantitative stability theorem, not an extension to higher-rank matrices.\end{quote}
\subsection{Threshold, coefficient energy, and main conclusions}\label{S10-sec-1}
Fix $m\ge3$. Let $\delta_m\in(0,1)$ be the unique solution of
\begin{equation}\label{S10-root}
 \delta_m\sum_{j=1}^{\lfloor m/2\rfloor}
 \frac{\binom m{2j}}{2j-\delta_m}=1,
 \qquad a_m=\frac{1-\delta_m}{2},\quad C_m=\frac{(2a_m)_m}{4^m}.
\end{equation}
Existence follows because the left side is strictly increasing from zero to a value greater than one. Here $(a)_m=\prod_{j=0}^{m-1}(a+j)$.
Every real rank-at-most-two correlation matrix can be written
\[
 A_{ij}=\cos\frac{\theta_i-\theta_j}{2}.
\]
Changing signs of its Gram vectors changes chosen angles by multiples of $2\pi$, and has no effect on the $\alpha$-permanent. Form
\[
 p(z)=\prod_{j=1}^m(z-e^{i\theta_j})=\sum_{k=0}^mc_kz^k,
 \qquad E(A)=\sum_{k=1}^{m-1}|c_k|^2.
\]
The energy $E$ is independent of the choice of planar orthonormal coordinates and the signs and order of the Gram vectors. Define $D(A)=\operatorname{per}_{a_m}(A)$, where
$\operatorname{per}_\alpha(A)=\sum_{\sigma\in S_m}\alpha^{c(\sigma)}\prod_iA_{i,\sigma(i)}$.
\begin{theorem}[Exact certificate]\label{S10-sos}
There are explicit positive numbers $d_{m,k}=d_{m,m-k}$, $1\le k<m$, such that
\begin{equation}\label{S10-SOS}
 D(A)=C_m\sum_{k=1}^{m-1}d_{m,k}|c_k|^2.
\end{equation}
They satisfy
\begin{equation}\label{S10-weightbounds}
 d_{m,1}=d_{m,m-1}=\min_{1\le k<m}d_{m,k}
 =\frac{\delta_m2^m}{m},\qquad
 d_{m,k}\le1+\frac{m\delta_m}{1-\delta_m}.
\end{equation}
Consequently $D=0$ precisely for regular projective polygons, and
\begin{equation}\label{S10-sandwich}
 C_m\frac{\delta_m2^m}{m}E(A)
 \le D(A)\le C_m\left(1+\frac{m\delta_m}{1-\delta_m}\right)E(A).
\end{equation}
Uniformly over all such matrices with $E(A)>0$,
\begin{equation}\label{S10-uniform}
 \frac{4^mD(A)}{m!E(A)}\longrightarrow1\qquad(m\to\infty).
\end{equation}
\end{theorem}
Let $\mathcal O_m$ be the finite set of matrices obtained from
$G_m=(\cos(\pi(i-j)/m))_{i,j=0}^{m-1}$ by simultaneous permutations and matching row-and-column sign switches. Distance below is Frobenius distance to this set.
\begin{theorem}[Quantitative rigidity]\label{S10-rigid}
Put $d_*:=\delta_m2^m/m$. If
\begin{equation}\label{S10-small}
 \frac{D(A)}{C_md_*}\le\frac1{64m^2(m-1)},
\end{equation}
then
\begin{equation}\label{S10-distance}
 \operatorname{dist}_F(A,\mathcal O_m)
 \le2\pi\sqrt{\frac{(m-1)D(A)}{C_md_*}}.
\end{equation}
For each fixed $m$, the exponent $1/2$ cannot be replaced by any larger exponent in a uniform local inequality of this kind.
\end{theorem}
The qualitative threshold and regular-polygon equality arise from finite-order inverse-Lorentz analysis. The present contribution is an exact coefficient certificate, its quantitative conditioning, and the inverse estimate. The all-size permanent problem and its relation to complete monotonicity are prior theory \cite{ref009,ref005}. No claim is made to resolve a new higher-rank positivity conjecture.

\subsection{Derivation of the exact coefficient identity}\label{S10-sec-2}
All roots of $p$ have modulus one, so self-inversivity implies
\[
 |c_k|^2=|c_{m-k}|^2,\qquad |c_0|^2=|c_m|^2=1.
\]
For $0<\delta<1$ and a polynomial $f$, set
\[
 \mathcal L_\delta(f)=f(0)-\delta\int_0^1
 (f(w)-f(0))w^{-1-\delta}\,dw.
\]
The integral converges. If $f(w)=\sum f_jw^j$, its value is
$f_0-\delta\sum_{j\ge1}f_j/(j-\delta)$.
Define
\[
 f_k(w)=(1+w)^{m-k}(1-w)^k,\quad h_k=f_k+f_{m-k},\quad
 g(w)=\sum_{k=0}^m|c_k|^2f_k(w).
\]
We first establish, for $0<\alpha<1/2$ and $\delta=1-2\alpha$,
\begin{equation}\label{S10-basic}
 \operatorname{per}_\alpha(A)=\frac{(2\alpha)_m}{4^m}\mathcal L_\delta(g).
\end{equation}
The details also guard against treating a signed continuation as a positive integral.

The squarefree coefficient identity
$[z_1\cdots z_m]\det(I-ZA)^{-\alpha}=\operatorname{per}_\alpha(A)$
follows by expanding $\exp(\alpha\sum_{\ell\ge1}\tr((ZA)^\ell)/\ell)$ into disjoint permutation cycles. Write $A=W^TW$ with unit columns $w_j=(\cos(\theta_j/2),\sin(\theta_j/2))$. Sylvester's identity converts the determinant to $\det(I_2-\sum z_jw_jw_j^T)$. Identifying
$\left(\begin{smallmatrix}t+x&y\\y&t-x\end{smallmatrix}\right)$ with $(t,x,y)$, its determinant is $Q=t^2-x^2-y^2$, and the rank-one matrices are the directions $v_j=\tfrac12(1,\cos\theta_j,\sin\theta_j)$. Thus the permanent is $(-1)^mD_{v_1}\cdots D_{v_m}Q^{-\alpha}(1,0,0)$.

For $\alpha>1/2$, let $V$ have disk density $(\alpha-1/2)\pi^{-1}(1-\|v\|^2)^{\alpha-3/2}$ and independently let $T$ have gamma shape $2\alpha$ and rate one. The beta moments
$\mathbb E V_1^{2j}=(1/2)_j/(\alpha+1/2)_j$
and the binomial theorem prove
\[
 Q(t,x,y)^{-\alpha}=\mathbb E e^{-T(t-1-xV_1-yV_2)}.
\]
Its mixed derivative is $(2\alpha)_m2^{-m}\mathbb E\prod_j(1-u_j\cdot V)$. At fixed radius $r$, put $w=\sqrt{1-r^2}$ and $\rho=r/(1+w)$. Factoring
$1-r\cos\phi=(1+w)|1-\rho e^{i\phi}|^2/2$
and using Fourier orthogonality, the angular average of that product equals $g(w)/2^m$. The radial variable $w$ has density $(2\alpha-1)w^{2\alpha-2}$. Hence the permanent equals
\[
 \frac{(2\alpha)_m}{4^m}(2\alpha-1)\int_0^1w^{2\alpha-2}g(w)\,dw.
\]
Expansion in powers of $w$ gives a polynomial identity in $\alpha$, since every denominator $2\alpha+j-1$ for $1\le j\le m$ divides $(2\alpha)_m$. Polynomial continuation proves \eqref{S10-basic}. The integral representation itself is used only for $\alpha>1/2$.

At $\delta=\delta_m$, equation \eqref{S10-root} says $\mathcal L_{\delta_m}(h_0)=0$. Using coefficient symmetry, set
\begin{equation}\label{S10-weights}
 d_{m,k}:=\frac12\mathcal L_{\delta_m}(h_k)
 =\frac{\delta_m}{2}\int_0^1
 \frac{h_0(w)-h_k(w)}{w^{1+\delta_m}}\,dw.
\end{equation}
The integrand is a polynomial times an integrable power, so this is an explicit finite sum of algebraic quantities determined by \eqref{S10-root}. Since $\mathcal L(g)=\tfrac12\sum|c_k|^2\mathcal L(h_k)$ and the endpoint terms vanish, \eqref{S10-SOS} follows.

\subsection{Exact least weight and uniform conditioning}\label{S10-sec-3}
For $x=1+w>y=1-w>0$, the function
$k\mapsto x^{m-k}y^k+x^ky^{m-k}$
is symmetric and strictly convex, decreasing from $k=0$ to $m/2$. Thus
\[
 h_0(w)>h_1(w)\ge h_k(w)\quad(0<w<1,\ 1\le k<m),
\]
where $h_1=h_{m-1}$ and the comparison uses symmetry when $k>m/2$. Equation \eqref{S10-weights} proves positivity and shows the least weight occurs at $k=1,m-1$.
Now
\[
 h_0-h_1=2w\bigl[(1+w)^{m-1}-(1-w)^{m-1}\bigr].
\]
Expanding odd powers and using $\binom{m-1}{2j-1}=(2j/m)\binom m{2j}$ yields
\begin{align*}
 d_{m,1}
 &=2\delta_m\sum_{j\ge1}\frac{\binom{m-1}{2j-1}}{2j-\delta_m}\\
 &=\frac{2\delta_m}{m}
   \left[\sum_{j\ge1}\binom m{2j}
       +\delta_m\sum_{j\ge1}\frac{\binom m{2j}}{2j-\delta_m}\right]
 =\frac{\delta_m2^m}{m}.
\end{align*}
For the upper bound, $h_k(w)/2\ge(1-w)^m\ge1-mw$, so
\[
 d_{m,k}=1+\delta_m\int_0^1(1-h_k(w)/2)w^{-1-\delta_m}\,dw
 \le1+\frac{m\delta_m}{1-\delta_m}.
\]
This proves \eqref{S10-weightbounds} and \eqref{S10-sandwich}. Vanishing of $D$ is equivalent to all interior coefficients being zero, hence to $p(z)=z^m-\zeta$ with $|\zeta|=1$: the roots form a regular polygon, and the original unoriented Gram lines form a regular projective polygon.

To verify the uniform asymptotic, write $S_m=\sum_{j\ge1}\binom m{2j}/(2j)$. A binomial variable $N\sim\operatorname{Bin}(m,1/2)$ gives
\[
 m2^{-m}S_m=\mathbb E[(m/N)\mathbf1_{\{N>0,\ N\text{ even}\}}]\longrightarrow1.
\]
Here the even probability is $1/2$, and the central region $|N-m/2|\le\epsilon m$ has $m/N\to2$ as $\epsilon\downarrow0$. The complementary expectation is $O_\epsilon(1/m)$, since $m/N\le m$ and the fourth central moment is $(3m^2-2m)/16$. The zero event is exponentially small. Equation \eqref{S10-root} and
\[
 S_m\le\sum_{j\ge1}\frac{\binom m{2j}}{2j-\delta_m}
 \le S_m/(1-\delta_m/2)
\]
therefore imply $\delta_m\sim m/2^m$. Both bounds in \eqref{S10-weightbounds} converge to one, uniformly in $k$. Finally
\[
 \frac{4^mC_m}{m!}=\prod_{j=1}^m\left(1-\frac{\delta_m}{j}\right)\longrightarrow1,
\]
because $\delta_m\sum_{j=1}^m1/j\to0$. The sandwich proves \eqref{S10-uniform}, with no restriction on how the matrix varies with $m$.

\subsection{From coefficients to matched roots and Gram matrices}\label{S10-sec-4}
We prove an explicit root-matching lemma.
\begin{lemma}\label{S10-roots}
Let $p(z)=z^m+c_0+\sum_{k=1}^{m-1}c_kz^k$ have all roots on the unit circle and $|c_0|=1$. Put $\eta=\sqrt{(m-1)\sum_{k=1}^{m-1}|c_k|^2}$. If $\eta\le1/(8m)$, its roots can be bijectively matched to the roots of $q(z)=z^m+c_0$ with chordal error at most $4\eta/m$.
\end{lemma}
\begin{proof}
The case $\eta=0$ is immediate. Put $\rho=4\eta/m\le1/(2m^2)$. The roots $\lambda_j$ of $q$ are separated by at least $2\sin(\pi/m)\ge4/m$. On $|z-\lambda_i|=\rho$, use $\prod_{j\ne i}|\lambda_i-\lambda_j|=|q'(\lambda_i)|=m$ and $\prod(1-a_j)\ge1-\sum a_j$ for $0\le a_j\le1$ to obtain
\[
 |q(z)|\ge m\rho\prod_{j\ne i}\left(1-\frac\rho{|\lambda_i-\lambda_j|}\right)
 \ge m\rho(1-m^2\rho/4)\ge\frac78m\rho.
\]
Cauchy--Schwarz gives
\[
 |p(z)-q(z)|\le\eta(1+\rho)^{m-1}\le\eta e^{m\rho}<2\eta.
\]
The preceding lower bound for $|q|$ is $7\eta/2>2\eta$. Rouch\'e's theorem gives exactly one root of $p$ in each of the disjoint disks, establishing the matching.
\end{proof}
By \eqref{S10-sandwich}, hypothesis \eqref{S10-small} implies the hypothesis of Lemma~\ref{S10-roots}. Since both matched roots lie on the circle, their angular distance is at most $\pi\rho/2=2\pi\eta/m$. Choosing angle representatives and signs of the original half-angle Gram vectors, each matrix entry changes by at most $2\pi\eta/m$, because cosine is one-Lipschitz. The Frobenius distance is at most $2\pi\eta$. Substituting $E\le D/(C_md_*)$ proves \eqref{S10-distance}.

\subsection{Hessian spectrum and optimality of the exponent}\label{S10-sec-5}
At the regular polygon take $\theta_j=2\pi j/m$ for $j=0,\ldots,m-1$ and perturb to $\theta_j+\epsilon_j$. Let
$\widehat\epsilon_k=\sum_{j=0}^{m-1}\epsilon_j e^{-2\pi i jk/m}$.
Differentiating $p(z)=\prod_j(z-e^{i\theta_j})$ at $p_0(z)=z^m-1$ gives, for $1\le k<m$,
\[
 c_k=-i\widehat\epsilon_k+O(\|\epsilon\|^2).
\]
Indeed $(z^m-1)/(z-e^{2\pi ij/m})$ is the geometric sum of its $m$ monomials, and coefficient extraction gives the formula. Inserting it into \eqref{S10-SOS} yields
\begin{equation}\label{S10-hessian}
 D(A(\epsilon))=C_m\sum_{k=1}^{m-1}d_{m,k}|\widehat\epsilon_k|^2
 +O(\|\epsilon\|^3).
\end{equation}
For fixed $m$, the remainder is uniform near zero because all quantities are analytic. Parseval gives $\sum_{k=1}^{m-1}|\widehat\epsilon_k|^2=m\sum_j|\epsilon_j-\overline\epsilon|^2$. Thus the real angular Hessian has one zero mode, common rotation, and its other eigenvalues are
\[
 2mC_md_{m,k},\qquad k=1,\ldots,m-1,
\]
with the usual conjugate-mode pairing. In particular its smallest positive eigenvalue is exactly $2^{m+1}C_m\delta_m$.

Perturb only one angle by $\epsilon$. Formula \eqref{S10-hessian} gives $D=\Theta(\epsilon^2)$. The Gram matrix has a nonzero first derivative, since some off-diagonal cosine has nonzero sine derivative when $m\ge3$. The target set $\mathcal O_m$ is finite, so its distance is $\Theta(|\epsilon|)$ for sufficiently small nonzero $\epsilon$. An inequality with distance bounded by a constant times $D^p$ for $p>1/2$ would imply $|\epsilon|=O(|\epsilon|^{2p})$, a contradiction. This proves optimality in Theorem~\ref{S10-rigid}.

The estimates are exact for rank at most two; they are not perturbation bounds for arbitrary matrices nearly of rank two. Nor does \eqref{S10-SOS} claim that the permanent is a sum of squares in unconstrained matrix entries. It is an explicit sum of squared coefficient moduli in the unit-circle Gram parametrization, with all weights independently specified by \eqref{S10-root} and \eqref{S10-weights}.




\clearpage
\part{Mixture Identifiability, Inverse Problems, and Minimax Recovery}


\clearpage
\section[M01 --- What a Stopped Coin Reveals: Exact Moment Rank of Sequential Mixture Experiments]{What a Stopped Coin Reveals: Exact Moment Rank of Sequential Mixture Experiments}\label{contrib-M01}
\markboth{Mixture Identifiability, Inverse Problems, and Minimax Recovery}{M01 --- Meiji}
\noindent\textsc{M01 \quad Meiji programme, source round 1}\par\medskip


\begin{quote}\small\textbf{Overview.} A latent probability is drawn once and a binary experiment is stopped according to a known finite tree. We prove that every tree of maximum depth $m$, however unbalanced, has exactly the same mixing-measure identifiability as $m$ fixed Bernoulli trials: its likelihood span is the entire polynomial space of degree at most $m$. A single deepest path yields an explicit inverse and a $2^m$ moment-stability bound. We extend the rank statement to adaptive affine sensors, prove full identifiability for almost surely terminating trees of unbounded depth, and give an exact finite-atomic threshold. Identifiability equivalence is sharply separated from statistical equivalence by an explicit pair of incomparable experiments. These are proved statements, not a claim to solve a previously named open problem.\end{quote}

\subsection{Problem and scope}\label{M01-sec-1}
Let $\mu\in\PP([0,1])$. One individual draws $p\sim\mu$ and, conditional on $p$, generates independent Bernoulli$(p)$ bits. A deterministic full binary tree tells the experimenter when to stop. Its leaves are finite bit words; every internal vertex has two children. The observed datum is the entire stopped word. The tree is known and does not depend on the unknown $p$ or $\mu$.

For a vertex $v$, let $s(v)$ and $f(v)$ count successes and failures, let $|v|=s(v)+f(v)$, and put
\[
 q_v(p)=p^{s(v)}(1-p)^{f(v)}.
\]
The probability of observing a leaf $v$ is $\int q_v\,d\mu$. The \emph{likelihood span} is the real linear span of these polynomials. Write $\PP_m$ for the vector space of real polynomials of degree at most $m$, and $M_j(\mu)=\int p^j\,d\mu$.

The classical fixed-size binomial experiment identifies precisely its first $m$ moments. Our question is whether early stopping destroys any of those moments. It need not preserve precision or information order; the assertion below concerns exact laws. The foundational mixture literature includes Teicher~\cite{ref012}; nothing here is an attribution of the stopping-tree theorem to that paper.

\subsection{The exact finite-tree theorem}\label{M01-sec-2}
\begin{theorem}[Depth equals moment rank]\label{M01-th:rank}
For every finite full binary stopping tree of maximum depth $m$, its likelihood span is $\PP_m$. Consequently, for arbitrary mixing probabilities $\mu,\nu$,
\[
 \mathcal L_\mu(\text{stopped word})=\mathcal L_\nu(\text{stopped word})
 \quad\Longleftrightarrow\quad
 M_j(\mu)=M_j(\nu)\quad(0\le j\le m).
\]
The rank is $m+1$, regardless of the number of leaves or their depths.
\end{theorem}
\begin{proof}
Every leaf polynomial has degree at most $m$. For every internal vertex $v$, all its descendant leaves partition the event of reaching $v$, so
\[
 q_v=\sum_{\ell\text{ leaf below }v}q_\ell.
\]
Choose one path $v_0,v_1,\ldots,v_m$ from the root to a deepest leaf. Each $q_{v_j}$ lies in the likelihood span, has degree exactly $j$, and has nonzero leading coefficient $(-1)^{f(v_j)}$. These $m+1$ polynomials are triangularly linearly independent and therefore form a basis of $\PP_m$. Integrating the equality of spans gives both implications.
\end{proof}

The theorem also holds when only the pair $(T,S)$ of stopping time and number of successes is recorded. All words with the same pair have the same likelihood polynomial, so grouping them multiplies that polynomial by a known positive integer and does not change its span. This statement uses the same fixed deterministic tree; arbitrary coarsening need not preserve the span.

\subsubsection{A constructive inverse and its conditioning}
Fix the deepest path in the proof, and abbreviate $q_j=q_{v_j}$. Its next bit $b_{j+1}\in\{0,1\}$ gives
\[
 p q_j=\begin{cases}q_{j+1},&b_{j+1}=1,\\q_j-q_{j+1},&b_{j+1}=0.\end{cases}\tag{M01.1}
\]
Starting from $p^0=q_0$, repeated use of (M01.1) gives coefficients $c_{rj}$ satisfying
\[
 p^r=\sum_{j=0}^r c_{rj}q_j,\qquad
 \sum_{j=0}^r|c_{rj}|\le 2^r.\tag{M01.2}
\]
Indeed, multiplying by $p$ replaces each basis term by one or two terms with coefficient absolute sum at most twice its former value. This is an exact integer-arithmetic algorithm; it requires only the selected path.

\begin{corollary}[Quantitative inverse]
If $P_\mu,P_\nu$ denote the stopped-word laws, then
\[
 |M_r(\mu)-M_r(\nu)|\le 2^r\TV(P_\mu,P_\nu),\quad 1\le r\le m,
\]
where $\TV(P,Q)=\sup_A|P(A)-Q(A)|$. More generally, estimates of all selected path-reach probabilities accurate within $\varepsilon$ produce moments accurate within $2^r\varepsilon$.
\end{corollary}
\begin{proof}
A path-reach probability is the probability of a subset of terminal observations, hence its change is at most total variation. Apply (M01.2). The root probability is exactly one and introduces no error.
\end{proof}
The bound is a uniform guarantee, not a sharp condition number for each tree. If the chosen path consists only of successes, $q_j=p^j$ and the coefficient bound improves to one.

\subsection{Adaptive calibrated sensors}\label{M01-sec-3}
At internal vertex $v$, allow the next bit to have success probability $r_v(p)=a_v+b_vp$, with $0\le r_v(p)\le1$ on $[0,1]$. Conditional on the past and $p$, this specifies the next experiment. Delete any branch whose probability is identically zero. Call an edge informative when its parent has $b_v\ne0$, and let $d$ be the maximum number of informative edges along a terminal path.

\begin{theorem}[Effective depth]
For a finite such sensor tree with complete transcript observed, the likelihood span is exactly $\PP_d$.
\end{theorem}
\begin{proof}
A path likelihood is the product of its affine success or failure factors. Noninformative edges contribute strictly positive constants after zero branches have been deleted. Thus its degree is the number of informative edges. Every leaf has degree at most $d$. On a path achieving $d$, select the root and the vertices immediately after the first, second, up to $d$th informative edge. Their likelihoods have distinct degrees $0,1,\ldots,d$ and nonzero leading coefficients. Each is the sum of its descendant terminal likelihoods. They give a basis as before.
\end{proof}
The theorem does not assume equal sensors across vertices. Its conditioning, unlike its rank, can be arbitrarily poor when some slopes approach zero. The $2^r$ bound above is specifically for ordinary Bernoulli trials and is not asserted for arbitrarily weak sensors.

\subsection{Infinite trees and compact moment determinacy}\label{M01-sec-4}
\begin{theorem}[Unbounded depth]
Consider a full binary stopping tree, possibly infinite, which terminates almost surely for every $p\in[0,1]$. Its stopped-word experiment identifies every $\mu\in\PP([0,1])$ if its finite vertices have unbounded depths. If depths are bounded by $m$, its identifiability is exactly that of the first $m$ moments.
\end{theorem}
\begin{proof}
The probability of reaching a finite vertex is still the sum of the probabilities of its descendant terminal words, since termination holds almost surely. Equal terminal laws therefore give equal integrals of every vertex polynomial. If the depths are unbounded, for each $r$ some path has vertices at depths $0,\ldots,r$. The finite triangular argument identifies $M_r$. Equality of all moments implies equality of measures: polynomials uniformly approximate every continuous function on $[0,1]$, so the two measures integrate all continuous functions equally. For completeness, uniform polynomial approximation follows directly from the Bernstein polynomials
\[
 B_nf(x)=\sum_{j=0}^n f(j/n)\binom{n}{j}x^j(1-x)^{n-j}.
\]
Uniform continuity, together with $\Pr(|\mathrm{Bin}(n,x)/n-x|>\delta)\le(4n\delta^2)^{-1}$, proves $\|B_nf-f\|_\infty\to0$. Equality of integrals of continuous functions determines a Borel probability on a compact interval, for example by continuous upper and lower approximations to interval indicators. The bounded-depth statement is Theorem~\ref{M01-th:rank}.
\end{proof}
A deterministic example has just one infinite path, with alternating bits $1,0,1,0,\ldots$; every deviation stops immediately. Its probability of reaching depth $2j$ is $[p(1-p)]^j$, and of reaching depth $2j+1$ is $p[p(1-p)]^j$. Thus it terminates almost surely even at the endpoints and has $\sup_p\E_pT\le8/3$. It nevertheless identifies every mixing measure. Identifiability alone says nothing about a useful finite-sample recovery rate.

\subsection{Exact finite-atomic threshold}\label{M01-sec-5}
\begin{theorem}
On the class of mixing measures with at most $k$ atoms, a finite tree of depth $m$ is injective if and only if $m\ge2k-1$.
\end{theorem}
\begin{proof}
If $m\ge2k-1$, write the signed difference of two such measures on its $r\le2k$ distinct support points $x_1,\ldots,x_r$. For each $i$, the Lagrange polynomial
\[
 L_i(x)=\prod_{j\ne i}\frac{x-x_j}{x_i-x_j}
\]
has degree at most $2k-1$. Equality of the identified moments forces its integral against the signed difference to be zero, which is precisely the signed mass at $x_i$.

Conversely suppose $m\le2k-2$. Put $r=m+1$, choose $r+1$ equally spaced points $x_j=j/r$, and set
\[
 \sigma=2^{1-r}\sum_{j=0}^{r}(-1)^j\binom{r}{j}\delta_{x_j}.
\]
Its positive and negative parts both have mass one and each has at most $k$ atoms. The $r$th finite difference annihilates every polynomial of degree at most $r-1=m$, so these distinct probability measures have identical observed laws. For $m=0$ this uses $r=1$ and remains valid.
\end{proof}
This threshold is classical for fixed binomial mixtures. The point here is its exact transfer to every stopping tree of the same maximum depth, not a new claim about the classical threshold.

\subsection{Equal rank is not equal statistical information}\label{M01-sec-6}
Let $A_m$ stop at the first success or at time $m$; let $B_m$ stop at the first failure or at time $m$, with $m\ge2$. Both have rank $m+1$. Neither is a parameter-independent randomization of the other.

\begin{proposition}[Blackwell incomparability]
The two parametric experiments $\{A_m(p):0<p<1\}$ and $\{B_m(p):0<p<1\}$ are incomparable under parameter-independent stochastic garbling.
\end{proposition}
\begin{proof}
For bounded stopping time $T$, the stopped log-likelihood score is
\[
 Z=\sum_{j=1}^{T}\frac{X_j-p}{p(1-p)}.
\]
The summands with indicators $\{T\ge j\}$ are orthogonal martingale differences, whence Fisher information equals $\E_pT/[p(1-p)]$. Here
\[
 \E_pT_{A_m}=\sum_{j=0}^{m-1}(1-p)^j,
 \qquad \E_pT_{B_m}=\sum_{j=0}^{m-1}p^j.
\]
The first is strictly larger for $p<1/2$ and strictly smaller for $p>1/2$. Under a parameter-independent garbling, the new score is the conditional expectation of the old score; conditional Jensen implies that Fisher information cannot increase. Either proposed dominance would contradict one of the strict inequalities.
\end{proof}
Fixed $m$ full trials dominate every bounded tree by simulating its rule and discarding unused bits. Thus the rank theorem cannot justify replacing a fixed experiment by early stopping without a separate precision analysis.

\subsection{Audit, contribution, and limitations}\label{M01-sec-7}
The proofs use only polynomial algebra, bounded-score differentiation, and compact approximation. No symbolic or numerical computation is required. The main algebraic claim has an explicit certificate: a deepest path and the leading coefficients of its $m+1$ reach polynomials. The infinite-tree result assumes almost sure termination at every parameter, including endpoints; dropping that assumption requires specifying an infinite-transcript outcome and changes the experiment.

A targeted primary-source search found the classical mixture-identifiability framework and the fixed-binomial threshold. It did not establish priority for the exact arbitrary-tree statement. Accordingly novelty is \emph{unverified}, not certified. This manuscript does not prove any of the prior tournament's ten proposed problems, does not certify those benchmark claims, and does not claim to exceed them in mathematical importance. It supplies a checkable theorem that can support a separate quantitative theory.





\clearpage
\section[M02 --- Logarithmic Recovery from Finite-Mean Coin Trees: A Sharp Minimax Theorem]{Logarithmic Recovery from Finite-Mean Coin Trees: A Sharp Minimax Theorem}\label{contrib-M02}
\markboth{Mixture Identifiability, Inverse Problems, and Minimax Recovery}{M02 --- Meiji}
\noindent\textsc{M02 \quad Meiji programme, source round 2}\par\medskip


\begin{quote}\small\textbf{Overview.} We estimate an arbitrary probability distribution on $[0,1]$ from independent individuals, each observed through a fixed deterministic binary stopping tree. For every almost surely terminating tree with unbounded depth, we prove a universal $O(1/\log n)$ expected Wasserstein recovery bound, without a lower bound on its path probabilities. If the stopping time has a uniform exponential tail, the rate is exactly $\Theta(1/\log n)$. A deterministic alternating-path tree realizes this sharp rate with uniformly bounded expected cost at most $8/3$ flips per individual. The proof includes an explicit moment estimator, a self-contained Jackson-kernel approximation lemma, and exact finite-difference testing alternatives. This is a statistical result distinct from a likelihood-rank identity; it neither asserts NPMLE optimality nor solves a named pre-existing open problem.\end{quote}

\subsection{Experiment and main result}\label{M02-sec-1}
A mixing law $\mu\in\PP([0,1])$ generates independent latent parameters $p_1,\ldots,p_n$. For each individual, conditional on its parameter, independent Bernoulli trials are run until a known deterministic full binary tree stops. The observed datum is the complete finite terminal word. The same tree is used for all individuals, and it terminates almost surely at every $p\in[0,1]$. Write $P_\mu$ for the one-individual law and
\[
 R_n(\mathcal T)=\inf_{\widehat\mu}\sup_{\mu\in\PP([0,1])}
 \E_\mu W_1(\widehat\mu,\mu).
\]
The infimum is over measurable probability-measure-valued estimators. The transportation distance is
\[
 W_1(\mu,\nu)=\inf_{\pi\in\Pi(\mu,\nu)}\int|x-y|\,d\pi(x,y).
\]
Define $q_m=\sup_{0\le p\le1}\Pr_p(T>m)$.

\begin{theorem}[Sharp logarithmic regime]\label{M02-th:main}
There is an absolute constant $C$ such that every such tree of unbounded depth satisfies
\[
 R_n(\mathcal T)\le\frac{C}{\log(en)}\qquad(n\ge1).
\]
For every tree, bounded or unbounded, and every integer $m\ge0$,
\[
 R_n(\mathcal T)\ge\frac{(1-nq_m)_+}{2(m+1)}.\tag{M02.1}
\]
Consequently, if $q_m\le A e^{-am}$ for all $m\ge0$, where $A\ge1$ and $a>0$, and the tree has unbounded depth, then
\[
 \frac{c(A,a)}{\log(en)}\le R_n(\mathcal T)
 \le\frac{C}{\log(en)}.
\]
\end{theorem}
The uniform upper bound is specific to deterministic ordinary-coin trees with their full transcripts. It is not asserted for arbitrarily rare externally randomized continuation, arbitrary sensor noise, or a coarsening which hides the stopping time.

\subsection{Three elementary lemmas}\label{M02-sec-2}
\subsubsection{Transport on the line}
\begin{lemma}\label{M02-lem:w}
For probabilities on $[0,1]$ with distribution functions $F,G$,
\[
 W_1(\mu,\nu)=\int_0^1|F(t)-G(t)|\,dt
 =\sup_{\Lip(f)\le1}\left|\int f\,d(\mu-\nu)\right|.
\]
The supremum may be restricted to functions with $f(0)=0$.
\end{lemma}
\begin{proof}
For any coupling, the identity
$|x-y|=\int_0^1|\1_{x\le t}-\1_{y\le t}|\,dt$
and Jensen give the lower bound by the CDF integral. Coupling inverse distribution functions at a common uniform random variable attains equality: their threshold events are nested for every $t$. Every one-Lipschitz function is absolutely continuous, and Fubini gives
$\int f\,d(\mu-\nu)=-\int f'(t)(F(t)-G(t))\,dt$.
Its absolute value is bounded by the same integral, attained by choosing $f'(t)=-\operatorname{sign}(F(t)-G(t))$ and $f(0)=0$.
\end{proof}

\subsubsection{Polynomial approximation, with coefficient control}
\begin{lemma}\label{M02-lem:jackson}
There is an absolute $J$ such that, for every integer $m\ge2$ and every one-Lipschitz $f:[0,1]\to\R$ with $f(0)=0$, a polynomial $Q(x)=\sum_{j=0}^m a_jx^j$ satisfies
\[
 \|f-Q\|_\infty\le J/m,\qquad
 \sum_{j=0}^m|a_j|\le3\cdot7^m.
\]
\end{lemma}
\begin{proof}
Put $g(\theta)=f((1+\cos\theta)/2)$, an even periodic function with circular Lipschitz constant at most $1/2$ and supremum norm at most one. For integer $k\ge2$ set
\[
 H_k(t)=\left(\frac{\sin(kt/2)}{\sin(t/2)}\right)^4,
 \quad D_k=\int_{-\pi}^{\pi}H_k(t)\,dt,
 \quad K_k=H_k/D_k.
\]
The removable singularities are filled continuously. It is a nonnegative even trigonometric polynomial of degree $2k-2$. For $|t|\le1/k$, the sine inequalities give $H_k(t)\ge(2k/\pi)^4$, whence $D_k\ge32k^3/\pi^4$. On $[-\pi,\pi]$,
\[
 H_k(t)\le\min\{k^4,\pi^4|t|^{-4}\}.
\]
The first bound follows from the geometric-series representation of the sine ratio. Splitting the integral at $1/k$ yields
\[
 \int_{-\pi}^{\pi}|t|H_k(t)\,dt
 \le(1+\pi^4)k^2.
\]
Therefore the convolution $h=g*K_k$, with convolution understood as integration against the probability density $K_k$, satisfies $\|h-g\|_\infty\le C/k$ and $\|h\|_\infty\le1$. It is even, so
$h(\theta)=\sum_{r=0}^{2k-2}b_r\cos(r\theta)$, with $|b_0|\le1$ and $|b_r|\le2$ for $r>0$ by the Fourier integral formula.

Choose $k=\lfloor m/2\rfloor+1$ and set $Q(x)=\sum b_rT_r(2x-1)$, where $T_r(\cos\theta)=\cos(r\theta)$. If $A_r$ is the sum of the absolute monomial coefficients of $T_r(2x-1)$, the recurrence gives
$A_{r+1}\le6A_r+A_{r-1}$, with $A_0=1,A_1=3$. Induction gives $A_r\le7^r$. Hence the coefficient sum of $Q$ is at most
$1+2\sum_{r=1}^m7^r\le3\cdot7^m$. The approximation estimate follows because $k>m/2$.
\end{proof}

\subsubsection{Exact moment-matching alternatives}
\begin{lemma}\label{M02-lem:pair}
For every $m\ge0$ there are probabilities $\mu_m,\nu_m$ supported on the grid $\{j/(m+1):0\le j\le m+1\}$ such that their moments agree through degree $m$ and
\[
 W_1(\mu_m,\nu_m)=\frac1{m+1}.
\]
\end{lemma}
\begin{proof}
Write $r=m+1$ and take the positive and negative parts of
\[
 \sigma_r=2^{1-r}\sum_{j=0}^r(-1)^j\binom{r}{j}\delta_{j/r}.
\]
Both parts have mass one. The $r$th finite difference of a polynomial of degree below $r$ vanishes, so the moments match. On the $j$th open grid interval, their CDF difference is
\[
 2^{1-r}\sum_{i=0}^j(-1)^i\binom{r}{i}
 =2^{1-r}(-1)^j\binom{r-1}{j},\quad 0\le j<r.
\]
Integrating its absolute value and using Lemma~\ref{M02-lem:w} gives exactly $1/r$.
\end{proof}
This standard finite-difference device is not claimed as a novel lower-bound construction; related moment matching appears in the fixed-size mixture literature and the supplied tournament benchmark.

\subsection{A universal estimator from one path}\label{M02-sec-3}
For each $m$, choose a finite path $v_0,\ldots,v_m$ in the unbounded tree. Let $q_j(p)=p^{s_j}(1-p)^{j-s_j}$ be its reach polynomial. Each reach event is observable from a finite terminal word: it occurs exactly when the word has that prefix. Its probability is $r_j(\mu)=\int q_j\,d\mu$ even when infinitely many terminal words lie below the prefix.

Multiplying $q_j$ by $p$ gives either $q_{j+1}$ or $q_j-q_{j+1}$ according to the next prescribed bit. Induction therefore constructs integers $c_{kj}$ with
\[
 p^k=\sum_{j=0}^k c_{kj}q_j(p),\qquad
 \sum_{j=0}^k|c_{kj}|\le2^k.\tag{M02.2}
\]
This proves the needed part of the stopping-tree rank theorem directly, without requiring the preceding round's manuscript.

For $n$ terminal words, let $\widehat r_j$ be the empirical frequency of reaching $v_j$, and define $\widehat M_k=\sum_jc_{kj}\widehat r_j$. Since each empirical frequency has variance at most $1/(4n)$,
\[
 \E|\widehat M_k-M_k(\mu)|\le\frac{2^k}{2\sqrt n},
 \qquad
 \E\Delta_m\le\frac{2^m}{\sqrt n},\quad
 \Delta_m=\max_{0\le k\le m}|\widehat M_k-M_k(\mu)|.\tag{M02.3}
\]
The second bound follows by summing the first for $k\ge1$; $\widehat M_0=M_0=1$ exactly.

Let $\widehat\mu$ approximately minimize
$\max_{k\le m}|M_k(\eta)-\widehat M_k|$ over probabilities $\eta$, to within $1/n$. A completely measurable choice is available: enumerate the probabilities with finitely many rational atoms and rational weights, a countable weakly dense family, and select the first whose objective is within $1/n$ of the infimum over that family. The infimum is a measurable real function and the first admissible index is measurable. Compactness and continuity show that this infimum agrees with the unrestricted one. As the true $\mu$ is a competitor,
\[
 \max_{k\le m}|M_k(\widehat\mu)-M_k(\mu)|\le2\Delta_m+1/n.\tag{M02.4}
\]
For any normalized one-Lipschitz $f$, Lemma~\ref{M02-lem:jackson} and (M02.4) imply
\[
 \left|\int f\,d(\widehat\mu-\mu)\right|
 \le\frac{2J}{m}+3\cdot7^m(2\Delta_m+1/n).
\]
The right side is independent of $f$. Taking the supremum and expectation gives
\[
 \sup_\mu\E W_1(\widehat\mu,\mu)
 \le\frac{2J}{m}+\frac{6\cdot14^m}{\sqrt n}
       +\frac{3\cdot7^m}{n}.\tag{M02.5}
\]
For sufficiently large $n$ choose $m=\lfloor\log n/(4\log14)\rfloor\ge2$. Then (M02.5) is $O(1/\log n)$ with an absolute constant. For smaller $n$, use $W_1\le1$ and increase the constant. This proves the universal upper bound.

The estimator is existence-level and computationally inefficient; no polynomial running-time or NPMLE assertion is made. It uses only $m=O(\log n)$ prefix events, even though terminal words may be long.

\subsection{Lower bound and a finite-cost realization}\label{M02-sec-4}
Apply Lemma~\ref{M02-lem:pair}. Every leaf of depth at most $m$ has polynomial likelihood of degree at most $m$, so $P_{\mu_m}$ and $P_{\nu_m}$ agree at each such leaf. They therefore also assign the same mass to the complementary event $T>m$. That common mass is at most $q_m$. The total variation distance of the two full laws is at most $q_m$, and the product distance is at most $nq_m$.

For completeness, if two data laws $P,Q$ have total variation $\delta$, their common measure has mass $1-\delta$. On that common part, the triangle inequality forces the sum of the two estimation losses for targets distance $d$ apart to be at least $d$. Thus the larger risk is at least $d(1-\delta)/2$. With $d=1/(m+1)$, this proves (M02.1).

Under $q_m\le Ae^{-am}$, take $m=\lceil\log(2An)/a\rceil$. Then $nq_m\le1/2$ and the lower bound is at least $1/[4(m+1)]$, proving the second part of Theorem~\ref{M02-th:main}.

\begin{corollary}[Bounded mean cost, sharp logarithmic recovery]
There is a deterministic stopping experiment with $\sup_p\E_pT\le8/3$ and $R_n=\Theta(1/\log(en))$.
\end{corollary}
\begin{proof}
Continue only along the infinite alternating word $1,0,1,0,\ldots$; stop at the first deviation. Then
\[
 \Pr_p(T>2j)=[p(1-p)]^j\le4^{-j},\quad
 \Pr_p(T>2j+1)=p[p(1-p)]^j\le4^{-j}.
\]
It terminates almost surely for every $p$, and
\[
 \E_pT=\frac{1+p}{1-p(1-p)}\le\frac83.
\]
It has unbounded depth and $q_m\le2\cdot2^{-m}$, so Theorem~\ref{M02-th:main} applies.
\end{proof}
The cost bound is deliberately simple, not the exact maximum of the displayed rational function. The risk statement fixes the protocol and counts individuals; it does not optimize an adaptive allocation under a deterministic total-flip budget.

\subsection{Relation to existing work and audit}\label{M02-sec-5}
Tian, Kong and Valiant~\cite{ref014} and Vinayak et al.~\cite{ref015} study population recovery from a fixed number of Bernoulli trials per individual. Their work supplies essential context for the classical moment-matching methodology. Our experiment is different: the unbounded stopping rule is fixed, may have uniformly finite expected cost, and has a countable output space. We do not claim a new proof of their full fixed-size minimax curve or of NPMLE optimality.

The supplied benchmark manuscript \emph{Observation and Recovery (Geometry, Mixtures, and Spectral Inversion)}, Part I, contains broader total-budget claims and the same standard finite-difference construction. Those claims are not assumed here. In particular, a total-budget design can outperform a fixed finite-mean protocol; there is no contradiction with a polynomial total-budget rate.

The main substantive claim is the universal tree upper bound together with its matching exponential-tail lower bound. The approximation and testing lemmas are written out to make the manuscript independent of external software and of preceding rounds. The literature search does not certify novelty of the exact theorem. No pre-existing open question is asserted solved, and superiority to the benchmark remains unestablished.





\clearpage
\section[M03 --- Finite Cost Does Not Control Inverse Difficulty: Random-Cap Mixtures with Prescribed Slow Rates]{Finite Cost Does Not Control Inverse Difficulty: Random-Cap Mixtures with Prescribed Slow Rates}\label{contrib-M03}
\markboth{Mixture Identifiability, Inverse Problems, and Minimax Recovery}{M03 --- Meiji}
\noindent\textsc{M03 \quad Meiji programme, source round 3}\par\medskip
\noindent\textit{Integration note.} The stretched-exponential rate is extended to the full positive exponent range in Contribution~\ref{contrib-M06}.\par\medskip


\begin{quote}\small\textbf{Overview.} A known random cap, independent of the latent success probability, is applied to each individual's Bernoulli sequence. We give explicit finite-sample bounds in terms of its survival probabilities. Caps with tails $\exp(-a m^\beta)$, $\beta\ge1$, have sharp minimax Wasserstein rate $(\log n)^{-1/\beta}$. More strongly, for every positive sequence tending to zero and every $\varepsilon>0$, one identifiable cap experiment with expected cost at most $1+\varepsilon$ has minimax risk exceeding that sequence by an unbounded factor along a subsequence. Every fixed cap with unbounded support nevertheless admits uniformly consistent recovery. The resulting quantifier separation shows why identifiability plus a bounded expected budget cannot yield a uniform statistical rate.\end{quote}

\subsection{Model, notation, and the distinction from a deterministic tree}\label{M03-sec-1}
For each individual draw $p\sim\mu\in\PP([0,1])$ and independently a positive integer cap $L$ with known law. Observe $L$ and $L$ conditionally independent Bernoulli$(p)$ bits. Different individuals are independent. Put
\[
 s_m=\Pr(L>m),\quad m\ge0,\qquad s_0=1,
 \quad R_n(s)=\inf_{\widehat\mu}\sup_\mu\E_\mu W_1(\widehat\mu,\mu).
\]
All caps below have $s_m>0$ for every $m$ and $s_m\downarrow0$. Observing $(L,S_L)$ is equivalent for inference: conditional on that pair the word is uniform over words with its success count, independently of $p$.

External cap randomization is essential. A very small known continuation probability multiplies a likelihood polynomial by a very small number. Algebraic identifiability survives this multiplication; quantitative inverse stability need not. This differs from a deterministic stopping tree where all path factors are $p$ or $1-p$ with no small external multiplier.

\subsection{A cap-dependent nonasymptotic theorem}\label{M03-sec-2}
\begin{theorem}\label{M03-th:bounds}
For every integer $m\ge2$,
\[
 \frac{(1-ns_m)_+}{2(m+1)}\le R_n(s)
 \le\frac{C}{m}+\frac{6m7^m}{\sqrt{n s_{m-1}}}
                  +\frac{3\cdot7^m}{n},\tag{M03.1}
\]
where $C$ is absolute. The lower bound also holds for $m=0,1$.
Every fixed cap with unbounded support identifies all mixing laws and satisfies $R_n(s)\to0$.
\end{theorem}

\subsubsection{Proof of the upper bound}
For $1\le j\le m$, use the unbiased statistic
\[
 Z_j=\frac{\1_{\{L\ge j,\ X_1=\cdots=X_j=1\}}}{s_{j-1}},
 \qquad \E_\mu Z_j=M_j(\mu)=\int p^j\,d\mu.
\]
Since $Z_j^2=Z_j/s_{j-1}$, $\Var Z_j\le1/s_{j-1}$. For its sample average $\widehat M_j$ and $\widehat M_0=1$, monotonicity of the survival probabilities gives
\[
 \E\Delta_m\le\sum_{j=1}^m\E|\widehat M_j-M_j|
 \le\frac{m}{\sqrt{ns_{m-1}}},\qquad
 \Delta_m=\max_{0\le j\le m}|\widehat M_j-M_j|.\tag{M03.2}
\]
If only $(L,S_L)$ is observed, replace the indicator of the first $j$ successes, when $L\ge j$, by $(S_L)_j/(L)_j$. This is its conditional expectation and has no larger variance; here $(z)_j=z(z-1)\cdots(z-j+1)$.

Choose a probability measure $\widehat\mu$ whose maximum discrepancy from the estimated moments is within $1/n$ of the minimum. A measurable implementation enumerates all finite rationally supported probabilities with rational weights, and takes the first within $1/n$ of the infimum over this countable dense family. It follows that
\[
 \max_{j\le m}|M_j(\widehat\mu)-M_j(\mu)|\le2\Delta_m+1/n.\tag{M03.3}
\]
The elementary approximation lemma in the appendix states that every one-Lipschitz $f$ with $f(0)=0$ has a degree-$m$ polynomial $Q=\sum a_jx^j$ with $\|f-Q\|_\infty\le C/(2m)$ and $\sum|a_j|\le3\cdot7^m$. Thus
\[
 W_1(\widehat\mu,\mu)\le C/m+3\cdot7^m(2\Delta_m+1/n).
\]
Here the identity $W_1=\sup_{\Lip(f)\le1}|\int f\,d(\mu-\nu)|$ follows by integrating the difference of the distribution functions; a short verification is included in the appendix. Taking expectations proves the upper bound in (M03.1).

\subsubsection{Proof of the lower bound}
Put $r=m+1$ and let $\mu_m,\nu_m$ be the positive and negative parts of
\[
 2^{1-r}\sum_{j=0}^r(-1)^j\binom{r}{j}\delta_{j/r}.
\]
Each part has mass one. The $r$th finite difference annihilates every polynomial of degree below $r$, so their first $m$ moments agree. Their CDF difference on $[j/r,(j+1)/r)$ is $2^{1-r}(-1)^j\binom{r-1}{j}$; consequently $W_1(\mu_m,\nu_m)=1/(m+1)$.

Conditional laws for $L\le m$ agree exactly. The complementary event has probability $s_m$ under both hypotheses, giving one-sample total variation at most $s_m$ and $n$-sample distance at most $ns_m$. On the common part of the two data laws, the triangle inequality forces the sum of the two estimation losses to be at least their target distance. Its mass is $1-\TV$, so the larger risk is at least half that distance times $1-\TV$. This proves the lower bound.

\subsubsection{Identification and consistency}
Equality of observed laws gives equality of $\E Z_j$ for every $j$, hence equality of all moments. Polynomials are uniformly dense in $C[0,1]$, for example by the elementary Bernstein-polynomial construction; therefore the mixing measures are equal. For consistency, fix $m$ in the upper bound, let $n\to\infty$, and then let $m\to\infty$. This proves $R_n(s)\to0$ without claiming a rate uniform in $s$.

One can also give a single explicit consistent estimator sequence. Choose the largest $m\le\log n$ such that $m\ge2$ and
\[
 \frac{m7^m}{\sqrt{ns_{m-1}}}+\frac{7^m}{n}\le1/m.
\]
If none exists, use an arbitrary probability estimate. Every fixed $m$ is eventually feasible, so the selected $m$ tends to infinity, and (M03.1) gives consistency.

\subsection{An exact continuum of severe inverse rates}\label{M03-sec-3}
\begin{theorem}[Super-exponential caps]\label{M03-th:beta}
Fix $a>0$ and $\beta\ge1$, and let $s_m=e^{-a m^\beta}$ for $m\ge0$. This is a valid cap law with finite expectation. There are positive constants depending only on $a,\beta$ such that
\[
 c_{a,\beta}(\log(en))^{-1/\beta}
 \le R_n(s)\le C_{a,\beta}(\log(en))^{-1/\beta}.
\]
\end{theorem}
\begin{proof}
A decreasing sequence $s_0=1$, $s_m\downarrow0$ specifies $\Pr(L=\ell)=s_{\ell-1}-s_\ell$. Its mean is $\sum_{m\ge0}s_m<\infty$.
For the lower bound take $m=\lceil(\log(2n)/a)^{1/\beta}\rceil$, so $ns_m\le1/2$ in (M03.1).

For the upper bound let $m=\lfloor c(\log n)^{1/\beta}\rfloor$, where $c>0$ is small enough that
\[
 a c^\beta/2\le1/8,\qquad c\log7\le1/8.
\]
For $\log n\ge1$ and $\beta\ge1$, $m\log7\le(\log n)/8$ and $a(m-1)^\beta/2\le(\log n)/8$. The second term in (M03.1) is therefore at most $6m n^{-1/4}$, and the third is at most $3n^{-7/8}$. They are dominated by a constant times $1/m$ for sufficiently large $n$. Enlarge constants for the finitely many remaining $n$.
\end{proof}
The restriction $\beta\ge1$ is intentional. For $0<\beta<1$, the lower and upper arguments here do not match, and long-binomial deconvolution can require a different analysis. No sharp claim is made in that regime.

\subsection{An arbitrarily slow yet uniformly consistent experiment}\label{M03-sec-4}
\begin{theorem}[No rate from a mean budget]\label{M03-th:slow}
Given any positive sequence $a_n\to0$ and any $\varepsilon>0$, there is one known cap law $s$, fixed for all sample sizes, such that
\begin{enumerate}[label=(\roman*)]
\item $s_m>0$ for every $m$, $\E L\le1+\varepsilon$, and the experiment identifies every mixing law;
\item $R_n(s)\to0$;
\item for a strictly increasing sequence $n_m$, $R_{n_m}(s)/a_{n_m}\to\infty$.
\end{enumerate}
\end{theorem}
\begin{proof}
Choose increasing integers $n_m$ so large that
$a_{n_m}\le1/[m^2(m+1)]$, for $m\ge1$. Let
\[
 s_0=1,\qquad
 s_m=\min\left\{\varepsilon2^{-m-1},\ 2^{-m},\
                    \frac1{2n_1},\ldots,\frac1{2n_m}\right\},\quad m\ge1.
\]
The sequence is positive, nonincreasing, tends to zero, and is at most one. It defines a cap law with
$\E L=1+\sum_{m\ge1}s_m\le1+\varepsilon/2$.
By construction $n_ms_m\le1/2$. The lower bound of Theorem~\ref{M03-th:bounds} gives
\[
 R_{n_m}(s)\ge\frac1{4(m+1)},\qquad
 \frac{R_{n_m}(s)}{a_{n_m}}\ge\frac{m^2}{4}\longrightarrow\infty.
\]
The same theorem supplies identification and uniform consistency for this fixed law.
\end{proof}
The adverse law is a distribution of \emph{observation effort}, not a noncompact or pathological mixing distribution. The parameter space remains the compact set $\PP([0,1])$ throughout.

\begin{corollary}[Failure of uniform consistency across protocols]
For each $\varepsilon>0$ and every $n$, the supremum of $R_n(s)$ over identifiable cap laws with $\E L\le1+\varepsilon$ is at least $1/8$.
\end{corollary}
\begin{proof}
For the given $n$, take $s_m=\delta2^{1-m}$ for $m\ge1$, where $0<\delta\le\min\{1/(2n),\varepsilon/2,1/2\}$. Then the mean is at most $1+\varepsilon$ and all tails are positive. Set $m=1$ in the lower bound: $R_n\ge(1-n\delta)/4\ge1/8$.
\end{proof}
This is a supremum over protocols at each $n$, whereas Theorem~\ref{M03-th:slow} constructs a single adverse protocol. Neither assertion contradicts consistency for each fixed protocol.

\subsection{Literature position and audit}\label{M03-sec-5}
The fixed-binomial population-estimation work of Tian, Kong and Valiant~\cite{ref014} and Vinayak et al.~\cite{ref015} provides context for moment matching. Slow estimation despite identifiability also appears in other mixture problems, such as mixture-proportion estimation~\cite{ref016}; that is a different parameter and experiment. We do not claim that the general phenomenon of arbitrary slowness is new. The precise cap-tail comparison and bounded-cost construction here have self-contained proofs; their publication priority has not been established.

This round changes the observation model and addresses a genuine failure of the deterministic-tree universal rate, rather than renaming that result. It gives sharp rates for a tail family, a diagonal impossibility construction, and the quantifier distinction needed to interpret the result. There is no claim about the NPMLE, optimal deterministic total-budget allocation, or an unchanged named open question.

\medskip\noindent\textit{Supplementary proof.}\par
\subsection{The approximation and transport facts used above}\label{M03-sec-6}
For probabilities with CDFs $F,G$, any coupling has expected distance at least $\int|F-G|$ by integrating threshold indicators. Coupling their quantile functions at one uniform variable attains equality. For absolutely continuous one-Lipschitz $f$, Fubini gives
$\int f\,d(\mu-\nu)=-\int f'(F-G)$; choose $f'=-\operatorname{sign}(F-G)$ to attain the CDF integral. This proves the transport dual formula.

Here are sufficient details for the polynomial bound, including its coefficient constant. For normalized one-Lipschitz $f$, put $g(\theta)=f((1+\cos\theta)/2)$. Convolve it with the probability kernel proportional on $[-\pi,\pi]$ to
$H_k(t)=(\sin(kt/2)/\sin(t/2))^4$.
This is an even nonnegative trigonometric polynomial of degree $2k-2$. On $|t|\le1/k$ it is at least $(2k/\pi)^4$, giving integral at least $32k^3/\pi^4$; everywhere it is at most $\min\{k^4,\pi^4|t|^{-4}\}$, giving first absolute moment of its unnormalized density at most $(1+\pi^4)k^2$. Consequently the normalized kernel has first moment $O(1/k)$, and the convolution approximates $g$ within $O(1/k)$ while remaining bounded by one. Its cosine coefficients have absolute value at most two. Replacing $\cos(r\theta)$ by $T_r(2x-1)$ produces an algebraic polynomial. The latter Chebyshev polynomials have coefficient sums $A_0=1,A_1=3$ and $A_{r+1}\le6A_r+A_{r-1}$, so $A_r\le7^r$. Choose $k=\lfloor m/2\rfloor+1$; the resulting degree is at most $m$, the approximation error is $O(1/m)$, and the coefficient sum is at most $1+2\sum_{r=1}^m7^r\le3\cdot7^m$.





\clearpage
\section[M04 --- Exact Robust Recovery of a Distributional Functional: Oscillation Duality and Small Atomic Certificates]{Exact Robust Recovery of a Distributional Functional: Oscillation Duality and Small Atomic Certificates}\label{contrib-M04}
\markboth{Mixture Identifiability, Inverse Problems, and Minimax Recovery}{M04 --- Meiji}
\noindent\textsc{M04 \quad Meiji programme, source round 4}\par\medskip


\begin{quote}\small\textbf{Overview.} For an unknown probability on a compact space, finitely many continuous moments are observed with arbitrary coordinatewise bounded errors. We prove an exact equality among minimax scalar recovery, half the largest compatible functional discrepancy, and a noise-penalized oscillation approximation problem. An optimal affine estimator exists even though the probability class is not centrally symmetric. Worst-case witnesses need at most $d+2$ atoms in total across two measures, and a necessary-and-sufficient primal-dual certificate specifies their contact sets and active noise coordinates. A fully solved second-moment example exhibits three exact noise regimes. The convex-analytic framework is classical; the paper makes no novelty claim for optimal-recovery duality.\end{quote}

\subsection{A deterministic inverse problem}\label{M04-sec-1}
Let $K$ be a nonempty compact metric space, $f,\phi_1,\ldots,\phi_d\in C(K)$, and $\varepsilon_i\ge0$. For an unknown $\mu\in\PP(K)$ we observe
\[
 y_i=\int\phi_i\,d\mu+e_i,\qquad |e_i|\le\varepsilon_i.
\]
The goal is to recover $F(\mu)=\int f\,d\mu$ under worst-case absolute loss. Define
\[
 \mathcal R=\inf_{A:\R^d\to\R}\ \sup_{\mu,\ |e_i|\le\varepsilon_i}
 |A(\Phi\mu+e)-F(\mu)|,
 \qquad \Phi\mu=(\textstyle\int\phi_i\,d\mu)_i.
\]
Allowing arbitrary functions $A$ only strengthens the lower bound. The estimator we produce is continuous and affine, so measurability causes no difficulty. Write $\operatorname{osc}(g)=\max_Kg-\min_Kg$.

\begin{theorem}[Three equal values]\label{M04-th:main}
The following quantities are equal, with the displayed maximum and minimum attained:
\begin{align}
 \mathcal R
 &=\frac12\max_{\substack{\mu,\nu\in\PP(K)\\
       |\int\phi_i\,d(\mu-\nu)|\le2\varepsilon_i}}
       \int f\,d(\mu-\nu)\label{M04-eq:primal}\\
 &=\min_{a\in\R^d}\left\{\frac12\operatorname{osc}
       \left(f-\sum_{i=1}^da_i\phi_i\right)
            +\sum_{i=1}^d\varepsilon_i|a_i|\right\}.\label{M04-eq:dual}
\end{align}
For any minimizing $a$, put $h=f-\sum a_i\phi_i$ and $c=(\max h+\min h)/2$. Then $A(y)=a\cdot y+c$ is minimax optimal. A pair achieving \eqref{M04-eq:primal} exists with at most $d+2$ atoms in total, counting the two measures separately.
\end{theorem}
This is not a stochastic minimax theorem. The errors may be adversarial and may depend on $\mu$. In particular, the exact affine optimality here must not be transferred unchanged to Gaussian mean-square risk. Classical optimal-recovery and statistical connections are treated by Donoho~\cite{ref017}.

\subsection{Exact duality without an infinite-dimensional duality gap}\label{M04-sec-2}
Let
\[
 H=\operatorname{conv}\{(\phi_1(x),\ldots,\phi_d(x),f(x)):x\in K\}
 \subset\R^{d+1},\qquad D=H-H.
\]
The convex hull is compact: Caratheodory's elementary elimination argument represents every point by at most $d+2$ points of the compact generating set. Continuous integration produces exactly $H$; arbitrary probability integrals lie in the closed convex hull, and finite convex combinations are probabilities. Thus $D$ is exactly the set of vectors
$(\Phi\mu-\Phi\nu,F(\mu)-F(\nu))$.

Put $B=\prod_i[-2\varepsilon_i,2\varepsilon_i]$ and
\[
 \omega=\max\{t:(u,t)\in D,\ u\in B\}.
\]
The maximum exists by compactness; $\omega\ge0$ because $D$ is symmetric and contains zero. For any $a\in\R^d$,
\[
 t\le\max_{(u',t')\in D}(t'-a\cdot u')+max_{u\in B}a\cdot u
 =\operatorname{osc}(f-a\cdot\phi)+2\sum_i\varepsilon_i|a_i|.
\tag{M04.3}
\]
This proves weak duality.

First assume every $\varepsilon_i>0$. Consider the compact convex set
\[
 G=\{(u-b,t):(u,t)\in D,\ b\in B\}.
\]
Its vertical section above the origin has maximum $\omega$. A supporting hyperplane at $(0,\omega)$ can be chosen with a strictly positive coefficient on the vertical coordinate. One direct justification is to separate $(0,\omega+\delta)$ from $G$ for $\delta>0$, normalize the normals to unit length, and take a convergent subsequence as $\delta\downarrow0$. The vertical coefficients are nonnegative. A zero limiting vertical coefficient would support the projection of $G$ at its interior point zero, impossible for a nonzero horizontal normal, since that projection contains the box $B$.

After normalizing the vertical coefficient to one, the supporting functional is $t-a\cdot(u-b)$ and has maximum $\omega$. Separating its two independent maximizations gives equality in (M04.3). Hence
\[
 \omega=\inf_a\{\operatorname{osc}(f-a\cdot\phi)+2\textstyle\sum_i\varepsilon_i|a_i|\}.
\tag{M04.4}
\]
If some error bounds vanish, replace all of them by $\varepsilon_i+\delta$. Compactness of $D$ shows that the resulting maxima $\omega_\delta$ decrease to $\omega$: a convergent subsequence of maximizers is feasible at the limit. The dual objective with the original errors is no larger than the perturbed one. Weak duality and $\omega_\delta\downarrow\omega$ therefore prove (M04.4) also in this boundary case.

\subsubsection{Attainment of the dual minimum}
For clarity, this does not follow merely from the preceding limit. Let $V$ be the finite-dimensional span of $f,1,\phi_1,\ldots,\phi_d$ modulo constant functions, with norm $\|[g]\|=\operatorname{osc}(g)/2$. In $V\times\R^d$ use the norm $\|([g],z)\|=\|[g]\|+\|z\|_1$. The right side of \eqref{M04-eq:dual} is the distance from $([f],0)$ to the linear subspace
\[
 \{([a\cdot\phi],(\varepsilon_i a_i)_i):a\in\R^d\}.
\]
Every linear subspace of a finite-dimensional normed space is closed, and a minimizing sequence can be restricted to a bounded set by the triangle inequality. Compactness then gives a nearest point, which has a preimage $a$. This proves attainment, including redundant moments and zero error coordinates.

\subsection{Why affine recovery is globally optimal}\label{M04-sec-3}
For $A(y)=a\cdot y+c$ with $c$ chosen as in Theorem~\ref{M04-th:main},
\[
 |F(\mu)-A(\Phi\mu+e)|
 =\left|\int h\,d\mu-c-a\cdot e\right|
 \le\tfrac12\operatorname{osc}(h)+\sum_i\varepsilon_i|a_i|.
\]
The minimizing affine estimator therefore has risk at most $\omega/2$.

Conversely take a maximizing compatible pair $\mu,\nu$. The data value $y=(\Phi\mu+\Phi\nu)/2$ is possible under both measures, since the two needed error vectors are $\pm(\Phi\mu-\Phi\nu)/2$ and obey the coordinate bounds. At the same $y$, any estimator makes an error of at least half $|F(\mu)-F(\nu)|=\omega/2$ on one of these measures. This proves both equalities and affine optimality.

For a particular data vector $y$, the shortest exact uncertainty interval is
\[
 [\min\{F(\mu):|\Phi\mu-y|\le\varepsilon\},\
   \max\{F(\mu):|\Phi\mu-y|\le\varepsilon\}].
\]
Its midpoint is conditionally optimal; the theorem says that a single affine estimator attains the best \emph{global} worst-case radius. It does not say these two estimators agree at every observation.

\subsection{Small atomic witnesses and complete certificates}\label{M04-sec-4}
\begin{proposition}[Total support bound]
A maximizing pair in \eqref{M04-eq:primal} can be chosen with
$|\supp\mu|+|\supp\nu|\le d+2$.
\end{proposition}
\begin{proof}
Starting from a maximizing pair, represent each of its vectors $(\Phi\mu,F(\mu))$ and $(\Phi\nu,F(\nu))$ by at most $d+2$ atoms using the finite-dimensional convex-hull representation above. If the combined number of positive weights exceeds $d+2$, the corresponding columns
\[
 (1,0,\phi(x))\quad\text{for atoms of }\mu,
 \qquad (0,1,-\phi(x))\quad\text{for atoms of }\nu
\]
in $\R^{d+2}$ are linearly dependent. Perturb weights in a nonzero dependence direction. Both signs of a sufficiently small perturbation retain nonnegativity, both total masses and all moment differences. Optimality forces the directional change in $F(\mu)-F(\nu)$ to be zero; otherwise one of the two signs improves it. Increase the perturbation until one positive weight becomes zero, leaving feasibility and optimality unchanged. Repeat until at most $d+2$ weights remain.
\end{proof}

\begin{theorem}[Primal-dual contact certificate]
Let $a\in\R^d$ and let $\mu,\nu$ satisfy the constraints of \eqref{M04-eq:primal}. Write $h=f-a\cdot\phi$. They are simultaneous primal and dual optimizers if and only if
\begin{enumerate}[label=(\roman*)]
\item $\mu$ is supported where $h$ attains its maximum, and $\nu$ where it attains its minimum;
\item for every $i$ with $a_i\ne0$,
$\int\phi_i\,d(\mu-\nu)=2\varepsilon_i\operatorname{sign}(a_i)$.
\end{enumerate}
\end{theorem}
\begin{proof}
The duality gap is the sum of the nonnegative quantities
\[
 (\max h-\textstyle\int h\,d\mu)
 +(\textstyle\int h\,d\nu-\min h)
 +\sum_i\left(2\varepsilon_i|a_i|-a_i\textstyle\int\phi_i\,d(\mu-\nu)\right).
\]
It vanishes precisely under the stated conditions. Continuity of $h$ identifies the contact supports. By strong duality, a zero gap is equivalent to simultaneous optimality.
\end{proof}
This certificate is finite once atoms and a residual polynomial are supplied. For polynomial data on an interval, verifying contact extrema reduces to checking endpoints and real zeros of the derivative, all representable by exact algebraic certificates.

\subsection{A complete three-regime example}\label{M04-sec-5}
Observe a mean $y=\int x\,d\mu+e$ on $K=[0,1]$, with $|e|\le\varepsilon$, and estimate the second moment $\int x^2\,d\mu$.
\begin{theorem}
The exact minimax radius is
\[
 \mathcal R(\varepsilon)=
 \begin{cases}
 \frac18+\varepsilon,&0\le\varepsilon\le\frac14,\\
 2\varepsilon-2\varepsilon^2,&\frac14\le\varepsilon\le\frac12,\\
 \frac12,&\varepsilon\ge\frac12.
 \end{cases}
\]
Optimal slopes are respectively $a=1$, $a=2-4\varepsilon$, and $a=0$.
\end{theorem}
\begin{proof}
For $0\le a\le1$, the residual $h(x)=x^2-ax$ has maximum $1-a$ at $1$ and minimum $-a^2/4$ at $a/2$. At $a=1$ it has equal maxima zero at $0,1$. With $a=0$, its extrema are zero and one. Rather than assume the optimal slope lies in these ranges, we certify it directly.

If $0\le\varepsilon\le1/4$, take
\[
 \mu=(1/2-2\varepsilon)\delta_0+(1/2+2\varepsilon)\delta_1,
 \qquad\nu=\delta_{1/2},\qquad a=1.
\]
The support contacts are correct, the mean difference is $2\varepsilon$, and the second-moment difference is $1/4+2\varepsilon$. The certificate gives radius $1/8+\varepsilon$.

If $1/4\le\varepsilon\le1/2$, take $a=2-4\varepsilon$, $\mu=\delta_1$, and $\nu=\delta_{1-2\varepsilon}$. These are the maximum and minimum contact points, their mean difference is $2\varepsilon$, and their second-moment difference is $4\varepsilon-4\varepsilon^2$. At $\varepsilon=1/2$, the slope is zero and the sign condition is vacuous. Finally, if $\varepsilon\ge1/2$, use $a=0$, $\mu=\delta_1$, $\nu=\delta_0$. The moment constraint is feasible and the radius is $1/2$.
\end{proof}
The optimal intercept always equals the midpoint of the residual extrema. At zero noise, it is $-1/8$, giving $A(y)=y-1/8$. This may fall outside the data-specific feasible second-moment interval near an endpoint; it is nonetheless globally minimax.

\subsection{What is and is not advanced here}\label{M04-sec-6}
The main principle belongs to classical optimal recovery and generalized moment duality. This round should not be ranked as the discovery of those subjects. Its deliverable is an explicit theorem allowing heterogeneous zero and positive noise, proof of dual attainment, a total $d+2$-atom certificate, and a fully certified nontrivial example. All arguments are finite-dimensional after taking one compact convex hull. No unresolved duality qualification, external computation, or claimed solution of a named open problem is hidden in the result.





\clearpage
\section[M05 --- Sign Variation Controls Moment Inversion: Sharp Transport Exponents and Robust Crossing Certificates]{Sign Variation Controls Moment Inversion: Sharp Transport Exponents and Robust Crossing Certificates}\label{contrib-M05}
\markboth{Mixture Identifiability, Inverse Problems, and Minimax Recovery}{M05 --- Meiji}
\noindent\textsc{M05 \quad Meiji programme, source round 5}\par\medskip


\begin{quote}\small\textbf{Overview.} For two probabilities on $[0,1]$, let their CDF difference have at most $q$ sign changes. We prove the explicit transport estimate $W_1\le(q+1)4^{q/(q+1)}\Delta_{q+1}^{1/(q+1)}$, where $\Delta_r$ is the largest discrepancy of the first $r$ power moments. The exponent is optimal for every $q$. A robust form allows a prescribed amount of CDF mass to violate the sign pattern and yields an explicit lower certificate for unavoidable oscillation. The familiar $1/(2k-1)$ exponent for two $k$-atomic measures follows as a corollary, not as a claimed new discovery. The proof uses a sign polynomial, one integration by parts, and a small-neighborhood estimate, and applies to non-atomic distributions as well.\end{quote}

\subsection{The structural variable is crossing count}\label{M05-sec-1}
For $\mu,\nu\in\PP([0,1])$, write $H=F_\mu-F_\nu$ and
\[
 W=\int_0^1|H(x)|\,dx=W_1(\mu,\nu),\qquad
 \Delta_r=\max_{1\le j\le r}\left|\int x^j\,d(\mu-\nu)(x)\right|.
\]
The transport identity follows by coupling quantile functions at a common uniform variable and integrating threshold indicators. We say that $H$ has at most $q$ sign changes if, after removing intervals where it vanishes almost everywhere, its nonzero signs can be partitioned into at most $q+1$ consecutive blocks of alternating sign. Values at isolated points are immaterial. Equivalently, there exist $0\le z_1\le\cdots\le z_q\le1$ and a sign $s\in\{-1,1\}$ such that
\[
 s\prod_{i=1}^q(x-z_i)H(x)\ge0\quad\text{for almost every }x,
\]
possibly after using fewer roots when fewer sign changes occur. In the theorem we use the actual number of roots of such a sign polynomial; extra roots are unnecessary.

\begin{theorem}[Sharp crossing exponent]\label{M05-th:main}
If $H$ has at most $q$ sign changes, then
\[
 W_1(\mu,\nu)\le(q+1)4^{q/(q+1)}\Delta_{q+1}^{1/(q+1)}.\tag{M05.1}
\]
For $q=0$, the sharper identity $W_1=|\int x\,d(\mu-\nu)|$ holds. For every integer $q\ge0$, the exponent $1/(q+1)$ in (M05.1) cannot be replaced by a larger exponent with a uniform finite constant, even for finite atomic probabilities.
\end{theorem}
A bound expressed only in terms of the number of atoms overlooks comparisons whose signed CDF crosses much less often. Conversely, many tiny oscillations are a concrete mechanism by which several small moments can coexist with nonnegligible transport distance.

\subsection{A sign-polynomial proof}\label{M05-sec-2}
\begin{lemma}[Moment integration]\label{M05-lem:ibp}
For every polynomial $Q$,
\[
 \int Q\,d(\mu-\nu)=-\int_0^1 Q'(x)H(x)\,dx.
\]
If $P(x)=s\prod_{i=1}^q(x-z_i)$ with $z_i\in[0,1]$ and $Q(x)=\int_0^xP(t)\,dt$, then the sum of the absolute coefficients of $Q$ is at most $2^q$.
\end{lemma}
\begin{proof}
Write $Q(x)=Q(1)-\int_x^1Q'(t)\,dt$ and apply Fubini. The constant term cancels because both measures have mass one, including when they have endpoint atoms. The coefficient sum of $P$ is at most $\prod_i(1+z_i)\le2^q$, and termwise integration from zero divides each coefficient by a positive integer, so cannot increase this sum.
\end{proof}

First suppose a sign polynomial with exactly $q$ roots satisfies $PH\ge0$ almost everywhere. From the lemma,
\[
 \int_0^1|P(x)||H(x)|\,dx
 =\left|\int Q\,d(\mu-\nu)\right|\le2^q\Delta_{q+1}.\tag{M05.2}
\]
Let $\delta>0$, and remove the $q$ neighborhoods $|x-z_i|<\delta$. Their union has Lebesgue measure at most $2q\delta$. Since $|H|\le1$, the removed $|H|$-mass is at most $2q\delta$. On the remainder, $|P|\ge\delta^q$. Hence
\[
 \int|P||H|\ge\delta^q(W-2q\delta).\tag{M05.3}
\]
If $W>0$, choose $\delta=W/[2(q+1)]$. Equations (M05.2)--(M05.3) give
\[
 \frac{W^{q+1}}{2^q(q+1)^{q+1}}\le2^q\Delta_{q+1},
\]
which is (M05.1). For $q=0$, take $P=\pm1$ and the integration identity gives equality directly. If $H$ has only $r<q$ changes, apply the proved inequality with $r$ and use $0\le\Delta_{q+1}\le1$. Both $r+1\le q+1$ and $4^{r/(r+1)}\le4^{q/(q+1)}$, while $\Delta_{r+1}^{1/(r+1)}\le\Delta_{q+1}^{1/(q+1)}$. This proves the asserted at-most formulation.

\subsection{Sharpness of every exponent}\label{M05-sec-3}
\begin{theorem}[Exact collapsing examples]\label{M05-th:sharp}
For each $r\ge1$ and $0<h\le1/r$, there are probabilities $\mu_h,\nu_h$ supported on $\{0,h,\ldots,rh\}$ such that
\[
 W_1(\mu_h,\nu_h)=h,\qquad
 \Delta_r=2^{1-r}r!\,h^r,
\]
and their CDF difference has exactly $r-1$ sign changes.
\end{theorem}
\begin{proof}
Take the positive and negative parts of
\[
 \sigma_h=2^{1-r}\sum_{j=0}^r(-1)^j\binom{r}{j}\delta_{jh}.
\]
They both have mass one. The finite-difference identities
\[
 \sum_{j=0}^r(-1)^j\binom{r}{j}j^\ell=0\quad(\ell<r),
 \qquad
 \left|\sum_{j=0}^r(-1)^j\binom{r}{j}j^r\right|=r!
\]
follow by applying $(z\frac{d}{dz})^\ell$ to $(1-z)^r$ at $z=1$, or by taking $r$ forward differences of $x^\ell$. They give the moment assertions. On the $j$th grid interval the signed CDF is
$2^{1-r}(-1)^j\binom{r-1}{j}$, so its sign changes at each of the $r-1$ internal grid points. Integrating its absolute value over intervals of length $h$ yields $h$.
\end{proof}
With $r=q+1$ and $h\downarrow0$, an inequality $W_1\le C\Delta_r^\alpha$ with $\alpha>1/r$ would force $h\le C' h^{r\alpha}$, impossible. This proves exponent sharpness, not sharpness of the numerical constant in (M05.1).

\subsection{Robust crossings and a certificate of hidden oscillation}\label{M05-sec-4}
Exact crossing restrictions can be unrealistic. Fix any sign polynomial $P(x)=s\prod_{i=1}^q(x-z_i)$ and define its violating transport mass
\[
 \eta(P,H)=\int_{\{x:P(x)H(x)<0\}}|H(x)|\,dx.
\]
This is an $L^1$ quantity with the same units as Wasserstein distance. It is not the number of violating points.

\begin{theorem}[Robust sign certificate]\label{M05-th:robust}
For any such $P$, with no crossing assumption on $H$,
\[
 W_1(\mu,\nu)
 \le(q+1)\left(4^q\Delta_{q+1}
                  +2^{q+1}\eta(P,H)\right)^{1/(q+1)}.\tag{M05.4}
\]
In particular, if $W_1(\mu,\nu)\ge w$ and $\Delta_{q+1}\le d$, then every sign polynomial of degree $q$ obeys
\[
 \eta(P,H)\ge2^{-q-1}\left[\left(\frac{w}{q+1}\right)^{q+1}-4^q d\right]_+.\tag{M05.5}
\]
\end{theorem}
\begin{proof}
Since $|P|\le1$ on $[0,1]$,
\[
 \int|P||H|=\int PH+2\int_{PH<0}|P||H|
 \le\left|\int PH\right|+2\eta(P,H)
 \le2^q\Delta_{q+1}+2\eta(P,H).
\]
Apply the same neighborhood bound (M05.3) and the same choice of $\delta$. This gives
$W^{q+1}\le(q+1)^{q+1}(4^q\Delta_{q+1}+2^{q+1}\eta)$.
Rearranging proves (M05.5). The proof also covers $q=0$, where it reduces to $W\le\Delta_1+2\eta$.
\end{proof}
Equation (M05.5) turns approximate moment matching into a quantitative obstruction to a low-crossing explanation. It remains valid for continuous, singular, or atomic distributions and requires no minimum atom weight or support separation.

A complementary linear regime is also immediate. If $H$ is supported, up to null sets, outside the $\rho$-neighborhoods of all roots of a correct sign polynomial, then $|P|\ge\rho^q$ wherever $H\ne0$, and
\[
 W_1(\mu,\nu)\le2^q\rho^{-q}\Delta_{q+1}.
\]
Thus the nonlinear exponent arises from CDF mass concentrating near crossings; removing that geometry restores a linear inverse estimate.

\subsection{Atomic mixtures: a known theorem recovered transparently}\label{M05-sec-5}
\begin{corollary}[Sparse moment comparison]
If $\mu,\nu$ have at most $k$ atoms each, then
\[
 W_1(\mu,\nu)\le4(2k-1)\Delta_{2k-1}^{1/(2k-1)}.
\]
The exponent is optimal, and equality of the first $2k-1$ moments implies equality of the measures.
\end{corollary}
\begin{proof}
The combined support has at most $2k$ distinct points. The signed CDF is zero outside the range of those points and is constant on at most $2k-1$ intervening intervals. It therefore has at most $2k-2$ sign changes. Theorem~\ref{M05-th:main} gives the bound. For sharpness take $r=2k-1$ in Theorem~\ref{M05-th:sharp}; each parity part then has exactly $k$ atoms.
\end{proof}
The sharp sparse exponent is established in existing moment-comparison work, notably Wu and Yang~\cite{ref018}; it is not a new result of this manuscript. The broader proof is organized around crossings rather than atom count, and the robust violating-mass certificate is recorded explicitly. The search performed here does not establish publication priority for either formulation.

\subsection{Audit and limitations}\label{M05-sec-6}
Every displayed bound is dimension one and uses compact support. The constant in (M05.1) is explicit but not optimal. The theorem is a pairwise structural statement: it does not give an algorithm to infer an unknown crossing count from noisy data, and it does not assert that a statistical estimator automatically adapts to it. The finite-difference examples certify only exponent optimality. No named open problem or benchmark claim is represented as resolved.





\clearpage
\section[M06 --- The Full Stretched-Exponential Cap Curve: Sharp Wasserstein Recovery Across a Tail Phase Transition]{The Full Stretched-Exponential Cap Curve: Sharp Wasserstein Recovery Across a Tail Phase Transition}\label{contrib-M06}
\markboth{Mixture Identifiability, Inverse Problems, and Minimax Recovery}{M06 --- Meiji}
\noindent\textsc{M06 \quad Meiji programme, source round 6}\par\medskip


\begin{quote}\small\textbf{Overview.} We determine the minimax expected Wasserstein risk for a latent Bernoulli mixing distribution when each individual is observed for an independent, known random number of trials with survival $\Pr(L>m)=\exp(-a m^\beta)$. For $0<\beta\le1$ the rate is $(\log n)^{-(1+\beta)/(2\beta)}$; for $\beta\ge1$ it is $(\log n)^{-1/\beta}$. The previously untreated subexponential regime is proved by combining a self-contained Bernstein-basis coefficient bound with local moment-matched testing alternatives. The lower bound accounts explicitly for arbitrarily long rare caps, rather than silently truncating the experiment. All estimators and approximation arguments are supplied, and the result extends to two-sided stretched-exponential tail bounds. Publication priority for this exact random-cap theorem is unverified.\end{quote}

\subsection{The rate and its two mechanisms}\label{M06-sec-1}
Let $p_1,\ldots,p_n$ be independent draws from an unknown $\mu\in\PP([0,1])$. Independently draw positive integer caps $L_i$ from a known common law. Conditional on $p_i,L_i$, observe $L_i$ independent Bernoulli$(p_i)$ trials, together with $L_i$. The full word or its success count may be used; they are statistically equivalent because the conditional ordering given the count is parameter-free.

Assume
\[
 s_m=\Pr(L>m)=e^{-a m^\beta},\qquad m\ge0,\quad a,\beta>0.
\]
This specifies $\Pr(L=\ell)=s_{\ell-1}-s_\ell$ and has finite mean for every $\beta>0$. Write
\[
 R_n=\inf_{\widehat\mu}\sup_{\mu\in\PP([0,1])}\E_\mu W_1(\widehat\mu,\mu),
 \qquad
 \gamma(\beta)=\begin{cases}(1+\beta)/(2\beta),&0<\beta\le1,\\1/\beta,&\beta\ge1.\end{cases}
\]

\begin{theorem}[Complete cap-tail curve]\label{M06-th:main}
For fixed $a,\beta>0$, constants $0<c\le C<\infty$ depending only on $a,\beta$ satisfy
\[
 c\{\log(en)\}^{-\gamma(\beta)}
 \le R_n\le C\{\log(en)\}^{-\gamma(\beta)},\qquad n\ge1.
\tag{M06.1}
\]
The same rate holds if, for all sufficiently large $m$, the known cap survival obeys
\[
 c_1e^{-a_1m^\beta}\le s_m\le c_2e^{-a_2m^\beta}
\]
for fixed positive constants, provided every finite survival probability is positive.
\end{theorem}
For $\beta>1$, the longest informative records impose a moment-degree bottleneck. For $\beta<1$, useful records are longer than $\log n$ and binomial smoothing rather than mere moment count becomes decisive. The two expressions coincide at $\beta=1$.

The proof does not assume the fixed-binomial NPMLE risk curve. That theory is important context~\cite{ref014,ref015}, but both sides of (M06.1) are derived here. No claim is made that a pre-existing named open problem has been solved.

\subsection{A quantitative Bernstein coefficient bound}\label{M06-sec-2}
Let $B_{t,j}(x)=\binom{t}{j}x^j(1-x)^{t-j}$, $0\le j\le t$.
\begin{lemma}\label{M06-lem:bernstein}
If $Q$ is a polynomial of degree at most $m<t$ and $\|Q\|_{[0,1]}\le1$, its degree-$t$ Bernstein coefficients, defined by $Q=\sum_{j=0}^t b_jB_{t,j}$, satisfy
\[
 \max_j|b_j|\le(2m+1)(t+1)
      \exp\left(\frac{m^2}{2(t-m)}\right).\tag{M06.2}
\]
For a single shifted Chebyshev polynomial $Q(x)=T_\ell(2x-1)$, $\ell<t$, the factor $2m+1$ can be omitted and $m$ replaced by $\ell$.
\end{lemma}
\begin{proof}
First take $Q=T_\ell(2x-1)$. The polynomial generating identity is
\[
 (1+z)^tQ\left(\frac z{1+z}\right)
 =\sum_{j=0}^t\binom{t}{j}b_jz^j.\tag{M06.3}
\]
For $z=\rho e^{i\theta}$, $\rho>0$, use
\[
 T_\ell\left(\frac{z-1}{z+1}\right)
 =\frac{(\sqrt z+i)^{2\ell}+(\sqrt z-i)^{2\ell}}
          {2(1+z)^\ell}.
\]
The right side is independent of the square-root branch. At any removable singularity, use continuity after multiplying as in (M06.3). Put
$u=2\sqrt\rho\,|\sin(\theta/2)|/(1+\rho)\in[0,1]$.
Then
\[
 |1+z|=(1+\rho)\sqrt{1-u^2},\qquad
 |\sqrt z\pm i|^2\le(1+\rho)(1+u).
\]
Consequently the absolute value of (M06.3)'s left side is at most
\begin{align*}
 (1+\rho)^t(1-u^2)^{(t-\ell)/2}(1+u)^\ell
 &\le(1+\rho)^t\exp\{-\tfrac12(t-\ell)u^2+\ell u\}\\
 &\le(1+\rho)^t\exp\{\ell^2/[2(t-\ell)]\}.
\end{align*}
For $1\le j\le t-1$ choose $\rho=j/(t-j)$. Cauchy's coefficient inequality in (M06.3) gives
\[
 |b_j|\le\frac{\exp\{\ell^2/[2(t-\ell)]\}}
 {\binom{t}{j}(j/t)^j(1-j/t)^{t-j}}.
\]
The denominator is the mass at a mode of a Binomial$(t,j/t)$ distribution. It is at least $1/(t+1)$, since there are $t+1$ masses summing to one. The endpoint coefficients are exactly $Q(0),Q(1)$ and have magnitude one, so the same bound holds uniformly, including $j=0,t$.

For general $Q$, expand $Q(x)=\sum_{\ell=0}^m c_\ell T_\ell(2x-1)$. The Fourier cosine coefficient formula and $\|Q\|\le1$ give $|c_0|\le1$, $|c_\ell|\le2$ for $\ell>0$. Sum the single-Chebyshev bounds and use monotonicity of $\ell^2/(t-\ell)$.
\end{proof}
The polynomial prefactor is not claimed optimal; it is harmless when $t$ is a power of $\log n$. The essential feature is exponential growth in $m^2/t$, not in $m$.

\subsection{Approximation and a general upper bound}\label{M06-sec-3}
\begin{lemma}\label{M06-lem:approx}
For every one-Lipschitz $f:[0,1]\to\R$ with $f(0)=0$ and every $m\ge2$, there is a degree-$m$ polynomial $Q$ with $\|Q\|_\infty\le1$ and $\|f-Q\|_\infty\le J/m$, where $J$ is absolute.
\end{lemma}
\begin{proof}
The even periodic function $g(\theta)=f((1+\cos\theta)/2)$ has Lipschitz constant at most $1/2$ and norm at most one. Convolve with the probability density proportional to
$H_k(t)=(\sin(kt/2)/\sin(t/2))^4$ on $[-\pi,\pi]$.
It is an even nonnegative trigonometric polynomial of degree $2k-2$. On $|t|\le1/k$, $H_k(t)\ge(2k/\pi)^4$, so its integral is at least $32k^3/\pi^4$. Everywhere $H_k(t)\le\min\{k^4,\pi^4|t|^{-4}\}$, giving $\int|t|H_k(t)\,dt\le(1+\pi^4)k^2$. The normalized kernel therefore has first absolute moment $O(1/k)$. Its convolution approximates $g$ within $O(1/k)$ and remains bounded by one. Replace each cosine term by its shifted Chebyshev polynomial, and choose $k=\lfloor m/2\rfloor+1$.
\end{proof}

Fix integers $2\le m<t$. Let $p_j(\mu)=\int B_{t,j}\,d\mu$ and $s=\Pr(L\ge t)=s_{t-1}$. Use the first $t$ bits whenever $L\ge t$ and form
\[
 \widehat p_j=\frac1{ns}\sum_{i=1}^n\1_{\{L_i\ge t,\ S_{i,t}=j\}}.
\]
If only $(L_i,S_{i,L_i})$ is stored, replace the indicator by its conditional hypergeometric probability; this is unbiased and has no larger covariance-based sum-of-variances bound. In either representation,
\[
 \E\|\widehat p-p(\mu)\|_1
 \le\sum_{j=0}^t\sqrt{p_j(\mu)/(ns)}
 \le\sqrt{(t+1)/(ns)}.\tag{M06.4}
\]
Choose a probability $\widehat\mu$ fitting $\widehat p$ in $\ell^1$ to within $1/n$ of the infimum. A measurable choice is obtained by enumerating finite rational probabilities and taking the first within $1/n$ of the countable infimum. Their density and continuity show that this infimum equals the unrestricted one. Hence
\[
 \|p(\widehat\mu)-p(\mu)\|_1
 \le2\|\widehat p-p(\mu)\|_1+1/n.
\]
The one-dimensional identity $W_1=\sup_{\Lip(f)\le1}|\int f\,d(\widehat\mu-\mu)|$ follows by integrating CDFs and using a common-uniform quantile coupling. Lemmas~\ref{M06-lem:bernstein}--\ref{M06-lem:approx} now yield
\[
 R_n\le\frac{2J}{m}+D_{m,t}\left(2\sqrt{\frac{t+1}{ns_{t-1}}}+\frac1n\right),
 \quad D_{m,t}=(2m+1)(t+1)e^{m^2/[2(t-m)]}.\tag{M06.5}
\]
This bound is valid for every independent cap law, not only stretched exponentials.

\subsubsection{The missing regime: $0<\beta<1$}
For sufficiently large $n$, choose
\[
 t=\left\lfloor\left(\frac{\log n}{4a}\right)^{1/\beta}\right\rfloor,
 \qquad m=\left\lfloor\frac14\sqrt{t\log n}\right\rfloor.
\]
Because $t/\log n\to\infty$, eventually $2\le m\le t/2$. Also $s_{t-1}\ge n^{-1/4}$ and
\[
 \exp\{m^2/[2(t-m)]\}\le\exp(m^2/t)\le n^{1/16}.
\]
The nonapproximation terms in (M06.5) are bounded by a fixed power of $\log n$ times $n^{-5/16}$, and hence are $o(1/m)$. As
$m\asymp(\log n)^{(1+\beta)/(2\beta)}$, this proves the desired upper bound.

\subsubsection{The regime $\beta\ge1$}
A simpler moment estimator suffices. For $j\le m$, the variable
$Z_j=\1_{\{L\ge j,X_1=\cdots=X_j=1\}}/s_{j-1}$ is unbiased for $\int p^j\,d\mu$ and has variance at most $1/s_{j-1}$. Fit a probability to the empirical first $m$ moments in maximum norm. The expected maximum moment error is at most $m/\sqrt{ns_{m-1}}$, and fitting at most doubles it, apart from $1/n$.

The polynomial in Lemma~\ref{M06-lem:approx} has monomial coefficient sum at most $3\cdot7^m$: its cosine coefficients are bounded by two, while $T_\ell(2x-1)$ has coefficient sum at most $7^\ell$, by the recurrence bound $A_{\ell+1}\le6A_\ell+A_{\ell-1}$, $A_0=1,A_1=3$. Thus
\[
 R_n\le\frac{2J}{m}+\frac{6m7^m}{\sqrt{ns_{m-1}}}+\frac{3\cdot7^m}{n}.
\tag{M06.6}
\]
Take $m=\lfloor c(\log n)^{1/\beta}\rfloor$, with $c$ small enough that $ac^\beta/2\le1/8$ and $c\log7\le1/8$. Equation (M06.6) is $O((\log n)^{-1/\beta})$. This includes the transition $\beta=1$.

\subsection{Local testing alternatives that survive the long records}\label{M06-sec-4}
\begin{lemma}[Local moment cancellation]\label{M06-lem:local}
Let two probabilities be supported on $[1/2-h,1/2+h]$, with $h\le1/2$, and agree in moments through degree $m$. For $\ell$ Bernoulli observations, their word laws $P_\ell,Q_\ell$ satisfy
\[
 \TV(P_\ell,Q_\ell)\le
 \left[\sum_{k=m+1}^{\ell}\binom{\ell}{k}(4h^2)^k\right]^{1/2},\tag{M06.7}
\]
with an empty sum interpreted as zero.
\end{lemma}
\begin{proof}
Under the fair-coin word law $U_\ell$, write bits as independent Rademacher variables $Y_i\in\{-1,1\}$. For $p=1/2+u$, the likelihood ratio is $\prod_{i=1}^{\ell}(1+2uY_i)$. Expand the product over subsets and integrate over the two mixing laws. The differences of $u$-moments vanish through degree $m$ and have magnitude at most $2h^k$ at degree $k$. Distinct subset products are orthogonal in $L^2(U_\ell)$, so
\[
 \int\left(\frac{dP_\ell}{dU_\ell}-\frac{dQ_\ell}{dU_\ell}\right)^2dU_\ell
 \le4\sum_{k=m+1}^{\ell}\binom{\ell}{k}(4h^2)^k.
\]
Cauchy--Schwarz and the half-$L^1$ formula for total variation prove (M06.7). This compares both laws to a positive reference law; it does not assert a relative chi-squared bound with one mixture in the denominator.
\end{proof}

For integers $r=m+1$ and $h>0$, take positive and negative parts of
\[
 2^{1-r}\sum_{j=0}^r(-1)^j\binom{r}{j}
     \delta_{1/2-h+2hj/r}.
\tag{M06.8}
\]
They are probabilities, match moments through $m$, and have exactly
\[
 W_1=\frac{2h}{m+1}.\tag{M06.9}
\]
The finite-difference identities prove matching. The cumulative alternating binomial sum is $(-1)^j\binom{r-1}{j}$, and integrating it on the grid proves (M06.9).

\subsubsection{Lower bound for $0<\beta<1$}
Set
\[
 t=\left\lceil\left(\frac{\log(8n)}a\right)^{1/\beta}\right\rceil,
 \qquad m=\left\lceil\frac{\log(16n)}{\log4}\right\rceil,
 \qquad h^2=\frac{m+1}{64e\,t}.
\]
For sufficiently large $n$, $h\le1/2$ because $m/t\to0$. For every $\ell\le t$, use $\binom{\ell}{k}\le(e t/k)^k$ in (M06.7). Since $4eth^2/(m+1)=1/16$,
\[
 \TV(P_\ell,Q_\ell)
 \le\left(\sum_{k=m+1}^{\infty}16^{-k}\right)^{1/2}
 \le2\cdot4^{-m-1}.
\tag{M06.10}
\]
In the full cap experiment, condition on $L$. For $L\le t$, (M06.10) applies; for $L>t$, use only the trivial bound one. Thus the full one-individual distance is at most
\[
 2\cdot4^{-m-1}+s_t\le\frac1{32n}+\frac1{8n}<\frac1{2n}.
\]
This step explicitly controls all longer records, however informative they might be. Product total variation for $n$ individuals is at most one half. The two-point triangle-inequality testing bound therefore gives risk at least $W_1/4$. By (M06.9),
\[
 R_n\ge\frac{h}{2(m+1)}\asymp\frac1{\sqrt{mt}}
 \asymp(\log n)^{-(1+\beta)/(2\beta)}.
\]
The numerical constants are deliberately conservative.

\subsubsection{Lower bound for $\beta\ge1$}
Use (M06.8) on the whole interval $[0,1]$, that is $h=1/2$, with
$m=\lceil(\log(2n)/a)^{1/\beta}\rceil$.
For $L\le m$, the observation laws agree exactly by moment matching. All other outcomes have total probability $s_m\le1/(2n)$. Hence their product total variation is at most one half, and the target distance $1/(m+1)$ gives $R_n\ge1/[4(m+1)]$. This proves the remaining lower bound.

\subsection{Robustness, reproduction, and status}\label{M06-sec-5}
The extension to two-sided tail bounds follows by using the lower survival bound in the upper estimates and the upper survival bound in the lower estimates, with slightly changed constants in the choices of $t,m$. Finitely many exceptional tail values are positive and affect only constants and the threshold beyond which the asymptotic choices work. Bounds for all smaller $n$ follow by enlarging the upper constant and using any fixed distinct mixing pair for a positive lower constant on that finite range.

All analytic estimates are contained in this single source. A finite sanity check of the single-Chebyshev coefficient bound for $2\le t\le30$ and every $0\le\ell<t$ passed; this is not used as proof. Coefficients can be checked exactly from
\[
 b_j=\sum_{k=0}^{\min(j,m)}a_k\frac{\binom{j}{k}}{\binom{t}{k}}
 \quad\text{when }Q(x)=\sum_{k=0}^ma_kx^k.
\]
No external file or numerical package is needed to verify the theorem.

The substantive new research step relative to the preceding cap study is the matched $0<\beta<1$ regime and its phase transition. The coefficient technology and local moment matching are related to established fixed-binomial population estimation~\cite{ref014,ref015}; their general principles are not claimed new. A targeted search did not establish the status or priority of the exact independent-cap minimax curve. It is a theorem for a specified model, not a certified resolution of an unchanged pre-existing open question, and it is not asserted to exceed the benchmark without independent assessment.





\clearpage
\section[M07 --- Polynomial Observation Tails: Multilevel Recovery and Sharp Noncritical Minimax Rates]{Polynomial Observation Tails: Multilevel Recovery and Sharp Noncritical Minimax Rates}\label{contrib-M07}
\markboth{Mixture Identifiability, Inverse Problems, and Minimax Recovery}{M07 --- Meiji}
\noindent\textsc{M07 \quad Meiji programme, source round 7}\par\medskip
\noindent\textit{Integration note.} The critical logarithmic gap mentioned in this contribution is resolved in Contributions~\ref{contrib-M09} and~\ref{contrib-M10}.\par\medskip


\begin{quote}\small\textbf{Overview.} When the independent number of Bernoulli trials per individual has polynomial survival $\Pr(L\ge t)=t^{-\alpha}$, exploiting only one record length loses information. We construct a coupled multilevel CDF estimator using every available dyadic prefix and prove a covariance-aware risk bound. Its rate is $n^{-1/(2\alpha+1)}$ for $\alpha>1/2$ and $n^{-1/2}$ for $0<\alpha<1/2$, and both are minimax optimal. The proof uses local moment-matched hypercubes with a shared atom to control Hellinger distance under the entire cap law. At the critical exponent $\alpha=1/2$ we obtain an upper bound $\sqrt{\log(en)/n}$ and a lower bound $n^{-1/2}$; the logarithmic gap is explicitly left unresolved. In particular, all finite-mean polynomial caps ($\alpha>1$) are covered sharply.\end{quote}

\subsection{Model and theorem}\label{M07-sec-1}
Independently for each of $n$ individuals, draw $p\sim\mu\in\PP([0,1])$ and a known cap $L$ with
\[
 \Pr(L\ge t)=t^{-\alpha},\qquad t=1,2,\ldots,\quad\alpha>0.
\]
Observe $L$ and conditionally independent Bernoulli$(p)$ bits up to $L$. Let $R_n(\alpha)$ denote minimax expected $W_1$ loss over all mixing probabilities. The cap has finite expectation exactly when $\alpha>1$.

\begin{theorem}\label{M07-th:main}
For constants depending only on $\alpha$,
\[
 R_n(\alpha)\asymp
 \begin{cases}
 n^{-1/2},&0<\alpha<1/2,\\
 n^{-1/(2\alpha+1)},&\alpha>1/2.
 \end{cases}\tag{M07.1}
\]
At $\alpha=1/2$,
\[
 c n^{-1/2}\le R_n(1/2)\le C\sqrt{\log(en)/n}.\tag{M07.2}
\]
The rates in (M07.1) also hold under two-sided constant-factor bounds on the stated cap survival, with positive survival at every finite depth.
\end{theorem}
The proof of the critical lower bound does not match the upper bound. No assertion is made that its square-root logarithm is necessary. The theorem fixes an observation law and counts individuals, not a deterministic total-flip budget.

\subsection{A multilevel CDF estimator with shared samples}\label{M07-sec-2}
Let $r_\ell=2^\ell$, $q_\ell=\Pr(L\ge r_\ell)=2^{-\alpha\ell}$, and let $U_\ell$ be the sample proportion in the first $r_\ell$ bits of one individual's underlying infinite sequence. Conditional on $p$,
\[
 \E|U_\ell-p|\le\frac1{2\sqrt{r_\ell}},\qquad
 \E|U_\ell-U_{\ell-1}|\le C2^{-\ell/2}.\tag{M07.3}
\]
For $x\in[0,1]$ put
\[
 D_0(x)=\1_{\{U_0\le x\}},\qquad
 D_\ell(x)=\1_{\{U_\ell\le x\}}-\1_{\{U_{\ell-1}\le x\}},\quad\ell\ge1.
\]
Although the underlying infinite sequence is a mathematical construction, the following summands are observable: multiply $D_\ell$ by $\1_{L\ge r_\ell}$, and take it to be zero when the prefix is unavailable. For a cutoff $J$, define
\[
 Y_J(x)=D_0(x)+\sum_{\ell=1}^J\frac{\1_{\{L\ge r_\ell\}}}{q_\ell}D_\ell(x),
 \qquad H_J(x)=\frac1n\sum_{i=1}^nY_{J,i}(x).
\tag{M07.4}
\]
Independence of cap and bits gives the telescoping expectation
$\E H_J(x)=\Pr(U_J\le x)$. The bias in integrated CDF distance is at most $C2^{-J/2}$ by (M07.3) and the coupling $(U_J,p)$.

\subsubsection{The covariance estimate}
The shared samples in (M07.4) are not independent across levels. Keeping their dependence improves the critical bound.
\begin{lemma}\label{M07-lem:var}
Uniformly in the mixing law,
\[
 \int_0^1\Var(Y_J(x))\,dx
 \le C\left(1+\sum_{\ell=1}^J2^{(\alpha-1/2)\ell}\right).
\tag{M07.5}
\]
\end{lemma}
\begin{proof}
For $k\ge\ell\ge1$, cap independence and nested survival events give
\[
 \E\int\left|\frac{\1_{L\ge r_\ell}}{q_\ell}D_\ell
                 \frac{\1_{L\ge r_k}}{q_k}D_k\right|
 =q_\ell^{-1}\E\int|D_\ell D_k|
 \le q_\ell^{-1}\E\int|D_k|
 \le C2^{\alpha\ell-k/2}.
\]
The last equality uses $\int|\1_{a\le x}-\1_{b\le x}|\,dx=|a-b|$ and (M07.3). Summing over $k\ge\ell$ gives at most $C2^{(\alpha-1/2)\ell}$. The base diagonal is at most one, and its cross term with level $k$ is at most $C2^{-k/2}$ because $|D_0|\le1$ and the survival factor cancels in expectation. Expand $Y_J^2$, bound absolute cross terms, and use $\Var Y_J\le\E Y_J^2$.
\end{proof}
Jensen and Cauchy--Schwarz now imply
\[
 \E\|H_J-\E H_J\|_{L^1[0,1]}
 \le\frac C{\sqrt n}\left(1+\sum_{\ell=1}^J2^{(\alpha-1/2)\ell}\right)^{1/2}.
\tag{M07.6}
\]
This does not replace the dependent sample contributions by an unjustified independence assumption.

\subsubsection{Projection to a probability and optimization}
The function $H_J$ need not be monotone. Choose a probability $\widehat\mu_J$ whose CDF is within $1/n$ of the best $L^1$ approximation to $H_J$ by a probability CDF. A measurable construction enumerates finite rational probabilities and chooses the first within $1/n$ of the countable infimum. Density in $W_1$ makes this infimum equal to the unrestricted one. Since the true CDF is a candidate,
\[
 \E W_1(\widehat\mu_J,\mu)
 \le C\left[2^{-J/2}
 +n^{-1/2}\left(1+\sum_{\ell=1}^J2^{(\alpha-1/2)\ell}\right)^{1/2}\right]+1/n.
\tag{M07.7}
\]
For $\alpha>1/2$, choose $2^J$ within a factor two of $n^{1/(\alpha+1/2)}$. For $\alpha<1/2$, take $2^J$ within a factor two of $n$; the sum is bounded independently of $J$. The same choice at $\alpha=1/2$ gives $J=O(\log(en))$. This proves all upper bounds.

If only total successes and cap are stored, condition the entire random function $Y_J$ on those statistics. Conditional expectation is computable from hypergeometric laws, leaves the mean unchanged, and decreases the integrated variance. Thus the same bound holds without recording bit order.

\subsection{A local Hellinger lemma for the full random cap}\label{M07-sec-3}
Use the convention $h^2(P,Q)=\int(\sqrt{dP}-\sqrt{dQ})^2$, so $\TV(P,Q)\le h(P,Q)$ and $h^2(P^{\otimes n},Q^{\otimes n})\le n h^2(P,Q)$.

Fix an integer $r>\alpha$. Let $\rho_+,\rho_-$ be the positive and negative parts of
\[
 2^{1-r}\sum_{j=0}^r(-1)^j\binom{r}{j}\delta_{-1+2j/r}.
\]
They match moments through degree $r-1$ and have $W_1(\rho_+,\rho_-)=2/r$, by finite differences and cumulative alternating binomial sums. Put
$\tau_\pm=\tfrac12\delta_0+\tfrac12\rho_\pm$.
Then $W_1(\tau_+,\tau_-)=1/r$, and both have a shared half-mass at zero.

For $x_0\in[1/4,3/4]$ and $0<h\le1/8$, let $\tau_{x_0,h,\pm}$ be their affine images under $z\mapsto x_0+hz$. Write $P_+,P_-$ for the complete one-individual cap laws generated by these two mixing distributions.
\begin{lemma}\label{M07-lem:hell}
There is $C_{\alpha,r}$, uniform in $x_0,h$, such that
\[
 h^2(P_+,P_-)\le C_{\alpha,r}h^{2\alpha}.
\tag{M07.8}
\]
\end{lemma}
\begin{proof}
First condition on $L=\ell$ and use the Bernoulli$(x_0)$ word law $U_\ell$ as reference. For one bit, the likelihood ratio at $x_0+u$ is $1+uZ$, where $\E Z=0$ and $\E Z^2=1/[x_0(1-x_0)]\le16/3=:v$. Expanding the product over subsets of independent bits, orthogonality gives
\[
 \int\left(\frac{dP_{+,\ell}}{dU_\ell}-\frac{dP_{-,\ell}}{dU_\ell}\right)^2dU_\ell
 \le\sum_{k=r}^{\ell}\binom{\ell}{k}(v h^2)^k.
\tag{M07.9}
\]
The moments below $r$ cancel; the difference of the $k$th centered moments is at most $h^k$ because of the common half-mass construction. If $\ell h^2\le1$, the sum is at most $C_r(\ell h^2)^r$, using $\binom{\ell}{k}\le\ell^k/k!$.

Each conditional mixture density is at least $\frac12dU_\ell$. Hence
$(\sqrt p-\sqrt q)^2\le(p-q)^2/(p+q)$ bounds its Hellinger distance by (M07.9). For larger $\ell$ use $h^2\le2$. Therefore the conditional distance is at most $C_r\min\{(\ell h^2)^r,1\}$.

Conditioning on a parameter-independent observed cap averages these squared Hellinger distances. For $H=h^{-2}\ge1$, the layer-cake identity gives
\[
 \E\min\{(L/H)^r,1\}
 =rH^{-r}\int_0^H t^{r-1}\Pr(L>t)\,dt
 \le C_{\alpha,r}H^{-\alpha},
\]
since $r>\alpha$ and $\Pr(L>t)\le C(1+t)^{-\alpha}$. This proves (M07.8), including all long caps.
\end{proof}
The shared atom is important: it supplies a valid density denominator. An $L^2$ difference relative to a reference law alone would not justify the same Hellinger bound without that domination.

\subsection{A hypercube lower bound}\label{M07-sec-4}
Let $h>0$ be small and $K=\lfloor1/(8h)\rfloor$. Choose centers
$x_j=1/4+(4j-2)h$, $1\le j\le K$. The intervals $[x_j-h,x_j+h]$ are disjoint with gaps, and lie in $[1/4,3/4]$. For $\theta\in\{-,+\}^K$, define
\[
 \mu_\theta=\frac1K\sum_{j=1}^K\tau_{x_j,h,\theta_j}.
\]
Each cell carries the same mass under every hypothesis. Its signed CDF vanishes outside that cell, so the exact distance is
\[
 W_1(\mu_\theta,\mu_{\theta'})=\frac{h}{rK}\,d_H(\theta,\theta').
\tag{M07.10}
\]
For neighboring vertices, reveal the cell label in addition to the actual data. This can only increase information. Its one-individual squared Hellinger distance is exactly $1/K$ times the differing cell's distance, at most $C h^{2\alpha}/K$. By data processing, the actual data distance is no larger. Thus neighbor $n$-sample total variation is at most
\[
 \sqrt{Cnh^{2\alpha}/K}\le\sqrt{C'nh^{2\alpha+1}}.
\]
Choose $h=c_\alpha n^{-1/(2\alpha+1)}$ with a sufficiently small positive constant. For large $n$, $K\ge1/(16h)$ and every neighbor distance is at most one half.

For any estimator, select a nearest $\mu_{\widehat\theta}$ in the finite hypercube. The triangle inequality and (M07.10) give
\[
 W_1(\widehat\mu,\mu_\theta)\ge\frac{h}{2rK}d_H(\widehat\theta,\theta).
\]
Under a uniform prior on cube vertices, testing each bit against its paired neighbor makes an error with probability at least $(1-\TV)/2\ge1/4$. Summing the $K$ bit errors yields worst-case risk at least $h/(8r)$. Thus
\[
 R_n(\alpha)\ge c_\alpha n^{-1/(2\alpha+1)}.\tag{M07.11}
\]
For $\alpha>1/2$ this matches (M07.7).

A second lower bound is always $c/\sqrt n$. Restrict to
$\mu_\vartheta=\vartheta\delta_1+(1-\vartheta)\delta_0$.
The first bit reveals the latent endpoint exactly; subsequent bits and the independent cap give no further information about $\vartheta$. Taking $\vartheta=1/2\pm c/\sqrt n$, the elementary product chi-squared bound for Bernoulli trials makes their $n$-sample total variation at most one half for a sufficiently small fixed $c$. Their $W_1$ distance is $2c/\sqrt n$, and the two-point bound gives the claimed root-$n$ lower bound. It matches (M07.7) for $\alpha<1/2$ and supplies (M07.2) at the boundary.

\subsection{Interpretation and remaining boundary}\label{M07-sec-5}
For $\alpha>1$, expected trials per individual are finite, and the complete rate is $n^{-1/(2\alpha+1)}$. The gain over using a single truncation level comes from telescoping fine-scale corrections against coarser estimates from the same latent individual. The estimator does not require smoothness or a finite number of atoms.

The multilevel principle is established methodology~\cite{ref019,ref020}; a related telescoping-CDF idea also appears in the supplied benchmark's total-budget discussion. This manuscript's work is the fixed random-cap weighting, explicit shared-level covariance bound, and matching cap-dependent hypercube lower bound. No priority claim for multilevel CDF estimation is made. Exact novelty of this cap-model minimax theorem remains unverified.

The critical case $\alpha=1/2$ is a real unresolved point of this manuscript, not hidden in asymptotic notation. The proof gives $n^{-1/2}$ versus $\sqrt{\log(en)/n}$, and resolving whether the logarithm is intrinsic requires a sharper multiscale lower bound or a better estimator. The general two-sided survival extension follows by replacing $q_\ell$ by its known values in (M07.4) and using constant-factor bounds in (M07.5), (M07.8). All proofs are unconditional under the stated assumptions, but the critical rate is not claimed settled.





\clearpage
\section[M08 --- The Information Spectrum of Random-Cap Binomial Mixtures: Exact Diagonalization and Sharp Tail Asymptotics]{The Information Spectrum of Random-Cap Binomial Mixtures: Exact Diagonalization and Sharp Tail Asymptotics}\label{contrib-M08}
\markboth{Mixture Identifiability, Inverse Problems, and Minimax Recovery}{M08 --- Meiji}
\noindent\textsc{M08 \quad Meiji programme, source round 8}\par\medskip
\noindent\textit{Integration note.} The diagonalization is classical; the random-cap transform and asymptotics are the objects of this contribution. The related heterogeneous-trial framework of Lee and Kennedy is discussed in the integrated introduction.\par\medskip


\begin{quote}\small\textbf{Overview.} At the uniform mixing distribution, the observation operator for binomial mixtures with a known independent random cap is diagonal in shifted Legendre polynomials. The fixed-cap diagonalization is classical Bernstein--Durrmeyer theory and is explicitly credited. We derive an exact survival-transform formula for the random-cap information eigenvalues and sharp asymptotic constants: polynomial survival $c t^{-\alpha}$ produces $c\Gamma(\alpha+1)k^{-2\alpha}$, while stretched-exponential survival has three distinct asymptotic regimes, including an explicit geometric-tail constant. The spectral identity is exact for chi-squared distance from the uniform reference, not merely a formal linearization. These results explain the observation-tail phase transition but do not by themselves resolve the critical Wasserstein minimax logarithm.\end{quote}

\subsection{The operator and its classical component}\label{M08-sec-1}
Let $w_t=\Pr(L=t)$, $t\ge1$, and $q_t=\Pr(L\ge t)$. Given latent $p\in[0,1]$ and cap $t$, observe a binomial count $J\sim\mathrm{Bin}(t,p)$ together with $t$. The cap is independent of $p$. Under the uniform mixing density $1$, the reference observation law is
\[
 Q_0(t,j)=\frac{w_t}{t+1},\qquad 0\le j\le t,
\]
since $\int_0^1\binom{t}{j}x^j(1-x)^{t-j}\,dx=1/(t+1)$.

For $f\in L^2[0,1]$, define the score operator
\[
 (Sf)(t,j)=(t+1)\int_0^1 B_{t,j}(x)f(x)\,dx,
 \qquad B_{t,j}(x)=\binom{t}{j}x^j(1-x)^{t-j}.
\]
It is the conditional expectation of $f(p)$ under the uniform-reference joint law, hence a contraction from $L^2[0,1]$ to $L^2(Q_0)$. Let $P_k$ be the orthonormal shifted Legendre polynomials, with $P_0=1$ and positive leading coefficients. Equivalently they are obtained by Gram--Schmidt from $1,x,x^2,\ldots$ in $L^2[0,1]$.

The fixed-$t$ operator
\[
 (K_tf)(x)=(t+1)\sum_{j=0}^tB_{t,j}(x)\int_0^1B_{t,j}(y)f(y)\,dy
\]
is the classical Bernstein--Durrmeyer operator. Its Legendre eigenstructure goes back to the established approximation literature, notably Derriennic~\cite{ref021}; the analogous conjugate-model operator perspective appears in Diaconis, Khare and Saloff-Coste~\cite{ref022}. We reproduce the short proof to isolate what randomization of $t$ changes.

\subsection{Exact information eigenvalues}\label{M08-sec-2}
\begin{theorem}\label{M08-th:diag}
The operator $S^*S$ is diagonal in the basis $P_k$ with eigenvalues
\[
 \Lambda_k=\sum_{t\ge k}w_t\lambda_{t,k},\qquad
 \lambda_{t,k}=\frac{t!(t+1)!}{(t-k)!(t+k+1)!},\quad 0\le k\le t,
\tag{M08.1}
\]
and $\lambda_{t,k}=0$ for $k>t$. In particular $\Lambda_0=1$. For $k\ge1$,
\[
 \boxed{\quad\Lambda_k=k(k+1)\sum_{t\ge k}
              \frac{q_t\lambda_{t,k}}{t(t+1)}.\quad}\tag{M08.2}
\]
For $f=\sum_{k\ge0}a_kP_k$,
\[
 \|Sf\|_{L^2(Q_0)}^2=\sum_{k\ge0}\Lambda_k a_k^2.\tag{M08.3}
\]
If $g=1+f\ge0$ and $\int f=0$, then the induced observation law $Q_g$ satisfies exactly
\[
 \chi^2(Q_g\Vert Q_0)=\sum_{k\ge1}\Lambda_k a_k^2.
\tag{M08.4}
\]
The singular values of $S$ are $\sqrt{\Lambda_k}$, not $\Lambda_k$.
\end{theorem}
\begin{proof}
The operator $K_t$ is self-adjoint. For $k\le t$, the elementary beta integral gives
\[
 (t+1)\int B_{t,j}(y)y^k\,dy
 =\frac{(j+1)(j+2)\cdots(j+k)}{(t+2)(t+3)\cdots(t+k+1)}.
\]
Taking expectation over $j\sim\mathrm{Bin}(t,x)$ shows that $K_t$ maps polynomials of degree at most $k$ into themselves. Its coefficient on the leading monomial $x^k$ is
$(t)_k/[(t+2)\cdots(t+k+1)]=\lambda_{t,k}$, because the leading term of the $k$th binomial moment is $(t)_kx^k$.

Self-adjointness now implies $K_tP_k$ is orthogonal to every polynomial of degree below $k$ while having degree at most $k$. It must be $\lambda_{t,k}P_k$. For $k>t$, all $B_{t,j}$ are orthogonal to $P_k$, giving zero. Averaging over $t$ gives (M08.1) and (M08.3), first on finite polynomials and then on $L^2$ by contraction and density.

For $t\ge k$, factorial cancellation gives
\[
 \lambda_{t,k}-\lambda_{t-1,k}
 =\frac{k(k+1)}{t(t+1)}\lambda_{t,k},\qquad \lambda_{k-1,k}=0.
\]
The differences are nonnegative. Expand $\lambda_{t,k}$ as their telescoping sum and apply Tonelli to $\sum_tw_t\lambda_{t,k}$; the coefficient of the difference at depth $t$ is $q_t$. This proves (M08.2). Finally $dQ_g/dQ_0=1+Sf$, so (M08.4) is an exact identity.
\end{proof}
If the cap has unbounded support, every $\Lambda_k>0$. If its maximum is $M$, exactly the modes $k>M$ have zero eigenvalue. Injectivity on $L^2$ follows in the former case; full identifiability on all mixing probabilities also follows by the power moments, independently of this Hilbert-space formulation.

\begin{corollary}[Metric comparison for bounded densities]
If $c\le g,h\le M$ almost everywhere, with $0<c\le M<\infty$, and $g-h=\sum a_kP_k$, then
\[
 \frac1{4M}\sum_k\Lambda_ka_k^2
 \le h^2(Q_g,Q_h)\le\frac1{4c}\sum_k\Lambda_ka_k^2,
\]
where $h^2(P,Q)=\int(\sqrt{dP}-\sqrt{dQ})^2$.
\end{corollary}
\begin{proof}
The data density ratios $dQ_g/dQ_0$ and $dQ_h/dQ_0$ lie in $[c,M]$ by conditional averaging. Divide their squared difference by the squared sum of their square roots, which lies in $[4c,4M]$, and use (M08.3).
\end{proof}

\subsection{Polynomial tails, including the exact leading constant}\label{M08-sec-3}
\begin{theorem}\label{M08-th:poly}
If $q_t\sim c t^{-\alpha}$ as $t\to\infty$, where $c,\alpha>0$, then
\[
 \Lambda_k\sim c\Gamma(\alpha+1)k^{-2\alpha}.
\tag{M08.5}
\]
\end{theorem}
\begin{proof}
The product representation gives
\[
 \log\lambda_{t,k}=\sum_{j=0}^{k-1}
 \log\left(1-\frac{2j+2}{t+j+2}\right).
\tag{M08.6}
\]
For $t=\lfloor k^2u\rfloor$, uniformly on compact subsets of $0<u<\infty$, expanding the logarithm yields $\log\lambda_{t,k}\to-1/u$. The error in replacing the logarithms by their linear terms is $O(k^3/t^2)$; the denominator replacement contributes the same order, which is $o(1)$ on this scale. Also, for all $t\ge k$,
\[
 \lambda_{t,k}\le\exp\left(-\frac{k(k+1)}{t+k+1}\right)
 \le e^{-k^2/(3t)}.\tag{M08.7}
\]
Insert these facts into the exact formula (M08.2), rescale $t=k^2u$, and divide by $c k^{-2\alpha}$. The sums converge on each fixed compact interval of $u$ to
\[
 \int_0^\infty e^{-1/u}u^{-\alpha-2}\,du
 =\int_0^\infty e^{-v}v^\alpha\,dv=\Gamma(\alpha+1).
\]
For a fully justified passage to the unbounded interval, use $q_t\le C t^{-\alpha}$ for all sufficiently large $t$. Equation (M08.7) bounds the rescaled summands by a constant times the Riemann-sum envelope $u^{-\alpha-2}e^{-1/(3u)}$. Its integrals near zero and infinity tend to zero. The corresponding discrete tails are uniformly bounded by constant multiples of those integrals, by splitting at the envelope's single turning point and comparing monotone sums to integrals. Thus no unaccounted small- or large-depth contribution remains.
\end{proof}
The critical polynomial cap exponent $\alpha=1/2$ therefore gives $\Lambda_k\sim c\Gamma(3/2)/k$. This exact statement alone does not settle a global $W_1$ minimax logarithm: the loss geometry and positivity constraints still matter.

\subsection{Stretched-exponential tails: three spectral regimes}\label{M08-sec-4}
Assume now $\Pr(L>m)=e^{-a m^\beta}$ exactly, with $a,\beta>0$.
\begin{theorem}\label{M08-th:stretch}
As $k\to\infty$,
\[
 -\log\Lambda_k\sim
 \begin{cases}
 C_{a,\beta}\,k^{2\beta/(1+\beta)},&0<\beta<1,\\
 C_a\,k,&\beta=1,\\
 a k^\beta,&\beta>1,
 \end{cases}\tag{M08.8}
\]
where
\[
 C_{a,\beta}=(1+\beta)a^{1/(1+\beta)}\beta^{-\beta/(1+\beta)},
 \qquad C_a=2\operatorname{arctanh}\sqrt{1-e^{-a}}.
\]
\end{theorem}

\subsubsection{Proof for $0<\beta<1$}
Put $T_k=k^{2/(1+\beta)}$ and $A_k=k^{2\beta/(1+\beta)}$, so $k^2/T_k=A_k$ and $T_k^\beta=A_k$. For $t=\lfloor uT_k\rfloor$ with $u$ in a fixed compact subinterval of $(0,\infty)$, $t/k\to\infty$. Formula (M08.6) yields uniformly
\[
 -\log\lambda_{t,k}=k^2/t+o(A_k)=A_k/u+o(A_k).
\]
The error $O(k/t+k^3/t^2)$ is $o(A_k)$ because $k/T_k\to0$. The cap mass is
\[
 w_t=e^{-a(t-1)^\beta}\left(1-e^{-a[t^\beta-(t-1)^\beta]}\right).
\]
Its logarithm satisfies $-\log w_t=a t^\beta+O(\log t)+o(A_k)$ uniformly on the same range: the second factor is asymptotic to $a\beta t^{\beta-1}$. Thus each summand in (M08.1) has exponential cost
\[
 A_k\{a u^\beta+1/u+o(1)\}.
\]
On a compact $u$-interval, the number of summands grows polynomially in $k$, so taking a sum changes its negative logarithm by $o(A_k)$. A term nearest the minimizer gives the matching lower estimate for the sum.

To justify restricting to such an interval, for $k\le t\le\delta T_k$, (M08.7)'s sharper first bound gives
$\lambda_{t,k}\le\exp\{-A_k/(\delta+o(1))\}$.
The total cap mass is at most one, so this entire range has that same upper bound. For $t\ge M T_k$, the sum is at most
$\Pr(L\ge\lceil M T_k\rceil)=\exp\{-a M^\beta A_k+o(A_k)\}$.
Taking small $\delta$ and large $M$ removes both tails. The minimum over $u>0$ is attained at $u=(a\beta)^{-1/(1+\beta)}$ and equals $C_{a,\beta}$, proving this case.

\subsubsection{Proof at $\beta=1$}
Here $w_t=(1-e^{-a})e^{-a(t-1)}$. Put $t=\lfloor uk\rfloor$, $u\ge1$. Uniform Stirling estimates on compact $u$-intervals, including $u=1$ with $0\log0=0$, give
\[
 -\frac1k\log\lambda_{t,k}\longrightarrow
 g(u)=(u-1)\log(u-1)+(u+1)\log(u+1)-2u\log u.
\]
The Stirling error is $O(\log k)$ uniformly because all factorial arguments lie between zero and a constant multiple of $k$. The sum in (M08.1) therefore has limiting exponential cost $\inf_{u\ge1}\{a u+g(u)\}$. As before, the cap tail beyond $Mk$ is exponentially negligible for sufficiently large $M$.

Since $g'(u)=\log(1-u^{-2})$, the unique minimizer is
$u_*=(1-e^{-a})^{-1/2}$. Moreover
\[
 g(u)=u\log(1-u^{-2})+\log\frac{u+1}{u-1},
\]
so $a u_*+g(u_*)=\log[(u_*+1)/(u_*-1)]=C_a$.
This proves the middle case and shows why replacing the exact fixed-cap spectrum by $e^{-k^2/t}$ at $t\asymp k$ would give the wrong constant.

\subsubsection{Proof for $\beta>1$}
Because $0\le\lambda_{t,k}\le1$ and it vanishes for $t<k$,
\[
 \Lambda_k\le\Pr(L\ge k)=e^{-a(k-1)^\beta}.
\]
Conversely the single summand at $t=k$ gives $\Lambda_k\ge w_k\lambda_{k,k}$. Here
$-\log w_k=a(k-1)^\beta+o(1)$ and
$-\log\lambda_{k,k}=2k\log2+O(\log k)$ by factorials. Divide by $k^\beta$ to obtain the matching limit $a$.

\subsection{Consequences, computation, and limitations}\label{M08-sec-5}
Formula (M08.2) is an exact nonnegative transform of the survival sequence; it is often numerically more stable than subtracting adjacent cap tails to form $w_t$. The fixed-cap coefficients can be generated recursively from
\[
 \lambda_{t,0}=1,\qquad
 \lambda_{t,k+1}=\lambda_{t,k}\frac{t-k}{t+k+2}.
\]
For a finite truncation of (M08.1), the omitted eigenvalue mass is at most the omitted cap probability. This is a rigorous numerical remainder bound, not an asymptotic heuristic.

The effective spectral degree satisfying $\Lambda_k\approx1/n$ has the same logarithmic exponents as the separately proved Wasserstein cap rates. This is an interpretation, not a substitute for minimax proofs. In particular, the singular-value formula does not remove the unresolved critical logarithm for arbitrary mixing measures under $\alpha=1/2$ polynomial survival.

The fixed-cap diagonalization and orthogonal-polynomial mechanism are classical. The research deliverable here is the exact survival transform and the sharp random-cap asymptotic constants, with all three stretched-exponential regimes distinguished. Publication priority for these exact tail-transform results has not been established. No computational output is needed for the proof, and no named pre-existing open question is claimed solved.





\clearpage
\section[M09 --- The Critical Random-Cap Mixture Rate: A Common-Prior Proof of the Wasserstein Logarithm]{The Critical Random-Cap Mixture Rate: A Common-Prior Proof of the Wasserstein Logarithm}\label{contrib-M09}
\markboth{Mixture Identifiability, Inverse Problems, and Minimax Recovery}{M09 --- Meiji}
\noindent\textsc{M09 \quad Meiji programme, source round 9}\par\medskip
\noindent\textit{Integration note.} This critical-case theorem is included in the more general information criterion and logarithmic hierarchy of Contribution~\ref{contrib-M10}.\par\medskip


\begin{quote}\small\textbf{Overview.} For binomial mixture observations with an independent observed cap satisfying $\Pr(L\ge t)\asymp t^{-1/2}$, we prove that the minimax expected $W_1$ error is exactly of order $\sqrt{\log(en)/n}$. The lower bound already holds for mixing densities between $1/2$ and $3/2$. Its central step uses one compact product prior on Legendre coefficients and a translation inequality for the joint parameter--data density. This supplies simultaneous pointwise CDF uncertainty without incorrectly adding coordinatewise Assouad losses under an $L^1$ norm. An elementary uniform bound for integrated Legendre polynomials converts this uncertainty into the logarithm. A multilevel upper bound retains all covariance terms. Thus the previously unresolved critical case of this random-cap model is closed; no claim about a named pre-existing open problem is made.\end{quote}

\subsection{Model and main result}\label{M09-sec-1}
For each of $n$ independent individuals, draw $p\sim\mu\in\PP([0,1])$ and, independently, a positive integer cap $L$ with known distribution. Conditional on $(p,L)$, observe $L$ independent Bernoulli$(p)$ variables. Equivalently, observe $(L,S)$, where $S$ is their sum: conditional on $(L,S)$, the word is uniform over all words with that count, independently of $\mu$.

Write $q_t=\Pr(L\ge t)$ and assume that fixed constants $0<c_0\le C_0<\infty$ satisfy
\[
 c_0t^{-1/2}\le q_t\le C_0t^{-1/2}\qquad(t\ge1).\tag{M09.1}
\]
The loss is $W_1(\nu,\mu)=\int_0^1|F_\nu(x)-F_\mu(x)|\,dx$. Let $R_n(\mathcal C)=\inf_{\widehat\mu}\sup_{\mu\in\mathcal C}\E_\mu W_1(\widehat\mu,\mu)$, and let $\mathcal G$ be the class of probability densities $g$ with $1/2\le g\le3/2$ almost everywhere.

\begin{theorem}\label{M09-th:main}
Under (M09.1), constants $0<c<C<\infty$, depending only on $c_0,C_0$, satisfy, for all $n\ge2$,
\[
 c\sqrt{\frac{\log(en)}n}
 \le R_n(\mathcal G)
 \le R_n(\PP([0,1]))
 \le C\sqrt{\frac{\log(en)}n}.\tag{M09.2}
\]
\end{theorem}
The upper bound is multilevel Monte Carlo in spirit; that general methodology is classical~\cite{ref019}. The missing lower logarithm is not obtained by summing separate single-coordinate testing bounds: $L^1$ cancellation makes that argument invalid. Sections~\ref{M09-sec-3}--\ref{M09-sec-5} give a common-prior argument that avoids the problem. Bayesian information and testing bounds are established tools~\cite{ref023,ref024}; the joint-translation lemma below is proved directly rather than invoked as a new general principle. Classical Bernstein--Durrmeyer diagonalization~\cite{ref021} is also proved in the form needed here.

\subsection{Upper bound with cross-level covariance}\label{M09-sec-2}
Let $r_\ell=2^\ell$, and imagine each individual has an infinite coin sequence. Write $U_\ell$ for its first-$r_\ell$ sample mean and, as functions of $x\in[0,1]$, put
\[
 D_0(x)=\1\{U_0\le x\},\qquad
 D_\ell(x)=\1\{U_\ell\le x\}-\1\{U_{\ell-1}\le x\}\quad(\ell\ge1).
\]
For one individual define the observable function
\[
 Y_J(x)=D_0(x)+\sum_{\ell=1}^J
             \frac{\1\{L\ge r_\ell\}}{q_{r_\ell}}D_\ell(x).
\tag{M09.3}
\]
The independent cap gives $\E Y_J=F_{U_J}$. Since $\E|U_\ell-p|\le2^{-\ell/2}/2$,
\[
 \int|F_{U_J}-F_\mu|\le\E|U_J-p|\le2^{-J/2}/2,
 \qquad \E\int|D_k|\le C2^{-k/2}.\tag{M09.4}
\]
Here and below unmarked integrals are over $[0,1]$. Set $Z_\ell=\1\{L\ge r_\ell\}D_\ell/q_{r_\ell}$. For $1\le\ell\le k$, nested cap events and independence give exactly
\[
 \E\int|Z_\ell Z_k|
 =q_{r_\ell}^{-1}\E\int|D_\ell D_k|
 \le C2^{\ell/2}2^{-k/2}.\tag{M09.5}
\]
We used $|D_\ell|\le1$, not independence across levels. Terms involving $D_0$ sum to a constant by (M09.4). Summing (M09.5) first over $k\ge\ell$ yields
\[
 \int\Var(Y_J(x))\,dx\le\E\int Y_J(x)^2\,dx\le C(J+1).\tag{M09.6}
\]
For the sample average $\overline Y_J$, Cauchy--Schwarz gives
\[
 \E\int|\overline Y_J-F_\mu|
 \le C\{2^{-J/2}+\sqrt{(J+1)/n}\}.\tag{M09.7}
\]
Project $\overline Y_J$ approximately in $L^1$ onto probability CDFs, within $1/n$ of the infimum. A measurable implementation enumerates the countable class of finite rationally weighted measures at rational points and chooses the first candidate within $1/n$ of that countable infimum. This class is $W_1$-dense. Triangle inequality bounds its error by twice (M09.7), plus $1/n$. Taking $J=\lfloor\log_2 n\rfloor$ proves the upper bound. For count-only data, condition (M09.3) on $(L,S)$; hypergeometric averaging preserves its expectation and decreases integrated variance. Thus no unobserved order information is used.

\subsection{Information spectrum and polynomial primitives}\label{M09-sec-3}
Let $P_k(x)=\sqrt{2k+1}\,L_k(2x-1)$, where $L_k$ is the ordinary Legendre polynomial normalized by $L_k(1)=1$. These form an orthonormal basis of $L^2[0,1]$, and $P_0=1$. Put
\[
 A_k(x)=\int_0^xP_k(u)\,du\qquad(k\ge1).
\]
At the uniform mixing density, the count observation law is $Q_0(t,j)=w_t/(t+1)$, where $w_t=\Pr(L=t)$. The score operator is
\[
 (Sf)(t,j)=(t+1)\int_0^1 \binom{t}{j}x^j(1-x)^{t-j}f(x)\,dx.
\]

\begin{lemma}\label{M09-lem:spectrum}
The scores $SP_k$ are orthogonal in $L^2(Q_0)$, with squared norms
\[
 \Lambda_k=\sum_{t\ge k}w_t\lambda_{t,k}
 =k(k+1)\sum_{t\ge k}\frac{q_t\lambda_{t,k}}{t(t+1)},\qquad
 \lambda_{t,k}=\frac{t!(t+1)!}{(t-k)!(t+k+1)!}.\tag{M09.8}
\]
Under (M09.1), $c/k\le\Lambda_k\le C/k$ for $k\ge1$.
\end{lemma}
\begin{proof}
For a fixed $t$, the self-adjoint operator $S_t^*S_t$ preserves polynomials of each degree at most $t$. Indeed, its action on $x^k$ has leading coefficient $(t)_k/[(t+2)\cdots(t+k+1)]$, by the beta integral and the binomial factorial moments. Orthogonality therefore makes $P_k$ an eigenfunction with eigenvalue $\lambda_{t,k}$; modes $k>t$ vanish. Average over $t$. Direct factorial cancellation gives
$\lambda_{t,k}-\lambda_{t-1,k}=k(k+1)\lambda_{t,k}/[t(t+1)]$, with $\lambda_{k-1,k}=0$. Nonnegative summation by parts gives (M09.8).

The product formula implies
\[
 \lambda_{t,k}\le\exp\{-k(k+1)/(t+k+1)\}\le e^{-k^2/(3t)}\quad(t\ge k).
\]
Using (M09.1) in (M09.8), the upper bound follows from
$k^2\sum_{t\ge k}t^{-5/2}e^{-k^2/(3t)}\le C/k$. To check the sum bound, rescale $t=k^2u$ and compare with the integrable function $u^{-5/2}e^{-1/(3u)}$, separately on its increasing and decreasing portions. For the lower bound, retain $2k^2\le t\le4k^2$. In this range the factors in the product have the form $1-u_j$, with $0\le u_j\le1/2$ and $\sum_j u_j\le C$; hence $\lambda_{t,k}\ge c$. There are order $k^2$ summands, each contributing at least order $k^{-3}$ after the factor $k(k+1)$ in (M09.8). This proves the lower bound, with finitely many small cases absorbed into constants.
\end{proof}

\begin{lemma}\label{M09-lem:primitive}
There are numerical constants $c,C>0$ such that
\[
 \|P_k\|_\infty\le\sqrt{2k+1},\qquad
 \|A_k\|_\infty\le C/k,\qquad
 c/k^2\le\int A_k^2\le C/k^2.\tag{M09.9}
\]
Consequently, for $K\ge2$,
\[
 \int_0^1\sqrt{\sum_{k=1}^K k A_k(x)^2}\,dx\ge c\sqrt{\log K}.\tag{M09.10}
\]
\end{lemma}
\begin{proof}
For $u=\cos\theta$ and $s=\sin\theta$, the elementary Laplace representation is
\[
 L_k(u)=\frac1\pi\int_0^\pi(u+i s\cos\phi)^k\,d\phi.\tag{M09.11}
\]
It follows by expanding the generating function: the integral of $[1-z(u+i s\cos\phi)]^{-1}$ equals $(1-2uz+z^2)^{-1/2}$, as is verified by the substitution $\tan(\phi/2)$. Since the integrand base has modulus at most one, $|L_k(u)|\le1$.

The polynomial identity $(L_{k+1}-L_{k-1})'=(2k+1)L_k$, checked from that same generating function, gives
\[
 A_k(x)=\frac{L_{k+1}(2x-1)-L_{k-1}(2x-1)}{2\sqrt{2k+1}},
\quad
 \int A_k^2=\frac1{4(2k+1)}\left(\frac1{2k+3}+\frac1{2k-1}\right).\tag{M09.12}
\]
The latter uses Legendre orthogonality and normalization; alternatively both follow directly from Rodrigues' formula by integration by parts.

To obtain the uniform, rather than interior-only, primitive bound, put $z_\phi=u+i s\cos\phi$. Then $|z_\phi^2-1|\le4s$ and $|z_\phi|^2=1-s^2\sin^2\phi$. For $k\ge2$, (M09.11) shows
\[
 |L_{k+1}(u)-L_{k-1}(u)|
 \le Cs\int_0^\pi e^{-(k-1)s^2\sin^2\phi/2}\,d\phi
 \le Cs\min\{1,(s\sqrt{k-1})^{-1}\}\le C/\sqrt k.
\]
The middle inequality follows from symmetry and $\sin\phi\ge2\phi/\pi$ on $[0,\pi/2]$; endpoints $s=0$ give zero directly. Equation (M09.12) proves $\|A_k\|_\infty\le C/k$; $k=1$ is immediate.

Finally $V_K=\sum_{k\le K}kA_k^2$ satisfies $\sup V_K\le C\log(eK)$ and $\int V_K\ge c\log(eK)$ by (M09.9). Since $\sqrt{V_K}\ge V_K/\sqrt{\sup V_K}$, (M09.10) follows. No pointwise oscillatory asymptotic is needed.
\end{proof}

\subsection{A common-prior absolute-error inequality}\label{M09-sec-4}
\begin{lemma}[Joint translation]\label{M09-lem:shift}
Let $q(\theta,z)$ be a joint probability density on $\R^K\times\mathcal Z$. Suppose its square root has square-integrable weak directional derivatives and a positive definite matrix $I$ satisfies
\[
 4\int |v\cdot\nabla_\theta\sqrt q|^2\,d\theta\,dz\le v^TIv
 \quad\text{for every }v\in\R^K.\tag{M09.13}
\]
For every estimator $T(z)$ and vector $a\in\R^K$,
\[
 \E_q|T(z)-a^T\theta|\ge\frac18\sqrt{a^TI^{-1}a}.\tag{M09.14}
\]
\end{lemma}
\begin{proof}
The case $a=0$ is trivial. Set $v=I^{-1}a/\sqrt{a^TI^{-1}a}$ and $d=a^Tv=\sqrt{a^TI^{-1}a}$. Let $q_v(\theta,z)=q(\theta-v,z)$. Translation in $L^2$, the fundamental theorem of calculus along lines, and (M09.13) give
\[
 \TV(q,q_v)\le\|\sqrt q-\sqrt{q_v}\|_2
 \le\|v\cdot\nabla_\theta\sqrt q\|_2\le1/2.\tag{M09.15}
\]
For weak derivatives, the same inequality follows by smooth approximation in the Sobolev norm.

Let $R=\E_q|T-a^T\theta|$ and $B=\{|T-a^T\theta|<d/2\}$. Markov's inequality gives $q(B)\ge1-2R/d$. Under $q_v$, substitute $\eta=\theta-v$: membership in $B$ implies $|T-a^T\eta-d|<d/2$, hence $|T-a^T\eta|>d/2$. Therefore $q_v(B)\le2R/d$. It follows that $1-4R/d\le\TV(q,q_v)\le1/2$, proving (M09.14). The translated joint law is a testing device only; its shifted parameter need not belong to the original model.
\end{proof}
This proof uses an absolute-error test, not a mean-square bound followed by an invalid reversal of Jensen's inequality. A single joint density $q$ may be used for all vectors $a$; this is essential when we later integrate over $x$.

\subsection{Lower bound: positivity, information, and transport}\label{M09-sec-5}
Choose an integer $K$ and define
\[
 g_\theta(x)=1+\sum_{k=1}^K\theta_kP_k(x),\qquad
 \sigma_k=(n\Lambda_k)^{-1/2}.
\]
Give the coefficients independent prior densities
\[
 \pi_k(u)=\frac1{\sigma_k}\cos^2\left(\frac{\pi u}{2\sigma_k}\right)
             \1\{|u|\le\sigma_k\}.\tag{M09.16}
\]
This is normalized and has Fisher information $\pi^2/\sigma_k^2$. By Lemmas~\ref{M09-lem:spectrum} and \ref{M09-lem:primitive}, throughout the prior box,
\[
 \|g_\theta-1\|_\infty
 \le\sum_{k\le K}\sigma_k\sqrt{2k+1}
 \le C K^2/\sqrt n.\tag{M09.17}
\]
Take $K=\lfloor\eta n^{1/4}\rfloor$ for a fixed sufficiently small $\eta>0$. For large $n$, every prior density then lies in $\mathcal G$. Each integrates to one because $P_k\perp1$ for $k\ge1$.

Let $q(\theta,z)=\pi(\theta)Q_{g_\theta}^{\otimes n}(z)$ be the joint prior--data density. Conditional averaging gives
\[
 \frac{dQ_{g_\theta}}{dQ_0}=1+S(g_\theta-1)\in[1/2,3/2].
\]
For a direction $v$, one-observation Fisher information is therefore bounded by
\[
 \int\frac{[S(\sum v_kP_k)]^2}{1+S(g_\theta-1)}\,dQ_0
 \le2\sum_{k\le K}\Lambda_kv_k^2.\tag{M09.18}
\]
The prior score and likelihood score have zero cross term, since each conditional likelihood score has mean zero. The $n$ independent likelihood scores also have zero cross terms. Thus
\[
 4\int|v\cdot\nabla_\theta\sqrt q|^2
 \le(\pi^2+2)n\sum_{k\le K}\Lambda_kv_k^2.\tag{M09.19}
\]
There is no boundary contribution: $\sqrt{\pi_k}$ is a cosine that vanishes at both ends of its support, extended by zero. It is globally in the first-order Sobolev space. On the support, the likelihood is smooth and bounded away from zero relative to $Q_0^{\otimes n}$; (M09.18) and the finite sum of prior information verify the square-integrable weak derivatives required in Lemma~\ref{M09-lem:shift}.

Apply that lemma with $I=(\pi^2+2)n\diag(\Lambda_1,\ldots,\Lambda_K)$, $a_k=A_k(x)$, and $T=F_{\widehat\mu}(x)-x$. For every $x$, under this \emph{same} prior,
\[
 \E_q|F_{\widehat\mu}(x)-F_{g_\theta}(x)|
 \ge\frac1{8\sqrt{\pi^2+2}}\sqrt{\frac1n\sum_{k\le K}\frac{A_k(x)^2}{\Lambda_k}}.
\tag{M09.20}
\]
Integrate using Tonelli. The upper eigenvalue bound $\Lambda_k\le C/k$ and (M09.10) yield
\[
 \sup_{g\in\mathcal G}\E_gW_1(\widehat\mu,g)
 \ge\E_qW_1(\widehat\mu,g_\theta)
 \ge c\sqrt{\frac{\log K}n}
 \ge c'\sqrt{\frac{\log n}n}.\tag{M09.21}
\]
Taking the infimum over estimators proves the lower bound for large $n$. For the remaining finite range, a two-point submodel $1\pm\varepsilon P_1$ with sufficiently small $\varepsilon>0$ has positive $W_1$ separation and overlapping data laws; reducing the constant completes (M09.2) for $n\ge2$.

\subsection{Scope and relation to preceding work}\label{M09-sec-6}
For polynomial cap tails $q_t\asymp t^{-\alpha}$, the neighboring noncritical rates in this model are $n^{-1/2}$ for $\alpha<1/2$ and $n^{-1/(2\alpha+1)}$ for $\alpha>1/2$. The present theorem proves that their junction has a genuine $\sqrt{\log n}$ penalty, even on a bounded positive density class. It therefore rules out a uniform root-$n$ improvement of the critical multilevel estimator, not merely a particular variance calculation.

Ordinary-smooth inverse problems and Wasserstein deconvolution have related critical phenomena; see, for example, Rousseau and Scricciolo~\cite{ref025}. The observation mechanism here is an observed random number of Bernoulli trials, not additive convolution. We do not transfer a published open-question label or claim publication priority for this exact model. The contribution claimed is the complete matched-rate theorem above, with the positivity, shared prior, full cap likelihood, and cross-level covariances checked explicitly. In particular, arbitrarily long rare caps are already included in every $\Lambda_k$ in (M09.18).





\clearpage
\section[M10 --- When Is Random-Cap Mixture Recovery Root-\texorpdfstring{$n$}{n}? An Exact Information Criterion and Critical Tail Hierarchies]{When Is Random-Cap Mixture Recovery Root-\texorpdfstring{$n$}{n}? An Exact Information Criterion and Critical Tail Hierarchies}\label{contrib-M10}
\markboth{Mixture Identifiability, Inverse Problems, and Minimax Recovery}{M10 --- Meiji}
\noindent\textsc{M10 \quad Meiji programme, source round 10}\par\medskip


\begin{quote}\small\textbf{Overview.} We characterize uniform root-$n$ Wasserstein recovery on a bounded positive mixing-density class for every known independent unbounded cap law. If $\Lambda_k$ are its Legendre information eigenvalues, the criterion is exactly $\sum_{k\ge1}(k^2\Lambda_k)^{-1}<\infty$. Sufficiency uses an explicit unbiased inverse-CDF kernel whose sum is finite at every observed cap. Necessity holds for all estimators, not merely linear or unbiased ones, by a compact-prior translation argument. For critical tails $\Pr(L\ge t)\asymp t^{-1/2}\ell(t)$ with a smoothly slowly varying multiplier, the same root-$n$ criterion holds over all mixing probabilities and becomes $\int^\infty dt/[t\ell(t)]<\infty$. We also derive matched minimax rates throughout a two-parameter logarithmic tail hierarchy, including genuine double- and triple-logarithmic penalties. Exact rates for arbitrary slowly varying multipliers are not claimed.\end{quote}

\subsection{Model, loss, and the information criterion}\label{M10-sec-1}
For each independent individual, draw $p\sim\mu\in\PP([0,1])$ and an independent observed positive integer cap $L$, then observe $S\sim\mathrm{Bin}(L,p)$. The cap law is known and has unbounded support. Write $w_t=\Pr(L=t)$, $q_t=\Pr(L\ge t)$, and
\[
 \Lambda_k=\sum_{t\ge k}w_t\frac{t!(t+1)!}{(t-k)!(t+k+1)!}>0
 \qquad(k\ge1).\tag{M10.1}
\]
For a class $\mathcal C$ of mixing measures, let
\[
 R_n(\mathcal C)=\inf_{\widehat\mu}\sup_{\mu\in\mathcal C}\E_\mu W_1(\widehat\mu,\mu),
 \qquad W_1(\nu,\mu)=\int_0^1|F_\nu-F_\mu|.
\]
Let $\mathcal G=\{g: \int g=1,\ 1/2\le g\le3/2\}$ and define
\[
 D_K=\sum_{k=1}^K\frac1{k^2\Lambda_k},\qquad
 B_K=\sum_{k=1}^K\sqrt{\frac{2k+1}{\Lambda_k}},\qquad
 D_\infty=\lim_KD_K.\tag{M10.2}
\]

\begin{theorem}[Exact uniform root-$n$ criterion]\label{M10-th:criterion}
For every known unbounded cap law,
\[
 R_n(\mathcal G)=O(n^{-1/2})
 \quad\Longleftrightarrow\quad D_\infty<\infty.\tag{M10.3}
\]
More precisely, there are numerical constants $c,C>0$ such that
\[
 B_K\le\tfrac12\sqrt n\quad\Longrightarrow\quad
 R_n(\mathcal G)\ge c\sqrt{D_K/n},\tag{M10.4}
\]
and, when $D_\infty<\infty$,
\[
 c n^{-1/2}\le R_n(\mathcal G)\le C\sqrt{D_\infty/n}.\tag{M10.5}
\]
When $D_\infty=\infty$, the stronger conclusion $\sqrt n R_n(\mathcal G)\to\infty$ holds.
\end{theorem}
The exact cap distribution may be highly irregular; no regular variation is assumed in this theorem. A separate result below extends the root-$n$ classification to \emph{all} mixing probabilities for a critical slowly varying family. We do not assert (M10.3) for arbitrary atomic mixing laws and arbitrary caps.

Information operators and inverse influence functions are classical statistical tools~\cite{ref026}; Bernstein--Durrmeyer diagonalization is classical approximation theory~\cite{ref021}. The point here is a global minimax equivalence for this experiment, with a necessity proof valid for every estimator, and its exact cap-tail consequences. Bayesian information/testing ideas are likewise classical~\cite{ref023}. All ingredients used in the theorem are proved below.

\subsection{Spectral identities and primitive geometry}\label{M10-sec-2}
Let $P_k(x)=\sqrt{2k+1}L_k(2x-1)$ be the orthonormal shifted Legendre polynomials, with $L_k(1)=1$, and put $A_k(x)=\int_0^xP_k(u)\,du$ for $k\ge1$. At the uniform mixing density the observation law is $Q_0(t,j)=w_t/(t+1)$, and define
\[
 (Sf)(t,j)=(t+1)\int_0^1\binom{t}{j}x^j(1-x)^{t-j}f(x)\,dx.
\]
Then $dQ_g/dQ_0=Sg$, and $S1=1$.

\begin{lemma}\label{M10-lem:ingredients}
The scores $SP_k$ are orthogonal, with $\|SP_k\|_{L^2(Q_0)}^2=\Lambda_k$. Moreover
\[
 \Lambda_k=k(k+1)\sum_{t\ge k}\frac{q_t\lambda_{t,k}}{t(t+1)},\qquad
 \lambda_{t,k}=\frac{t!(t+1)!}{(t-k)!(t+k+1)!},\tag{M10.6}
\]
and numerical constants satisfy
\[
 \|P_k\|_\infty\le\sqrt{2k+1},\quad
 \|A_k\|_\infty\le C/k,\quad
 \int_0^1 A_k^2=\frac1{4(2k+1)}\left(\frac1{2k+3}+\frac1{2k-1}\right)\asymp k^{-2}.
\tag{M10.7}
\]
In particular, for all positive weights $b_1,\ldots,b_K$,
\[
 \int_0^1\sqrt{\sum_{k\le K}b_kA_k(x)^2}\,dx
 \ge c\sqrt{\sum_{k\le K}b_k/k^2}.\tag{M10.8}
\]
\end{lemma}
\begin{proof}
For a fixed cap $t$, $S_t^*S_t$ is self-adjoint and preserves polynomials of each degree $k\le t$. The beta integral gives its leading coefficient on $x^k$ as $(t)_k/[(t+2)\cdots(t+k+1)]$. Thus orthogonality makes $P_k$ an eigenfunction with eigenvalue $\lambda_{t,k}$; for $k>t$ the score vanishes. Averaging over caps proves the score assertion. Factorial cancellation gives
$\lambda_{t,k}-\lambda_{t-1,k}=k(k+1)\lambda_{t,k}/[t(t+1)]$, so nonnegative summation by parts proves (M10.6).

For completeness, the classical integral representation
\[
 L_k(\cos\theta)=\pi^{-1}\int_0^\pi
       (\cos\theta+i\sin\theta\cos\phi)^k\,d\phi
\]
follows by integrating its geometric-series generating function; the result is $(1-2uz+z^2)^{-1/2}$. It immediately gives $|L_k|\le1$. The derivative identity $(L_{k+1}-L_{k-1})'=(2k+1)L_k$ gives
$A_k=[L_{k+1}(2x-1)-L_{k-1}(2x-1)]/[2\sqrt{2k+1}]$, and orthogonality gives the exact squared norm in (M10.7).

The uniform sup bound requires more than $|L_k|\le1$. Put $s=\sin\theta$ and $z=\cos\theta+i s\cos\phi$. Then $|z^2-1|\le4s$, while $|z|^{k-1}\le e^{-(k-1)s^2\sin^2\phi/2}$. Integrating $z^{k-1}(z^2-1)$ bounds $|L_{k+1}-L_{k-1}|$ by
$Cs\min(1,(s\sqrt{k-1})^{-1})\le C/\sqrt k$ for $k\ge2$; use $\sin\phi\ge2\phi/\pi$ on $[0,\pi/2]$ and symmetry. Endpoints and $k=1$ are immediate. This proves (M10.7).

Finally let $V=\sum b_kA_k^2$ and $D=\sum b_k/k^2$. Then $\int V\ge cD$ and $\sup V\le CD$, so $\int\sqrt V\ge\int V/\sqrt{\sup V}\ge c'\sqrt D$, proving (M10.8).
\end{proof}

\subsection{An explicit inverse-CDF estimator}\label{M10-sec-3}
Assume $D_\infty<\infty$. For a data point $z=(t,j)$ define
\[
 H(x;t,j)=\sum_{k=1}^{t}\frac{A_k(x)(SP_k)(t,j)}{\Lambda_k}.\tag{M10.9}
\]
This is a \emph{finite} sum for every observation. It is also the $L^2(dx\,dQ_0)$ limit of the infinite score series, because orthogonality gives
\[
 \int_0^1\E_{Q_0}H(x;Z)^2\,dx
 =\mathcal D:=\sum_{k\ge1}\frac{\int A_k^2}{\Lambda_k}\asymp D_\infty<\infty.
\tag{M10.10}
\]
The two definitions agree: the $k$th score vanishes whenever $k>t$, so the partial sums stabilize at each data point. The convergent Hilbert-space series therefore has that pointwise-in-data version.

For $g\in\mathcal G$, $dQ_g/dQ_0\le3/2$. The same series converges in $L^2(dx\,dQ_g)$, and
\[
 \E_g H(x;Z)=F_g(x)-x\quad\text{for almost every }x.\tag{M10.11}
\]
Indeed, if $g-1=\sum a_kP_k$, then $\E_g SP_k=\Lambda_ka_k$. The expectations of the truncated sums are $\sum_{k\le K}a_kA_k$, which converge to $F_g-x$, even uniformly by Cauchy--Schwarz applied to the $L^2$ approximation of $g-1$.

Define $\widetilde F_n(x)=x+n^{-1}\sum_{i=1}^nH(x;Z_i)$. Equations (M10.10)--(M10.11) yield
\[
 \E_g\int|\widetilde F_n-F_g|
 \le\sqrt{\frac1n\int\Var_g H(x;Z)\,dx}
 \le\sqrt{\frac{3\mathcal D}{2n}}.\tag{M10.12}
\]
Project approximately onto probability CDFs in $L^1$, within $1/n$ of the infimum. One measurable construction enumerates finite rationally weighted measures with rational locations, a $W_1$-dense countable class, and takes the first candidate within $1/n$ of its infimum. Triangle inequality doubles (M10.12) and adds $1/n$. Since $D_\infty\ge1/\Lambda_1\ge1$, this proves the upper bound in (M10.5).

\begin{proposition}[Minimum uniform-reference variance]
For fixed $x$, any $\phi_x\in L^2(Q_0)$ with $\E_{Q_0}\phi_x=0$ and
$\E_g\phi_x=F_g(x)-x$ for every $g\in\mathcal G$ must satisfy
\[
 \E_{Q_0}\phi_x^2\ge\sum_{k\ge1}\frac{A_k(x)^2}{\Lambda_k}.\tag{M10.13}
\]
When the right side is finite, the score series in (M10.9) attains equality.
\end{proposition}
\begin{proof}
Vary $g=1+\varepsilon P_k$ for sufficiently small positive and negative $\varepsilon$. The unbiasedness identity forces $\langle\phi_x,SP_k\rangle=A_k(x)$. Bessel's inequality for the orthonormal scores $SP_k/\sqrt{\Lambda_k}$ proves (M10.13), and the series attains its bound. This is the usual minimum-norm inverse principle, not a new general efficiency theory.
\end{proof}
The necessity in Theorem~\ref{M10-th:criterion} is substantially stronger than (M10.13): it applies to arbitrary biased, nonlinear, and sample-size-dependent estimators.

\subsection{A finite-dimensional lower certificate for all estimators}\label{M10-sec-4}
We give the information/testing argument in full to make the result independent of earlier manuscripts.

\begin{lemma}[Absolute-error information bound]\label{M10-lem:translation}
Let a joint parameter--data density $f(\theta,z)$ satisfy
$4\int|v\cdot\nabla\sqrt f|^2\le v^TIv$ for a positive definite matrix $I$, with weak derivatives in $L^2$. Then every estimator $T(z)$ obeys
\[
 \E_f|T-a^T\theta|\ge\tfrac18\sqrt{a^TI^{-1}a}.\tag{M10.14}
\]
\end{lemma}
\begin{proof}
For $a\ne0$, put $v=I^{-1}a/\sqrt{a^TI^{-1}a}$ and $d=a^Tv$. The translated joint density $f_v(\theta,z)=f(\theta-v,z)$ has total variation distance at most
$\|\sqrt f-\sqrt{f_v}\|_2\le\|v\cdot\nabla\sqrt f\|_2\le1/2$.
The translation inequality follows by integrating derivatives along lines, and by Sobolev approximation for weak derivatives. If $R=\E_f|T-a^T\theta|$ and $B=\{|T-a^T\theta|<d/2\}$, then $f(B)\ge1-2R/d$. Substitution $\eta=\theta-v$ shows $f_v(B)\le2R/d$, because the shifted interval lies outside $|T-a^T\eta|\le d/2$. Thus $1-4R/d\le\TV(f,f_v)\le1/2$, proving (M10.14). The translated density is only a testing device, so its shifted parameter need not be an admissible mixing density.
\end{proof}

Suppose $B_K\le\sqrt n/2$. Set $\sigma_k=(n\Lambda_k)^{-1/2}$ and use the independent coefficient prior
\[
 \pi_k(u)=\sigma_k^{-1}\cos^2(\pi u/(2\sigma_k))\1\{|u|\le\sigma_k\},
 \qquad g_\theta=1+\sum_{k=1}^K\theta_kP_k.
\tag{M10.15}
\]
The prior lies in $\mathcal G$, since $\|g_\theta-1\|_\infty\le B_K/\sqrt n\le1/2$. Its coordinate Fisher information is $\pi^2/\sigma_k^2$. One observation has directional Fisher information
\[
 \int\frac{[S(\sum v_kP_k)]^2}{Sg_\theta}\,dQ_0
 \le2\sum_{k\le K}\Lambda_kv_k^2.\tag{M10.16}
\]
The joint prior--sample information is consequently bounded by
$I=(\pi^2+2)n\diag(\Lambda_1,\ldots,\Lambda_K)$: all prior/likelihood and different-observation cross terms vanish by conditional mean-zero scores. The square root of each cosine-squared prior vanishes at its support endpoints; extended by zero it has a square-integrable weak derivative. Equation (M10.16) verifies the remaining derivative integrability, so Lemma~\ref{M10-lem:translation} applies without a boundary term.

Apply it, under this single prior, at every $x$ with $a_k=A_k(x)$ and $T=F_{\widehat\mu}(x)-x$. Tonelli and (M10.8) give
\[
 \sup_{g\in\mathcal G}\E_g W_1(\widehat\mu,g)
 \ge\frac{c}{\sqrt n}\int\sqrt{\sum_{k\le K}\frac{A_k(x)^2}{\Lambda_k}}\,dx
 \ge c'\sqrt{D_K/n}.\tag{M10.17}
\]
This proves (M10.4). If $D_\infty=\infty$, choose any fixed $K$ first and then let $n$ exceed $4B_K^2$; (M10.17) holds for every larger $n$. Letting $K\to\infty$ proves $\sqrt n R_n(\mathcal G)\to\infty$.

Finally a two-point submodel $g_\pm=1\pm c n^{-1/2}P_1$ supplies the universal lower bound $c'n^{-1/2}$ in (M10.5). These densities are in $\mathcal G$ for a small numerical $c$, their $W_1$ separation is $2c n^{-1/2}\int|A_1|$, and the data Hellinger distance is bounded by a constant times $c\sqrt{\Lambda_1}\le Cc$, including the $n$-fold product. Taking $c$ small keeps the total variation below $1/2$. This completes Theorem~\ref{M10-th:criterion}.

\subsection{Critical slowly varying tails: full-class classification}\label{M10-sec-5}
Assume now
\[
 q_t\asymp t^{-1/2}\ell(t),\tag{M10.18}
\]
where $\ell:[1,\infty)\to(0,\infty)$ is continuously differentiable and
\[
 \frac{t\ell'(t)}{\ell(t)}\longrightarrow0.\tag{M10.19}
\]
Two-sided comparability permits a finite initial modification to make (M10.18) a genuine survival law. Condition (M10.19) is imposed explicitly; no assertion is made here for all measurable slowly varying functions.

\begin{theorem}\label{M10-th:slow}
Under (M10.18)--(M10.19),
\[
 \Lambda_k\asymp k^{-1}\ell(k^2),\tag{M10.20}
\]
and the following are equivalent:
\begin{enumerate}[label=(\roman*)]
\item $R_n(\PP([0,1]))=O(n^{-1/2})$;
\item $R_n(\mathcal G)=O(n^{-1/2})$;
\item $\displaystyle\int_1^\infty\frac{dt}{t\ell(t)}<\infty$.
\end{enumerate}
For all integers $J\ge0$, one also has the full-class upper bound
\[
 R_n(\PP([0,1]))\le C\left\{2^{-J/2}
 +\sqrt{\frac{1+\sum_{j=1}^J\ell(2^j)^{-1}}n}\right\}+n^{-1}.
\tag{M10.21}
\]
\end{theorem}
\begin{proof}
For any $\varepsilon>0$, integrating (M10.19) above a sufficiently large fixed threshold, and absorbing the compact initial interval into a constant, gives
\[
 \frac{\ell(t)}{\ell(s)}\le C_\varepsilon
       \max\{(t/s)^\varepsilon,(s/t)^\varepsilon\}\quad(t,s\ge1).\tag{M10.22}
\]
The product formula in (M10.6) gives $\lambda_{t,k}\le e^{-k^2/(3t)}$. Substituting (M10.18) into (M10.6), rescaling $t=k^2u$, and using (M10.22) with $\varepsilon<1$ bounds $\Lambda_k/[k^{-1}\ell(k^2)]$ by a constant times the integral of
\[
 u^{-5/2}\max(u^\varepsilon,u^{-\varepsilon})e^{-1/(3u)}.
\]
This is integrable at both ends. The discrete sum comparison follows by splitting into intervals on which these finitely many power-exponential terms are monotone. The lower bound retains $2k^2\le t\le4k^2$, where $\lambda_{t,k}\ge c$ by $\log(1-v)\ge-2v$ for the product factors. In this range (M10.22) also gives $\ell(t)\asymp\ell(k^2)$. This proves (M10.20).

It follows that $D_\infty<\infty$ exactly when $\sum_k[k\ell(k^2)]^{-1}<\infty$. A sum--integral comparison is valid since (M10.19) bounds the logarithmic derivative on each unit or dyadic interval. Substituting $t=u^2$ makes this equivalent to (iii). Theorem~\ref{M10-th:criterion} proves (ii)$\Leftrightarrow$(iii); trivially (i)$\Rightarrow$(ii).

To prove (M10.21), let $U_j$ be the first-$2^j$ sample mean, and put
\[
 D_j(x)=\1\{U_j\le x\}-\1\{U_{j-1}\le x\}\quad(j\ge1),\qquad
 D_0(x)=\1\{U_0\le x\}.
\]
Use the observable telescoping estimator
\[
 Y_J=D_0+\sum_{j=1}^J\frac{\1\{L\ge2^j\}}{q_{2^j}}D_j.
\]
Its expectation is $F_{U_J}$ and its integrated bias is at most $2^{-J/2}/2$. Since $\E\int|D_k|\le C2^{-k/2}$, nested cap events give, for $1\le j\le k$,
\[
 \E\int\left|\frac{\1\{L\ge2^j\}D_j}{q_{2^j}}
                \frac{\1\{L\ge2^k\}D_k}{q_{2^k}}\right|
 \le C\ell(2^j)^{-1}2^{-(k-j)/2}.\tag{M10.23}
\]
Summing all diagonal and off-diagonal terms yields
$\int\Var Y_J\le C[1+\sum_{j\le J}\ell(2^j)^{-1}]$; base-level cross terms are summable separately. Average over individuals, apply Cauchy--Schwarz, and approximately project onto CDFs as before. For count-only observations, condition on $(L,S)$; conditional word order is uniform, and this conditioning preserves the mean while reducing variance. This proves (M10.21).

Finally (iii) implies $\sum_{j\ge1}\ell(2^j)^{-1}<\infty$, by comparison on the intervals $[2^j,2^{j+1}]$ using (M10.19). Take $J=\lceil\log_2n\rceil$ in (M10.21) to get (i).
\end{proof}
Thus the criterion is not merely that the cap exceed a particular power exponent. At the critical exponent, accumulated information across infinitely many scales determines root-$n$ feasibility.

\subsection{Matched logarithmic and iterated-logarithmic rates}\label{M10-sec-6}
Define positive iterated logarithms
\[
 L_1(t)=\log(et),\qquad L_2(t)=\log(eL_1(t)),\qquad
 L_3(t)=\log(eL_2(t)).
\]
For fixed real $\kappa,\eta$, suppose
\[
 q_t\asymp t^{-1/2}L_1(t)^\kappa L_2(t)^\eta.\tag{M10.24}
\]
The right side is eventually decreasing and tends to zero, so every such pair is realizable by a cap law after finite modification and rescaling.

\begin{theorem}\label{M10-th:logs}
Under (M10.24), both $R_n(\mathcal G)$ and $R_n(\PP([0,1]))$ are, up to positive constants, $n^{-1/2}$ times
\[
 \begin{cases}
 L_1(n)^{(1-\kappa)/2}L_2(n)^{-\eta/2},&\kappa<1,\\
 L_2(n)^{(1-\eta)/2},&\kappa=1,\ \eta<1,\\
 L_3(n)^{1/2},&\kappa=1,\ \eta=1,\\
 1,&\kappa>1\text{ or }(\kappa=1,\eta>1).
 \end{cases}\tag{M10.25}
\]
\end{theorem}
\begin{proof}
Put $H(J)=1+\sum_{j=1}^J(1+j)^{-\kappa}[\log(e+j)]^{-\eta}$. The upper bound (M10.21) with $J=\lceil\log_2n\rceil$ is $C\sqrt{H(J)/n}$: the bias is at most $n^{-1/2}$ and $H\ge1$.

For the lower bound, take $K=\lfloor n^{1/8}\rfloor$. By (M10.20) and (M10.22), $\Lambda_k^{-1/2}\le Ck^{3/4}$ for all $k$ after changing the constant. Hence $B_K\le C K^{9/4}=O(n^{9/32})=o(\sqrt n)$, so (M10.4) applies for all sufficiently large $n$. Moreover
\[
 D_K\asymp\sum_{k\le K}
 \frac1{k\,L_1(k)^\kappa L_2(k)^\eta}
 \asymp H(\lfloor\log K\rfloor).\tag{M10.26}
\]
The first equivalence uses $L_1(k^2)\asymp L_1(k)$ and $L_2(k^2)\asymp L_2(k)$; the second follows by substituting $u=\log k$ in the corresponding integral and absorbing the initial compact interval. Standard integral comparisons here are elementary: integrate $u^{-\kappa}(\log u)^{-\eta}$ for $\kappa\ne1$, and substitute $v=\log u$ for $\kappa=1$. They give
\[
 H(J)\asymp
 \begin{cases}
 J^{1-\kappa}(\log(eJ))^{-\eta},&\kappa<1,\\
 (\log(eJ))^{1-\eta},&\kappa=1,\eta<1,\\
 \log\log(e^eJ),&\kappa=1,\eta=1,\\
 1,&\kappa>1\text{ or }(\kappa=1,\eta>1).
 \end{cases}
\]
In all four cases, replacing $J$ by a fixed positive multiple of $J$ changes $H(J)$ only by a bounded factor. Since $\log K\asymp\log n$, (M10.26) and (M10.4) match the upper bound and give (M10.25). Finite small sample sizes are absorbed into constants using the two-point lower bound from Section~\ref{M10-sec-4}.
\end{proof}
For example, the cap $q_t\asymp t^{-1/2}\log(et)$ has risk $\sqrt{L_2(n)/n}$, not $n^{-1/2}$. Multiplying its tail by an additional $L_2(t)$ changes the penalty to $\sqrt{L_3(n)}$, while a power $L_2(t)^{1+\varepsilon}$ removes the penalty entirely. All these cap laws have infinite mean; this boundary classification does not contradict finite-mean lower rates.

\subsection{What is, and is not, resolved}\label{M10-sec-7}
The general criterion (M10.3) separates statistical recovery from identifiability: every unbounded cap identifies the mixing law, but root-$n$ recovery requires a summable inverse-information trace. The finite-dimensional certificate (M10.4) is useful even when an explicit rate is unavailable, and the inverse-CDF kernel (M10.9) exhibits why the trace condition is sufficient on bounded densities. For the critical smooth slow-variation class, Theorem~\ref{M10-th:slow} makes the root-$n$ boundary exact for arbitrary mixing probabilities. The two-log family then has completely matched rates.

For a general multiplier obeying (M10.19), (M10.21) and (M10.4) need not immediately match in rate: truncating the spectral degree at a fixed power of $n$ can change a slowly varying accumulation by an unbounded factor. We therefore claim the exact root-$n$ criterion, but not a universal sharp non-root-$n$ formula in that generality. We also do not claim that the bounded-density criterion automatically covers every irregular cap on the unrestricted mixing class.

Related additive-noise Wasserstein theory is developed by Dedecker, Fischer and Michel~\cite{ref027}; it is not the present random-cap experiment. The Bayesian translation method, Hilbert-space inverse principle, and fixed-cap polynomial diagonalization are not claimed new. Publication priority for the exact random-cap criteria and tail hierarchy has not been established. No named pre-existing open problem, NPMLE optimality, or superiority to the benchmark is asserted.





\clearpage
\part{Integral Presentations and Critical Groups of Graph Blowups}


\clearpage
\section[K01 --- Integral Compression of Block Constant Perturbations]{Integral Compression of Block Constant Perturbations}\label{contrib-K01}
\markboth{Integral Presentations and Critical Groups of Graph Blowups}{K01 --- Keio}
\noindent\textsc{K01 \quad Keio programme, source round 1}\par\medskip


\begin{quote}\small\textbf{Overview.} For an integer matrix with scalar diagonal blocks and constant off-diagonal blocks, we give an explicit integral reduction to a matrix of order twice the number of blocks, together with scalar summands. No division by a block size, eigenvalue, or prime is used. The construction works over every commutative ring, including polynomial rings, and therefore preserves all determinantal ideals. A weighted-network application and an example where spectral splitting gives the wrong group are included. The basic mechanism has published special cases; we do not claim a previously open problem has been solved.\end{quote}

\subsection{Statement and scope}\label{K01-sec-1}
Let $R$ be a commutative ring with identity. Fix $k\geq1$, integers $n_i\geq2$, elements $a_i\in R$, and $B=(b_{ij})\in R^{k\times k}$. Put $N=\sum_i n_i$, $D=\diag(a_1,\ldots,a_k)$, and $H=\diag(n_1,\ldots,n_k)$, where integers act through the structure map $\Z\to R$. Define the $N\times N$ matrix
\[
 A_{ij}=\delta_{ij}a_iI_{n_i}+b_{ij}J_{n_i,n_j}.
\]
Here the first summand is absent for $i\ne j$, and $J$ is an all-ones matrix. Matrices are \emph{equivalent} if related by invertible row and column operations. All operations below have determinant $1$ or $-1$ already over $\Z$.
\begin{theorem}\label{K01-main}
There is an equivalence over $R$
\[
 A\ \sim\ \left(\bigoplus_{i=1}^k a_iI_{n_i-2}\right)\oplus C,
 \qquad C=\begin{pmatrix}D&B\\0&D+HB\end{pmatrix}.                 \tag{K01.1}
\]
It is valid when $A$, $B$, or any $a_i$ is singular. In particular, over $\Z$,
\[
 \coker A\cong \bigoplus_i(\Z/a_i\Z)^{n_i-2}\oplus\coker C.       \tag{K01.2}
\]
We use $\Z/0\Z=\Z$ and omit summands with exponent zero.
\end{theorem}
This is an integral, rather than spectral, compression. The residual matrix is not generally diagonalizable integrally into $D$ and $D+HB$.

\subsection{Constructive proof}\label{K01-sec-2}
Use target generators $x_{i,t}$, $1\leq t\leq n_i$. The column indexed by $(j,t)$ is the relation
\[
 a_jx_{j,t}+\sum_i b_{ij}S_i=0,
 \qquad S_i=\sum_{s=1}^{n_i}x_{i,s}.
\]
Change generators to $z_i=x_{i,1}$ and $y_{i,t}=x_{i,t}-x_{i,1}$ for $t\geq2$. This is a unimodular change. Subtract the relation indexed by $(j,1)$ from each relation $(j,t)$, $t\geq2$. The resulting relations are
\[
 a_jy_{j,t}=0,\qquad
 a_jz_j+\sum_i b_{ij}\left(n_iz_i+\sum_{s=2}^{n_i}y_{i,s}\right)=0. \tag{K01.3}
\]
In each difference block use generators
\[
 u_i=\sum_{s=2}^{n_i}y_{i,s},\quad
 v_{i,2}=y_{i,2},\ldots,v_{i,n_i-1}=y_{i,n_i-1}.
\]
The coefficient of $y_{i,n_i}$ in $u_i$ is one. Hence these form a basis, including when $n_i=2$. The scalar relation matrix $a_iI$ remains scalar after the inverse change of relation basis. Consequently all $v_{i,s}$ occur only in their own relations $a_iv_{i,s}=0$. The remaining relations are
\[
 a_iu_i=0,\qquad
 a_jz_j+\sum_i b_{ij}u_i+\sum_i n_ib_{ij}z_i=0.
\]
With row order $(u_1,\ldots,u_k,z_1,\ldots,z_k)$ these are exactly $C$ in (K01.1). Every generator and relation change was explicitly invertible over $\Z$, so the proof establishes equivalence over every $R$, not merely equality of abstract cokernels. This proves Theorem~\ref{K01-main}.

\subsection{Consequences that retain exceptional primes}\label{K01-sec-3}
\begin{corollary}
For integer data,
\[
 \det A=\left(\prod_i a_i^{n_i-1}\right)\det(D+HB).               \tag{K01.4}
\]
If all block sizes are units in a commutative ring $R$, then there is the stronger equivalence
\[
 A\sim\left(\bigoplus_i a_iI_{n_i-1}\right)\oplus(D+HB).          \tag{K01.5}
\]
\end{corollary}
\begin{proof}
Equation (K01.4) follows from the triangular core in (K01.1). For (K01.5), split each block into the kernel of coordinate sum and the constant vector. The splitting is possible since $n_i$ is a unit. Both submodules are invariant under $A$. Its action on the former is scalar. On constant vectors the quotient matrix is $D+BH$, which is similar to $D+HB$ through $H$. The basis changes are invertible over $R$.
\end{proof}
For $k=1$, write $a=a_1$, $b=b_{11}$, $n=n_1$ and suppose $a(a+nb)\ne0$. The $2\times2$ core gives
\[
 \coker(aI_n+bJ_n)\cong
 (\Z/a\Z)^{n-2}\oplus\Z/g\Z\oplus\Z/\bigl(|a(a+nb)|/g\bigr)\Z,
 \quad g=\gcd(a,b).                                           \tag{K01.6}
\]
Indeed the first Smith factor is the gcd of the entries, $\gcd(a,b,a+nb)=g$, and the product is the absolute determinant. For $a=2,b=1,n=2$, the group is $\Z/8\Z$, not $\Z/2\Z\oplus\Z/4\Z$, despite the two eigenvalues being $2,4$. Thus prime $2$, which divides the block size, cannot simply be ignored.

Let $I_t(M)$ be the ideal generated by all $t$-minors, with $I_0=R$, $I_t=0$ outside the possible sizes. Equivalence preserves $I_t$. If $\mathcal D$ is the multiset containing $n_i-2$ copies of $a_i$, direct-sum expansion gives the exact identity
\[
 I_t(A)=\sum_{S\subseteq\mathcal D}\left(\prod_{a\in S}a\right) I_{t-|S|}(C). \tag{K01.7}
\]
A subset here chooses labeled copies. To prove (K01.7), any nonzero minor of a direct sum uses the same number of selected rows and columns from each summand and factors accordingly; conversely each product of minors occurs as a minor of the direct sum. This formula is valid over rings that are not principal ideal domains. It is not a claim that a Smith form exists over those rings.

\subsection{Weighted twin networks}\label{K01-sec-4}
Let $k$ clusters have sizes $n_i\geq2$ and add one sink. Give every pair inside cluster $i$ edge multiplicity $c_i\geq0$, every cross-cluster pair multiplicity $w_{ij}=w_{ji}\geq0$, and every vertex in cluster $i$ sink multiplicity $s_i\geq0$. Suppose the resulting multigraph is connected. Its reduced Laplacian has the form above with
\[
 b_{ii}=-c_i,\quad b_{ij}=-w_{ij}\ (i\ne j),\qquad
 a_i=s_i+n_ic_i+\sum_{j\ne i}n_jw_{ij}.                         \tag{K01.8}
\]
Therefore its critical group is given by (K01.2) without any restriction on the primes dividing $n_i$ or $a_i$. Equation (K01.4) counts spanning trees by the matrix-tree theorem. The computation needs the Smith form of just a $2k$-order core, irrespective of how large the twin classes are. For fixed $k$ the scalar factors and the core may be stored with multiplicities; expanding the final invariant-factor list still takes time proportional to its length.

\subsection{Literature and limitations}\label{K01-sec-5}
Dong, Guo and Jiang explicitly reduce the unweighted path-of-clusters case to a $2k$-order matrix \cite{ref028}. Alfaro, Corrales and Valencia study substantially richer multivariate critical ideals for graphs with twin vertices \cite{ref029}. The argument here abstracts a repeated-block reduction to arbitrary coefficient rings and arbitrary, not necessarily symmetric, block couplings. It does not purport to supersede those results or establish priority for this generalization. The scalar specialization $aI+bJ$ and the good-prime decomposition are classical. General Smith-form background is surveyed by Stanley \cite{ref030}.

The result is a complete reduction, not a closed list of invariant factors for arbitrary $B$. Computing those factors still requires a Smith calculation on $C$. No assertion is made about a named pre-existing open conjecture or superiority to the tournament benchmark.

\medskip\noindent\textit{Supplementary proof.}\par
\subsection{An exact reproducibility certificate}\label{K01-sec-6}
The proof provides a certificate independent of software: in each block first use $x_{i,t}-x_{i,1}$, subtract the first relation from the others, then replace the final difference generator by the sum of all difference generators. Every transformation has determinant $\pm1$. To check an instance, write (K01.3), verify the inverse formulas $y_{i,n_i}=u_i-\sum_{s=2}^{n_i-1}v_{i,s}$, and read off $C$. No numerical diagonalization or externally generated data is needed.





\clearpage
\section[K02 --- Critical Groups of Uniform Blowups of Matched Trees]{Critical Groups of Uniform Blowups of Matched Trees}\label{contrib-K02}
\markboth{Integral Presentations and Critical Groups of Graph Blowups}{K02 --- Keio}
\noindent\textsc{K02 \quad Keio programme, source round 2}\par\medskip
\noindent\textit{Integration note.} The star example below incorporates the supplied correction: the two groups have the same order but different generator ranks.\par\medskip


\begin{quote}\small\textbf{Overview.} Let $T$ be an arbitrary finite tree with a perfect matching, and replace each vertex by $n\geq2$ independent twins. We determine the entire critical group, at every prime, as an explicit direct sum depending only on $n$ and the degree multiset of $T$. No remaining Smith calculation is required. The proof combines an integral twin reduction, unimodularity of the tree adjacency matrix, and a lattice quotient. The matching assumption is essential: a star gives a different structure. We also obtain the exact exponent and minimal number of generators. Priority for this family formula has not been established.\end{quote}

\subsection{The family and the classification}\label{K02-sec-1}
For a finite simple tree $T$ on $m\geq2$ vertices let $d_v$ be its degrees. For $n\geq2$, define $T[n]$ to have vertex set $V(T)\times\{1,\ldots,n\}$, with $(u,i)$ adjacent to $(v,j)$ precisely when $uv\in E(T)$. No two twins inside the same cluster are adjacent. A \emph{perfect matching} is a collection of disjoint edges covering every vertex. The critical group $K(G)$ of a connected graph is the torsion subgroup of $\coker L_G$, where $L_G$ is its integer Laplacian.

\begin{theorem}\label{K02-classification}
Suppose $T$ has a perfect matching, and let $\ell$ be any leaf. Then
\[
 K(T[n])\cong
 \bigoplus_{v\in V(T)}\left(\Z/(nd_v)\Z\right)^{n-2}
 \ \oplus\ \bigoplus_{v\ne\ell}\Z/(n^2d_v)\Z.                 \tag{K02.1}
\]
In particular the isomorphism type is determined by the degree multiset of $T$ and $n$, not by the locations of its branching vertices. Formula (K02.1) includes all primes, even those dividing $n$ or one or more $d_v$.
\end{theorem}
The expression is a cyclic decomposition, not necessarily the divisibility-ordered invariant-factor decomposition. The latter follows by sorting the prime valuations of the displayed cyclic orders. Omission of $\ell$ is independent of the chosen leaf because $d_\ell=1$.

\subsection{Three integral lemmas}\label{K02-sec-2}
\begin{lemma}\label{K02-tree}
For a tree $T$, the lattice image of $L_T$ is
\[
 W=\{x\in\Z^m:\textstyle\sum_vx_v=0\},
\]
and $\ker L_T=\Z\mathbf1$. If $T$ has a perfect matching, its adjacency matrix $A_T$ is unimodular.
\end{lemma}
\begin{proof}
The real identity $x^tL_Tx=\sum_{uv\in E(T)}(x_u-x_v)^2$ gives the kernel. Orient the tree and let $F$ be its $m\times(m-1)$ incidence matrix. Deleting any one row gives a square matrix of determinant $\pm1$, by repeatedly removing a leaf distinct from the deleted vertex. Thus $F$ maps $\Z^{m-1}$ isomorphically onto $W$, and $F^t:\Z^m\to\Z^{m-1}$ is surjective. Hence $L_T=FF^t$ has image $W$.

In the permutation expansion of $\det A_T$, any nonzero term must be a product of disjoint transpositions: longer permutation cycles would give a cycle in $T$, while fixed points have zero diagonal entries. Thus terms correspond precisely to perfect matchings. Two distinct perfect matchings would have a union containing an alternating cycle, impossible in a tree. There is exactly one term, equal to $(-1)^{m/2}$.
\end{proof}

\begin{lemma}\label{K02-compression}
Let $D=\diag(d_v)$, $A=A_T$, and $L=D-A$. Then
\[
 L_{T[n]}\sim\left(\bigoplus_v nd_vI_{n-2}\right)
 \oplus\begin{pmatrix}nD&-A\\0&nL\end{pmatrix}.              \tag{K02.2}
\]
\end{lemma}
\begin{proof}
Present the cokernel with generators $x_{v,j}$. Their relations are $nd_vx_{v,j}-\sum_{u\sim v}\sum_tx_{u,t}=0$. Put $z_v=x_{v,1}$ and $y_{v,j}=x_{v,j}-z_v$. Subtracting the first relation in a cluster from the others gives $nd_vy_{v,j}=0$. Replace the difference generators by their sum $u_v$ and $n-2$ of the original differences. This is unimodular because the omitted difference has coefficient one in the sum. The remaining differences split off. The remaining relations have matrix in (K02.2), exactly as in the universal reduction \cite{ref031}; this direct proof makes the paper independent of that submission.
\end{proof}

\begin{lemma}\label{K02-product}
If $A_T$ is unimodular, then
\[
 M=L_TA_T^{-1}D\sim \diag\bigl((d_v)_{v\ne\ell},0\bigr).       \tag{K02.3}
\]
\end{lemma}
\begin{proof}
Set $X=A_T^{-1}D$, an integral matrix. Since $L_T$ maps $\Z^m$ onto $W$ with kernel $\Z\mathbf1$, the first isomorphism theorem yields
\[
 W/L_TX\Z^m\cong \Z^m/(X\Z^m+\Z\mathbf1).
\]
Multiplication by the unimodular matrix $A_T$ identifies the last quotient with
\[
 \Z^m/(D\Z^m+\Z A_T\mathbf1).
\]
But $A_T\mathbf1=(d_v)_v=D\mathbf1$, which is already in $D\Z^m$. Therefore
\[
 W/M\Z^m\cong\bigoplus_v\Z/d_v\Z.                          \tag{K02.4}
\]
Because $W$ is primitive in $\Z^m$, the exact sequence from $W/M\Z^m$ to $\coker M$ and then to $\Z^m/W\cong\Z$ splits. Thus $\coker M\cong\Z\oplus\bigoplus_v\Z/d_v\Z$. Remove the trivial factor indexed by $\ell$. Both sides of (K02.3) are $m\times m$ integer matrices with the same Smith form, which proves their equivalence.
\end{proof}

\subsection{Proof of the classification}\label{K02-sec-3}
In the core in (K02.2), interchange the two block columns. Its top-left block is now $-A$, unimodular. Integral block elimination gives
\[
 \begin{pmatrix}-A&nD\\nL&0\end{pmatrix}
 \sim I_m\oplus n^2LA^{-1}D.                                \tag{K02.5}
\]
For clarity, the Schur complement of $-A$ is $0-nL(-A)^{-1}nD=n^2LA^{-1}D$; all eliminations use the integral matrix $A^{-1}$. Multiplying the equivalence in (K02.3) by $n^2$ preserves it. Consequently (K02.2) and (K02.5) have exactly one zero Smith factor, $m$ unit factors, and the cyclic factors in (K02.1). Taking torsion proves Theorem~\ref{K02-classification}.

\subsection{Arithmetic consequences and an obstruction}\label{K02-sec-4}
\begin{corollary}
For the graphs in Theorem~\ref{K02-classification},
\[
 |K(T[n])|=n^{mn-2}\prod_vd_v^{n-1},\qquad
 \exp K(T[n])=n^2\lcm_vd_v.                                  \tag{K02.6}
\]
The minimal number of generators is exactly $m(n-1)-1$.
\end{corollary}
\begin{proof}
The order and exponent follow by multiplying and taking the lcm of the cyclic orders in (K02.1). There are $m(n-2)+(m-1)=m(n-1)-1$ displayed factors. Every one is divisible by every prime dividing $n$, so reduction modulo any such prime has precisely this dimension. No smaller generating set is possible, and the displayed factors give a generating set of that size.
\end{proof}
For the path on four vertices, degrees are $1,2,2,1$, and for $n=2$ the formula gives $\Z/4\Z\oplus(\Z/8\Z)^2$. For a branching example, start with any tree $S$ and attach a new pendant leaf to each vertex. The resulting tree $T$ has the evident perfect matching; if $S$ has degrees $e_1,\ldots,e_q$, the degree multiset of $T$ is $e_1+1,\ldots,e_q+1$ together with $q$ ones. Thus (K02.1) covers arbitrarily complicated branching and unbounded degrees, not just paths.

The matching hypothesis cannot be dropped. For $T=K_{1,r}$ with $r\geq2$, the graph $T[n]$ is $K_{n,rn}$. Its familiar critical-group formula is
\[
 K(K_{n,rn})\cong(\Z/n\Z)^{rn-2}\oplus(\Z/(rn)\Z)^{n-2}
                    \oplus\Z/(rn^2)\Z.                       \tag{K02.7}
\]
For $n=2,r=2$, this is $(\Z/2\Z)^2\oplus\Z/8\Z$. Naively applying (K02.1) instead gives $\Z/4\Z\oplus\Z/8\Z$. Both groups have order $32$, but their minimal numbers of generators are $3$ and $2$, respectively. Their group structures therefore differ, and the matching hypothesis cannot simply be omitted.

\subsection{Research position and reproducibility}\label{K02-sec-5}
The path-of-clusters critical-group calculations of Dong, Guo and Jiang \cite{ref028} overlap the even-path subfamily. The full complete multipartite family, including (K02.7), was determined by Jacobson, Niedermaier and Reiner \cite{ref032}. Twin reductions and critical ideals are studied in \cite{ref029}. The present contribution is the explicit matched-tree family formula and its integral proof. Searches for critical groups of blowups, lexicographic products and matched trees did not establish an identical prior theorem, but this is not proof of novelty. No named open problem or formal verification is claimed.

All examples can be checked directly from (K02.2), whose order is $2m$. A reproducibility check does not require computer algebra: verify the unique matching determinant in Lemma~\ref{K02-tree}, the incidence-lattice argument, and the integral block eliminations (K02.5). The determinant and the prime ranks in (K02.6) provide independent numerical sanity checks, but are not substitutes for the full group proof.





\clearpage
\section[K03 --- Matching Numbers Control the Prime Parts of Tree Blowup Groups]{Matching Numbers Control the Prime Parts of Tree Blowup Groups}\label{contrib-K03}
\markboth{Integral Presentations and Critical Groups of Graph Blowups}{K03 --- Keio}
\noindent\textsc{K03 \quad Keio programme, source round 3}\par\medskip
\noindent\textit{Integration note.} The common-prime cases left open at this stage are resolved in Contribution~\ref{contrib-K04}; two independent cluster sizes are treated in Contribution~\ref{contrib-K07}.\par\medskip


\begin{quote}\small\textbf{Overview.} For an arbitrary finite tree $T$, we classify the prime-primary critical group of its uniform independent-set blowup at every prime not dividing both the blowup size and a vertex degree. When the size and degree product are coprime, this is a complete classification. The answer is governed by the maximum matching number and the degree multiset. A local algebra lemma proves that only two torsion exponents occur in the exceptional-to-spectral regime; its proof counts congruence kernels rather than assuming diagonalizability. Trees with identical degree multisets but different matching numbers give different critical groups after the same blowup.\end{quote}

\subsection{Classification}\label{K03-sec-1}
Let $T$ be a simple tree on $m\geq2$ vertices, with degrees $d_v$ and maximum matching size $\nu$. Let $T[n]$, $n\geq2$, replace each vertex by $n$ independent twins and each edge by a complete bipartite graph. For a finite abelian group $G$, write $G_{(p)}$ for its $p$-primary subgroup, and use $v_p$ for the $p$-adic valuation.
\begin{theorem}\label{K03-main}
If $p\nmid n$, then
\[
 K(T[n])_{(p)}\cong\bigoplus_v
       \left(\Z/p^{v_p(d_v)}\Z\right)^{n-1}.                  \tag{K03.1}
\]
If $p\mid n$, $p\nmid\prod_vd_v$, and $\alpha=v_p(n)$, then
\[
 K(T[n])_{(p)}\cong
 (\Z/p^\alpha\Z)^{mn-4\nu}\oplus
 (\Z/p^{2\alpha}\Z)^{2\nu-1}.                               \tag{K03.2}
\]
In particular, if $\gcd(n,\prod_vd_v)=1$, the full group is
\[
 K(T[n])\cong (\Z/n\Z)^{mn-4\nu}\oplus
 (\Z/n^2\Z)^{2\nu-1}\oplus
 \bigoplus_v(\Z/d_v\Z)^{n-1}.                               \tag{K03.3}
\]
\end{theorem}
Trivial cyclic factors are omitted. Formula (K03.2) includes primes dividing the cluster size, where the usual constant-versus-difference splitting is not integral. The only primes left untreated by this theorem are the explicitly specified common divisors of $n$ and the degree product.

\subsection{A local algebra lemma}\label{K03-sec-2}
Work over the discrete valuation ring $R=\Z_{(p)}$; $R/p^hR\cong\Z/p^h\Z$. A free cokernel means free as an $R$-module.
\begin{lemma}\label{K03-local}
Let $P\in R^{m\times m}$, $Q=I-P$, and suppose $\coker P$ is free of rank $\delta$ and $\coker Q$ is free of rank $q$. Let $n=p^\alpha u$ with $u\in R^\times$, $\alpha\geq1$. For
\[
 C=\begin{pmatrix}nI&-P\\0&nQ\end{pmatrix},
\]
there is an isomorphism
\[
 \coker C\cong R^q\oplus(R/p^\alpha R)^{2\delta}
                 \oplus(R/p^{2\alpha}R)^{m-\delta-q}.          \tag{K03.4}
\]
\end{lemma}
\begin{proof}
The rational nullity of $C$ is $q$. First show that $n^2$ annihilates all torsion. If $[(v,w)]$ is torsion, its image in $\coker Q$ is torsion and hence zero. Thus $w=Qt$ for some $t\in R^m$. Direct multiplication gives
\[
 C\binom{nv+Pt}{nt}=n^2\binom vw.                            \tag{K03.5}
\]
It remains to determine which valuations between $1$ and $2\alpha$ occur.

Let $\kappa(h)=\log_p|\ker(C\bmod p^h)|$. For $1\leq h\leq\alpha$, the only nonzero condition is $Py=0\bmod p^h$, so $x$ is arbitrary and
\[
 \kappa(h)=(m+\delta)h.                                      \tag{K03.6}
\]
Here $P$ has Smith factors consisting of $m-\delta$ units and $\delta$ zeros.

For $\alpha\leq h\leq2\alpha$, the condition $nx=Py$ first requires $Py=0\bmod p^\alpha$. Every such residue of $y$ lifts from $\ker P$, because the cokernel of $P$ is free. Write
\[
 y=y_0+p^\alpha z,\qquad Py_0=0,
\]
where $y_0$ is chosen modulo $p^\alpha$ in the rank-$\delta$ kernel and $z$ varies modulo $p^{h-\alpha}$. Since $Qy_0=y_0$ and $h\leq2\alpha$, the second equation becomes $ny_0=0\bmod p^h$. There are $p^{\delta(2\alpha-h)}$ eligible choices for $y_0$, and $p^{m(h-\alpha)}$ choices for $z$. For each such $y$, the first equation has $p^{\alpha m}$ solutions for $x$. Consequently
\[
 \kappa(h)=(m-\delta)h+2\alpha\delta.                         \tag{K03.7}
\]
For a Smith factor $p^e$, its contribution to $\kappa(h)$ is $\min(e,h)$, and each zero factor contributes $h$. Thus the first differences of (K03.6)--(K03.7) show that exactly $2\delta$ nonzero valuations equal $\alpha$, and every remaining nonunit factor has valuation at least $2\alpha$. By (K03.5) no torsion valuation exceeds $2\alpha$. There are $q$ zero factors; the remaining $m-\delta-q$ factors have valuation $2\alpha$. This proves (K03.4).
\end{proof}
No assertion that $P$ is diagonalizable has entered the proof. In particular, nilpotent Jordan behavior modulo $p$ does not invalidate the kernel count.

\subsection{Tree lattices and matching rank}\label{K03-sec-3}
\begin{lemma}\label{K03-forest}
For a forest with $m$ vertices and maximum matching size $\nu$, its integer adjacency matrix has Smith form $I_{2\nu}\oplus0_{m-2\nu}$. For a tree, $\coker L_T\cong\Z$.
\end{lemma}
\begin{proof}
Choose a leaf $a$ with neighbor $b$ in any nontrivial component. In vertex order $(a,b,\text{rest})$, the adjacency matrix is
\[
 \begin{pmatrix}0&1&0\\1&0&t^t\\0&t&A'\end{pmatrix}.
\]
The first $2\times2$ block is unimodular. Its Schur complement is exactly $A'$, since $(0,t)\left(\begin{smallmatrix}0&1\\1&0\end{smallmatrix}\right)(0,t)^t=0$. Thus this pair contributes two unit Smith factors. A maximum matching can always be chosen to contain $ab$: if $b$ is already matched elsewhere replace that edge by $ab$, while if it is unmatched adding $ab$ improves the matching. Hence $\nu=1+\nu(T-\{a,b\})$. Induction proves the claim, including isolated vertices. For the Laplacian, orient the tree. Its incidence matrix maps the edge lattice isomorphically onto the zero-sum lattice, and its transpose is surjective, as follows by deleting one row and peeling leaves. Thus $L_T=FF^t$ has image the full zero-sum lattice and cokernel $\Z$.
\end{proof}

The integral twin reduction, proved by subtracting the first generator and first relation within every cluster, gives
\[
 L_{T[n]}\sim \bigoplus_v nd_vI_{n-2}\ \oplus C_T,
 \quad C_T=\begin{pmatrix}nD&-A\\0&n(D-A)\end{pmatrix},        \tag{K03.8}
\]
where $D=\diag(d_v)$ and $A$ is the adjacency matrix. For completeness, after the subtractions, the $n-1$ differences in a cluster obey scalar relations. Replace one by their sum, a unimodular move; the other $n-2$ then split off. The retained sum and the reference vertex give exactly (K03.8).

\subsection{Proof of the classification and a separating example}\label{K03-sec-4}
If $p\mid n$ but all degrees are units in $R$, multiply each block row of $C_T$ by $D^{-1}$. This yields the matrix in Lemma~\ref{K03-local}, with $P=D^{-1}A$ and $Q=D^{-1}L_T$. Lemma~\ref{K03-forest} gives $\delta=m-2\nu$ and $q=1$. The scalar factors in (K03.8) contribute $m(n-2)$ copies of $R/p^\alpha R$. The total number of these factors is $m(n-2)+2(m-2\nu)=mn-4\nu$, proving (K03.2).

If $p\nmid n$, divide coordinate sums by $n$ to split each cluster into a constant submodule and its sum-zero submodule. The latter contributes $(nd_v)I_{n-1}$; the constant submodule contributes $nL_T$. Since $n$ is a unit and $\coker L_T$ is free, this proves (K03.1). Primary decomposition then proves (K03.3).

Here is a structural distinction not visible in the degree multiset. Take a tree with one degree-three center and three arms, where an arm length counts edges from the center. The arm lengths $(3,1,1)$ and $(2,2,1)$ both give degree multiset $\{3,2,2,1,1,1\}$, but their matching sizes are $2$ and $3$ respectively. For $n=5$, their $5$-primary groups are
\[
 (\Z/5\Z)^{22}\oplus(\Z/25\Z)^3,
 \qquad (\Z/5\Z)^{18}\oplus(\Z/25\Z)^5.                      \tag{K03.9}
\]
Both have order $5^{28}$, but their minimal generator numbers differ. At other primes their critical groups agree by (K03.1). Thus the new information is genuinely group-theoretic, not just another spanning-tree count.

\subsection{Scope and relation to earlier work}\label{K03-sec-5}
Round 2 classified all primes for trees with a perfect matching, using an invertible adjacency matrix. This round removes that hypothesis but explicitly isolates common degree-and-size primes. It supplies a different mechanism: congruence-kernel counting and forest matching rank. The complete multipartite results of Jacobson, Niedermaier and Reiner \cite{ref032} include the star subfamily; Dong, Guo and Jiang \cite{ref028} treat path clusters, and Alfaro, Corrales and Valencia \cite{ref029} study twin critical ideals. These are relevant predecessors, not newly proved discoveries of this manuscript. Priority for the general family statement remains unestablished. The paper does not solve a named open problem.

\medskip\noindent\textit{Supplementary proof.}\par
\subsection{How to check the local calculation}\label{K03-sec-6}
For any integer matrix with Smith factors $s_i$ and $q$ zeros, compute $\kappa(h)=qh+\sum_i\min(h,v_p(s_i))$. Then $\kappa(h)-\kappa(h-1)$ counts the zero factors plus the factors with valuation at least $h$. Equations (K03.6)--(K03.7), together with the explicit annihilator (K03.5), determine every factor without large matrix computation. This supplies a finite exact certificate for (K03.2). Cases excluded by the degree condition must not be inferred from the formula.




\clearpage
\section[K04 --- The Critical Group of Every Uniform Tree Blowup]{The Critical Group of Every Uniform Tree Blowup}\label{contrib-K04}
\markboth{Integral Presentations and Critical Groups of Graph Blowups}{K04 --- Keio}
\noindent\textsc{K04 \quad Keio programme, source round 4}\par\medskip


\begin{quote}\small\textbf{Overview.} For every finite simple tree $T$ and every integer $n\geq2$, we determine the full critical group of the graph obtained by replacing each vertex by $n$ independent twins. The cyclic decomposition depends only on the degree multiset and the maximum matching number of $T$. The formula has no exceptional primes and no residual Smith calculation. The key step is a primitive-coupling lemma for integer block presentations, followed by a leaf-rooted reduction using the unimodular reduced Laplacian of a tree. This completes the family classification pursued in the preceding rounds. An identical prior theorem has not been established by the literature search; no novelty certification or named open-problem claim is made.\end{quote}

\subsection{The complete formula}\label{K04-sec-1}
Let $T$ be a finite simple tree on $m\geq2$ vertices. Write $d_v$ for its degrees, $\ell$ for its number of leaves, $\nu$ for its maximum matching size, and $\delta=m-2\nu$. Let $T[n]$ replace each vertex by $n\geq2$ independent vertices, with complete bipartite adjacency between clusters corresponding to edges of $T$. The critical group is the torsion of the full integer Laplacian cokernel.
\begin{theorem}\label{K04-main}
Set $h=2\nu+\ell-m-1=\ell-1-\delta$. Then $h\geq0$, and
\begin{align}
 K(T[n])\cong{}&\bigoplus_{v\in V(T)}(\Z/(nd_v)\Z)^{n-2}
 \oplus(\Z/n\Z)^{2\delta} \notag\\
 &\oplus(\Z/n^2\Z)^h
 \oplus\bigoplus_{d_v>1}\Z/(n^2d_v)\Z.                       \tag{K04.1}
\end{align}
This formula holds over all primes with no coprimality assumptions.
\end{theorem}
The factors in (K04.1) need not be in divisibility order. Sorting their prime valuations gives the invariant factors. In particular, trees with the same degree multiset and matching number have isomorphic critical groups after every uniform blowup. The classification does not assert that their canonical pairings or symmetry actions agree.

\subsection{A primitive-coupling lemma}\label{K04-sec-2}
A square integer matrix $K$ has \emph{primitive image} if its cokernel is free abelian, equivalently its nonzero Smith factors are all one.
\begin{lemma}\label{K04-coupling}
Let $D\in\Z^{q\times q}$ be nonsingular, and let $K\in\Z^{q\times q}$ have primitive image and rank $r=q-\delta$. Suppose the map with matrix $[D\ K]:\Z^{2q}\to\Z^q$ is surjective. There is a nonsingular integer $r\times r$ matrix $E$ with $\coker E\cong\coker D$ such that, for every positive integer $n$,
\[
 \begin{pmatrix}nD&K\\0&nI_q\end{pmatrix}
 \sim I_r\oplus nI_{2\delta}\oplus n^2E.                     \tag{K04.2}
\]
All equivalences are integral. Empty matrices are permitted when $r=0$.
\end{lemma}
\begin{proof}
Choose unimodular $U,V$ with $UKV=\diag(I_r,0_\delta)$. The bottom $\delta$ rows of $UD$ are surjective onto $\Z^\delta$, because $[UD\ UKV]$ is surjective and the last $\delta$ rows of $UKV$ are zero. A basis of the kernel together with lifts of the standard basis of $\Z^\delta$ therefore gives a unimodular $W$ with
\[
 UDW=\begin{pmatrix}E&F\\0&I_\delta\end{pmatrix}.             \tag{K04.3}
\]
The determinant shows $E$ is nonsingular. Eliminating $F$ in (K04.3) also shows $\coker E\cong\coker D$.

Multiply the block matrix in (K04.2) on the left by $\diag(U,V^{-1})$ and on the right by $\diag(W,V)$. In block sizes $(r,\delta,r,\delta)$ the result is
\[
 \begin{pmatrix}
 nE&nF&I_r&0\\
 0&nI_\delta&0&0\\
 0&0&nI_r&0\\
 0&0&0&nI_\delta
 \end{pmatrix}.                                               \tag{K04.4}
\]
Eliminate the unit block in the first block row and third block column. Its Schur complement has first row blocks $-n^2E,-n^2F,0$. Subtract suitable $nF$ multiples of the row block containing $nI_\delta$ to eliminate $-n^2F$. After sign changes and permutations the result is (K04.2). No inversion of $n$ or $D$ was used.
\end{proof}

\subsection{Two facts about trees}\label{K04-sec-3}
\begin{lemma}\label{K04-treefacts}
The adjacency matrix $A$ of $T$ has Smith form $I_{2\nu}\oplus0_\delta$. Every reduced Laplacian $\bar L$ of $T$ is unimodular. Finally $\delta\leq\ell-1$.
\end{lemma}
\begin{proof}
For the adjacency statement, select a leaf and its neighbor. The corresponding principal block is $\left(\begin{smallmatrix}0&1\\1&0\end{smallmatrix}\right)$, and its Schur complement is the adjacency matrix of the forest obtained by deleting the pair. Indeed the leaf has no other neighbor, so the Schur correction is zero. A maximum matching can be chosen to contain the leaf edge, by replacing the edge incident to its neighbor if necessary. Induction proves the Smith statement for forests, including isolated vertices.

For the Laplacian, orient the tree and delete one row from its incidence matrix. The resulting square matrix $F$ has determinant $\pm1$, by leaf elimination; $\bar L=FF^t$ has determinant one.

To prove the matching bound, choose a maximum matching minimizing the number of unmatched nonleaves. If a nonleaf is unmatched, start an alternating path there. At each even step, an internal vertex has a next edge other than the one just used. At an odd step the next vertex must be matched, or the path would augment the matching. Since the graph is a finite tree, the path eventually ends at a matched leaf at an even step. Reversing its matched and unmatched edges moves the unmatched vertex from an internal vertex to a leaf, a contradiction. Thus all $\delta$ unmatched vertices are leaves. At least one leaf is matched: otherwise start the same alternating-path procedure at any leaf; it cannot end at an unmatched odd-step vertex by maximality, or at a matched leaf by assumption, contradicting finiteness. Hence $\delta\leq\ell-1$.
\end{proof}

\subsection{Integral reduction to primitive coupling}\label{K04-sec-4}
Write $D_0=\diag(d_v)$, $A=A_T$ and $L=D_0-A$. Subtract the first generator from all other generators in each twin cluster, and do the analogous subtraction on the relations. The $n-1$ differences have scalar relation $nd_v$. Replace one of them by the sum of all differences, a unimodular operation. The other $n-2$ split off, leaving
\[
 L_{T[n]}\sim\bigoplus_v nd_vI_{n-2}\ \oplus C,
 \qquad C=\begin{pmatrix}nD_0&-A\\0&nL\end{pmatrix}.          \tag{K04.5}
\]
In a presentation of $\coker C$ use generators $u_v,z_v$ and relations
\[
 nd_vu_v=0,\qquad nLz-Au=0.                                  \tag{K04.6}
\]
The second expression is a vector of relations; symmetry of $A,L$ ensures it agrees with the column presentation.

Choose a leaf $a$ with neighbor $b$, put $V'=V(T)\setminus\{a\}$, and let $D=\diag(d_v:v\in V')$. Let $\bar A,\bar L$ be obtained by deleting row and column $a$, let $e$ denote the basis vector of $b$ in $\Z^{V'}$, and let $d=D\mathbf1$. Then
\[
 \bar L\mathbf1=e,\quad \bar L^{-1}e=\mathbf1.               \tag{K04.7}
\]
Replace the relation indexed by $a$ in the second family of (K04.6) by the sum of that family. It becomes $-\sum_vd_vu_v=0$. Since $d_a=1$, eliminate $u_a=-d^tu'$. Its scalar relation becomes redundant, because it is the negative sum of the other scalar relations. Also use $x_v=z_v-z_a$; the generator $z_a$ is free. The remaining second-family relations are
\[
 n\bar Lx-(\bar A-ed^t)u'=0.                                 \tag{K04.8}
\]
Multiplying the relation columns by the unimodular matrix $\bar L^{-1}$ yields
\[
 C\sim I_1\oplus S\oplus0_1,
 \quad S=\begin{pmatrix}nD&K\\0&nI_{m-1}\end{pmatrix},
 \quad K=I-DG,\quad G=\bar L^{-1}-J.                         \tag{K04.9}
\]
Here $J=\mathbf1\mathbf1^t$. To verify the upper-right block explicitly, it is
\[
 -(\bar A-de^t)\bar L^{-1}
 =I-D\bar L^{-1}+d\mathbf1^t=I-DG.
\]
Thus $[D\ K]$ is surjective, since $DG+K=I$.

It remains to show $K$ has primitive image. In the adjacency matrix, replace column $a$ by the sum of all columns; its first entry is $d_a=1$ and the remaining entries are $d$. Eliminate below this unit pivot. The residual matrix is $\bar A-de^t$. Multiplication on the right by $\bar L^{-1}$ gives $-K$. Hence
\[
 A\sim I_1\oplus K.                                         \tag{K04.10}
\]
By Lemma~\ref{K04-treefacts}, $K$ has $2\nu-1$ unit Smith factors and $\delta$ zeros.

\subsection{Completion of the proof}\label{K04-sec-5}
Apply Lemma~\ref{K04-coupling} with $q=m-1$, $r=2\nu-1$. It replaces $S$ by $I_r\oplus nI_{2\delta}\oplus n^2E$, where $\coker E\cong\coker D$. There are $\ell-1$ literal unit entries in the diagonal matrix $D$. Since $\delta\leq\ell-1$, remove $\delta$ of those entries to form an $r\times r$ diagonal matrix
\[
 E_0=I_{\ell-1-\delta}\oplus\diag(d_v:d_v>1).
\]
The nonsingular matrices $E$ and $E_0$ have the same size and isomorphic cokernels, so uniqueness of Smith normal form makes them integrally equivalent. The same equivalence persists after multiplication by $n^2$. Combine this with (K04.5) and (K04.9), and take torsion. This proves (K04.1).

The reasoning also specifies a certificate: a leaf, the inverse of its reduced tree Laplacian, Smith bases for the primitive matrix $K$, and a basis for the kernel of the surjective bottom block in (K04.3). All existence claims are ordinary integral basis extensions; no rational basis is treated as integral.

\subsection{Consequences and checks}\label{K04-sec-6}
For every tree and every $n\geq2$, (K04.1) gives
\[
 |K(T[n])|=n^{mn-2}\prod_vd_v^{n-1},\qquad
 \exp K(T[n])=n^2\lcm_vd_v,\qquad
 d(K(T[n]))=mn-2\nu-1,                                      \tag{K04.11}
\]
where $d(G)$ is the minimal number of generators. For the exponent, if there is a nonleaf its degree contributes a factor $n^2d_v$; the two-vertex tree instead has $h=1$. For generator number, every displayed factor is divisible by every prime dividing $n$, and there are exactly $mn-2\nu-1$ of them.

For a star $K_{1,s}$ one has $\nu=1$, $\delta=s-1$, $\ell=s$, $h=0$. Formula (K04.1) becomes the standard complete-bipartite answer for $K_{n,sn}$. For a tree with a perfect matching, $\delta=0$ and the last two terms are one $n^2d_v$ factor per vertex after omitting any leaf, recovering Round 2. Formula (K04.1) also specializes to the nonexceptional-prime classification in Round 3. Both earlier papers are strictly subsumed at the level of abstract groups.

Exact exploratory checks used 450 random recursive-tree instances with $2\leq m\leq16$, and blowup sizes cycling through $2,3,4$. The Smith forms of the $2m$-order cores agreed with (K04.1). These finite computations motivated but do not prove the theorem; the proof above applies to every tree.

\subsection{Literature and claim boundary}\label{K04-sec-7}
Dong, Guo and Jiang \cite{ref028} compute incomplete multipartite path-of-cluster groups and give explicit cases with two to six parts. Jacobson, Niedermaier and Reiner \cite{ref032} determine complete multipartite groups. Alfaro, Corrales and Valencia \cite{ref029} develop twin critical ideals. Searches of these primary sources and for tree-blowup critical groups did not establish an identical arbitrary-tree formula, but do not certify novelty. We claim a complete written proof of (K04.1), not formal verification, solution of a named pre-existing open problem, or any benchmark ranking. The formula concerns unweighted trees and equal independent cluster sizes; arbitrary cluster sizes, graph weights, and cyclic base graphs are not covered.




\clearpage
\section[K05 --- Smith Groups of Surjective Coupled Presentations]{Smith Groups of Surjective Coupled Presentations}\label{contrib-K05}
\markboth{Integral Presentations and Critical Groups of Graph Blowups}{K05 --- Keio}
\noindent\textsc{K05 \quad Keio programme, source round 5}\par\medskip


\begin{quote}\small\textbf{Overview.} We classify the cokernel of the integer block matrix $\left(\begin{smallmatrix}aD&K\\0&bI\end{smallmatrix}\right)$ whenever $D$ is nonsingular and $[D\ K]$ is surjective. The answer depends only on the separate Smith factors of $D$ and $K$, not on their relative integral bases. Explicit gcd factors handle every prime and unequal scales $a,b$. The primitive-coupling result of the previous round is a special case. As an application we classify uniform mixed clique/independent blowups of any tree possessing an independent leaf cluster, including looped adjacency matrices with nontrivial torsion.\end{quote}

\subsection{A classification from two separate Smith forms}\label{K05-sec-1}
For an integer matrix $M$, its Smith factors are taken nonnegative, in divisibility order, with zeros last. Write $\Z/1\Z=0$.
\begin{theorem}\label{K05-main}
Let $D,K\in\Z^{q\times q}$, with $D$ nonsingular, and suppose $[D\ K]:\Z^{2q}\to\Z^q$ is surjective. Let $d_1\mid\cdots\mid d_q$ be the positive Smith factors of $D$, and $k_1\mid\cdots\mid k_q$ those of $K$, allowing zeros. For positive integers $a,b$ put $g_i=\gcd(a,b,k_i)$, where $\gcd(a,b,0)=\gcd(a,b)$. Then
\[
 \coker\begin{pmatrix}aD&K\\0&bI_q\end{pmatrix}
 \cong\bigoplus_{i=1}^q\left(
 \Z/g_i\Z\ \oplus\ \Z/(ab\,d_{q+1-i}/g_i)\Z\right).          \tag{K05.1}
\]
In addition, the hypotheses imply
\[
 \gcd(k_i,d_{q+1-i})=1\quad(1\leq i\leq q).                  \tag{K05.2}
\]
Equation (K05.1) is a cyclic decomposition, not necessarily a globally sorted Smith list.
\end{theorem}
Thus no relative-position invariant of the two lattices is needed under the stated surjectivity hypothesis. That hypothesis is essential and is stronger than mere nonsingularity of the block matrix.

\subsection{Local reduction}\label{K05-sec-2}
Fix a prime $p$ and work over $R=\Z_{(p)}$. Write $\alpha=v_p(a)$ and $\beta=v_p(b)$. Let $s$ be the number of Smith factors of $K$ divisible by $p$, counting each zero as divisible, and let $r=q-s$. Choose invertible $U,V$ over $R$ so that
\[
 UKV=\diag(I_r,H),\qquad
 H=\diag(p^{\gamma_1},\ldots,p^{\gamma_s}),                    \tag{K05.3}
\]
where $\gamma_j\geq1$ or $\gamma_j=\infty$, with $p^\infty=0$. Unit parts are absorbed into $U,V$.

The bottom $s$ rows of $UD$ are surjective over $R$. Indeed modulo $p$, the bottom rows of $UKV$ vanish and surjectivity of $[D\ K]$ forces full row rank for that block of $UD$; a unit minor supplies a basis extension over the local ring. Hence an invertible $W$ gives
\[
 UDW=\begin{pmatrix}E&F\\0&I_s\end{pmatrix}.                 \tag{K05.4}
\]
The matrix $E$ is nonsingular and $\coker E\cong\coker D$ over $R$. In particular $D$ has at least $s$ unit Smith factors. This proves (K05.2), prime by prime, because the $p$-divisible $k_i$ have indices $i>q-s$ and the reversed $d_{q+1-i}$ have indices at most $s$.

\begin{lemma}\label{K05-local}
Over $R$, the block presentation in Theorem~\ref{K05-main} is equivalent to
\[
 I_r\ \oplus abE\ \oplus\bigoplus_{j=1}^s
       \begin{pmatrix}a&p^{\gamma_j}\\0&b\end{pmatrix}.       \tag{K05.5}
\]
\end{lemma}
\begin{proof}
Multiply on the left by $\diag(U,V^{-1})$ and on the right by $\diag(W,V)$. In block sizes $(r,s,r,s)$ the matrix becomes
\[
 \begin{pmatrix}
 aE&aF&I_r&0\\
 0&aI_s&0&H\\
 0&0&bI_r&0\\
 0&0&0&bI_s
 \end{pmatrix}.                                               \tag{K05.6}
\]
Eliminate the unit block. The first remaining block row is $(-abE,-abF,0)$. Adding $bF$ times the row block $(0,aI_s,H)$ removes the middle block and introduces $bFH$ in the last block. Subtracting $FH$ times the row block $(0,0,bI_s)$ removes it. This leaves (K05.5), after signs and permutations. Every coefficient is in $R$; no division by $a$ or $b$ occurs.
\end{proof}

\subsection{Proof of the global formula}\label{K05-sec-3}
A $2\times2$ block in (K05.5) has first Smith valuation
\[
 t_j=\min(\alpha,\beta,\gamma_j)
\]
and determinant valuation $\alpha+\beta$, so its second valuation is $\alpha+\beta-t_j$. If $\gamma_j=\infty$, this is just the diagonal block $(a,b)$.

Write the valuations of $D$ as $\eta_1\leq\cdots\leq\eta_q$. The first $s$ are zero by (K05.4), and those of $E$ are $\eta_{s+1},\ldots,\eta_q$. The $abE$ part therefore contributes valuations $\alpha+\beta+\eta_{s+1},\ldots,\alpha+\beta+\eta_q$. For indices with $p\nmid k_i$, $v_p(g_i)=0$, and the reversed $d$ indices are precisely $s+1,\ldots,q$. For the remaining indices the reversed $d$ valuations are zero and $v_p(g_i)=\min(\alpha,\beta,v_p(k_i))$. Thus the $p$-primary factors of (K05.1) agree exactly with (K05.5). This holds for every prime, proving Theorem~\ref{K05-main} by the classification of finite abelian groups.

Two useful checks are
\[
 |\coker S|=(ab)^q|\det D|,
\]
and, when $\gcd(a,b)=1$, the simpler expression
\[
 \coker S\cong\bigoplus_i\Z/(ab\,d_i)\Z.                    \tag{K05.7}
\]
The coupling is erased by coprime scales, despite possibly large torsion in $\coker K$.

Surjectivity cannot be omitted. Take $D=2,K=2,a=1,b=2$. The actual matrix $\left(\begin{smallmatrix}2&2\\0&2\end{smallmatrix}\right)$ has Smith factors $2,2$, whereas (K05.1) would predict $1,4$. Here $[D\ K]$ has image $2\Z$, not $\Z$.

\subsection{Mixed blowups of a tree}\label{K05-sec-4}
Let $T$ be a tree, and choose $c_v\in\{0,1\}$. Replace vertex $v$ by an independent set of size $n$ if $c_v=0$, or a clique of size $n$ if $c_v=1$. Connect two clusters completely when their base vertices are adjacent. Assume at least one leaf $a$ has $c_a=0$. Put
\[
 e_v=d_v+c_v,\quad B=A_T+\diag(c_v),\quad L=L_T.
\]
The graph Laplacian has the integral reduction
\[
 L_{\rm mixed}\sim\bigoplus_v ne_vI_{n-2}\oplus
 \begin{pmatrix}n\diag(e_v)&-B\\0&nL\end{pmatrix}.             \tag{K05.8}
\]
This follows by subtracting reference generators in each cluster and isolating their sum, exactly as for independent twins.

Delete the independent leaf $a$ and write $E=\diag(e_v:v\ne a)$, $\bar L=L\setminus a$, and $G=\bar L^{-1}-J$. The leaf has $e_a=1$, so the sum of the second-family relations eliminates its difference generator with unit coefficient. The remaining free reference generator and redundant relation split off. The remaining finite block is
\[
 S=\begin{pmatrix}nE&K\\0&nI\end{pmatrix},\qquad K=I-EG.     \tag{K05.9}
\]
For a direct verification, $\bar B=E-\bar L$ and $\bar L^{-1}e_b=\mathbf1$, where $b$ is the leaf's neighbor. These identities give the upper-right block $I-E(\bar L^{-1}-J)$. Moreover $EG+K=I$, so $[E\ K]$ is surjective. The same leaf-column pivot in $B$ yields
\[
 B\sim I_1\oplus K.                                         \tag{K05.10}
\]
All required matrices are integral because a reduced tree Laplacian has determinant one.

\begin{corollary}\label{K05-mixed}
Let $d_1\mid\cdots\mid d_{m-1}$ be the Smith factors of $E$, and let $k_1\mid\cdots\mid k_{m-1}$ be the Smith factors of $B$ after removing one initial unit factor. Put $g_i=\gcd(n,k_i)$. Then the mixed blowup has critical group
\[
 \bigoplus_v(\Z/(ne_v)\Z)^{n-2}\ \oplus
 \bigoplus_{i=1}^{m-1}\left(\Z/g_i\Z\oplus
       \Z/(n^2d_{m-i}/g_i)\Z\right).                         \tag{K05.11}
\]
\end{corollary}
This classifies every blowup size using two Smith invariants of the \emph{fixed base} tree and its chosen clique pattern. Unlike the independent case, $B$ may have nonunit nonzero Smith factors. No primitive-image assumption remains.

For example, take a four-vertex star whose four clusters are cliques. Attach a path of length two to one star leaf; let the new intermediate cluster and terminal leaf be independent. The terminal leaf supplies the required root. Eliminating its adjacent pair from $B$ leaves $I+A_{K_{1,3}}$, whose determinant is $-2$; its Smith factors are $1,1,1,2$. Thus the full $B$ has factors $1,1,1,1,1,2$. The matrix $E$ has diagonal entries $4,3,2,2,2$ and Smith factors $1,2,2,2,12$. Formula (K05.11) introduces a genuine $\gcd(n,2)$ transition, absent from the independent-tree theorem.

\subsection{Scope and sources}\label{K05-sec-5}
Smith-form background is standard \cite{ref030}. Twin-vertex critical ideals \cite{ref029} and path-cluster groups \cite{ref028} are relevant graph predecessors. The new claim here is the complete coupled-presentation formula under the precise surjectivity hypothesis, and its mixed-tree application. Searches did not establish an identical theorem, but priority is unverified. No named open problem is claimed solved. Exact exploratory checks of (K05.1) used 180 instances with $1\leq q\leq6$, $K=I-DG$, and integer diagonal $D$; all agreed with direct Smith forms. Those checks are supplementary to the local proof.




\clearpage
\section[K06 --- Potentials and Exact Exponents of Nonuniform Tree Blowups]{Potentials and Exact Exponents of Nonuniform Tree Blowups}\label{contrib-K06}
\markboth{Integral Presentations and Critical Groups of Graph Blowups}{K06 --- Keio}
\noindent\textsc{K06 \quad Keio programme, source round 6}\par\medskip


\begin{quote}\small\textbf{Overview.} For an arbitrary tree whose vertices are replaced by independent clusters of arbitrary sizes at least two, we give a closed formula for the canonical monodromy pairing, a denominator certificate for the order of every divisor class, and an exact least-common-multiple formula for the exponent of the critical group. The exponent involves only internal vertices and internal edges: apparent denominators from leaf edges cancel. Explicit two-vertex divisors certify every factor in the formula. No Smith-form computation or coprimality assumption is required.\end{quote}

\subsection{Statement and conventions}\label{K06-sec-1}
Let $T$ be a tree on $m\geq3$ vertices. Replace each vertex $i$ by an independent cluster $V_i$ of size $n_i\geq2$, and each edge by all edges between its two clusters. Denote the resulting connected graph by $G$. Put
\[
 D_i=\sum_{j\sim i}n_j,\quad I=\{i:\deg_T(i)>1\},\quad
 E_I=\{ij\in E(T):i,j\in I\}.
\]
A divisor is an integer vector indexed by vertices of $G$; $\operatorname{Div}^0(G)$ consists of the vectors with coordinate sum zero. If $L_G$ is the Laplacian, its critical group is
\[
 K(G)=\operatorname{Div}^0(G)/L_G\Z^{V(G)}.
\]
For a finite abelian group, its exponent is the least positive integer annihilating it.
\begin{theorem}\label{K06-exponent}
For every such nonuniform tree blowup,
\[
 \boxed{\ \exp K(G)=\lcm\left(
       \{n_iD_i:i\in I\}\ \cup\ \{n_in_j:ij\in E_I\}
                         \right).\ }                          \tag{K06.1}
\]
Each integer in the displayed set divides the order of an explicitly specified divisor supported on two vertices of $G$.
\end{theorem}
The internal-vertex set is nonempty. The two-vertex base tree is exceptional only in notation: $G=K_{n_1,n_2}$ has exponent $n_1n_2$, as proved below. Formula (K06.1) does not list individual leaf-edge products $n_in_j$. In general those products do not divide the exponent.

\subsection{An exact potential for every divisor}\label{K06-sec-2}
For $z\in\operatorname{Div}^0(G)$ write
\[
 S_i(z)=\sum_{v\in V_i}z_v.
\]
Orient the edges of $T$ away from a chosen root. For an oriented edge $e=(i,j)$ let $T_j$ be the component containing $j$ after deleting $e$, and define
\[
 F_e(z)=\sum_{k\in T_j}S_k(z).
\]
Define rational numbers $Z_i$, up to a common constant, by
\[
 Z_j-Z_i=\frac{F_e(z)}{n_in_j}\qquad(e=(i,j)).                \tag{K06.2}
\]
Because $T$ has no cycles, these equations determine $Z$ consistently. Their signed sum at each vertex says
\[
 \sum_{j\sim i}n_in_j(Z_i-Z_j)=S_i(z).                       \tag{K06.3}
\]
\begin{lemma}\label{K06-potential}
The vector
\[
 x_v=Z_i+\frac{z_v}{D_i}-\frac{S_i(z)}{n_iD_i},
                 \qquad v\in V_i,                           \tag{K06.4}
\]
satisfies $L_Gx=z$. If $\ell$ is a leaf with neighbor $p$, then it simplifies to
\[
 x_v=Z_p+\frac{z_v}{n_p}\qquad(v\in V_\ell).                 \tag{K06.5}
\]
\end{lemma}
\begin{proof}
Summing (K06.4) over $V_i$ gives $\sum_{v\in V_i}x_v=n_iZ_i$. Thus
\[
 (L_Gx)_v=D_ix_v-\sum_{j\sim i}n_jZ_j
 =z_v+D_iZ_i-\sum_{j\sim i}n_jZ_j-S_i(z)/n_i=z_v
\]
by (K06.3). At a leaf, (K06.3) reads
$Z_\ell-Z_p=S_\ell(z)/(n_\ell n_p)$. Since $D_\ell=n_p$, substitution cancels the two terms containing $S_\ell(z)$ and proves (K06.5).
\end{proof}

\begin{proposition}\label{K06-order}
For any choice of a vertex $o$ of $G$,
\[
 \ord([z])=\lcm_{v\in V(G)}\operatorname{den}(x_v-x_o),       \tag{K06.6}
\]
where $x$ is given by (K06.2)--(K06.4), and $\operatorname{den}$ is the positive reduced denominator, with $\operatorname{den}(0)=1$.
\end{proposition}
\begin{proof}
The equality $tz=L_Gu$ for an integer vector $u$ holds exactly when $L_G(tx-u)=0$. Connectedness gives $\ker_{\Q}L_G=\Q\boldsymbol1$. Hence $tz$ is principal exactly when all differences $t(x_v-x_o)$ are integers. This is precisely (K06.6).
\end{proof}
This is an exact order certificate involving only rational arithmetic on the base tree. It applies to arbitrary divisors, not only to the special divisors used to prove the exponent.

\subsection{The paired group, with leaf cancellation}\label{K06-sec-3}
The canonical monodromy pairing is
\[
 \langle[z],[w]\rangle=x^{\mathsf T}w\pmod{\Z},\qquad L_Gx=z.
\]
Changing $x$ by a constant has no effect because $w$ has degree zero. Changing either divisor by a principal divisor changes the answer by an integer. The pairing is symmetric because the Laplacian is symmetric.
\begin{theorem}\label{K06-pairing}
Choose any orientations on the internal edges, and use the associated component sums $F_e$. Then
\begin{align*}
\langle[z],[w]\rangle={}&
 \sum_{ij\in E_I}\frac{F_{ij}(z)F_{ij}(w)}{n_in_j}\\
 &+\sum_{i\in I}\left(
 \frac{\sum_{v\in V_i}z_vw_v}{D_i}
 -\frac{S_i(z)S_i(w)}{n_iD_i}\right)\\
 &+\sum_{\ell\text{ leaf}}
       \frac{\sum_{v\in V_\ell}z_vw_v}{n_{p(\ell)}}
                 \pmod{\Z}.                                 \tag{K06.7}
\end{align*}
This pairing is nondegenerate.
\end{theorem}
\begin{proof}
Substituting (K06.4) in $x^{\mathsf T}w$ gives the within-cluster terms plus $\sum_iZ_iS_i(w)$. Discrete summation by parts using (K06.3) transforms the last expression into
\[
 \sum_{ij\in E(T)}n_in_j(Z_i-Z_j)(W_i-W_j)
 =\sum_{ij\in E(T)}\frac{F_{ij}(z)F_{ij}(w)}{n_in_j},
\]
where $W$ is the potential associated with $w$. On a leaf edge the summand is $S_\ell(z)S_\ell(w)/(n_\ell n_p)$, which cancels the negative within-cluster term because $D_\ell=n_p$. This proves (K06.7). If $\langle[z],[w]\rangle=0$ for all degree-zero integer $w$, take $w=e_v-e_o$. All potential differences are then integers, and Proposition~\ref{K06-order} implies $[z]=0$. This proves nondegeneracy without invoking a general duality theorem.
\end{proof}

\subsection{Proof of the exact exponent}\label{K06-sec-4}
Let $E$ denote the right side of (K06.1), and root the base tree at an internal vertex $r$, taking $Z_r=0$. The internal vertices induce a connected subtree: a path between two internal vertices cannot pass through a leaf. For every internal $i$, equation (K06.2) therefore expresses $Z_i$ as a sum with denominators $n_un_v$ only from internal edges. Consequently $EZ_i\in\Z$. Equation (K06.4) gives $Ex_v\in\Z$ in every internal cluster. For a leaf, (K06.5) gives the same conclusion: its neighbor $p$ is internal and $n_p\mid n_pD_p\mid E$. Thus $E$ annihilates every class by Proposition~\ref{K06-order}.

We prove the reverse divisibility factor by factor. Fix $i\in I$. Choose a neighbor $j$ and a leaf $\ell$ in a different component of $T-i$; such choices exist since $i$ has at least two neighbors, and every component of $T-i$ contains a leaf of $T$. Choose a vertex $v\in V_i$ and a vertex $w\in V_\ell$, and set $z=e_v-e_w$. In the component through $j$ every cluster sum is zero, so (K06.2) implies $Z_j=Z_i$. Take another vertex $v'\in V_i$, possible since $n_i\geq2$, and any vertex $u\in V_j$. Neither $u$ nor its cluster is charged. Equations (K06.4)--(K06.5) yield
\[
 x_u=Z_i,\qquad x_{v'}=Z_i-\frac1{n_iD_i}.
\]
Proposition~\ref{K06-order} shows $n_iD_i\mid\ord([e_v-e_w])$.

Next fix an internal edge $ij$. Each component of $T-ij$ contains a leaf of $T$ distinct from $i,j$. Choose one leaf on each side, and put one positive chip at a vertex over the first and one negative chip over the second. Then $S_i=S_j=0$ and the flow across $ij$ has magnitude one. Every clone in $V_i$ has potential $Z_i$ and every clone in $V_j$ has potential $Z_j$, so their potential difference is $\pm1/(n_in_j)$. Hence $n_in_j$ divides the order of this two-vertex class. All factors in (K06.1) divide the exponent, completing the proof.

For $m=2$, choose source $v\in V_1$ and sink $w\in V_2$. Set $Z_2=0$ and $Z_1=1/(n_1n_2)$. An uncharged clone in $V_1$ has potential zero, and an uncharged clone in $V_2$ has potential $1/(n_1n_2)$. Thus the class order is divisible by $n_1n_2$. Equation (K06.4) also bounds all denominators by $n_1n_2$, proving the exceptional formula.

\subsection{Examples and scope}\label{K06-sec-5}
For a star, let the center have size $a$ and let the leaf sizes sum to $B$. The blowup is $K_{a,B}$, and (K06.1) gives $aB$, independently of the partition of $B$ among leaf clusters. This explains why retaining the separate leaf-edge products would be wrong. For instance $a=2$ and leaf sizes $2,3$ give exponent $10$, not a multiple of the leaf-edge product $6$.

For a four-vertex path with sizes $(a,b,c,d)$, the formula is
\[
 \exp K(G)=\lcm\bigl(b(a+c),\ c(b+d),\ bc\bigr).             \tag{K06.8}
\]
At $(a,b,c,d)=(2,3,4,5)$ this is $\lcm(18,32,12)=288$. The internal-edge term is essential: at $(a,b,c,d)=(3,2,2,3)$ the two internal-vertex terms both equal $10$, whereas the internal-edge term is $4$, giving exponent $20$ rather than $10$.

For uniform size $n$, (K06.1) reduces to $n^2\lcm_i\deg_T(i)$, consistent with the full abstract-group classification of Round 4. Here neither sizes nor degrees need be coprime, and they may vary independently. Exact Smith computations on 275 randomly generated tree cores, with $2\leq m\leq12$ and sizes from $2$ through $8$, were supplementary checks, not steps of the proof.

\subsection{Relation to existing work and limitations}\label{K06-sec-6}
The monodromy pairing and its general computational use are established theory; see Shokrieh \cite{ref033}. Complete bipartite and multipartite critical groups are classical \cite{ref032}. The contribution here is the explicit leaf-cancelled pairing and the closed exponent formula for arbitrary nonuniform independent-cluster tree blowups, with lower-bound certificates for every factor. This is not a complete Smith classification in the nonuniform setting. It is not claimed to settle a previously named open problem, and bibliographic priority has not been certified.




\clearpage
\section[K07 --- Two Independent Scales in the Critical Groups of Tree Blowups]{Two Independent Scales in the Critical Groups of Tree Blowups}\label{contrib-K07}
\markboth{Integral Presentations and Critical Groups of Graph Blowups}{K07 --- Keio}
\noindent\textsc{K07 \quad Keio programme, source round 7}\par\medskip


\begin{quote}\small\textbf{Overview.} Let a tree have bipartition $U\sqcup V$. Replace every vertex in $U$ by $a$ independent twins and every vertex in $V$ by $b$ independent twins, where $a,b\geq2$ are arbitrary. We determine the entire critical group, including primes simultaneously dividing both scales and vertex degrees. The answer uses the degree multiset, the two color-class sizes, and the matching number. The proof counts kernels over every finite local ring in four valuation ranges. This supplies a two-parameter classification rather than an inference from the equal-scale case.\end{quote}

\subsection{The complete group}\label{K07-sec-1}
Let $T$ be a tree on $m\geq2$ vertices, with bipartition $U\sqcup V$, $r=|U|$, $s=|V|$, degrees $d_i$, number of leaves $\ell$, and maximum matching size $\nu$. Write
\[
 \delta_U=r-\nu,\quad \delta_V=s-\nu,\quad
 \delta=m-2\nu,\quad h_0=\ell-1-\delta.
\]
Let $T[a,b]$ replace vertices in $U$ by independent clusters of size $a\geq2$ and vertices in $V$ by independent clusters of size $b\geq2$, with complete adjacency along each base edge.
\begin{theorem}\label{K07-main}
One has $h_0\geq0$, and the critical group is
\begin{align*}
 K(T[a,b])\cong{}&
 \bigoplus_{u\in U}(\Z/(bd_u)\Z)^{a-2}
 \oplus\bigoplus_{v\in V}(\Z/(ad_v)\Z)^{b-2}\\
 &\oplus(\Z/b\Z)^{2\delta_U}
 \oplus(\Z/a\Z)^{2\delta_V}
 \oplus(\Z/(ab)\Z)^{h_0}
 \oplus\bigoplus_{d_i>1}\Z/(ab d_i)\Z.                      \tag{K07.1}
\end{align*}
There are no coprimality assumptions. The displayed cyclic factors need not be in divisibility order.
\end{theorem}
In particular, the assignments of nonleaf degrees to the two colors matter in the first line, while the final large factors use the uncolored degree multiset. Unequal scales distinguish the two matching deficiencies.

\subsection{Integral preparation}\label{K07-sec-2}
Write the adjacency matrix and degree matrix in bipartite order as
\[
 A=\begin{pmatrix}0&B\\B^{\mathsf T}&0\end{pmatrix},\quad
 D=\diag(d_i),\quad L=D-A,\quad
 R=\diag(bI_r,aI_s),\quad H=\diag(aI_r,bI_s).
\]
Thus $RH=abI$, $HA=AR$, and $RD-HA=LR$.
\begin{lemma}\label{K07-tree}
The image of $B:\Z^s\to\Z^r$ is primitive and its rank is $\nu$. Both $[D_U\ B]$ and $[D_V\ B^{\mathsf T}]$ are surjective over $\Z$. Over every ring $\Z/p^k\Z$, the kernel of $L$ is precisely the constant vectors.
\end{lemma}
\begin{proof}
Every minor of $B$ is zero or $\pm1$: two distinct nonzero determinant terms would give two perfect matchings on the selected rows and columns, whose symmetric difference contains a cycle, impossible in a forest. A maximum matching supplies a nonzero $\nu$-minor, and no larger minor has a determinant term. This proves the first assertion.

Every reduced tree Laplacian has determinant one: a reduced oriented incidence matrix has determinant $\pm1$ by leaf elimination, and its product with its transpose is the reduced Laplacian. Therefore the image of $L$ over $\Z$ is the lattice of vectors with coordinate sum zero, and its kernel over any $\Z/p^k\Z$ consists of constants. The image of $[D\ A]$ contains the image of $L=D-A$, and it contains $e_\ell$ for any leaf $\ell$, because $d_\ell=1$. Hence $[D\ A]$ is surjective. Its two color blocks are exactly the two surjections in the statement.
\end{proof}
To see $h_0\geq0$, choose a maximum matching with as few unmatched nonleaves as possible. An alternating path starting at an unmatched nonleaf can be continued until it ends at a matched leaf: an unmatched endpoint at an odd step would augment the matching, and a tree prevents revisiting vertices. Flipping the path contradicts the chosen minimality. Thus every unmatched vertex is a leaf. At least one leaf is matched, since otherwise the same alternating-path procedure starting at a leaf could not terminate. Consequently $\delta\leq\ell-1$.

Subtracting one chosen generator from the others in each cluster, and subtracting the corresponding relations, gives scalar relations $bd_u$ or $ad_v$ on the difference generators. Replace one difference by their sum; the other $a-2$ or $b-2$ differences split off. The remaining two generators per cluster have presentation
\[
 C=\begin{pmatrix}RD&-A\\0&LR\end{pmatrix},                 \tag{K07.2}
\]
and therefore
\[
 L_{T[a,b]}\sim
 \bigoplus_{u\in U}bd_u I_{a-2}\ \oplus
 \bigoplus_{v\in V}ad_v I_{b-2}\ \oplus C.                  \tag{K07.3}
\]
These are integer basis operations, not division by cluster sizes. Equivalently, for a general block matrix with diagonal blocks $q_iI_{n_i}$ and constant block perturbation $b_{ij}J$, the same operations leave the core $\left(\begin{smallmatrix}Q&B_0\\0&Q+N B_0\end{smallmatrix}\right)$, with $Q=\diag(q_i)$ and $N=\diag(n_i)$; here $B_0=-A$ and $Q=RD$.

\subsection{Two local counting facts}\label{K07-sec-3}
Fix a prime $p$. Work over $\Z_{(p)}$, and put $\alpha=v_p(a)$, $\beta=v_p(b)$. Relabel the colors if needed so that $\alpha\leq\beta$. Unit changes of variables reduce $a,b$ to $p^\alpha,p^\beta$. For completeness, if $a=up^\alpha$, $b=vp^\beta$ and $W=\diag(vI_r,uI_s)$, multiply (K07.2) on the left by $\diag(uvW^{-1},I)$ and on the right by $\diag((uv)^{-1}I,W^{-1})$. The result is (K07.2) with the unit parts removed.

We use two elementary consequences of Lemma~\ref{K07-tree}.
\begin{lemma}\label{K07-count}
Let $D_0$ be a diagonal $q\times q$ matrix and $B_0$ a $q\times t$ matrix of primitive rank $\rho$, with $[D_0\ B_0]$ surjective.
\begin{enumerate}
\item The nonzero Smith factors of $[p^cD_0\ B_0]$ are $1$ repeated $\rho$ times and $p^c$ repeated $q-\rho$ times.
\item The index of the image of $[D_0\ p^cB_0]$ is
$p^{\sum_{i=1}^q\min(v_p((D_0)_{ii}),c)}$.
\end{enumerate}
\end{lemma}
\begin{proof}
For the first assertion, put $B_0$ in primitive Smith form. Surjectivity says that the projection of $D_0$ onto the quotient by $\im B_0$ is surjective. The image of the displayed matrix is therefore all of $\im B_0$ together with $p^c$ times its free complementary lattice. For the second assertion,
\[
 D_0\Z_{(p)}^q+p^cB_0\Z_{(p)}^t
 =D_0\Z_{(p)}^q+p^c\Z_{(p)}^q,
\]
since $[D_0\ B_0]$ is surjective. The diagonal entries give the stated index.
\end{proof}
If a matrix with $N$ columns and full row rank $q$ has Smith valuations $\gamma_i$, its kernel over $\Z/p^k\Z$ has size $p^{(N-q)k+\sum_i\min(\gamma_i,k)}$. This follows by counting one scalar congruence at a time.

\subsection{The four kernel ranges}\label{K07-sec-4}
Let
\[
 \kappa(k)=\log_p|\ker(C\bmod p^k)|,\qquad
 Q(c)=\sum_{i=1}^m\min(v_p(d_i),c).
\]
We claim the following complete formula:
\[
\kappa(k)=
\begin{cases}
 (m+\delta)k,&0<k\leq\alpha,\\
 2rk+2\alpha\delta_V,&\alpha<k\leq\beta,\\
 2\nu k+2\beta\delta_U+2\alpha\delta_V,
       &\beta<k\leq\alpha+\beta,\\
 c+2r\beta+2s\alpha+Q(c),\quad c=k-\alpha-\beta,
       &k>\alpha+\beta.
\end{cases}                                                   \tag{K07.4}
\]
Empty ranges cause no difficulty, including when one or both valuations are zero.

A kernel vector $(x,y)$ satisfies
\[
 RDx=Ay,\qquad LRy=0.                                        \tag{K07.5}
\]
By Lemma~\ref{K07-tree}, $Ry=t\boldsymbol1$ for some scalar $t$ modulo $p^k$. Multiplying the first equation by $H$ gives the crucial constraint
\[
 D(abx-t\boldsymbol1)=0.                                    \tag{K07.6}
\]
Since a leaf degree is one, $t$ is a multiple of $ab$ modulo $p^k$. Thus $t=0$ when $k\leq\alpha+\beta$; when $k>\alpha+\beta$, it has precisely $p^{k-\alpha-\beta}$ possible values.

\subsubsection*{Range 1: $k\leq\alpha$}
Here $R=0$ modulo $p^k$. The first equation is $Ay=0$, while $x$ is arbitrary. The matrix $A$ has $2\nu$ unit Smith factors and $\delta$ zeros, so $\kappa(k)=(m+\delta)k$.

\subsubsection*{Range 2: $\alpha<k\leq\beta$}
Now $t=0$, $y_U$ is arbitrary and $y_V=p^{k-\alpha}z_V$, with $z_V$ taken modulo $p^\alpha$. The $U$ equation requires $Bz_V=0$ modulo $p^\alpha$, giving $p^{\alpha\delta_V}$ choices. The $V$ equation is
\[
 p^\alpha D_Vx_V-B^{\mathsf T}y_U=0\pmod{p^k}.
\]
By Lemma~\ref{K07-count}(K07.1), its kernel has $p^{rk+\alpha\delta_V}$ elements. Finally $x_U$ is arbitrary, contributing $p^{rk}$. Their product proves the second line of (K07.4).

\subsubsection*{Range 3: $\beta<k\leq\alpha+\beta$}
Again $t=0$, so write $y_U=p^{k-\beta}z_U$ and $y_V=p^{k-\alpha}z_V$, with $z_U$ modulo $p^\beta$ and $z_V$ modulo $p^\alpha$. Put $c=\alpha+\beta-k$. Dividing the two color equations by their common powers of $p$ gives
\[
 p^cD_Ux_U-Bz_V=0\pmod{p^\alpha},\qquad
 p^cD_Vx_V-B^{\mathsf T}z_U=0\pmod{p^\beta}.
\]
Here $0\leq c\leq\alpha\leq\beta$. The first congruence has $p^{s\alpha+c\delta_U}$ solutions for its reduced variables, and $x_U$ has $p^{r(k-\alpha)}$ lifts. The second similarly contributes $p^{r\beta+c\delta_V+s(k-\beta)}$. Adding exponents gives
\[
 mk+(r-s)(\beta-\alpha)+c\delta
 =2\nu k+2\beta\delta_U+2\alpha\delta_V,
\]
as claimed.

\subsubsection*{Range 4: $k>\alpha+\beta$}
Put $c=k-\alpha-\beta>0$. There are $p^c$ choices of $t$. For each, choose $u$ with $t=ab u$ and subtract the solution $(u\boldsymbol1,Hu\boldsymbol1)$ from $(x,y)$. This bijects its fiber with the fiber $t=0$. In that fiber write $y_U,y_V$ as in Range 3. The $U$ equation becomes
\[
 D_Ux_U-p^cBz_V=0\pmod{p^{\alpha+c}},\qquad
 z_V\in(\Z/p^\alpha\Z)^s.                                  \tag{K07.7}
\]
Temporarily allow $z_V$ modulo $p^{\alpha+c}$. Lemma~\ref{K07-count}(K07.2) shows that the kernel has exponent $s(\alpha+c)+\sum_{u\in U}\min(v_p(d_u),c)$. Each original $z_V$ is counted $p^{cs}$ times. The lifts of $x_U$ from modulus $p^{\alpha+c}$ to $p^k$ contribute $p^{r\beta}$. Thus (K07.7) contributes
\[
 p^{s\alpha+r\beta+\sum_{u\in U}\min(v_p(d_u),c)}.
\]
The $V$ equation similarly contributes
$p^{r\beta+s\alpha+\sum_{v\in V}\min(v_p(d_v),c)}$.
Multiplication and the $p^c$ choices of $t$ prove the last line of (K07.4).

\subsection{Reading the group from the kernel counts}\label{K07-sec-5}
A matrix with one zero Smith factor and nonzero valuations $e_j$ satisfies
\[
 \kappa(k)=k+\sum_j\min(k,e_j).
\]
The differences $\kappa(k)-\kappa(k-1)$ determine the number of valuations at least $k$, and hence the entire primary group. Let $E_0$ be the diagonal matrix containing all nonleaf degrees together with $h_0$ entries equal to one. Its size is $2\nu-1$. The candidate primary decomposition for the core is
\[
 \Z\ \oplus(\Z/p^\beta\Z)^{2\delta_U}
 \oplus(\Z/p^\alpha\Z)^{2\delta_V}
 \oplus\bigoplus_{e\in\diag E_0}\Z/p^{\alpha+\beta+v_p(e)}\Z.
                                                                    \tag{K07.8}
\]
Its kernel counts agree with the first three lines of (K07.4) by direct substitution. In the last range they equal
\[
 k+2\beta\delta_U+2\alpha\delta_V
 +(\alpha+\beta)(2\nu-1)+Q(c)
 =c+2r\beta+2s\alpha+Q(c).
\]
Thus (K07.8) is the actual core group at every prime. Assembling primary parts and adjoining the scalar factors in (K07.3) proves (K07.1).

\subsection{Consequences, checks, and boundary}\label{K07-sec-6}
If the tree has a perfect matching, both deficiencies vanish and the final line of (K07.1) consists of $ab d_i$ for all vertices except any one leaf. If $a=b=n$, (K07.1) becomes the complete uniform classification. For a star whose center lies in $U$, one has $r=1$, $s=t$, $\nu=1$; (K07.1) reduces to the classical critical group of $K_{a,bt}$.

The order obtained from (K07.1) is
\[
 |K(T[a,b])|=b^{ar-1}a^{bs-1}
       \prod_{u\in U}d_u^{a-1}\prod_{v\in V}d_v^{b-1}.       \tag{K07.9}
\]
The group exponent is $ab\lcm_i d_i$, including the two-vertex tree. These are consistency consequences, not substitutes for the primary-group proof.

Before the proof was found, 280 exact core Smith calculations on random recursive trees with $2\leq m\leq15$ and scale pairs $(2,3),(2,4),(3,6),(4,6)$ agreed with (K07.1). The four-range proof is independent of that finite evidence.

The block compression and complete multipartite special cases have predecessors \cite{ref028,ref032}; Smith-form background is surveyed in \cite{ref030}. The new claim here is a complete written two-scale arbitrary-tree classification with a finite-ring proof. No named open problem is claimed solved, and an identical prior theorem has not been ruled out by an exhaustive literature review. Fully arbitrary vertex-dependent sizes remain outside this abstract-group formula, although Round 6 gives their pairing and exact exponent.




\clearpage
\section[K08 --- Pendant Twin Attachments Transfer Critical Groups to Weighted Bases]{Pendant Twin Attachments Transfer Critical Groups to Weighted Bases}\label{contrib-K08}
\markboth{Integral Presentations and Critical Groups of Graph Blowups}{K08 --- Keio}
\noindent\textsc{K08 \quad Keio programme, source round 8}\par\medskip
\noindent\textit{Integration note.} The finalized source already reports 120 completed verification instances with base size at most six. The earlier count of 200 is not used.\par\medskip


\begin{quote}\small\textbf{Overview.} We prove an exact integral transfer theorem for graphs built from an arbitrary connected base: each base vertex is replaced by an independent set or clique, and receives its own independent pendant cluster. All cluster sizes may vary. The critical group splits into explicit local cyclic groups and the critical group of the base with edge multiplicities equal to products of core-cluster sizes. For constant core size this determines the answer directly from the original base critical group. For tree bases it gives a complete closed formula with no residual Smith computation.\end{quote}

\subsection{Graph construction and theorem}\label{K08-sec-1}
Let $H$ be a connected simple graph on vertices $1,\ldots,m$, allowing $m=1$. Choose integers $n_i,b_i\geq2$ and indicators $c_i\in\{0,1\}$. Build a graph $G$ as follows:
\begin{itemize}
\item Vertex $i$ becomes a core cluster $V_i$ of size $n_i$, independent if $c_i=0$ and complete if $c_i=1$.
\item For each edge $ij$ of $H$, join every vertex of $V_i$ to every vertex of $V_j$.
\item Add an independent cluster $W_i$ of size $b_i$, adjacent to every vertex of $V_i$ and to no other vertices.
\end{itemize}
Set
\[
 q_i=b_i+\sum_{j\sim i}n_j+c_in_i.
\]
Let $H_{\mathbf n}$ be the weighted multigraph with vertex set $V(H)$ and $n_in_j$ parallel edges for every edge $ij$ of $H$. Its Laplacian is
\[
 L_w=\diag\left(n_i\sum_{j\sim i}n_j\right)-NAN,
 \qquad N=\diag(n_i),                                       \tag{K08.1}
\]
where $A$ is the adjacency matrix of $H$. Its critical group is the torsion of $\coker L_w$; for $m=1$ it is trivial.
\begin{theorem}\label{K08-transfer}
There is an abstract-group isomorphism
\begin{align*}
 K(G)\cong K(H_{\mathbf n})\ \oplus
 \bigoplus_{i=1}^m\bigl(& (\Z/q_i\Z)^{n_i-2}
 \oplus(\Z/n_i\Z)^{b_i-2}\\
 &\oplus\Z/(n_iq_i)\Z\bigr).                               \tag{K08.2}
\end{align*}
All parameters are unrestricted apart from the stated positivity and graph assumptions. In particular, no prime is excluded.
\end{theorem}
The right side is explicit once the weighted-base critical group is known. This does not claim a closed formula for the critical group of every arbitrary weighted graph. Rather, it separates the entire attachment contribution by an integral direct sum, including all possible extension interactions.

\subsection{Two elementary integral reductions}\label{K08-sec-2}
We first record the matrix operations used in the proof. Matrices are equivalent if related by invertible integer row and column transformations.
\begin{lemma}\label{K08-compress}
Suppose a block matrix has blocks $d_iI_{t_i}\delta_{ij}+b_{ij}J_{t_i,t_j}$, with all $t_i\geq2$. Put $D=\diag(d_i)$, $B=(b_{ij})$, $T=\diag(t_i)$. Then it is equivalent to
\[
 \bigoplus_i d_iI_{t_i-2}\ \oplus
 \begin{pmatrix}D&B\\0&D+TB\end{pmatrix}.                   \tag{K08.3}
\]
\end{lemma}
\begin{proof}
In a presentation use one distinguished generator $x_i$ from block $i$ and differences $u_{ij}=x_{ij}-x_i$ for the other generators. Subtracting the distinguished relation from every other relation in that block gives $d_i u_{ij}=0$. The remaining relation depends on these differences only through their sum. Replacing one difference by the sum of all $t_i-1$ differences is unimodular. The other $t_i-2$ scalar generators split off. Transposing the remaining presentation if necessary gives (K08.3); an integer matrix and its transpose have the same Smith form. No scalar $t_i$ is inverted.
\end{proof}
\begin{lemma}\label{K08-pivot}
If $B$ is unimodular, then for arbitrary integer square $D,Q$ of the same size,
\[
 \begin{pmatrix}D&-B\\0&Q\end{pmatrix}
 \sim I\oplus QB^{-1}D.                                    \tag{K08.4}
\]
\end{lemma}
\begin{proof}
Permute the two block columns to put $-B$ in the upper-left corner. Use it as an integer unit block. Its Schur complement is $QB^{-1}D$, and $-B$ itself is equivalent to the identity. Every entry of $B^{-1}$ is integral.
\end{proof}

\subsection{Proof of the transfer theorem}\label{K08-sec-3}
Put $Q=\diag(q_i)$, $C=\diag(c_i)$ and $P=\diag(b_i)$. The block-constant part of the graph is encoded by the $2m\times2m$ matrix
\[
 B_0=\begin{pmatrix}A+C&I_m\\I_m&0\end{pmatrix},\qquad
 B_0^{-1}=\begin{pmatrix}0&I_m\\I_m&-A-C\end{pmatrix}.        \tag{K08.5}
\]
Thus $B_0$ is unimodular regardless of cycles or adjacency torsion in $H$. The Laplacian scalar part is $D_0=\diag(Q,N)$, and the cluster-size matrix is $T_0=\diag(N,P)$. Its compressed quotient block is
\[
 F=D_0-T_0B_0
 =\begin{pmatrix}Q-N(A+C)&-N\\-P&N\end{pmatrix}.            \tag{K08.6}
\]
The $c_i n_i$ contribution in $Q$ cancels in its top-left block, as it must: constant vectors on a clique have zero internal Laplacian contribution.

Lemma~\ref{K08-compress} splits off $(\Z/q_i)^{n_i-2}$ and $(\Z/n_i)^{b_i-2}$, leaving
$\left(\begin{smallmatrix}D_0&-B_0\\0&F\end{smallmatrix}\right)$.
By Lemma~\ref{K08-pivot}, its only nonunit block is $FB_0^{-1}D_0$. Direct multiplication gives
\[
 FB_0^{-1}D_0
 =\begin{pmatrix}
 -NQ&NQ\\
 NQ&-PN-N(A+C)N
 \end{pmatrix}.                                             \tag{K08.7}
\]
Add the first block row to the second. The lower-right block becomes
\[
 NQ-PN-N(A+C)N=L_w,                                         \tag{K08.8}
\]
using the definition of $q_i$. Next add the first block column to the second. This leaves
\[
 FB_0^{-1}D_0\sim (-NQ)\oplus L_w.                          \tag{K08.9}
\]
The diagonal block $NQ$ gives the $n_iq_i$ cyclic factors, and $L_w$ has one free factor because the weighted base is connected. Taking torsion proves (K08.2). This is an integral block diagonalization, not merely equality of determinants or an exact sequence whose splitting is left open.

\subsection{Closed consequences}\label{K08-sec-4}
\begin{corollary}[Arbitrary nonuniform decorated trees]\label{K08-tree}
If $H$ is a tree, then
\[
 K(G)\cong
 \bigoplus_i\left((\Z/q_i)^{n_i-2}
 \oplus(\Z/n_i)^{b_i-2}\oplus\Z/(n_iq_i)\right)
 \oplus\bigoplus_{ij\in E(H)}\Z/(n_in_j).                    \tag{K08.10}
\]
In particular
\[
 \exp K(G)=\lcm\bigl(\{n_iq_i\}_i\cup\{n_in_j:ij\in E(H)\}\bigr).
                                                                    \tag{K08.11}
\]
\end{corollary}
\begin{proof}
Choose an orientation of $H$ and delete one row from its incidence matrix to obtain a unimodular square matrix $F$. The reduced weighted Laplacian is $F\diag(n_in_j)F^{\mathsf T}$, so its cokernel is the direct sum of the edge-weight cyclic groups. Formula (K08.11) follows since the smaller local factors divide $n_iq_i$.
\end{proof}

\begin{corollary}[Constant core size, arbitrary base]\label{K08-constant}
Suppose $n_i=n\geq2$ for every $i$, and let $s_1\mid\cdots\mid s_{m-1}$ be the nonzero Smith factors of $L_H$, including its unit factors. Then the last term of (K08.2) contributed by the weighted base is
\[
 K(H_{\mathbf n})\cong\bigoplus_{j=1}^{m-1}\Z/(n^2s_j)\Z.   \tag{K08.12}
\]
The local factors use $q_i=b_i+n(d_i+c_i)$.
\end{corollary}
\begin{proof}
The weighted Laplacian is exactly $n^2L_H$, so the same unimodular Smith transformations multiply each nonzero invariant by $n^2$.
\end{proof}
Thus the ordinary critical group of $H$, together with its number of vertices and degrees, suffices to compute the entire new group in this regime. The count of unit Smith factors matters and is fixed by $m$.

For an arbitrary base, the order in (K08.2) is
\[
 |K(G)|=|K(H_{\mathbf n})|\prod_iq_i^{n_i-1}n_i^{b_i-1}.     \tag{K08.13}
\]
For a cycle base of length $m\geq3$ and constant core size $n$, its ordinary critical group is $\Z/m$, with full nonzero Smith list $1^{m-2},m$. Consequently (K08.12) becomes $(\Z/n^2)^{m-2}\oplus\Z/(mn^2)$. The local $q_i$ may still vary through $b_i,c_i$.

For a concrete nonuniform tree example, take $H$ to be one edge, $(n_1,n_2)=(2,3)$, $(b_1,b_2)=(2,4)$ and $(c_1,c_2)=(0,1)$. Then $(q_1,q_2)=(5,9)$ and (K08.10) reads
\[
 K(G)\cong \Z/9\oplus(\Z/3)^2\oplus\Z/10\oplus\Z/27
               \oplus\Z/6.
\]
Its order is $131220$ and its exponent is $270$. Every factor is covered without selecting favorable primes.

\subsection{Verification and literature boundary}\label{K08-sec-5}
One hundred twenty completed random connected base instances with $1\leq m\leq6$, independent sizes between two and six, and independent clique indicators were checked by exact Smith arithmetic on the compressed core. They agreed with (K08.9). The block proof, not these finite checks, establishes the theorem.

Twin compression has antecedents in \cite{ref028,ref029}; general Smith background is \cite{ref030}. A directly relevant 2018 paper by Tan, Wang and Qian is titled \emph{The Critical Group of the Corona of Trees and Cycles} \cite{ref034}. Its accessible abstract reports corona results, but its complete theorem statements were not available in the consulted landing page. Therefore we do not claim that all special cases above are new or that our formula has been exhaustively separated from that literature. The precise contribution submitted for checking is the fully written integral identity (K08.9) and its stated generality: arbitrary connected bases, varying core and pendant sizes, and mixed clique/independent core clusters. No named open problem or benchmark superiority is asserted.




\clearpage
\section[K09 --- The Exceptional Prime in Mixed Clique--Independent Tree Blowups]{The Exceptional Prime in Mixed Clique--Independent Tree Blowups}\label{contrib-K09}
\markboth{Integral Presentations and Critical Groups of Graph Blowups}{K09 --- Keio}
\noindent\textsc{K09 \quad Keio programme, source round 9}\par\medskip


\begin{quote}\small\textbf{Overview.} We complete the primewise reduction of uniform tree blowups whose clusters may independently be cliques or independent sets. At every prime where some shifted degree is a unit, a surjective-coupling theorem determines the group from separate degree and looped-adjacency Smith data. The only remaining situation is the prime two with every shifted degree even. In that case we prove an explicit local equivalence to a symmetric arrowhead matrix and evaluate all its determinantal divisors by subset products and sums. This removes the independent-leaf hypothesis and resolves the previously untreated arithmetic case.\end{quote}

\subsection{Setup and exact output}\label{K09-sec-1}
Let $T$ be a tree on $m\geq2$ vertices, with adjacency matrix $A$, degrees $d_i$, and Laplacian $L$. Choose $c_i\in\{0,1\}$, and replace vertex $i$ by $n\geq2$ vertices forming a clique if $c_i=1$, or an independent set if $c_i=0$. Adjacent clusters are completely joined. Put
\[
 e_i=d_i+c_i,\quad E=\diag(e_i),\quad B=A+\diag(c_i)=E-L.
\]
Integer twin compression gives
\[
 K(G)\cong\bigoplus_i(\Z/(ne_i)\Z)^{n-2}
           \oplus\operatorname{tors}\coker C,
 \qquad C=\begin{pmatrix}nE&-B\\0&nL\end{pmatrix}.          \tag{K09.1}
\]
To verify the compression, choose one generator in each cluster and subtract it from the others. The corresponding relation differences are scalar $ne_i$ relations. Replace one difference by their sum; the other $n-2$ split off. The remaining core is (K09.1), since $E-B=L$. These changes are unimodular.

Every leaf has shifted degree one or two. Consequently
\[
 \gcd_i e_i\in\{1,2\}.                                    \tag{K09.2}
\]
We give the complete core group at every prime. The unit case uses ordinary Smith data of two $m$-order base matrices. The exceptional case uses only a rational inverse with odd determinant and explicit subset arithmetic. This is a constructive classification; it is not asserted that looped-adjacency Smith data themselves depend only on the degree multiset.

\subsection{All primes with a unit shifted degree}\label{K09-sec-2}
Fix a prime $p$, put $\alpha=v_p(n)$, and work over $R=\Z_{(p)}$. Assume $p\nmid e_a$ for some vertex $a$. From the Smith valuation list of $E$, remove one zero to obtain
$\eta_1\leq\cdots\leq\eta_q$, where $q=m-1$. From the Smith valuation list of $B$, remove one zero to obtain $\beta_1\leq\cdots\leq\beta_q$, allowing $\beta_i=\infty$ for a zero Smith factor. Define $t_i=\min(\alpha,\beta_i)$.
\begin{theorem}\label{K09-unit}
In this case the $p$-primary torsion of the core is
\[
 \bigoplus_{i=1}^q\left(
 \Z/p^{t_i}\Z\ \oplus\ \Z/p^{2\alpha+\eta_{q+1-i}-t_i}\Z
                       \right).                            \tag{K09.3}
\]
An exponent zero denotes a trivial factor. This applies to every odd prime, and also to two unless all $e_i$ are even.
\end{theorem}

We give the needed algebraic fact rather than referring to an earlier round.
\begin{lemma}\label{K09-couple}
Suppose $D,K$ are square over $R$, $D$ is nonsingular and $[D\ K]$ is surjective. Let the Smith valuations of $D,K$ be $\eta_i,\beta_i$ in increasing order. Then the cokernel of $\left(\begin{smallmatrix}nD&K\\0&nI\end{smallmatrix}\right)$ has the factors in (K09.3).
\end{lemma}
\begin{proof}
Let $s$ be the number of positive or infinite $\beta_i$, and $r=q-s$. Put $K$ into the form $\diag(I_r,H)$, with diagonal entries of $H$ equal to powers $p^{\beta_i}$, permitting zero. Surjectivity forces the bottom $s$ rows of the transformed $D$ to be surjective. A column basis change makes that matrix
$\left(\begin{smallmatrix}D_1&F\\0&I_s\end{smallmatrix}\right)$.
The full presentation becomes, in block sizes $(r,s,r,s)$,
\[
 \begin{pmatrix}
 nD_1&nF&I&0\\0&nI&0&H\\0&0&nI&0\\0&0&0&nI
 \end{pmatrix}.
\]
Eliminate the unit block, then remove the residual $-n^2F$ using $nF$ times the row containing $nI$; the introduced $nFH$ is removed with $FH$ times the last row. The result is $I_r\oplus n^2D_1$ and $s$ blocks $\left(\begin{smallmatrix}n&p^{\beta_i}\\0&n\end{smallmatrix}\right)$. Each small block has valuations $\min(\alpha,\beta_i)$ and $2\alpha-\min(\alpha,\beta_i)$. The first $s$ valuations of $D$ are zero, and the others are the valuations of $D_1$. Reversing their indices pairs the positive $\beta_i$ with these zero degree valuations and proves (K09.3).
\end{proof}

\begin{proof}[Proof of Theorem~\ref{K09-unit}]
The sum of the second-family relations in (K09.1) is $-\sum_i e_i u_i=0$, where $u_i$ are the top generators. Since $e_a$ is a unit, eliminate $u_a=-\sum_{i\ne a}e_i u_i/e_a$. Its scalar relation becomes redundant. Replace bottom generators by their differences from the one at $a$, which is free. Let $E'$ be $E$ with row and column $a$ deleted, and let $\bar L$ be the reduced tree Laplacian. This matrix is unimodular: it is the product of a reduced tree incidence matrix and its transpose. The remaining finite presentation is
\[
 S=\begin{pmatrix}nE'&K\\0&nI_q\end{pmatrix},\qquad
 K=I-E'\left(\bar L^{-1}-J/e_a\right).                      \tag{K09.4}
\]
Indeed if $b$ is column $a$ of $B$ with its $a$ entry deleted, then $b=\bar L\boldsymbol1$. The upper-right block before multiplying by $\bar L^{-1}$ is $-(\bar B-e'b^{\mathsf T}/e_a)$, where $e'=E'\boldsymbol1$, giving (K09.4). Thus $[E'\ K]$ is surjective.

Replace column $a$ of $B$ by the sum of all its columns. That column is $e$, with unit $e_a$ at the pivot. Its Schur complement is $\bar B-e'b^{\mathsf T}/e_a$, which is equivalent to $-K$. Hence $B\sim1\oplus K$ over $R$, and $E'$ has the degree valuation list with one zero removed. Lemma~\ref{K09-couple} proves the theorem.
\end{proof}

\subsection{The sole exceptional situation}\label{K09-sec-3}
Assume now that every $e_i$ is even and $p=2$. Put $R=\Z_{(2)}$ and $\alpha=v_2(n)$. Choose any root indexed by $0$, and write $q=m-1$ and
\[
 B=\begin{pmatrix}\beta&b^{\mathsf T}\\b&\bar B\end{pmatrix},
 \quad E'=\diag(e_1,\ldots,e_q),\quad
 v=\bar B^{-1}b,\quad f_i=e_iv_i,
 \quad s=\beta-b^{\mathsf T}\bar B^{-1}b.                    \tag{K09.5}
\]
These expressions lie in $R$. Indeed $B\equiv-L\pmod2$, so the reduced matrix $\bar B$ has odd determinant. Moreover
\[
 \bar B\boldsymbol1+b=e',\qquad v\equiv-\boldsymbol1\pmod2,
 \qquad s+\sum_{i=1}^qf_i=e_0.                              \tag{K09.6}
\]
Thus every $v_i$ is a unit and $v_2(f_i)=v_2(e_i)$. The scalar $s$ may vanish.
\begin{theorem}\label{K09-arrow}
Over $R$, the core is equivalent to $I_q\oplus S\oplus0_1$, where
\[
 S=\begin{pmatrix}
 n^2\diag(f_1,\ldots,f_q)&-n f\\
 -n f^{\mathsf T}&-s
 \end{pmatrix},\qquad f=(f_1,\ldots,f_q)^{\mathsf T}.         \tag{K09.7}
\]
In particular, $S$ is nonsingular and has order $m$, with
\[
 v_2(\det S)=(2m-2)\alpha+\sum_{i=0}^qv_2(e_i).             \tag{K09.8}
\]
\end{theorem}
\begin{proof}
Pivot on the unit block $-\bar B$ in the upper-right portion of $C$. The residual first row, in columns consisting of the original top generator at $0$, the other top generators, and the second-family column at $0$, is
\[
 (ne_0,\ -n b^{\mathsf T}\bar B^{-1}E',\ -s).
\]
The lower rows are
\[
 (0,\ n^2 L_{*,\{1,\ldots,q\}}\bar B^{-1}E',\
 n(L_{*,0}-L_{*,\{1,\ldots,q\}}\bar B^{-1}b)).
\]
The $q$ indicated columns of $L$ form an integral basis for the zero-sum lattice, because a reduced tree Laplacian is unimodular. A bottom row basis change turns them into $\left(\begin{smallmatrix}I_q\\0\end{smallmatrix}\right)$. Since $L_{*,0}=-L_{*,\{1,\ldots,q\}}\boldsymbol1$ and $\boldsymbol1+\bar B^{-1}b=\bar B^{-1}e'$, multiplying the nonzero bottom rows by $\bar B$ leaves
\[
 \begin{pmatrix}
 ne_0&-n v^{\mathsf T}E'&-s\\
 0&n^2E'&-n e'
 \end{pmatrix}                                               \tag{K09.9}
\]
and one zero row. The sum of the first column, all middle columns, and $n$ times the last column is zero, by (K09.6). Replacing the first column by that sum is unimodular and produces a zero column. Delete this zero row and column, multiply bottom row $i$ by the unit $v_i$, and put the first row last. The result is (K09.7). Its determinant is
$-n^{2q}(\prod_i f_i)(s+\sum_i f_i)=-n^{2q}e_0\prod_i f_i$, giving (K09.8).
\end{proof}

\subsection{Explicit invariant factors at two}\label{K09-sec-4}
Set $\eta_i=v_2(e_i)$ and use $v_2(0)=\infty$. Let $\Delta_k$ be the valuation of the gcd of the $k$-minors of $S$, with $\Delta_0=0$. The following formulas require no Smith calculation:
\begin{align*}
 \Delta_1&=\min\left(v_2(s),\ \alpha+\min_{1\leq i\leq q}\eta_i\right),               \tag{K09.10}\\
 \Delta_k&=(2k-2)\alpha+\min\Biggl\{
 \min_{\substack{J\subseteq[q]\\|J|=k}}\sum_{j\in J}\eta_j,\\[-2mm]
 &\hspace{35mm}\min_{\substack{I\subseteq[q]\\|I|=k-1}}
 \left[\sum_{i\in I}\eta_i+v_2\!\left(s+\sum_{i\in I}f_i\right)\right]
 \Biggr\}\quad(2\leq k\leq q),                              \tag{K09.11}\\
 \Delta_m&=(2m-2)\alpha+\sum_{i=0}^q\eta_i.                 \tag{K09.12}
\end{align*}
For $m=2$ the middle range is empty. The core's two-primary group is
\[
 \bigoplus_{k=1}^m\Z/2^{\Delta_k-\Delta_{k-1}}\Z.            \tag{K09.13}
\]
All sums in (K09.10)--(K09.12) are rational numbers with odd denominators, so their two-adic valuations are obtained by exact integer arithmetic.

\begin{proof}
For $k=1$, the entries of (K09.7) generate the ideal $(s,nf_1,\ldots,nf_q)$, proving (K09.10). For larger minors classify their row and column sets by whether they contain the border index. Without either border, nonzero minors are $n^{2k}\prod_{j\in J}f_j$. With exactly one border they are, up to sign, $n^{2k-1}\prod_{j\in J}f_j$. With both borders and identical nonborder sets $I$, they are
\[
 -n^{2k-2}\left(\prod_{i\in I}f_i\right)
               \left(s+\sum_{i\in I}f_i\right).
\]
With both borders but different nonborder sets, a nonzero minor requires the sets to differ by exactly one index. It is then, up to sign, $n^{2k-2}\prod_{j\in J}f_j$, with $|J|=k$. Every such product occurs for $2\leq k\leq q$. If the sets differ by more than one index the minor vanishes. Thus the last two types generate the ideal of all $k$-minors; taking valuations gives (K09.11). The full determinant gives (K09.12). Successive determinantal-divisor ratios are the Smith factors, proving (K09.13).
\end{proof}

\subsection{Completeness, checks, and limitations}\label{K09-sec-5}
Equations (K09.1), (K09.3), and (K09.10)--(K09.13) determine the primary decomposition for every choice of clique indicators and every $n\geq2$. The exceptional case occurs exactly when $c_i\equiv d_i\pmod2$ at all vertices. An independent leaf prevents it, but the converse restriction is unnecessary in the unit case. Thus this paper removes the independent-leaf condition of the earlier coupling application rather than merely repeating it.

For the two-vertex tree with both clusters cliques, $e_0=e_1=2$, $v_1=1$, $s=0$. The arrowhead matrix is $\left(\begin{smallmatrix}2n^2&-2n\\-2n&0\end{smallmatrix}\right)$; its two nonzero Smith factors are both $2n$. With the scalar factors in (K09.1), this gives $(\Z/(2n))^{2n-2}$, as required for the complete graph on $2n$ vertices.

The determinant identity in the exceptional proof yields the core order $n^{2m-2}\prod_i e_i$. Combining with (K09.1) gives the usual spanning-tree count $n^{mn-2}\prod_i e_i^{n-1}$. More importantly, the subset sums in (K09.11) retain cancellations that an order calculation alone cannot see. The formula is exponential in the number of subsets if evaluated naively; no optimal complexity claim is made.

Exact checks on 135 random trees with $2\leq m\leq10$, the parity-forced exceptional indicators, and $n\in\{2,3,4,6,8\}$ confirmed the local arrowhead equivalence at the level of Smith valuations. They are supplementary to the written proof. The looped-adjacency input is not a new invariant theory: Corrales and Valencia \cite{ref035} describe generalized tree-Laplacian determinantal ideals through minimal $2$-matchings. Specializing their diagonal variables to $-c_i$ gives the determinantal ideals of $-B$, and hence can supply the Smith data used in (K09.3). Our exceptional arrowhead calculation is a specialized arithmetic reduction, not a replacement claim for that general theory. The local coupling and twin-ideal viewpoint have precedents \cite{ref029,ref028}; Smith-form background is \cite{ref030}. No named open problem, exhaustive novelty certification, or benchmark superiority is claimed.




\clearpage
\section[K10 --- What All Uniform Blowup Critical Groups Can and Cannot Recover]{What All Uniform Blowup Critical Groups Can and Cannot Recover}\label{contrib-K10}
\markboth{Integral Presentations and Critical Groups of Graph Blowups}{K10 --- Keio}
\noindent\textsc{K10 \quad Keio programme, source round 10}\par\medskip


\begin{quote}\small\textbf{Overview.} We characterize exactly when two trees have isomorphic critical groups after every uniform independent blowup. The complete invariant is the number of vertices, the maximum matching number, and the finite abelian group formed by the vertex degrees, not the degree multiset itself. A single blowup at any prime larger than both maximum degrees already detects this invariant. We then construct, for every $r\geq1$, $2^r$ trees on $15r^2+19r$ vertices with pairwise different degree multisets but identical critical groups for every blowup size. An explicit subfamily also has different Laplacian spectra at every blowup size. All required local-counting arguments are included.\end{quote}

\subsection{An inverse classification}\label{K10-sec-1}
For a tree $T$ on $m\geq2$ vertices write $\nu(T)$ for its maximum matching number, and define its degree group
\[
 \mathcal D(T)=\bigoplus_{v\in V(T)}\Z/\deg_T(v)\Z.
\]
Unit cyclic factors are trivial. Let $T[n]$ replace every vertex by $n\geq2$ independent twins and every base edge by a complete bipartite graph.
\begin{theorem}\label{K10-inverse}
For two finite trees $T,S$, the following conditions are equivalent:
\begin{enumerate}
\item $K(T[n])\cong K(S[n])$ for every integer $n\geq2$.
\item $|V(T)|=|V(S)|$, $\nu(T)=\nu(S)$, and $\mathcal D(T)\cong\mathcal D(S)$.
\item $K(T[p])\cong K(S[p])$ for some prime $p$ larger than every vertex degree of both trees.
\end{enumerate}
In condition 3, ``some'' can be replaced by ``every'' without changing the equivalence.
\end{theorem}
Thus the family of all abstract critical groups does not generally recover the degree multiset. The result concerns abstract groups, not their canonical monodromy pairings, marked divisor classes, or automorphism actions.

\subsection{A self-contained local-counting lemma}\label{K10-sec-2}
Let $A,D,L=D-A$ denote the adjacency, degree and Laplacian matrices of a tree, and put $\delta=m-2\nu$. Twin generator differences give the integral compression
\[
 L_{T[n]}\sim\bigoplus_i nd_iI_{n-2}\ \oplus C_n,
 \qquad C_n=\begin{pmatrix}nD&-A\\0&nL\end{pmatrix}.         \tag{K10.1}
\]
Indeed, subtract one distinguished generator and its relation in each cluster from the others. The difference relations are scalar $nd_i$ relations. Replace one difference by the sum of the $n-1$ differences. The other $n-2$ split off, leaving (K10.1); all changes are integral basis operations.

We need three elementary tree facts. First, $A$ has $2\nu$ unit Smith factors and $\delta$ zeros. To prove this, pivot on the adjacency block of a leaf and its neighbor, $\left(\begin{smallmatrix}0&1\\1&0\end{smallmatrix}\right)$. Its Schur correction is zero because the leaf has no other neighbor. Delete the pair and induct; a maximum matching can always contain that leaf edge. Second, the image of $L$ is the zero-sum lattice, and its kernel over every $\Z/p^k\Z$ consists of constants. A reduced incidence matrix of a tree is unimodular, so every reduced Laplacian has determinant one, proving both assertions. Third, $[D\ A]$ is surjective: its image contains that of $L$, and it contains the coordinate vector at any leaf since its degree is one.

Fix a prime $p$, put $\alpha=v_p(n)$, and define
\[
 Q_p(c)=\sum_{i=1}^m\min(v_p(d_i),c),\qquad
 \kappa_p(k)=\log_p|\ker(C_n\bmod p^k)|.
\]
\begin{lemma}\label{K10-counts}
For every $k\geq1$,
\[
 \kappa_p(k)=
 \begin{cases}
 (m+\delta)k,&k\leq\alpha,\\
 2\nu k+2\alpha\delta,&\alpha<k\leq2\alpha,\\
 (k-2\alpha)+2m\alpha+Q_p(k-2\alpha),&k>2\alpha.
 \end{cases}                                                 \tag{K10.2}
\]
\end{lemma}
\begin{proof}
Unit rescalings allow us to replace $n$ locally by $p^\alpha$. A kernel vector $(x,y)$ satisfies
\[
 nDx=Ay,\qquad L(ny)=0.
\]
Thus $ny=t\boldsymbol1$. Multiplying the first equation by $n$ gives $D(n^2x-t\boldsymbol1)=0$. A leaf degree is one, so $t$ is a multiple of $n^2$ modulo $p^k$.

For $k\leq\alpha$, $x$ is arbitrary and $Ay=0$; the primitive adjacency Smith form gives the first line. For $\alpha<k\leq2\alpha$, one has $t=0$ and $y=p^{k-\alpha}z$, with $z$ modulo $p^\alpha$. Put $c=2\alpha-k$. The first equation becomes
\[
 p^cDx-Az=0\pmod{p^\alpha}.
\]
Since $A$ has primitive image and $[D\ A]$ is surjective, the matrix $[p^cD\ -A]$ has $2\nu$ unit factors and $\delta$ factors $p^c$: quotient by $\im A$ and use surjectivity of the induced $D$. Its kernel has $p^{m\alpha+c\delta}$ elements, and $x$ has $p^{m(k-\alpha)}$ lifts. The exponent is $mk+c\delta=2\nu k+2\alpha\delta$.

Finally let $k>2\alpha$ and $c=k-2\alpha$. There are $p^c$ possible $t$. For $t=n^2u$, subtract the kernel vector $(u\boldsymbol1,nu\boldsymbol1)$, reducing its fiber to $t=0$. Then $y=p^{k-\alpha}z$ with $z$ modulo $p^\alpha$, and the top equation becomes
\[
 Dx-p^cAz=0\pmod{p^{\alpha+c}}.
\]
The image of $[D\ p^cA]$ is $D\Z_{(p)}^m+p^c\Z_{(p)}^m$, by surjectivity of $[D\ A]$, and has index $p^{Q_p(c)}$. If $z$ is temporarily allowed modulo $p^{\alpha+c}$, the kernel has size $p^{m(\alpha+c)+Q_p(c)}$. Divide by $p^{mc}$ for the duplicated lifts of $z$, multiply by $p^{m\alpha}$ for lifts of $x$ to modulus $p^k$, and finally by $p^c$ for the choices of $t$. This gives the final line of (K10.2).
\end{proof}
A matrix with one zero Smith factor and finite valuations $e_j$ has kernel exponent $k+\sum_j\min(k,e_j)$. Consecutive differences determine the count of valuations at least $k$, so (K10.2) determines the full primary core group. In particular, the right side depends only on $m,\nu$ and the elementary-divisor multiplicities of $\mathcal D(T)$. The scalar factors in (K10.1) have the same dependence. This proves that condition 2 implies condition 1 in Theorem~\ref{K10-inverse}.

\subsection{One prime is enough}\label{K10-sec-3}
Let $p$ exceed every degree of $T$, and take blowup size $n=p$. Since all degree valuations at $p$ vanish, (K10.2) yields
\[
 K(T[p])_{(p)}\cong
 (\Z/p\Z)^{mp-4\nu}\oplus(\Z/p^2\Z)^{2\nu-1}.             \tag{K10.3}
\]
If $a$ is the number of $p$ factors and $b$ the number of $p^2$ factors, then
\[
 \nu=(b+1)/2,\qquad m=(a+4\nu)/p.                           \tag{K10.4}
\]
Thus this primary group recovers both integers exactly.

For a different prime $q$, the valuation of $n=p$ is zero. Lemma~\ref{K10-counts} gives one copy of the $q$-primary degree group in the core. The scalar factors contribute another $p-2$ copies. Hence the prime-to-$p$ subgroup is
\[
 K(T[p])_{(p')}\cong\mathcal D(T)^{\oplus(p-1)}.             \tag{K10.5}
\]
A finite abelian group has at most one $(p-1)$-fold direct-sum root: in each primary elementary-divisor list, divide every multiplicity by $p-1$. Therefore (K10.3)--(K10.5) recover $\mathcal D(T)$. Condition 3 implies condition 2. Condition 1 plainly implies condition 3, since primes larger than any fixed degree bound exist. This completes the theorem.

The degree bound matters to the simple extraction procedure: at a prime dividing a vertex degree, the large factors in (K10.3) acquire additional valuations. The theorem does not claim that an arbitrary single blowup size always determines the complete invariant.

\subsection{An arithmetic trade realized by trees}\label{K10-sec-4}
Let $t\geq3$ be divisible by three. The triples
\[
 (2,2t+1,3t+3),\qquad (3,t+1,4t+2)                          \tag{K10.6}
\]
have the same sum $5t+6$ and the same degree group. Indeed
\[
 \Z/2\oplus\Z/(2t+1)\cong\Z/(4t+2),\qquad
 \Z/3\oplus\Z/(t+1)\cong\Z/(3t+3),                        \tag{K10.7}
\]
because the respective pairs are coprime.

Construct a tree $P_t$ with two adjacent nonleaf vertices of degrees $2t$ and $3t+2$, by attaching $2t-1$ leaves to the first and $3t+1$ leaves to the second. Construct $Q_t$ with three nonleaf vertices in a path, of degrees $2,t,4t+1$ in that order, by attaching $1,t-2,4t$ leaves, respectively. Both trees have $5t+2$ vertices.

Let $\widehat P_t$ and $\widehat Q_t$ be obtained by attaching one new leaf to every vertex of $P_t,Q_t$. These are trees on $10t+4$ vertices with perfect matchings consisting of all newly attached leaf edges. Their degree multisets are
\begin{align*}
 \widehat P_t:&\quad 1^{5t+2},\ 2^{5t},\ 2t+1,\ 3t+3,\\
 \widehat Q_t:&\quad 1^{5t+2},\ 2^{5t-1},\ 3,\ t+1,\ 4t+2.
                                                                    \tag{K10.8}
\end{align*}
After cancelling the common entries, (K10.6)--(K10.7) show their degree groups agree. Theorem~\ref{K10-inverse} proves
\[
 K(\widehat P_t[n])\cong K(\widehat Q_t[n])\qquad(n\geq2).   \tag{K10.9}
\]
Their maximum degrees are $3t+3$ and $4t+2$, so even the maximum degree is not recovered.

These pairs are not Laplacian cospectral. Their degree-square sums differ by
\[
 3^2+(t+1)^2+(4t+2)^2-
       \bigl(2^2+(2t+1)^2+(3t+3)^2\bigr)=4t(t-1)>0.          \tag{K10.10}
\]
For any simple graph, $\operatorname{tr}(L^2)=\sum_vd_v^2+\sum_vd_v$. The two trees have equal order and edge count, so (K10.10) separates their spectra. After uniform blowup, the degree-square sum is multiplied by $n^3$ and the degree sum by $n^2$. Thus the Laplacian spectra remain different for every $n\geq2$, despite (K10.9).

\subsection{A superpolynomial family with different degree multisets}\label{K10-sec-5}
\begin{theorem}\label{K10-family}
For every integer $r\geq1$, there are $2^r$ trees on
\[
 M_r=15r^2+19r
\]
vertices whose degree multisets are pairwise different, but whose critical groups after uniform independent blowup are mutually isomorphic for every size $n\geq2$.
\end{theorem}
\begin{proof}
For each $j=1,\ldots,r$, put $t_j=3j$ and independently choose a block $P_{t_j}$ or $Q_{t_j}$. Join the blocks in a path by adding one edge between a leaf of block $j$ and a leaf of block $j+1$. Use distinct leaves for the two connections of an internal block. Both block types have at least $5t_j-1$ leaves, so these ports can be chosen among a common pool of leaves in either type.

The connected result $H$ is a tree. The added edges change precisely $2(r-1)$ port degrees from one to two, identically for every choice of blocks. Now attach one new leaf to every vertex of $H$, obtaining a tree $T$. It has a perfect matching and
\[
 |V(T)|=2\sum_{j=1}^r(5t_j+2)=15r^2+19r.
\]
In the degree multiset of $T$, each port change replaces a common entry $2$ by $3$. These identical replacements preserve the degree-group equality proved in (K10.7) for each block. More explicitly, all choices have degree group
\[
 (\Z/2)^{\sum_j(5t_j-1)-2(r-1)}\oplus(\Z/3)^{2(r-1)}
 \oplus\bigoplus_{j=1}^r
       \left(\Z/(4t_j+2)\oplus\Z/(3t_j+3)\right).           \tag{K10.11}
\]
Theorem~\ref{K10-inverse} now gives equality of every blowup critical group.

To distinguish the degree multisets, observe that choosing $Q_{3j}$ creates the degree $12j+2$, while choosing $P_{3j}$ does not. These markers cannot arise from another kind of degree: the $P$ exceptional degrees are $6k+1$ (odd) and $9k+3$ (divisible by three), while the other $Q$ exceptional degree is $3k+1$ (one modulo three). The common degrees are one, two and three. Among the marker degrees themselves, $12j+2=12k+2$ only if $j=k$. Therefore the degree multiset determines every block choice, giving exactly $2^r$ distinct multisets.
\end{proof}
Since $M_r$ is quadratic in $r$, the family size is $\exp(\Omega(\sqrt{M_r}))$, in particular larger than every fixed power of $M_r$ along this sequence. This statement is about degree-distinct examples, not just multiple tree shapes with one degree sequence.

\subsection{Scope and provenance}\label{K10-sec-6}
The forward local-counting formula develops the uniform classification established earlier in this tournament; its complete proof is reproduced here so the inverse theorem is standalone. The inverse equivalence, one-prime reconstruction, arithmetic trade, and growing degree-distinct families are the new questions addressed in this round. The construction is explicit and uses no search certificate. Four parameter checks at $t=3,6,9,12$ independently verified the degree sums, elementary-divisor signatures, and second-moment identity; they are not needed for the proof.

Smith methods and twin critical ideals have established literature \cite{ref030,ref029,ref035}; multipartite path-of-cluster groups are studied in \cite{ref028}. A literature search did not establish an identical inverse classification, but was not exhaustive. No named pre-existing open question or benchmark superiority is claimed. In particular, the result makes no assertion that the canonical pairings are indistinguishable, or that arbitrary nonuniform blowup groups fail to recover more information.




\clearpage
\part{Electrical Networks, Spectral Design, and Greedy Augmentation}


\clearpage
\section[T01 --- Weighted Star Defects Extremize Resistance, Heat Traces, and Tree Sums]{Weighted Star Defects Extremize Resistance, Heat Traces, and Tree Sums}\label{contrib-T01}
\markboth{Electrical Networks, Spectral Design, and Greedy Augmentation}{T01 --- Taisho}
\noindent\textsc{T01 \quad Taisho programme, source round 1}\par\medskip



\begin{quote}\small\textbf{Overview.} Prescribe the magnitudes of at most $n-1$ conductance losses in an initially homogeneous complete network. We prove that arranging the losses on a star simultaneously maximizes every inverse spectral moment and every nonconstant heat trace, and minimizes the weighted spanning-tree sum, whenever the star remains connected. The result permits arbitrary unequal loss magnitudes. A single entrywise comparison of weighted incidence Gram matrices yields all these conclusions, their equality cases, and a quantitative deficit away from the star. We give closed formulas for the worst total effective resistance and tree sum. The framework of edge-defect matrices and several unweighted special cases predate this paper; the precise weighted simultaneous comparison is the result established here, without a certified priority claim.\end{quote}
\subsection{Problem and scope}\label{T01-sec-1}
Let $n\ge3$, $c>0$, and $1\le k\le n-1$. Fix $\alpha_1,\ldots,\alpha_k\in(0,c]$. Choose distinct edges $e_1,\ldots,e_k$ of $K_n$ and reduce the conductance of $e_i$ from $c$ to $c-\alpha_i$. Other conductances stay $c$. Call the resulting network $G$. The \emph{star arrangement} $G_\star$ puts all $k$ losses at edges sharing one endpoint. We assume
\begin{equation}\label{T01-eq:condition}
 k\le n-2,\qquad\text{or}\qquad k=n-1\text{ and }\min_i\alpha_i<c.
\end{equation}
Write $L_G$ for the weighted Laplacian and $\lambda_2(G),\ldots,\lambda_n(G)$ for its eigenvalues on $\cH=\one^\perp$. The quantities of interest are
\[
 Z_s(G)=\sum_{j=2}^n\lambda_j(G)^{-s}\quad(s>0),\qquad
 H_t(G)=\sum_{j=2}^n e^{-t\lambda_j(G)}\quad(t>0),
\]
and the weighted tree sum $\tau(G)=n^{-1}\prod_{j=2}^n\lambda_j(G)$. The Kirchhoff index is $nZ_1(G)$. We prove the following theorem without assuming that the loss graph is a forest.
\begin{theorem}[Simultaneous weighted star extremality]\label{T01-thm:main}
Under \eqref{T01-eq:condition}, every arrangement is connected and
\[
 Z_s(G)\le Z_s(G_\star),\qquad H_t(G)\le H_t(G_\star),\qquad
 \tau(G)\ge\tau(G_\star)
\]
for every $s,t>0$. If $k\ge2$, equality in any one of these inequalities holds exactly when the loss edges have a common endpoint. For $k=1$ every arrangement is isomorphic to the star.
\end{theorem}
The theorem is about prescribed \emph{loss magnitudes}, not a prescribed survivor spectrum or an optimization of their magnitudes. It does not compare each ordered eigenvalue individually.

\subsection{A Gram-matrix comparison}\label{T01-sec-2}
Orient the lost edges arbitrarily and let $B$ have columns $b_i=e_{u_i}-e_{v_i}$. Put
\[
 A=\diag(\alpha_1,\ldots,\alpha_k),\quad Q=A^{1/2}B^TBA^{1/2},\quad
 v=(\sqrt{\alpha_i})_{i=1}^k,\quad C=A+vv^T.
\]
The matrix $C$ is the Gram matrix of an outward-oriented weighted star.
\begin{lemma}\label{T01-lem:moments}
For every integer $r\ge1$,
\[
 \tr Q^r\le\tr C^r,\qquad
 \lambda_{\max}(Q)\le\lambda_{\max}(C)\le\max_i\alpha_i+\sum_i\alpha_i<nc.
\]
The last strict inequality uses \eqref{T01-eq:condition}. Moreover,
\begin{equation}\label{T01-eq:second}
 \tr C^2-\tr Q^2=2\sum_{i<j:\,e_i\cap e_j=\varnothing}\alpha_i\alpha_j.
\end{equation}
If the loss edges are pairwise intersecting but have no common endpoint, then $k=3$, the edges form a triangle, and
\begin{equation}\label{T01-eq:third}
 \tr C^3-\tr Q^3=12\alpha_1\alpha_2\alpha_3.
\end{equation}
\end{lemma}
\begin{proof}
The diagonal of $Q$ is $2\alpha_i$, equal to that of $C$. For $i\ne j$, the entry of $Q$ is $0$ for disjoint edges and is $\pm\sqrt{\alpha_i\alpha_j}$ otherwise. Hence $|Q_{ij}|\le C_{ij}$. Expand each trace as a sum over closed index walks. Every summand for $Q$ is at most its absolute value, which is at most the corresponding nonnegative summand for $C$. This proves the moment inequality. The Rayleigh quotient gives $x^TQx\le |x|^TC|x|$, proving the first spectral bound. The second follows from $A\preceq(\max_i\alpha_i)I$ and $\norm v^2=\sum_i\alpha_i$. Condition \eqref{T01-eq:condition} makes their sum strictly below $nc$.

Squaring entries proves \eqref{T01-eq:second}. To classify pairwise intersecting edges, start with $ab,ac$. Any edge not incident to $a$ must be $bc$. An edge intersecting all of $ab,ac,bc$ is one of those three, so a nonstar family consists exactly of a triangle. For a triangle, the product of the three distinct off-diagonal entries of $Q$ is $-\alpha_1\alpha_2\alpha_3$, irrespective of orientations; for $C$ it is positive. All other length-three closed walks have the same contribution. There are six orders of the three indices, giving \eqref{T01-eq:third}.
\end{proof}
\begin{proposition}[Positive-coefficient functional comparison]\label{T01-prop:function}
Let $f(x)=\sum_{r\ge0}a_rx^r$ converge at $\lambda_{\max}(C)$ with $a_r\ge0$ for $r\ge1$. Then $\tr f(Q)\le\tr f(C)$. If $a_2,a_3>0$, equality holds exactly for star loss support. More quantitatively,
\begin{equation}\label{T01-eq:quant}
 \tr f(C)-\tr f(Q)\ge
 2a_2\sum_{i<j:\,e_i\cap e_j=\varnothing}\alpha_i\alpha_j.
\end{equation}
For the exceptional pairwise-intersecting nonstar support, the right side can instead be taken as $12a_3\alpha_1\alpha_2\alpha_3$.
\end{proposition}
\begin{proof}
Apply Lemma~\ref{T01-lem:moments} to partial sums, then pass to the convergent series. Each moment gap is nonnegative. Retaining the second or third gap proves the bounds and strictness. Star support admits orientations for which $Q=C$, proving the converse.
\end{proof}

\subsection{Proof of the spectral conclusions}\label{T01-sec-3}
Put $a=nc$ and $X=BA^{1/2}$. On $\cH$,
\[
 L_G=aI-XX^T.
\]
The nonzero eigenvalues of $XX^T$ are those of $Q=X^TX$, with zeros added as necessary. Therefore, for every function $g$ defined on the relevant spectra,
\begin{equation}\label{T01-eq:transfer}
 \tr_{\cH}g(L_G)=(n-1-k)g(a)+\tr g(aI_k-Q).
\end{equation}
This remains valid when $Q$ is singular: the zero eigenvalues on both sides cancel with the displayed multiplicity. Lemma~\ref{T01-lem:moments} also proves positive definiteness on $\cH$, hence connectivity.

For $s>0$ the series
\[
 (a-x)^{-s}=a^{-s}\sum_{r\ge0}\frac{(s)_r}{r!}(x/a)^r
\]
has strictly positive coefficients; $(s)_r=s(s+1)\cdots(s+r-1)$. Similarly $e^{-t(a-x)}=e^{-ta}\sum_{r\ge0}t^rx^r/r!$ does. For the tree sum use
\[
 -\log(a-x)=-\log a+\sum_{r\ge1}\frac{x^r}{ra^r}.
\]
All three series converge at $\lambda_{\max}(C)<a$. Proposition~\ref{T01-prop:function} and \eqref{T01-eq:transfer} prove Theorem~\ref{T01-thm:main}, including all equality cases.

For example, if two lost edges are disjoint, the inverse-moment deficit is at least
\[
 Z_s(G_\star)-Z_s(G)\ge
 \frac{s(s+1)}{a^{s+2}}
 \sum_{i<j:\,e_i\cap e_j=\varnothing}\alpha_i\alpha_j.
\]
For a triangle of losses it is at least
$2s(s+1)(s+2)\alpha_1\alpha_2\alpha_3/a^{s+3}$.

\subsection{Closed worst-case formulas}\label{T01-sec-4}
Define $d_i=a-\alpha_i$ and
\[
 q=1-\sum_{i=1}^k\frac{\alpha_i}{d_i}>0.
\]
Positivity follows because $aI-C\succ0$. The diagonal-plus-rank-one determinant and inverse identities give
\begin{align}
 Z_1(G_\star)&=\frac{n-1-k}{a}+\sum_i\frac1{d_i}
 +\frac{\sum_i\alpha_i/d_i^2}{q},\label{T01-eq:closed1}\\
 \tau(G_\star)&=\frac{a^{n-1-k}}n\left(\prod_i d_i\right)q.\label{T01-eq:closed2}
\end{align}
For completeness, write $D=\diag(d_i)$. Then
$(D-vv^T)^{-1}=D^{-1}+D^{-1}vv^TD^{-1}/(1-v^TD^{-1}v)$, verified by multiplication. Factoring $D^{1/2}$ and diagonalizing $I-D^{-1/2}vv^TD^{-1/2}$ proves its determinant is $\det(D)(1-v^TD^{-1}v)$. These identities prove \eqref{T01-eq:closed1}--\eqref{T01-eq:closed2} directly.

If every loss is a complete deletion, $\alpha_i=c$, and $k\le n-2$, the surviving star-defect spectrum is
\[
 c(n-k-1)\ (1\text{ time}),\quad c(n-1)\ (k-1\text{ times}),\quad
 cn\ (n-k-1\text{ times}).
\]
Consequently the sharp upper bound becomes
\[
 Z_s(G)\le c^{-s}\left((n-k-1)^{-s}+(k-1)(n-1)^{-s}+(n-k-1)n^{-s}\right).
\]
This unweighted specialization is not asserted as an original discovery.

\subsection{Boundary and limitations}\label{T01-sec-5}
If $k=n-1$ and every loss is $c$, the star isolates its center. Inverse moments are then infinite; the finite comparison above deliberately excludes this case. For $k\ge n$, a simple star cannot realize the prescribed number of loss edges, and the theorem makes no assertion. Neither the entrywise proof nor the conclusion permits negative loss magnitudes. The result also does not assert full eigenvalue majorization: it uses only an explicitly identified cone of trace functions.

\subsection{Prior work and reproducibility}\label{T01-sec-6}
Ghosh, Boyd and Saberi established the convex resistance-design framework \cite{ref036}. Xu, Das and Zhang studied extremal Kirchhoff indices after unweighted complete-graph deletions \cite{ref037}. Tamura explicitly develops the edge-defect resolvent and its stability theory \cite{ref038}; we do not claim the dimension reduction as new. The present proof establishes the stated weighted comparison and functional/equality extension, but a short search cannot certify its priority. No existing named conjecture is claimed resolved.

The proof is finite-dimensional and self-contained. An exact verification consists of checking $Q_{ii}=2\alpha_i$, the three possibilities for $Q_{ij}$, and the closed-walk expansion. No computational experiment is required. The supplied benchmark manuscripts on weighted-Laplacian cones and one missing edge concern inverse spectral feasibility; none is used as a lemma here.




\clearpage
\section[T02 --- Exact Minimax Conductance Design under Fractional Edge Failures]{Exact Minimax Conductance Design under Fractional Edge Failures}\label{contrib-T02}
\markboth{Electrical Networks, Spectral Design, and Greedy Augmentation}{T02 --- Taisho}
\noindent\textsc{T02 \quad Taisho programme, source round 2}\par\medskip



\begin{quote}\small\textbf{Overview.} A designer distributes a fixed conductance budget over the complete graph; an adversary applies a prescribed list of fractional edge losses at distinct edges. For every inverse-power spectral loss we determine the unique minimax design, its exact value, and its worst attack. The design is uniform, the attack has star support, and a global quadratic inequality controls the distance of any near-optimal design from uniformity. The same unique design optimizes the worst heat trace and worst weighted spanning-tree sum. Symmetrization is classical; the main content is the exact adversarial value, simultaneous criteria, and an explicit stability certificate that survives complete edge deletion.\end{quote}
\subsection{The design game}\label{T02-sec-1}
Let $n\ge3$, $m=\binom n2$, $B>0$, and $\mathcal W=\{w\in\R_+^m:\sum_e w_e=B\}$. Fix $k$ loss fractions $\theta_1,\ldots,\theta_k\in(0,1]$, where
\[
 k\le n-2,\quad\text{or}\quad k=n-1\text{ and }\min_i\theta_i<1.
\]
An attack $\sigma$ assigns the $k$ labeled fractions to distinct edges. Put $q_e^\sigma=1-\theta_i$ if $i$ is assigned to $e$, and $q_e^\sigma=1$ otherwise. The surviving Laplacian is
\[
 L_\sigma(w)=\sum_e q_e^\sigma w_e b_eb_e^T,\qquad b_{uv}=e_u-e_v.
\]
All spectral matrices below are restricted to $\cH=\one^\perp$. For $s>0$, define
\[
 J_s(w)=\max_\sigma\tr_{\cH}L_\sigma(w)^{-s},
\]
with value $+\infty$ if some attack disconnects the support. The fractions multiply the designed weights. This is not the different problem of subtracting fixed absolute conductances from a nonuniform design.

\begin{theorem}[Exact minimax solution and quantitative uniqueness]\label{T02-thm:main}
For every $s>0$ the unique minimizer is $u_e=B/m$. Set $v_i=\sqrt{\theta_i}$, $C=\diag(\theta_i)+vv^T$, and $c=B/m$. Then
\begin{equation}\label{T02-eq:value}
 \min_{w\in\mathcal W}J_s(w)=c^{-s}\left((n-1-k)n^{-s}+\tr(nI_k-C)^{-s}\right).
\end{equation}
An attack is worst at $u$ exactly when its support is a star (with the trivial interpretation for $k=1$). Every feasible $w$ satisfies
\begin{equation}\label{T02-eq:stability}
 J_s(w)-J_s(u)\ge
 \frac{s(s+1)}{(2B)^{s+2}}
 \frac{m-\sum_i(2\theta_i-\theta_i^2)}{m}
 \sum_e(w_e-B/m)^2.
\end{equation}
\end{theorem}
The coefficient is positive because an attack leaves at least one edge unaffected. The inequality is global, not a local Hessian assertion.

\subsection{A trace convexity lemma with an explicit modulus}\label{T02-sec-2}
\begin{lemma}\label{T02-lem:convex}
If a scalar function $f$ is convex on an interval $I$, then $M\mapsto\tr f(M)$ is convex on symmetric matrices with spectrum in $I$. If $f''\ge\mu>0$, then, for matrices $M_j$ in that set and a probability vector $(p_j)$,
\[
 \sum_jp_j\tr f(M_j)-\tr f(\overline M)
 \ge\frac\mu2\sum_jp_j\norm{M_j-\overline M}_F^2,
 \qquad \overline M=\sum_jp_jM_j.
\]
\end{lemma}
\begin{proof}
Diagonalize $\overline M$ in an orthonormal basis $z_i$. Scalar Jensen first gives
$f(z_i^T\overline Mz_i)\le\sum_jp_j f(z_i^TM_jz_i)$. Diagonalizing each $M_j$ and applying scalar Jensen to the squared overlaps gives $f(z_i^TM_jz_i)\le z_i^Tf(M_j)z_i$. Summing proves trace convexity. Apply this to the convex scalar function $f(x)-\mu x^2/2$ and use
$\sum_jp_j\tr M_j^2-\tr\overline M^2=\sum_jp_j\norm{M_j-\overline M}_F^2$.
\end{proof}
For $f(x)=x^{-s}$ and $0<x\le2B$, we may take $\mu=s(s+1)/(2B)^{s+2}$. Every surviving Laplacian has trace at most $2B$, so this interval contains its positive eigenvalues.

\subsection{Symmetrization with a stability remainder}\label{T02-sec-3}
Let $\pi$ range uniformly over vertex permutations, and let $\pi w$ denote the permuted edge weights. Edge transitivity gives $\mathbb E_\pi\pi w=u$. Also $J_s(\pi w)=J_s(w)$ because the attack family is permutation invariant. Fix a worst attack $\sigma_0$ at $u$. If $J_s(w)=+\infty$, \eqref{T02-eq:stability} is immediate; otherwise every matrix below is positive definite. By Lemma~\ref{T02-lem:convex},
\begin{align*}
 J_s(w)&\ge\mathbb E_\pi\tr L_{\sigma_0}(\pi w)^{-s}\\
 &\ge J_s(u)+\frac\mu2\mathbb E_\pi
 \norm{L_{\sigma_0}(\pi w)-L_{\sigma_0}(u)}_F^2.
\end{align*}
The Frobenius norm may be taken on $\R^n$: the difference annihilates $\one$, so its norm equals the restricted norm. Its off-diagonal entries imply
\[
 \norm{L_{\sigma_0}(\pi w)-L_{\sigma_0}(u)}_F^2
 \ge2\sum_e(q_e^{\sigma_0})^2((\pi w)_e-B/m)^2.
\]
For every fixed edge $e$, the random preimage $\pi^{-1}e$ is uniform on the $m$ edges. Thus, writing $V=\sum_e(w_e-B/m)^2$,
\[
 \mathbb E_\pi\norm{L_{\sigma_0}(\pi w)-L_{\sigma_0}(u)}_F^2
 \ge\frac{2V}{m}\sum_e(q_e^{\sigma_0})^2
 =\frac{2V}{m}\left(m-\sum_i(2\theta_i-\theta_i^2)\right).
\]
This proves \eqref{T02-eq:stability}, and hence existence and uniqueness of the stated optimizer, independently of the calculation of the worst attack.

\subsection{Computing the adversarial value}\label{T02-sec-4}
We include the requisite extremal argument to keep this paper independent of the preceding round.
\begin{lemma}[Weighted star comparison]\label{T02-lem:star}
Let $D$ be an oriented incidence matrix for $k$ distinct simple edges, and $Q=\diag(\sqrt{\theta})D^TD\diag(\sqrt{\theta})$. Then
\[
 \tr Q^r\le\tr C^r\quad(r\ge1),\qquad \lambda_{\max}(Q)\le\lambda_{\max}(C)<n.
\]
For a power series with nonnegative coefficients and strictly positive quadratic and cubic coefficients, equality of its traces occurs precisely for star support.
\end{lemma}
\begin{proof}
The diagonal entries of both matrices are $2\theta_i$, and $|Q_{ij}|\le C_{ij}=\sqrt{\theta_i\theta_j}$ off the diagonal. Closed-walk expansion proves every moment inequality; absolute-value Rayleigh comparison proves the spectral bound. Since $C\preceq(\max_i\theta_i)I+vv^T$, its norm is at most $\max_i\theta_i+\sum_i\theta_i<n$ under the hypotheses. At second order every disjoint edge pair contributes a strict gap $2\theta_i\theta_j$. A pairwise-intersecting family of distinct edges is a star or a triangle: start with $ab,ac$ and inspect an edge not meeting $a$. A triangle has a third-order gap $12\theta_1\theta_2\theta_3$ because the product of its three incidence Gram off-diagonal entries is negative. A star oriented away from its center has $Q=C$.
\end{proof}
At $u$, the surviving Laplacian on $\cH$ is $c(nI-XX^T)$ with $X=D\diag(\sqrt\theta)$. Nonzero eigenvalues of $XX^T$ agree with those of $Q$, so
\[
 \tr_{\cH}L_\sigma(u)^{-s}
 =c^{-s}\big((n-1-k)n^{-s}+\tr(nI-Q)^{-s}\big).
\]
Apply Lemma~\ref{T02-lem:star} to the positive-coefficient binomial series of $(n-x)^{-s}$. This proves \eqref{T02-eq:value} and all attack equality cases, completing Theorem~\ref{T02-thm:main}.

\begin{corollary}[Complete deletion of $k$ edges]
If $1\le k\le n-2$ and every $\theta_i=1$, then
\[
 \min_wJ_s(w)=\left(\frac{\binom n2}{B}\right)^s
 \left((n-k-1)^{-s}+(k-1)(n-1)^{-s}+(n-k-1)n^{-s}\right).
\]
Moreover, the coefficient in \eqref{T02-eq:stability} simplifies to
$s(s+1)(m-k)/(m(2B)^{s+2})$.
\end{corollary}
\begin{proof}
The eigenvalues of $C=I_k+\one\one^T$ are $k+1$ once and $1$ with multiplicity $k-1$.
\end{proof}

\subsection{Heat and tree criteria}\label{T02-sec-5}
\begin{theorem}
The uniform design is also the unique minimizer of
$\max_\sigma\tr_{\cH}e^{-tL_\sigma(w)}$ for each $t>0$, and the unique maximizer of $\min_\sigma\tau(L_\sigma(w))$. At the uniform design the worst attacks are precisely the stars.
\end{theorem}
\begin{proof}
For heat traces use Lemma~\ref{T02-lem:convex} with $f(x)=e^{-tx}$ on $[0,2B]$, whose second derivative is at least $t^2e^{-2Bt}>0$. The same symmetrization gives strict uniqueness, even when a survivor is disconnected. For tree sums, maximizing the minimum positive determinant is equivalent to minimizing the maximum of $-\log\det_{\cH}L_\sigma(w)$. Disconnection gives $+\infty$. Use $f(x)=-\log x$, with $f''\ge(2B)^{-2}$. The star calculation follows from the positive series for $e^{tx}$ and $-\log(n-x)$ as in Lemma~\ref{T02-lem:star}.
\end{proof}
For the tree criterion the optimal value has a closed form. With $d_i=n-\theta_i$ and $q=1-\sum_i\theta_i/d_i$ it is
\[
 \max_w\min_\sigma\tau(L_\sigma(w))
 =\frac{c^{n-1}}n\,n^{n-1-k}\left(\prod_i d_i\right)q.
\]
This is the determinant of a diagonal matrix minus $vv^T$.

\subsection{Context, novelty, and boundaries}\label{T02-sec-6}
The general convex optimization of graph spectra is classical \cite{ref039,ref036}; neither that observation nor group averaging is new. The preceding tournament paper \cite{ref040} supplies the weighted star comparison, reproduced here with proof. This round adds a two-player design problem, a unique simultaneous minimax solution, and a uniform quantitative rigidity inequality. It is not simply another evaluation of a fixed damaged graph. Exact priority of these combined formulas is not established, and no named open question is claimed solved.

There are no claims about sparse prescribed host graphs, attacks with $k\ge n$, adaptive budget reallocation, or heterogeneous installation costs. Those changes destroy part of the symmetry used in the proof. The stated results require no numerical computation or external data.




\clearpage
\section[T03 --- An Exponent-Independent Support Transition in Optimal Split-Graph Design]{An Exponent-Independent Support Transition in Optimal Split-Graph Design}\label{contrib-T03}
\markboth{Electrical Networks, Spectral Design, and Greedy Augmentation}{T03 --- Taisho}
\noindent\textsc{T03 \quad Taisho programme, source round 3}\par\medskip



\begin{quote}\small\textbf{Overview.} On the complete split graph $K_k\vee\overline K_m$, we solve conductance-budget optimization for every inverse power of the positive Laplacian spectrum. The optimum has two edge weights, determined explicitly. More surprisingly, the entire clique is assigned zero weight if and only if $m\ge k+1$, independently of the positive inverse-power exponent. Thus the optimal support undergoes one exponent-independent integer transition, from the full split graph to the complete bipartite graph. We also determine the determinant-optimal design, the algebraic-connectivity limit, and the exact cost for total effective resistance. The paper proves these statements from scalar inequalities and records the distinction from previously studied fixed-weight split-graph resistance formulas.\end{quote}
\subsection{Setup and main result}\label{T03-sec-1}
Let $k\ge2$, $m\ge1$, $n=k+m$, and $G=K_k\vee\overline K_m$. The first part is a clique of $k$ vertices and the second an independent set of $m$ vertices; every cross edge is present. Give the available edges nonnegative conductances of total $B>0$. A disconnected positive support has infinite inverse-power objective. For $s>0$, minimize
\[
 F_s(w)=\sum_{j=2}^n\lambda_j(L_w)^{-s}.
\]
Set
\[
 \gamma=\frac{k+1}{k-1},\quad A_s=(k-1)k^s,\quad
 D_s=(m-1)m^s+\left(\frac{km}{n}\right)^s,
 \quad t_s=\left(\frac{\gamma A_s}{D_s}\right)^{1/(s+1)}.
\]
\begin{theorem}\label{T03-thm:main}
The minimizer is unique. If $m\ge k+1$, every clique edge has weight zero and every cross edge has weight $B/(km)$. If $m\le k$, the optimal weights are
\begin{equation}\label{T03-eq:weights}
 a=\frac{2Bx_s}{k(k-1)}\quad\text{on clique edges},\qquad
 b=\frac{B(1-x_s)}{km}\quad\text{on cross edges},
 \qquad x_s=\frac{t_s-1}{t_s+\gamma}>0.
\end{equation}
The optimal value is
\begin{equation}\label{T03-eq:min}
 \min F_s=B^{-s}
 \begin{cases}
 A_s+D_s,&m\ge k+1,\\[2mm]
 \displaystyle\frac{\big(A_s^{1/(s+1)}+\gamma^{s/(s+1)}D_s^{1/(s+1)}\big)^{s+1}}
 {(1+\gamma)^s},&m\le k.
 \end{cases}
\end{equation}
\end{theorem}
No minimizer has disconnected support: in the zero-clique regime, the surviving $K_{k,m}$ remains connected. The phase transition concerns optimal support, not the presence of edges in the given host.

\subsection{Why a two-weight calculation is globally valid}\label{T03-sec-2}
\begin{lemma}\label{T03-lem:strict}
For $s>0$, $M\mapsto\tr M^{-s}$ is strictly convex on positive definite symmetric matrices. Consequently $F_s$ is strictly convex in the conductances of a simple graph on its connected domain.
\end{lemma}
\begin{proof}
Let $f(x)=x^{-s}$. In an orthonormal eigenbasis $z_i$ of $tM+(1-t)N$, scalar Jensen gives
\[
 \tr f(tM+(1-t)N)\le t\sum_if(z_i^TMz_i)+(1-t)\sum_if(z_i^TNz_i)
 \le t\tr f(M)+(1-t)\tr f(N).
\]
The second inequality follows by diagonalizing $M,N$ separately and applying scalar Jensen to squared overlaps. Equality, with $0<t<1$, forces each $z_i$ to lie in an eigenspace of $M$ and of $N$, and forces the corresponding two eigenvalues to coincide. Thus $M=N$. For graph Laplacians, equality on $\one^\perp$ implies equality on all of $\R^n$, and off-diagonal entries recover every individual conductance.
\end{proof}
Average any design under the automorphism group $S_k\times S_m$. Convexity cannot increase $F_s$, and strictness shows that every optimum is invariant. Hence clique edges have a common weight $a$, and cross edges a common weight $b$. Existence follows either from the explicit minimum below or from lower semicontinuity on the compact budget simplex with infinite boundary values at disconnection.

Let $x$ be the fraction of the total budget on clique edges. Then $a=2Bx/[k(k-1)]$, $b=B(1-x)/(km)$, and $0\le x<1$. Decompose $\one^\perp$ into vectors supported on the clique with sum zero, vectors supported on the independent part with sum zero, and the one-dimensional space constant on each part. Direct multiplication gives eigenvalues
\begin{equation}\label{T03-eq:spectrum}
 \frac{B(1+\gamma x)}k\quad(k-1\text{ times}),\qquad
 \frac{B(1-x)}m\quad(m-1\text{ times}),\qquad
 \frac{Bn(1-x)}{km}\quad(1\text{ time}).
\end{equation}
The second multiplicity is zero when $m=1$. Therefore
\begin{equation}\label{T03-eq:scalar}
 B^sF_s(x)=\frac{A_s}{(1+\gamma x)^s}+\frac{D_s}{(1-x)^s}.
\end{equation}
This strictly convex scalar function has its minimum at $x=0$ exactly when $D_s\ge\gamma A_s$. Otherwise its derivative vanishes exactly at \eqref{T03-eq:weights}.

\subsection{The exponent-independent threshold}\label{T03-sec-3}
\begin{lemma}\label{T03-lem:threshold}
For every real $s>0$,
\[
 D_s<(k+1)k^s\ \Longleftrightarrow\ m\le k.
\]
\end{lemma}
\begin{proof}
For fixed $k,s$, both summands in $D_s$ are nondecreasing in $m$, and their sum is strictly increasing. At $m=k$,
\[
 D_s=(k-1)k^s+(k/2)^s<(k+1)k^s.
\]
At $m=k+1$, divide by $k^s$ and define
\[
 h(s)=k(1+1/k)^s+\left(\frac{k+1}{2k+1}\right)^s.
\]
We have $h(0)=k+1$ and $h''(s)>0$. Moreover
\[
 h'(0)=\log\left((1+1/k)^k\frac{k+1}{2k+1}\right)>0.
\]
Indeed $(1+1/k)^k>2$ for $k\ge2$ by the binomial theorem, while $(2k+1)/(k+1)<2$. Thus $h(s)>k+1$ for all $s>0$. Monotonicity in the integer $m$ completes the proof.
\end{proof}
Since $\gamma A_s=(k+1)k^s$, the lemma proves the support transition in Theorem~\ref{T03-thm:main}. To evaluate the interior minimum, put $u=1+\gamma x$, $v=1-x$, so $u+\gamma v=1+\gamma$. For positive $u,v$, scalar H\"older gives
\[
 \frac{A_s}{u^s}+\frac{D_s}{v^s}\ge
 \frac{\big(A_s^{1/(s+1)}+\gamma^{s/(s+1)}D_s^{1/(s+1)}\big)^{s+1}}
 {(1+\gamma)^s}.
\]
Equality is at the derivative solution. This completes the proof of Theorem~\ref{T03-thm:main}.

\subsection{Resistance, determinant, and bottleneck endpoints}\label{T03-sec-4}
For $s=1$, $D_1=m^2(n-1)/n$ and $t_1=\sqrt{k(k+1)n/[m^2(n-1)]}$. The minimal Kirchhoff index is $n\min F_1$. In the regime $m\ge k+1$ it is
\[
 \frac nB\left(k(k-1)+\frac{m^2(n-1)}n\right).
\]
This provides an exact algebraic formula without solving a semidefinite program.

\begin{proposition}[Maximum weighted tree sum]
The unique determinant-optimal design is also two-weight. Its clique budget fraction is
\[
 x_D=\max\left\{0,\frac{(k-1)(k+1-m)}{(k+1)(k+m-1)}\right\}.
\]
Its support has the same integer threshold $m\ge k+1$.
\end{proposition}
\begin{proof}
The function $-\log\det(L|_{\one^\perp})$ is strictly convex by the proof of Lemma~\ref{T03-lem:strict}, so symmetry again applies. From \eqref{T03-eq:spectrum}, maximizing the determinant is equivalent to maximizing
$(k-1)\log(1+\gamma x)+m\log(1-x)$. Its derivative is
$(k+1)/(1+\gamma x)-m/(1-x)$, giving the claimed expression and clipping at zero.
\end{proof}
\begin{proposition}[A symmetric algebraic-connectivity optimizer]
For $m\ge2$, a design maximizing $\lambda_2$ is given by
\[
 x_E=\max\left\{0,\frac{k-m}{k+m\gamma}\right\}.
\]
It is the limit of $x_s$ as $s\to\infty$. For $m=1$ the graph is complete and the uniform design, $x_E=(k-1)/(k+1)$, is optimal.
\end{proposition}
\begin{proof}
The minimum eigenvalue of a symmetric matrix is concave because it is the infimum of linear Rayleigh quotients. Averaging therefore yields a symmetric maximizer, not necessarily uniqueness. For $m\ge2$ the mixed eigenvalue in \eqref{T03-eq:spectrum} is larger than the independent-part eigenvalue. We maximize the minimum of an increasing and a decreasing affine function, obtaining their intersection if nonnegative, otherwise $x=0$. The limit follows from $t_s\to k/m$ for $m\ge2$. When $m=1$, $t_s=k+1$ exactly and all edge weights are equal.
\end{proof}

\subsection{Context and exact status}\label{T03-sec-5}
Convex optimization of graph Laplacian eigenvalues and effective resistance is classical \cite{ref039,ref036}. Fixed or vertex-induced weighted complete split graphs have also been studied \cite{ref041}. The present task optimizes \emph{all} available edge conductances under a budget and proves an exponent-independent support threshold, rather than merely computing resistance for a chosen weighting. Priority of this exact formula has not been certified. No named open conjecture or benchmark superiority is asserted. All cases, including $m=1$ and $m=k+1$, are covered analytically, and no numerical certificate is needed.




\clearpage
\section[T04 --- Fractional-Power Optimal Designs and a Sharp Universal Laplacian Threshold]{Fractional-Power Optimal Designs and a Sharp Universal Laplacian Threshold}\label{contrib-T04}
\markboth{Electrical Networks, Spectral Design, and Greedy Augmentation}{T04 --- Taisho}
\noindent\textsc{T04 \quad Taisho programme, source round 4}\par\medskip



\begin{quote}\small\textbf{Overview.} For an arbitrary positive definite demand covariance on the zero-sum subspace, we solve the trace-budget relaxation of demand-weighted inverse-power network design for every positive exponent. The optimizer is a normalized fractional power of the demand matrix. We then characterize when this matrix is an admissible nonnegative-conductance Laplacian. For demand Laplacians it is always admissible and has complete support. For arbitrary covariances we prove a sharp universal condition-number threshold: every eigenbasis is admissible exactly when $\kappa^{1/(s+1)}\le n/(n-2)$. A rank-one construction proves necessity. Thus sufficiently large inverse-power exponents always make the unconstrained optimizer physically realizable, with an explicit sharp uniform exponent. The square-root special case and fractional-Laplacian positivity are acknowledged as existing ingredients.\end{quote}
\subsection{Demand-weighted design}\label{T04-sec-1}
Let $n\ge3$, $\cH=\one^\perp$, and let $C$ be symmetric with $C\one=0$ and $C|_{\cH}\succ0$. It represents a covariance of zero-sum injections, not necessarily a graph Laplacian. Let $B>0$ and $s>0$. A complete-host design is any nonnegative-conductance Laplacian $L$ with connected support and $\tr L=2B$. Its objective is
\[
 \mathcal J_s(L)=\tr_{\cH}(C L^{-s}).
\]
For $s=1$, this is expected dissipated energy under an injection with covariance $C$. If $C=\sum d_{ij}(e_i-e_j)(e_i-e_j)^T$ with $d_{ij}\ge0$, it is a nonnegative demand-weighted sum of effective resistances.

Write $r=1/(s+1)$, $T=\tr C^r$, and $P=I-\one\one^T/n$.
\begin{theorem}[Unconstrained optimum and graph feasibility]\label{T04-thm:opt}
Among all symmetric $L$ with $L\one=0$, $L|_{\cH}\succ0$, and $\tr L=2B$, the unique minimizer is
\[
 L_\star=\frac{2B}{T}C^r,\qquad
 \min\mathcal J_s=\frac{T^{s+1}}{(2B)^s}.
\]
The same value is attained by a nonnegative-conductance graph if and only if every off-diagonal entry of $C^r$ is nonpositive. If $C$ is a connected demand Laplacian, all these entries are strictly negative, for every $s>0$.
\end{theorem}
The statement applies to all $s>0$ without asserting that $L\mapsto\tr(C L^{-s})$ is convex for every such $s$. Its proof uses scalar spectral Jensen instead.

\subsection{A sharp matrix-budget inequality}\label{T04-sec-2}
\begin{lemma}\label{T04-lem:budget}
For positive definite $d\times d$ matrices $C,M$ and $s>0$,
\[
 \tr(CM^{-s})(\tr M)^s\ge(\tr C^{1/(s+1)})^{s+1},
\]
with equality exactly when $M$ is a positive scalar multiple of $C^{1/(s+1)}$.
\end{lemma}
\begin{proof}
Take an orthonormal eigenbasis $v_i$ of $C$, with positive eigenvalues $\mu_i$, and put $x_i=v_i^TMv_i$. Diagonalizing $M$ and applying strict scalar convexity of $x^{-s}$ gives
\[
 v_i^TM^{-s}v_i\ge x_i^{-s}.
\]
Equality holds exactly when $v_i$ lies within one eigenspace of $M$. Next, H\"older with exponents $s+1$ and $(s+1)/s$ gives
\[
 \sum_i\mu_i^{1/(s+1)}
 =\sum_i(\mu_i x_i^{-s})^{1/(s+1)}x_i^{s/(s+1)}
 \le\left(\sum_i\mu_i x_i^{-s}\right)^{1/(s+1)}
       \left(\sum_i x_i\right)^{s/(s+1)}.
\]
Equality in H\"older requires $x_i$ proportional to $\mu_i^{1/(s+1)}$. Since all $\mu_i>0$, equality in the first step is required for every $i$, forcing $M$ to have precisely those eigenspaces and eigenvalues. Conversely the proportional fractional power gives equality directly.
\end{proof}
Apply the lemma on $\cH$. A symmetric matrix annihilating $\one$ is a nonnegative-conductance Laplacian exactly when its off-diagonals are nonpositive; its edge conductances are $-L_{ij}$. This proves the first two assertions of Theorem~\ref{T04-thm:opt}.

\subsection{Why fractional demand Laplacians are feasible}\label{T04-sec-3}
\begin{lemma}\label{T04-lem:fractional}
If $C$ is a connected nonnegative-conductance Laplacian and $0<r<1$, then $C^r\one=0$ and $(C^r)_{ij}<0$ for $i\ne j$.
\end{lemma}
\begin{proof}
Zero row sums follow from spectral calculus. For $t>0$, the resolvent $(tI+C)^{-1}$ has strictly positive entries. Indeed choose $d$ at least the maximum weighted degree and write $A=dI-C\ge0$ entrywise. Its row sums are $d$. Then
\[
 (tI+C)^{-1}=\frac1{t+d}\sum_{\ell\ge0}\left(\frac A{t+d}\right)^\ell.
\]
Connectivity supplies a positive path contribution between every pair of vertices, so the sum is strictly positive.

The scalar identity
\[
 x^r=\frac{\sin(\pi r)}\pi\int_0^\infty t^{r-1}\frac{x}{t+x}\,dt
\]
follows by substituting $t=xu$ and evaluating the beta integral; it also holds at $x=0$. Diagonalization therefore gives
\[
 C^r=\frac{\sin(\pi r)}\pi\int_0^\infty t^{r-1}C(tI+C)^{-1}\,dt.
\]
For $i\ne j$, the integrand entry is $-t^r(tI+C)^{-1}_{ij}<0$. The integral converges by the scalar formula, and is strictly negative.
\end{proof}
This proves the remaining assertion of Theorem~\ref{T04-thm:opt}. In particular, even sparse demand graphs lead to dense optimal conductance designs when every pair is available for installation.

\subsection{A sharp universal sign threshold}\label{T04-sec-4}
Let $\lambda_{\min},\lambda_{\max}$ be the extreme eigenvalues of $C|_{\cH}$ and $\kappa=\lambda_{\max}/\lambda_{\min}$.
\begin{theorem}[Sharp uniform feasibility threshold]\label{T04-thm:threshold}
Every such $C$ with condition number at most $\kappa$ has a feasible fractional-power optimizer whenever
\begin{equation}\label{T04-eq:threshold}
 \kappa^{1/(s+1)}\le\frac n{n-2}.
\end{equation}
Strict inequality guarantees strictly positive conductance on every edge. Conversely, if \eqref{T04-eq:threshold} fails, there exists a covariance $C$ of condition number exactly $\kappa$ for which the unconstrained optimizer has a positive off-diagonal entry and hence is infeasible. Thus the threshold is sharp uniformly over eigenvectors.
\end{theorem}
\begin{proof}
We first prove a general bound. Suppose $M\one=0$ and the eigenvalues on $\cH$ lie in $[a,b]$, with $0<a\le b$. Write $M=aP+(b-a)K$, where $K\one=0$ and $0\preceq K\preceq P$. For $i\ne j$, put $u=Pe_i$, $v=Pe_j$. Among positive contractions on $\cH$,
\begin{equation}\label{T04-eq:contraction}
 \max_K u^TKv=\frac{u^Tv+\norm u\norm v}{2}=\frac{n-2}{2n}.
\end{equation}
To see this, let $S=(uv^T+vu^T)/2$. Its only possibly nonzero eigenvalues are $(u^Tv\pm\norm u\norm v)/2$. The maximum of $\tr(KS)$ over $0\preceq K\preceq I$ is the positive eigenvalue, achieved by the projection onto its eigenspace. This proves \eqref{T04-eq:contraction}. Therefore
\begin{equation}\label{T04-eq:entry}
 M_{ij}\le -\frac an+(b-a)\frac{n-2}{2n}
 =\frac{(n-2)b-na}{2n}.
\end{equation}
Apply this with $M=C^r$, $a=\lambda_{\min}^r$, $b=\lambda_{\max}^r$ to prove sufficiency and strictness.

For necessity, fix $i\ne j$ and take the unit vector
\[
 z=\frac{Pe_i+Pe_j}{\norm{Pe_i+Pe_j}}.
\]
Its $i$ and $j$ coordinates are both $\sqrt{(n-2)/(2n)}$. Set
$C=P+(\kappa-1)zz^T$. Its spectrum on $\cH$ is $\kappa$ once and $1$ with multiplicity $n-2$, so
$C^r=P+(\kappa^r-1)zz^T$. Its $(i,j)$ entry is exactly $((n-2)\kappa^r-n)/(2n)$, positive when the threshold fails.
\end{proof}
\begin{corollary}[Eventual realizability]
For every fixed positive covariance $C|_{\cH}$, the unconstrained optimum is realizable for all $s>0$ satisfying
\[
 s+1\ge\frac{\log\kappa}{\log(n/(n-2))}.
\]
The bound is the best possible one depending only on $n$ and $\kappa$. As $s\to\infty$, $L_\star\to 2BP/(n-1)$, the uniform complete-graph design.
\end{corollary}
\begin{proof}
The inequality is \eqref{T04-eq:threshold}. On $\cH$, $C^{1/(s+1)}\to I$, and its trace tends to $n-1$.
\end{proof}
The threshold is sufficient for a particular covariance, not necessary: a demand Laplacian remains feasible for every $s$, even far outside this universal condition-number range.

\subsection{An explicit non-Laplacian covariance}\label{T04-sec-5}
For $n=3$ take $z=(1,1,-2)^T/\sqrt6$ and $C=P+(\kappa-1)zz^T$. Its $(1,2)$ entry is $(\kappa-3)/6$, so $C$ itself is not a graph Laplacian when $\kappa>3$. Nevertheless, its optimal inverse-power design is feasible exactly when $\kappa^{1/(s+1)}\le3$. Thus a valid injection covariance need not be a demand Laplacian, and increasing the inverse-power exponent can remove a genuine sign obstruction.

\subsection{Provenance and limitations}\label{T04-sec-6}
The square-root optimization mechanism occurs in the supplied benchmark's block-demand designs; it is not presented as a new discovery. This paper extends the calculation to every positive inverse power, arbitrary positive injection covariances, and the sharp eigenbasis-uniform feasibility threshold. Fractional graph Laplacians and their nonlocal positive conductances are well-established \cite{ref042,ref043}. General convex resistance design is also classical \cite{ref036}. Exact priority of the combined theorem is unverified, and no named open problem is claimed solved.

The complete host is essential: feasibility of $C^r$ does not make it feasible on a sparse prescribed host. The assumption $C|_{\cH}\succ0$ is used for uniqueness and to avoid nonattained infima associated with unused demand directions. All proofs are analytic and need no data or numerical certificates.




\clearpage
\section[T05 --- Exact Resistance-Diameter Design on Trees with Heterogeneous Costs]{Exact Resistance-Diameter Design on Trees with Heterogeneous Costs}\label{contrib-T05}
\markboth{Electrical Networks, Spectral Design, and Greedy Augmentation}{T05 --- Taisho}
\noindent\textsc{T05 \quad Taisho programme, source round 5}\par\medskip



\begin{quote}\small\textbf{Overview.} For an arbitrary tree with positive heterogeneous installation costs, we give an exact radical formula and a linear-arithmetic algorithm for the smallest possible resistance diameter under a conductance budget. The formula combines rooted subtree coefficients with a sharp balance test on every edge; an optimal metric center may lie strictly inside an edge, so minimizing over vertex-centered radius designs alone is insufficient. All optimal edge lengths can be reconstructed by local formulas. The rooted-radius recursion is already present in the budget-allocation literature. The additional result proved here is the diameter optimization with movable interior centers, its exact edge-center criterion, and the reduction to a finite set of explicit candidates. This is a tree-case result, not a solution of the general-graph budget-diameter problem.\end{quote}
\subsection{Problem and finite formula}\label{T05-sec-1}
Let $T=(V,E)$ be a tree with $|V|=n\ge2$ and positive costs $c_e$. Choose conductances $w_e>0$ with $\sum_ec_ew_e=B$, where $B>0$. The effective resistance between $u,v$ is the sum of $r_e=1/w_e$ along their unique path. The resistance diameter is
\[
 D(w)=\max_{u,v\in V}\sum_{e\in P(u,v)}r_e.
\]
Equivalently, we choose positive edge lengths $r_e$ with cost $\sum_ec_e/r_e=B$.

For each oriented edge $uv$, let $T_{u\mid v}$ be the component containing $u$ when $uv$ is removed, rooted at $u$. Define directed coefficients recursively by
\begin{equation}\label{T05-eq:rec}
 K_{u\mid v}=\sum_{z\sim u,\ z\ne v}
 \left(\sqrt{c_{uz}}+\sqrt{K_{z\mid u}}\right)^2.
\end{equation}
An empty sum is zero. Also define
\[
 K_u=\sum_{z\sim u}\left(\sqrt{c_{uz}}+\sqrt{K_{z\mid u}}\right)^2.
\]
An edge $e=uv$ is \emph{balanced} if
\begin{equation}\label{T05-eq:balance}
 \left|\sqrt{K_{u\mid v}}-\sqrt{K_{v\mid u}}\right|\le\sqrt{c_e}.
\end{equation}
\begin{theorem}[Exact budget diameter]\label{T05-thm:main}
Put
\[
 Q=\min\left(\{2K_u:u\in V\}\ \cup\
 \left\{\left(\sqrt{K_{u\mid v}}+\sqrt{K_{v\mid u}}+\sqrt{c_{uv}}\right)^2:
 uv\in E\text{ is balanced}\right\}\right).
\]
Then
\[
 \min_{\sum c_ew_e=B}D(w)=Q/B.
\]
All coefficients, candidate values, and a minimizing allocation can be computed in $O(n)$ arithmetic operations involving addition, multiplication, division, and square roots. This is an exact real-arithmetic bound, not a bit-complexity claim for nested radicals.
\end{theorem}

\subsection{Rooted radius costs}\label{T05-sec-2}
For a rooted tree $S$ and $R>0$, let $F_S(R)$ be the minimum of $\sum c_e/r_e$ over positive edge lengths such that every root-to-vertex distance is at most $R$.
\begin{lemma}[Rooted recursion]\label{T05-lem:root}
If the root has child subtrees with coefficients $K_i$ and incident costs $c_i$, then
\[
 F_S(R)=\frac{K_S}{R},\qquad
 K_S=\sum_i(\sqrt{c_i}+\sqrt{K_i})^2.
\]
For a child with $K_i>0$, the optimal incident length and remaining radius are
\begin{equation}\label{T05-eq:reconstruct}
 t_i=R\frac{\sqrt{c_i}}{\sqrt{c_i}+\sqrt{K_i}},\qquad
 R_i=R\frac{\sqrt{K_i}}{\sqrt{c_i}+\sqrt{K_i}}.
\end{equation}
For a leaf child, $K_i=0$, the incident length is $R$ and no remaining subtree is present.
\end{lemma}
\begin{proof}
The claim is immediate for a one-vertex tree. For each branch choose its incident length $t\in(0,R)$ and use induction on the child subtree. The branch cost is
\[
 \frac{c_i}{t}+\frac{K_i}{R-t}\ge\frac{(\sqrt{c_i}+\sqrt{K_i})^2}{R},
\]
by Cauchy--Schwarz, with equality precisely at \eqref{T05-eq:reconstruct}. For $K_i=0$ the minimum is $c_i/R$, attained at $t=R$. Different branches have separate root-radius constraints, so their costs add. Every edge length obtained is positive.
\end{proof}
This recursion is the established budget-radius mechanism \cite{ref044}. It is reproduced for completeness, not claimed as new. Applying it at a vertex $u$ with radius $D/2$ gives a feasible diameter-$D$ design of cost $2K_u/D$.

\subsection{The movable center inside an edge}\label{T05-sec-3}
\begin{lemma}[Center theorem for a metric tree]\label{T05-lem:center}
A tree with positive edge lengths has diameter at most $D$ if and only if some point of its geometric realization is within distance $D/2$ of every vertex.
\end{lemma}
\begin{proof}
The reverse implication is the triangle inequality. For the forward direction, take a longest vertex-to-vertex path of length $d\le D$ and its midpoint $o$. For any vertex $z$, let $p$ be the point where its path first meets that diameter path. If $z$ were farther than $d/2$ from $o$, then the path from $z$ to the farther endpoint of the diameter path would have length greater than $d$, a contradiction. Thus every vertex is within $d/2\le D/2$ of $o$.
\end{proof}
Suppose the center is in edge $uv$. Write $R=D/2$, and let its distances to $u,v$ be $t_u,t_v\ge0$. The edge length is $t_u+t_v$. The radii available for the two rooted components are $x=R-t_u$ and $y=R-t_v$. Thus $0\le x,y\le R$, and the cost to minimize is
\begin{equation}\label{T05-eq:edgecost}
 \frac{c_{uv}}{2R-x-y}+\frac{K_{u\mid v}}x+\frac{K_{v\mid u}}y.
\end{equation}
A term $0/0$ means an empty rooted component and is zero; a positive numerator with zero denominator is forbidden. The edge denominator must be positive.

Put $A=\sqrt{K_{u\mid v}}$, $H=\sqrt{K_{v\mid u}}$, $q=\sqrt{c_{uv}}$. Ignoring the two upper bounds for the moment, Cauchy--Schwarz gives
\begin{equation}\label{T05-eq:edgeformula}
 \frac{q^2}{2R-x-y}+\frac{A^2}x+\frac{H^2}y
 \ge\frac{(A+H+q)^2}{2R}.
\end{equation}
Equality requires
\begin{equation}\label{T05-eq:edgex}
 x=\frac{2RA}{A+H+q},\quad
 y=\frac{2RH}{A+H+q},\quad
 r_{uv}=\frac{2Rq}{A+H+q}.
\end{equation}
These values satisfy $x,y\le R$ exactly when $|A-H|\le q$, the balance test \eqref{T05-eq:balance}. The center is strictly inside the edge exactly when both inequalities are strict. An equality corresponds to a center at an endpoint.

\begin{lemma}[No missing unbalanced candidate]\label{T05-lem:unbalanced}
If $A>H+q$, the minimum of \eqref{T05-eq:edgecost} under $x,y\le R$ occurs at $x=R$ and equals $(A^2+(H+q)^2)/R=K_u/R$. The symmetric statement holds if $H>A+q$.
\end{lemma}
\begin{proof}
For a fixed $x$, first minimize in $y$ without the bound $y\le R$; the remaining cost is
$A^2/x+(H+q)^2/(2R-x)$. Its unconstrained minimizing $x$ is $2RA/(A+H+q)>R$, so it is decreasing up to $R$. At $x=R$, the minimizing $y=RH/(H+q)$ satisfies $y\le R$, so the relaxed minimum is feasible and globally optimal. The value agrees with the root recursion at $u$ because its branch toward $v$ contributes $(H+q)^2$.
\end{proof}
The formulas also apply when one rooted side is a single vertex: its coefficient and allocated radius are both zero, while all actual edge lengths remain positive.

\subsection{Global proof, reconstruction, and algorithm}\label{T05-sec-4}
Fix a diameter bound $D>0$. By Lemma~\ref{T05-lem:center}, every feasible length assignment has a center at a vertex or inside some edge. Lemma~\ref{T05-lem:root}, \eqref{T05-eq:edgeformula}, and Lemma~\ref{T05-lem:unbalanced} therefore show that its cost is at least $Q/D$. Conversely, each candidate in $Q$ has an explicitly feasible realization: use \eqref{T05-eq:reconstruct} for a vertex center, or first use \eqref{T05-eq:edgex} and then reconstruct the two rooted components. Hence the least cost for diameter at most $D$ is exactly $Q/D$.

Set $D=Q/B$. The construction has cost exactly $B$ and diameter at most $D$. No smaller diameter is possible, since its minimum cost would exceed $B$. This proves the value in Theorem~\ref{T05-thm:main}.

For the algorithm, root the combinatorial tree arbitrarily. A bottom-up pass computes all child-to-parent coefficients by \eqref{T05-eq:rec}. A top-down pass computes each parent's coefficient toward a child by taking its full sum and subtracting that child's contribution. Each directed message is used a constant number of times; sums over all degrees cost $O(n)$. Evaluate the $n$ vertex candidates and $n-1$ edge candidates, retain a minimum, and propagate the reconstruction formulas once. No optimization oracle is needed.

\subsection{Examples and the importance of interior centers}\label{T05-sec-5}
On a path with $h$ edges, repeated recursion gives $\sqrt{K}$ equal to the sum of $\sqrt{c_e}$ on the corresponding subpath. A balanced center exists at the weighted midpoint, and
\[
 Q=\left(\sum_{e\in E}\sqrt{c_e}\right)^2,\qquad
 w_e=\frac{B}{\sqrt{c_e}\sum_f\sqrt{c_f}}.
\]
This agrees with direct minimization of the sum of all resistances.

For a two-edge path with costs $1$ and $9$, the optimal diameter coefficient is $(1+3)^2=16$. The internal vertex-centered coefficient is $2(1+9)=20$, and the leaf-centered coefficients are $32$. Thus searching only vertex-centered radius designs gives a strictly suboptimal answer. The correct metric center lies inside the cost-$9$ edge.

For a unit-cost star with $h\ge2$ leaves, $K$ at the hub is $h$, hence $Q=2h$, attained by equal conductances. If $h>2$, every edge is unbalanced because $\sqrt{h-1}>1$, so the hub is the only possible optimal-center type among the candidates.

\subsection{Research status and limitations}\label{T05-sec-6}
Ben-Moshe, Elkin, Omri and Gottlieb introduced the inverse-budget edge-length model and solved rooted and vertex-centered radius optimization on trees \cite{ref044}. Their conclusion also raises budget diameter as a further graph-optimization direction. This paper settles the fixed-tree diameter calculation with explicit interior-center formulas; it does \emph{not} claim to settle their unrestricted general-graph direction. The generic effective-resistance formulation is classical \cite{ref036}. A comprehensive priority search for the precise tree diameter formula remains open.

All costs are strictly positive, conductances are continuous, and there are no lower or upper capacity bounds. Zero costs, integer conductances, cycles in the host, and additional geometric constraints require different arguments. The proof and reconstruction are contained in this single source file.




\clearpage
\section[T06 --- Sharp Random-Tree Energy Inflation:\ Exact Bond Ranges and Theta-Core Rigidity]{Sharp Random-Tree Energy Inflation:\ Exact Bond Ranges and Theta-Core Rigidity}\label{contrib-T06}
\markboth{Electrical Networks, Spectral Design, and Greedy Augmentation}{T06 --- Taisho}
\noindent\textsc{T06 \quad Taisho programme, source round 6}\par\medskip



\begin{quote}\small\textbf{Overview.} Fix a connected graph and a nonzero positive semidefinite covariance of zero-sum electrical demands. Sample a spanning tree with probability proportional to the product of its conductances, retaining those conductances in the tree. We determine the exact range, over all positive conductances, of the ratio between expected tree energy and original network energy. Its endpoints are the minimum and maximum cardinalities of bonds on which the demand has nonzero variance. A two-forest formula gives the exact defect from the cycle-rank bound. For one terminal pair, equality in that bound holds precisely for trees or a theta core joining the terminals with trees attached. For total pairwise resistance, equality holds precisely for trees and cycles. The underlying tree--forest bijection and resistance enumeration formula are classical; the focus is their sharp conductance-extremal and equality consequences.\end{quote}
\subsection{Statement and conventions}\label{T06-sec-1}
Let $G=(V,E)$ be a finite connected simple graph, $n\ge2$, $m=|E|$, and $\beta=m-n+1$. Each edge has conductance $w_e>0$. Write
\[
 L=\sum_{e\in E}w_e b_e b_e^T,\qquad \tau=\sum_{T\text{ spanning tree}}w(T),\qquad w(H)=\prod_{e\in H}w_e,
\]
where $b_e$ is either orientation of the incidence vector. The Moore--Penrose inverse $L^+$ is the inverse on $\one^\perp$ and zero on constants. Fix $C\succeq0$, $C\one=0$, $C\ne0$, and put
\[
 \mathcal E_C(G,w)=\tr(CL^+),\qquad
 \rho_C(w)=\frac{\mathbb E\,\tr(CL_T^+)}{\tr(CL^+)},
 \quad \mathbb P(T)=w(T)/\tau.
\]
A \emph{bond partition} $S\mid S^c$ has nonempty connected induced sides. It is unordered. Set $q_C(S)=\one_S^TC\one_S=q_C(S^c)$, and call the bond active if $q_C(S)>0$. Let
\[
 b_-(C)=\min_{\text{active bonds}}|\delta(S)|,\qquad
 b_+(C)=\max_{\text{active bonds}}|\delta(S)|.
\]
\begin{theorem}[Exact range]\label{T06-range}
Active bonds exist. If $b_-<b_+$, then
\[
 \{\rho_C(w):w\in(0,\infty)^E\}=(b_-,b_+).
\]
If $b_-=b_+$, the ratio is identically that common integer. In particular $\rho_C(w)\le b_+\le\beta+1$, and $b_+$ is the best constant uniform over conductances.
\end{theorem}
For $C=(e_s-e_t)(e_s-e_t)^T$, active bonds are exactly the bonds separating $s,t$. Thus the lower endpoint is the ordinary minimum $s$--$t$ edge-cut cardinality and the upper endpoint is the maximum $s$--$t$ bond cardinality. The theorem concerns arbitrary positive conductances, not only unit conductances, and does not assert efficient computation of a maximum bond.

\subsection{An exact forest identity}\label{T06-sec-2}
Let $\mathcal F_2$ be the spanning forests with exactly two components. If their vertex sets are $S_F,S_F^c$, abbreviate $q_F=q_C(S_F)$ and $d_F=|\delta(S_F)|$.
\begin{lemma}[Demand-weighted forest formulas]\label{T06-forest}
One has
\begin{align}
 \tau\,\tr(CL^+)&=\sum_{F\in\mathcal F_2}w(F)q_F,\label{T06-den}\\
 \tau\,\mathbb E\tr(CL_T^+)&=\sum_{F\in\mathcal F_2}w(F)q_Fd_F.\label{T06-num}
\end{align}
\end{lemma}
\begin{proof}
First recall the weighted matrix-tree identity. Delete a row of the oriented incidence matrix $B$ to obtain $\widetilde B$. Cauchy--Binet applied to $\widetilde B\diag(w)\widetilde B^T$ gives its determinant as $\tau$: an $(n-1)$-column minor is zero unless it is a tree, and a tree minor is $\pm1$, by successively removing leaves.

Add a new edge of conductance $z$ between $i,j$, allowing a parallel edge for this calculation. Expanding tree weights according to whether the new edge is present gives $\tau(z)=\tau+z\sum_{F:i\mid j}w(F)$. The determinant lemma on a reduced Laplacian gives $\tau(z)=\tau(1+zR_{ij})$, where $R_{ij}=(e_i-e_j)^TL^+(e_i-e_j)$. Consequently
\[
 \tau R_{ij}=\sum_{F:i\mid j}w(F).
\]
For a zero-sum vector $a$, the identity $a^TL^+a=-\frac12\sum_{i,j}a_i a_jR_{ij}$ now implies
\[
 \tau a^TL^+a=\sum_F w(F)\Big(\sum_{i\in S_F}a_i\Big)^2.
\]
Indeed the contribution of one separating partition is
$-\sum_{i\in S_F,j\in S_F^c}a_i a_j=(\sum_{S_F}a_i)^2$.
Decompose $C$ into sums $aa^T$ with $a\perp\one$ to obtain \eqref{T06-den}.

In a tree, conservation of current forces the current on an edge $e$ to equal the total injection on one component of $T-e$. Its dissipated energy is that current squared divided by $w_e$. Thus
\[
 \tr(CL_T^+)=\sum_{e\in T}\frac{q_C(S_{T-e})}{w_e}.
\]
In the expectation, $w(T)/w_e=w(T-e)$. Removing $e$ gives a two-forest; conversely each $F$ has exactly $d_F$ crossing edges that can be added to make a tree. This proves \eqref{T06-num}.
\end{proof}
The tree--forest edge-addition bijection is classical, as is the separating-forest resistance formula; see \cite{ref045,ref046}. Here they are combined with general demand weights and conductance degeneration.

\begin{proof}[Proof of Theorem \ref{T06-range}]
The denominator in \eqref{T06-den} is positive since $L^+$ is positive definite on $\one^\perp$ and $C$ is nonzero positive semidefinite there. Hence active bonds exist. Each bond partition occurs as the partition of a two-forest by choosing spanning trees on both induced sides. Equations \eqref{T06-den}--\eqref{T06-num} express $\rho_C$ as a strictly positive weighted average of the cardinalities of all active bonds. Thus it lies strictly between distinct extreme cardinalities, and is constant if the extremes coincide.

For any chosen active bond $S\mid S^c$, assign weight $1$ to its internal edges and $\varepsilon$ to its crossing edges. Two-forests using no crossing edge are precisely a spanning tree of $G[S]$ together with one of $G[S^c]$. Their total weight is a positive constant, their demand weight is $q_C(S)>0$, and their cut cardinality is $|\delta(S)|$. Every other two-forest has weight $O(\varepsilon)$. Therefore $\rho_C(w(\varepsilon))\to|\delta(S)|$. Both endpoints are in the closure of the range. The positive orthant is connected and $\rho_C$ is continuous, so its image contains every number strictly between the endpoints. Finally, each connected side has at least its vertex count minus one internal edges, implying $|\delta(S)|\le m-(n-2)=\beta+1$.
\end{proof}
For completeness, a minimum-cardinality $s$--$t$ cut can be chosen with both sides connected: if a component on the $s$-side does not contain $s$, moving it to the other side strictly decreases the cut; apply the same argument to the other side. Thus the lower endpoint interpretation above is exact.

\subsection{Cycle defect and equality classification}\label{T06-sec-3}
For a connected graph $H$, write $\beta(H)=|E(H)|-|V(H)|+1$.
\begin{theorem}[Exact nonnegative defect]\label{T06-defect}
For all the stated demands and conductances,
\[
 (\beta+1)\mathcal E_C(G,w)-\mathbb E\mathcal E_C(T,w)
 =\frac1\tau\sum_{F\in\mathcal F_2}w(F)q_F
 \bigl[\beta(G[S_F])+\beta(G[S_F^c])\bigr].
\]
Equality holds if and only if the two induced sides of every active bond are trees. The equality condition depends on $G,C$, not on the positive conductances.
\end{theorem}
\begin{proof}
For each bond, edge counting gives
$\beta+1-d_F=\beta(G[S_F])+\beta(G[S_F^c])$.
Subtract \eqref{T06-num} from $(\beta+1)$ times \eqref{T06-den}. All summands are nonnegative; every two-forest has positive weight, and both induced sides are connected.
\end{proof}

\begin{theorem}[Two-terminal rigidity]\label{T06-theta}
Fix distinct $s,t$. The equality
\[
 \mathbb E R_T(s,t)=(\beta+1)R_G(s,t)
\]
holds for one (equivalently every) positive conductance assignment exactly when every simple cycle of $G$ contains both $s$ and $t$. Equivalently, either $G$ is a tree, or its cyclic part is the union of $q\ge2$ internally vertex-disjoint $s$--$t$ paths, with arbitrary trees attached to vertices of this union. In the latter case $\beta=q-1$.
\end{theorem}
\begin{proof}
If every cycle contains both terminals, neither side of an $s$--$t$ bond can contain a cycle, so Theorem \ref{T06-defect} gives equality.

Conversely suppose a simple cycle $D$ omits at least one terminal. Extend $D$ to a connected spanning unicyclic subgraph $H$: repeatedly attach a new vertex by one edge until all vertices are included. Since the two terminals are not both on its unique cycle, the path between their attachments in the tree-of-components of $H$ uses a bridge separating them. Removing that bridge gives a connected partition separating $s,t$, with $D$ on one side. The same induced sides in $G$ are connected, and one is cyclic. This yields a positive term in the defect.

For the structural description, decompose a connected graph into maximal two-connected blocks and bridges. Distinct nontrivial blocks share at most one vertex, and every cycle lies in one block. The cycle condition forces all cyclic blocks to contain both terminals, so there is exactly one, say $B$. All other blocks are bridges and hence form attached trees. Both $B-s$ and $B-t$ are connected and acyclic. If $N=|V(B)|$, $M=|E(B)|$, it follows that
\[
 M-\deg_B(s)=N-2=M-\deg_B(t).
\]
Put their common degree equal to $q$. Then $M=N-2+q$ and
$\sum_{v\ne s,t}\deg_B(v)=2M-2q=2(N-2)$.
Every vertex of a two-connected block has degree at least two, so all vertices other than $s,t$ have degree exactly two. Following each edge out of $s$ therefore leads along an internally disjoint path to $t$; no closed detour is possible in a two-connected block. These paths exhaust $B$. Conversely such a union has only cycles passing through both terminals.
\end{proof}
The block facts used here can also be seen directly by splitting at each cut vertex: a simple cycle cannot cross two resulting pieces, and pieces that contain a common cycle merge into the same maximal block.

\begin{corollary}[Total resistance and weighted pair demands]
Let $K(G)=\sum_{s<t}R_G(s,t)$. Then
$\mathbb E K(T)\le(\beta+1)K(G)$, with equality exactly when $G$ is a tree or a spanning cycle. More generally, for nonnegative pair weights $d_{st}$ not all zero, equality for $\sum d_{st}R(s,t)$ holds precisely when every cycle contains both endpoints of every pair of positive weight.
\end{corollary}
\begin{proof}
Sum the nonnegative pairwise defects. For all pairs, the condition is that every cycle contain every vertex. If a cycle exists, it is spanning. Any extra edge is a chord and creates a smaller cycle, a contradiction. The converse is immediate.
\end{proof}

\subsection{Examples, interpretation, and scope}\label{T06-sec-4}
\paragraph{Complete graphs.}
For $K_n$ and one terminal pair, connected bipartitions are all nontrivial partitions. The minimum and maximum sizes of cuts separating the terminals are $n-1$ and $\lfloor n^2/4\rfloor$. For $n\ge4$ the exact conductance range is therefore
\[
 n-1<\frac{\mathbb E R_T(s,t)}{R_G(s,t)}<\lfloor n^2/4\rfloor.
\]
For $n=2,3$ the ratio is constantly $1,2$, respectively. The upper endpoint for $n\ge4$ is approached by weak conductances across a balanced separating cut. For $n\ge5$ it improves strictly on $\beta+1=(n-1)(n-2)/2+1$.

\paragraph{Theta cores.}
If the $q$ terminal paths have series resistances $r_1,\ldots,r_q$, their parallel resistance is $R=(\sum_i r_i^{-1})^{-1}$. A spanning tree has exactly one intact terminal path. Its index has probability $R/r_i$: the tree weight contribution of choosing path $i$ intact is proportional to $\prod_{j\ne i}r_j$. Therefore $\mathbb E R_T=\sum_i (R/r_i)r_i=qR$, independently of conductances. Attached trees do not affect the terminal resistance.

\paragraph{What is and is not asserted.}
The theorem is a sharp finite-graph statement, not a claim that random trees are low-stretch in all metrics. The extremal constant may be large, and maximizing a bond is a distinct combinatorial optimization problem. The forest enumeration machinery is established in the literature. The contribution proposed here is the exact conductance range for arbitrary covariance, its cut-concentration construction, and the complete cycle-rank equality classification. A targeted search of primary sources \cite{ref045,ref046,ref047} and the supplied graph-design benchmark did not establish an identical result; this is not a comprehensive priority certification or a solution to a named open conjecture.





\clearpage
\section[T07 --- Exact Sign-Robustness Certificates for Resistance Reconstruction]{Exact Sign-Robustness Certificates for Resistance Reconstruction}\label{contrib-T07}
\markboth{Electrical Networks, Spectral Design, and Greedy Augmentation}{T07 --- Taisho}
\noindent\textsc{T07 \quad Taisho programme, source round 7}\par\medskip



\begin{quote}\small\textbf{Overview.} How much symmetric error can resistance data tolerate before direct inversion produces a negative conductance? We give an exact answer in a centered spectral norm, and more generally for anisotropic Loewner uncertainty. The inverse image of the uncertainty set is an explicit matrix order interval. Its support function yields simultaneous necessary and sufficient sign certificates, exact individual conductance intervals, and constructive adverse perturbations. For a uniform complete network the sharp relative radius is $1/(n-1)$; every sparse connected network has radius zero for unrestricted centered noise. The result concerns sign-validity of direct inversion, not statistical recovery under a prescribed sparse model. Classical inverse-M-matrix robustness results are distinguished from this symmetric fixed-kernel problem.\end{quote}
\subsection{Setup and the exact certificate}\label{T07-sec-1}
Put $P=I-\one\one^T/n$ and $\mathcal H=\one^\perp$. A centered matrix annihilates $\one$. Throughout, inverses and positive definiteness of centered matrices are understood on $\mathcal H$, with zero extension to constants. Let $A$ be centered, symmetric and positive definite on $\mathcal H$. Let $W$ have the same properties. For $0\le\delta<\lambda_{\min}(W^{-1/2}AW^{-1/2}|_{\mathcal H})$, consider the error set
\[
 \mathcal U_\delta=\{E=E^T:E\one=0,\ -\delta W\preceq E\preceq\delta W\}.
\]
For $W=P$, this is exactly $\norm{E}_2\le\delta$. Define
\[
 L_\delta=(A+\delta W)^{-1},\quad U_\delta=(A-\delta W)^{-1},
 \quad D_\delta=U_\delta-L_\delta,\quad M_\delta=(L_\delta+U_\delta)/2.
\]
\begin{theorem}[Exact inverse-entry intervals]\label{T07-entries}
For every $i,j$, the range of $[(A+E)^{-1}]_{ij}$ over $E\in\mathcal U_\delta$ is exactly
\[
 \left[(M_\delta)_{ij}-\frac12\sqrt{(D_\delta)_{ii}(D_\delta)_{jj}},\quad
 (M_\delta)_{ij}+\frac12\sqrt{(D_\delta)_{ii}(D_\delta)_{jj}}\right].
\]
Consequently every perturbed inverse is a connected nonnegative-conductance graph Laplacian if and only if
\begin{equation}\label{T07-criterion}
 (L_\delta+U_\delta)_{ij}+\sqrt{(D_\delta)_{ii}(D_\delta)_{jj}}\le0
 \qquad(i\ne j).
\end{equation}
If one inequality fails, an explicit perturbation in $\mathcal U_\delta$ produces a positive inverse off-diagonal entry, hence a negative reconstructed conductance. The entry intervals need not be jointly attainable at their endpoints.
\end{theorem}
The zero radius of sparse networks, proved below, explains why ordinary norm-wise stability of an inverse does not by itself ensure topological or sign validity.

\subsection{Order geometry and a constructive extremizer}\label{T07-sec-2}
\begin{lemma}[Order-interval inversion]\label{T07-order}
There is an exact set equality
\[
 \{(A+E)^{-1}:E\in\mathcal U_\delta\}
 =\{X=X^T:X\one=0,\ L_\delta\preceq X\preceq U_\delta\}.
\]
For $\delta>0$, every matrix on the right can uniquely be written
$X=L_\delta+D_\delta^{1/2}KD_\delta^{1/2}$, where $0\preceq K\preceq P$ on $\mathcal H$.
\end{lemma}
\begin{proof}
Inversion reverses positive-definite order: $B\preceq C$ is equivalent to $C^{-1}\preceq B^{-1}$, by congruence with $B^{-1/2}$ and inversion of eigenvalues. Apply this fact in both directions to
$A-\delta W\preceq A+E\preceq A+\delta W$.
The remaining parametrization follows by congruence with $D_\delta^{-1/2}$; $D_\delta$ is positive definite on $\mathcal H$ for $\delta>0$.
\end{proof}
\begin{lemma}[Support function]\label{T07-support}
For centered symmetric $H$, let $Z=D_\delta^{1/2}HD_\delta^{1/2}$. Then
\[
 \max_{E\in\mathcal U_\delta}\tr\bigl(H(A+E)^{-1}\bigr)
 =\tr(HL_\delta)+\tr(Z_+),
\]
where $Z_+$ is the positive spectral part. A maximizer is obtained by taking $K$ to be the orthogonal projection onto the positive eigenspace of $Z$, setting $X=L_\delta+D_\delta^{1/2}KD_\delta^{1/2}$, and then $E=X^{-1}-A$. The minimum uses the negative spectral part instead.
\end{lemma}
\begin{proof}
Diagonalize $Z$. Since $0\le K_{aa}\le1$ in this basis, $\tr(ZK)\le\sum_{\lambda_a(Z)>0}\lambda_a(Z)$, with equality for the stated projection. Lemma \ref{T07-order} certifies the resulting error.
\end{proof}
\begin{proof}[Proof of Theorem \ref{T07-entries}]
Set $u=D_\delta^{1/2}Pe_i$ and $v=D_\delta^{1/2}Pe_j$. The two possibly nonzero eigenvalues of $(uv^T+vu^T)/2$ are
$(u^Tv\pm\norm{u}\norm{v})/2$.
Thus the maximum of $u^TKv$ for $0\preceq K\preceq P$ is
$(u^Tv+\norm{u}\norm{v})/2$, and the minimum is the same expression with the minus sign. Add $(L_\delta)_{ij}$ and use $u^Tv=(D_\delta)_{ij}$. Continuity on the connected order interval fills the whole scalar interval. The case $\delta=0$ is immediate.

A centered symmetric matrix positive definite on $\mathcal H$ is a graph Laplacian exactly when its off-diagonal entries are nonpositive: assign conductance $-X_{ij}$, use zero row sums for the diagonal, and use the one-dimensional kernel for connectedness. This proves the certificate.
\end{proof}
For $n\ge3$, $i\ne j$, $\delta>0$, the maximizing $K$ is rank one. Explicitly, with $u,v$ as above, normalize
\[
 q=\frac{u/\norm{u}+v/\norm{v}}{\norm{u/\norm{u}+v/\norm{v}}},\qquad K=qq^T.
\]
The numerator cannot vanish because $Pe_i,Pe_j$ are independent and $D_\delta^{1/2}$ is invertible on $\mathcal H$. For a violating pair, this gives a checkable adverse matrix without solving a nonconvex optimization problem.

\begin{proposition}[Spectral-ball extremizers are reflections]\label{T07-reflection}
Suppose $W=P$, $\delta>0$, and use the rank-one $K=qq^T$ above. Put
\[
 v=(A+\delta P)D_\delta^{1/2}q,\qquad h=v/\norm v.
\]
Then the adverse perturbation has the particularly simple form
\[
 E=\delta(P-2hh^T),\qquad \norm E_2=\delta.
\]
\end{proposition}
\begin{proof}
Woodbury gives $E=\delta P-vv^T/d$ with
$d=1+q^TD_\delta^{1/2}(A+\delta P)D_\delta^{1/2}q$.
Here $A$ and $D_\delta=2\delta(A^2-\delta^2P)^{-1}$ commute. Consequently
\[
 \norm v^2=2\delta q^T(A+\delta P)(A-\delta P)^{-1}q
 =2\delta\bigl[1+2\delta q^T(A-\delta P)^{-1}q\bigr]=2\delta d.
\]
Substitution proves the formula; it has eigenvalues $-\delta$ on $h$ and $+\delta$ on its orthogonal complement in $\mathcal H$.
\end{proof}

\subsection{Sharp robustness radii}\label{T07-sec-3}
Assume now that $L=A^{-1}$ already is a connected graph Laplacian. Define the sign-validity radius
\[
 \begin{split}
r_*(A)=\sup\{r\ge0:\ &(A+E)^{-1}\text{ is a graph Laplacian}\\
 &\text{for every centered symmetric }E\text{ with }\norm E_2<r\}.
\end{split}
\]
The definition includes existence and positive definiteness of the inverse. Write $a=\lambda_{\min}(A|_{\mathcal H})$ and, for $0\le\delta<a$,
\[
 g_{ij}(\delta)=(M_\delta)_{ij}+\tfrac12\sqrt{(D_\delta)_{ii}(D_\delta)_{jj}},\quad W=P.
\]
\begin{corollary}[One-dimensional exact certificate]
For $\delta<a$, the closed spectral ball is sign-valid exactly when $\max_{i<j}g_{ij}(\delta)\le0$. Each $g_{ij}$ is continuous and nondecreasing. In particular
\[
 r_*(A)=\min\left\{a,\ \inf\{\delta\in[0,a):\max_{i<j}g_{ij}(\delta)>0\}\right\},
\]
where the infimum of the empty set is $+\infty$. Testing the inequalities gives an exact decision certificate and monotone bisection for the radius, with any desired certified arithmetic accuracy.
\end{corollary}
\begin{proof}
The certificate is Theorem \ref{T07-entries}. Each $g_{ij}$ is a maximum over nested compact error balls, giving monotonicity; its displayed formula gives continuity. At radius $a$, the perturbation $-a uu^T$ for a minimum eigenvector makes the covariance singular, bounding the radius by $a$.
\end{proof}
\begin{theorem}[Sparse fragility and dense interior]\label{T07-sparse}
For $n\ge3$, $r_*(A)>0$ if and only if every off-diagonal entry of $L=A^{-1}$ is strictly negative. Thus a connected network with any missing edge has radius zero under unrestricted centered spectral errors.
\end{theorem}
\begin{proof}
Strict negativity is preserved in a sufficiently small neighborhood by continuity and finiteness of the entries. Conversely expand on $\mathcal H$:
\[
 M_\delta=L+\delta^2L^3+O(\delta^4),\qquad
 D_\delta=2\delta L^2+O(\delta^3).
\]
Since $(L^2)_{ii}>0$ for every vertex, Theorem \ref{T07-entries} gives
\[
 g_{ij}(\delta)=L_{ij}+\delta\sqrt{(L^2)_{ii}(L^2)_{jj}}+O(\delta^2).
\]
If $L_{ij}=0$, this is positive for every sufficiently small positive $\delta$.
\end{proof}
\begin{theorem}[Uniform complete network]
If $A=\alpha P$, $\alpha>0$, then
\[
 r_*(A)=\frac{\alpha}{n-1}.
\]
For $n\ge3$, the inverse remains positive definite throughout the larger range $\delta<\alpha$, so loss of graph signs occurs strictly before loss of definiteness.
\end{theorem}
\begin{proof}
Direct substitution yields, for $i\ne j$,
\[
 g_{ij}(\delta)=\frac{-\alpha+(n-1)\delta}{n(\alpha^2-\delta^2)}.
\]
This is nonpositive precisely when $\delta\le\alpha/(n-1)$. For $n=2$ the signs persist up to the singularity threshold $\alpha$. For $n\ge3$ and any selected pair, the extremizer in Proposition \ref{T07-reflection} uses
$h=(Pe_i+Pe_j)/\norm{Pe_i+Pe_j}$ and witnesses sharpness.
\end{proof}
The reciprocal-Laplacian condition can therefore fail even when the centered data matrix is very well conditioned: at the complete-network threshold its spectral condition number is only $n/(n-2)$.

\subsection{Resistance data and limitations}\label{T07-sec-4}
For a resistance matrix $R$, with $R_{ij}=(e_i-e_j)^TL^+(e_i-e_j)$, double centering gives
\[
 A=L^+=-\tfrac12 PRP.
\]
Conversely every centered symmetric $A$ corresponds to a unique symmetric hollow matrix
$R(A)_{ij}=A_{ii}+A_{jj}-2A_{ij}$. These formulas are inverse linear maps, as is checked using $A\one=0$ and $R_{ii}=0$. Thus the exact noise model on hollow symmetric data is
\[
 \tfrac12\norm{P\Delta R P}_2\le\delta.
\]
Every admissible centered $E$ occurs from such a $\Delta R$, namely $\Delta R=R(E)$. If $A+E$ is positive definite on $\mathcal H$, its associated hollow matrix is a squared Euclidean distance matrix; the sign certificate decides the additional, stronger property of being a resistance matrix of a nonnegative network. This distinguishes metric geometry from electrical realizability.

The familiar centering reconstruction is classical; see \cite{ref049}. Hlad\'ik \cite{ref048}, Theorem 24 and Corollary 4, gives exact \emph{unstructured nonsingular} inverse-M-matrix radii using singular values of cofactors. Here errors must be symmetric and preserve a prescribed nullspace, and we compute complete entry intervals over the error ball, with explicit extremizers. The unstructured cofactor formula cannot simply be applied to the singular full matrix $A$. No claim is made that the elementary order-interval support-function lemma is new. The intended refinement is its exact fixed-kernel reconstruction certificate and sharp graph consequences; comprehensive priority remains unverified.

This result does not cover independent entrywise error bounds, missing resistance measurements, or projection onto a prescribed graph support. Such structured models can have different and nonzero sparse-recovery radii. It also does not claim simultaneous attainability of different worst entries.




\clearpage
\section[T08 --- Universal Submodularity of Rank-One Inverse Gains:\ A Commutation Classification and a Three-Vertex Obstruction]{Universal Submodularity of Rank-One Inverse Gains:\ A Commutation Classification and a Three-Vertex Obstruction}\label{contrib-T08}
\markboth{Electrical Networks, Spectral Design, and Greedy Augmentation}{T08 --- Taisho}
\noindent\textsc{T08 \quad Taisho programme, source round 8}\par\medskip



\begin{quote}\small\textbf{Overview.} Resistance reduction and A-optimal information gain are not generally submodular. We classify exactly when submodularity holds uniformly over all positive semidefinite demand matrices: the whitened rank-one update directions must lie in mutually orthogonal one-dimensional groups, equivalently their rank-one matrices must commute. The necessity is constructive and uses a rank-one demand that gives strict complementarity at every pair of positive update strengths. For a fixed demand and two directions, we give a sharp scalar criterion for diminishing returns at all strengths. A three-vertex electrical example shows that the submodularity ratio can tend to zero on a fixed topology with a single terminal-pair demand. General non-submodularity and vanishing-ratio phenomena are known; the purpose here is the exact universal classification and a small weighted-network obstruction, not a new claim that these phenomena exist.\end{quote}
\subsection{Universal submodularity: a complete classification}\label{T08-sec-1}
Let $M\succ0$ be a real $d\times d$ matrix, let $u_1,\ldots,u_k\ne0$ be update directions, and fix strengths $x_1,\ldots,x_k>0$. For $C\succeq0$ define
\[
 G_C(S)=\tr(CM^{-1})-\tr\left(C\left(M+\sum_{i\in S}x_i u_i u_i^T\right)^{-1}\right),\qquad S\subseteq[k].
\]
This is a nonnegative monotone gain function. Submodularity means
$G_C(S\cup\{i\})-G_C(S)\ge G_C(T\cup\{i\})-G_C(T)$ whenever $S\subseteq T$ and $i\notin T$.
Put $p_i=M^{-1/2}u_i$.
\begin{theorem}[Universal demand classification]\label{T08-main}
The following conditions are equivalent.
\begin{enumerate}
\item For the fixed positive strengths, $G_C$ is submodular for every $C\succeq0$.
\item For every choice of positive strengths and every $C\succeq0$, $G_C$ is submodular.
\item Each pair $p_i,p_j$ is either parallel or orthogonal.
\item The rank-one matrices $p_i p_i^T$ commute pairwise.
\end{enumerate}
If these conditions fail, a rank-one $C$ and a pair of indices violate the submodular inequality already at the empty set, strictly for every positive strength on that pair. In particular testing rank-one demands is enough for the universal classification.
\end{theorem}
In an electrical network, work on $\one^\perp$, take $M=L$ for a connected baseline Laplacian, and let $u_i$ be edge incidence vectors. Then $u_i^TL^+u_j$ is the electrical inner product. For distinct simple edges, the directions are never parallel, so the universal property is equivalent to pairwise electrical orthogonality. There can be at most $n-1$ such distinct candidates. Reinforcing the edges of a baseline tree supplies an example of an electrically orthogonal family.

\subsection{An exact two-update identity}\label{T08-sec-2}
Fix nonzero $u,v$, set
\begin{align*}
 r&=u^TM^{-1}u,&s&=v^TM^{-1}v,&t&=u^TM^{-1}v,\\
 a&=u^TM^{-1}CM^{-1}u,&b&=v^TM^{-1}CM^{-1}v,&c&=u^TM^{-1}CM^{-1}v.
\end{align*}
For $x,y>0$, let $G(x,y)$ denote the gain from $xuu^T+yvv^T$, and set
$D=(1+xr)(1+ys)-xy t^2>0$.
\begin{lemma}[Exact complementarity formula]\label{T08-two}
\begin{align}
 G(x,y)&=\frac{xa(1+ys)+yb(1+xr)-2xy tc}{D},\label{T08-wood}\\
 G(x,y)-G(x,0)-G(0,y)
 &=\frac{xy t^2}{D}\left(\frac{xa}{1+xr}+\frac{yb}{1+ys}-\frac{2c}{t}\right)\quad(t\ne0).\label{T08-interaction}
\end{align}
If $t=0$, the interaction is identically zero, for every $C$ and every $x,y$.
\end{lemma}
\begin{proof}
The Woodbury identity, verified by multiplication, is
\[
 (M+UD_xU^T)^{-1}=M^{-1}-M^{-1}U(D_x^{-1}+U^TM^{-1}U)^{-1}U^TM^{-1},
\]
with $U=[u\ v]$ and $D_x=\diag(x,y)$. Inverting its two-by-two middle matrix and taking the trace against $C$ gives \eqref{T08-wood}. For one update the gain is $xa/(1+xr)$ or $yb/(1+ys)$. Subtracting these two expressions from \eqref{T08-wood}, with a common denominator, gives \eqref{T08-interaction}. The formula without division by $t$ is
\[
 \frac{xy}{D}\left[t^2\left(\frac{xa}{1+xr}+\frac{yb}{1+ys}\right)-2tc\right],
\]
which also proves the assertion for $t=0$.
\end{proof}
\begin{theorem}[Fixed-demand all-strength criterion]\label{T08-fixed}
If $t\ne0$, the two-element gain function is submodular for every $x,y>0$ if and only if
\[
 \frac{2c}{t}\ge\frac ar+\frac bs.
\]
If this fails, sufficiently large finite $x,y$ violate submodularity. If $tc<0$, every positive $x,y$ violates it. If $t=0$, the gain is modular.
\end{theorem}
\begin{proof}
The two summands in parentheses in \eqref{T08-interaction} increase from zero to $a/r$ and $b/s$. Their supremum is $a/r+b/s$. The prefactor is positive. If $tc<0$, then $c/t<0$, so the parenthesis is positive even when one or both of $a,b$ vanish; in fact $tc<0$ forces both positive by positive semidefiniteness.
\end{proof}
This criterion is for a fixed background $M$. It cannot establish global submodularity for a larger family merely by testing all pairs at the empty set: the background changes after other updates. The universal theorem avoids this pitfall by identifying a structure preserved under all updates.

\subsection{Proof of the universal classification}\label{T08-sec-3}
\begin{proof}[Proof of Theorem \ref{T08-main}]
Conditions 3 and 4 are equivalent because
$(pp^T)(qq^T)=(p^Tq)pq^T$: for $p^Tq\ne0$, commutation forces $pq^T=qp^T$, hence parallelism; orthogonal or parallel pairs plainly commute.

Assume condition 3. Partition the $p_i$ into parallel classes and choose a unit vector $z_\ell$ on each class. Different $z_\ell$ are orthogonal. Write $p_i=\alpha_i z_{\ell(i)}$ and let $K=M^{-1/2}CM^{-1/2}\succeq0$. Whitening and inverting in these orthogonal directions gives
\[
 G_C(S)=\sum_\ell (z_\ell^TKz_\ell)\frac{q_\ell(S)}{1+q_\ell(S)},\qquad
 q_\ell(S)=\sum_{i\in S:\ell(i)=\ell}x_i\alpha_i^2.
\]
Every coefficient is nonnegative and $q/(1+q)$ is increasing and concave on $[0,\infty)$. Directly, its marginal for an addition $h\ge0$ is
$h/[(1+q)(1+q+h)]$, which decreases with $q$. Thus condition 2 holds and implies condition 1.

For necessity, suppose $p_i,p_j$ are independent and not orthogonal. There exists $z\in\R^d$ with
\[
 p_i^Tz=1,\qquad p_j^Tz=-\operatorname{sgn}(p_i^Tp_j),
\]
because two independent linear equations are solvable. Choose
$C=M^{1/2}zz^TM^{1/2}\succeq0$. Then $a=b=1$ and
$c=-\operatorname{sgn}(t)$, so $tc<0$. Lemma \ref{T08-two} gives a strict violation at $S=\varnothing$ for every positive pair of strengths. Therefore condition 1 fails.
\end{proof}
\begin{corollary}[Exact modularity]
If no two whitened directions are parallel, universal submodularity is equivalent to universal modularity, and to pairwise electrical orthogonality. With parallel classes, all interactions between distinct classes are zero, and within each class the scalar gain is concave.
\end{corollary}
To verify the tree example stated earlier, send unit current between the endpoints of an edge of a tree. Only that edge carries current, so the voltage drop across any different edge is zero. This says $u_i^TL^+u_j=0$ for distinct tree edges. By contrast, shared vertices alone do not determine the electrical inner product on a cyclic graph.

\subsection{A fixed three-vertex obstruction with vanishing ratio}\label{T08-sec-4}
The classification allows arbitrary positive semidefinite demands, which are more general than nonnegative sums of pair demands. The next obstruction uses only one ordinary terminal pair.

Consider vertices $1,2,3$ with baseline conductances $w_{12}=w_{23}=1$, $w_{13}=M>0$. Candidate $e$ increases $w_{12}$ by $x>0$ and candidate $f$ increases $w_{23}$ by the same $x$. Let $G(S)=R_{13}(\varnothing)-R_{13}(S)$.
\begin{theorem}[Three-vertex unbounded complementarity]\label{T08-triangle}
The ratio of the sum of individual gains to the joint gain is
\[
 \frac{G(\{e\})+G(\{f\})}{G(\{e,f\})}
 =\frac{2M+x+1}{M(x+2)+x+1}.
\]
It is strictly less than one for every $M,x>0$. With $x=M^2$, it equals
\[
 \frac{(M+1)^2}{M^3+M^2+2M+1}\longrightarrow0.
\]
Thus no positive lower bound on the submodularity ratio can hold uniformly over positive weights and reinforcement strengths, even on this fixed three-vertex topology for one terminal pair.
\end{theorem}
\begin{proof}
The edge $13$ is in parallel with the series path $1$--$2$--$3$. Therefore
\[
 R_0=\frac2{2M+1},\quad R_e=R_f=\frac{x+2}{M(x+2)+x+1},\quad R_{ef}=\frac2{2M+x+1}.
\]
Subtracting gives
\[
 G(\{e\})=\frac{x}{(2M+1)[M(x+2)+x+1]},\qquad
 G(\{e,f\})=\frac{2x}{(2M+1)(2M+x+1)}.
\]
Their ratio is as stated. Its denominator minus numerator is $Mx>0$. Substitution proves the limit. Under the usual ratio defined as an infimum over backgrounds and nonempty groups of additions, this empty-background two-element ratio is an upper bound, so the last conclusion follows.
\end{proof}
For $M=x=1$, the three resistances are $2/3$, $3/5$, $1/2$. The individual gains sum to $2/15$, whereas the joint gain is $1/6$, a strict gap of $1/30$.

\subsection{Prior work, interpretation, and limits}\label{T08-sec-5}
Non-submodularity of A-optimal objectives and effective-resistance reduction is established; see \cite{ref051,ref052}. Achterberg--Kooij \cite{ref052} also construct graph families with submodularity ratios tending to zero for an all-pairs link-addition objective. Our three-vertex construction changes both the admissible action (reinforcement) and objective (one pair). It is not a replacement proof of their theorem under unchanged assumptions. Likewise a vanishing submodularity ratio removes a particular generic guarantee; it does \emph{not} by itself prove that greedy solutions have arbitrarily poor approximation ratios.

The main structural result quantifies over all positive semidefinite $C$. A fixed $C$, including the usual all-pairs covariance, may exhibit submodularity for a much wider class of updates. The exact two-update criterion separates this fixed-demand question from the universal one. The rank-one commutation criterion and its constructive necessity are proved independently here, but exhaustive priority is not certified. No named open conjecture is claimed solved.





\clearpage
\section[T09 --- Higher Spectral Volumes of Random Trees:\ Sharp Quotient-Partition Ranges]{Higher Spectral Volumes of Random Trees:\ Sharp Quotient-Partition Ranges}\label{contrib-T09}
\markboth{Electrical Networks, Spectral Design, and Greedy Augmentation}{T09 --- Taisho}
\noindent\textsc{T09 \quad Taisho programme, source round 9}\par\medskip



\begin{quote}\small\textbf{Overview.} We determine the exact conductance-extremal range of every inverse spectral-volume coefficient under weighted random spanning-tree sampling. With an arbitrary positive semidefinite demand matrix, the relevant combinatorial objects are connected vertex partitions whose block-constant demand compression has positive determinant. The sharp inflation endpoints are the minimum and maximum unweighted spanning-tree counts of the corresponding quotient multigraphs. This yields all-order cycle-rank bounds, their strictness and equality structure, and explicit sharp endpoints for complete graphs in terms of balanced integer partitions. The matrix-forest and volume-sampling mechanisms are classical; the topology-sensitive exact range and complete-graph extremizers are the focus. This is a higher-exterior-volume extension of the two-forest energy identity, not a claim that the classical forest theorem is new.\end{quote}
\subsection{Spectral volumes and quotient partitions}\label{T09-sec-1}
Let $G=(V,E)$ be finite, connected and simple, $n\ge2$, with positive conductances $w_e$. Write $\beta=|E|-n+1$, $L=\sum_e w_e b_e b_e^T$, $\mathcal H=\one^\perp$, and $P=I-\one\one^T/n$. Let $\tau_w(G)$ denote the weighted spanning-tree sum. A random tree has law $\mathbb P(T)=w(T)/\tau_w(G)$ and inherits the original conductances.

Fix $C\succeq0$, $C\one=0$, and $1\le k\le\rank C$. Define
\[
 V_k(L,C)=[t^k]\det_{\mathcal H}(I+tL^{-1/2}CL^{-1/2}).
\]
Here inverse powers act on $\mathcal H$. This is positive. For $C=P$, it is $e_k(L^+)$, the $k$th elementary symmetric function of the nonzero inverse eigenvalues.

A connected $(k+1)$-partition $\pi=\{S_1,\ldots,S_{k+1}\}$ consists of nonempty blocks inducing connected subgraphs. Let $G/\pi$ be the loopless multigraph formed by contracting each block; keep the multiplicity of crossing edges, but give each edge unit weight. Write $\tau_0(G/\pi)$ for its unweighted tree count. Put
\[
 \mathcal K_\pi=\{x\in\mathcal H:x\text{ is constant on each }S_i\},\quad
 q_C(\pi)=\left(\prod_i|S_i|\right)\det(Z_\pi^TCZ_\pi),
\]
where $Z_\pi$ is any orthonormal basis matrix for the $k$-dimensional space $\mathcal K_\pi$. This definition is independent of that basis. Call $\pi$ active if $q_C(\pi)>0$.

\begin{theorem}[Exact higher-volume range]\label{T09-range}
Active connected partitions exist. Define
\[
 a_k=\min_{\pi\text{ active}}\tau_0(G/\pi),\qquad
 b_k=\max_{\pi\text{ active}}\tau_0(G/\pi).
\]
As the positive conductances vary, the ratio
\[
 \frac{\mathbb E V_k(L_T,C)}{V_k(L,C)}
\]
has range exactly $(a_k,b_k)$ if $a_k<b_k$, and is identically $a_k$ otherwise. Both sharp endpoints are approached by weakening all edges crossing an appropriate partition while fixing its internal conductances.
\end{theorem}
For $k=1$, the quotient has two vertices, and its tree count is the cut cardinality. For $k>1$, the tree count of the entire quotient replaces a single cut size. It records how the components can be connected, not merely how many crossing edges there are.

\subsection{A demand-compressed matrix-forest expansion}\label{T09-sec-2}
Let $\mathcal F_{k+1}$ be the spanning forests with exactly $k+1$ components, and let $\pi(F)$ be their vertex partition.
\begin{lemma}[Forest coefficient formula]\label{T09-forest}
\[
 n\tau_w(G)V_k(L,C)=\sum_{F\in\mathcal F_{k+1}}w(F)q_C(\pi(F)).
\]
\end{lemma}
\begin{proof}
Choose orthonormal coordinates on $\mathcal H$. The determinant factorization gives
$\det_{\mathcal H}(L+tC)=\det_{\mathcal H}(L)\sum_j V_j(L,C)t^j$.
The weighted matrix-tree identity gives $\det_{\mathcal H}(L)=n\tau_w(G)$. One proof uses Cauchy--Binet on a reduced incidence matrix: maximal nonzero incidence minors are exactly tree minors and have absolute value one.

Now apply Cauchy--Binet to the Gram matrix of the column family
$[\,\sqrt{w_e}b_e\ (e\in E)\mid\sqrt t\,C^{1/2}\,]$ in $\mathcal H$.
A term of degree $k$ in $t$ uses $n-1-k$ edge columns. They are independent exactly when their edge set is a forest $F\in\mathcal F_{k+1}$; cyclic sets contribute zero. Their span has orthogonal complement $\mathcal K_{\pi(F)}$ in $\mathcal H$.

By orthogonally triangularizing these edge columns, the sum over all choices of the remaining $k$ columns is
\[
 w(F)\det(B_F^TB_F)\det(Z_{\pi(F)}^TCZ_{\pi(F)}).
\]
Indeed the squared volume factors into the edge-column Gram determinant and the squared projected volume of the remaining columns, whose sum is the last determinant by Cauchy--Binet. Finally $\det(B_F^TB_F)=\prod_i|S_i|$: for a tree on $s$ vertices its incidence Gram determinant is $s$, either by Cauchy--Binet over the $s$ possible deleted rows, each minor being $\pm1$, or by leaf elimination. Different components give separate blocks. This proves the formula, including singleton components through the empty determinant convention.
\end{proof}
For $C=P$ this is the usual rooted-forest coefficient formula. The compression determinant is the additional demand factor.

\begin{theorem}[Exact random-tree identity]\label{T09-expectation}
\[
 n\tau_w(G)\,\mathbb E V_k(L_T,C)
 =\sum_{F\in\mathcal F_{k+1}}w(F)q_C(\pi(F))\tau_0(G/\pi(F)).
\]
\end{theorem}
\begin{proof}
Apply Lemma \ref{T09-forest} to each tree $T$, for which $\tau_w(T)=w(T)$. Multiplication by its sampling probability cancels $w(T)$ in the denominator. A fixed forest $F$ occurs in exactly $\tau_0(G/\pi(F))$ spanning trees: adding the remaining $k$ edges is precisely the choice of a spanning tree of the quotient multigraph. Notice that this multiplicity is \emph{unweighted}; the sampling weight cancellation is essential.
\end{proof}
\begin{proof}[Proof of Theorem \ref{T09-range}]
The two identities express the ratio as a positive weighted average over all active connected partitions. Positivity of $V_k$ guarantees that at least one is active. Every connected partition occurs by choosing a tree inside each block.

Fix an active partition $\pi$ and set all internal conductances to one and all crossing conductances to $\varepsilon$. The forests in $\mathcal F_{k+1}$ using no crossing edge are exactly the unions of one spanning tree in every block of $\pi$. Their total coefficient is positive. Every other forest has weight $O(\varepsilon)$. The ratio therefore tends to $\tau_0(G/\pi)$. The function is continuous on the connected positive orthant; its image is an interval and contains every value strictly between its infimum and supremum. No endpoint can occur when active quotient counts differ, because every corresponding forest has positive weight.
\end{proof}

\subsection{Cycle-rank bounds and their equality cases}\label{T09-sec-3}
\begin{theorem}[All-order universal bound]\label{T09-bound}
For all the stated demands,
\[
 \mathbb E V_k(L_T,C)\le \binom{\beta+k}{k}V_k(L,C).
\]
More precisely, the exact defect is the nonnegative forest sum
\[
 \frac1{n\tau_w(G)}\sum_{F\in\mathcal F_{k+1}}w(F)q_C(\pi(F))
 \left[\binom{\beta+k}{k}-\tau_0(G/\pi(F))\right].
\]
If $k\ge2$ and $\beta\ge2$, this defect is strictly positive for every positive conductance assignment and every demand of rank at least $k$.
\end{theorem}
\begin{proof}
A connected partition has at least $\sum_i(|S_i|-1)=n-k-1$ internal edges, so its quotient has $\ell\le\beta+k$ edges. A quotient tree chooses $k$ of them; hence its tree count is at most $\binom{\ell}{k}\le\binom{\beta+k}{k}$. The defect formula follows by subtraction.

If equality held for one quotient with $k\ge2$, $\beta\ge2$, then $\ell=\beta+k\ge k+2$ and every $k$-edge subset would be a spanning tree. Thus the quotient would have no cycle of length at most $k$, including parallel-edge cycles. It is therefore simple and has girth at least $k+1$, on exactly $k+1$ vertices. It contains a cycle because $\ell>k$; such a cycle must be spanning, and any additional edge is a chord creating a shorter cycle. Hence $\ell\le k+1$, a contradiction. Every active summand is consequently strictly positive.
\end{proof}
\begin{corollary}[Isotropic equality, all higher orders]\label{T09-eq}
Let $C=P$ and $2\le k\le n-1$. Equality in Theorem \ref{T09-bound} holds if and only if $G$ is a tree or a cycle. On a cycle, for every positive weighting,
\[
 \mathbb E e_k(L_T^+)=(k+1)e_k(L^+),\qquad 0\le k\le n-1.
\]
\end{corollary}
\begin{proof}
All connected partitions are active. Trees give quotient trees. If $G$ is a cycle, each connected proper block is an arc and every quotient is a cycle on $k+1$ vertices (two parallel edges when $k=1$), with $k+1$ spanning trees. The case $k=0$ is trivial.

Theorem \ref{T09-bound} excludes $\beta\ge2$. If $\beta=1$ and $G$ is not a cycle, it has a bridge. Choose a spanning tree and remove this bridge together with any other $k-1$ edges. The resulting partition has a quotient containing a bridge. Its tree count is strictly less than $k+1$: it has at most $k+1$ edges and is either a tree, or is unicyclic with cycle length at most $k$. This partition contributes a strictly positive defect.
\end{proof}
For $k=1$, the same isotropic equality classification follows directly: equality requires both sides of every connected bipartition to be acyclic. Any nonspanning cycle extends to a spanning unicyclic graph with a bridge separating a vertex outside the cycle, contradicting that requirement. Thus every cycle must be spanning, and an extra chord is impossible.

\subsection{Explicit sharp complete-graph ranges}\label{T09-sec-4}
\begin{theorem}[Complete-graph formula]\label{T09-complete}
Let $G=K_n$, $C=P$, and $1\le k\le n-1$. Set $r=k+1$, $q=\lfloor n/r\rfloor$, and $h=n-rq$. The sharp lower and upper inflation endpoints are
\[
 A_{n,k}=n^{k-1}(n-k),\qquad
 B_{n,k}=n^{k-1}q^{\,r-h}(q+1)^h.
\]
The ratio takes every value strictly between them when they differ, and is constant when they coincide. The lower endpoint is approached by a partition of sizes $(n-k,1,\ldots,1)$; the upper endpoint by block sizes as equal as possible.
\end{theorem}
\begin{proof}
For block sizes $s_1,\ldots,s_r$, the quotient has $s_i s_j$ parallel edges between vertices $i,j$. Its Laplacian is $n\diag(s)-ss^T$. Delete the last row and column and apply the determinant lemma:
\[
 \tau_0(K_n/\pi)
 =n^{r-1}\prod_{i<r}s_i\left(1-\frac1n\sum_{i<r}s_i\right)
 =n^{r-2}\prod_{i=1}^r s_i.
\]
With positive integer $s_i$ summing to $n$, their product is minimized by $(n-r+1,1,\ldots,1)$: replacing $a,b\ge2$ by $a+b-1,1$ decreases it. It is maximized by sizes differing by at most one: replacing $a,b$ with $a-1,b+1$ increases the product when $a\ge b+2$. Theorem \ref{T09-range} now gives the entire range and the stated degenerations.
\end{proof}
For example, with $n=6$, the endpoints for $k=1,2,3,4,5$ are respectively
\[
 (5,9),\quad(24,48),\quad(108,144),\quad(432,432),\quad(1296,1296).
\]
In particular some high-order ratios are identically constant even though lower-order ratios vary substantially. At $k=n-1$, every connected graph has the exact conductance-independent identity
\[
 \frac{\mathbb E\det_{\mathcal H}(L_T^{-1/2}CL_T^{-1/2})}
 {\det_{\mathcal H}(L^{-1/2}CL^{-1/2})}=\tau_0(G)
\]
when $C$ is positive definite on $\mathcal H$. There is only one partition into $n$ singleton blocks.

\subsection{Polynomial consequences and prior work}\label{T09-sec-5}
Coefficientwise, for $t\ge0$,
\[
 \mathbb E\det_{\mathcal H}(I+tL_T^{-1/2}CL_T^{-1/2})
 \le\sum_{k=0}^{\rank C}\binom{\beta+k}{k}V_k(L,C)t^k.
\]
This is an expectation bound, not a claim that one tree simultaneously realizes each coefficientwise minimum. On a cycle with $C=P$, writing $p(t)=\det_{\mathcal H}(I+tL^+)$, the exact expected polynomial is $p(t)+tp'(t)$.

The classical matrix-forest theorem is developed in \cite{ref053}; inverse expectation identities in volume sampling are developed in \cite{ref054}. Graphic incidence configurations are not generally in full general position, so equalities for full-spark volume sampling must not be imported without checking their hypotheses. Our forest multiplicities retain the exact topology-dependent defect. The $k=1$ mechanism was developed in Taisho round 6; all proofs needed here are included. The genuine higher-order ingredients are the demand-compression determinant, quotient-tree multiplicities, strict higher-rank graphic obstruction, and the all-order complete-graph range. No claim is made that the binomial bound alone is new, or that a named open problem has been solved.




\clearpage
\section[T10 --- Arbitrarily Bad Greedy Link Addition on Unweighted Networks\ with Three Candidate Edges]{Arbitrarily Bad Greedy Link Addition on Unweighted Networks\ with Three Candidate Edges}\label{contrib-T10}
\markboth{Electrical Networks, Spectral Design, and Greedy Augmentation}{T10 --- Taisho}
\noindent\textsc{T10 \quad Taisho programme, source round 10}\par\medskip



\begin{quote}\small\textbf{Overview.} We construct unweighted simple connected graphs with only three specified candidate nonedges and a budget of two additions for which the exact greedy algorithm is arbitrarily bad for the ordinary all-pairs Kirchhoff index. For an integer parameter $M\ge10$, the greedy resistance-reduction gain is at most $9/M$ times optimal, while its final Kirchhoff index is at least $M/10$ times optimal. Every original and added edge has unit conductance, every action has equal cost, and the first greedy choice is strict. The construction has $\Theta(M^4)$ vertices and uses a unicyclic core plus pendant leaves. The proof separates a small electrical complementarity mechanism from an exact terminal-demand amplification identity. This is a candidate-constrained theorem: the allowed set contains three nonedges, not every missing edge. Accordingly it does not resolve the unrestricted candidate-set question posed in the existing literature.\end{quote}
\subsection{Problem and result}\label{T10-sec-1}
For a connected unweighted graph $G$, let
\[
 K(G)=\sum_{\{x,y\}\subseteq V(G)}R_G(x,y)
\]
be its Kirchhoff index, where every edge has resistance one. Given a specified set $\mathcal A$ of absent edges, write $G+S$ for adding $S\subseteq\mathcal A$, and let
\[
 F(S)=K(G)-K(G+S).
\]
With budget two, exact greedy first chooses an edge maximizing $F(\{e\})$, then an unused edge maximizing its marginal gain. An optimal set $S_*\subseteq\mathcal A$, $|S_*|=2$, maximizes $F$, equivalently minimizes the final Kirchhoff index.

\begin{theorem}[Unbounded failure with three candidate edges]\label{T10-main}
For each integer $M\ge10$ there is an unweighted simple connected graph $G_M$ on
\[
 n_M=(3M^2+M)(1+200M^2)
\]
vertices, and a set of exactly three absent edges, such that every exact greedy output $S_g$ with budget two satisfies
\[
 \frac{F(S_g)}{F(S_*)}\le\frac9M,
 \qquad
 \frac{K(G_M+S_g)}{K(G_M+S_*)}\ge\frac M{10}.
\]
The first greedy choice is unique. Thus greedy has no positive constant guarantee for gain, and no finite constant multiplicative guarantee for final Kirchhoff index, on candidate-constrained unweighted link-addition instances. Along this family the bounds have orders $O(n_M^{-1/4})$ and $\Omega(n_M^{1/4})$, respectively.
\end{theorem}
These are actual algorithmic performance bounds, not merely a failure of a submodularity-ratio certificate. The distinction is important: a vanishing submodularity ratio alone does not prove poor greedy performance.

\subsection{The unweighted construction}\label{T10-sec-2}
The core $H_M$ has four distinguished vertices $s,u,t,v$. Form three internally vertex-disjoint paths
\[
 s\longleftrightarrow u\quad(M^2\text{ edges}),\qquad
 u\longleftrightarrow t\quad(M^2\text{ edges}),\qquad
 t\longleftrightarrow s\quad(M\text{ edges}).
\]
Their union is a cycle. Attach to $s$ a further path of $M^2$ edges ending at $v$, with all its other vertices new. Thus the core has
\[
 q=3M^2+M
\]
vertices. The three allowed additions are
\[
 e=\{s,u\},\qquad f=\{u,t\},\qquad h=\{v,t\}.
\]
They are distinct nonedges. All three added edges have unit conductance and cost one. No other addition is allowed.

Finally attach $N=100qM^2$ new pendant vertices to $s$ and another $N$ new pendant vertices to $t$, by unit edges. This gives $G_M$ and the stated vertex count $q+2N=q(1+200M^2)$. Pendant vertices are not endpoints of candidate edges. The original graph is unicyclic; its maximum degree is large only because of the two leaf collections.

\subsection{The terminal mechanism}\label{T10-sec-3}
Let $R_S$ be the effective resistance from $s$ to $t$ inside $H_M+S$, and put $g(S)=R_\varnothing-R_S$. Pendant leaves do not change these terminal resistances.

\begin{lemma}[Core estimates]\label{T10-core}
For every integer $M\ge10$,
\begin{align}
 R_\varnothing&=\frac{2M^2}{2M+1}>M-\tfrac12,\label{T10-base}\\
 g(\{e\})=g(\{f\})&\le\tfrac12,\label{T10-single}\\
 g(\{h\})&>\tfrac45,\label{T10-trap}\\
 g(\{e,f\})&\ge M/4,\label{T10-good}\\
 g(\{h,e\}),\ g(\{h,f\})&\le2.\label{T10-bad}
\end{align}
Consequently the trap edge $h$ has a strictly larger individual terminal gain than either good edge, but any pair containing $h$ has bounded total terminal gain, while the two good edges have gain growing linearly with $M$.
\end{lemma}
\begin{proof}
All internal vertices of the four defining paths have degree two, except after the stated candidate additions at their endpoints. A path of length $\ell$ has effective conductance $1/\ell$. Suppressing these internal vertices for the terminal calculation gives baseline triangle conductances
\[
 w_{su}=w_{ut}=M^{-2},\qquad w_{st}=M^{-1},
\]
with a pendant path of resistance $M^2$ from $s$ to $v$. Adding $e$ or $f$ increases the corresponding triangle conductance by one. Adding $h$ creates an additional $s$--$t$ branch with series resistance $M^2+1$, so it increases $w_{st}$ by $1/(M^2+1)$.

Multiply all reduced conductances by $M^2$. The resulting triangle has baseline weights $1,1,M$, side increments $x=M^2$, and direct-edge increment
\[
 \delta=\frac{M^2}{M^2+1}\le1
\]
when $h$ is selected. Resistances in the original core are $M^2$ times those in this scaled triangle.

Parallel and series laws now give \eqref{T10-base} and
\[
 g(\{e\})=g(\{f\})
 =\frac{M^4}{(2M+1)(M^3+M^2+2M+1)}\le\frac12.
\]
For the trap,
\[
 g(\{h\})=\frac{4M^2\delta}{(2M+1)(2M+1+2\delta)}.
\]
For $M\ge10$, $\delta\ge100/101$, $2M+1\le(21/10)M$, and $2M+1+2\delta\le(23/10)M$. Therefore
\[
 g(\{h\})\ge\frac{40000}{48783}>\frac45.
\]
Both good edges give
\[
 R_{\{e,f\}}=\frac{2M^2}{(M+1)^2}<2,\qquad
 g(\{e,f\})=\frac{2M^4}{(2M+1)(M+1)^2}\ge\frac M4.
\]
For the last bound, $2M+1\le3M$ and $M+1\le3M/2$ for $M\ge2$ imply the stronger lower bound $8M/27$.

Replacing $\delta$ by $1$ can only increase the terminal gain. In that larger network, a trap plus one good edge gives the exact gain
\[
 \frac{M^2(3M^2+4)}{(2M+1)(M^3+2M^2+2M+3)}
 \le\frac32+\frac2{M^2}\le2.
\]
This proves \eqref{T10-bad} for the actual $\delta\le1$, for either side. All formulas involve only series/parallel reduction; they do not assume that suppression preserves all-pairs resistance involving the suppressed vertices.
\end{proof}

\subsection{Exact amplification from one pair to all pairs}\label{T10-sec-4}
We state the identity for an arbitrary connected core $H$ on $q$ vertices with distinguished terminals $s,t$. Add $N$ pendant leaves at each terminal, and allow modifications only within the core. Write $R_S(x,y)$ for core resistances after additions $S$, and define
\[
 A(S)=\sum_{x\in V(H)}\bigl(R_S(s,x)+R_S(t,x)\bigr),\qquad
 B(S)=\sum_{\{x,y\}\subseteq V(H)}R_S(x,y).
\]
\begin{lemma}[Leaf amplification identity]\label{T10-amplify}
The Kirchhoff index of the augmented graph is exactly
\begin{equation}\label{T10-Kexact}
 K(S)=4N^2+2(q-1)N+N^2R_S(s,t)+NA(S)+B(S).
\end{equation}
In particular its gain satisfies
\begin{equation}\label{T10-Fexact}
 F(S)=N^2g(S)+E(S),\qquad
 E(S)=N[A(\varnothing)-A(S)]+B(\varnothing)-B(S)\ge0.
\end{equation}
If every baseline core resistance is at most $D$ and $N\ge q$, then
\[
 0\le E(S)\le\tfrac52qND.
\]
The same upper bound applies to $NA(S)+B(S)$ itself.
\end{lemma}
\begin{proof}
A pendant edge carries current only when its leaf is a source or sink. Thus a leaf at $s$ has resistance $1+R_S(s,x)$ to a core vertex $x$; leaves attached to different terminals have resistance $2+R_S(s,t)$; distinct leaves at the same terminal have resistance two. Summing all pair types gives
\[
 4\binom{N}{2}+N^2(2+R_S(s,t))+2qN+NA(S)+B(S),
\]
which simplifies to \eqref{T10-Kexact}.

Adding an edge cannot increase effective resistance: in the minimum-energy unit-flow characterization, the old flow remains feasible with zero flow on the new edge. Hence the differences in \eqref{T10-Fexact} are nonnegative and at most their baseline resistances. Also $A(S)\le2qD$ and $B(S)\le q(q-1)D/2$. If $N\ge q$, their weighted sum is at most $(5/2)qND$.
\end{proof}
This identity is also useful independently of the construction: a fixed terminal objective can be made the leading term of the ordinary all-pairs objective, with a controlled error. The additive constants matter when comparing final objective values rather than gains.

\begin{lemma}[Uniform core bound]\label{T10-diam}
Every baseline effective resistance in $H_M$ is at most $3M^2$. Consequently, with the chosen $N=100qM^2$,
\[
 0\le E(S)\le8qNM^2=\frac{2}{25}N^2,
 \qquad NA(S)+B(S)\le\frac{2}{25}N^2.
\]
\end{lemma}
\begin{proof}
Delete the edge adjacent to $t$ on the short length-$M$ path from $s$ to $t$. The remaining core is a tree: the long $s$--$u$--$t$ path has length $2M^2$, the $s$--$v$ branch has length $M^2$, and the remaining short-path branch from $s$ has length $M-1$. Its diameter is $3M^2$. Sending unit current along a tree path proves that every resistance in this tree is at most $3M^2$, and restoring the deleted edge only lowers resistance. Apply Lemma \ref{T10-amplify} with $D=3M^2$, then round $(15/2)qNM^2$ upward to $8qNM^2$. The choice of $N$ gives the displayed value.
\end{proof}

\subsection{Greedy failure for both objective conventions}\label{T10-sec-5}
\begin{proof}[Proof of Theorem \ref{T10-main}]
By Lemmas \ref{T10-core} and \ref{T10-diam}, for $a=e$ or $f$,
\[
 F(\{h\})-F(\{a\})
 \ge N^2\left(\frac45-\frac12\right)-\frac{2}{25}N^2
 =\frac{11}{50}N^2>0.
\]
Thus greedy must choose $h$ first. Its second choice is one of $e,f$, regardless of any tie rule. For either output,
\[
 F(S_g)\le2N^2+\frac{2}{25}N^2=\frac{52}{25}N^2.
\]
The feasible pair $\{e,f\}$ has gain at least $N^2M/4$, so
\[
 \frac{F(S_g)}{F(S_*)}\le\frac{208}{25M}<\frac9M.
\]
This argument does not need to identify the exact optimal pair.

For final objective values, \eqref{T10-base} and \eqref{T10-bad} yield
\[
 R_{S_g}(s,t)>M-\frac12-2\ge\frac34M.
\]
All terms in \eqref{T10-Kexact} are nonnegative, giving
$K(S_g)\ge(3/4)MN^2$.
On the other hand, $R_{\{e,f\}}<2$ and Lemma \ref{T10-diam} give
\[
 K(\{e,f\})<6N^2+2qN+\frac{2}{25}N^2<7N^2.
\]
The last inequality uses $2q/N=1/(50M^2)\le1/5000$ for $M\ge10$. Optimality gives $K(S_*)\le K(\{e,f\})$, hence
\[
 \frac{K(S_g)}{K(S_*)}>\frac{3M}{28}\ge\frac M{10}.
\]
The vertex count is explicit, $n_M=\Theta(M^4)$, proving the asymptotic statements.
\end{proof}
For a concrete parameter, $M=100$ gives the certified bounds $F(S_g)/F(S_*)\le0.09$ and $K(S_g)/K(S_*)\ge10$. This explicit proof-oriented instance is very large; no small-vertex extremality is claimed. The constants and leaf count were chosen for transparent inequalities, not optimized size.

\subsection{Reproducibility and the scope of the obstruction}\label{T10-sec-6}
The following finite description is an exact construction certificate and requires no external data.
\begin{enumerate}
\item Choose an integer $M\ge10$.
\item Create internally disjoint paths of lengths $M^2,M^2,M$ around the marked cycle $s,u,t,s$, and a length-$M^2$ tail $s,v$.
\item Set $q=3M^2+M$, $N=100qM^2$, and attach $N$ new unit leaves at each of $s,t$.
\item Permit only the three additions $su,ut,vt$, each with unit conductance, and take budget two.
\item Verify the terminal formulas in Lemma \ref{T10-core}, the exact leaf sum \eqref{T10-Kexact}, and the uniform remainder bound. These certify the theorem without constructing the large Laplacian.
\end{enumerate}
The construction is simple and unweighted before and after every allowed addition. Restriction of the candidate set is essential to the theorem as proved. If every absent edge is allowed, extra shortcuts between terminals, path vertices, or leaf collections can change the first greedy choice and invalidate the argument. We make no claim for that larger action set.

Achterberg and Kooij \cite{ref052} distinguish vanishing submodularity ratios from actual greedy inefficiency and pose the search for examples below the classical submodular efficiency threshold in their unrestricted setting. Their introduction also discusses constrained candidate sets, while their default problem permits all nonedges. Our result settles an explicit constrained variant with unbounded failure; it does \emph{not} settle their unchanged unrestricted formulation. Recent positive bicriteria approximation results \cite{ref057} concern different algorithms and may allow extra additions; there is no contradiction. Efficient implementations of greedy \cite{ref058} likewise do not imply a uniform worst-case approximation guarantee.

The complementarity seed is related to the two-update analysis in Taisho round 8. The new work here is the actual greedy trap, realization by equal unit nonedges in an unweighted graph, exact amplification to ordinary all-pairs resistance, and the simultaneous proof of vanishing gain ratio and unbounded final-value ratio. Every claim is proved here. A targeted primary-source search did not locate this exact construction or theorem, but comprehensive priority and independent expert review remain outstanding.





\clearpage
\part{Entropy Inequalities, Marginal Compatibility, and Extremal Laws}


\clearpage
\section[R01 --- Overlap-Affinity Certificates for Hypergraph Product Deficits]{Overlap-Affinity Certificates for Hypergraph Product Deficits}\label{contrib-R01}
\markboth{Entropy Inequalities, Marginal Compatibility, and Extremal Laws}{R01 --- Reiwa}
\noindent\textsc{R01 \quad Reiwa programme, source round 1}\par\medskip


\begin{quote}\small\textbf{Overview.} For functions on overlapping coordinate sets, the classical Finner inequality ignores whether their normalized densities prescribe contradictory laws on shared coordinates. We prove a computable exponential correction: a fractional matching of pairs of factors packs their negative log Hellinger affinities into the logarithmic Finner deficit. An asymmetric version uses weighted geometric affinities. The proof supplies an exact decomposition into marginal relative entropies and conditional mutual informations, establishes a sharp two-factor case, and gives a positive full-support example showing that all overlap certificates can nevertheless miss a genuinely global inconsistency. These are finite-space, unconditional statements. The underlying entropy duality is classical; priority for the particular certificate optimization is not established.\end{quote}

\subsection{Question and scope}\label{R01-sec-1}
Let $[n]=\{1,\ldots,n\}$ and let $\mu=\bigotimes_{i=1}^n\mu_i$ on a finite product $\Omega=\prod_i\Omega_i$, with every $\mu_i$ strictly positive. A factor family is an indexed list $\mathcal E$ of nonempty subsets of $[n]$; repetitions are permitted. Fix $w_e>0$ such that
\begin{equation}\label{R01-cover}c_i:=\sum_{e\ni i}w_e\leq1\qquad(i\in[n]).\end{equation}
Let $f_e>0$, $a_e=\E_{\mu_e}f_e$, and define probability densities
\[
\nu_e(x_e)=\mu_e(x_e)f_e(x_e)/a_e,\qquad
z=\E_\mu\prod_e(f_e/a_e)^{w_e}.
\]
Finner's inequality says $z\leq1$. Our question is narrower than general stability of its extremizers: can contradictions between the explicitly given laws $\nu_e$ certify a quantitative deficit without finding an extremizing joint distribution?

All logarithms are natural. For probabilities $p,q$ with $q>0$, write
\[\D(p\Vert q)=\sum_xp(x)\log(p(x)/q(x)),\]
with $0\log0=0$. The overlap law $\nu_e^S$ is the marginal of $\nu_e$ on $S\subseteq e$. On the empty overlap it is the one-point probability law.

\subsection{The overlap packing theorem}\label{R01-sec-2}
For each pair of distinct factor indices $e,f$, let $S_{ef}=e\cap f$ and put
\[
A_{ef}=\sum_{u\in\Omega_{S_{ef}}}\sqrt{\nu_e^{S_{ef}}(u)\nu_f^{S_{ef}}(u)},
\qquad b_{ef}=-2\log A_{ef}\geq0.
\]
Define the finite linear program
\begin{equation}\label{R01-LP}
 B=\max\left\{\sum_{\{e,f\}}\lambda_{ef}b_{ef}:\lambda_{ef}\geq0,
 \ \sum_{f\ne e}\lambda_{ef}\leq w_e\ \text{for every }e\right\}.
\end{equation}
\begin{theorem}\label{R01-main}
Under the hypotheses above,
\begin{equation}\label{R01-bound}
\E_\mu\prod_e f_e^{w_e}\leq e^{-B}\prod_e a_e^{w_e}.
\end{equation}
Thus pairwise overlap inconsistency yields a certified multiplicative improvement. The coefficient $2$ in $b_{ef}$ is optimal uniformly over the stated class.
\end{theorem}

The LP has one variable per factor pair, not per point of $\Omega$. In particular it remains small when the joint alphabet is large and each overlap alphabet is small. Empty or identical overlap laws have zero cost and can be omitted.

\subsubsection{An exact entropic identity}
Define the positive probability $P$ by
\[
P(x)=z^{-1}\mu(x)\prod_e(f_e(x_e)/a_e)^{w_e}.
\]
For an ordered list of coordinates, write $[i-1]=\{1,\ldots,i-1\}$.
\begin{lemma}\label{R01-identity}
For this $P$,
\begin{equation}\label{R01-identityeq}
-\log z=\sum_e w_e\D(P_e\Vert\nu_e)+\mathcal S_w(P),
\qquad
\mathcal S_w(P)=\D(P\Vert\mu)-\sum_e w_e\D(P_e\Vert\mu_e).
\end{equation}
Moreover, for every probability $Q$,
\begin{align}\label{R01-chain}
\mathcal S_w(Q)
={}&\sum_i(1-c_i)\E_Q\D(Q_{i\mid[i-1]}\Vert\mu_i)\\
&+\sum_e w_e\sum_{i\in e}
 I_Q\bigl(X_i;X_{[i-1]\setminus e}\mid X_{e\cap[i-1]}\bigr)\geq0.\notag
\end{align}
\end{lemma}
\begin{proof}
From the definition of $P$,
\[
\D(P\Vert\mu)=\sum_e w_e\E_P\log(f_e/a_e)-\log z.
\]
On the other hand $\D(P_e\Vert\nu_e)=\D(P_e\Vert\mu_e)-\E_P\log(f_e/a_e)$. Subtraction gives \eqref{R01-identityeq}.
For \eqref{R01-chain}, the chain rule for relative entropy gives
\[
\D(Q\Vert\mu)=\sum_i d_i,\quad
 d_i=\E_Q\D(Q_{i\mid[i-1]}\Vert\mu_i).
\]
Likewise the edge divergence is the sum of terms $d_{i,e}=\E_Q\D(Q_{i\mid e\cap[i-1]}\Vert\mu_i)$. Expanding the logarithms shows
\[
d_i-d_{i,e}=H_Q(X_i\mid X_{e\cap[i-1]})-H_Q(X_i\mid X_{[i-1]})
=I_Q(X_i;X_{[i-1]\setminus e}\mid X_{e\cap[i-1]}).
\]
Regroup $\sum_i d_i-\sum_e w_e\sum_{i\in e}d_{i,e}$. Nonnegativity follows from \eqref{R01-cover}, relative entropy, and conditional mutual information. Their nonnegativity follows in turn from the log-sum inequality (or Jensen applied to $-\log$).
\end{proof}
This also proves the finite Finner inequality, rather than assuming it.

\subsubsection{Pairwise geometric affinity}
\begin{lemma}\label{R01-affinity}
For strictly positive probabilities $p,q$, $a,b>0$, and any probability $r$,
\begin{equation}\label{R01-geom}
a\D(r\Vert p)+b\D(r\Vert q)\geq
-(a+b)\log\sum_u p(u)^{a/(a+b)}q(u)^{b/(a+b)}.
\end{equation}
Equality holds exactly when $r$ is the normalized geometric mean inside the sum.
\end{lemma}
\begin{proof}
Let $G=\sum p^{a/(a+b)}q^{b/(a+b)}$ and $g=p^{a/(a+b)}q^{b/(a+b)}/G$. Direct expansion identifies the left side as $(a+b)\D(r\Vert g)-(a+b)\log G$. The assertion follows. Notice that the divergence direction is essential: minimizing $\D(r\Vert p)$ gives a geometric, not arithmetic, barycenter.
\end{proof}
\begin{proof}[Proof of Theorem \ref{R01-main}]
Marginalization decreases relative entropy: the chain rule gives $\D(P_e\Vert\nu_e)\geq\D(P_S\Vert\nu_e^S)$. Thus each pair obeys
\[
\D(P_e\Vert\nu_e)+\D(P_f\Vert\nu_f)\geq b_{ef}
\]
by Lemma \ref{R01-affinity} with $a=b=1$. Multiply by a feasible $\lambda_{ef}$ and sum. Each marginal divergence is charged at most $w_e$ times. Lemma \ref{R01-identity} then yields $-\log z\geq\sum\lambda_{ef}b_{ef}$. Maximizing proves \eqref{R01-bound}.

For sharpness take $n=1$, two copies $e=f=\{1\}$ with weights $1/2$, any strictly positive $p,q$, and $f_e=p/\mu$, $f_f=q/\mu$. Then $z=\sum\sqrt{pq}$ and the LP optimum is $B=\frac12(-2\log z)=-\log z$. Any larger universal coefficient would fail when $p\ne q$.
\end{proof}

\subsection{Asymmetric budgets and usable certificates}\label{R01-sec-3}
The symmetric LP is not the only way to spend the available divergence. Assign to every unordered pair two nonnegative endpoint budgets $a_{ef}$ and $a_{fe}$, with $\sum_{f\ne e}a_{ef}\leq w_e$. Define, when both budgets are positive,
\[
C_{ef}(a,b)=-(a+b)\log\sum_u(\nu_e^{S_{ef}}(u))^{a/(a+b)}
(\nu_f^{S_{ef}}(u))^{b/(a+b)},
\]
and set $C_{ef}(a,0)=C_{ef}(0,b)=0$.
\begin{corollary}\label{R01-asym}
For every such budget assignment, $-\log z\geq\sum_{\{e,f\}}C_{ef}(a_{ef},a_{fe})$.
The right side is concave in the budgets, so maximizing it over the budget polytope is a convex optimization problem in the usual ``maximize concave'' convention.
\end{corollary}
\begin{proof}
Apply \eqref{R01-geom} after marginalization and sum. Concavity follows also without differentiation: $C_{ef}(a,b)=\inf_r\{a\D(r\Vert p)+b\D(r\Vert q)\}$, an infimum of affine functions of $(a,b)$. Continuity at the boundary follows from positivity and the finite alphabet.
\end{proof}
A submitted feasible budget is already a proof certificate: it needs no assertion that a numerical optimizer found the maximum. For \eqref{R01-LP}, LP duality additionally supplies
\[
B=\min\left\{\sum_e w_e y_e:y_e\geq0,\quad y_e+y_f\geq b_{ef}\right\}.
\]
A matching feasible primal and dual pair certifies the exact optimum. To avoid relying on numerical logarithms, interval lower bounds on each $b_{ef}$ suffice for a rigorous lower certificate.

\subsection{A blind spot: globally inconsistent binary marginals}\label{R01-sec-4}
Pairwise compatibility is not global compatibility. Let $n=3$, factors $12,23,13$, and $w_e=1/2$. For $0<\epsilon<1/3$, let $\nu_{12}$ and $\nu_{23}$ prefer equal bits with probability $1-\epsilon$, and let $\nu_{13}$ prefer unequal bits with probability $1-\epsilon$. In each case assign half the stated mass to each of the two outcomes. Every one-coordinate marginal is uniform, hence $A_{ef}=1$ and every certificate above is zero.

Nevertheless no joint law can have these three pair marginals. For every bit triple at least one preferred relation fails; expectations would give at least one failure, whereas the proposed marginals give $3\epsilon<1$. With uniform product $\mu$ and $f_e=\nu_e/\mu_e$, direct enumeration gives
\[
z=\frac{\sqrt2}{2}\sqrt\epsilon\,(3-2\epsilon)<1.
\]
There are six assignments with one failure and two with three failures. The strict inequality also follows from Lemma \ref{R01-identity}: equality would require $P_e=\nu_e$ for all three edges. This example shows exactly what is not proved: overlap certificates do not characterize all equality or near-equality cases.

\subsection{Literature position and verification}\label{R01-sec-5}
The Finner inequality is classical \cite{ref059}. Entropy-functional duality and equality correspondence belong to the framework of Carlen--Cordero-Erausquin \cite{ref060}; broader entropic Brascamp--Lieb formulations are in \cite{ref061}. Current stability work includes Bennett--Oh \cite{ref062}. We do not claim these general principles as discoveries, nor claim to settle a named open problem. The candidate contribution is the explicitly budgeted overlap-affinity certificate and its full-support obstruction example. A limited primary-source search does not establish that this formulation is absent from the literature.

Every proof is finite and reproduced here. No computer algebra, numerical search, external data, or unlisted lemma is required. The manuscripts in the tournament benchmark were screened for related entropy themes; no benchmark theorem is invoked. The present result is narrower in reach than a complete solution of a long-standing open classification problem.




\clearpage
\section[R02 --- Global Marginal Inconsistency: Attained Dual Certificates and Quantitative Recovery]{Global Marginal Inconsistency: Attained Dual Certificates and Quantitative Recovery}\label{contrib-R02}
\markboth{Entropy Inequalities, Marginal Compatibility, and Extremal Laws}{R02 --- Reiwa}
\noindent\textsc{R02 \quad Reiwa programme, source round 2}\par\medskip


\begin{quote}\small\textbf{Overview.} Given positive local probability laws on overlapping subsets, we study the least weighted relative-entropy cost of replacing them by marginals of one joint law. We give a finite, attained dual certificate involving local test functions and a global minimum; prove uniqueness of the optimal marginal tuple without claiming uniqueness of the joint law; derive an exact Bregman remainder and a posteriori total-variation error bounds; and prove tensorization. We then insert this global obstruction into Finner's inequality and evaluate a frustrated triangle where every pairwise overlap test is identically zero. The general convex-analytic framework is classical. The paper is a self-contained certificate-oriented synthesis, not a claimed solution of a previously named open problem.\end{quote}

\subsection{The global optimization problem}\label{R02-sec-1}
Let $\Omega=\prod_{i=1}^n\Omega_i$ be finite, $\mathcal E$ an indexed family of nonempty coordinate subsets, $w_e>0$, and $\nu_e$ strictly positive probabilities on $\Omega_e$. For a joint probability $P$, put
\[
F(P)=\sum_e w_e\D(P_e\Vert\nu_e),\qquad
C(\nu,w)=\min_{P\in\Pp(\Omega)}F(P).
\]
Here $\D(p\Vert q)=\sum p\log(p/q)$, with natural logarithms. No fractional-cover hypothesis is needed until Section \ref{R02-finner}. The image of the joint probability simplex under the marginal maps is the \emph{marginal polytope}. The problem is a separable entropic projection onto this polytope.

The distinction between local consistency and a common joint extension is important. For a cycle, consistent overlaps alone need not be extendible. The quantity $C$ detects exactly the latter obstruction, but is not promised to be easy to compute on arbitrary hypergraphs.

\subsection{An attained dual without a separation theorem}\label{R02-sec-2}
For arbitrary real functions $h_e$ on $\Omega_e$, define
\begin{equation}\label{R02-dual}
L(h)=\min_{x\in\Omega}\sum_e w_eh_e(x_e)
      -\sum_e w_e\log\E_{\nu_e}e^{h_e}.
\end{equation}
\begin{theorem}\label{R02-duality}
The minimum defining $C$ is attained. All minimizers have the same marginal tuple $(q_e)_e$, and every $q_e$ is strictly positive. Moreover,
\begin{equation}\label{R02-strong}
C=\max_{(h_e)}L(h),
\end{equation}
and one maximizing choice is $h_e=\log(q_e/\nu_e)$. Finally $C=0$ if and only if the laws $\nu_e$ have a joint extension.
\end{theorem}
\begin{proof}
Continuity of $F$ and compactness of the joint simplex give attainment. Strict convexity of $p\mapsto\D(p\Vert\nu_e)$ implies uniqueness of the marginal tuple: if two minimizing tuples differ on an edge, their midpoint has strictly smaller objective.

To prove positivity, suppose an optimal $P^*$ has $(P^*)_e(u)=0$ for some $e,u$. Let $U$ be the uniform joint law and $P_t=(1-t)P^*+tU$. At each zero marginal cell the change in $s\log(s/\nu)$ is a positive multiple of $t\log t+O(t)$; at each positive cell it is $O(t)$. There are finitely many cells and every weight is positive. Hence $F(P_t)-F(P^*)<0$ for all sufficiently small $t>0$, a contradiction.

The elementary variational inequality
\begin{equation}\label{R02-var}
\D(p\Vert\nu)\geq\E_p h-\log\E_\nu e^h
\end{equation}
follows by subtracting the right side and recognizing relative entropy to $\nu e^h/\E_\nu e^h$. Summing \eqref{R02-var} and bounding the expectation of $\sum w_eh_e$ below by its minimum proves $F(P)\geq L(h)$ for every $P,h$.

Write $q_e=(P^*)_e$ and $h_e=\log(q_e/\nu_e)$. Differentiability at the positive tuple gives, for every point mass $\delta_x$, the one-sided optimality condition
\[
0\leq\left.\frac{d}{dt}F((1-t)P^*+t\delta_x)\right|_{t=0+}
 =\sum_e w_eh_e(x_e)-\sum_e w_e\E_{q_e}h_e.
\]
The constants $1$ in the derivatives cancel since each marginal has mass one. The last expectation equals $C$. Thus $g(x):=\sum_e w_eh_e(x_e)\geq C$ everywhere. Since $\E_{P^*}g=C$, its minimum is $C$, and it equals $C$ on the support of $P^*$. Also $\E_{\nu_e}e^{h_e}=1$ for each edge. Hence $L(h)=C$. Nonnegativity of every divergence proves the final equivalence.
\end{proof}

\begin{corollary}[Finite exact certificates]\label{R02-cert}
A joint probability $P$ and local test functions $h_e$ certify optimality if
\[
P_e=\frac{\nu_e e^{h_e}}{\E_{\nu_e}e^{h_e}}\quad\text{for every }e,
\qquad
\supp P\subseteq\argmin_x\sum_e w_eh_e(x_e).
\]
Then $F(P)=L(h)=C$. Conversely, every minimizing joint law admits such a certificate.
\end{corollary}
\begin{proof}
These are exactly the equality conditions in the two inequalities proving weak duality. The maximizing functions constructed in Theorem \ref{R02-duality} prove the converse.
\end{proof}
A lower bound on the global minimum in \eqref{R02-dual} is sufficient for a valid lower certificate. A numerical local minimum is not: its use would reverse the needed inequality. Exhaustive enumeration is always a finite verification method, although possibly exponential.

\subsection{A quantitative remainder and certified approximation}\label{R02-sec-3}
Let $(q_e)_e$ be the unique minimizing tuple and choose any minimizer $P^*$. Set $g(x)=\sum_e w_e\log(q_e(x_e)/\nu_e(x_e))$.
\begin{theorem}\label{R02-pyth}
For every joint law $P$,
\begin{equation}\label{R02-pytheq}
F(P)-C=\sum_e w_e\D(P_e\Vert q_e)+\E_P[g-C]
\geq\sum_e w_e\D(P_e\Vert q_e).
\end{equation}
Consequently, if $\ell\leq C\leq F(P)=u$, then
\[
\sum_e w_e\TV(P_e,q_e)^2\leq\frac{u-\ell}{2},\qquad
\TV(P_e,q_e)\leq\sqrt{\frac{u-\ell}{2w_e}}.
\]
Here $\TV(p,q)=\frac12\sum|p-q|$.
\end{theorem}
\begin{proof}
For each edge, expansion gives the exact identity
\[
\D(P_e\Vert\nu_e)-\D(q_e\Vert\nu_e)
=\D(P_e\Vert q_e)+\sum_u(P_e(u)-q_e(u))\log\frac{q_e(u)}{\nu_e(u)}.
\]
Sum and use $\E_{P^*}g=C$ and $g\geq C$ from the preceding proof. For completeness, Pinsker's inequality $\D(p\Vert q)\geq2\TV(p,q)^2$ follows as follows. Let $A=\{p\geq q\}$ and apply log-sum to the two cells $A,A^c$. The resulting binary divergence $d(a\Vert b)$ satisfies $\partial_a^2d=1/(a(1-a))\geq4$, $d(b\Vert b)=\partial_ad(b\Vert b)=0$, so $d(a\Vert b)\geq2(a-b)^2$. Here $a-b=\TV(p,q)$. Endpoint cases follow by limits. Inserting this into \eqref{R02-pytheq} proves the bounds.
\end{proof}
These bounds concern the identifiable marginal tuple, not the full joint law. No analogous recovery of $P^*$ is possible in general: when only one-coordinate marginals are prescribed, many different joint laws share the same optimal tuple.

\subsection{Tensorization and monotonicity}\label{R02-sec-4}
Consider two systems $A,B$ with the same edge labels and weights, on alphabets $\Omega_i^A,\Omega_i^B$. Define product local laws $\nu_e^{AB}=\nu_e^A\otimes\nu_e^B$.
\begin{theorem}\label{R02-tensor}
Their inconsistency costs satisfy
\[
C(\nu^{AB},w)=C(\nu^A,w)+C(\nu^B,w).
\]
The unique minimizing marginal on each combined edge is $q_e^A\otimes q_e^B$.
\end{theorem}
\begin{proof}
For any combined joint law $P$, a direct entropy expansion gives
\[
\D(P_{A_eB_e}\Vert\nu_e^A\otimes\nu_e^B)
=\D(P_{A_e}\Vert\nu_e^A)+\D(P_{B_e}\Vert\nu_e^B)+I_P(A_e;B_e).
\]
Summing proves the lower bound by nonnegativity of mutual information. A product of minimizing joint laws gives equality and the asserted marginal tuple. Its uniqueness was proved above.
\end{proof}
If some local coordinates are erased from each edge, keeping a consistent global coordinate erasure, data processing shows that $C$ cannot increase. Similarly, adding additional positively weighted marginal requirements cannot decrease $C$. These statements follow directly by applying the relevant marginal map to every feasible joint law; they do not require any asserted computational algorithm.

\subsection{A global correction to Finner}\label{R02-sec-5}\label{R02-finner}
Suppose now that $\mu$ is a strictly positive product measure and $\sum_{e\ni i}w_e\leq1$. Given $f_e>0$, put $a_e=\E_{\mu_e}f_e$ and $\nu_e=\mu_e f_e/a_e$.
\begin{corollary}\label{R02-fincor}
\[
\E_\mu\prod_e f_e^{w_e}\leq\exp(-C(\nu,w))\prod_e a_e^{w_e}.
\]
Every dual certificate $L(h)$ may replace $C$ on the right.
\end{corollary}
\begin{proof}
Let $z$ be the left side divided by $\prod a_e^{w_e}$, and $P=z^{-1}\mu\prod(f_e/a_e)^{w_e}$. Direct logarithmic expansion gives
\[
-\log z=F(P)+\D(P\Vert\mu)-\sum_e w_e\D(P_e\Vert\mu_e).
\]
The final difference is nonnegative. Indeed, put $d_i=\E\D(P_{i\mid[i-1]}\Vert\mu_i)$. For $i\in e$, the corresponding chain-rule term $d_{i,e}$ using conditioning only on $e\cap[i-1]$ is at most $d_i$; the difference is a conditional mutual information. Thus
$\sum_e w_e\D(P_e\Vert\mu_e)\leq\sum_i(\sum_{e\ni i}w_e)d_i\leq\sum_i d_i=\D(P\Vert\mu)$.
Now $F(P)\geq C\geq L(h)$.
\end{proof}
The discarded entropy difference can be positive. Therefore $C$ is the exact marginal-inconsistency cost, not asserted to be the entire Finner deficit.

\subsection{Exact frustrated-triangle calculation}\label{R02-sec-6}
Take three binary variables and edges $12,23,13$, all with weight $1/2$. Let the first two edges prefer equality and the third prefer inequality. For each edge let the preferred relation have probability $1-\epsilon$ and the other probability $\epsilon$, splitting each between its two outcomes equally. Assume $0<\epsilon<1/3$. Define
\[
d(p\Vert\epsilon)=p\log(p/\epsilon)+(1-p)\log((1-p)/(1-\epsilon)).
\]
\begin{proposition}\label{R02-tri}
The exact global cost is
\[
C=\frac32d(1/3\Vert\epsilon)>0.
\]
All pairwise shared-coordinate marginals nevertheless agree.
\end{proposition}
\begin{proof}
Let $p_e$ be a joint law's probability of violating edge $e$. At least one edge is violated by every assignment, hence $\sum_ep_e\geq1$. Data processing to the violation indicator gives
\[
F(P)\geq\frac12\sum_ed(p_e\Vert\epsilon)
\geq\frac32d\left(\frac{p_{12}+p_{23}+p_{13}}3\,\middle\Vert\,\epsilon\right)
\geq\frac32d(1/3\Vert\epsilon).
\]
The middle step is Jensen. The last uses $\epsilon<1/3$ and the derivative
$\partial_pd(p\Vert\epsilon)=\log\frac{p(1-\epsilon)}{\epsilon(1-p)}>0$ for $p>\epsilon$.
Take $P$ uniform on the six assignments with exactly one violation. Every edge then has violation probability $1/3$, and its two assignments in each relation are equally likely. Both data processing and Jensen are equalities. Each original one-coordinate marginal is uniform, proving the last assertion.
\end{proof}
A concise dual certificate is $h_e=\log(q_e/\nu_e)$ for these attained marginals. Its sum has its minimum precisely on the six one-violation configurations because $\epsilon<1/3$. This independently checks the optimality conditions in Corollary \ref{R02-cert}.

\subsection{Status and relation to earlier work}\label{R02-sec-7}
Marginal polytopes and entropy variational representations are standard \cite{ref063}; entropic Brascamp--Lieb duality is classical \cite{ref060}. This manuscript supplies elementary attainment, a certificate remainder, tensorization, and a concrete globally frustrated application in one finite formulation. It builds conceptually on the author's round-1 observation that pairwise overlap tests can miss global inconsistency, but none of its proofs require that manuscript. It is a genuine change from pairwise certificate packing to a global convex projection and recovery theorem. No claim of first discovery of convex duality or entropic projection is made; a complete priority search remains open. No pre-existing named open question is claimed resolved.




\clearpage
\section[R03 --- Exact Weighted Entropic Projection onto Parity Polytopes]{Exact Weighted Entropic Projection onto Parity Polytopes}\label{contrib-R03}
\markboth{Entropy Inequalities, Marginal Compatibility, and Extremal Laws}{R03 --- Reiwa}
\noindent\textsc{R03 \quad Reiwa programme, source round 3}\par\medskip


\begin{quote}\small\textbf{Overview.} For either binary parity polytope and arbitrary positive coordinate weights, we determine the projection of every interior point under the separable Bernoulli relative entropy. If the point is infeasible, precisely one forbidden-vertex facet is violated. After coordinate reflections, the full optimization reduces to one strictly monotone scalar logistic equation; the projected point lies in the simplex of neighbors of the forbidden vertex. We supply an elementary proof of the parity-polytope description, a linear-time facet-identification procedure, an explicit attaining distribution, and a quadratic boundary expansion. Applied to any heterogeneous noisy binary cycle, this gives the exact global marginal-inconsistency cost for all noise parameters, together with a closed partition-function formula. The parity-polytope framework is established prior art; priority for the weighted entropic projection formula is unverified.\end{quote}

\subsection{Problem and main classification}\label{R03-sec-1}
For $m\geq3$ and $s\in\{0,1\}$ let
\[
\mathcal P_s=\conv\{v\in\{0,1\}^m:\ \sum_iv_i\equiv s\pmod2\}.
\]
Let $a\in(0,1)^m$ and $w_i>0$. Define
\[
J_{a,w}(p)=\sum_{i=1}^mw_i d(p_i\Vert a_i),\qquad
 d(p\Vert a)=p\log(p/a)+(1-p)\log((1-p)/(1-a)).
\]
The minimizer on $\mathcal P_s$ exists and is unique by compactness and strict convexity. Our goal is its explicit classification, not just a generic convex-program formulation.
For a binary vertex $u$, let
\[
\ell_u(p)=\sum_{i:u_i=0}p_i+\sum_{i:u_i=1}(1-p_i)=\|p-u\|_1.
\]
Call $u$ forbidden if its parity is $1-s$.

\begin{theorem}\label{R03-main}
Exactly one of the following cases holds.
\begin{enumerate}[label=(\roman*)]
\item If $\ell_u(a)\geq1$ for every forbidden $u$, then $a\in\mathcal P_s$ and the minimum is zero at $p=a$.
\item Otherwise there is a unique forbidden vertex $u$ with $\ell_u(a)<1$. Set $b_i=a_i$ when $u_i=0$ and $b_i=1-a_i$ when $u_i=1$. There is a unique $\lambda>0$ satisfying
\begin{equation}\label{R03-root}
\sum_{i=1}^m \frac{b_i e^{\lambda/w_i}}{1-b_i+b_i e^{\lambda/w_i}}=1.
\end{equation}
Write the summands as $r_i$. The unique optimizer is $p_i^*=r_i$ if $u_i=0$ and $p_i^*=1-r_i$ if $u_i=1$, and
\begin{align}\label{R03-value}
\min_{p\in\mathcal P_s}J_{a,w}(p)
&=\sum_iw_id(r_i\Vert b_i)\\
&=\lambda-\sum_iw_i\log(1-b_i+b_ie^{\lambda/w_i}).\notag
\end{align}
The point $p^*$ is realized by the probability law assigning mass $r_i$ to the allowed vertex $u\mathbin\oplus e_i$.
\end{enumerate}
\end{theorem}
Here $\oplus$ denotes addition modulo two and $e_i$ is the $i$th unit vector. Thus in the infeasible case the optimal realizing distribution needs only $m$ support points, irrespective of the $2^{m-1}$ vertices of the parity polytope.

\subsection{The polyhedral fact, with proof}\label{R03-sec-2}
\begin{lemma}\label{R03-poly}
\begin{equation}\label{R03-facets}
\mathcal P_s=\{p\in[0,1]^m:\ell_u(p)\geq1\text{ for every forbidden }u\}.
\end{equation}
At most one of these inequalities can fail at any point of the cube.
\end{lemma}
\begin{proof}
Every allowed binary vertex has odd Hamming distance from every forbidden vertex, hence satisfies all displayed inequalities. Let $K$ denote their intersection with the cube. We show that every extreme point of $K$ is integral.

Suppose $p$ has a nonempty set $T$ of fractional coordinates. If two distinct forbidden vertices $u,v$ give active constraints, then
\[
2=\ell_u(p)+\ell_v(p)\geq\|u-v\|_1\geq2.
\]
Their Hamming distance is even, so equality forces distance exactly two. At every coordinate where $u_i=v_i$, equality forces $p_i=u_i$. Thus $T$ is contained in the two coordinates where $u,v$ differ.

If $|T|\geq3$, at most one forbidden constraint is active. On the fractional coordinates there is consequently a nonzero direction preserving that active affine equation (if any). A sufficiently small perturbation in both signs preserves all inactive strict inequalities and cube constraints. Hence $p$ is not extreme.
If $|T|=2$, the restrictions of any two active forbidden-constraint normals to $T$ are opposites, by the preceding argument. Their span has dimension at most one, so the same two-sided perturbation exists. If $|T|=1$, no forbidden constraint can be active: its value would be an integer plus a number strictly between zero and one, and so could not equal one. Again a perturbation exists. This proves that every extreme point is integral. Integral points of $K$ have parity $s$, since a forbidden vertex violates its own constraint. A bounded polytope is the convex hull of its extreme points, proving \eqref{R03-facets}.

Finally two violated constraints would imply $\ell_u(p)+\ell_v(p)<2$, contrary to the Hamming-distance inequality above.
\end{proof}
The forbidden-set description is classical; this proof is included to make the projection argument independently checkable.

\subsection{Proof of the projection formula}\label{R03-sec-3}
Coordinate reflections $p_i\mapsto1-p_i$ preserve Bernoulli relative entropy when the target is reflected as well. Reflect at the unique forbidden $u$. This maps the allowed parity class to odd parity, maps $u$ to zero, and maps the violated inequality to $\sum r_i\geq1$. The reflected target $b$ satisfies $\sum b_i<1$.

First minimize $\sum_i w_i d(r_i\Vert b_i)$ over the larger set $r\in[0,1]^m$, $\sum r_i\geq1$. For each $\lambda\geq0$, elementary one-variable minimization gives
\begin{equation}\label{R03-conj}
 w_i d(t\Vert b_i)\geq\lambda t-w_i\log(1-b_i+b_i e^{\lambda/w_i}),
\end{equation}
with equality at $t=b_i e^{\lambda/w_i}/(1-b_i+b_ie^{\lambda/w_i})$. Indeed, subtracting the right side gives $w_i d(t\Vert r_i(\lambda))$.
The sum in \eqref{R03-root} is continuous and strictly increasing, has value $\sum b_i<1$ at zero, and tends to $m>1$ as $\lambda\to\infty$. Thus there is exactly one positive root. At that root $r_i>0$ and $\sum r_i=1$, so $r$ belongs to the simplex $\conv\{e_1,\ldots,e_m\}\subset\mathcal P_1$. Summing \eqref{R03-conj} shows that it minimizes the objective even on the larger halfspace and attains \eqref{R03-value}. Therefore it also minimizes on the parity polytope. Reflecting back proves Theorem \ref{R03-main}, including the explicit realizing law. Case (i) follows from Lemma \ref{R03-poly} and nonnegativity of divergence.

\subsection{Finding the violated facet and certifying the root}\label{R03-sec-4}
There is no need to enumerate the forbidden vertices. Choose a nearest cube vertex $v$ by rounding each $a_i$ at $1/2$; break ties arbitrarily. Then $\ell_v(a)=\sum_i\min(a_i,1-a_i)$. If $v$ is forbidden, it is a nearest forbidden vertex. If not, a nearest forbidden vertex is obtained by flipping a coordinate minimizing $|1-2a_i|$.

To justify the latter rule, changing a set of rounded coordinates adds $\sum_{i\text{ changed}}|1-2a_i|$ to the distance. The parity change requires an odd number of flips. Since the summands are nonnegative, the cheapest odd set is one cheapest coordinate. This establishes an $O(m)$ arithmetic procedure for finding the minimum forbidden distance. A distance at least one certifies feasibility; a distance less than one identifies the unique violated facet.

In the infeasible case, bisection for \eqref{R03-root} is globally valid. An explicit upper bracket is
\[
\Lambda=\max_i\left\{w_i\log\frac{1-b_i}{b_i}\right\}_+ +1.
\]
At $\lambda=\Lambda$ every logistic summand exceeds $1/2$, so their sum exceeds one. Interval arithmetic can certify a bracket and the value without assuming convergence of an unconstrained numerical solver. The scalar concave function in \eqref{R03-value}, evaluated at any $\lambda\geq0$, is a rigorous lower bound. A feasible simplex point gives an upper bound through its explicit objective.

\begin{proposition}[Quadratic onset at a single facet]\label{R03-quad}
Fix positive weights. Consider reflected targets $b$ in a compact subset of $(0,1)^m$ with $\delta=1-\sum_i b_i>0$ tending to zero. Define
\[
V(b,w)=\sum_i\frac{b_i(1-b_i)}{w_i}.
\]
Then, uniformly on that compact subset,
\[
\lambda=\frac{\delta}{V(b,w)}+O(\delta^2),\qquad
\min J=\frac{\delta^2}{2V(b,w)}+O(\delta^3).
\]
\end{proposition}
\begin{proof}
Taylor expansion of the logistic sum at zero gives $1-\delta+V\lambda+O(\lambda^2)$. Compactness bounds $V$ above and away from zero and controls the remainder uniformly. The implicit equation yields the first formula. Expanding $d(b_i+t\Vert b_i)=t^2/(2b_i(1-b_i))+O(t^3)$ and $r_i-b_i=\lambda b_i(1-b_i)/w_i+O(\lambda^2)$ gives the second.
\end{proof}

\subsection{Complete noisy-cycle marginal classification}\label{R03-sec-5}
Let $X_1,\ldots,X_m\in\{0,1\}$, cyclic indices, and prescribed edge labels $c_i\in\{0,1\}$. Edge $i$ is violated when
\[
Y_i=X_i\oplus X_{i+1}\oplus c_i=1,
\qquad s=\bigoplus_i c_i.
\]
Always $\bigoplus_iY_i=s$. Conversely each vector of parity $s$ comes from exactly two spin assignments, related by global bit flip.
Give edge $i$ the strictly positive target law $\nu_i$ with violation probability $a_i\in(0,1)$, distributing each relation's mass equally between its two spin pairs. Let
\[
C=\min_P\sum_iw_i\D(P_{i,i+1}\Vert\nu_i).
\]
\begin{theorem}\label{R03-cycle}
For all positive weights and all noise parameters,
\[
C=\min_{p\in\mathcal P_s}\sum_iw_id(p_i\Vert a_i).
\]
Consequently Theorem \ref{R03-main} computes $C$ completely. In the infeasible case an attaining joint law is obtained by sampling $Y=u\oplus e_i$ with probabilities $r_i$ and then choosing uniformly between its two spin lifts.
\end{theorem}
\begin{proof}
Marginalization to $Y_i$ gives $\D(P_{i,i+1}\Vert\nu_i)\geq d(p_i\Vert a_i)$, and the vector $p$ belongs to $\mathcal P_s$. Conversely any $p$ in the parity polytope is the mean of a law on parity-$s$ vectors. Lift it using a uniform global bit flip. For each edge, conditional on its violation indicator the two spin pairs are equiprobable, exactly as in $\nu_i$. Equality holds in data processing on every edge. The displayed formula and the construction follow.
\end{proof}
This covers both frustrated and unfrustrated cycles and arbitrary, not necessarily small, noise. All adjacent one-vertex marginals of the targets are uniform, so pairwise overlap tests detect none of these obstructions.

For completeness the associated product integral can also be computed exactly. With uniform product measure $\mu$ on the spins and $f_i=\nu_i/\mu_{i,i+1}$, put $W=\sum_iw_i$. Direct expansion over the parity class yields
\begin{align*}
\E_\mu\prod_i f_i^{w_i}
=2^{W-m}\bigg[&\prod_i\big((1-a_i)^{w_i}+a_i^{w_i}\big)\\
&+(-1)^s\prod_i\big((1-a_i)^{w_i}-a_i^{w_i}\big)\bigg].
\end{align*}
Indeed the indicator of parity $s$ is $(1+(-1)^{s+\sum y_i})/2$, and every violation vector has two spin lifts. If $w_{i-1}+w_i\leq1$, this integral is at most $e^{-C}$ by the elementary entropic proof of Finner's inequality. The formula above also permits independent exact checks.

\subsection{Prior art and claim limits}\label{R03-sec-6}
Parity polytopes and their forbidden-set inequalities are established tools of coding theory; see Feldman--Wainwright--Karger \cite{ref064}. Marginal-polytope variational methods are classical \cite{ref063}. The specific object here is a weighted separable Bernoulli KL projection, and the mathematical contribution sought is the full all-parameter scalar classification and its joint-law realization, rather than the known polytope itself. Priority for this exact formulation is unverified. This is not advertised as a solution of an unchanged named open question. All formulas and proofs needed to verify the claims appear in this single source.




\clearpage
\section[R04 --- Sharp Graph-Supported Entropy Stability and Positive Equality Cases]{Sharp Graph-Supported Entropy Stability and Positive Equality Cases}\label{contrib-R04}
\markboth{Entropy Inequalities, Marginal Compatibility, and Extremal Laws}{R04 --- Reiwa}
\noindent\textsc{R04 \quad Reiwa programme, source round 4}\par\medskip


\begin{quote}\small\textbf{Overview.} For an arbitrary nonnegative weighted graph, we identify the exact largest ratio of weighted pairwise mutual information to total correlation over all finite joint laws. It is the weighted fractional arboricity, namely the maximum induced-edge weight divided by one less than the number of vertices. This yields the optimal total-correlation remainder for every graph-supported fractional entropy partition. We give a self-contained forest-polytope proof, an exact relative-entropy remainder, binary extremizers, and a classification of all strictly positive equality cases: they are precisely the forest-Markov laws supported on edges whose individual weight attains the extremal density. In particular many sharp nontrivial extremizers must lie on the boundary of the probability simplex. Forest entropy bounds and tree reweighting are classical; priority for this exact stability/equality formulation is not established.\end{quote}

\subsection{The extremal constant}\label{R04-sec-1}
Let $G=(V,E)$ be a finite simple graph, $|V|=n\geq2$, with weights $a_e\geq0$, not all zero. Delete zero-weight edges when referring to the support graph. Set
\[
\rho(G,a)=\max_{S\subseteq V,\ |S|\geq2}\frac{a(E(S))}{|S|-1},
\qquad a(A)=\sum_{e\in A}a_e.
\]
For finite random variables $X_V$, write
\[
T(X_V)=\sum_{i\in V}H(X_i)-H(X_V),\qquad I_{ij}=H(X_i)+H(X_j)-H(X_i,X_j).
\]
Entropy is in natural units. The alphabets need not be equal and the law may have zeros.
\begin{theorem}\label{R04-sharp}
For every joint law,
\begin{equation}\label{R04-mi}
\sum_{ij\in E}a_{ij}I_{ij}\leq\rho(G,a)T(X_V).
\end{equation}
The coefficient is best possible, already for binary variables:
\[
\sup_{T(X_V)>0}\frac{\sum_ea_eI_e}{T(X_V)}=\rho(G,a).
\]
The same supremum results if the laws are required to be strictly positive on a full binary product alphabet.
\end{theorem}
The final supremum need not be attained. Section \ref{R04-equality} determines exactly when a positive law attains equality.

Suppose additionally that $d_a(i)=\sum_{j:ij\in E}a_{ij}\leq1$ and put $b_i=1-d_a(i)$. The edge sets with weights $a_e$ and singleton sets with weights $b_i$ form a fractional partition. Its entropy gap is
\[
\Delta_a=\sum_{ij\in E}a_{ij}H(X_i,X_j)+\sum_i b_iH(X_i)-H(X_V)
=T(X_V)-\sum_ea_eI_e.
\]
\begin{corollary}\label{R04-stab}
The exact largest constant $\kappa$ for which $\Delta_a\geq\kappa T(X_V)$ holds for all finite joint laws is
\[
\kappa=1-\rho(G,a).
\]
In this fractional-partition regime $0\leq\rho\leq1$.
\end{corollary}
Indeed $a(E(S))\leq|S|/2\leq|S|-1$ for $|S|\geq2$. The bound concerns the full dependence $T$, not merely any one coordinate's mutual information with its complement.

\subsection{The needed forest-polytope fact}\label{R04-sec-2}
A forest may be disconnected and may have isolated vertices. Write $\mathbf1_F$ for its edge-incidence vector.
\begin{lemma}\label{R04-forestpoly}
A vector $x\in\R_{\geq0}^E$ is a convex combination of forest-incidence vectors if and only if
\begin{equation}\label{R04-forestconstr}
x(E(S))\leq|S|-1\quad\text{for every }S\subseteq V,\ |S|\geq2.
\end{equation}
\end{lemma}
\begin{proof}
Necessity follows since a forest restricted to $S$ has at most $|S|-1$ edges. To prove sufficiency, first note that for every edge set $A$, its nontrivial connected components have vertex sets $S_1,\ldots,S_k$, so
\[
x(A)\leq\sum_jx(E(S_j))\leq\sum_j(|S_j|-1)=r(A),
\]
where $r(A)=n-c(V,A)$ is the size of a maximum forest in $A$.

Consider any nonnegative edge costs $c_e$. Order edges $e_1,\ldots,e_m$ by decreasing cost, put $A_j=\{e_1,\ldots,e_j\}$ and $c_{e_{m+1}}=0$. A greedy forest $F$ is obtained by scanning the edges and retaining an edge exactly when it does not create a cycle. For every $j$, its retained edges in $A_j$ form a maximal forest in $A_j$, so $|F\cap A_j|=r(A_j)$. Abel summation gives
\[
c\cdot x=\sum_{j=1}^m(c_{e_j}-c_{e_{j+1}})x(A_j)
\leq\sum_{j=1}^m(c_{e_j}-c_{e_{j+1}})r(A_j)=c\cdot\mathbf1_F.
\]
This shows that no nonnegative linear functional separates $x$ from the convex hull of forests. Neither can a functional with negative entries: zero those entries, optimize by a forest avoiding them, and use $x\geq0$. Since a compact convex polytope and an exterior point admit a separating affine hyperplane, $x$ belongs to that convex hull. For completeness, separation follows by taking a closest point $y$ in the hull and the normal $x-y$; the squared-distance minimum implies $(x-y)\cdot(z-y)\leq0$ for every point $z$ in the hull.
\end{proof}
Applying the lemma to $x_e=a_e/\rho$ gives a finite probability distribution $\pi$ on forests such that
\begin{equation}\label{R04-decomp}
\sum_F\pi_F=1,\qquad\sum_{F\ni e}\pi_F=a_e/\rho.
\end{equation}
This is a certificate of the density upper bound. Conversely a maximizing set $S$ certifies its lower bound. Enumerating the finitely many forests and solving these linear equations is sufficient for exact rational verification when the weights are rational; no faster algorithm is assumed here.

\subsection{The entropy identity and sharpness}\label{R04-sec-3}
For each forest $F$, choose one root in each component and orient edges away from the roots. Define a probability $Q_F$ using the marginals of $P$, the law of $X_V$:
\[
Q_F(x)=\prod_{r\text{ root}}P_r(x_r)\prod_{i\text{ nonroot}}P(x_i\mid x_{\operatorname{pa}(i)}).
\]
Rows conditioned on a zero-probability parent may be filled arbitrarily. Summing leaf variables shows normalization, and $P\ll Q_F$ because every positive joint atom has positive relevant marginal and conditional probabilities.
\begin{lemma}\label{R04-kl}
\[
\D(P\Vert Q_F)=T(X_V)-\sum_{ij\in F}I_{ij}.
\]
\end{lemma}
\begin{proof}
Expanding $\E_P\log Q_F$ gives $-\sum_rH(X_r)-\sum_{i\text{ nonroot}}H(X_i\mid X_{\operatorname{pa}(i)})$. Subtract $H(X_V)$ with the correct sign and use $H(X_i\mid X_j)=H(X_i)-I_{ij}$.
\end{proof}
Averaging according to \eqref{R04-decomp} gives the exact remainder
\begin{equation}\label{R04-remainder}
\rho T(X_V)-\sum_ea_eI_e=\rho\sum_F\pi_F\D(P\Vert Q_F)\geq0.
\end{equation}
This proves \eqref{R04-mi}. To prove sharpness, choose a maximizing set $S$ and a nondegenerate binary variable $Z$. Set $X_i=Z$ for $i\in S$, and take variables outside $S$ independent of this block and each other. Then
\[
T(X_V)=(|S|-1)H(Z),\quad
I_{ij}=\begin{cases}H(Z),&i,j\in S,\\0,&\text{otherwise},\end{cases}
\]
so the ratio is $\rho$. For the positive-law supremum, mix this binary joint law with an arbitrarily small positive uniform product law. Finite-alphabet entropy is continuous and the limiting denominator is positive, so the ratios converge to $\rho$.

Equation \eqref{R04-remainder} also gives a posteriori structure: some forest in the support of $\pi$ satisfies
\[
\D(P\Vert Q_F)\leq T(X_V)-\rho^{-1}\sum_ea_eI_e.
\]
Pinsker's inequality then bounds $\TV(P,Q_F)$ by the square root of half this right side. This asserts proximity to at least one forest model, not to a uniquely identifiable graph.

\subsection{All strictly positive equality cases}\label{R04-sec-4}\label{R04-equality}
Let $J$ be the graph with edge set
\[
E(J)=\{e\in E:a_e=\rho\}.
\]
\begin{theorem}\label{R04-eqthm}
The graph $J$ is a forest. If $P$ is strictly positive on its product alphabet, equality in \eqref{R04-mi} holds if and only if $P=Q_J$, that is, $P$ is a forest-Markov law on $J$ with independent components.
In particular, if every edge satisfies $a_e<\rho$, the only positive equality laws are product laws.
\end{theorem}
\begin{proof}
A cycle of $k$ edges in $J$ would have total weight $k\rho$ on $k$ vertices, contradicting the definition of $\rho$; thus $J$ is a forest. Fix any decomposition \eqref{R04-decomp}. An edge appears in every forest with positive $\pi_F$ exactly when $a_e/\rho=1$. Hence the intersection of these forests is $J$.

If equality holds, \eqref{R04-remainder} and positivity imply $P=Q_F$ for every $F$ in the support. We show that a positive law that factors over several forests factors over their intersection. Fix a reference point $x^0$. The unique anchored expansion of $\log P$ is
\[
\log P(x)=\sum_{A\subseteq V}\phi_A(x_A),\qquad
\phi_A(x_A)=\sum_{B\subseteq A}(-1)^{|A|-|B|}
\log P(x_B,x^0_{V\setminus B}).
\]
M\"obius inversion verifies this identity directly. If $\log P$ is a sum of vertex and edge functions over a forest $F$, then $\phi_A=0$ unless $A$ is empty, a singleton, or an edge of $F$; this follows by cancellation for a function depending on fewer coordinates. Applying this to all the forests leaves only vertices and edges in $J$. Thus $P$ is a normalized product of positive vertex and edge factors on $J$. Repeated leaf summation, component by component, gives the rooted conditional factorization $P=Q_J$.

Conversely suppose $P=Q_J$. A positive-weight edge $e$ outside $J$ cannot join two vertices of the same component of $J$. Otherwise the unique $J$-path between its endpoints, with $k$ edges, together with $e$ gives weight at least $k\rho+a_e>k\rho$ on $k+1$ vertices, contradicting the density bound. Therefore the endpoints of such $e$ lie in independent components and $I_e=0$. Lemma \ref{R04-kl} for $J$ gives $T=\sum_{e\in J}I_e$. Multiplying by $\rho=a_e$ on $J$ proves equality.
\end{proof}
For laws with zeros, the exact equality criterion remains $P=Q_F$ for every positively weighted forest in one (equivalently any) decomposition. The clean intersection conclusion can fail without positivity; deterministic common-variable extremizers show why the assumption cannot simply be removed.

\subsection{Examples and interpretation}\label{R04-sec-5}
\begin{enumerate}
\item If the support graph is a forest, then $\rho=\max_ea_e$. The sharp coefficient is $1-\max_ea_e$ whenever singleton fillers make a fractional partition. Positive equality laws may retain dependence only on maximum-weight edges.
\item For a cycle on $n\geq3$ vertices with $a_e=1/2$, every proper induced subgraph is a forest and the full graph has density $n/(2(n-1))$. Hence
\[
\Delta_a\geq\frac{n-2}{2(n-1)}T.
\]
Nontrivial exact extremizers are necessarily nonpositive somewhere, since $1/2<\rho$.
\item For the complete graph with $a_e=1/(n-1)$, the density of a $k$-vertex set is $k/(2(n-1))$, maximized at $k=n$. The same sharp coefficient $(n-2)/(2(n-1))$ results, but the sparse and dense certificate decompositions are quite different.
\end{enumerate}
Generic separation-based stability bounds can be much smaller than this optimal graph constant. The value here is determined by all induced densities, not only the minimum separation weight of a pair of coordinates.

\subsection{Attribution and unresolved issues}\label{R04-sec-6}
The forest-polytope principle is standard matroid optimization; the proof above is supplied rather than attributed to this paper. Tree-reweighted entropy bounds are classical in Wainwright--Jaakkola--Willsky \cite{ref065}. General fractional-subadditivity stability and equality results appear in Jakhar--Kurri--Chillara--Prabhakaran \cite{ref066}. The present target is the exact scalar constant for graph-supported partitions plus the full positive-law equality classification. This formulation's priority remains unverified, and the manuscript does not claim a famous conjecture or an unchanged named open question. In particular the use of classical ingredients must be distinguished from an independent breakthrough.




\clearpage
\section[R05 --- The Exact Total-Correlation Remainder for Every Fractional Entropy Partition]{The Exact Total-Correlation Remainder for Every Fractional Entropy Partition}\label{contrib-R05}
\markboth{Entropy Inequalities, Marginal Compatibility, and Extremal Laws}{R05 --- Reiwa}
\noindent\textsc{R05 \quad Reiwa programme, source round 5}\par\medskip


\begin{quote}\small\textbf{Overview.} We determine the optimal coefficient of total correlation in every finite fractional-partition entropy inequality. It is the minimum, over subsets with at least two coordinates, of the total weight of sets meeting the subset minus one, divided by the subset size minus one. More generally, the exact universal ratio of weighted local total correlations to global total correlation is a finite hypergraph density. The proof uses a known normal-polymatroid minorant theorem, for which we provide a short constructive induction and explicit attribution. Binary common-variable models attain every sharp obstruction. An exact nonnegative remainder identifies critical subsets, and a minimum-cut construction verifies candidate constants without enumerating all subsets. The result is presented as a sharp corollary and synthesis of established polymatroid theory, with priority for the stated stability formula unverified.\end{quote}

\subsection{An exact universal constant}\label{R05-sec-1}
Let $V=[n]$, $n\geq2$, and let $\mathcal E$ be a finite indexed family of nonempty subsets, with weights $w_e\geq0$. For finite random variables define
\[
T(S)=\sum_{i\in S}H(X_i)-H(X_S),\qquad T(\emptyset)=0.
\]
No assumption on the joint law or alphabets is made. Define
\begin{equation}\label{R05-rho}
\rho(\mathcal E,w)=\max_{S\subseteq V,\ |S|\geq2}
\frac{\sum_e w_e(|e\cap S|-1)_+}{|S|-1}.
\end{equation}
\begin{theorem}\label{R05-universal}
For every joint law,
\begin{equation}\label{R05-tcineq}
\sum_e w_eT(e)\leq\rho(\mathcal E,w)T(V).
\end{equation}
The coefficient $\rho$ is exact, already over binary joint laws. The supremum remains $\rho$ if laws must be strictly positive. The same inequality holds when $H(X_S)$ is replaced by any normalized submodular set function.
\end{theorem}
A normalized function has value zero at the empty set. Its total-correlation transform is $T_h(S)=\sum_{i\in S}h(\{i\})-h(S)$. This transform is always nonnegative for submodular $h$, even when $h$ is not monotone.

Now assume the weights form a fractional partition:
\begin{equation}\label{R05-partition}
\sum_{e\ni i}w_e=1\qquad(i\in V).
\end{equation}
Let $W(S)=\sum_{e:e\cap S\ne\emptyset}w_e$ and define
\begin{equation}\label{R05-kappa}
\kappa(\mathcal E,w)=\min_{|S|\geq2}\frac{W(S)-1}{|S|-1}.
\end{equation}
\begin{corollary}\label{R05-sharp}
The largest universally valid coefficient in
\begin{equation}\label{R05-stability}
\sum_e w_eH(X_e)-H(X_V)\geq\kappa\left(\sum_iH(X_i)-H(X_V)\right)
\end{equation}
is exactly \eqref{R05-kappa}. In particular $\kappa=1-\rho$ and $0\leq\kappa\leq1$.
\end{corollary}
Indeed \eqref{R05-partition} gives $\sum_e w_eT(e)=T(V)-[\sum_e w_eH(X_e)-H(X_V)]$ and
\[
\sum_e w_e(|e\cap S|-1)_+=|S|-W(S).
\]
Also $1\leq W(S)\leq|S|$ for nonempty $S$. Thus the sharp remainder is a hypergraph invariant, not a distribution-dependent optimization over the entropy cone.

\subsection{A classical minorant theorem with a constructive proof}\label{R05-sec-2}
A \emph{polymatroid} is a normalized, monotone, submodular function $h:2^V\to\R_{\geq0}$. Call a function normal if it has the form
\begin{equation}\label{R05-normal}
h'(A)=\sum_{\emptyset\ne S\subseteq V}c_S\mathbf1\{A\cap S\ne\emptyset\},\qquad c_S\geq0.
\end{equation}
These functions are coverage functions and can be realized as entropy functions by independent common sources shared by the coordinates in $S$.
\begin{lemma}[Normal minorant; Abo Khamis--Kolaitis--Ngo--Suciu]\label{R05-minorant}
Every polymatroid $h$ has a normal minorant $h'$ such that
\[
h'(A)\leq h(A)\quad\text{for all }A,\qquad
h'(V)=h(V),\qquad h'(\{i\})=h(\{i\})\quad\text{for all }i.
\]
\end{lemma}
This is established prior art \cite[Lemma 3.10]{ref067}, not a discovery claimed here. The following induction gives a direct version adequate for the present application.
\begin{proof}
The one-coordinate case is immediate. Choose $v\in V$ and put $U=V\setminus\{v\}$. The contraction
$k_0(A)=h(A\cup\{v\})-h(\{v\})$ is a polymatroid on $U$. By induction let $k$ be its normal minorant preserving full-set and singleton values. For $i\in U$ set
\[
a_i=h(\{i\})+h(\{v\})-h(\{i,v\}),\qquad a_v=h(\{v\}).
\]
Submodularity and monotonicity give $0\leq a_i\leq a_v$. The function $m(A)=\max_{i\in A}a_i$, with $m(\emptyset)=0$, is normal: if $0=t_0<t_1<\cdots<t_r$ are its distinct positive levels, then
\[
m(A)=\sum_{j=1}^r(t_j-t_{j-1})\mathbf1\{A\cap\{i:a_i\geq t_j\}\ne\emptyset\}.
\]
Define $h'(A)=k(A\cap U)+m(A)$, which is normal. If $v\in A$, then $m(A)=h(\{v\})$, hence $h'(A)\leq h(A)$ by contraction. If $v\notin A$, submodularity implies that
\[
I_h(A;v):=h(A)+h(\{v\})-h(A\cup\{v\})\geq I_h(\{i\};v)=a_i
\]
for every $i\in A$. Therefore $m(A)\leq I_h(A;v)$ and
$h'(A)\leq k_0(A)+I_h(A;v)=h(A)$.
The preserved singleton values follow from $k(\{i\})=h(\{i,v\})-h(\{v\})$; the preserved full value follows from $k(U)=h(V)-h(\{v\})$.
\end{proof}
This construction has finite depth and produces explicit nonnegative coefficients $c_S$. Different choices of eliminated coordinates may produce different minorants, all valid.

\subsection{Proof of sharpness and a remainder decomposition}\label{R05-sec-3}
For a normal $h'$ as in \eqref{R05-normal}, direct calculation gives
\[
T_{h'}(A)=\sum_S c_S(|A\cap S|-1)_+,
\quad
T_{h'}(V)=\sum_S c_S(|S|-1).
\]
The minorant preserves the full and singleton values, so $T_{h'}(V)=T_h(V)$; and $h'\leq h$ gives $T_h(e)\leq T_{h'}(e)$. Consequently,
\[
\sum_e w_eT_h(e)\leq\sum_S c_S\sum_e w_e(|e\cap S|-1)_+
\leq\rho\sum_Sc_S(|S|-1)=\rho T_h(V).
\]
This proves Theorem \ref{R05-universal} for polymatroids. Any normalized submodular $h$ can be made monotone by adding $M|A|$ for a sufficiently large finite $M$; neither $T_h$ nor the balanced fractional-partition gap changes. The general submodular statement follows.

To attain the bound, choose a maximizing $S$ in \eqref{R05-rho}. Let $X_i=Z$ on $S$, where $Z$ is a fair bit, and let variables outside $S$ be deterministic. Then
\[
T(e)=(|e\cap S|-1)_+H(Z),\qquad T(V)=(|S|-1)H(Z).
\]
Thus the ratio is $\rho$. Mixing the resulting binary joint law with an arbitrarily small uniform product law gives strictly positive laws whose entropy ratios converge to $\rho$.

More precisely, for any chosen minorant the slack has the exact expansion
\begin{align}\label{R05-remainder}
\rho T_h(V)-\sum_e w_eT_h(e)
={}&\sum_e w_e\bigl(h(e)-h'(e)\bigr)\\
&+\sum_{|S|\geq2}c_S(|S|-1)\bigl(\rho-R(S)\bigr),\notag
\end{align}
where $R(S)$ denotes the ratio in \eqref{R05-rho}. Both terms are nonnegative.

\begin{corollary}[Critical-subset certificate]\label{R05-critical}
Let $\mathcal C=\{S:|S|\geq2,\ R(S)=\rho\}$. Equality in \eqref{R05-tcineq} holds exactly when, for any one valid normal minorant, both
\[
h(e)=h'(e)\text{ for every }e\text{ with }w_e>0,
\qquad c_S=0\text{ for }S\notin\mathcal C,\ |S|\geq2
\]
hold. If the slack is $\eta$ and
$\gamma=\min_{S\notin\mathcal C,\ |S|\geq2}(\rho-R(S))>0$, then
\[
\sum_{S\notin\mathcal C}c_S(|S|-1)\leq\eta/\gamma,
\qquad\sum_e w_e(h(e)-h'(e))\leq\eta.
\]
\end{corollary}
This is an entropy-vector certificate. It is not a claim that the original random variables are coupled to the minorant's independent sources, nor a distributional reconstruction theorem. If no noncritical subsets exist, the first quantitative bound is vacuous.

\subsection{When the sharp coefficient vanishes}\label{R05-sec-4}
\begin{proposition}\label{R05-separation}
For a fractional partition, $\kappa=0$ exactly when two distinct coordinates have identical incidence in all positive-weight sets. Otherwise $\kappa>0$.
\end{proposition}
\begin{proof}
If $i,j$ always occur together, $W(\{i,j\})=1$ and \eqref{R05-kappa} vanishes. Conversely, if $W(S)=1$ for some $|S|\geq2$, then
\[
|S|=\sum_e w_e|e\cap S|\leq|S|\sum_{e:e\cap S\ne\emptyset}w_e=|S|.
\]
Equality forces each positive-weight set meeting $S$ to contain all of $S$. Any two coordinates of $S$ therefore have identical incidence. Since the minimum in \eqref{R05-kappa} ranges over finitely many sets, failure of this event gives strict positivity.
\end{proof}
For a separating fractional partition this yields the quantitative distributional consequence
\[
\TV\left(P,\bigotimes_iP_i\right)
\leq\sqrt{\frac{\sum_e w_eH(X_e)-H(X_V)}{2\kappa}},
\]
since $T(V)=\D(P\Vert\bigotimes_iP_i)$ and Pinsker's inequality applies. The coefficient $\kappa$ in the preceding KL/entropy inequality is optimal; this does not assert that the displayed TV constant is separately optimal.

\subsection{Minimum-cut verification and weight design}\label{R05-sec-5}
The sharp coefficient appears to require $2^n$ subsets, but testing a proposed $q\geq0$ can be reduced to ordinary finite minimum cuts. The condition $q\leq\kappa$ is equivalent to
\begin{equation}\label{R05-test}
W(S)-q|S|\geq1-q\qquad\text{for all }|S|\geq2.
\end{equation}
For each pair $i\ne j$, minimize the left side over $S\supseteq\{i,j\}$. Construct a directed network with source $s$, sink $t$, coordinate nodes $v$, and factor nodes $e$. Add arcs $s\to v$ of capacity $q$, arcs $v\to e$ of capacity $M$ whenever $v\in e$, and arcs $e\to t$ of capacity $w_e$. Force $i,j$ to the source side by additional capacity-$M$ arcs from $s$. Choose
$M>nq+\sum_ew_e$, so no minimum cut uses an $M$ arc.

A finite cut selecting coordinates $S$ must select all factor nodes meeting $S$, and a minimum cut need select no others. Its capacity is
\[
q(n-|S|)+W(S)=nq+W(S)-q|S|.
\]
Thus checking all $\binom n2$ forced-pair minimum cuts verifies \eqref{R05-test}. Any failure gives an explicit violating set. Max-flow/min-cut algorithms therefore supply practical exact certificates for rational data. This network construction uses only coverage costs, not an oracle for the entropy cone.

For a prescribed family admitting a fractional partition, optimizing its stability is the linear program
\[
\max q\quad\text{subject to}\quad
w_e\geq0,\quad\sum_{e\ni i}w_e=1,\quad
\sum_{e:e\cap S\ne\emptyset}w_e\geq1+q(|S|-1)\ (|S|\geq2).
\]
It is bounded by $0\leq q\leq1$ and attains its optimum: each nonempty edge weight is at most one, so the feasible set can be restricted to a compact box. The same forced-pair cuts test all subset constraints for any candidate solution. This is a precise design problem with checkable optimality certificates, not a claim that a particular implementation has been run.

\subsection{Examples, attribution and limits}\label{R05-sec-6}
For the family of all $(n-1)$-subsets with weight $1/(n-1)$, every set $S$ of size at least two meets every edge. Thus $W(S)=n/(n-1)$ and
\[
\kappa=\frac1{(n-1)^2},
\]
attained by the common-bit model on all coordinates. For a graph with singleton fillers, \eqref{R05-rho} reduces to the induced-edge density $a(E(S))/(|S|-1)$, recovering the forest-arboricity result of the preceding round. For general hypergraphs the normal-minorant proof replaces the forest decomposition and covers arbitrary arities.

Fractional entropy inequalities are classical \cite{ref068}. The key normal-minorant theorem is due to \cite{ref067}; its role here must not be confused with a newly proved external lemma. General quantitative fractional-subadditivity results were recently studied in \cite{ref066}. The present submission extracts an explicit optimal constant, sharp witnesses, a nonnegative remainder, and cut certificates from these principles. A limited primary-source search has not certified priority for this combination. No named pre-existing open question is claimed solved, and no ranking over the tournament benchmark is asserted.




\clearpage
\section[R06 --- Local Entropy Stability Has a Different Hypergraph Spectrum]{Local Entropy Stability Has a Different Hypergraph Spectrum}\label{contrib-R06}
\markboth{Entropy Inequalities, Marginal Compatibility, and Extremal Laws}{R06 --- Reiwa}
\noindent\textsc{R06 \quad Reiwa programme, source round 6}\par\medskip


\begin{quote}\small\textbf{Overview.} For a product reference law and perturbations preserving all one-coordinate marginals, we diagonalize the weighted marginal-observation operator into interaction subspaces. Its exact chi-square contraction coefficient is the largest total weight of sets containing a pair of coordinates. The same number is the sharp infinitesimal contraction coefficient for every smooth strictly convex divergence. For relative entropy we give a uniform nonlinear bound under a two-sided density constraint. We then prove an arbitrarily large separation between local and global fractional-partition stability: for all coordinate-deletion projections the local coefficient is $1/(n-1)$, while the global sharp coefficient is $1/(n-1)^2$. Thus Hessian analysis alone cannot identify global near-extremizers. The underlying Hoeffding decomposition and local divergence geometry are classical; the explicit hypergraph comparison is the object of this synthesis.\end{quote}

\subsection{Product perturbations with fixed marginals}\label{R06-sec-1}
Let $\mu=\bigotimes_{i=1}^n\mu_i$ on finite alphabets, each with at least two points and full support. Let $\mathcal E$ be a family of coordinate sets with nonnegative weights $w_e$. For $A\subseteq[n]$, write $\mathsf E_Au=\E_\mu[u\mid X_A]$. Put
\[
\theta_S=\sum_{e\supseteq S}w_e,\qquad
\beta=\max_{i\ne j}\theta_{\{i,j\}}.
\]
Consider probabilities $P$ with one-coordinate marginals $P_i=\mu_i$ and density
$P/\mu=1+u$. Then $\E_\mu u=0$ and $\mathsf E_{\{i\}}u=0$ for all $i$. The density of $P_A$ relative to $\mu_A$ is $1+\mathsf E_Au$.

\begin{theorem}\label{R06-chi}
For every such probability,
\[
\sum_e w_e\chi^2(P_e\Vert\mu_e)\leq\beta\chi^2(P\Vert\mu),
\qquad
\chi^2(P\Vert\mu)=\E_\mu u^2.
\]
The coefficient $\beta$ is optimal, even among arbitrarily small nonzero positive perturbations of $\mu$.
\end{theorem}
Unlike a generic data-processing bound, this constant excludes first-order changes of the marginals. That exclusion is precisely why pair containment rather than single-coordinate containment appears.

\subsection{Exact interaction decomposition}\label{R06-sec-2}
For each $S\subseteq[n]$, define
\begin{equation}\label{R06-hoeff}
u_S=\sum_{A\subseteq S}(-1)^{|S|-|A|}\mathsf E_Au.
\end{equation}
The symbol $u_S$ denotes an interaction component, not the density of a marginal.
\begin{lemma}\label{R06-anova}
The components in \eqref{R06-hoeff} are mutually orthogonal in $L^2(\mu)$, and
\[
u=\sum_Su_S,\qquad \mathsf E_eu=\sum_{S\subseteq e}u_S.
\]
If $P_i=\mu_i$, the empty and singleton components vanish. Consequently
\begin{equation}\label{R06-spectrum}
\sum_e w_e\|\mathsf E_eu\|_2^2
=\sum_{|S|\geq2}\theta_S\|u_S\|_2^2.
\end{equation}
\end{lemma}
\begin{proof}
On a product space, $\mathsf E_A\mathsf E_B=\mathsf E_{A\cap B}$, as one sees by summing out coordinates. M\"obius inversion gives $u=\sum_Su_S$. Summing \eqref{R06-hoeff} over a coordinate in $S$ cancels pairs of terms, so $\mathsf E_{S\setminus\{i\}}u_S=0$ for every $i\in S$. If $S\ne T$, take $i\in S\setminus T$ (interchange the sets if needed). Since $u_T$ is independent of coordinate $i$, averaging the product over $i$ gives zero. Thus orthogonality holds. The same cancellation proves $\mathsf E_eu_S=u_S$ when $S\subseteq e$ and zero otherwise. Equation \eqref{R06-spectrum} follows by Pythagoras and summation. The singleton claim is immediate from the marginal constraints.
\end{proof}
Because weights are nonnegative, $\theta_S\leq\theta_{\{i,j\}}\leq\beta$ for any pair in $S$. Lemma \ref{R06-anova} proves Theorem \ref{R06-chi}. For sharpness choose a maximizing pair $i,j$ and nonzero mean-zero functions $g_i,g_j$. Then $u=\epsilon g_ig_j$ has only this pair interaction and $1+u>0$ for sufficiently small $\epsilon$. Equality holds in \eqref{R06-spectrum}.

\subsection{Uniform nonlinear relative-entropy control}\label{R06-sec-3}
Let $D_A=\D(P_A\Vert\mu_A)$ and $D=\D(P\Vert\mu)$, with natural logarithms.
\begin{theorem}\label{R06-nonlinear}
If $\|u\|_\infty\leq\eta<1$, then
\begin{equation}\label{R06-nonlinearbound}
\sum_e w_eD_e\leq\beta\frac{1+\eta}{1-\eta}\,D.
\end{equation}
If additionally $\sum_{e\ni i}w_e=1$ for all $i$, the entropy gap obeys
\begin{equation}\label{R06-gap}
\sum_e w_eH(P_e)-H(P)
\geq\left(1-\beta\frac{1+\eta}{1-\eta}\right)T(P),
\end{equation}
where $T(P)=\D(P\Vert\bigotimes_iP_i)=D$.
\end{theorem}
\begin{proof}
Set $\phi(t)=(1+t)\log(1+t)-t$. Since $\phi(0)=\phi'(0)=0$ and $\phi''(t)=1/(1+t)$, Taylor's formula with integral remainder gives, for $|t|\leq\eta$,
\[
\frac{t^2}{2(1+\eta)}\leq\phi(t)\leq\frac{t^2}{2(1-\eta)}.
\]
Conditional expectation contracts the supremum norm, so the same bounds apply to every $\mathsf E_eu$. Therefore
\[
\sum_e w_eD_e\leq\frac{\sum_ew_e\|\mathsf E_eu\|_2^2}{2(1-\eta)}
\leq\frac{\beta\|u\|_2^2}{2(1-\eta)}
\leq\beta\frac{1+\eta}{1-\eta}D.
\]
For a fractional partition, expanding $D$ and the $D_e$ against the product reference gives $\sum_e w_eH(P_e)-H(P)=D-\sum_ew_eD_e$. This proves \eqref{R06-gap}.
\end{proof}
The finite-$\eta$ constant is a safe uniform bound, not claimed optimal. It is informative when $\eta<(1-\beta)/(1+\beta)$. A more detailed second-order formula is
\[
D-\sum_ew_eD_e
=\frac12\sum_{|S|\geq2}(1-\theta_S)\|u_S\|_2^2+R,
\]
with, whenever $\beta\leq1$,
\[
|R|\leq\frac{\eta(1+\beta)}{6(1-\eta)^2}\|u\|_2^2.
\]
Indeed $|\phi^{(3)}(t)|\leq(1-\eta)^{-2}$, and the cubic remainder is bounded by $\eta\|u\|_2^2/[6(1-\eta)^2]$ for the full law and by the same expression with $\sum_ew_e\|\mathsf E_eu\|_2^2\leq\beta\|u\|_2^2$ for the weighted marginals.

\subsection{The exact local coefficient for smooth divergences}\label{R06-sec-4}
Let $f$ be twice continuously differentiable near one, with $f(1)=f'(1)=0$ and $f''(1)>0$. Define $D_f(P\Vert\mu)=\E_\mu f(P/\mu)$ whenever the density is in that neighborhood. This includes the usual smooth convex divergence generators after adding an irrelevant affine function.
\begin{theorem}\label{R06-local}
With fixed one-coordinate marginals,
\[
\lim_{\eta\downarrow0}\ \sup_{0<\|P/\mu-1\|_\infty\leq\eta}
\frac{\sum_e w_eD_f(P_e\Vert\mu_e)}{D_f(P\Vert\mu)}=\beta.
\]
For relative entropy and a fractional partition, the corresponding sharp local gap coefficient is $1-\beta$.
\end{theorem}
\begin{proof}
Uniform continuity of $f''$ implies numbers $m_\eta,M_\eta>0$ tending to $f''(1)$ such that
$\frac12m_\eta t^2\leq f(1+t)\leq\frac12M_\eta t^2$ for $|t|\leq\eta$. Lemma \ref{R06-anova} bounds every ratio by $\beta M_\eta/m_\eta$, proving the upper limit. For the pair perturbation $u=\epsilon g_ig_j$ from Theorem \ref{R06-chi}, a marginal density is $1+u$ on every edge containing the pair and is one otherwise. Thus its divergence ratio is exactly $\theta_{ij}=\beta$, for every sufficiently small nonzero $\epsilon$. This proves the matching lower bound without a limiting calculation.
\end{proof}
The ratio is unchanged if unused coordinates are added. The assumption that each alphabet has at least two points ensures that every candidate pair interaction is actually available; deterministic coordinates should instead be removed before maximizing.

\subsection{An unbounded local--global separation}\label{R06-sec-5}
Consider all $n$ coordinate-deletion sets $e_i=[n]\setminus\{i\}$ with weight $1/(n-1)$ and $n\geq3$. Every pair is contained in $n-2$ sets, so
\[
\beta=\frac{n-2}{n-1},\qquad \kappa_{\rm local}=\frac1{n-1}.
\]
By contrast the exact unrestricted coefficient is
\begin{equation}\label{R06-global}
\kappa_{\rm global}=\frac1{(n-1)^2}.
\end{equation}
Here is a self-contained verification of the lower bound in \eqref{R06-global}. For a submodular entropy function $h$, take a normal minorant
$h'(A)=\sum_{S\ne\emptyset}c_S\mathbf1\{A\cap S\ne\emptyset\}$ preserving singletons and the full set. Such a minorant exists by the following induction: contract one coordinate $v$, apply induction to the contraction, and add the coverage function $A\mapsto\max_{i\in A}a_i$, where $a_v=h(v)$ and $a_i=h(i)+h(v)-h(iv)$. Submodularity ensures this sum is below $h$ and preserves the stated values.
For a set $S$ of size $s\geq2$, its contribution to $T(V)$ is $(s-1)c_S$, while its contribution to the deletion entropy gap is $c_S/(n-1)$, because every deletion set meets $S$. Since $s-1\leq n-1$, the gap for $h'$ is at least $T(V)/(n-1)^2$. The gap for $h$ is no smaller. This proves the lower bound.

For sharpness take all $X_i$ equal to one fair bit. Then $T=(n-1)\log2$ and the deletion gap is $\log2/(n-1)$. Uniform binary one-coordinate marginals are preserved. Thus the local coefficient exceeds the global one by the unbounded factor $n-1$ even at a fixed reference with uniform binary marginals.

One can approach the global coefficient through strictly positive laws. Let $R$ put probability $1/2$ on each of the all-zero and all-one configurations, and let $P_t=(1-t)\mu+tR$, $0<t<1$, where $\mu$ is uniform product. Its $m$-coordinate marginal divergence is
\begin{align*}
D_m(t)={}&2^{1-m}\bigl[1+(2^{m-1}-1)t\bigr]
\log\bigl[1+(2^{m-1}-1)t\bigr]\\
&+(1-2^{1-m})(1-t)\log(1-t).
\end{align*}
The entropy-gap ratio is $1-\frac n{n-1}D_{n-1}(t)/D_n(t)$ and tends to \eqref{R06-global} as $t\uparrow1$. The density deviation is
$\|P_t/\mu-1\|_\infty=t(2^{n-1}-1)$, so this route does not remain in a small density neighborhood. There is no contradiction with Theorem \ref{R06-local}.

\subsection{Status and questions left open}\label{R06-sec-6}
Orthogonal product-space decompositions are classical Hoeffding analysis \cite{ref069}. The connection between local smooth divergences and chi-square geometry is well established \cite{ref070,ref071}. The normal-minorant argument used for the global comparison is due to \cite{ref067}. The contribution sought here is the explicit hypergraph spectrum under fixed marginals, a uniform nonlinear certificate, and the quantified failure of local analysis to determine global stability. Priority for this formulation remains unverified. We do not determine the exact optimal finite-$\eta$ KL constant, nor classify every globally extremizing positive law, nor claim a named open problem has been solved.




\clearpage
\section[R07 --- Sharp Same-Family Comparisons of Symmetric Fractional Entropy Gaps]{Sharp Same-Family Comparisons of Symmetric Fractional Entropy Gaps}\label{contrib-R07}
\markboth{Entropy Inequalities, Marginal Compatibility, and Extremal Laws}{R07 --- Reiwa}
\noindent\textsc{R07 \quad Reiwa programme, source round 7}\par\medskip


\begin{quote}\small\textbf{Overview.} For every permutation-symmetric fractional partition, we determine the exact two-sided comparison of its upper and lower entropy gaps. The optimal constants are a finite rank spectrum, are reciprocal, and are attained by explicit finite-field distributions. We also determine the complete closed region of the two gaps after normalization by total correlation, and give an exact slack decomposition identifying the ranks responsible for near equality. Mixed subset sizes can have strictly interior extremizing ranks even when the total fractional weight is two. These statements are applications of the classical symmetric polymatroid cone, not a new characterization of that cone or a solution of the unrestricted entropy-region problem.\end{quote}

\subsection{The comparison problem and its scope}\label{R07-sec-1}
Let $n\ge2$, and let $\alpha_1,\ldots,\alpha_{n-1}\ge0$ satisfy
\begin{equation}\label{R07-cover}\sum_{k=1}^{n-1}k\alpha_k=n.\end{equation}
Give each $k$-element subset weight $\alpha_k/\binom nk$. Thus the sum of weights on subsets containing any fixed coordinate is one. Empty and full subsets are omitted; full subsets contribute zero to both gaps and can be removed by a common rescaling when the remaining incidence is nonzero. Write $W=\sum_k\alpha_k>1$.

For any normalized submodular function $h$ on $[n]$, put
\[
a_k=\binom nk^{-1}\sum_{|A|=k}h(A),\quad a_0=0,
\]
and define
\begin{align}
 U(h)&=\sum_k\alpha_k a_k-a_n,\label{R07-upper}\\
 L(h)&=\sum_k\alpha_k a_{n-k}-(W-1)a_n,\label{R07-lower}\\
 T(h)&=na_1-a_n.
\end{align}
For $h(A)=H(X_A)$, these are the upper fractional-partition gap, the lower gap, and total correlation, respectively. In particular $L=h([n])-\sum_Aw_A h(A\mid A^c)$.

Madiman--Tetali's complement duality relates gaps for complementary families, but does not by itself compare the two gaps of the same family \cite{ref072}. We give the complete sharp answer for permutation-symmetric families. The symmetric entropy cone and its uniform-matroid generators are classical; see Han's work as discussed in \cite{ref073}. Our contribution here is an explicit comparison and normalized-region formulation, with finite witnesses and slack certificates. Priority of this exact formulation has not been established. The question for arbitrary nonsymmetric families is not answered.

\subsection{The rank spectrum}\label{R07-sec-2}
For $1\le r\le n-1$, define
\begin{equation}\label{R07-ur}
 u_r=\sum_{k=1}^{n-1}\alpha_k\min(k,r)-r.
\end{equation}
These numbers are strictly positive. Indeed, using \eqref{R07-cover},
\[
u_r=\sum_k\alpha_k\left(\min(k,r)-\frac{kr}{n}\right)>0,
\]
since each summand with $\alpha_k>0$ is positive. In particular $u_1=W-1$ and $u_{n-1}=1$.

\begin{theorem}[Exact same-family constants]\label{R07-main}
Set
\[
 K=\max_{1\le r\le n-1}\frac{u_r}{u_{n-r}}\ge1.
\]
Every normalized submodular function satisfies
\begin{equation}\label{R07-comp} K^{-1}L(h)\le U(h)\le K L(h).\end{equation}
Both constants are optimal even if one restricts to entropy functions of finite-valued random variables. Every extremal rank is attained over any finite field with at least $n$ elements. If $T(h)=0$, both gaps vanish.
\end{theorem}

We include the elementary symmetric-cone argument, so that the result does not depend on an external polyhedral description.

\begin{lemma}[Averaged concavity and rank expansion]\label{R07-conc}
The increments $d_k=a_k-a_{k-1}$ satisfy $d_k\ge d_{k+1}$. Consequently, with $\delta_r=d_r-d_{r+1}\ge0$,
\begin{equation}\label{R07-rankexp}
 a_k=d_n k+\sum_{r=1}^{n-1}\delta_r\min(k,r).
\end{equation}
Moreover,
\begin{equation}\label{R07-rankgap}
 T=\sum_r(n-r)\delta_r,\qquad
 U=\sum_r\delta_r u_r,\qquad L=\sum_r\delta_r u_{n-r}.
\end{equation}
\end{lemma}
\begin{proof}
Choose uniformly a set $S$ of size $k-1$ and distinct $i,j\notin S$. Diminishing returns gives
\[
h(S+i)-h(S)\ge h(S+i+j)-h(S+j).
\]
On averaging, the two sides are $d_k$ and $d_{k+1}$. The expansion follows by summing $d_k=d_n+\sum_{r\ge k}\delta_r$. Modular terms cancel in $U,L,T$ because of \eqref{R07-cover}. The rank-$r$ value of $U$ is \eqref{R07-ur}. Its lower-gap value equals $u_{n-r}$: use
\[
\min(n-k,r)=r-k+\min(k,n-r)
\]
and \eqref{R07-cover}. The total-correlation value is $n-r$.
\end{proof}

\begin{proof}[Proof of Theorem \ref{R07-main}]
Equation \eqref{R07-rankgap} writes $U/L$, whenever $L>0$, as a weighted average of $u_r/u_{n-r}$ with weights $\delta_r u_{n-r}$. The spectrum is invariant under taking reciprocals, by replacing $r$ with $n-r$. This proves \eqref{R07-comp}. If $T=0$, every $\delta_r=0$.

For sharpness, choose distinct $t_1,\ldots,t_n$ in a finite field $\mathbb F_q$, and choose $Z_0,\ldots,Z_{r-1}$ independently and uniformly in that field. Set
\begin{equation}\label{R07-MDS}
 X_i=\sum_{j=0}^{r-1}t_i^jZ_j.
\end{equation}
Every collection of at most $r$ of these linear forms has full row rank, by the Vandermonde determinant; every larger collection has rank $r$. Hence
\[
 H(X_A)=\min(|A|,r)\log q.
\]
Its gap ratio is exactly $u_r/u_{n-r}$. Finite fields with $q\ge n$ exist, for example fields of sufficiently large prime order.
\end{proof}

\subsection{The complete normalized region and stability}\label{R07-sec-3}
\begin{theorem}[Two-gap region]\label{R07-region}
Let $\mathcal R$ be the set of pairs $(U/T,L/T)$ obtained from finite-valued distributions with $T>0$. Then
\begin{equation}\label{R07-hull}
 \overline{\mathcal R}
 =\conv\left\{\left(\frac{u_r}{n-r},\frac{u_{n-r}}{n-r}\right):1\le r\le n-1\right\}.
\end{equation}
For arbitrary normalized submodular functions with $T>0$, the same convex hull is the exact region, without closure. If $R_* =\{r:u_r=K u_{n-r}\}$, the upper-comparison slack is
\begin{equation}\label{R07-slack}
 KL-U=\sum_r\delta_r(Ku_{n-r}-u_r).
\end{equation}
Thus equality holds precisely when $\delta_r=0$ outside $R_*$. If this set has nonempty complement and
\[
 \gamma=\min_{r\notin R_*}\frac{Ku_{n-r}-u_r}{n-r}>0,
\]
then
\begin{equation}\label{R07-stability}
 \sum_{r\notin R_*}\frac{(n-r)\delta_r}{T}
 \le\frac{KL-U}{\gamma T}.
\end{equation}
\end{theorem}
\begin{proof}
Use the weights $(n-r)\delta_r/T$ in \eqref{R07-rankgap} for the inclusion in \eqref{R07-hull}. Any point of the hull is realized by a nonnegative linear combination of the rank functions $\min(|A|,r)$, which is submodular.

For the entropic density assertion, fix a single field $\mathbb F_q$ with $q\ge n$. For each rank $r$, take $m_r$ independent copies of construction \eqref{R07-MDS}; at coordinate $i$ concatenate all its copies. Independence makes entropy vectors add. The normalized coefficients are
\[
 \frac{m_r(n-r)}{\sum_s m_s(n-s)}.
\]
As the nonnegative integers $m_r$ vary, these coefficients are dense in the simplex: approximate any desired coefficient vector by taking $m_r$ proportional to its $r$th component divided by $n-r$ and rounding at a scale tending to infinity. This proves the reverse inclusion after closure. Equation \eqref{R07-slack} follows from \eqref{R07-rankgap}; positivity yields the equality statement and \eqref{R07-stability}.
\end{proof}

The stability assertion concerns the averaged entropy profile. It does not assert that the original distribution is close in total variation to a linear code or a product of code distributions. Such a probabilistic conclusion does not follow from this proof.

\subsection{Examples and the failure of endpoint-only tests}\label{R07-sec-4}
\subsubsection{A single subset size}
If $\alpha_k=n/k$ and all other weights vanish, direct evaluation gives
\[
 u_r=\begin{cases}(n-k)r/k,&r\le k,\\n-r,&r\ge k.\end{cases}
\]
Consequently
\begin{equation}\label{R07-pure}
 K=\max\left\{\frac{k}{n-k},\frac{n-k}{k}\right\}.
\end{equation}
For completeness, suppose $k\le n/2$. On $r\le k$ the ratio is $(n-k)/k$; on $k\le r\le n-k$ it is $(n-r)/r$; on $r\ge n-k$ it is $k/(n-k)$. The other case follows by complementation. Thus rank-one common randomness and rank-$(n-1)$ parity laws already witness the endpoints in this special case.

\subsubsection{A mixture whose extremizers have interior ranks}
Take $n=6$ and
\[
 \alpha_1=\frac25,\qquad \alpha_2=\frac25,\qquad \alpha_4=\frac65,
\]
with all other weights zero. Then $\sum k\alpha_k=6$ and $W=2$, but
\[
 (u_1,u_2,u_3,u_4,u_5)=\left(1,\frac85,\frac95,2,1\right).
\]
The ratios $u_r/u_{6-r}$ are $1,4/5,1,5/4,1$. Thus the exact comparison is
\[
 \frac45 L\le U\le\frac54L,
\]
with ranks two and four as the extremizers. Checking only common randomness and a single parity relation would incorrectly suggest equality of the two gaps. In particular $W=2$ alone does not imply $U=L$.

\subsubsection{Exactly when are the two gaps identical?}
\begin{proposition}\label{R07-selfdual}
The identity $U(h)=L(h)$ for every finite entropy function holds if and only if
\begin{equation}\label{R07-sym}\alpha_k=\alpha_{n-k}\quad(1\le k\le n-1).\end{equation}
In that case $W=2$ and the fractional partition is its own complement dual.
\end{proposition}
\begin{proof}
Necessity of $u_r=u_{n-r}$ follows by testing \eqref{R07-MDS}. Set $u_0=u_n=0$. The second difference satisfies
\[
 2u_r-u_{r-1}-u_{r+1}=\alpha_r.
\]
Reflection symmetry of the sequence therefore gives \eqref{R07-sym}. Conversely \eqref{R07-sym} and \eqref{R07-cover} imply $nW=2\sum k\alpha_k=2n$, so $W=2$; equations \eqref{R07-upper}--\eqref{R07-lower} then coincide.
\end{proof}

\subsection{Computation, limits, and provenance}\label{R07-sec-5}
The $u_r$ can be computed in $O(n)$ arithmetic operations using prefix sums of $\alpha_k$ and $k\alpha_k$. The exact comparison constant then requires $n-1$ ratios. For rational input weights all constants, extremizing ranks, and slack coefficients are rational and exactly checkable. A hull algorithm on the $n-1$ displayed planar points gives the normalized region.

The sharpness examples use alphabet sizes that can grow with $n$; optimal constants with a prescribed binary alphabet are not determined here. Neither the symmetric-cone argument nor the Vandermonde realization is claimed new. This round changes the target from the local/global total-correlation remainders of preceding rounds to a complete two-functional comparison and its attainable normalized region. No result of another competitor is used. In particular no comparison with the benchmark manuscripts' mathematical importance is asserted.





\clearpage
\section[R08 --- Entropic Frustration on Signed Graphs: Exact Cactus Solutions and a Complexity Barrier]{Entropic Frustration on Signed Graphs: Exact Cactus Solutions and a Complexity Barrier}\label{contrib-R08}
\markboth{Entropy Inequalities, Marginal Compatibility, and Extremal Laws}{R08 --- Reiwa}
\noindent\textsc{R08 \quad Reiwa programme, source round 8}\par\medskip


\begin{quote}\small\textbf{Overview.} We study the least sum of edge relative entropies required to reconcile noisy binary parity specifications with a global law. Global flip symmetry identifies the problem exactly with separable Bernoulli relative entropy on a translated cut polytope. On every cactus graph, arbitrary positive edge weights and arbitrary positive noise levels reduce to independent scalar equations, with explicit global realizing laws. For general graphs, a constant-additive approximation is NP-hard when the noise is part of the rational input. A separate low-temperature theorem identifies the selected point on the minimum-frustration face and gives an explicit exponential bound on the value remainder. The general entropic-penalty phenomenon and the cut-polytope connection are classical; the unified formulas and their precise constants are the claims developed here.\end{quote}

\subsection{The global reconciliation problem}\label{R08-sec-1}
Let $G=(V,E)$ be a finite simple graph with $m\ge1$ edges. On each edge $e=ij$, prescribe a parity $c_e\in\{0,1\}$, a noise $a_e\in(0,1)$, and weight $w_e>0$. The target law on two bits is
\[
 \nu_e(x_i,x_j)=\begin{cases}(1-a_e)/2,&x_i\oplus x_j=c_e,\\a_e/2,&x_i\oplus x_j\ne c_e.\end{cases}
\]
The symbol $\oplus$ denotes addition modulo two. Define
\begin{equation}\label{R08-C}
 C(G,c,a,w)=\min_{P\text{ on }\{0,1\}^V}\sum_{e\in E}w_e\D(P_e\Vert\nu_e).
\end{equation}
Write $d(p\Vert a)=p\log(p/a)+(1-p)\log((1-p)/(1-a))$, with continuous endpoint values, and set
\[
 \mathcal Y=\{(x_i\oplus x_j\oplus c_e)_{e\in E}:x\in\{0,1\}^V\},\qquad K=\conv\mathcal Y.
\]
Thus $K$ is a coordinate reflection of the cut polytope.

\begin{theorem}[Exact reduction]\label{R08-reduction}
For all the preceding data,
\begin{equation}\label{R08-cutopt}
 C(G,c,a,w)=\min_{p\in K}\sum_e w_e d(p_e\Vert a_e).
\end{equation}
The minimizer $p$ is unique. Every feasible $p$ has a realizing global law with uniform one-bit marginals, for which all edge divergences equal $d(p_e\Vert a_e)$.
\end{theorem}
\begin{proof}
Map each pair $(X_i,X_j)$ to its violation bit $Y_e$. Grouping the terms in relative entropy by the value of $Y_e$ proves $\D(P_e\Vert\nu_e)\ge d(P(Y_e=1)\Vert a_e)$. Also $\E Y\in K$.

Conversely, express $p$ as the expectation of a random vector $Y\in\mathcal Y$. Choose for each such $Y$ one vertex assignment producing it. Independently flip all bits in each connected component with a fair bit. This leaves $Y$ unchanged and makes the two configurations of each edge parity equally likely. The edge law then has precisely the target's uniform conditional law given parity, and equality holds in the divergence bound. Strict convexity of the sum of Bernoulli divergences gives uniqueness of $p$.
\end{proof}

The polyhedral connection between binary marginals and cuts is established background, not a new observation; see \cite{ref074}. The point here is the exact solvable-versus-hard boundary for this particular objective and a quantitative low-noise analysis. The preceding tournament's single-cycle calculation is fully reproved below, rather than used as an external dependency.

\subsection{Cactus graphs: independent parity blocks}\label{R08-sec-2}
A cactus is a graph in which each edge belongs to at most one simple cycle. Let its cycles be $B_1,\ldots,B_t$, and let $B_0$ be its bridge edges. Put $s_j=\bigoplus_{e\in B_j}c_e$ and
\[
 P_s^{\ell}=\conv\{v\in\{0,1\}^{\ell}:\textstyle\bigoplus_i v_i=s\}.
\]
\begin{lemma}\label{R08-product}
For a cactus,
\[
 K=[0,1]^{B_0}\times\prod_{j=1}^t P_{s_j}^{|B_j|}.
\]
\end{lemma}
\begin{proof}
Any violation vector has parity $s_j$ around cycle $B_j$, since vertex bits cancel in pairs. Conversely, select edge bits meeting these cycle conditions. Remove one edge from each cycle to obtain a spanning forest. Choose a root bit in each component and propagate along that forest using $x_i\oplus x_j=Y_e\oplus c_e$. The removed-edge equations hold by their cycle conditions. Hence the set of allowed binary vectors itself is the product of the stated blocks; taking convex hulls proves the result.
\end{proof}

\begin{lemma}[A parity block]\label{R08-parity}
For $\ell\ge3$, the polytope $P_s^{\ell}$ consists exactly of $p\in[0,1]^\ell$ satisfying
\begin{equation}\label{R08-facet}
 \ell_u(p):=\sum_{u_i=0}p_i+\sum_{u_i=1}(1-p_i)\ge1
 \quad\text{for every }u\text{ with parity }1-s.
\end{equation}
At most one such inequality is violated at any point of the cube.
\end{lemma}
\begin{proof}
Every allowed vertex differs from a forbidden $u$ in at least one coordinate. Distinct forbidden vertices have Hamming distance at least two, while $\ell_u(p)+\ell_v(p)\ge\|u-v\|_1$. Thus two strict violations are impossible.

To show sufficiency, consider an extreme point of the indicated intersection and let $F$ be its set of fractional coordinates. If $|F|=1$, choose a forbidden vertex agreeing on all integral coordinates; then its distance is strictly less than one, impossible. If $|F|\ge2$ and two distinct forbidden inequalities are active, their vertices must differ in exactly two coordinates and $p$ must agree with them on every other coordinate. Thus $|F|=2$. On those two coordinates their active normals are opposites. If more active inequalities exist, the same argument shows their restrictions have rank at most one. Consequently, with $|F|\ge2$, there is a nonzero direction supported on $F$ annihilating all active inequalities (also when zero or one is active). A sufficiently small move in both directions preserves every inactive inequality and cube bound, contradicting extremality. All extreme points are therefore allowed binary vertices.
\end{proof}

\begin{theorem}[Explicit cactus projection]\label{R08-cactus}
Bridge coordinates incur zero cost and have $p_e=a_e$. On a cycle $B$ of parity $s$, if $a_B\in P_s^{|B|}$ the cost is zero and $p_B=a_B$. Otherwise let $u$ be its unique violated forbidden vertex and set
\[
 b_e=\begin{cases}a_e,&u_e=0,\\1-a_e,&u_e=1.\end{cases}
 \qquad\left(\sum_{e\in B}b_e<1\right).
\]
There is a unique $\lambda>0$ solving
\begin{equation}\label{R08-root}
 \sum_{e\in B}\frac{b_e\exp(\lambda/w_e)}{1-b_e+b_e\exp(\lambda/w_e)}=1.
\end{equation}
Let $r_e$ be its summands and reflect back: $p_e=r_e$ if $u_e=0$, and $p_e=1-r_e$ otherwise. The cycle cost is
\begin{equation}\label{R08-cost}
 C_B=\lambda-\sum_{e\in B}w_e\log(1-b_e+b_e\exp(\lambda/w_e)),
\end{equation}
and the total cost is $\sum_B C_B$.
\end{theorem}
\begin{proof}
Lemmas \ref{R08-product}--\ref{R08-parity} separate the objective into blocks. In an infeasible block, every feasible point satisfies $\ell_u(p)\ge1$. In reflected coordinates, minimize the strictly convex objective $\sum w_e d(r_e\Vert b_e)$ over $\sum r_e\ge1$. The unconstrained minimum $b$ has sum less than one, so the constrained minimum has sum one. Lagrange differentiation gives \eqref{R08-root}; its left side increases continuously from $\sum b_e<1$ to $|B|>1$. Each $r_e\in(0,1)$.

The resulting point is a convex combination of allowed neighbors $u\oplus\one_e$, with coefficients $r_e$. It therefore belongs to the full parity polytope, not just the relaxed half-space. Substitution of the logistic formula gives \eqref{R08-cost}.
\end{proof}

The violated vertex can be found in linear time per cycle: first choose each bit of $u$ nearest to $a_e$; if its parity is wrong, flip a coordinate minimizing $|1-2a_e|$. Test whether its distance is below one. Equation \eqref{R08-root} admits bisection; a finite upper bracket is
\[
 1+\max_{e\in B}\left[w_e\log\frac{1-b_e}{b_e}\right]_+,
\]
since then each summand exceeds $1/2$. These are real-arithmetic formulas, not a claim of constant-cost transcendental evaluation. A realizing global law is obtained by independent block sampling and the forest lift in Lemma \ref{R08-product}; in an active block the neighbor mixture is explicit. Feasible blocks can be realized by any convex decomposition, whose existence was proved above.

\subsection{General graphs encode a hard integer problem}\label{R08-sec-3}
Now set all $w_e=1$ and all noises $a_e=(1+q)^{-1}$, where $q>1$. Put $\beta=\log q$, $A_q=\log(1+q^{-1})$, and $h_2(t)=-t\log t-(1-t)\log(1-t)$. Then
\begin{equation}\label{R08-free}
 C_q=mA_q+\min_{p\in K}\left[\beta\sum_ep_e-\sum_eh_2(p_e)\right].
\end{equation}
Let $\tau=\min_{y\in\mathcal Y}\sum_e y_e$ be the frustration index.
\begin{proposition}\label{R08-rough}
For every graph and $q>1$,
\[
 \tau\log q-m\log2+mA_q\le C_q\le\tau\log q+mA_q.
\]
In particular, approximating $C_q$ to absolute error at most one for rational inputs $(G,c,q)$ is NP-hard.
\end{proposition}
\begin{proof}
The lower bound follows from $\sum p_e\ge\tau$ and $\sum h_2(p_e)\le m\log2$; use a minimum-frustration binary vector for the upper bound.

For the hardness assertion, take any instance of unweighted MAX-CUT with $m\ge1$, set every $c_e=1$, and choose the integer $q=2^{16(m+1)}$. Then $\tau=m-\operatorname{MaxCut}(G)$. The binary encoding of $q$ has $O(m)$ bits. If $|\widetilde C-C_q|\le1$, the preceding bounds imply
\[
 \left|\frac{\widetilde C}{\log q}-\tau\right|
 \le\frac{m\log2+1+m/q}{16(m+1)\log2}<\frac14.
\]
Approximating the logarithm to sufficient polynomially many bits and rounding therefore recovers the integer $\tau$. This is a polynomial-time reduction from unweighted MAX-CUT, whose NP-completeness is classical \cite{ref075}. No claim that the real-valued approximation problem itself lies in NP is needed.
\end{proof}

The input dependence of $q$ is essential to this proof. It does not establish hardness for one fixed noise bounded away from zero, nor does it rule out algorithms on other structured graph classes. It also does not contradict the cactus formulas.

\subsection{The low-temperature limit selects a marginal-entropy center}\label{R08-sec-4}
The logarithmic leading term in Proposition \ref{R08-rough} does not describe which minimum-frustration mixtures are selected. Define the exposed face
\[
 K_0=\{p\in K:\textstyle\sum_ep_e=\tau\},\quad
 H(p)=\sum_eh_2(p_e),\quad H_* =\max_{p\in K_0}H(p).
\]
Strict concavity gives a unique maximizer $p_*\in K_0$. This is a center for the sum of coordinate entropies, not in general the edge-marginal vector of the uniform distribution over optimal vertex assignments.

\begin{theorem}[Exponential value remainder and selection]\label{R08-temperature}
For every $\beta=\log q>0$,
\begin{equation}\label{R08-exp}
 0\le \beta\tau-H_*+mA_q-C_q
 \le m\log(1+2e^{-\beta/m}).
\end{equation}
The unique optimizer $p_\beta$ in \eqref{R08-free} converges to $p_*$ as $\beta\to\infty$. More quantitatively, take any convex decomposition of $p_\beta$ into vectors of $\mathcal Y$, and let $t$ be its total coefficient on vectors of weight greater than $\tau$. Then
\begin{equation}\label{R08-tbound}t\le\min\{1,2e\,e^{-\beta/m}\}.
\end{equation}
If $t<1$, let $p_0\in K_0$ be the conditional mean of its optimal vectors. Then
\begin{equation}\label{R08-dist}
 \|p_\beta-p_*\|_2\le\sqrt m\,t+
 \sqrt{\frac{m\{t\log2+h_2(t)\}}{2(1-t)}}.
\end{equation}
When $t=0$, interpret the right side as zero.
\end{theorem}
\begin{proof}
Write $p=(1-t)p_0+tp_1$ by separating optimal and nonoptimal vertices; if one group is empty its mean may be chosen arbitrarily. Integrality gives $\sum_ep_e\ge\tau+t$. Introduce a Bernoulli selector with probability $t$. Applying the entropy chain rule to each coordinate of this mixture gives
\[
 H(p)\le(1-t)H(p_0)+tH(p_1)+mh_2(t)
 \le H_*+mt\log2+mh_2(t).
\]
Therefore
\[
 \beta\sum_ep_e-H(p)\ge\beta\tau-H_*+
 (\beta-m\log2)t-mh_2(t).
\]
The minimum of the final two terms over $0\le t\le1$ is
$-m\log(1+2e^{-\beta/m})$; this follows by differentiation, including endpoint continuity. Testing $p_*$ supplies the opposite bound, proving \eqref{R08-exp}.

At the optimizer, comparison with $p_*$ also yields
$(\beta-m\log2)t\le mh_2(t)$. If $t>0$, use $h_2(t)\le t\log(e/t)$ to get \eqref{R08-tbound}. Further, optimality implies $H(p_\beta)\ge H_*$. The sharper mixture bound therefore gives
\[
 H_*-H(p_0)\le\frac{m\{t\log2+h_2(t)\}}{1-t}.
\]
The function $H$ is $4$-strongly concave in the Euclidean norm, since $-h_2^{\prime\prime}(x)=1/[x(1-x)]\ge4$. Hence maximality of $p_*$ on the convex set $K_0$ implies
$H_*-H(p_0)\ge2\|p_0-p_*\|_2^2$. This remains true on the boundary by continuity; alternatively apply strong concavity to $(1-s)p_*+sp_0$ and let $s\downarrow0$. Finally $\|p_\beta-p_0\|_2\le\sqrt m\,t$, proving \eqref{R08-dist}. Since $t\to0$, convergence follows.
\end{proof}

\subsection{Research boundary and provenance}\label{R08-sec-5}
Exponential convergence of entropic penalization for linear programs is known, including explicit bounds \cite{ref076}. Theorem \ref{R08-temperature} is a direct specialized derivation for binary marginal entropy; neither exponential convergence itself nor the general penalty method is claimed new. The cactus theorem gives an exact construction for heterogeneous data, whereas the hardness reduction explains why arbitrary graphs should not be expected to admit the same simple extension. Novelty of these precise combined formulations is unverified. The work does not solve a named pre-existing open problem. No peer competitor's result is used, and no formal verification is claimed.





\clearpage
\section[R09 --- Common Blocks and Quantitative Intersection in Pairwise Entropy Extremizers]{Common Blocks and Quantitative Intersection in Pairwise Entropy Extremizers}\label{contrib-R09}
\markboth{Entropy Inequalities, Marginal Compatibility, and Extremal Laws}{R09 --- Reiwa}
\noindent\textsc{R09 \quad Reiwa programme, source round 9}\par\medskip


\begin{quote}\small\textbf{Overview.} We give a complete finite-alphabet classification of equality in the sharp bound $\sum_{i<j}I(X_i;X_j)\le(n/2)\operatorname{TC}(X_1,\ldots,X_n)$, including singular laws. Equality means precisely that the variables are independent conditional on one label that is a deterministic function of every coordinate; equivalently, the joint law is a disjoint mixture of product blocks. We then prove a quantitative intersection lemma and use it to obtain a global strict remainder from bounds on pairwise likelihood ratios. No positivity assumption on the full joint law is needed. The equality argument uses the classical double-Markov/common-part mechanism, and the quantitative argument is an elementary entropy-comparison application; the displayed remainder is not asserted optimal.\end{quote}

\subsection{The deficit and its exact certificates}\label{R09-sec-1}
Let $X_1,\ldots,X_n$ be finite-valued random variables, $n\ge3$. All logarithms are natural. Set
\[
 T=\sum_{i=1}^nH(X_i)-H(X_1,\ldots,X_n),\qquad
 S=\frac n2T-\sum_{i<j}I(X_i;X_j).
\]
For each center $i$, define the star law
\[
 Q_i(x)=P_i(x_i)\prod_{j\ne i}P(X_j=x_j\mid X_i=x_i).
\]
Conditional laws at zero-mass values can be defined arbitrarily. The resulting $Q_i$ is a probability law and is positive wherever $P$ is positive.

\begin{lemma}[Star certificates]\label{R09-stars}
The nonnegative quantities
\begin{equation}\label{R09-Di}
 D_i=\D(P\Vert Q_i)=T-\sum_{j\ne i}I(X_i;X_j)
 =\sum_{j\ne i}H(X_j\mid X_i)-H(X_{-i}\mid X_i)
\end{equation}
satisfy $\sum_iD_i=2S$. Thus $S\ge0$. Equality holds if and only if, for every $i$, the variables other than $X_i$ are mutually independent conditionally on $X_i$.
\end{lemma}
\begin{proof}
Expand the logarithm of $Q_i$ under $P$ to obtain \eqref{R09-Di}. Relative entropy is nonnegative and vanishes exactly when its two laws agree. Summing counts each pairwise mutual information twice.
\end{proof}

The inequality itself is classical and follows, for example, by averaging tree entropy inequalities. The factor $n/2$ is sharp, as is already witnessed by $X_1=\cdots=X_n$ equal to one nonconstant bit. The issue here is the full singular equality structure and a quantitative nondegeneracy refinement, not a new discovery of that factor.

\subsection{All finite equality cases}\label{R09-sec-2}
\begin{theorem}[Disjoint product blocks]\label{R09-blocks}
The following conditions are equivalent:
\begin{enumerate}[label=(\roman*)]
\item $S=0$.
\item There is a finite random variable $Z$ such that $H(Z\mid X_i)=0$ for every $i$, and $X_1,\ldots,X_n$ are mutually independent conditionally on $Z$.
\item After discarding zero-mass symbols, each alphabet has a partition into nonempty sets $\mathcal X_i(z)$, indexed by the same finite set, and
\begin{equation}\label{R09-blocklaw}
 P(x_1,\ldots,x_n)=\sum_z\pi_z\prod_{i=1}^np_{i,z}(x_i),
\end{equation}
where $\pi_z>0$, each $p_{i,z}$ is a probability law with support exactly $\mathcal X_i(z)$, and the sets $\mathcal X_i(z)$ are disjoint for fixed $i$.
\end{enumerate}
In these cases $I(X_i;X_j)=H(Z)$ for all $i\ne j$ and $T=(n-1)H(Z)$. The label can be chosen canonically as the connected-component label of the bipartite support graph of $(X_1,X_2)$.
\end{theorem}

We give the support argument explicitly. It is the classical common-part construction of G\'acs and K\"orner \cite{ref077}, rather than a new notion of common information.

\begin{lemma}[Double Markov]\label{R09-double}
Suppose finite random variables $A,B,Y$ satisfy
$I(Y;A\mid B)=I(Y;B\mid A)=0$. Form the bipartite graph whose vertices are positive-mass symbols of $A$ and $B$, with an edge $(a,b)$ when $P(a,b)>0$. Let $Z$ be its connected-component label. Then $Z$ is a function of $A$ and also of $B$, and $Y$ is independent of $(A,B)$ conditionally on $Z$.
\end{lemma}
\begin{proof}
On every support edge, the two conditional independence assumptions imply
\[
 P(Y=\cdot\mid A=a,B=b)=P(Y=\cdot\mid A=a)=P(Y=\cdot\mid B=b).
\]
The common conditional probability vector is therefore constant along graph paths, hence on every component. This is exactly the asserted conditional independence.
\end{proof}

\begin{proof}[Proof of Theorem \ref{R09-blocks}]
Assume $S=0$, and put $A=X_1$, $B=X_2$, $Y=(X_3,\ldots,X_n)$. The stars centered at $1$ and $2$ imply the two hypotheses of Lemma \ref{R09-double}. Let $Z$ be the resulting label.

For every $j\ge3$, the star centered at $j$ gives $I(A;B\mid X_j)=0$. Since $Z$ is a function of both $A$ and $B$, the chain rule gives
\begin{equation}\label{R09-commonidentity}
 I(A;B\mid X_j)=H(Z\mid X_j)+I(A;B\mid Z,X_j).
\end{equation}
The double-Markov conclusion makes the second term equal to $I(A;B\mid Z)$. Nonnegativity now proves both $H(Z\mid X_j)=0$ and $I(A;B\mid Z)=0$.

It remains to check joint, not just pairwise, independence within each block. By the star centered at $1$, the coordinates of $Y$ are mutually independent given $A$. The double-Markov conclusion and the fact that $Z$ is a function of $A$ give $P(Y\mid A)=P(Y\mid Z)$ and the analogous identity for every coordinate of $Y$. Hence $P(Y\mid Z)=\prod_{j\ge3}P(X_j\mid Z)$. Together with $Y\perp(A,B)\mid Z$ and $A\perp B\mid Z$, this proves (ii).

The deterministic label partitions each coordinate's positive-mass alphabet, and conditional independence gives (iii). Conversely (iii) implies (ii). Under (ii),
\[
 H(X_i)=H(Z)+H(X_i\mid Z),\qquad
 H(X_1,\ldots,X_n)=H(Z)+\sum_iH(X_i\mid Z).
\]
These identities give $T=(n-1)H(Z)$ and $I(X_i;X_j)=H(Z)$, so $S=0$. Finally each block's bipartite marginal support is a complete bipartite graph, and different blocks do not meet, verifying the canonical description.
\end{proof}

\begin{corollary}\label{R09-rectangular}
If even one pair $(X_a,X_b)$ has full support on the product of its marginal supports, then equality $S=0$ forces all variables to be mutually independent. Full support of the entire joint law is unnecessary.
\end{corollary}
\begin{proof}
That pair's bipartite support graph is connected, so the common label in Theorem \ref{R09-blocks} is constant. The theorem can be applied after relabeling coordinates.
\end{proof}

\subsection{A quantitative intersection lemma}\label{R09-sec-3}
Write $\operatorname{Ent}_\pi f=\E_\pi[f\log f]-\E_\pi f\log\E_\pi f$ for $f\ge0$ on a finite space. The following elementary comparison is a two-coordinate instance of approximate tensorization, a standard subject in functional inequalities; see \cite{ref078}.

\begin{lemma}[Entropy comparison]\label{R09-comparison}
Let $\pi$ be a probability law on a product $\mathcal A\times\mathcal B$, let $\mu=\pi_A\otimes\pi_B$, and assume
\[
 0<a\le\frac{\pi(x,y)}{\mu(x,y)}\le b<\infty
\]
on the product of the marginal supports. Then every $f\ge0$ satisfies
\begin{equation}\label{R09-AT}
 \operatorname{Ent}_\pi f\le\frac ba\left\{
 \E_\pi\operatorname{Ent}_{\pi(A\mid B)}f+
 \E_\pi\operatorname{Ent}_{\pi(B\mid A)}f\right\}.
\end{equation}
Consequently, for any jointly distributed finite $Y$,
\begin{equation}\label{R09-QI}
 I(Y;A,B)\le\frac ba\{I(Y;A\mid B)+I(Y;B\mid A)\}.
\end{equation}
\end{lemma}
\begin{proof}
For $t>0$ put $\phi(s,t)=s\log(s/t)-s+t\ge0$. The variational identity
\[
 \operatorname{Ent}_\pi f=\inf_{t>0}\E_\pi\phi(f,t)
\]
(with limiting interpretations at zero) gives $\operatorname{Ent}_\pi f\le b\operatorname{Ent}_\mu f$. Product entropy tensorization gives
\[
 \operatorname{Ent}_\mu f\le
 \E_\mu\operatorname{Ent}_{\mu_A}f+
 \E_\mu\operatorname{Ent}_{\mu_B}f.
\]
For clarity, this last inequality follows from the entropy chain rule and convexity of the entropy functional: $\operatorname{Ent}_{\mu_B}(\E_{\mu_A}f)\le\E_{\mu_A}\operatorname{Ent}_{\mu_B}f$; convexity is the log-sum inequality.

The conditional variational identity gives, for example,
\[
 \E_\pi\operatorname{Ent}_{\pi(A\mid B)}f
 =\inf_{g:\mathcal B\to(0,\infty)}\E_\pi\phi(f,g(B))
 \ge a\E_\mu\operatorname{Ent}_{\mu_A}f.
\]
The analogous inequality holds with the coordinates reversed. Combining proves \eqref{R09-AT}.

For every $y$ with $P(Y=y)>0$, apply \eqref{R09-AT} to
$f_y(a,b)=P(Y=y\mid a,b)/P(Y=y)$. Multiply by $P(Y=y)$ and sum. The left side becomes $I(Y;A,B)$; the two conditional entropy terms become $I(Y;A\mid B)$ and $I(Y;B\mid A)$ by direct expansion.
\end{proof}

\subsection{A global strict remainder from pairwise data}\label{R09-sec-4}
\begin{theorem}[Pairwise nondegeneracy remainder]\label{R09-quant}
Suppose a pair $(X_a,X_b)$ satisfies the bounds in Lemma \ref{R09-comparison}, and put $C_{ab}=b/a$. Then
\begin{equation}\label{R09-onepair}
 S\ge\frac{T}{2\{1+(n-1)C_{ab}\}}.
\end{equation}
If every pair has such bounds and $C=\max_{a<b}C_{ab}$, then
\begin{equation}\label{R09-allpairs}
 S\ge\frac{n}{4\{1+(n-1)C\}}T.
\end{equation}
In the second case,
\begin{equation}\label{R09-TVbound}
 \TV\!\left(P,\bigotimes_iP_i\right)
 \le\sqrt{\frac{2\{1+(n-1)C\}}{n}\,S}.
\end{equation}
These constants are explicit, but are not asserted optimal.
\end{theorem}
\begin{proof}
Fix distinct $a,b,c$, and let $Y=X_{[n]\setminus\{a,b\}}$. The conditional total correlations in \eqref{R09-Di} dominate the conditional mutual informations between $X_b$ and $Y$ given $X_a$, and between $X_a$ and $Y$ given $X_b$. Thus \eqref{R09-QI} yields
\[
 J:=I(Y;X_a,X_b)\le C_{ab}(D_a+D_b).
\]
For $j\notin\{a,b\}$, data processing gives $I(X_a;X_j)\le J$. Also, by the mutual information chain rule,
\[
 I(X_a;X_b)\le I(X_a;X_b\mid X_c)+I(X_a;X_c)\le D_c+J.
\]
Using \eqref{R09-Di} once more,
\begin{equation}\label{R09-key}
 T=D_a+\sum_{j\ne a}I(X_a;X_j)
 \le D_a+D_c+(n-1)C_{ab}(D_a+D_b).
\end{equation}
Since $C_{ab}\ge1$ and $\sum_iD_i=2S$, this proves \eqref{R09-onepair}. If every pair has constant at most $C$, average \eqref{R09-key} over all ordered triples of distinct indices. Each $D_i$ appears with average coefficient $2\{1+(n-1)C\}/n$, proving \eqref{R09-allpairs}. Finally $T=\D(P\Vert\bigotimes_iP_i)$ and Pinsker's inequality gives \eqref{R09-TVbound}.
\end{proof}

Only two-coordinate distributions enter $C$. The joint law may vanish on most configurations. For example, a uniform parity law on at least three bits has pairwise independent uniform marginals and hence $C=1$, although its joint support is not a product.

\subsection{Why a conditioning parameter cannot be omitted}\label{R09-sec-5}
Let $R$ be the law of $n$ identical fair bits, let $U$ be the uniform law on all binary strings, and put $P_\varepsilon=(1-\varepsilon)R+\varepsilon U$, $0<\varepsilon<1$. Every $P_\varepsilon$ has full support. By continuity of finite entropy,
\[
 T(P_\varepsilon)\longrightarrow(n-1)\log2>0,
 \qquad S(P_\varepsilon)\longrightarrow0.
\]
Thus there is no positive constant depending only on $n$ such that $S\ge c_nT$ for all strictly positive laws. Here each pair has density ratio $2-\varepsilon$ on equal bits and $\varepsilon$ on unequal bits, so $C=(2-\varepsilon)/\varepsilon\to\infty$, as the bound correctly permits.

Theorem \ref{R09-blocks} is a full equality classification, whereas Theorem \ref{R09-quant} deliberately addresses a nondegenerate regime rather than asserting closeness to some unknown singular block mixture. A dimension-free quantitative rounding theorem to the class \eqref{R09-blocklaw} is not proved. The classical common-part and approximate-tensorization ingredients are acknowledged; priority of this exact combination and its constants is unverified. This paper uses no result of another competitor and makes no benchmark-superiority or named-open-problem claim.





\clearpage
\section[R10 --- Constructive Rounding of Entropy Near-Extremizers to Disjoint Product Blocks]{Constructive Rounding of Entropy Near-Extremizers to Disjoint Product Blocks}\label{contrib-R10}
\markboth{Entropy Inequalities, Marginal Compatibility, and Extremal Laws}{R10 --- Reiwa}
\noindent\textsc{R10 \quad Reiwa programme, source round 10}\par\medskip


\begin{quote}\small\textbf{Overview.} For finite-valued variables we construct an exact extremizer of the complete-graph pairwise-information inequality near any given near-extremizer. The bound is explicit in the entropy deficit and a computable constant of a single pair marginal. That constant is finite even when the pair marginal has zeros: its connected support components and the spectral gaps of their two-coordinate heat-bath operators suffice. The proof first controls conditional total correlation within the common-component label, then decodes that label in every coordinate and truncates a conditional product law. The result is a genuine total-variation rounding theorem, not merely a statement about entropy vectors. Constants are not optimal, and they can diverge as support bottlenecks develop.\end{quote}

\subsection{Statement and the rounding target}\label{R10-sec-1}
Let $X=(X_1,\ldots,X_n)$ have a law $P$ on a finite product, with $n\ge3$. Remove zero-mass symbols from each alphabet. Write
\begin{align*}
 T&=\sum_iH(X_i)-H(X),\\
 D_i&=\sum_{j\ne i}H(X_j\mid X_i)-H(X_{-i}\mid X_i),\\
 S&=\frac n2T-\sum_{i<j}I(X_i;X_j)=\frac12\sum_iD_i\ge0.
\end{align*}
A \emph{disjoint product-block law} is a law under which there exists a finite label $Z$ that is a deterministic function of each coordinate and the coordinates are mutually independent given $Z$. Its support is a union of Cartesian boxes, disjoint in every coordinate, and its density on each box is a product times a block weight. Such laws have $S=0$: each pair has mutual information $H(Z)$ and $T=(n-1)H(Z)$.

Fix two distinct coordinates $a,b$. Form the bipartite support graph of $\pi=P_{X_aX_b}$, and let $Z$ be its connected-component label. This is automatically a deterministic function of $X_a$ and of $X_b$, without any small-deficit assumption. In Section \ref{R10-spectralsection} we define explicitly a finite constant $C_{ab}\ge1$ using only $\pi$.

\begin{theorem}[Constructive block rounding]\label{R10-main}
For every pair $a,b$, there exists a disjoint product-block law $R$ on the original alphabets such that
\begin{equation}\label{R10-simple}
 \TV(P,R)\le\min\left\{1,\sqrt{\{1+(n-1)C_{ab}\}S}+2S\right\}.
\end{equation}
More precisely, put
\begin{equation}\label{R10-BL}
 B_{ab}=\min(D_a,D_b)+\min_{c\notin\{a,b\}}D_c
 +(n-1)C_{ab}(D_a+D_b),\qquad
 L_{ab}=\sum_{j\notin\{a,b\}}D_j.
\end{equation}
If $L_{ab}<1$, an explicit construction using the component label $Z$ yields
\begin{equation}\label{R10-refined}
 \TV(P,R)\le \sqrt{B_{ab}/2}+L_{ab}.
\end{equation}
The construction uses at most as many product blocks as the chosen pair's support graph has components. No positivity of the joint law or of the pair law on a rectangle is required.
\end{theorem}

The result closes a specific structural loop: a small entropy deficit yields an actual nearby equality law, subject to a stated pair-dependent conditioning constant. It does not give a uniform dimension-free bound, an optimal constant, or a new solution to a named open problem. Common-component labels are classical \cite{ref077}, and spectral-to-entropy comparison and approximate tensorization are established methods \cite{ref078,ref079}. All estimates needed for this finite construction are proved below.

\subsection{A finite support-sensitive intersection constant}\label{R10-sec-2}\label{R10-spectralsection}
For a positive-mass component $z$, let $\Omega_z$ be its support edges and let $\pi_z=\pi(\cdot\mid Z=z)$. On the Hilbert space $L^2(\pi_z)$ let $E_a,E_b$ be conditional expectations given the first and second coordinate. They are orthogonal projections. Define
\[
 \mathcal L_z=(I-E_a)+(I-E_b).
\]
If $|\Omega_z|>1$, let $\lambda_z$ be its smallest positive eigenvalue and let $p_z=\min_{\omega\in\Omega_z}\pi_z(\omega)$. Set
\begin{equation}\label{R10-Cdef}
 C_z=\frac{4}{p_z\lambda_z}\quad(|\Omega_z|>1),\qquad
 C_z=1\quad(|\Omega_z|=1),\qquad C_{ab}=\max_zC_z.
\end{equation}
These quantities are finite. Indeed,
\[
 \langle g,\mathcal L_zg\rangle
 =\E_{\pi_z}\operatorname{Var}(g\mid X_a)
 +\E_{\pi_z}\operatorname{Var}(g\mid X_b).
\]
A vector in the kernel is constant on all support edges meeting a given first or second coordinate; connectivity propagates this to the whole component. Thus the kernel consists exactly of constants. Finite-dimensional spectral theory gives $\lambda_z>0$ and
\begin{equation}\label{R10-Poincare}
 \operatorname{Var}_{\pi_z}g\le\lambda_z^{-1}
 \left\{\E\operatorname{Var}(g\mid X_a)+\E\operatorname{Var}(g\mid X_b)\right\}.
\end{equation}
Also $\lambda_z\le2$, since $\mathcal L_z$ is a sum of two projection complements, so $C_{ab}\ge1$.

\begin{lemma}[Entropy versus square-root variance]\label{R10-sqrt}
For any probability law $\nu$ on a finite support, and $f\ge0$,
\begin{equation}\label{R10-sqrlower}
 \operatorname{Ent}_\nu f\ge\operatorname{Var}_\nu\sqrt f.
\end{equation}
If $p=\min_\omega\nu(\omega)>0$, then
\begin{equation}\label{R10-sqrupper}
 \operatorname{Ent}_\nu f\le\frac4p\operatorname{Var}_\nu\sqrt f.
\end{equation}
Here $\operatorname{Ent}_\nu f=\E[f\log f]-\E f\log\E f$.
\end{lemma}
\begin{proof}
Both sides are homogeneous in $f$, so assume $\E f=1$; the zero-mean case is immediate. Let $q(\omega)=f(\omega)\nu(\omega)$. Jensen's inequality applied to $\log$ under $q$ yields
\[
 \D(q\Vert\nu)\ge-2\log\sum_\omega\sqrt{q(\omega)\nu(\omega)}
 =-2\log\E\sqrt f.
\]
Zeros follow by approximation. With $r=\E\sqrt f\le1$, we have $-2\log r\ge1-r^2=\operatorname{Var}\sqrt f$, proving \eqref{R10-sqrlower}.

For the upper bound, $\log x\le x-1$ gives $\operatorname{Ent}f\le\operatorname{Var}f$. Also $\max f\le1/p$. For independent copies $g,g'$ of $\sqrt f$,
\[
 \operatorname{Var}f=\tfrac12\E(g^2-g'^2)^2
 \le4\max(g^2)\operatorname{Var}g
 \le\frac4p\operatorname{Var}g.
\]
This proves \eqref{R10-sqrupper}.
\end{proof}

\begin{lemma}[Conditional intersection]\label{R10-intersection}
For any finite random variable $Y$ jointly distributed with $(X_a,X_b)$,
\begin{equation}\label{R10-intersectioneq}
 I(Y;X_a,X_b\mid Z)
 \le C_{ab}\{I(Y;X_a\mid X_b)+I(Y;X_b\mid X_a)\}.
\end{equation}
\end{lemma}
\begin{proof}
On a component with more than one state, combine \eqref{R10-sqrupper}, \eqref{R10-Poincare}, and the conditional version of \eqref{R10-sqrlower} to obtain
\[
 \operatorname{Ent}_{\pi_z}f\le C_z\left\{
 \E_{\pi_z}\operatorname{Ent}(f\mid X_a)+
 \E_{\pi_z}\operatorname{Ent}(f\mid X_b)\right\}.
\]
For a singleton component both sides vanish. Apply this inequality to
\[
 f_{z,y}(x_a,x_b)=
 \frac{P(Y=y\mid x_a,x_b,Z=z)}{P(Y=y\mid Z=z)}
\]
for positive conditional masses, multiply by $P(Z=z,Y=y)$, and sum. Expanding the entropy functionals gives the left side of \eqref{R10-intersectioneq} and the conditional mutual informations $I(Y;X_b\mid X_a,Z)$ and $I(Y;X_a\mid X_b,Z)$. Since $Z$ is already a function of each of $X_a,X_b$, these are the two terms displayed in \eqref{R10-intersectioneq}.
\end{proof}

The constant is computable from a finite matrix, not an oracle over arbitrary auxiliary variables. In the ordinary Euclidean inner product, conjugate $\mathcal L_z$ by the diagonal matrix with entries $\sqrt{\pi_z(\omega)}$ to obtain a symmetric matrix with the same eigenvalues. Any rigorously certified lower bound on its first positive eigenvalue may replace $\lambda_z$, at the cost of increasing $C_{ab}$.

\subsection{Entropy control within the support label}\label{R10-sec-3}
Let
\[
 T_Z=\sum_iH(X_i\mid Z)-H(X\mid Z).
\]
\begin{proposition}\label{R10-control}
With the preceding notation,
\begin{equation}\label{R10-TZ}
 T_Z\le B_{ab},\qquad H(Z\mid X_j)\le D_j\quad(j\notin\{a,b\}).
\end{equation}
\end{proposition}
\begin{proof}
Write $A=X_a$, $B=X_b$, and $Y=X_{[n]\setminus\{a,b\}}$. Conditional total correlation dominates the conditional mutual information across any partition of its variables. In particular
\[
 I(Y;B\mid A)\le D_a,\qquad I(Y;A\mid B)\le D_b.
\]
Lemma \ref{R10-intersection} gives
\begin{equation}\label{R10-J}
 J:=I(Y;A,B\mid Z)\le C_{ab}(D_a+D_b).
\end{equation}
For every $c\notin\{a,b\}$, the fact that $Z$ is a function of both $A$ and $B$ gives the exact identity
\[
 I(A;B\mid X_c)=H(Z\mid X_c)+I(A;B\mid Z,X_c)\le D_c.
\]
This proves the leakage bound in \eqref{R10-TZ}. Further,
\[
 I(A;B\mid Z)\le I(A;B\mid Z,X_c)+I(A;X_c\mid Z)\le D_c+J.
\]
For $j\notin\{a,b\}$, data processing gives $I(A;X_j\mid Z)\le J$. Since $Z$ is a function of $A$, the conditional form of the star identity is
\[
 T_Z=D_a+\sum_{j\ne a}I(A;X_j\mid Z)
 \le D_a+D_c+(n-1)J.
\]
Interchange $a,b$ to replace $D_a$ by their minimum, minimize over $c$, and use \eqref{R10-J}. This yields the bound $B_{ab}$.
\end{proof}

\subsection{Decoding and truncation: proof of the rounding theorem}\label{R10-sec-4}
Construct the conditional-product law on $(Z,X)$,
\begin{equation}\label{R10-Q}
 Q(z,x)=P_Z(z)\prod_{i=1}^nP(X_i=x_i\mid Z=z).
\end{equation}
It has the same $(Z,X_i)$ marginals as $P$. Moreover $Z$ remains a deterministic function of $X_a$ and $X_b$ under $Q$, because their component alphabets are disjoint. Direct expansion therefore gives
\begin{equation}\label{R10-KLQ}
 \D(P_X\Vert Q_X)=T_Z\le B_{ab}.
\end{equation}
Alternatively one may first compute the divergence on $(Z,X)$ and then use that both laws recover the same $Z$ from $X_a$.

For $j\notin\{a,b\}$, choose a maximum-posterior decoder
\[
 d_j(x_j)\in\arg\max_z P(Z=z\mid X_j=x_j),
\]
breaking ties deterministically. Set $d_a,d_b$ equal to the exact component maps. If a posterior's largest mass is $r$, its entropy is at least $-\log r\ge1-r$. Averaging this inequality proves
\[
 P\{d_j(X_j)\ne Z\}\le H(Z\mid X_j)\le D_j.
\]
The same error probability holds under $Q$, because the relevant pair marginals match. For
\[
 E=\{d_i(X_i)=Z\text{ for every }i\},
\]
the union bound gives $Q(E^c)\le L_{ab}$.

If $L_{ab}<1$, define $R_X$ to be the $X$-marginal of $Q(\cdot\mid E)$. For fixed $z$, the event $E$ is a product of coordinate restrictions $d_i(x_i)=z$. Thus conditioning $Q$ on $E$ preserves conditional product structure and makes the label recoverable from every coordinate. Blocks with zero surviving weight are discarded. Consequently $R$ is a disjoint product-block law.

Pinsker's inequality and \eqref{R10-KLQ} give $\TV(P_X,Q_X)\le\sqrt{B_{ab}/2}$. For completeness, Pinsker follows by grouping the positive and negative parts of $P-Q$, reducing to Bernoulli relative entropy, and using $d''(p\Vert q)=1/[p(1-p)]\ge4$ to get $d(p\Vert q)\ge2(p-q)^2$. Also conditioning a probability law on an event of positive mass changes it in total variation by exactly the event's complement mass. Marginalization cannot increase total variation. Hence
\[
 \TV(P_X,R_X)\le\sqrt{B_{ab}/2}+L_{ab},
\]
proving \eqref{R10-refined}.

Finally $L_{ab}\le2S$ and, since $C_{ab}\ge1$,
\[
 B_{ab}\le\{1+(n-1)C_{ab}\}\sum_iD_i
 =2\{1+(n-1)C_{ab}\}S.
\]
If $2S<1$, the construction applies and gives \eqref{R10-simple}. If $2S\ge1$, the right side of \eqref{R10-simple} is one and any point-mass product law suffices. Total variation is always at most one. This completes the proof of Theorem \ref{R10-main}.

\subsection{What the theorem does and does not make uniform}\label{R10-sec-5}
\subsubsection{Exact equality and support selection}
If $S=0$, then $B_{ab}=L_{ab}=0$ for every pair. The construction returns $R=P$, recovering the full singular equality classification directly. If the chosen pair's support graph is connected, the label is constant and the constructed $R$ is a product law. If it has several components, the construction can preserve genuine common information rather than forcing independence.

The choice of pair is free. One may compute the displayed bounds for every pair and select the smallest. The recipe uses only finite conditional distributions, connected components, a finite symmetric eigenvalue problem, and deterministic maximum-posterior choices; it is constructive given the full probability table. No polynomial-time claim in the number of variables with an implicit input representation is made.

\subsubsection{A necessary warning about bottlenecks}
Let $P_\varepsilon=(1-\varepsilon)R_0+\varepsilon U$, where $R_0$ consists of identical fair bits and $U$ is uniform on all binary strings. Then $S(P_\varepsilon)\to0$, but $P_\varepsilon$ does not approach the product of its marginals. Every positive $\varepsilon$ makes the pair support graph connected, so the exact support label is constant. Accordingly the pair constant in \eqref{R10-Cdef} diverges as $\varepsilon\downarrow0$. The theorem is valid but can be uninformative in this regime. Recovering nearly disconnected labels by a thresholded support analysis with uniform optimal bounds is a different problem, not solved here.

The constants $4/p_z$, the spectral comparison, and the decoding union bound are deliberately transparent rather than optimized. The novelty of this precise combined rounding statement has not been certified by an exhaustive search. The spectral/entropy framework is established literature; this manuscript supplies an explicit finite proof and a concrete extremizer construction. The prior tournament round treated equality and rectangular-pair nondegeneracy; the present round handles arbitrary singular pair supports and constructs an exact nearby block law. No result from another competitor is used.





\clearpage
\part{Polynomial and Rational Differential Spectra}


\clearpage
\section[H01 --- A Degree-Filtration Bound for Algebraic Schr\"odinger Energies]{A Degree-Filtration Bound for Algebraic Schr\"odinger Energies}\label{contrib-H01}
\markboth{Polynomial and Rational Differential Spectra}{H01 --- Heisei}
\noindent\textsc{H01 \quad Heisei programme, source round 1}\par\medskip






\begin{quote}\small\textbf{Overview.} For a polynomial Schr\"odinger potential of degree $2m\ge4$, we prove a bound
$\lfloor n/(m-1)\rfloor+1$ for the number of energies possessing a solution with rational logarithmic derivative. Here $n$ is an explicitly determined nonnegative integer read from the coefficients at infinity; if it does not exist, the spectrum in question is empty. The proof identifies one admissible exponential gauge, then uses the degree filtration of an invariant polynomial space. This improves the elementary $n+1$ estimate when $m>2$ and is attained for every admissible polynomial degree in the sextic family. The rational-logarithmic-derivative theorem is self-contained. Its identification with the Liouvillian spectrum is kept separate and attributed to the established polynomial-potential criterion.\end{quote}
\subsection{Statement and comparison}\label{H01-sec-1}
Let $V\in\C[x]$ have degree $2m$, $m\ge2$, with leading coefficient $a^2\ne0$. Fix one square root $a$. There is a unique polynomial $A$ of degree $m$, leading coefficient $a$, such that
\begin{equation}\label{H01-split} V=A^2+R,\qquad \deg R\le m-1.\end{equation}
It is obtained successively by comparing coefficients of degrees $2m,2m-1,\ldots,m$. Put $r=[x^{m-1}]R$ and
\[ n_+=\frac{r-ma}{2a},\qquad n_-=\frac{-r-ma}{2a}. \]
In particular $n_++n_-=-m$. Define
\[\Sigma_{\rm rat}(V)=\{E\in\C:\ w'+w^2=V-E\text{ has a solution }w\in\C(x)\}.\]
This concerns formal complex energies, not necessarily square-integrable states.
\begin{theorem}\label{H01-main}
If neither $n_+$ nor $n_-$ is a nonnegative integer, then $\Sigma_{\rm rat}(V)=\varnothing$. Otherwise exactly one is a nonnegative integer, denoted $n$, and
\begin{equation}\label{H01-bound}|\Sigma_{\rm rat}(V)|\le\left\lfloor\frac{n}{m-1}\right\rfloor+1.\end{equation}
Every rational solution is uniquely expressible as
\[w=\varepsilon A+P'/P,\qquad \deg P=n,\]
where $P$ is monic and squarefree and $\varepsilon\in\{1,-1\}$ is the unique admissible sign. The bound is sharp for $m=3$ and every $n\ge0$.
\end{theorem}
Acosta-Hum\'anez, Bl\'azquez-Sanz and Venegas-G\'omez give an $n+1$ estimate for the Liouvillian energy set in Theorem 4.2 of the arXiv version of \cite{ref080}. Our estimate is a degree-filtration refinement for the rational logarithmic derivative set; the published polynomial-potential criterion identifies that set with the Liouvillian set. We do not reprove that differential-Galois criterion. The comparison is a comparison with that displayed estimate, not a claim that no stronger bound occurs elsewhere.
\subsection{Rational logarithmic derivatives}\label{H01-sec-2}
\begin{lemma}\label{H01-normal}
If $w\in\C(x)$ satisfies $w'+w^2\in\C[x]$, every finite pole of $w$ is simple and has residue $1$. Thus $w=B+P'/P$ for a polynomial $B$ and a monic squarefree polynomial $P$.
\end{lemma}
\begin{proof}
At a pole of order $k>1$, $w^2$ has order $2k>k+1$, so its leading term cannot be canceled by $w'$. At a simple pole with residue $\rho\ne0$, the coefficient of $(x-z)^{-2}$ is $\rho^2-\rho$, so $\rho=1$. Subtracting the logarithmic derivative of the product of all distinct pole factors leaves a rational function without finite poles, hence a polynomial.
\end{proof}
Write $\deg P=d$. At infinity, $P'/P=d/x+O(x^{-2})$. Consequently
\[ w'+w^2=B^2+B'+2B\frac{P'}P+\frac{P''}P
 =B^2+(m+2d)b_mx^{m-1}+O(x^{m-2})\]
if $\deg B=m$, $[x^m]B=b_m$. Since $V-E$ has degree $2m$, $B$ must have degree $m$ and leading coefficient $\pm a$. All terms other than $B^2$ have degree at infinity at most $m-1$; coefficient matching in \eqref{H01-split} gives $B=\varepsilon A$. Matching degree $m-1$ gives $d=n_\varepsilon$. Since their sum is $-m$, at most one sign is admissible. This proves all normal-form assertions except the bound. The corresponding polynomial equation is
\begin{equation}\label{H01-gauge}
TP=EP,\qquad T=-D^2-2BD+(V-B^2-B'),\quad D=\dd.
\end{equation}
Conversely a nonzero polynomial solution of \eqref{H01-gauge} yields a rational Riccati solution. Its zeros are necessarily simple: if $P(z)=P'(z)=0$, the differential equation and repeated differentiation force all derivatives to vanish. Therefore no extra squarefreeness test is needed after the equation is verified.
\subsection{A degree-filtration lemma}\label{H01-sec-3}
\begin{lemma}\label{H01-ladder}
Suppose a linear operator $T:\C[x]\to\C[x]$ has, for every nonzero polynomial $p$ of degree $d\ne n$,
\[\deg(Tp)=d+s,\qquad s\ge1,\]
and the coefficient at that degree is $c(n-d)$ times the leading coefficient of $p$, with $c\ne0$. Let $W\ne0$ be a finite-dimensional invariant polynomial space all of whose elements have degree at most $n$. Then for some $r=\dim W$ the degrees occurring as leading degrees of a degree-adapted basis are exactly
\[n-(r-1)s,n-(r-2)s,\ldots,n.\]
In particular $r\le\lfloor n/s\rfloor+1$.
\end{lemma}
\begin{proof}
Gaussian elimination on coefficients supplies a basis of distinct increasing degrees; denote that finite set by $S$. It is also precisely the set of degrees of nonzero elements of $W$. If $d\in S$ and $d\ne n$, invariance implies $d+s\in S$. Repetition must eventually stop, and the only possible stopping degree is $n$. Hence every element of $S$ is congruent to $n$ modulo $s$. Starting from its minimum, the same closure condition fills the entire arithmetic interval up to $n$. Nonnegativity of the minimum proves the bound.
\end{proof}
For \eqref{H01-gauge}, the coefficient of $x^{d+m-1}$ in $T(x^d)$ is $2b_m(n-d)$; all remaining terms have lower degree. If $E_1,\ldots,E_k$ are distinct elements of $\Sigma_{\rm rat}(V)$, choose their monic polynomial factors $P_i$. Eigenvectors belonging to distinct eigenvalues are linearly independent: apply $\prod_{j\ne i}(T-E_j)$ to a putative relation. Their span is invariant, all their degrees equal $n$, and Lemma \ref{H01-ladder} applies with $s=m-1$. Since the resulting bound is independent of the chosen finite subset, it bounds the entire set and proves \eqref{H01-bound}.
\begin{corollary}
If $0\le n<m-1$, there is at most one such energy. If $n=m-1$, there are at most two. Thus large potential degree makes a fixed polynomial factor degree more restrictive, rather than merely changing an auxiliary resultant.
\end{corollary}
\subsection{Sharpness in every sextic degree}\label{H01-sec-4}
Choose $a>0$, $b\in\R$, $p\in\{0,1\}$ and $N\ge0$, and set $n=2N+p$. Consider
\begin{equation}\label{H01-sextic}
V(x)=a^2x^6+2abx^4+\bigl(b^2-a(4N+2p+3)\bigr)x^2.
\end{equation}
Set $F=ax^4/4+bx^2/2$. Direct differentiation shows that $e^F(-D^2+V)e^{-F}$ preserves
\[W=x^p\Span\{1,x^2,\ldots,x^{2N}\}\]
and acts on $x^{p+2k}$ as
\begin{align*}
T x^{p+2k}={}&-2k(2k+2p-1)x^{p+2k-2}\\
&+(4k+2p+1)b x^{p+2k}-4a(N-k)x^{p+2k+2}.
\end{align*}
The lower coefficient at $k=0$ is zero. Both off-diagonal coefficients linking neighbors have negative sign. A positive diagonal change of basis therefore makes this finite matrix real symmetric tridiagonal with off-diagonal magnitudes
\[\sqrt{8a(N-k)(k+1)(2k+2p+1)},\qquad0\le k<N.\]
They are strictly positive. Every eigenvector of such a matrix is determined by its first coordinate through the row recurrence; if that coordinate is zero, all coordinates vanish. Thus every eigenspace is one-dimensional. Symmetry gives a complete orthogonal eigenbasis, so there are $N+1$ distinct real eigenvalues. Each supplies a rational logarithmic derivative. Theorem \ref{H01-main} gives the reverse bound $\lfloor n/2\rfloor+1=N+1$. Hence equality holds. Moreover $P e^{-F}\in L^2(\R)$, so these are genuine bound states as well as algebraic energies.
\subsection{Exact verification and limits}\label{H01-sec-5}
No root-finding is required to verify a proposed state: expand
\[P''+2BP'+(B'+B^2-V+E)P\]
and check all coefficients are zero. To search, write $P=x^n+\sum_{j<n}c_jx^j$ and solve that finite coefficient system in $c_j,E$. The theorem bounds the number of distinct possible energies, not the algebraic multiplicity of an eliminant or the number of real zeros of an eigenfunction.

For example, $a=1,b=0,N=1,p=0$ gives $V=x^6-7x^2$. In the basis $(1,x^2)$ the matrix is
\[\begin{pmatrix}0&-2\\-4&0\end{pmatrix},\]
so $E=\pm2\sqrt2$, with $P=x^2-E/4$. Substitution gives $TP=EP$ exactly. The opposite exponential sign is excluded by the negative degree sum.

\paragraph{Scope.} The quadratic case $m=1$ is intentionally excluded: its coefficient controlling polynomial degree contains $E$, and infinitely many energies are possible. Odd-degree polynomial potentials have no rational Riccati solution, by the leading-degree comparison in Lemma \ref{H01-normal}. No statement here solves a named pre-existing open question or certifies priority. The elementary refinement and its sharpness are the submitted claims.




\clearpage
\section[H02 --- Pole Budgets and Finite Certificates for Rational Schr\"odinger Spectra]{Pole Budgets and Finite Certificates for Rational Schr\"odinger Spectra}\label{contrib-H02}
\markboth{Polynomial and Rational Differential Spectra}{H02 --- Heisei}
\noindent\textsc{H02 \quad Heisei programme, source round 2}\par\medskip






\begin{quote}\small\textbf{Overview.} We give an explicit finite certificate and a degree-sensitive counting theorem for rational logarithmic derivatives of Schr\"odinger eigenfunctions when the potential has double finite poles and polynomial growth of degree at least four. Each local exponent choice and each sign at infinity yields an arithmetic degree budget. The number of distinct admissible energies in that sector is at most the budget divided by the growth shift, rounded down, plus one. A rectangular polynomial matrix gives an exact, terminating search without genericity assumptions, including resonant poles and repeated exponent roots. The result is elementary and does not identify every possible Liouvillian solution for a rational potential.\end{quote}
\subsection{Input data and theorem}\label{H02-sec-1}
Let
\[V(x)=V_0(x)+\sum_{j=1}^{q}\left(\frac{g_j}{(x-z_j)^2}+\frac{h_j}{x-z_j}\right),\]
where the $z_j$ are pairwise distinct, $g_j\ne0$, and $V_0$ has degree $2m\ge4$. Write uniquely
\[V_0=A^2+R,\quad\deg A=m,\quad [x^m]A=a\ne0,\quad\deg R\le m-1,\]
and put $r=[x^{m-1}]R$. At each pole let
\[\mathcal E_j=\{\rho\in\C:\rho(\rho-1)=g_j\};\]
this is a set, so a repeated root is counted once. For $\varepsilon=\pm1$ and $\rho=(\rho_1,\ldots,\rho_q)\in\prod_j\mathcal E_j$, define
\begin{align}
B_{\varepsilon,\rho}(x)&=\varepsilon A(x)+\sum_j\frac{\rho_j}{x-z_j},\label{H02-B}\\
n_{\varepsilon,\rho}&=\frac{\varepsilon r/a-m}{2}-\sum_j\rho_j.\label{H02-n}
\end{align}
Only sectors with $n_{\varepsilon,\rho}\in\mathbb Z_{\ge0}$ will be called admissible.
\begin{theorem}\label{H02-count}
Every rational solution of $w'+w^2=V-E$ has the form
\[ w=B_{\varepsilon,\rho}+P'/P \]
in an admissible sector, with $P$ monic of degree $n_{\varepsilon,\rho}$, squarefree and nonvanishing at every $z_j$. Conversely any polynomial $P\ne0$ satisfying
\begin{equation}\label{H02-poly}
-P''-2BP'+(V-B^2-B')P=EP
\end{equation}
produces such a rational solution, whether or not $P(z_j)=0$.
For a fixed admissible sector of budget $n$, the number of distinct energies certified by \eqref{H02-poly} is at most
\[\left\lfloor\frac n{m-1}\right\rfloor+1.\]
Consequently
\begin{equation}\label{H02-total}
|\Sigma_{\rm rat}(V)|\le
\sum_{(\varepsilon,\rho)\,\mathrm{admissible}}
\left(\left\lfloor\frac{n_{\varepsilon,\rho}}{m-1}\right\rfloor+1\right).
\end{equation}
Empty sums mean that no rational logarithmic derivative exists.
\end{theorem}
The potentially overlapping sector spectra explain why \eqref{H02-total} is an upper bound rather than an exact addition formula. The coefficients $h_j$ constrain the certificate but not the arithmetic budget; they can therefore destroy admissible states without increasing the bound.
\subsection{Local-to-global normal form}\label{H02-sec-2}
At an ordinary point of $V$, a pole of $w$ can only be simple with residue one: a pole of order $k>1$ makes $w^2$ dominate $w'$, and a simple-pole residue obeys $\rho^2-\rho=0$. At a double pole $z_j$ of $V$, $w$ must have a simple pole, and comparison of the double coefficient yields $\rho_j(\rho_j-1)=g_j$. A higher-order pole is impossible by the same order comparison. Since $g_j\ne0$, neither root is zero.

Subtract the prescribed simple fractions and the logarithmic derivative of the product of all remaining pole factors. The result has no finite poles and hence is a polynomial $C$. Writing $d$ for the degree of the latter product gives
\[w=C+\frac{d+\sum_j\rho_j}{x}+O(x^{-2}).\]
The leading-degree comparison forces $\deg C=m$ and $C=\varepsilon A$: all remaining terms in $w'+w^2-C^2$ have order at infinity at most $m-1$. The coefficient at $x^{m-1}$ is
\[\varepsilon a\left(m+2d+2\sum_j\rho_j\right),\]
which proves \eqref{H02-n}. All ordinary pole factors are distinct, giving the stated squarefreeness and disjointness.

The converse follows from direct multiplication:
\[\left(B+\frac{P'}P\right)'+\left(B+\frac{P'}P\right)^2
=B'+B^2+\frac{P''+2BP'}P.\]
If a certificate has $P(z_j)=0$, its actual residue is $\rho_j+\operatorname{ord}_{z_j}P$; the equation forces this to be another permitted local exponent. Thus allowing these certificates creates no false positive and is convenient for linear algebra.
\subsection{The counting argument with rational coefficients}\label{H02-sec-3}
Put $T=-D^2-2BD+C$, where $C=V-B^2-B'$. In a sector of budget $n$, $C$ has growth at infinity $2\varepsilon a n x^{m-1}+O(x^{m-2})$. Thus for every polynomial $p$ of degree $d\ne n$,
\begin{equation}\label{H02-lead}
Tp=2\varepsilon a(n-d)\operatorname{lc}(p)x^{d+m-1}
+O(x^{d+m-2})
\end{equation}
at infinity. Although $Tp$ may be rational for a general $p$, it is polynomial whenever $p$ lies in the span of polynomial eigenfunctions of $T$.

Choose any $k$ distinct eigenvalues in the sector and eigenpolynomials $P_1,\ldots,P_k$. Every one has degree exactly $n$, by \eqref{H02-lead}. Their span $W$ is a $k$-dimensional invariant space of polynomials. A coefficient-echelon basis has distinct degrees; let $S$ be this set. For $d\in S\setminus\{n\}$, \eqref{H02-lead} and invariance force $d+m-1\in S$. Repetition reaches $n$, so $S$ is exactly one uninterrupted arithmetic progression ending at $n$, of step $m-1$. Therefore $k\le\lfloor n/(m-1)\rfloor+1$. This works for every finite subset of the spectrum, proving the sector bound and finiteness.
\subsection{A finite exact certificate}\label{H02-sec-4}
Let $Q=\prod_j(x-z_j)$. The coefficients of
\[\mathcal M_E(P)=Q^2\bigl(-P''-2BP'+(C-E)P\bigr)\]
are polynomial. For an admissible budget $n$, use the domain basis $1,x,\ldots,x^n$ and the codomain basis through degree $n+m-1+2q$. This gives an explicit rectangular matrix $M(E)$ with $n+1$ columns whose entries are affine functions of $E$.
\begin{theorem}\label{H02-algorithm}
The energy set in a fixed sector equals
\[\{E\in\C:\operatorname{rank}M(E)<n+1\}.\]
Equivalently it is the common zero set of all $(n+1)$-rowed maximal minors. A nonzero kernel vector gives a polynomial certificate. All nonzero polynomial certificates have degree $n$. Taking the union over the finitely many admissible sectors gives exactly $\Sigma_{\rm rat}(V)$.
\end{theorem}
\begin{proof}
Multiplication by the nonzero polynomial $Q^2$ does not change equality in $\C(x)$. Kernel vectors are therefore exactly polynomial solutions of \eqref{H02-poly}. A nonzero polynomial of degree $d<n$ cannot solve it, because the uncanceled leading term in \eqref{H02-lead} has degree $d+m-1>d$. The rank criterion is elementary linear algebra; the local-to-global theorem proves exhaustion.
\end{proof}
Over an effective exact coefficient field, compute the greatest common divisor of the maximal minors as univariate polynomials in $E$. At least one minor is nonzero: otherwise every $E$ would be in the spectrum, contradicting the already proved finite bound. The gcd, made monic, has exactly the required roots. If its value is a nonzero constant the sector is empty. Repeated roots of the gcd do not imply multiple distinct eigenvalues. This procedure is finite, but no polynomial-time complexity claim is made.
\subsection{An exactly solvable test with a genuine pole}\label{H02-sec-5}
Choose $a>0$, $b\in\R$, $N\ge0$ and $\rho>1$. Let
\begin{align*}
V(x)={}&a^2x^6+2abx^4+
\bigl(b^2-a(4N+2\rho+3)\bigr)x^2+\frac{\rho(\rho-1)}{x^2}.
\end{align*}
For $\psi=x^\rho e^{-ax^4/4-bx^2/2}P(x^2)$ on $x>0$, the gauged operator acts on $t^k$, $t=x^2$, as
\[Tt^k=-2k(2k+2\rho-1)t^{k-1}
 +(4k+2\rho+1)bt^k-4a(N-k)t^{k+1}.
\]
Thus the $N+1$ dimensional polynomial space in $t$ of degree at most $N$ is invariant. Its off-diagonal products are
\[8a(N-k)(k+1)(2k+2\rho+1)>0.\]
A positive diagonal similarity makes the matrix real symmetric irreducible tridiagonal, hence its $N+1$ eigenvalues are distinct and real. The sector has degree budget $n=2N$, so the bound is attained. Here the gauge in \eqref{H02-B} is $B=-ax^3-bx+\rho/x$. The singular exponent affects both the local behavior and the off-diagonal entries, illustrating information absent from the polynomial case.
\subsection{What is and is not claimed}\label{H02-sec-6}
The local exponent enumeration resembles the rational branch of Kovacic's algorithm \cite{ref082}. The submitted contribution is the growth-shift sector bound and its integration with a complete rectangular-matrix certificate. No blanket assertion that a Liouvillian solution of an arbitrary rational equation must have rational logarithmic derivative is made; that assertion would be false in general. Simple finite poles can be added with a forced residue one, but are outside the stated input class. Irregular finite poles require a larger local principal-part enumeration and are also not claimed here.




\clearpage
\section[H03 --- The Monomial-Module Dichotomy for Polynomial Schr\"odinger Gauges]{The Monomial-Module Dichotomy for Polynomial Schr\"odinger Gauges}\label{contrib-H03}
\markboth{Polynomial and Rational Differential Spectra}{H03 --- Heisei}
\noindent\textsc{H03 \quad Heisei programme, source round 3}\par\medskip






\begin{quote}\small\textbf{Overview.} We classify, without assuming that the exponents are consecutive, every finite monomial invariant space of dimension at least two for a polynomial Schr\"odinger gauge operator $-D^2+2A(x)D+B(x)$ with $\deg A\ge2$. There are exactly two possibilities: a full polynomial interval with quadratic drift, or a full parity interval with cubic odd drift. All one-dimensional exceptions are also classified. This gives an elementary support-closure proof of the quartic/sextic dichotomy within the explicitly stated monomial class, and explains precisely why it is not a classification of arbitrary quasi-exactly solvable potentials. The familiar potential families are classical; the exhaustive support argument is the focus.\end{quote}
\subsection{Definitions and classification}\label{H03-sec-1}
For a nonempty finite set $S\subseteq\mathbb Z_{\ge0}$ put
\[M_S=\Span_\C\{x^k:k\in S\}.\]
Throughout $A,B\in\C[x]$, $m=\deg A\ge2$, and
\[T=-D^2+2A(x)D+B(x),\qquad D=\dd.\]
Invariance means $T(M_S)\subseteq M_S$; the monomial coordinate is the given $x$, with no unmentioned change of variable.
\begin{theorem}\label{H03-class}
If $|S|\ge2$, invariance holds if and only if one of the following alternatives holds.
\begin{enumerate}[label=(\roman*)]
\item $A=a_2x^2+a_1x+a_0$, $a_2\ne0$,
\[S=\{0,1,\ldots,n\},\quad n\ge1,\qquad B=-2na_2x+b_0.\]
\item $A=a_3x^3+a_1x$, $a_3\ne0$,
\[S=\{p,p+2,\ldots,p+2N\},\quad p\in\{0,1\},\quad N\ge1,\]
\[B=-2(p+2N)a_3x^2+b_0.\]
\end{enumerate}
In particular no polynomial drift of degree at least four preserves a monomial space of dimension at least two.
\end{theorem}
The restriction to monomial spaces is essential. An arbitrary invariant polynomial space can have basis vectors containing many lower terms; the support argument below does not apply to it. In particular the final sentence is not a nonexistence theorem for all high-degree quasi-exactly solvable potentials.
\begin{theorem}\label{H03-single}
If $S$ has one element, the complete list is:
\begin{enumerate}[label=(\roman*)]
\item $S=\{0\}$, with $B$ constant and $A$ arbitrary of degree at least two;
\item $S=\{1\}$, with $A(0)=0$ and $B=c-2A/x$ for a constant $c$.
\end{enumerate}
\end{theorem}
\subsection{The degree graph}\label{H03-sec-2}
Write $A=\sum_{j=0}^m a_jx^j$ and let $n=\max S$. If $\deg B>m-1$, the highest term of $B x^n$ cannot be canceled by $2A D x^n$ or by $-D^2x^n$, and has degree greater than $n$. Invariance therefore forces $\deg B\le m-1$. Comparing the coefficient at degree $n+m-1$ next gives
\begin{equation}\label{H03-top}b_{m-1}=-2na_m.\end{equation}
Hence for every $k\in S$ with $k<n$,
\begin{equation}\label{H03-up} k+m-1\in S,\end{equation}
since the corresponding coefficient is $2a_m(k-n)\ne0$.

There is also an uncancelable lower edge. The coefficient of $x^{k-2}$ in $T(x^k)$ is $-k(k-1)$: polynomial drift terms have degree at least $k-1$, and multiplication terms have degree at least $k$. Consequently
\begin{equation}\label{H03-down}k\in S,\ k\ge2\quad\Longrightarrow\quad k-2\in S.\end{equation}
These two exact closure rules are the entire combinatorial mechanism.
\subsection{Excluding higher drift}\label{H03-sec-3}
Assume first that $n\ge2$. Rule \eqref{H03-down} puts $n-2$ in $S$. Rule \eqref{H03-up} then puts $n+m-3$ in $S$, so maximality of $n$ implies $m\le3$. If $n=1$ and $|S|\ge2$, then $S=\{0,1\}$. Rule \eqref{H03-up} applied to zero gives $m-1\le1$, hence $m=2$. Thus only degrees two and three remain.

If $m=2$, rule \eqref{H03-up} fills every integer from $\min S$ to $n$. Rule \eqref{H03-down} forces $\min S\le1$. If $\min S=1$ and the dimension is at least two, then $2\in S$ and \eqref{H03-down} forces $0\in S$, a contradiction. Thus $S=\{0,\ldots,n\}$. Formula \eqref{H03-top} says exactly $B=-2na_2x+b_0$. Conversely the three possible output degrees $k-2,k-1,k,k+1$ all lie in the interval unless a coefficient vanishes at the boundary; the coefficients at the lower boundary do vanish, and the raising coefficient vanishes at $k=n$. This proves sufficiency.

If $m=3$, rule \eqref{H03-up} has step two. Every element of $S$ must have the same parity as $n$: an element of opposite parity would eventually reach $n-1$ and then force $n+1\in S$. Rule \eqref{H03-down} fills that parity interval down to $p=n\bmod2$. Thus $S=\{p,p+2,\ldots,n\}$.

The coefficient of $x^{k+1}$ is $2k a_2+b_1$. It must vanish for every $k\in S$, since that exponent has opposite parity. There are at least two distinct $k$, so $a_2=b_1=0$. The coefficient of $x^{k-1}$ is $2k a_0$ and must also vanish; some $k>0$ is present, giving $a_0=0$. The remaining coefficient condition is $b_2=-2na_3$. Conversely, all surviving shifts are $-2,0,+2$, and the two end coefficients vanish exactly as required. This completes the proof of Theorem \ref{H03-class}.
\subsection{One-dimensional spaces}\label{H03-sec-4}
For $S=\{n\}$ with $n\ge2$, the term $-n(n-1)x^{n-2}$ prevents invariance, so no such case exists. For $n=0$, $T1=B$ gives the first alternative. For $n=1$,
\[Tx=2A+Bx=cx\]
for some scalar $c$. Polynomiality of $B=c-2A/x$ is equivalent to $A(0)=0$. This proves Theorem \ref{H03-single}, including sufficiency.
\subsection{Translation into potentials}\label{H03-sec-5}
Let $F'=A$ and $H=-D^2+V$. A direct product-rule computation gives
\[e^FHe^{-F}=-D^2+2AD+V-A^2+A'.\]
Thus the associated potential is $V=A^2-A'+B$.

The quadratic-drift family gives
\[V=(a_2x^2+a_1x+a_0)^2-(2a_2x+a_1)-2na_2x+b_0.\]
It preserves $e^{-F}\Pcal_n$. When $A$ has real coefficients and $a_2\ne0$, $F$ has odd degree three, so no nonzero member of this module is square-integrable on the full real line. This is a boundary-condition obstruction, not an algebraic failure.

The cubic-drift family gives
\begin{align*}
V={}&a_3^2x^6+2a_3a_1x^4
+\bigl(a_1^2-(2n+3)a_3\bigr)x^2+(b_0-a_1),\\
n={}&p+2N.
\end{align*}
For real $a_3>0$, $a_1,b_0$, the gauge is decaying at both ends and every module vector lies in $L^2(\R)$. In the basis $x^{p+2k}$ the tridiagonal coefficients are
\begin{align*}
T x^{p+2k}={}&-(p+2k)(p+2k-1)x^{p+2k-2}\\
&+\bigl(2a_1(p+2k)+b_0\bigr)x^{p+2k}
-4a_3(N-k)x^{p+2k+2}.
\end{align*}
The products of paired off-diagonal entries are strictly positive, so diagonal symmetrization supplies $N+1$ distinct real algebraic energies. The resulting sextic oscillator is the classical quasi-exactly solvable family, not a new potential.
\subsection{A support test that does not guess the module}\label{H03-sec-6}
For reproducibility, the proof yields the following finite decision procedure once $A,B$ are given.
\begin{enumerate}
\item If $\deg A\ge4$, test only the two singleton possibilities of Theorem \ref{H03-single}.
\item If $\deg A=2$, read $n=-b_1/(2a_2)$, require $n\in\mathbb Z_{\ge1}$ and $\deg B\le1$, then the only multidimensional candidate is $\Pcal_n$.
\item If $\deg A=3$, require $a_2=a_0=b_1=0$, $\deg B\le2$, and $n=-b_2/(2a_3)\in\mathbb Z_{\ge2}$. The only multidimensional candidate is the full parity interval ending at $n$.
\item Verify the displayed coefficient identities. No numerical spectral computation is part of the decision.
\end{enumerate}
The one-dimensional test should always be performed separately; otherwise high-degree exact single states would be incorrectly discarded.
\paragraph{Position in the literature.} Polynomial invariant modules and their Lie-algebraic realization are classical \cite{ref083,ref084}. This paper gives a self-contained classification for a restricted but completely specified class, with singleton and physical-boundary exceptions included. It does not claim a new classification of all finite-dimensional polynomial subspaces or resolution of an open problem.




\clearpage
\section[H04 --- Maximal Polynomial Modules and Their Hidden Jacobi Matrices]{Maximal Polynomial Modules and Their Hidden Jacobi Matrices}\label{contrib-H04}
\markboth{Polynomial and Rational Differential Spectra}{H04 --- Heisei}
\noindent\textsc{H04 \quad Heisei programme, source round 4}\par\medskip






\begin{quote}\small\textbf{Overview.} For a polynomial differential operator with a single positive degree shift and a unique terminating degree, we construct its largest finite polynomial invariant space using one finite block-kernel calculation. This detects generalized eigenvectors as well as eigenvectors and does not assume a monomial module. In the real confining Schr\"odinger case, the maximal space carries a canonical irreducible Jacobi matrix ordered by polynomial degree. We obtain exact spectral completeness, simplicity, interlacing, positive spectral weights, and a finite characteristic-polynomial certificate. The invariant-subspace and Jacobi ingredients are standard linear algebra; their degree-sensitive combination is proved here without genericity assumptions.\end{quote}
\subsection{Setup}\label{H04-sec-1}
Let
\[T=-D^2+2A(x)D+B(x),\quad \deg A=m\ge2,\quad [x^m]A=a\ne0,\]
where $\deg B\le m-1$ and $[x^{m-1}]B=-2an$ for some integer $n\ge0$. Set $s=m-1$ and $\Pcal_n=\{p:\deg p\le n\}$. For every polynomial of degree $d\ne n$,
\begin{equation}\label{H04-degree}
Tp=2a(d-n)\operatorname{lc}(p)x^{d+s}+O(x^{d+s-1}).
\end{equation}
Let $\pi$ denote coefficient truncation onto $\Pcal_n$, and define
\[L=\pi T|_{\Pcal_n},\qquad C=(I-\pi)T|_{\Pcal_n}.\]
Thus $L$ is an $(n+1)$-square matrix and $C$ has at most $s$ output coefficients, the degrees $n+1,\ldots,n+s$. Define
\begin{equation}\label{H04-W}
W_* = \bigcap_{j=0}^{n}\ker(CL^j).
\end{equation}
\begin{theorem}\label{H04-max}
The space $W_*$ is the unique largest finite-dimensional $T$-invariant polynomial space. Every finite-dimensional invariant polynomial space lies in $\Pcal_n$ and in $W_*$. In particular every polynomial eigenvector and every polynomial Jordan chain lies in $W_*$. If $r=\dim W_*>0$, its degree sequence is
\[d_0,d_0+s,\ldots,d_0+(r-1)s=n,\qquad d_0\ge0,\]
so $r\le\lfloor n/s\rfloor+1$.
\end{theorem}
\subsection{Proof and finite termination}\label{H04-sec-2}
If $W$ is a finite invariant polynomial space and $d$ is its largest degree, formula \eqref{H04-degree} forces $d=n$; otherwise the degree increases beyond $d$. Hence $W\subseteq\Pcal_n$. For $p\in W$, all $T^j p$ lie in $\Pcal_n$, so $C L^j p=0$ for all $j\ge0$. Thus $W\subseteq W_*$.

Conversely, by Cayley--Hamilton, every power $L^{n+1}$ is a linear combination of $I,L,\ldots,L^n$. Therefore $CL^{n+1}$ vanishes on $W_*$, and the defining equations show $L(W_*)\subseteq W_*$. Also $C(W_*)=0$, so $T=L$ there and $W_*$ is invariant. This proves maximality and uniqueness. A polynomial Jordan chain spans a finite invariant space and is therefore included.

For the degree assertion, a coefficient-echelon basis has a set $S$ of distinct degrees. For each $d\in S$ below $n$, invariance and \eqref{H04-degree} force $d+s\in S$. Every chain ends at $n$, so all degrees have the same residue modulo $s$, and the least one generates the whole arithmetic interval. This proves Theorem \ref{H04-max}.

An equivalent descending algorithm is
\[U_0=\Pcal_n,\qquad U_{j+1}=\{p\in U_j:Tp\in U_j\}.\]
At a strict step dimension decreases. Once a step is unchanged, the space is invariant and the sequence stabilizes. At most $n+1$ strict steps occur, and the stable value is $W_*$. The block matrix in \eqref{H04-W} avoids symbolic quantifiers in this description.
\subsection{A positive inner product changes the conclusion}\label{H04-sec-3}
Now assume $A,B$ have real coefficients, $m$ is odd, and $a>0$. Choose a real antiderivative $F$ of $A$ and put $\omega=e^{-2F}$. Since $F$ has even degree and positive leading coefficient, every polynomial and its derivatives are square-integrable against $\omega$. Set
\[\langle p,q\rangle=\int_{\R}p(x)q(x)\omega(x)\,dx.\]
An integration by parts, with vanishing boundary terms, gives
\begin{equation}\label{H04-sym}
\langle p,Tq\rangle=\int_{\R}(p'q'+Bpq)\omega\,dx=\langle Tp,q\rangle.
\end{equation}
The weight is strictly positive, so this is an inner product on polynomials. Therefore the restriction to $W_*$ is a real symmetric operator. This excludes Jordan chains of length greater than one, but simplicity will require the degree argument too.
\begin{theorem}\label{H04-jac}
If $r=\dim W_*>0$, Gram--Schmidt on a degree-echelon basis produces an orthonormal basis $\phi_0,\ldots,\phi_{r-1}$ with degrees $d_j=d_0+js$ and positive leading coefficients $\kappa_j$. The matrix of $T$ is tridiagonal, with adjacent entries
\[\beta_j=2a(d_j-n)\frac{\kappa_j}{\kappa_{j+1}}<0,
\qquad 0\le j<r-1.\]
After alternating basis signs, it becomes a Jacobi matrix with strictly positive off-diagonal entries $b_j=-\beta_j$. Consequently $T|_{W_*}$ has exactly $r$ distinct real eigenvalues, each eigenspace is one-dimensional, and every eigenpolynomial has degree $n$.
\end{theorem}
\begin{proof}
For $j<r-1$, $T\phi_j$ has degree $d_{j+1}$; thus it is orthogonal to $\phi_i$ for $i>j+1$. Symmetry gives the corresponding zeros below the first subdiagonal. Comparing leading coefficients gives the formula for $\beta_j$, which is nonzero. A tridiagonal eigenvector is determined by its first coordinate through the first row and successive row equations. If that coordinate vanishes, every coordinate vanishes, so eigenspaces have dimension one. The real spectral theorem supplies $r$ eigenvectors; hence all eigenvalues are distinct. Finally, a polynomial eigenvector of degree below $n$ contradicts \eqref{H04-degree}.
\end{proof}
This theorem does not assume the space is a parity module, or even that $d_0$ is zero or one. Such assumptions would exclude exceptional polynomial spaces without justification.
\subsection{Spectral certificates and interlacing}\label{H04-sec-4}
Let $\alpha_j=\langle\phi_j,T\phi_j\rangle$. Write $J_k$ for the leading $k$-square block of the positive-off-diagonal Jacobi matrix, and $D_k(E)=\det(EI-J_k)$. Expansion of the last row gives
\[D_0=1,\quad D_1=E-\alpha_0,\quad
D_{k+1}=(E-\alpha_k)D_k-b_{k-1}^2D_{k-1}.\]
\begin{corollary}
The roots of $D_k$ are real and simple and strictly interlace those of $D_{k+1}$. The full algebraic energy set is exactly the root set of $D_r$.
\end{corollary}
\begin{proof}
The min--max characterization of the eigenvalues of a real symmetric matrix gives weak interlacing for a principal compression: the admissible subspaces are restricted by at most one dimension. Equality would be a common root of $D_k,D_{k+1}$; the recurrence then forces $D_{k-1}=0$ and ultimately $D_0=0$, impossible. This also proves simplicity inductively. Maximality of $W_*$ proves exhaustion of polynomial eigenvalues.
\end{proof}
At a full eigenvalue $\lambda$, the normalized spectral weight at the first basis vector is
\begin{equation}\label{H04-weight}
\mu_\lambda=\left[\sum_{k=0}^{r-1}
\frac{D_k(\lambda)^2}{b_0^2\cdots b_{k-1}^2}\right]^{-1}>0,
\end{equation}
where the empty product is one. Indeed the eigenvector recurrence gives its $k$th coordinate as its first coordinate times $D_k(\lambda)/(b_0\cdots b_{k-1})$. Normalization yields \eqref{H04-weight}. Completeness of the eigenbasis gives $\sum\mu_\lambda=1$ and
\[\langle e_0,J^\ell e_0\rangle=\sum_\lambda\mu_\lambda\lambda^\ell,\qquad\ell\ge0.\]
Thus the construction produces a positive finite spectral measure, not only an eliminant with unclassified roots.
\subsection{Implementation independent of unknown integrals}\label{H04-sec-5}
For exact coefficient data one need not evaluate any weighted integral to find the energy polynomial. Compute a basis matrix $U$ for the kernel in \eqref{H04-W}; solve $TU=UK$ for the finite restriction matrix $K$ in that basis. Then
\[\det(EI-K)=D_r(E).\]
The integrals prove positivity and simplicity but are not required for the algebraic certificate. Verify $TU=UK$ coefficient by coefficient, then verify its characteristic polynomial by exact arithmetic. If $W_*=0$, the same kernel computation is a certificate that no polynomial eigenfunction exists.

For a reproducible small example, $T=-D^2+2x^3D-4x^2$ has $n=2,s=2$. In $\Pcal_2$ the leakage of $x$ is $-2x^3$, while $1$ and $x^2$ have no leakage. The kernel calculation gives $W_*=\Span\{1,x^2\}$ and
\[K=\begin{pmatrix}0&-2\\-4&0\end{pmatrix},\qquad\det(EI-K)=E^2-8.\]
The associated potential is $x^6-7x^2$. In contrast, adding a generic linear term to $B$ changes both the leakage kernel and its $L$-iterates; checking only the first leakage equation is not a valid invariance test.
\subsection{Prior work and boundary of the result}\label{H04-sec-6}
Orthogonal spectral polynomials in quasi-exact solvability are classical \cite{ref081}; invariant polynomial spaces are classical as well \cite{ref083}. The present paper isolates a finite maximal-space calculation that does not guess the space and combines it with the forced arithmetic degree ladder. It makes no general new-classification or named-open-problem claim. The positivity conclusion requires the full-line decay assumptions; over complex coefficients, the same kernel algorithm remains correct but Jordan blocks and nonreal eigenvalues can occur.




\clearpage
\section[H05 --- Two-State Polynomial Modules: Divisibility, Decatic Classification, and Exact Nodal Indices]{Two-State Polynomial Modules: Divisibility, Decatic Classification, and Exact Nodal Indices}\label{contrib-H05}
\markboth{Polynomial and Rational Differential Spectra}{H05 --- Heisei}
\noindent\textsc{H05 \quad Heisei programme, source round 5}\par\medskip






\begin{quote}\small\textbf{Overview.} We classify all polynomial Schr\"odinger gauges having a two-dimensional polynomial invariant space containing the constant function. The classification is an exact divisibility condition on one polynomial. In the real confining case it determines the ground state, the exact index of the other algebraic state, and the complete rational-logarithmic-derivative spectrum. We solve the condition completely for even quartic generators, obtaining a nonmonomial decatic family with exactly two algebraic energies. The second state is the second excited state on one parameter component and the fourth excited state on the other. This is not a claim that two-state inverse constructions or decatic quasi-exact solvability are new; the explicit iff classification and exhaustive energy count are the precise submitted results.\end{quote}
\subsection{Two-dimensional modules containing one}\label{H05-sec-1}
Let $T=-D^2+2AD+B$ with $A,B\in\C[x]$, $\deg A\ge2$. Suppose a two-dimensional invariant polynomial space $W$ contains $1$. Then $B$ is nonconstant: if $B$ were constant, a nonconstant element of maximal degree $d$ in $W$ would be sent to degree $d+\deg A-1>d$. Since $T1=B$, necessarily $W=\Span\{1,B\}$. There are unique scalars $u,v$ such that $TB=uB+v$. Recenter by
\[C=B-u/2,\qquad T_0=T-u/2=-D^2+2AD+C,\qquad\Lambda=v+u^2/4.\]
Then $T_0 1=C$ and $T_0 C=\Lambda$.
\begin{theorem}\label{H05-general}
A nonconstant polynomial $C$ gives such a centered module if and only if
\begin{equation}\label{H05-div} C'\ \text{divides}\ C''-C^2+\Lambda
\end{equation}
for some scalar $\Lambda$. In that case
\begin{equation}\label{H05-drift}
A=\frac{C''-C^2+\Lambda}{2C'},\qquad
V=A^2-A'+C,
\end{equation}
and, for either square root $\eta^2=\Lambda$,
\[\psi_\pm=(C\pm\eta)e^{-F},\qquad F'=A,\qquad H\psi_\pm=\pm\eta\,\psi_\pm,
\quad H=-D^2+V.
\]
If $n=\deg C$, then $\deg A=n+1$, and $\operatorname{lc}(A)=-\operatorname{lc}(C)/(2n)$. Every zero of $C'$ is simple. If $\Lambda\ne0$, the complete rational-logarithmic-derivative energy set is exactly $\{-\eta,\eta\}$.
\end{theorem}
\begin{proof}
The equation $T_0C=\Lambda$ is exactly $-C''+2AC'+C^2=\Lambda$, which proves necessity, sufficiency, and the quotient formula. The leading degree of the quotient is $n+1$ with the stated coefficient. The two eigenfunction formulas follow by applying $T_0$ to $C\pm\eta$.

If $z$ is a zero of $C'$ of multiplicity $\ell\ge2$, write $C=C(z)+c(x-z)^{\ell+1}+\cdots$, $c\ne0$. Divisibility first requires $C(z)^2=\Lambda$. Then $C''$ has a nonzero term of order $\ell-1$, whereas $C^2-\Lambda$ has order at least $\ell+1$. Their difference is not divisible by $(x-z)^\ell$, a contradiction.

For exhaustiveness, recall the elementary rational Riccati normal form: for polynomial $V$ of degree $2m\ge4$, any rational solution of $w'+w^2=V-E$ is $w=G+P'/P$, with $G$ polynomial of degree $m$ and $P$ squarefree. The high coefficients determine $G$ up to sign, and the two possible polynomial degrees have sum $-m$, so at most one is nonnegative. Here the displayed solutions give $G=-A$, $\deg P=n$, $m=n+1$. For distinct energies their polynomial factors span an invariant space. Its degree set is closed under $d\mapsto d+m-1$ until reaching $n$, because the leading coefficient of the gauged operator is a nonzero multiple of $d-n$. With $m-1=n$, the degree set is contained in $\{0,n\}$ and the energy count is at most two. The displayed two distinct values exhaust it. This argument concerns rational logarithmic derivatives and uses no differential-Galois converse.
\end{proof}
\subsection{Nodal indices in the real confining case}\label{H05-sec-2}
Assume $C$ is real of even degree $n\ge2$ with negative leading coefficient and satisfies \eqref{H05-div} for real $\Lambda$. Then $A$ has odd degree and positive leading coefficient, so $e^{-F}$ decays at both ends. On polynomials define
\[\langle f,g\rangle=\int_\R fg e^{-2F}\,dx.\]
Integration by parts makes $T_0$ symmetric. Therefore
\[\Lambda\langle1,1\rangle=\langle1,T_0C\rangle
=\langle T_01,C\rangle=\langle C,C\rangle>0.
\]
In particular $\eta=\sqrt\Lambda>0$.
\begin{theorem}\label{H05-nodes}
Let $r$ be the number of real roots of $C'$. Then $r$ is odd, all those roots are simple, $C-\eta$ is strictly negative on $\R$, and $C+\eta$ has exactly $r+1$ simple real roots. Thus $-\eta$ is the ground-state energy and $+\eta$ is the $(r+1)$th excited energy, numbering the ground state by zero. There are no other rational-logarithmic-derivative energies.
\end{theorem}
\begin{proof}
At a critical point $z$, \eqref{H05-div} gives $C''(z)=C(z)^2-\eta^2$. At a local maximum, simplicity of $C'$ gives $C''(z)<0$, so $|C(z)|<\eta$. Since $C\to-\infty$ at both ends, the global maximum is one of these local maxima; hence $C<\eta$ everywhere. At a local minimum, $C''>0$, and the already proved upper bound forces $C(z)<-\eta$. Thus successive extrema alternate strictly above and below $-\eta$, with both tails below it. There is exactly one crossing in each interval between successive extrema and one in each exterior interval: $r+1$ crossings in total. None is critical, so they are simple. The eigenfunctions have the same zeros as their polynomial factors, and the classical one-dimensional oscillation theorem identifies their indices. Exhaustiveness is Theorem \ref{H05-general}.
\end{proof}
The oscillation theorem is the sole standard analytic spectral input used for the numerical state index; all zero counts and all algebraic assertions were proved above.
\subsection{Complete even quartic-generator classification}\label{H05-sec-3}
Let $C=\beta x^4+\gamma x^2+\delta$, with $\beta\ne0$. Condition \eqref{H05-div} has a solution if and only if
\begin{equation}\label{H05-parameters}
\gamma\ne0,\qquad
\delta=\frac{12\beta}{\gamma}+\frac{\gamma^2}{8\beta},\qquad
\Lambda=\delta^2-2\gamma.
\end{equation}
The associated drift is
\begin{equation}\label{H05-Aexplicit}
A=-\frac\beta8x^5-\frac{3\gamma}{16}x^3
-\frac{48\beta^2+\gamma^3}{16\beta\gamma}x.
\end{equation}
\begin{proof}
If $\gamma=0$, $C'=4\beta x^3$, but the numerator $C''-C^2+\Lambda$ has coefficient $12\beta x^2$, so divisibility fails. For $\gamma\ne0$, $C'$ has the three distinct roots $0$ and $x^2=-\gamma/(2\beta)$. At zero, divisibility gives $\Lambda=\delta^2-2\gamma$. At either other critical point it gives
\[\Lambda=\left(\delta-\frac{\gamma^2}{4\beta}\right)^2+4\gamma.\]
Equating these values gives \eqref{H05-parameters}. Since the derivative is squarefree, vanishing at all its roots is equivalent to divisibility. Polynomial division gives \eqref{H05-Aexplicit}.
\end{proof}
Now let $\beta<0$, $\gamma\in\R\setminus\{0\}$. The potential $V=A^2-A'+C$ is a real even polynomial of degree ten, with positive leading coefficient $\beta^2/64$. Also
\[\Lambda=
\frac{(\gamma^3+32\beta^2)^2+8192\beta^4}{64\beta^2\gamma^2}>0.
\]
If $\gamma<0$, $C'$ has only one real root, so the exact states have indices $0$ and $2$. If $\gamma>0$, it has three real roots, so the exact states have indices $0$ and $4$. The entire rational-logarithmic-derivative spectrum is these two energies in both components.
\subsection{An integer-coefficient certificate}\label{H05-sec-4}
For $\beta=-8$, $\gamma=-8$ one obtains
\[C=-8x^4-8x^2+11,\qquad\Lambda=137,\qquad
A=x^5+\frac32x^3-\frac52x.
\]
Thus
\[V=x^{10}+3x^8-\frac{11}{4}x^6
-\frac{41}{2}x^4-\frac{25}{4}x^2+\frac{27}{2}.
\]
The exact energies are $\pm\sqrt{137}$ and
\[\psi_\pm=(-8x^4-8x^2+11\pm\sqrt{137})
\exp\left(-\frac{x^6}{6}-\frac{3x^4}{8}+\frac{5x^2}{4}\right).
\]
The identities to check are only
\[2AC'=C''-C^2+137,\qquad V=A^2-A'+C.
\]
The negative-energy eigenfunction has no real zeros and the positive-energy one has two. No decimal approximation or external code is needed to check these statements.
\subsection{Comparison and limits}\label{H05-sec-5}
General constructions with two known eigenstates predate this work \cite{ref085}; high-degree polynomial QES examples also predate it. The 2025 nilpotent-group construction \cite{ref086} considers the restricted form $V=A^2+\alpha A'$ up to its parametrizations. The present family is not of that form with the same odd quintic gauge: matching its quartic coefficient would require $C$ to be proportional to $A'$ with factor $-8/5$, while the quadratic coefficients would then be $9\gamma/10$ rather than $\gamma$. Since $\gamma\ne0$, that restriction fails, even allowing an additive energy constant.

This comparison does not establish novelty against all two-state inverse methods. The exact contributions submitted are the divisibility iff condition, the complete even-quartic-generator solution, the exhaustive two-energy statement, and the nodal classification. We do not claim that all decatic two-state potentials contain the constant in their gauged module, nor that a named prior open problem has been solved.




\clearpage
\section[H06 --- A Spectral-Index Gap in Degree-Fourteen Two-State Polynomial Potentials]{A Spectral-Index Gap in Degree-Fourteen Two-State Polynomial Potentials}\label{contrib-H06}
\markboth{Polynomial and Rational Differential Spectra}{H06 --- Heisei}
\noindent\textsc{H06 \quad Heisei programme, source round 6}\par\medskip






\begin{quote}\small\textbf{Overview.} We solve the two-state divisibility problem completely for even sextic polynomial generators. The classification is rational in the two nonzero squared critical points. For real confining Schr\"odinger potentials of degree fourteen arising from this class, the algebraic states can have indices only $(0,2)$ or $(0,6)$: the intermediate possibility $(0,4)$ is impossible. Exact parameter chambers and an explicit rational-coefficient example with the sixth excited state are provided. The statement is restricted to two-dimensional constant-containing even polynomial gauge modules; no general exclusion for all degree-fourteen potentials is asserted.\end{quote}
\subsection{The class and the result}\label{H06-sec-1}
Let $C$ be an even polynomial of degree six and suppose
\begin{equation}\label{H06-certificate}
C''-C^2+\Lambda=2AC'
\end{equation}
for polynomials $A$ and a constant $\Lambda$. Then $A$ is odd of degree seven, and the potential
\[V=A^2-A'+C\]
has degree fourteen. The functions $(C\pm\sqrt\Lambda)e^{-\int A}$ solve the Schr\"odinger equation at energies $\pm\sqrt\Lambda$.

For distinct nonzero $u,v$, define
\begin{align*}
f(t)&=t^3-\frac32(u+v)t^2+3uvt,\\
\Delta&=uv(u-3v)(u-v)^2(3u-v),\\
N&=3u^2-14uv+3v^2,\\
Q&=3u^4-20u^3v+42u^2v^2-20uv^3+3v^4.
\end{align*}
\begin{theorem}\label{H06-classification}
Every even sextic solution of \eqref{H06-certificate} is, uniquely up to interchanging $u,v$, given by
\begin{align}\label{H06-param}
\beta&=-\frac{24N}{\Delta},\qquad
\delta=-\frac{6(u+v)Q}{\Delta},\\
C(x)&=\beta f(x^2)+\delta,\qquad
\Lambda=\delta^2-6\beta uv,\nonumber
\end{align}
where $\Delta N\ne0$. Conversely every such pair $u,v$ gives a solution, with
\[A=\frac{C''-C^2+\Lambda}{2C'}\in\C[x].\]
\end{theorem}
The formula is a polynomial quotient certificate: expanding $A$ is optional, and no external computer algebra is required to define or check the potential.
\begin{theorem}\label{H06-physical}
Suppose $C$ has real coefficients and $\beta<0$. Then $u,v$ are real and have the same sign. Put $t=\max(|u|,|v|)/\min(|u|,|v|)>1$ and
\[t_*=(7+2\sqrt{10})/3.\]
The confining condition is exactly
\[1<t<3\qquad\text{or}\qquad t>t_*.
\]
If $u,v<0$, the two algebraic states have indices $0,2$. If $u,v>0$, they have indices $0,6$. In particular index pair $0,4$ never occurs in this class. In each case the complete rational-logarithmic-derivative spectrum consists of the two displayed energies.
\end{theorem}
\subsection{Critical points are necessarily simple}\label{H06-sec-2}
We first justify a fact needed to avoid silently discarding degenerate cases. If a nonconstant polynomial satisfies \eqref{H06-certificate}, every root of $C'$ is simple. Indeed, a zero of order $\ell\ge2$ in $C'$ yields $C=C(z)+c(x-z)^{\ell+1}+\cdots$. Divisibility requires $C(z)^2=\Lambda$ but leaves a nonzero term of order $\ell-1$ from $C''$, which cannot be divisible by $(x-z)^\ell$.

Since $C$ is even sextic, its derivative is odd quintic. Simplicity therefore forces
\[C'(x)=6\beta x(x^2-u)(x^2-v),\quad uv(u-v)\ne0,\]
and integration gives $C=\beta f(x^2)+\delta$.
\subsection{Solving every local compatibility equation}\label{H06-sec-3}
At the critical point zero, equation \eqref{H06-certificate} gives
\[\Lambda=\delta^2-6\beta uv.\]
At either point with $x^2=u$ it gives
\[\beta f(u)^2+2\delta f(u)=6u(2u-3v),\]
and at either point with $x^2=v$ it gives
\[\beta f(v)^2+2\delta f(v)=6v(2v-3u).\]
Here
\[f(u)=\frac{u^2(3v-u)}2,\quad f(v)=\frac{v^2(3u-v)}2,
\quad f(u)-f(v)=-\frac{(u-v)^3}{2}.\]
If $u=3v$, the first left side vanishes while its right side is $54v^2\ne0$, a contradiction. Similarly $v=3u$ is impossible. Thus the two linear equations for $\beta,\delta$ have nonzero determinant. Cramer's rule gives exactly \eqref{H06-param}. If $N=0$, it would force $\beta=0$, contrary to degree six. These observations prove necessity with every degenerate case addressed.

Conversely, the formulas solve the compatibility equations at all five simple critical points. Hence $C'$ divides $C''-C^2+\Lambda$ over $\C[x]$. Its quotient is odd of degree seven, proving sufficiency and Theorem \ref{H06-classification}.
\subsection{The real sign obstruction}\label{H06-sec-4}
For real $C$, the numbers $u+v$ and $uv$ are real. Thus either $u,v$ are both real or they are a nonreal conjugate pair.

Suppose they are nonreal conjugates. Write $u=re^{i\theta}$, $v=re^{-i\theta}$ with $r>0$, $\sin\theta\ne0$. Then
\begin{align*}
N&=r^2(6\cos2\theta-14)<0,\\
(u-v)^2&=-4r^2\sin^2\theta<0,\\
(u-3v)(3u-v)&=r^2(6\cos2\theta-10)<0.
\end{align*}
Hence $\Delta>0$ and $\beta=-24N/\Delta>0$, so this case is not confining.

If $u,v$ are real and of opposite signs, $N>0$, $uv<0$, and $(u-3v)(3u-v)>0$. Therefore $\Delta<0$, again giving $\beta>0$. Thus confinement forces the same sign. By interchange and simultaneous sign reversal, the sign calculation reduces to $u>0$, $v=tu$, $t>1$. Up to a positive factor,
\[\beta\ \text{has sign}\ -\frac{3t^2-14t+3}{(1-3t)(3-t)}.
\]
The larger numerator root is $t_*$, while the denominator changes sign at $t=3$. A direct sign chart yields $\beta<0$ precisely on $(1,3)\cup(t_*,\infty)$, as asserted. Endpoints are excluded either by a zero denominator, a repeated critical point, or loss of degree.
\subsection{Positivity, nodal count, and spectral completeness}\label{H06-sec-5}
When $\beta<0$, the leading coefficient of $A$ is $-\beta/12>0$. Thus $F'=A$ has positive even leading term and $e^{-F}$ decays at both ends. With $T=-D^2+2AD+C$, equation \eqref{H06-certificate} says $T1=C$ and $TC=\Lambda$. Integration by parts with weight $e^{-2F}$ yields
\[\Lambda\int e^{-2F}dx=\int C^2e^{-2F}dx>0.
\]
Put $\eta=\sqrt\Lambda>0$.

At every real local maximum of $C$, the compatibility identity gives $C''=C^2-\eta^2<0$, hence $|C|<\eta$. The global maximum is therefore below $\eta$. At each local minimum, $C''>0$ and the global upper bound force $C<-\eta$. Since both tails tend to $-\infty$, $C-\eta$ has no real zeros and $C+\eta$ has exactly one more real zero than $C'$ has. If $u,v<0$, only zero is a real critical point; if both are positive, all five critical points are real. The zero counts are therefore two and six. The standard Sturm oscillation theorem gives the stated indices.

For completeness of the algebraic spectrum, the polynomial Riccati normal form gives a unique admissible exponential sign and polynomial factor degree six. Any distinct algebraic energies give independent degree-six eigenpolynomials. Their invariant span has degree set closed under the shift $m-1=6$ below degree six. Therefore only degrees zero and six can occur and its dimension is at most two. The two displayed real distinct energies exhaust the rational-logarithmic-derivative set.
\subsection{A fully rational example}\label{H06-sec-6}
Take $u=1$, $v=2$. Then
\[\beta=-156/5,\quad\delta=171/5,\quad\Lambda=38601/25,
\]
\[C(x)=\frac{-156x^6+702x^4-936x^2+171}{5}.
\]
Define $A$ by the exact quotient in Theorem \ref{H06-classification}, and $V=A^2-A'+C$. The quotient is a polynomial because its numerator vanishes at $0,\pm1,\pm\sqrt2$, the five distinct roots of $C'$. This is a certificate using only rational arithmetic and one square root.

The degree-fourteen potential has exactly the two rational-logarithmic-derivative energies
\[E_0=-\sqrt{38601}/5,\qquad E_6=\sqrt{38601}/5,
\]
with wavefunctions $(C+E_j)e^{-\int A}$. The ground state has no zeros; the other has six. Replacing $u,v$ by $-1,-2$ leaves $\beta$ unchanged and gives the index pair $0,2$ instead.
\subsection{Status and limitations}\label{H06-sec-7}
Two-state inverse constructions are known \cite{ref087}; the novelty question here is the specific complete sextic-generator parametrization and its real sign obstruction. A targeted primary-source search located general two-state methods and recent high-degree polynomial constructions \cite{ref086}, but did not establish an identical formula or an absence of one. Therefore priority is not certified.

The exclusion of $0,4$ is not a theorem about every even degree-fourteen potential. It requires an invariant space containing $1$ after a common polynomial exponential gauge and an even sextic centered generator. Other polynomial modules, rational gauges, or non-even generators are outside the statement.




\clearpage
\section[H07 --- Symmetry-Breaking Branches in the Complete Quartic-Generator Classification]{Symmetry-Breaking Branches in the Complete Quartic-Generator Classification}\label{contrib-H07}
\markboth{Polynomial and Rational Differential Spectra}{H07 --- Heisei}
\noindent\textsc{H07 \quad Heisei programme, source round 7}\par\medskip






\begin{quote}\small\textbf{Overview.} We classify all degree-four polynomial generators of a two-dimensional constant-containing Schr\"odinger gauge module, without imposing reflection symmetry. After translation, the parameter space has an even branch and a genuinely asymmetric branch meeting at an explicit bifurcation locus. In the real confining region every asymmetric solution has exactly two rational-logarithmic-derivative energies, with state indices zero and four. A rational-coefficient asymmetric potential with energies $\pm9\sqrt{65}$ is given by a compact polynomial certificate. The theorem is an inverse classification within a specified module class, not a classification of all decatic quasi-exactly solvable potentials.\end{quote}
\subsection{Problem and normal form}\label{H07-sec-1}
Consider polynomials $A,C$ and a constant $\Lambda$ satisfying
\begin{equation}\label{H07-identity}C''-C^2+\Lambda=2AC'.\end{equation}
Then $T=-D^2+2AD+C$ preserves $\Span\{1,C\}$, with $T1=C$, $TC=\Lambda$. With $F'=A$ and $V=A^2-A'+C$, the Schr\"odinger eigenstates are
\[(C\pm\sqrt\Lambda)e^{-F},\qquad E=\pm\sqrt\Lambda.
\]
Every two-dimensional polynomial invariant space containing $1$ can be centered into this form: $T1=B$ forces the space to be $\Span\{1,B\}$, and subtracting half the trace of its restriction gives $C=B-u/2$.

If $C$ has degree four, translate $x$ uniquely to remove its cubic coefficient. Thus write
\[C(x)=\beta x^4+\gamma x^2+\zeta x+\delta,\qquad\beta\ne0.
\]
The remaining coordinate freedom is reflection and scaling; neither is silently used in the classification.
\begin{theorem}\label{H07-branches}
Identity \eqref{H07-identity} holds for a polynomial $A$ if and only if one of the following holds.
\begin{enumerate}[label=(\roman*)]
\item The even branch:
\[\zeta=0,\quad\gamma\ne0,\quad
\delta=\frac{12\beta}{\gamma}+\frac{\gamma^2}{8\beta},\quad
\Lambda=\delta^2-2\gamma.
\]
\item The asymmetric branch:
\[\zeta\ne0,\quad
\delta=\frac{7\gamma^2}{24\beta},\quad
\zeta^2=\frac{8(72\beta^2-\gamma^3)}{27\beta},\quad
\Lambda=\frac{9\gamma^4}{64\beta^2}-6\gamma.
\]
\end{enumerate}
In both branches the drift is
\begin{equation}\label{H07-A}
A=-\frac\beta8x^5-\frac{3\gamma}{16}x^3
-\frac{7\zeta}{32}x^2-\frac{8\beta\delta+\gamma^2}{32\beta}x
-\frac{3\gamma\zeta}{32\beta}.
\end{equation}
The closures of the two branches meet exactly where
\[\zeta=0,\qquad\gamma^3=72\beta^2,\qquad
\delta=7\gamma^2/(24\beta).
\]
\end{theorem}
\subsection{An exact remainder proof}\label{H07-sec-2}
Polynomial division of $C''-C^2+\Lambda$ by $2C'$ gives the quotient \eqref{H07-A}; its remainder is
\begin{align}\label{H07-rem}
\frac1{16\beta}\bigl[&
(192\beta^2-16\beta\gamma\delta-9\beta\zeta^2+2\gamma^3)x^2\nonumber\\
&+\zeta(-24\beta\delta+7\gamma^2)x
+32\beta\gamma-16\beta\delta^2+16\beta\Lambda+3\gamma\zeta^2\bigr].
\end{align}
This displayed division identity can be checked simply by multiplying the quotient by $2C'$; it is a complete finite certificate for the calculation.

If $\zeta=0$, vanishing of the quadratic coefficient forces $\gamma\ne0$ and gives the even formula for $\delta$. The constant coefficient then gives $\Lambda$. If $\zeta\ne0$, the linear coefficient forces $\delta=7\gamma^2/(24\beta)$. The quadratic coefficient gives the formula for $\zeta^2$, and the constant coefficient gives the displayed $\Lambda$. Conversely these substitutions annihilate every coefficient in \eqref{H07-rem}. This proves the iff assertion, with no genericity assumption and no discarded special case. The branch intersection follows by setting $\zeta=0$ in the asymmetric relation and comparing the two formulas for $\delta$.
\subsection{Real confinement and symmetry breaking}\label{H07-sec-3}
Assume all coefficients are real and $\beta<0$. The leading coefficient of $A$ is positive, so $F$ has positive degree-six leading term and the two states are square-integrable whenever $\Lambda>0$.
\begin{theorem}\label{H07-real}
On the asymmetric branch, real solutions exist exactly when
\[\gamma>(72\beta^2)^{1/3},\qquad
\zeta=\pm\sqrt{\frac{8(72\beta^2-\gamma^3)}{27\beta}}.
\]
Both have $\Lambda>0$ and their algebraic state indices are $0$ and $4$. They are related by reflection $x\mapsto-x$. At equality the branches join an even potential; below it there is no real asymmetric branch.
\end{theorem}
\begin{proof}
With $\beta<0$, positivity of $\zeta^2$ is equivalent to $\gamma^3>72\beta^2$. This implies
\[\Lambda=\frac{3\gamma(3\gamma^3-128\beta^2)}{64\beta^2}>0.
\]
The critical points are roots of $4\beta x^3+2\gamma x+\zeta$. Dividing by $4\beta$ makes this a depressed monic cubic. Its discriminant is
\[-\frac{\gamma^3}{2\beta^3}-\frac{27\zeta^2}{16\beta^2}
=-\frac{36}{\beta}>0.
\]
Thus $C'$ has three distinct real roots. The nodal argument in the next section proves that the excited algebraic state has four real zeros and the other has none. Reflection changes the sign of $\zeta$ and sends the corresponding drift and potential to their reflected versions.
\end{proof}
The word ``bifurcation'' here describes the exact algebraic branch geometry, not a claim about a thermodynamic phase transition or spontaneous breaking in a dynamical system.
\subsection{Why the states are exactly the ground and fourth states}\label{H07-sec-4}
Put $\eta=\sqrt\Lambda$. At a critical point of $C$, identity \eqref{H07-identity} gives $C''=C^2-\eta^2$. Since the critical points are simple, local maxima have values strictly between $-\eta$ and $\eta$. Every local minimum has value below $-\eta$: its square exceeds $\eta^2$, but the global maximum is already below $\eta$.

The polynomial tends to $-\infty$ at both ends. Its three critical points are therefore a maximum, minimum, maximum. The horizontal line $-\eta$ is crossed once on each of the four monotonic intervals, while the line $+\eta$ is never crossed. Thus $C+\eta$ has four simple real zeros and $C-\eta$ none. Multiplication by the positive gauge does not change zeros. The standard full-line Sturm oscillation theorem identifies energies $-\eta$ and $+\eta$ as levels zero and four.

The rational Riccati spectrum has no additional points. To see this without an eliminant, any rational logarithmic derivative for a polynomial potential is a polynomial gauge plus the logarithmic derivative of a squarefree polynomial. Comparing the highest coefficients selects the same gauge $-A$ and factor degree four; the other sign would require a negative degree. Distinct energies yield independent eigenpolynomials whose invariant span has degrees closed upward by the shift $\deg A-1=4$ until degree four. Only degrees zero and four can occur, so there are at most two energies. The two displayed energies attain that count.
\subsection{A compact rational example}\label{H07-sec-5}
Choose
\[\beta=-216,\quad\gamma=216,\quad\zeta=96,\quad\delta=-63.
\]
Then
\begin{align*}
C&=-216x^4+216x^2+96x-63,\\
A&=27x^5-\frac{81}{2}x^3-21x^2+\frac{45}{2}x+9,\\
\Lambda&=5265=81\cdot65.
\end{align*}
The identity $2AC'=C''-C^2+5265$ is checked directly by multiplication. Define the even-degree but non-even potential
\[V=A^2-A'+C.
\]
It is a degree-ten real polynomial with positive leading coefficient $729$; keeping it in this factored certificate form is less error-prone than listing eleven expanded coefficients. Its only rational-logarithmic-derivative energies are
\[E_0=-9\sqrt{65},\qquad E_4=9\sqrt{65}.
\]
The corresponding square-integrable eigenfunctions are
\[(C+E_j)\exp\left(-\frac92x^6+\frac{81}{8}x^4+7x^3-\frac{45}{4}x^2-9x\right).
\]
Every exponent coefficient follows by integrating $A$ and changing sign. The reflected choice $\zeta=-96$ gives the reflected potential and the same two algebraic energies.
\subsection{Scope and priority}\label{H07-sec-6}
Two-state inverse constructions are longstanding \cite{ref087}; general polynomial-potential algebraic methods are also established \cite{ref080,ref086}. The submitted result is the complete quartic-generator parameter decomposition, including the non-even branch, exact discriminant, nodal indices, and spectral exhaustiveness. A targeted literature search does not certify novelty. No assertion is made that every two-state decatic potential has this module, or that the branch geometry solves a previously named open question.




\clearpage
\section[H08 --- A Complete Three-State Decatic Module with Exact Levels Zero, Four, and Eight]{A Complete Three-State Decatic Module with Exact Levels Zero, Four, and Eight}\label{contrib-H08}
\markboth{Polynomial and Rational Differential Spectra}{H08 --- Heisei}
\noindent\textsc{H08 \quad Heisei programme, source round 8}\par\medskip
\noindent\textit{Integration note.} The polynomial congruence notation in the discriminant computation incorporates the supplied typographical correction.\par\medskip






\begin{quote}\small\textbf{Overview.} We classify three-dimensional constant-containing invariant polynomial modules for even decatic Schr\"odinger gauges. After scaling the leading drift to one and removing an additive energy constant, there is exactly one real parameter solution. A rational representative has a complete algebraic energy spectrum given by an explicit cubic and exact state indices $0,4,8$. All three polynomial-exponential states are certified by finite polynomial identities; the nodal counts are certified by rational interval signs and a reduced quartic discriminant. The module is not a full monomial parity space. No novelty certification or claim about all three-state decatic potentials is made.\end{quote}
\subsection{The inverse classification}\label{H08-sec-1}
Consider
\[T=-D^2+2(x^5+bx^3+cx)D-16x^4+dx^2,
\qquad b,c,d\in\R.
\]
An additive constant in the multiplication term is removed by translating all energies. The highest degree of any finite invariant polynomial space must be eight: on $x^k$ the leading raising coefficient is $2(k-8)$ and the shift is four.
\begin{theorem}\label{H08-classification}
The operator $T$ has a three-dimensional polynomial invariant space containing $1$ if and only if
\begin{equation}\label{H08-bc}
 b^3=-\frac{999}{14},\qquad c=\frac{b^2}{3},\qquad d=-\frac{32b}{3}.
\end{equation}
The real solution is unique. Its invariant space is exactly
\[W=\Span\{1,B,P\},\qquad B=-16x^4+dx^2,\quad P=TB.
\]
It is the maximal polynomial invariant space and has three distinct real eigenvalues. The corresponding polynomial potential is $V=A^2-A'+B$, $A=x^5+bx^3+cx$.
\end{theorem}
The assumption that the module contains $1$ is substantive. No claim is made that all polynomial modules with three eigenvalues have that property.
\subsection{Elimination certificate for completeness}\label{H08-sec-2}
If such a space exists, $T1=B$ lies in it, and $P=TB$ has degree eight with nonzero leading coefficient. Hence $1,B,P$ are linearly independent and must form its basis. Invariance is therefore equivalent to
\begin{equation}\label{H08-closure}TP=uP+vB+w\end{equation}
for constants $u,v,w$.
Comparing coefficients of $x^{10}$ and $x^8$ gives
\[d=-32b/3,\qquad u=-8(b^2-9c)/3.
\]
The remaining coefficients, after nonzero scalar factors are removed, are exactly
\begin{align*}
0={}&28b^3-126bc-999,\\
0={}&320b^4-2016b^2c-11376b+3456c^2+27v,\\
0={}&-32b^3c-160b^2+240bc^2+3bv-648c,\\
0={}&-512b^3+3840bc+9w+3456.
\end{align*}
These four displayed equations, together with the two top equations, are a complete coefficient certificate for \eqref{H08-closure}.

The first equation forces $b\ne0$ and gives $c=(28b^3-999)/(126b)$. The second solves for $v$. Substituting both expressions into the third leaves exactly
\[-\frac{192(14b^3+999)}{49b}=0.
\]
Thus \eqref{H08-bc} follows. Back-substitution gives
\[u=16b^2/3,\qquad v=10624b/21,\qquad w=39936/7.
\]
Conversely these values satisfy every displayed equation, proving necessity and sufficiency, including all potential exceptional cases.

For maximality, any invariant polynomial space has largest degree eight. Its degree set is closed by $k\mapsto k+4$ below eight, hence is a subset of $\{0,4,8\}$. Its dimension is at most three. Our space attains this bound. The weight $e^{-2\int A}$ is positive and rapidly decreasing; integration by parts makes $T$ symmetric on the module. In a degree-ordered orthonormal basis its matrix is irreducible tridiagonal, since degrees increase from zero to four to eight with nonzero leading coefficients. Hence the three eigenvalues are real and simple.
\subsection{A rational representative}\label{H08-sec-3}
A positive dilation of the coordinate gives the particularly convenient data
\begin{align*}
A(x)&=\frac{37}{14}(x^5-3x^3+3x),\\
B(x)&=\frac{296}{7}(2x^2-x^4),\\
P(x)=TB(x)&=\frac{1184}{49}(37x^8-148x^6+74x^4+132x^2-7).
\end{align*}
The potential $V=A^2-A'+B$ is an explicit real even polynomial of degree ten, with positive leading coefficient. The exact closure identity is
\[TP=\frac{888}{7}P-\frac{196544}{49}B+\frac{738816}{49}.
\]
All coefficients are rational; direct differentiation and multiplication verify the entire module.

Let
\begin{equation}\label{H08-chi}
\chi(E)=49E^3-6216E^2+196544E-738816.
\end{equation}
In the basis $1,B,P$, the restriction matrix is\[K=\begin{pmatrix}0&0&738816/49\\1&0&-196544/49\\0&1&888/7\end{pmatrix}.\] Its characteristic polynomial is $\chi(E)/49$. For each root $E$ of \eqref{H08-chi}, an eigenpolynomial is
\begin{equation}\label{H08-p}
p_E=P+(E-888/7)B+E^2-(888/7)E+196544/49.
\end{equation}
Indeed $Tp_E=Ep_E$ follows immediately from the closure identity and $\chi(E)=0$.
\subsection{Three real energies with exact isolating intervals}\label{H08-sec-4}
The endpoint values are
\begin{align*}
\chi(4)&=-48960,&\chi(5)&=94629,\\
\chi(44)&=48960,&\chi(45)&=-16611,\\
\chi(77)&=-89475,&\chi(78)&=26520.
\end{align*}
The intermediate value theorem gives one root in each of $(4,5)$, $(44,45)$ and $(77,78)$. Since $\chi$ is cubic, these are all roots, each simple. Denote them $E_0<E_4<E_8$; the subscripts will be justified by exact zero counts, not by numerical plots.

Write $t=x^2$ and $f_E(t)=49p_E(x)$. Expansion of \eqref{H08-p} gives
\begin{align}\label{H08-f}
f_E(t)={}&43808t^4-175232t^3+(350464-2072E)t^2\nonumber\\
&+(4144E-369408)t+49E^2-6216E+188256.
\end{align}
\subsection{Exact nodal certificates}\label{H08-sec-5}
\subsubsection*{The lowest root}
Putting $y=t-1$ rewrites \eqref{H08-f} as
\[43808y^4+(87616-2072E)y^2-18944y+49E^2-4144E+37888.
\]
For $4<E<5$, this is bounded below by
\[77256y^2-18944y+18393.
\]
The latter quadratic is strictly positive because its leading coefficient is positive and
\[18944^2<4\cdot77256\cdot18393.
\]
Thus $p_{E_0}$ has no real roots.
\subsubsection*{The middle root}
For $44<E<45$, direct substitution gives
\[f_E(0)>0,\qquad f_E(1/10)<0,\qquad f_E(3)>0.
\]
The endpoint pairs are\begin{align*}t=0:&\quad(9616,7761),\\t=1/10:&\quad(-4168182/625,-5081507/625),\\t=3:&\quad(599248,591177).\end{align*} Each expression is monotone on the indicated energy interval, as follows by differentiating its quadratic expression in $E$.

To exclude two additional real roots, the exact discriminant reduction is
\begin{equation}\label{H08-disc}
\operatorname{disc}_t f_E\equiv
41395941145477467607138304(49E^2-4088E+22080)\pmod{\chi(E)}.
\end{equation}
This is an identity of polynomial remainders and can be verified from the Sylvester determinant definition of the discriminant. For $44<E<45$, the parenthesis is at most
\[49\cdot45^2-4088\cdot44+22080=-58567<0.
\]
A real quartic with negative nonzero discriminant has exactly two real roots and one nonreal conjugate pair: this follows directly by pairing conjugate factors in the product-of-squared-differences formula. The signs above locate both real roots on the positive $t$-axis. Hence $p_{E_4}(x)$ has exactly four simple real roots.
\subsubsection*{The highest root}
For $77<E<78$, the five values at
\[t=0,\quad1/10,\quad1/2,\quad2,\quad3
\]
have signs $+,-,+,-,+$, respectively. Exact endpoint pairs are
\begin{align*}
&(145,1524),\quad(-1967907/625,-859982/625),\quad(3549,6482),\\
&(-37743,-36364),\quad(384649,379812).
\end{align*}
Again each quadratic is monotone on this interval. Thus there are at least four distinct positive roots. Degree four makes this exact. Consequently $p_{E_8}$ has eight simple real roots.

All three wavefunctions are $p_Ee^{-F}$ with
\[F=\frac{37}{14}\left(\frac{x^6}{6}-\frac{3x^4}{4}+\frac{3x^2}{2}\right).
\]
The gauge is strictly positive and decays rapidly. The standard Sturm oscillation theorem now identifies the levels as $0,4,8$.
\subsection{Exhausting the algebraic energy set}\label{H08-sec-6}
A rational Riccati solution for a polynomial potential has the form $G+Q'/Q$, with $G$ polynomial and $Q$ squarefree. Leading-coefficient comparison determines $G$ up to sign; the two candidate degrees of $Q$ sum to minus the degree of $G$, so at most one is nonnegative. The displayed states show that this admissible sign is $G=-A$ and degree is eight.

Any set of distinct additional algebraic energies would provide linearly independent degree-eight polynomial eigenfunctions. Their invariant span has degree sequence closed upward by four until eight, so dimension at most three. Our three states already attain it. Thus \eqref{H08-chi} gives the complete rational-logarithmic-derivative spectrum, not merely three selected exact levels.
\subsection{Status and literature}\label{H08-sec-7}
The differential-Galois and polynomial-spectral framework is established \cite{ref080}. General non-Lie-algebraic modules are also established \cite{ref088}. The precise inverse classification and fully certified three-state example here are candidate refinements; no exhaustive novelty certification is claimed. In particular we do not assert the first existence of a decatic QES potential or solve a named pre-existing open question.




\clearpage
\section[H09 --- A Complex Nonexistence Certificate for Four-State Even Decatic Modules]{A Complex Nonexistence Certificate for Four-State Even Decatic Modules}\label{contrib-H09}
\markboth{Polynomial and Rational Differential Spectra}{H09 --- Heisei}
\noindent\textsc{H09 \quad Heisei programme, source round 9}\par\medskip






\begin{quote}\small\textbf{Overview.} We prove that an even decatic Schr\"odinger gauge cannot have a four-dimensional polynomial invariant space containing the constant function. The obstruction holds over the complex numbers, not only in the self-adjoint regime. After normalizing the leading drift and shifting energy, the assertion reduces to two incompatible cubic-in-parameter conditions; a one-line polynomial remainder identity is the final nonexistence certificate. This separates an attainable three-state nonmonomial sector from the next proposed sector. The theorem excludes dimension four specifically and does not claim nonexistence of every larger or constant-free module.\end{quote}
\subsection{Precise assertion}\label{H09-sec-1}
Let
\[A=x^5+bx^3+cx,\quad B=-24x^4+dx^2,\quad
T=-D^2+2AD+B,\qquad b,c,d\in\C.
\]
The constant term of $B$ may be omitted since it only shifts all eigenvalues. The normalization of the leading coefficient of $A$ can be achieved over $\C$ by a nonzero dilation; in the real confining case a positive dilation suffices.
\begin{theorem}\label{H09-no4}
There is no four-dimensional polynomial space $W$ such that $1\in W$ and $T(W)\subseteq W$.
\end{theorem}
Equivalently, no even decatic polynomial Schr\"odinger potential has a four-dimensional constant-containing polynomial module after a common polynomial exponential gauge with odd quintic logarithmic derivative. The equivalence includes the necessary leading coefficient $-24$ in $B$ and does not assume that $W$ has a monomial basis.
\subsection{Why the displayed family is exhaustive for this question}\label{H09-sec-2}
For a gauged operator $-D^2+2AD+B$ with $\deg A=5$ and $\deg B\le4$, let $a=[x^5]A\ne0$ and $\beta=[x^4]B$. A polynomial $p$ of degree $k$ is sent to degree $k+4$ with coefficient $(2ak+\beta)\operatorname{lc}(p)$, unless this coefficient vanishes. A finite invariant polynomial space must therefore have maximum degree $n=-\beta/(2a)$.

Its degree set is closed by $k\mapsto k+4$ below $n$. Thus it is one uninterrupted arithmetic progression ending at $n$. If it contains $1$ and has dimension four, its degrees are exactly $0,4,8,12$, so $n=12$ and $\beta=-24a$. After normalizing $a=1$, the even coefficient structure is exactly that in the theorem. No other even quartic multiplication term is missed.
\subsection{The forced cyclic basis}\label{H09-sec-3}
Since $T1=B$, define
\[P=TB,\qquad Q=TP.
\]
The polynomials $1,B,P,Q$ have degrees $0,4,8,12$, respectively, and nonzero leading coefficients. They are independent, and any proposed four-dimensional space containing $1$ must therefore equal their span. Invariance is equivalent to
\begin{equation}\label{H09-closure}TQ=uQ+vP+wB+z\end{equation}
for constants $u,v,w,z\in\C$.

The proof below is a complete elimination of this finite coefficient identity. All divisions are by explicitly nonzero quantities, and the only nonzero parameter assertion is derived before it is used.
\subsection{Coefficient elimination}\label{H09-sec-4}
Comparing coefficients of $x^{14}$ and $x^{12}$ in \eqref{H09-closure} gives
\begin{equation}\label{H09-du}
d=-16b,\qquad u=-\frac{16}{3}(b^2-9c).
\end{equation}
After these substitutions, the coefficients at $x^{10},x^8,x^6,x^4,x^2$ vanish exactly when the following expressions vanish:
\begin{align}
F_{10}={}&28b^3-126bc-1431,\label{H09-f10}\\
F_8={}&496b^4-3288b^2c-22860b+6336c^2+9v,\nonumber\\
F_6={}&48b^5-496b^3c-2076b^2+1788bc^2+3bv-2322c,\nonumber\\
F_4={}&-1536b^4c+3072b^3+12160b^2c^2+24b^2v\nonumber\\
&-58752bc-7680c^3-24cv-3w-103968,\nonumber\\
F_2={}&-2816b^4-256b^3c^2+21120b^2c+2112bc^3\nonumber\\
&+12bcv+3bw+45792b-25056c^2-54v.\nonumber
\end{align}
The actual coefficients are nonzero scalar multiples of these expressions: in descending order the multiples are $-512$, $-128/3$, $-512/3$, $-8$, $16/3$. This information makes the reduction directly checkable by differentiating $B,P,Q$; the constant coefficient only determines $z$ and is not needed for contradiction.

Equation \eqref{H09-f10} forces $b\ne0$, and therefore
\begin{align*}
c&=\frac{28b^3-1431}{126b},\\
v&=-\frac{176(196b^6-43848b^3+2047761)}{3969b^2}.
\end{align*}
Substituting these into the actual $x^6$ coefficient gives
\begin{equation}\label{H09-first}
\frac{90112(14b^3+1431)}{49b}=0.
\end{equation}
Next $F_4=0$ determines
\[w=\frac{128(2744b^9-1534680b^6+66481884b^3-2930345991)}{83349b^3}.
\]
The actual $x^2$ coefficient now reduces to
\begin{equation}\label{H09-second}
-\frac{4096(7448b^6+903798b^3+10238805)}{3087b^2}=0.
\end{equation}
These are rational-function identities obtained by the displayed substitutions, so they require no numerical root-finding or unreported Gr\"obner basis.
\subsection{The final inconsistency certificate}\label{H09-sec-5}
Put $y=b^3$. Since $b\ne0$, equations \eqref{H09-first}--\eqref{H09-second} require
\[14y+1431=0,\qquad 7448y^2+903798y+10238805=0.
\]
But direct multiplication gives
\begin{equation}\label{H09-bezout}
7448y^2+903798y+10238805
=(532y+10179)(14y+1431)-4327344.
\end{equation}
Thus the first equation makes the second left side equal the nonzero rational integer $-4327344$. This contradiction proves Theorem \ref{H09-no4} over $\C$.

For an independent short check after \eqref{H09-first}, one may substitute $b^3=-1431/14$ and $c=b^2/3$ directly into the complete cyclic reduction. The remaining leakage is $-393216b x^2/7$, again nonzero. Identity \eqref{H09-bezout} is preferable because it exposes the exact algebraic obstruction without reducing powers implicitly.
\subsection{Interpretation and boundaries}\label{H09-sec-6}
The degree-filtration estimate alone permits dimension four when the polynomial factor degree is twelve and the drift degree is five. The theorem shows that this upper bound is not attainable by the constant-containing even decatic module of that degree. It is therefore a structural obstruction, not a numerical failure to locate parameters.

The next smaller case is genuinely nonempty. For comparison, a three-state example is obtained from
\[A=\frac{37}{14}(x^5-3x^3+3x),\qquad B=\frac{296}{7}(2x^2-x^4),
\]
whose module is $\Span\{1,B,TB\}$ and whose three energies are roots of
\[49E^3-6216E^2+196544E-738816.
\]
This comparator is included only to distinguish a nonexistence result from a vacuous class; the theorem's proof does not depend on it.

Several stronger statements remain unproved here. The theorem does not rule out a four-state module not containing $1$, a non-even quintic gauge, a rational rather than polynomial gauge, or a module of dimension at least five. In particular nonexistence in dimension four cannot be extrapolated to all larger dimensions: an invariant space containing $1$ need not have a smaller invariant subspace that still contains $1$.
\subsection{Literature status}\label{H09-sec-7}
Polynomial invariant-space and quasi-exact solvability methods are classical \cite{ref083,ref088}. The exact complex nonexistence certificate here is the submitted result. A targeted primary-source review did not certify whether this restricted obstruction was previously recorded. It is not presented as resolution of a named open problem.




\clearpage
\section[H10 --- Sharp Resonance-Chain Bounds for Arbitrary-Order Rational Differential Spectra]{Sharp Resonance-Chain Bounds for Arbitrary-Order Rational Differential Spectra}\label{contrib-H10}
\markboth{Polynomial and Rational Differential Spectra}{H10 --- Heisei}
\noindent\textsc{H10 \quad Heisei programme, source round 10}\par\medskip
\noindent\textit{Integration note.} The maximal-space theorem below incorporates the substantive scope correction: every member of a Jordan chain must remain polynomial. A rational-ambient counterexample is recorded after its proof.\par\medskip






\begin{quote}\small\textbf{Overview.} For an arbitrary-order rational differential operator with positive degree shift at infinity, we bound the dimension of every finite polynomial invariant space by a finite union of arithmetic resonance chains. This bounds the number of distinct polynomial eigenvalues and is sharper than using only the largest indicial root. The bound is optimal for every prescribed finite resonance set and every positive shift, with explicit eigenpolynomials for an extremal family. We also construct the unique maximal polynomial invariant space by a finite block-kernel calculation that handles finite poles and Jordan chains. Indicial degree bounds and invariant-space methods are classical; the exact chain count, universal sharpness construction, and unified certificate are the claims proved here.\end{quote}
\subsection{Degree data at infinity}\label{H10-sec-1}
Let
\[T=\sum_{j=0}^{q}a_j(x)D^j,\qquad a_j\in\C(x),\qquad D=\dd,
\]
be nonzero. For a nonzero rational function, define its degree at infinity as numerator degree minus denominator degree; set the degree of zero to $-\infty$. Suppose
\[s=\max_j(\deg_\infty a_j-j)>0.
\]
For each $j$ attaining the maximum, write
\[a_j(x)=\alpha_jx^{j+s}+O(x^{j+s-1}).\]
Define the leading multiplier polynomial
\begin{equation}\label{H10-h}
h(z)=\sum_{\deg_\infty a_j-j=s}\alpha_j(z)_j,
\qquad (z)_j=z(z-1)\cdots(z-j+1),\quad(z)_0=1.
\end{equation}
It is nonzero, since the falling factorials of different degrees are linearly independent. For every polynomial $p$ of degree $d$,
\begin{equation}\label{H10-lead}
Tp=h(d)\operatorname{lc}(p)x^{d+s}+O(x^{d+s-1})
\end{equation}
at infinity. When $h(d)\ne0$, the displayed degree is exact.

Let $H=\{n\in\mathbb Z_{\ge0}:h(n)=0\}$. This is finite. Define its downward resonance-chain set
\begin{equation}\label{H10-Aset}
\mathcal A_s(H)=\{d\in\mathbb Z_{\ge0}:d+ks\in H\text{ for some }k\ge0\}.
\end{equation}
For each residue $r\in\{0,\ldots,s-1\}$ that occurs in $H$, let $M_r$ be its largest member of that residue. Then
\begin{equation}\label{H10-count}
N_s(H):=|\mathcal A_s(H)|=
\sum_{r:\,M_r\text{ exists}}\left(\frac{M_r-r}{s}+1\right).
\end{equation}
Repeated roots of $h$ do not alter $H$ or this count.
\subsection{The resonance-chain bound}\label{H10-sec-2}
\begin{theorem}\label{H10-main}
Every finite-dimensional $T$-invariant subspace $W\subseteq\C[x]$ satisfies
\[\dim W\le N_s(H).
\]
More precisely, each degree occurring among nonzero members of $W$ belongs to $\mathcal A_s(H)$, and its degree set is closed under $d\mapsto d+s$ whenever $d\notin H$. If $H=\varnothing$, no nonzero finite polynomial invariant space exists. In particular
\[|\{\lambda\in\C:\ Tp=\lambda p\text{ for some }p\in\C[x]\setminus\{0\}\}|\le N_s(H).
\]
\end{theorem}
\begin{proof}
Coefficient elimination gives a basis of $W$ with distinct leading degrees. Its degree set $S$ is exactly the set of degrees of nonzero members of $W$: in a linear combination the largest basis degree with nonzero coefficient cannot cancel. Thus $|S|=\dim W$.

If $d\in S$ and $h(d)\ne0$, invariance makes $Tp$ a polynomial in $W$, and \eqref{H10-lead} forces $d+s\in S$. Repetition in a finite set must eventually hit $H$; otherwise degrees increase indefinitely. Hence every $d\in S$ lies in \eqref{H10-Aset}. This proves the dimension bound and the empty-resonance assertion.

Polynomial eigenvectors belonging to distinct eigenvalues are independent, and their span is invariant. Applying the dimension bound to every finite subset of the eigenvalue set proves the stated cardinality bound and finiteness.
\end{proof}
The result does not assert that the support exponents of every polynomial lie in $\mathcal A_s(H)$. It constrains leading degrees in the invariant space. Lower terms can have other exponents, a distinction crucial for nonmonomial examples.
\subsection{Universal optimality}\label{H10-sec-3}
\begin{theorem}\label{H10-sharp}
Given any finite nonempty set $H\subseteq\mathbb Z_{\ge0}$ and any integer $s\ge1$, there is a polynomial-coefficient differential operator of order $|H|$ whose leading multiplier has exactly the nonnegative integer roots $H$, and which has exactly $N_s(H)$ distinct polynomial eigenvalues. Every degree in the chain set occurs as an eigenvalue.
\end{theorem}
\begin{proof}
Let $h(z)=\prod_{n\in H}(z-n)$ and $\theta=xD$. Set
\[T_{H,s}=x^s h(\theta)+\theta.
\]
This has order $|H|$ and positive shift $s$, with leading multiplier $h$. It acts by
\[T_{H,s}x^d=h(d)x^{d+s}+d x^d.
\]
The monomial space $M=\Span\{x^d:d\in\mathcal A_s(H)\}$ is invariant: either $h(d)=0$ or $d+s$ is still in the chain set. In degree order, the matrix is triangular with pairwise distinct diagonal entries $d\in\mathcal A_s(H)$. Hence it has exactly $N_s(H)$ distinct eigenvalues. Theorem \ref{H10-main} excludes any further polynomial eigenvalues.
\end{proof}
There is an explicit certificate for every eigenpolynomial. For $d\in\mathcal A_s(H)$, let $\ell$ be the least nonnegative integer with $d+\ell s\in H$. Define
\begin{equation}\label{H10-eigen}
P_d(x)=\sum_{j=0}^{\ell}\frac{(-1)^j}{s^j j!}
\left(\prod_{i=0}^{j-1}h(d+is)\right)x^{d+js}.
\end{equation}
The empty product is one. Comparing the coefficient at $x^{d+js}$ gives
\[js\,c_j+h(d+(j-1)s)c_{j-1}=0\quad(j\ge1),
\]
and the final raising coefficient vanishes. Therefore $T_{H,s}P_d=dP_d$. Formula \eqref{H10-eigen} also shows explicitly why eigenpolynomials for different energies can have the same top degree.
\subsection{An exact maximal-space certificate with finite poles}\label{H10-sec-4}
Suppose $H\ne\varnothing$ and put $M=\max H$. Every finite invariant polynomial space has maximum degree in $H$, so it is contained in $\Pcal_M$.

Choose a polynomial $Q\ne0$ such that every $Qa_j$ is polynomial. For $p\in\Pcal_M$, take the polynomial part of the rational function $Tp$, then retain its coefficients through degree $M$; call the result $Lp$. This defines a linear endomorphism $L$ of $\Pcal_M$. Define the leakage map
\[Cp=Q(Tp-Lp).
\]
Its values are polynomials in a fixed finite-dimensional space. It records both the proper rational part and polynomial terms of degree greater than $M$; it is not merely a high-degree truncation.
\begin{theorem}\label{H10-algorithm}
The space
\begin{equation}\label{H10-kernel}
W_* =\bigcap_{j=0}^{M}\ker(CL^j)
\end{equation}
is the unique largest finite-dimensional polynomial invariant space for $T$. It contains every Jordan chain all of whose members are polynomials, and $\dim W_*\le N_s(H)$. If a basis matrix $U$ is computed and $TU=UK$, then the complete polynomial eigenvalue set is the root set of $\det(\lambda I-K)$.
\end{theorem}
\begin{proof}
Any finite invariant space lies in $\Pcal_M$ and has zero leakage at every iterate. On that space $T=L$, so it lies in \eqref{H10-kernel}. Conversely, Cayley--Hamilton for the $(M+1)$-dimensional operator $L$ shows that $CL^{M+1}$ is a linear combination of $CL^0,\ldots,CL^M$. Thus $L(W_*)\subseteq W_*$. Since $C(W_*)=0$ and $Q\ne0$, one has $T=L$ on $W_*$, proving invariance and maximality. A polynomial Jordan chain spans a finite invariant space, so it is included. The chain bound applies to $W_*$. Finally an eigenvector of $K$ gives a polynomial eigenfunction, and every polynomial eigenfunction lies in $W_*$.
\end{proof}
\begin{remark}[Corrected Jordan-chain scope]
If $(T-\lambda)^k p=0$ and $(T-\lambda)^j p\in\C[x]$ for every $0\le j<k$, the span of these iterates is a polynomial invariant space and is contained in $W_*$. Requiring only the initial vector to be polynomial is insufficient. For
\[
T=x^4D^2+(2x^3+1)D+1/x
\]
one has $T1=1/x$ and $T^2 1=0$ in $\C(x)$, while the maximal polynomial invariant space is zero. Indeed $s=2$, $h(z)=z(z+1)$, $M=0$, and with $Q=x$ the leakage map satisfies $C1=1$. This correction does not change the dimension bound, sharpness construction, or maximality certificate.
\end{remark}

Over an effective exact coefficient field, every step uses finite polynomial division and linear algebra. The method can be slow for large $M$; no polynomial-time complexity guarantee is asserted. Characteristic-polynomial multiplicity describes the finite restriction and must not be confused with the number of distinct eigenvalues.
\subsection{A sparse sharp example}\label{H10-sec-5}
Take $s=3$ and $H=\{2,7\}$. Then
\[\mathcal A_3(H)=\{1,2,4,7\},\qquad N_3(H)=4,
\]
whereas the crude largest-degree estimate gives eight. The extremal operator is
\begin{align*}
T&=x^3(\theta-2)(\theta-7)+\theta\\
 &=x^5D^2-8x^4D+14x^3+xD.
\end{align*}
Its complete polynomial spectrum and eigenpolynomials are
\begin{align*}
\lambda=1:&\quad P_1=x-2x^4-2x^7,\\
\lambda=2:&\quad P_2=x^2,\\
\lambda=4:&\quad P_4=x^4+2x^7,\\
\lambda=7:&\quad P_7=x^7.
\end{align*}
Direct differentiation checks every identity. This example also warns against bounding the number of polynomial eigenvalues by the differential order: the order is two but there are four distinct polynomial eigenvalues.
\subsection{Schr\"odinger specialization and limitations}\label{H10-sec-6}
For a polynomial gauge
\[T=-D^2-2A D+C,\qquad\deg A=m\ge2,
\]
with the highest multiplication coefficient chosen so that the leading multiplier is $2a(n-z)$, the resonance set is the singleton $H=\{n\}$ and $s=m-1$. The chain count becomes
\[N_s(H)=\left\lfloor\frac{n}{m-1}\right\rfloor+1.
\]
Thus the Schr\"odinger energy bound is one case of a sharp all-order theorem, while the multiple-resonance formula shows exactly what changes beyond a linear leading multiplier.

The hypothesis $s>0$ is essential. For $\theta=xD$, the shift is zero and every monomial is an eigenfunction, giving infinitely many energies. Finite poles are allowed in Theorem \ref{H10-algorithm}, but the theorem concerns polynomial eigenfunctions, not all rational, algebraic, or Liouvillian solutions. Passing to such classes requires separate denominator or gauge information.
\subsection{Literature and mathematical status}\label{H10-sec-7}
Degree bounds from indicial data are classical in rational-solution algorithms \cite{ref089}; monomial invariant differential operators were systematically studied by Post and Turbiner \cite{ref084}. This manuscript does not rebrand those general programs as new. Its exact statements are the resonance-chain dimension and spectral bound, its optimality for every finite resonance set, and the finite maximal-space certificate including poles and Jordan chains. All are proved above; exhaustive novelty and formal verification are not claimed.



\clearpage\part{Further Deductions and Sharper Boundaries}

\noindent The following results were developed while integrating the sixty contributions. Each is stated with a complete argument and a provenance discussion. They are not counted among the sixty source contributions. A claim to settle a pre-existing named problem is made only if its original formulation and the exact scope of the proof are explicitly identified.


\clearpage
\markboth{Further deductions}{Showa extension}
\section{A separate minimization transition for rank-two alpha-permanents}
\label{ext-showa-section}

The positivity threshold and the minimization threshold are different questions.
The preceding rank-two results determine when all alpha-permanents are
nonnegative. We now determine the strictly larger parameter range in which a
regular projective polygon minimizes the alpha-permanent, even when that
minimum is negative. At the new boundary the minimizer ceases to be unique:
an explicit continuous family of nonregular configurations attains the same
minimum. This extension also covers every exponent above one half, a range
not covered by the earlier minimization statement.

For $m\ge3$, let $\mathcal C_m$ be the real correlation matrices of rank at
most two, and put
\[
 G_m=\bigl(\cos(\pi(i-j)/m)\bigr)_{i,j=0}^{m-1},\qquad
 \operatorname{per}_{\alpha}(A)=
 \sum_{\sigma\in S_m}\alpha^{c(\sigma)}\prod_{j=1}^m A_{j,\sigma(j)}.
\]
Here $c(\sigma)$ is the number of cycles, and $(u)_m=u(u+1)\cdots(u+m-1)$
is the rising factorial. Matrix equivalence below means
simultaneous permutation and matching row-and-column sign switches.
Pate's sharp ordinary-permanent lower bound on this class is prior work
\cite{ext-showa-Pate}; the case $\alpha=1$ below recovers that bound and is
not presented as a new solution of that problem.

Define $\mu_3=\mu_4=\mu_5=0$. For $m\ge6$, let $\eta_m\in(0,1)$ solve
\begin{equation}\label{ext-showa-minroot}
 \eta_m(2-\eta_m)
 \sum_{j=1}^{\lfloor(m-2)/2\rfloor}
 \frac{\binom{m-2}{2j}}{(2j-\eta_m)(2j+2-\eta_m)}=1,
 \qquad \mu_m=\frac{1-\eta_m}{2}.
\end{equation}
The existence and uniqueness asserted in this definition are included in
the proof.

\begin{theorem}[Exact polygon-minimization range]
\label{ext-showa-minimum}
Fix $m\ge3$ and $\alpha>0$. The inequality
\begin{equation}\label{ext-showa-minineq}
 \operatorname{per}_{\alpha}(A)\ge
 \operatorname{per}_{\alpha}(G_m)\qquad(A\in\mathcal C_m)
\end{equation}
holds if and only if $\alpha\ge\mu_m$.
If $\alpha>\mu_m$, equality holds exactly for matrices equivalent to $G_m$.

Suppose $m\ge6$ and $\alpha=\mu_m$. Write
\[
 A_{ij}=\cos((\theta_i-\theta_j)/2),\qquad
 p_A(z)=\prod_{j=1}^m(z-e^{i\theta_j})
       =\sum_{k=0}^m c_kz^k.
\]
Then equality in \eqref{ext-showa-minineq} holds exactly when
$c_2=\cdots=c_{m-2}=0$. In particular, every polynomial
\begin{equation}\label{ext-showa-blaschke-family}
 p_b(z)=z^m-bz^{m-1}+\overline b\,z-1,
 \qquad |b|<1,
\end{equation}
has $m$ distinct unit-circle roots and gives an equality configuration.
For $b\ne0$ this configuration is not a regular projective polygon.

For $0<\alpha<\mu_m$, the same family with any $0<|b|<1$ gives a strict
counterexample to \eqref{ext-showa-minineq}.
\end{theorem}

Thus the theorem classifies exactly when the polygon is a global minimizer;
it does not classify all global minimizers below $\mu_m$.

\subsection{An exact gap formula at every positive exponent}

For a polynomial $f$ and a real $\delta<1$, define
\[
 \mathcal L_\delta(f)=f(0)-\delta\int_0^1
       (f(w)-f(0))w^{-1-\delta}\,dw.
\]
This integral converges. At $\delta=0$ it means $f(0)$. If $\delta<0$,
direct integration of the constant term gives the positive representation
\begin{equation}\label{ext-showa-positivefunctional}
 \mathcal L_\delta(f)=-\delta\int_0^1 f(w)w^{-1-\delta}\,dw.
\end{equation}
For $\alpha>0$, set $\delta=1-2\alpha$ and
\[
 C_m(\alpha)=\frac{(2\alpha)_m}{4^m},\qquad
 h_{m,k}(w)=(1+w)^{m-k}(1-w)^k+(1+w)^k(1-w)^{m-k},
\]
\begin{equation}\label{ext-showa-weights}
 d_{m,k}(\alpha)=\frac12\mathcal L_{1-2\alpha}(h_{m,k}).
\end{equation}
All these quantities are explicit finite sums: if
$h_{m,k}(w)/2=\sum_j a_jw^j$, then
\[
 d_{m,k}(\alpha)=a_0-(1-2\alpha)
        \sum_{j\ge1}\frac{a_j}{j-1+2\alpha}.
\]

\begin{lemma}[Exact gap identity]
\label{ext-showa-gap}
For every $\alpha>0$ and $A\in\mathcal C_m$,
\begin{equation}\label{ext-showa-gapformula}
 \operatorname{per}_{\alpha}(A)-\operatorname{per}_{\alpha}(G_m)
 =C_m(\alpha)\sum_{k=1}^{m-1}d_{m,k}(\alpha)|c_k|^2.
\end{equation}
For $\alpha\ge1/2$ all interior weights are strictly positive. For
$0<\alpha<1/2$, their minimum is $d_{m,1}=d_{m,m-1}$; when
$2\le k\le m-2$, it is strictly smaller than $d_{m,k}$.
\end{lemma}

\begin{proof}
We provide the coefficient calculation to fix the normalization and the
range of the integral argument. The squarefree determinant identity is
\[
 [z_1\cdots z_m]\det(I-ZA)^{-\alpha}
       =\operatorname{per}_{\alpha}(A),\qquad Z=\diag(z_1,\ldots,z_m).
\]
Indeed, expanding the exponential of
$\alpha\sum_{\ell\ge1}\tr((ZA)^\ell)/\ell$ enumerates disjoint
permutation cycles. Write $A=W^TW$ with unit columns
$w_j=(\cos(\theta_j/2),\sin(\theta_j/2))$. Sylvester's determinant
identity turns this coefficient into
$(-1)^mD_{v_1}\cdots D_{v_m}Q^{-\alpha}(1,0,0)$, where
$Q=t^2-x^2-y^2$ and
$v_j=(1,\cos\theta_j,\sin\theta_j)/2$.

Initially let $\alpha>1/2$. Take a disk variable $V$ with density
$(\alpha-1/2)\pi^{-1}(1-\|v\|^2)^{\alpha-3/2}$ and, independently,
a gamma variable $T$ of shape $2\alpha$ and rate one. The beta moments
$\mathbb E V_1^{2j}=(1/2)_j/(\alpha+1/2)_j$ show by the binomial series
that
\[
 Q(t,x,y)^{-\alpha}
   =\mathbb E\exp\{-T(t-1-xV_1-yV_2)\}.
\]
Differentiation therefore gives
$(2\alpha)_m2^{-m}\mathbb E\prod_j(1-u_j\cdot V)$, with
$u_j=(\cos\theta_j,\sin\theta_j)$.
Writing $w=\sqrt{1-r^2}$ and $\rho=r/(1+w)$, factor
$1-r\cos\phi=(1+w)|1-\rho e^{i\phi}|^2/2$.
Fourier orthogonality makes the angular average equal to
\[
 2^{-m}g(w),\qquad
 g(w)=\sum_{k=0}^m|c_k|^2(1+w)^{m-k}(1-w)^k.
\]
The radial variable $w$ has density $(2\alpha-1)w^{2\alpha-2}$.
Consequently the permanent is $C_m(\alpha)\mathcal L_{1-2\alpha}(g)$.

For a polynomial of degree at most $m$, every denominator
$2\alpha+j-1$, $1\le j\le m$, in this expression divides
$(2\alpha)_m$. Thus the resulting identity is polynomial in $\alpha$
and extends from $\alpha>1/2$ to every $\alpha>0$. This extends an
algebraic identity, not the positive disk measure.

Unit-circle roots imply $|c_k|^2=|c_{m-k}|^2$ and
$|c_0|^2=|c_m|^2=1$. Pairing these coefficients yields
\[
 \operatorname{per}_{\alpha}(A)
 =C_m(\alpha)\left(2d_{m,0}(\alpha)
       +\sum_{k=1}^{m-1}d_{m,k}(\alpha)|c_k|^2\right).
\]
For the regular polygon all interior coefficients vanish, proving
\eqref{ext-showa-gapformula}.

For $\alpha>1/2$, positivity follows from
\eqref{ext-showa-positivefunctional}, since $h_{m,k}>0$ on $(0,1)$.
At $\alpha=1/2$ every weight is one. For $0<\alpha<1/2$, put
$x=1+w>y=1-w>0$. The function
$k\mapsto x^{m-k}y^k+x^ky^{m-k}$ is symmetric and strictly convex,
and decreases from $0$ to $m/2$. Hence
$h_{m,1}\ge h_{m,k}$ for every interior $k$, with strict inequality
for $2\le k\le m-2$ and $0<w<1$. These functions all equal two at
zero. The defining integral of $\mathcal L_\delta$, with $\delta>0$,
reverses this comparison and proves the claims about the least weight.
\end{proof}

\subsection{The scalar transition and its equality family}

The least weight admits a particularly simple formula. Since
\[
 \frac12h_{m,1}(w)
  =(1-w^2)\sum_{j=0}^{\lfloor(m-2)/2\rfloor}\binom{m-2}{2j}w^{2j},
\]
direct evaluation gives, for $0<\delta<1$,
\begin{equation}\label{ext-showa-leastweight}
 d_{m,1}\bigl((1-\delta)/2\bigr)
 =\frac2{2-\delta}\left[1-\delta(2-\delta)
  \sum_{j=1}^{\lfloor(m-2)/2\rfloor}
  \frac{\binom{m-2}{2j}}{(2j-\delta)(2j+2-\delta)}\right].
\end{equation}
Each summand multiplied by $\delta(2-\delta)$ is strictly increasing
on $(0,1)$: the numerator increases and the two positive denominator
factors decrease. At $\delta=1$ the sum in square brackets after the
initial one is $0$, $1/3$, and $1$ for $m=3,4,5$, respectively.
For $m=6$ it is $31/15>1$, and it increases strictly with $m$.
This proves existence and uniqueness in \eqref{ext-showa-minroot}.
It also proves that every interior weight is positive if
$\alpha>\mu_m$, and that at $\alpha=\mu_m>0$ exactly the endpoint
interior weights $d_{m,1},d_{m,m-1}$ vanish. The gap identity proves
the sufficiency and equality assertions of
Theorem~\ref{ext-showa-minimum}, apart from construction of its family.

\begin{lemma}[Unit-circle family]
\label{ext-showa-blaschke}
For $|b|<1$, all zeros of the polynomial
\eqref{ext-showa-blaschke-family} are distinct and lie on the unit circle.
The associated matrix $A_b\in\mathcal C_m$ satisfies, for every $\alpha>0$,
\begin{equation}\label{ext-showa-exactfamilygap}
 \operatorname{per}_{\alpha}(A_b)-\operatorname{per}_{\alpha}(G_m)
       =2C_m(\alpha)d_{m,1}(\alpha)|b|^2.
\end{equation}
\end{lemma}
\begin{proof}
On the unit circle the rational function
\[
 B_b(z)=z^{m-1}\frac{z-b}{1-\overline b z}
\]
has modulus one. Its argument derivative at $z=e^{it}$ is
\[
 m-1+\frac{1-|b|^2}{|e^{it}-b|^2}>0.
\]
Its winding number is $m$, so $B_b(z)=1$ has exactly $m$ distinct
unit-circle solutions. Its denominator is nonzero there, and this
equation is exactly $p_b(z)=0$. These roots supply the doubled Gram
angles of $A_b$. The only nonzero interior coefficients are $-b$ and
$\overline b$; substitution in \eqref{ext-showa-gapformula} proves
\eqref{ext-showa-exactfamilygap}.
\end{proof}

For $m\ge6$ and $0<\alpha<\mu_m$, formula
\eqref{ext-showa-leastweight} is negative. The last lemma therefore
supplies strict counterexamples and completes the necessity proof of
Theorem~\ref{ext-showa-minimum}. At the boundary it supplies the asserted
continuous equality family. No optimization over arbitrary matrices is
left unresolved in deciding whether the polygon itself is a global
minimizer.

\subsection{Separation from positivity and asymptotic size}

\begin{theorem}[Two different thresholds]
\label{ext-showa-asymptotic}
The sequence $\mu_m$ is strictly increasing for $m\ge5$, tends to
$1/2$, and satisfies
\begin{equation}\label{ext-showa-muasym}
 \frac12-\mu_m\sim\frac{m^2}{2^{m+1}}.
\end{equation}
Its first nonzero values are
\[
 \mu_6=\frac{\sqrt6-2}{2},\qquad
 \mu_7=\frac{\sqrt{19}-3}{4}.
\]
Let $a_m=(1-\delta_m)/2$ denote the positivity threshold defined by
\[
 \delta_m\sum_{j=1}^{\lfloor m/2\rfloor}
       \frac{\binom m{2j}}{2j-\delta_m}=1.
\]
Then $\mu_m<a_m$ for every $m\ge6$. Thus on
$[\mu_m,a_m)$ the polygon is a global minimizer with negative
alpha-permanent, whereas at $a_m$ its minimum first becomes zero.
\end{theorem}

\begin{proof}
For fixed $\delta\in(0,1)$ the left side of
\eqref{ext-showa-minroot} increases strictly with $m$; its unique root
therefore decreases, proving monotonicity.
Put
\[
 T_m(0)=\sum_{j=1}^{\lfloor(m-2)/2\rfloor}
       \frac{\binom{m-2}{2j}}{(2j)(2j+2)}.
\]
For $N\sim\operatorname{Bin}(m-2,1/2)$,
\[
 \frac{m^2T_m(0)}{2^{m-1}}
  =\frac12\mathbb E\left[
    \frac{m^2}{N(N+2)}\mathbf1_{\{N>0,\ N\text{ even}\}}\right]
  \longrightarrow1.
\]
Here the even probability is $1/2$, and the displayed ratio inside the
expectation tends uniformly to four on a shrinking relative central
region. To justify uniform integrability, put $n=m-2$. For
$N\ge n/4$ the ratio is bounded, so concentration of $N/n$ suffices.
For $0<N<n/4$ it is at most $m^2/3$, while the elementary exponential
Markov estimate
\[
 \mathbb P(N\le n/4)
 \le3^{n/4}\mathbb E3^{-N}
 =\left(\frac2{3^{3/4}}\right)^n
\]
is exponentially small. The event $N=0$ is omitted. This proves
$T_m(0)\sim2^{m-1}/m^2$.

Writing $T_m(\eta)$ for the sum in \eqref{ext-showa-minroot}, one has
\[
 T_m(0)\le T_m(\eta)
 \le(1-\eta/2)^{-2}T_m(0),\qquad 0\le\eta<1.
\]
The root equation gives $\eta_m\le1/T_m(0)\to0$ and consequently
$\eta_m\sim1/(2T_m(0))\sim m^2/2^m$. This proves
\eqref{ext-showa-muasym}. Substitution in the finite equation for
$m=6,7$ gives the stated radicals.

For comparison with $a_m$, write
$S_m(\delta)=\sum_{j\ge1}\binom m{2j}/(2j-\delta)$.
The identity
\[
 h_{m,0}-h_{m,1}
  =2w\bigl((1+w)^{m-1}-(1-w)^{m-1}\bigr)
\]
and $\delta_mS_m(\delta_m)=1$ give
\[
 d_{m,1}(a_m)
 =\frac{2\delta_m}{m}
   \left(\sum_{j\ge1}\binom m{2j}
      +\delta_mS_m(\delta_m)\right)
 =\frac{\delta_m2^m}{m}>0.
\]
The unique zero of the least weight therefore occurs at a strictly
smaller exponent. Finally
$\operatorname{per}_{\alpha}(G_m)=
2C_m(\alpha)[1-(1-2\alpha)S_m(1-2\alpha)]$
is negative for $0<\alpha<a_m$, proving the last assertion.
\end{proof}

At $\alpha=1$ the exact polygon value simplifies to
$\operatorname{per}(G_m)=m!/2^{m-1}$: in the positive integral,
$\int_0^1 h_{m,0}(w)\,dw=2^{m+1}/(m+1)$.
The contribution here is the full positive-exponent minimization
transition, its nonregular boundary equality family, and its distinct
$m^2 2^{-m}$ scale, not the previously known ordinary-permanent bound.
Neither a universal fourth-order theorem for all real-stable polynomials
nor a higher-rank permanent classification follows from this argument.
Historical priority of the positive-exponent refinement has not been
certified by an exhaustive literature search.



\clearpage
\markboth{Further deductions}{Meiji extension}
\section{Budget-optimal cap design and the strict cost of fixing a protocol}
\label{ext-meiji-section}

The preceding random-cap results fix the observation law before taking
the sample-size limit. We now optimize that law under an expected budget,
and prove that no single finite-mean law attains the resulting boundary
rate. The distinction between a fixed protocol and a sequence of designed
protocols is essential. The multilevel allocation and testing methods are
classical; the assertions below are additional theorems within the present
model, with historical priority unverified. They are not represented as
solutions of a named published open problem.

\subsection{The observation-design problem}

There are $n\ge1$ independent latent probabilities $p_i\sim\mu$, where
$\mu\in\PP([0,1])$. Before any coin outcomes are seen, choose a random vector
of positive integer caps $(L_1,\ldots,L_n)$, independent of all latent
probabilities and of their underlying coin sequences. The caps may be
dependent on one another and need not have the same marginal law. Their
joint law is known and every cap is observed. Conditional on the caps and
latent probabilities, the data are independent binomial counts
$S_i\sim\mathrm{Bin}(L_i,p_i)$. Recording complete words is equivalent,
since their conditional ordering given the counts is parameter-free.

One initial observation per individual is compulsory. Let $b>0$ be the
total expected budget for additional observations, and impose
\[
 \sum_{i=1}^n\E(L_i-1)\le b.
\]
Write $\mathfrak R(n,b)$ for the infimum, over all these designs and all
measurable probability-valued estimators, of their worst-case expected
$W_1$ error over $\PP([0,1])$. The design is allowed to depend on $n,b$.

\begin{theorem}[Exact sample--budget tradeoff]
\label{ext-meiji-budget}
There are numerical constants $0<c<C<\infty$ such that
\[
 c\min\{1,\max(n^{-1/2},b^{-1/3})\}
 \le \mathfrak R(n,b)
 \le C\min\{1,\max(n^{-1/2},b^{-1/3})\}.
\]
The upper bound can be attained, up to its numerical constant, using
independent identically distributed caps. It remains true if these caps
are required to have unbounded support, so that each designed experiment
is identifiable. The lower bound permits arbitrary dependence among the
caps and arbitrary nonidentical cap marginals, subject to the stated
independence from the latent variables and coin outcomes.
\end{theorem}

Thus a mean cap budget $\E L\le1+B$ gives the rate
$\min\{1,\max(n^{-1/2},(nB)^{-1/3})\}$. This theorem concerns an expected
budget and a design chosen without observing outcomes; it makes no claim
about outcome-adaptive sampling or an almost-sure budget constraint.

\subsection{A local testing bound that charges only additional observations}

Use squared Hellinger distance
$\mathsf h^2(P,Q)=\int(\sqrt{dP}-\sqrt{dQ})^2$. For
$x\in[1/4,3/4]$ and $0<u\le1/8$, consider
\[
 \tau_{x,u,-}=\delta_x,
 \qquad
 \tau_{x,u,+}=\tfrac14\delta_{x-u}
                 +\tfrac12\delta_x+\tfrac14\delta_{x+u}.
\]
They have the same mean and $W_1(\tau_{x,u,+},\tau_{x,u,-})=u/2$.

\begin{lemma}[One-cell information cost]
\label{ext-meiji-cell}
Let $P_{\ell,+},P_{\ell,-}$ be their laws for $\ell$ Bernoulli observations.
For a numerical constant $C$,
\[
 \mathsf h^2(P_{1,+},P_{1,-})=0,
 \qquad
 \mathsf h^2(P_{\ell,+},P_{\ell,-})
 \le C\min\{(\ell u^2)^2,1\}\quad(\ell\ge2).
\]
In particular, for every integer $\ell\ge1$,
\[
 \mathsf h^2(P_{\ell,+},P_{\ell,-})
 \le C(\ell-1)u^2.
\]
\end{lemma}
\begin{proof}
Use the Bernoulli$(x)$ word law $U_\ell$ as reference. The likelihood ratio
of one Bernoulli$(x+v)$ observation is $1+vZ$, where $\E Z=0$ and
$\E Z^2=[x(1-x)]^{-1}\le16/3$. Products of distinct subsets of the
independent variables $Z_i$ are orthogonal in $L^2(U_\ell)$. The signed
mixing difference has vanishing moments of degrees zero and one; every
centered moment of degree $k\ge2$ has magnitude at most $u^k$. Expanding
the likelihood products therefore gives
\[
 \int\left(\frac{dP_{\ell,+}}{dU_\ell}
            -\frac{dP_{\ell,-}}{dU_\ell}\right)^2dU_\ell
 \le\sum_{k=2}^{\ell}\binom{\ell}{k}(16u^2/3)^k.
\]
When $\ell u^2\le1$, the right side is at most $C(\ell u^2)^2$, by
$\binom{\ell}{k}\le\ell^k/k!$. Both mixture laws dominate
$U_\ell/2$, so
$(\sqrt a-\sqrt b)^2\le(a-b)^2/(a+b)$ bounds squared Hellinger distance
by this display. When $\ell u^2>1$, use $\mathsf h^2\le2$. The first-bit
laws coincide because the means coincide. Finally,
$\min\{z^2,1\}\le z$ for $z\ge0$, and
$\ell\le2(\ell-1)$ for $\ell\ge2$.
\end{proof}

\begin{proof}[Lower bound in Theorem~\ref{ext-meiji-budget}]
Choose $u\le1/32$, let $K=\lfloor1/(8u)\rfloor$, and put
$x_j=1/4+(4j-2)u$ for $1\le j\le K$. The cells
$[x_j-u,x_j+u]$ are disjoint and contained in $[1/4,3/4]$. For
$\theta\in\{-,+\}^K$, set
\[
 \mu_\theta=\frac1K\sum_{j=1}^K\tau_{x_j,u,\theta_j}.
\]
Every cell has the same mass under all hypotheses. The signed CDF from
each cell vanishes outside it, and consequently
\[
 W_1(\mu_\theta,\mu_{\theta'})
       =\frac{u}{2K}d_H(\theta,\theta').
\]
For neighboring vertices, additionally reveal the latent cell label of
each individual. Given the cap vector, one individual's squared
Hellinger distance in this enlarged experiment is at most
$C(L_i-1)u^2/K$ by Lemma~\ref{ext-meiji-cell}. Conditional independence,
the product inequality for squared Hellinger distance, and averaging over
the common cap-vector law yield
\[
 \mathsf h^2(\text{neighboring full data laws})
 \le \frac{Cu^2}{K}\sum_i\E(L_i-1)
 \le C'b u^3.
\]
This argument allows dependent caps: conditional on their entire vector,
the individuals are independent, and the cap-vector law itself is the
same under both hypotheses. Discarding the revealed labels can only
decrease distance.

Take $u=c_0\min\{1,b^{-1/3}\}$, with a sufficiently small numerical
$c_0$. Then $K\ge1/(16u)$ and every neighboring total variation distance
is at most $1/2$. Given an arbitrary estimator, select a nearest member
of the finite family $\{\mu_\theta\}$. Triangle inequality bounds its
original loss below by $u/(4K)$ times the Hamming loss of this selection.
Under the uniform prior on the cube, every bit has error probability at
least $(1-\TV)/2\ge1/4$. Summing gives risk at least
$u/16\ge c\min\{1,b^{-1/3}\}$.

There is also a design-independent lower bound $c n^{-1/2}$. On the
submodel $\mu_\vartheta=\vartheta\delta_1+(1-\vartheta)\delta_0$, the first
bit reveals the latent endpoint and all subsequent bits repeat it. The
cap vector is ancillary. Thus this subexperiment is exactly $n$
Bernoulli$(\vartheta)$ observations. The alternatives
$\vartheta=1/2\pm c_1n^{-1/2}$ have $W_1$ separation $2c_1n^{-1/2}$ and
product total variation bounded away from one, by the elementary
Bernoulli chi-squared product formula. Combining these two lower bounds
proves the assertion.
\end{proof}

\subsection{An explicit budget-feasible multilevel design}

\begin{proof}[Upper bound in Theorem~\ref{ext-meiji-budget}]
The bound $W_1\le1$ handles $b<1$ and every case in which
$T_*:=\min\{n,b^{2/3}\}<2$. Otherwise choose $J\ge1$ so that
$T=2^J\le T_*<2T$, and write $B=b/n$, $r_j=2^j$. For a fixed sufficiently
small numerical $a_0>0$, prescribe the dyadic survival probabilities
\[
 q_j^0=\min\{1,a_0B T^{-1/4}r_j^{-3/4}\},\qquad 1\le j\le J.
\]
This decreasing sequence defines a cap $L^0$ supported on
$1,2,\ldots,2^J$ with only dyadic values assigned positive mass:
$\Pr(L^0=1)=1-q_1^0$, $\Pr(L^0=2^j)=q_j^0-q_{j+1}^0$ for $j<J$, and
$\Pr(L^0=2^J)=q_J^0$. Since
\[
 \E(L^0-1)
 =\sum_{j=1}^J(r_j-r_{j-1})q_j^0
 \le a_0B T^{-1/4}\sum_{j=1}^Jr_j^{1/4}\le B/2
\]
for sufficiently small $a_0$, the budget has slack.

If unbounded support is desired, let $G$ be independent with
$\Pr(G=\ell)=2^{-\ell}$ for $\ell\ge1$, and use the cap-law mixture
$(1-\delta)\mathcal L(L^0)+\delta\mathcal L(G)$, where
$\delta=\min\{1/2,B/4\}$. Its mean excess is at most $B/2+\delta\le B$,
its support is unbounded, and its actual dyadic survivals satisfy
$q_j\ge q_j^0/2$. Independent copies of this law give the desired
identically distributed design. The following argument also works
without the optional mixture, in which case $q_j=q_j^0$.

For one individual let $U_j$ be its first-$r_j$ sample proportion, and set
\[
 D_0(x)=\mathbf1\{U_0\le x\},\qquad
 D_j(x)=\mathbf1\{U_j\le x\}-\mathbf1\{U_{j-1}\le x\}\quad(j\ge1).
\]
The observable random function
\[
 Y_J(x)=D_0(x)+\sum_{j=1}^J
           \frac{\mathbf1\{L\ge r_j\}}{q_j}D_j(x)
\]
has expectation $F_{U_J}(x)$. The coupling with the latent $p$ gives
$\int|F_{U_J}-F_\mu|\le(2\sqrt T)^{-1}$ and
$\E\int|D_j|\le C r_j^{-1/2}$. For $1\le j\le k$, nested cap events
and cap independence give
\[
 \E\int\left|
 \frac{\mathbf1\{L\ge r_j\}}{q_j}D_j
 \frac{\mathbf1\{L\ge r_k\}}{q_k}D_k\right|
 \le Cq_j^{-1}r_k^{-1/2}.
\]
Summing first over $k\ge j$ and treating base-level terms separately
proves
\[
 \int\Var Y_J\le C\left(1+\sum_{j=1}^J\frac{r_j^{-1/2}}{q_j}\right)
 \le C\left(1+\frac{T^{1/2}}{B}\right).
\]
Indeed, $q_j^{-1}\le2[1+a_0^{-1}B^{-1}T^{1/4}r_j^{3/4}]$, and the
two resulting geometric sums are bounded by a constant and a constant
times $T^{1/4}$, respectively. No independence across levels is used.

Average $Y_J$ over individuals and project approximately in $L^1$ onto
probability CDFs, within $1/n$ of the infimum. A measurable projection
is obtained by enumerating finite rationally supported and rationally
weighted probability measures and choosing the first within $1/n$ of
the countable infimum. Their $W_1$ density makes this infimum equal to
the unrestricted one. Cauchy--Schwarz and triangle inequality give
\[
 \sup_\mu\E W_1(\widehat\mu,\mu)
 \le C\left(T^{-1/2}+n^{-1/2}+\sqrt{\frac{T^{1/2}}b}\right)+n^{-1}.
\]
If $T_* =b^{2/3}$ this is at most
$C(n^{-1/2}+b^{-1/3})$; if $T_*=n$, then $b\ge n^{3/2}$ and it is at
most $C n^{-1/2}$. The sum and maximum differ by at most a factor two.
For count-only data, condition $Y_J$ on $(L,S)$: hypergeometric averaging
preserves expectation and decreases integrated variance. This completes
the upper bound for the stated observation model.
\end{proof}

\subsection{No fixed finite-mean cap attains the budget boundary}

\begin{theorem}[Strict fixed-protocol lower rate]
\label{ext-meiji-fixed}
Fix one cap law, used independently and identically for all sample sizes,
and let $R_n(L)$ be its minimax expected $W_1$ risk over $\PP([0,1])$.
If $\E L<\infty$, then
\[
                  n^{1/3}R_n(L)\longrightarrow\infty.
\]
More generally, for every $\alpha\ge1/2$,
\[
 \E L^\alpha<\infty
 \quad\Longrightarrow\quad
 n^{1/(2\alpha+1)}R_n(L)\longrightarrow\infty.
\]
Both conclusions hold whether or not the cap has bounded support. For
unbounded support the experiment is identifiable and uniformly consistent;
the statements concern its unavoidable rate of recovery.
\end{theorem}
\begin{proof}
For the finite-mean assertion, Lemma~\ref{ext-meiji-cell} gives the local
one-individual squared Hellinger bound
\[
 \psi(u)=C\E\min\{(Lu^2)^2,1\}=o(u^2).
\]
To justify the little-oh statement, divide by $u^2$: the integrand tends
to zero at every fixed $L$ and is bounded by $L$, since
$\min\{z^2,1\}\le z$. Dominated convergence applies.
The same disjoint-cell hypercube used in the budget lower bound then has
neighboring $n$-sample squared Hellinger distance at most
$n\psi(u)/K\le Cn u\psi(u)=o(nu^3)$.
For any fixed $A>0$, choose $u=A n^{-1/3}$. For all sufficiently large
$n$ the cells are admissible and neighboring total variation is at most
$1/2$, giving $R_n(L)\ge c A n^{-1/3}$ with a numerical $c>0$ independent
of $A$. Since $A$ is arbitrary, the first limit follows.

For the general statement, choose an integer $r>\alpha$. Let
$\rho_+,\rho_-$ be the positive and negative parts of
\[
 2^{1-r}\sum_{j=0}^r(-1)^j\binom{r}{j}\delta_{-1+2j/r},
\]
and put $\tau_\pm=\tfrac12\delta_0+\tfrac12\rho_\pm$. These are
probabilities with matching moments through degree $r-1$, a shared
half-mass at zero, and $W_1(\tau_+,\tau_-)=1/r$. The last identity follows
from the cumulative alternating binomial sum
$\sum_{i=0}^j(-1)^i\binom ri=(-1)^j\binom{r-1}j$ and integration over the
grid intervals. After translation to $x\in[1/4,3/4]$ and scaling by $u$,
the word-likelihood expansion in Lemma~\ref{ext-meiji-cell}, now starting
at degree $r$, gives
\[
 \mathsf h^2(P_{\ell,+},P_{\ell,-})
       \le C_r\min\{(\ell u^2)^r,1\}.
\]
The shared atom again supplies the positive reference denominator.
For $z\ge0$ one has $\min\{z^r,1\}\le z^\alpha$.
Since $r>\alpha$, dominated convergence gives
\[
 \E\min\{(Lu^2)^r,1\}=o(u^{2\alpha}).
\]
The cell distance is $u/r$. The squared Hellinger distance between
neighboring full-data laws is $o(nu^{2\alpha+1})$.
For each fixed $A>0$ set $u=A n^{-1/(2\alpha+1)}$.
The same testing argument gives, eventually,
\[
 R_n(L)\ge c_r A n^{-1/(2\alpha+1)}.
\]
Letting $A$ grow proves the stated divergence.
\end{proof}

For every fixed $B>0$, Theorem~\ref{ext-meiji-budget} permits an
identifiable, sample-size-dependent cap sequence with mean at most
$1+B$ and risk $O((nB)^{-1/3})$. Theorem~\ref{ext-meiji-fixed} proves
that no single finite-mean cap realizes this boundary rate. It does not
contradict the fixed polynomial-tail rates: tails of exponent
$\alpha>1$ have finite mean and rate $n^{-1/(2\alpha+1)}$, which is
strictly slower than $n^{-1/3}$.

\paragraph{Scope and antecedents.}
The population-mixture problem and moment-based methods have established
antecedents\cite{ref014,ref015}; multilevel CDF estimation
and allocation are established methodology\cite{ref020}.
The fixed-law covariance estimate and local testing constructions in the
earlier chapters are reused transparently. The additional conclusions are
the exact two-resource design theorem, its lower bound under dependent
preassigned caps, and the strict separation between fixed finite-moment
protocols and sample-size-dependent designs. Related total-budget claims
were mentioned in the supplied benchmark, whose full content is not
assumed or certified here. The unrestricted irregular-cap root-$n$
criterion on all mixing probabilities remains unresolved by this section.


\clearpage
\markboth{Further deductions}{Keio extension}
\section{Inverse classification from two cluster scales and reconstruction from nonuniform sizes}
\label{ext-keio-section}

The complete two-scale classification proved earlier admits an inverse theorem
which is stronger than the inverse statement for uniform blowups.  In this
section the bipartition is ordered: a bipartite tree is written
$T=(U,V;E)$, and $T[a,b]$ replaces vertices of $U$ by independent clusters of
size $a$ and vertices of $V$ by independent clusters of size $b$.  All trees
have at least two vertices, and all cluster sizes are at least two.  Write
\[
 r_T=|U|,\qquad s_T=|V|,\qquad \nu_T=\nu(T),\qquad
 \mathcal D_U(T)=\bigoplus_{u\in U}\Z/d_u\Z,\qquad
 \mathcal D_V(T)=\bigoplus_{v\in V}\Z/d_v\Z.
\]
Factors of order one are omitted.  These degree groups generally contain less
information than the corresponding degree multisets.  The results below
concern abstract critical groups; no identification of monodromy pairings or
distinguished divisor classes is asserted.

\subsection{The complete invariant for all two-scale groups}

We first make explicit how the forward classification depends on the two
degree groups.  For a finite abelian group $H$, a prime $q$, and $j\geq1$, let
$m_{q,j}(H)$ be the multiplicity of $\Z/q^j\Z$ in its elementary-divisor
decomposition.  These multiplicities are uniquely determined by $H$.

\begin{lemma}[Primary multiplicity formula]
\label{ext-keio-profile}
Fix $a,b\geq2$ and a prime $q$.  Put
$\alpha=v_q(a)$, $\beta=v_q(b)$, $m=r_T+s_T$, and
$\delta=m-2\nu_T$.  Define
\[
 P_U(X)=\sum_{u\in U}X^{v_q(d_u)},\qquad
 P_V(X)=\sum_{v\in V}X^{v_q(d_v)}.
\]
If $[F(X)]_{>0}$ means deletion of the constant term, then
\begin{align}
 \sum_{j\geq1}m_{q,j}(K(T[a,b]))X^j
 =\Bigl[{}&(a-2)X^\beta P_U(X)+(b-2)X^\alpha P_V(X)
       \notag\\
 &+2(r_T-\nu_T)X^\beta+2(s_T-\nu_T)X^\alpha
       \notag\\
 &+X^{\alpha+\beta}\bigl(P_U(X)+P_V(X)-1-\delta\bigr)
 \Bigr]_{>0}.                                             \label{ext-keio-profile-eq}
\end{align}
Consequently all the groups $K(T[a,b])$ are determined by
$(r_T,s_T,\nu_T,\mathcal D_U(T),\mathcal D_V(T))$.
\end{lemma}

\begin{proof}
We use the complete two-scale cyclic decomposition established in the forward
classification.  If $\ell$ is the number of leaves, its factors are
\[
 \begin{gathered}
 (\Z/(bd_u)\Z)^{a-2}\ (u\in U),\qquad
 (\Z/(ad_v)\Z)^{b-2}\ (v\in V),\\
 (\Z/b\Z)^{2(r_T-\nu_T)},\qquad
 (\Z/a\Z)^{2(s_T-\nu_T)},\qquad
 (\Z/(ab)\Z)^{\ell-1-\delta},\qquad
 \Z/(abd_i)\Z\ (d_i>1).
 \end{gathered}
\]
Taking $q$-valuations gives the first two lines of
\eqref{ext-keio-profile-eq}.  The final two kinds of factors contribute
\[
 X^{\alpha+\beta}
 \left(\ell-1-\delta+\sum_{d_i>1}X^{v_q(d_i)}\right)
 =X^{\alpha+\beta}\bigl(P_U(X)+P_V(X)-1-\delta\bigr).
\]
Valuation zero means a trivial primary factor, explaining the deletion of the
constant term.  The positive-degree coefficients of $P_U$ are exactly the
elementary-divisor multiplicities of $\mathcal D_U(T)_{(q)}$, and its constant
coefficient is $r_T$ minus their sum.  The same holds for $P_V$.  Thus the
five displayed invariants determine every primary group, and hence every
finite abelian group in question.
\end{proof}

\begin{theorem}[Two-scale inverse theorem and a single asymmetric probe]
\label{ext-keio-inverse}
Let $T=(U_T,V_T;E_T)$ and $S=(U_S,V_S;E_S)$ be two trees with ordered
bipartitions.  The following are equivalent:
\begin{enumerate}
\item $K(T[a,b])\cong K(S[a,b])$ for every $a,b\geq2$.
\item
\[
 r_T=r_S,\quad s_T=s_S,\quad \nu_T=\nu_S,\quad
 \mathcal D_U(T)\cong\mathcal D_U(S),\quad
 \mathcal D_V(T)\cong\mathcal D_V(S).
\]
\item For some prime $p$ and integer $k\geq2$ satisfying
\[
 p>\max(\Delta(T),\Delta(S)),\qquad
 R_k:=1+p+\cdots+p^{k-1}>\max(r_T,r_S),                      \label{ext-keio-probe-condition}
\]
one has $K(T[p,p^k])\cong K(S[p,p^k])$.
\end{enumerate}
Moreover, for every pair $(p,k)$ satisfying these inequalities, the single
group in item 3 determines all five invariants in item 2.  In particular one
may take $k=2$ and any prime larger than both tree orders.
\end{theorem}

\begin{proof}
Lemma~\ref{ext-keio-profile} proves $2\Rightarrow1$.  Primes larger than any
fixed integer exist, and choosing one larger than both orders permits $k=2$;
therefore $1\Rightarrow3$.

It remains to recover the invariants from a single probe.  Abbreviate
$r=r_T$, $s=s_T$, $\nu=\nu_T$.  Since $p$ divides no degree, the $p$-primary
part of the forward decomposition is exactly
\begin{equation}
 K(T[p,p^k])_{(p)}\cong
 (\Z/p\Z)^{s p^k-2\nu}\oplus
 (\Z/p^k\Z)^{r p-2\nu}\oplus
 (\Z/p^{k+1}\Z)^{2\nu-1}.                                 \label{ext-keio-probe-primary}
\end{equation}
The three exponents $1,k,k+1$ are distinct.  If their multiplicities are
$A_1,A_k,A_{k+1}$, respectively, then
\begin{equation}
 \nu=\frac{A_{k+1}+1}{2},\qquad
 r=\frac{A_k+2\nu}{p},\qquad
 s=\frac{A_1+2\nu}{p^k}.                                   \label{ext-keio-recover-counts}
\end{equation}
These are exact identities, including when one of the first two
multiplicities is zero.

Fix now a prime $q\ne p$ and $j\geq1$.  Set
\[
 x=m_{q,j}(\mathcal D_U(T)),\qquad
 y=m_{q,j}(\mathcal D_V(T)),\qquad
 M=m_{q,j}(K(T[p,p^k])).
\]
At $q$, both scales are units.  Formula~\eqref{ext-keio-profile-eq} reduces to
\[
 M=(p-1)x+(p^k-1)y=(p-1)(x+R_k y).
\]
Thus $z=M/(p-1)$ is an integer.  Since $0\leq x\leq r<R_k$, ordinary Euclidean
division recovers the two unknown multiplicities without ambiguity:
\begin{equation}
 x=z-R_k\lfloor z/R_k\rfloor,\qquad
 y=\lfloor z/R_k\rfloor.                                  \label{ext-keio-radix}
\end{equation}
Doing this for every nonzero elementary-divisor multiplicity recovers the
two degree groups.  Neither has a $p$-primary part because $p$ exceeds every
degree.  Consequently the single probe determines item 2 and proves
$3\Rightarrow2$.
\end{proof}

The bound on $R_k$ is a sufficient condition for this explicit decoding
procedure, not a claim of an optimal probe size.  The reconstruction concerns
the ordered color classes: interchanging them also interchanges the roles of
$a$ and $b$.

\subsection{Degree-distinct collisions surviving every two-scale blowup}

\begin{lemma}[Balanced doubling]
\label{ext-keio-doubling}
Let $X$ be a tree on $M$ vertices with a perfect matching, and choose a leaf
$x$ of $X$.  Take two disjoint copies of $X$, give them opposite bipartition
orientations, and join the two copies of $x$ by one edge.  Denote the resulting
ordered bipartite tree by $Y(X)$.  Then
\[
 r_{Y(X)}=s_{Y(X)}=\nu_{Y(X)}=M,
 \qquad
 \mathcal D_U(Y(X))\cong\mathcal D_V(Y(X))
       \cong\mathcal D(X)\oplus\Z/2\Z,
\]
where $\mathcal D(X)=\bigoplus_{v\in V(X)}\Z/d_v\Z$.
Furthermore the transformation from the degree multiset of $X$ to that of
$Y(X)$ is injective.
\end{lemma}

\begin{proof}
Joining two trees by one edge gives a tree.  The opposite orientations put the
two selected leaves in different color classes.  Each resulting color class
contains one copy of every vertex degree of $X$.  Within each color exactly
one chosen leaf has its degree changed from one to two.  This proves the
degree-group statements and the two color-class sizes.  The union of the two
original perfect matchings remains a perfect matching, so its size is $M$.

If $c_j(X)$ counts vertices of degree $j$, then
\[
 c_j(Y(X))=2c_j(X)-2\mathbf1_{j=1}+2\mathbf1_{j=2}.
\]
This formula uniquely recovers every $c_j(X)$ from the output multiset.
\end{proof}

\begin{theorem}[A superpolynomial two-scale collision family]
\label{ext-keio-collisions}
For every integer $h\geq1$, there are $2^h$ ordered bipartite trees on
\[
 N_h=30h^2+38h
\]
vertices, with pairwise different degree multisets, such that all of their
critical groups $K(T[a,b])$ are mutually isomorphic for every $a,b\geq2$.
Each tree has a perfect matching and both color classes have size
$15h^2+19h$.  In particular the collision multiplicity is
$\exp(\Omega(\sqrt{N_h}))$ along this sequence of orders.
\end{theorem}

\begin{proof}
We recall the arithmetic construction behind the uniform collision family,
so that the strengthened assertion can be checked directly.  For $t\geq3$
divisible by three, let $P_t$ be a tree with two adjacent nonleaf vertices of
degrees $2t$ and $3t+2$, and let $Q_t$ have three nonleaf vertices in a path
with degrees $2,t,4t+1$ in that order.  Attach respectively $2t-1,3t+1$
leaves in the first case and $1,t-2,4t$ leaves in the second.  Both trees have
$5t+2$ vertices.

Attach a new leaf to every vertex of either tree.  After cancellation of
common degree entries, the resulting degree multisets differ by the trade
\[
 (2,2t+1,3t+3)\quad\longleftrightarrow\quad(3,t+1,4t+2).
\]
The two triples have the same degree group, because
\[
 \Z/2\Z\oplus\Z/(2t+1)\Z\cong\Z/(4t+2)\Z,
 \qquad
 \Z/3\Z\oplus\Z/(t+1)\Z\cong\Z/(3t+3)\Z.
\]
The coprimality required in the second identity follows from $3\mid t$.

For $j=1,\ldots,h$, independently select $P_{3j}$ or $Q_{3j}$.  Join these
blocks in a path by joining chosen leaves of consecutive blocks, using
distinct leaves for the two ports of an interior block.  Each block has
enough leaves that the ports can be taken from a common pool in either
choice.  Attach one new leaf to every vertex of the joined tree, and call the
result $X$.  The new leaf edges form a perfect matching, and
\[
 |V(X)|=2\sum_{j=1}^h(15j+2)=15h^2+19h=:M_h.
\]
The joining operation changes $2(h-1)$ common degree entries from two to
three after the final leaf attachment.  Hence all $2^h$ choices have the same
degree group by the displayed arithmetic trade.

Their degree multisets are different.  Choosing $Q_{3j}$ creates a degree
$12j+2$, while choosing $P_{3j}$ does not.  This marker cannot be supplied by
another kind of degree: the other exceptional degrees are $6i+1$, $9i+3$,
and $3i+1$, whose residues modulo three are zero or one.  The common entries
one, two, and three are smaller than every marker.  Two markers agree only
for the same index.
Thus the multiset records every binary block choice.

Apply Lemma~\ref{ext-keio-doubling} to each of these trees $X$.  The resulting
trees $Y(X)$ have $2M_h=N_h$ vertices, the same ordered color-class sizes,
the same matching number, and the same two color degree groups.  Their degree
multisets remain distinct by injectivity of the doubling operation.
Theorem~\ref{ext-keio-inverse} now proves equality of the critical groups for
every pair of scales.  Finally $N_h$ is quadratic in $h$, giving the stated
growth rate.
\end{proof}

\subsection{Arbitrary nonuniform sizes recover the labeled tree}

The preceding loss of information persists even when two scales vary
independently.  Allowing a separate cluster size at every labeled vertex has
a different outcome: the orders of the groups alone suffice to recover the
base tree.  This is an elementary consequence of the spanning-tree formula,
and is presented as such rather than as a new general reconstruction method.

\begin{lemma}[Order formula for a nonuniform tree blowup]
\label{ext-keio-nonuniform-order}
Let $T$ be a tree on the labeled vertex set $\{1,\ldots,m\}$, let
$\mathbf n=(n_1,\ldots,n_m)$ with all $n_i\geq2$, and write
$T[\mathbf n]$ for its independent-cluster blowup.  Put
$D_i=\sum_{j\sim i}n_j$.  Then
\begin{equation}
 |K(T[\mathbf n])|
 =\prod_{i=1}^m n_i^{d_i-1}D_i^{n_i-1}.                    \label{ext-keio-nonuniform-order-eq}
\end{equation}
\end{lemma}

\begin{proof}
Let $N=\sum_i n_i$.  On the sum-zero subspace within cluster $i$, the
blowup Laplacian acts as $D_iI$, contributing the eigenvalue $D_i$ with
multiplicity $n_i-1$.  In the orthonormal basis of normalized cluster
indicators, its remaining restriction is the symmetric $m$-order matrix
$Q$ with diagonal $D_i$ and off-diagonal entries $-\sqrt{n_i n_j}$ along
base edges.  Its kernel is spanned by
$w=(\sqrt{n_i})_{i=1}^m$, with $w^{\mathsf T}w=N$.

Let $\rho$ be the product of the nonzero eigenvalues of $Q$.  The adjugate
identity for this symmetric rank-$m-1$ matrix gives
$\det Q^{(i)}=\rho n_i/N$.  On the other hand, congruence by
$\diag(\sqrt{n_i})$ turns $Q$ into the weighted base Laplacian with edge
weights $n_i n_j$.  Thus the matrix-tree theorem gives
\[
 \det Q^{(i)}=\frac{\tau_w}{\prod_{j\ne i}n_j},\qquad
 \rho=\frac{N\tau_w}{\prod_j n_j}.
\]
The base is a tree, so
$\tau_w=\prod_{ij\in E(T)}n_i n_j=\prod_i n_i^{d_i}$.
Applying the matrix-tree theorem to the full blowup spectrum yields
\[
 |K(T[\mathbf n])|=\frac{1}{N}\rho\prod_iD_i^{n_i-1},
\]
which is \eqref{ext-keio-nonuniform-order-eq}.
\end{proof}

\begin{theorem}[Recovery from nonuniform group orders]
\label{ext-keio-labeled-recovery}
Let $T,S$ be trees on the same labeled vertex set $\{1,\ldots,m\}$.
If
\[
 |K(T[\mathbf n])|=|K(S[\mathbf n])|
 \quad\hbox{for every }\mathbf n\in\{2,3,\ldots\}^m,
\]
then $T=S$ as labeled graphs.  In fact only the following $m$ one-parameter
families are needed.  Fix $b_j=2^j$, replace $n_i$ by $t$, keep
$n_j=b_j$ for $j\ne i$, and let the resulting group order be $\tau_i(t)$.
Then
\begin{equation}
 \lim_{t\to\infty}\tau_i(t)^{1/t}
       =\sum_{j\sim i}2^j.                                 \label{ext-keio-neighborhood-limit}
\end{equation}
The binary digits of this integer recover the complete neighbor set of $i$.
\end{theorem}

\begin{proof}
Fix $i$ and put $C_i=\sum_{j\sim i}b_j$.  In the specialization just
described, $D_i=C_i$ is constant.  For $v\ne i$, the quantity $D_v(t)$ is
either constant or a polynomial of degree one with leading coefficient one,
according as $v$ is not or is adjacent to $i$.  The order formula becomes
\[
 \tau_i(t)=C_i^{t-1}P_i(t),\qquad
 P_i(t)=t^{d_i-1}
       \prod_{v\ne i}b_v^{d_v-1}D_v(t)^{b_v-1}.
\]
This is a nonzero polynomial positive for $t\geq2$.  Its degree is exactly
\[
 d_i-1+\sum_{v\sim i}(b_v-1)=C_i-1.
\]
Therefore $P_i(t)^{1/t}\to1$, proving
\eqref{ext-keio-neighborhood-limit}.  Distinct subsets of
$\{2^1,\ldots,2^m\}$ have distinct sums, so the limit determines precisely
which vertices are adjacent to $i$.  Equality of the indicated group orders
for $T,S$ gives the same neighbor set at every vertex.
\end{proof}

The parameter correspondence in the last theorem is essential.  It does not
assert recovery from an unordered set of group values, from one arbitrary
nonuniform blowup, or from a bounded number of noisy measurements.  Together
with Theorem~\ref{ext-keio-collisions}, it separates two exact questions:
two globally chosen color scales can leave superpolynomial degree-distinct
families indistinguishable, whereas independently labeled cluster sizes
determine the tree through group orders alone.

\paragraph{Contribution boundary.}
The new deductions in this section are the two-color inverse equivalence,
the single asymmetric prime-power decoding rule, and the extension of the
degree-distinct collision construction to all two-scale blowups.  They use
the already proved forward two-scale classification and the elementary
arithmetic trade underlying the uniform construction.  The nonuniform
reconstruction theorem is a further consequence of the standard
matrix-tree method.  No named pre-existing open problem or historical
priority claim is made.


\clearpage
\markboth{Further deductions}{Taisho extension}
\section{A bounded-degree obstruction for candidate-constrained greedy resistance reduction}
\label{ext-taisho-section}

The three-candidate counterexample developed in the tournament uses large
collections of pendant leaves to amplify one terminal resistance.  The present
extension removes their unbounded degrees.  Perfect binary pendant trees preserve
the gain amplification exactly, while their additional distances cost only a
logarithmic factor in the comparison of final Kirchhoff indices.  The resulting
counterexamples have maximum degree at most five and treewidth at most two.

For a connected unit-conductance graph $G$, write
\[
 K(G)=\sum_{\{x,y\}\subseteq V(G)}R_G(x,y),\qquad
 F(S)=K(G)-K(G+S).
\]
The candidate set $\mathcal A$ is prescribed in advance.  With a budget of two
unit-cost additions, exact greedy chooses an element with maximum current
marginal gain at each step.  An optimal pair is denoted by $S_*$ and a greedy
pair by $S_g$.

\begin{theorem}[Bounded-degree, width-two greedy obstruction]
\label{ext-taisho-main}
For every integer $M\ge10$, put
\[
 q=3M^2+M,\qquad
 h=\left\lceil\log_2(100qM^2+1)\right\rceil,\qquad N=2^h-1.
\]
There is a simple connected unweighted graph $\widehat G_M$ on $q+2N$ vertices,
with exactly three specified candidate nonedges, such that every exact greedy
output with a budget of two additions satisfies
\[
 \frac{F(S_g)}{F(S_*)}\le\frac9M,
 \qquad
 \frac{K(\widehat G_M+S_g)}{K(\widehat G_M+S_*)}
 \ge\frac{M}{10h}.
\]
Every graph obtained by adding any subset of the three candidates has maximum
degree at most five and treewidth at most two.  The original graph is unicyclic,
all original and added edges have unit conductance, and all additions have equal
cost.  The first greedy choice is unique.

In particular, within this bounded-degree, bounded-treewidth class, greedy has
neither a positive constant approximation guarantee for gain nor a finite
constant multiplicative guarantee for the final Kirchhoff index.  Along this
family, writing $n=|V(\widehat G_M)|$, the two bounds have orders
$O(n^{-1/4})$ and $\Omega(n^{1/4}/\log n)$, respectively.
\end{theorem}

The theorem concerns exactly three allowed nonedges.  It does not assert the
same failure when every absent edge is available.  This restriction is retained
throughout the construction and proof.

\subsection{An attachment identity independent of the pendant geometry}

Let $H$ be a connected core with $q$ vertices and distinct terminals $s,t$.
Attach a finite unweighted tree at $s$ and another at $t$, each meeting the core
only at its designated terminal.  Let $U_s,U_t$ be the new vertex sets, with
cardinalities $N_s,N_t$.  Define
\[
 d_s=\sum_{z\in U_s}d(z,s),\qquad
 W_s=\sum_{\{z,z'\}\subseteq U_s}d(z,z'),
\]
and define $d_t,W_t$ analogously.  Distances here are in the attached trees.
Allow additions $S$ only within the core and write
\[
 A_s(S)=\sum_{x\in V(H)}R_{H+S}(s,x),\qquad
 A_t(S)=\sum_{x\in V(H)}R_{H+S}(t,x),\qquad B(S)=K(H+S).
\]

\begin{lemma}[Exact pendant-tree amplification]
\label{ext-taisho-amplification}
The Kirchhoff index of the augmented graph is
\begin{align*}
 K(S)={}&W_s+W_t+(q+N_t)d_s+(q+N_s)d_t\\
       &+N_sN_tR_{H+S}(s,t)+N_sA_s(S)+N_tA_t(S)+B(S).
\end{align*}
Consequently, when $N_s=N_t=N$, the gain has the form
\[
 F(S)=N^2g(S)+E(S),\qquad
 g(S)=R_H(s,t)-R_{H+S}(s,t),
\]
where
\[
 E(S)=N\bigl[A_s(\varnothing)+A_t(\varnothing)-A_s(S)-A_t(S)\bigr]
       +B(\varnothing)-B(S)\ge0.
\]
If every baseline core resistance is at most $D$ and $N\ge q$, then
\[
 E(S)\le\tfrac52qND,
 \qquad N[A_s(S)+A_t(S)]+B(S)\le\tfrac52qND.
\]
\end{lemma}
\begin{proof}
An attached tree meets the rest of the graph at one vertex.  A unit flow between
core vertices therefore sends no current into that tree.  Between a new vertex
$z\in U_s$ and a core vertex $x$, the effective resistance is
$d(z,s)+R_{H+S}(s,x)$.  Between $z\in U_s$ and $z'\in U_t$, it is
$d(z,s)+R_{H+S}(s,t)+d(z',t)$.  Two vertices in the same attached tree have
effective resistance equal to their tree distance: the component attached at its
root cannot carry any net current.  Summing these four types of unordered pairs
gives the displayed identity.

All terms involving only the attached trees cancel in a gain.  Adding an edge
cannot increase effective resistance, since a feasible unit flow in the old
network remains feasible with zero current on the new edge.  This proves
$E(S)\ge0$.  Also $A_s(S)+A_t(S)\le2qD$ and
$B(S)\le q(q-1)D/2$.  Since $q\le N$, both stated upper bounds follow.
\end{proof}

This identity shows that replacing pendant stars by other rooted trees of the
same orders leaves every gain comparison unchanged, when the core is fixed.
The final objective retains the attachment-dependent additive constant, which
must be controlled separately.

\subsection{The core and its strict greedy trap}

Construct $H_M$ from internally vertex-disjoint paths of lengths $M^2,M^2,M$
around the cycle $s,u,t,s$, and attach a further path of length $M^2$ from $s$
to a new endpoint $v$.  Thus $|V(H_M)|=q$.  The allowed nonedges are exactly
\[
 e=\{s,u\},\qquad f=\{u,t\},\qquad a=\{v,t\}.
\]
Each is added with unit conductance.  For $S\subseteq\{e,f,a\}$ let $R_S$ be
the core resistance between $s,t$, and let $g(S)=R_\varnothing-R_S$.

\begin{lemma}[Core estimates]
\label{ext-taisho-core}
For $M\ge10$,
\[
 R_\varnothing> M-\tfrac12,\quad
 g(\{e\})=g(\{f\})\le\tfrac12,\quad
 g(\{a\})>\tfrac45,
\]
\[
 R_{\{e,f\}}<2,\qquad
 g(\{e,f\})\ge M/4,\qquad
 g(\{a,e\}),\ g(\{a,f\})\le2.
\]
Every baseline core resistance, between any two core vertices, is at most
$3M^2$.
\end{lemma}
\begin{proof}
For the terminal calculation, suppress the degree-two internal path vertices.
The triangle has conductances $M^{-2},M^{-2},M^{-1}$.  Adding $a$ completes a
parallel $s$--$t$ branch of resistance $M^2+1$.  After multiplying the reduced
conductances by $M^2$, the triangle has baseline weights $1,1,M$, each good
addition increases a side weight by $M^2$, and the trap increases the direct
weight by $\delta=M^2/(M^2+1)$.  Original resistances are $M^2$ times the
resistances of this scaled triangle.  Series and parallel laws give
\[
 R_\varnothing=\frac{2M^2}{2M+1}>M-\tfrac12,
 \qquad
 g(\{e\})=\frac{M^4}{(2M+1)(M^3+M^2+2M+1)}\le\tfrac12,
\]
and the same single-edge formula for $f$.  For the trap,
\[
 g(\{a\})=\frac{4M^2\delta}{(2M+1)(2M+1+2\delta)}
 \ge\frac{40000}{48783}>\frac45.
\]
Here $\delta\ge100/101$, $2M+1\le(21/10)M$, and
$2M+1+2\delta\le(23/10)M$.  The good pair gives
\[
 R_{\{e,f\}}=\frac{2M^2}{(M+1)^2}<2,
 \qquad
 g(\{e,f\})=\frac{2M^4}{(2M+1)(M+1)^2}
 \ge\frac{8M}{27}>\frac M4.
\]
Increasing $\delta$ to one can only increase the gain of a trap-containing pair.
At $\delta=1$ its gain is
\[
 \frac{M^2(3M^2+4)}{(2M+1)(M^3+2M^2+2M+3)}
 \le\frac32+\frac2{M^2}\le2.
\]
These rational inequalities follow by multiplying the positive denominators.

Finally, delete the edge adjacent to $t$ on the short length-$M$ path.
The remaining tree has diameter $3M^2$: its longest arms from $s$ have lengths
$2M^2,M^2,M-1$.  Unit current along a tree path bounds every resistance by its
length.  Restoring the deleted edge can only decrease resistances, proving the
last assertion.
\end{proof}

\subsection{Binary attachments and proof of the theorem}

At each of $s,t$, attach one full binary tree of $h$ levels of new vertices by
one edge from the terminal to the tree root.  There are $N=2^h-1$ new vertices
at each terminal, and their depths from that terminal are at most $h$.
This defines $\widehat G_M$.  By the choice of $h$,
\[
 100qM^2\le N<200qM^2+1.
\]
Lemma~\ref{ext-taisho-amplification}, with $D=3M^2$, yields
\begin{equation}
\label{ext-taisho-remainder}
 0\le E(S)\le8qNM^2\le\frac{2}{25}N^2,
 \qquad N[A_s(S)+A_t(S)]+B(S)\le\frac{2}{25}N^2.
\end{equation}

\begin{proof}[Proof of Theorem~\ref{ext-taisho-main}]
For $b=e$ or $f$, the first-step gain difference satisfies
\[
 F(\{a\})-F(\{b\})
 >N^2\left(\frac45-\frac12-\frac2{25}\right)
 =\frac{11}{50}N^2>0.
\]
Thus greedy chooses $a$ uniquely first.  Whatever its second choice,
\[
 F(S_g)\le\frac{52}{25}N^2,
 \qquad F(S_*)\ge F(\{e,f\})\ge\frac M4N^2.
\]
It follows that $F(S_g)/F(S_*)\le208/(25M)<9/M$.

For the final index, Lemma~\ref{ext-taisho-core} gives
$R_{S_g}>M-5/2\ge3M/4$.  Every term in the attachment identity is nonnegative,
so
\[
 K(\widehat G_M+S_g)\ge\tfrac34 MN^2.
\]
Let the two identical binary trees have depth sum $d$ and internal new-vertex
distance sum $W$.  Since their depths are at most $h$,
\[
 d\le Nh,\qquad W\le hN(N-1).
\]
Hence the attachment-dependent constant obeys
\[
 2W+2(q+N)d\le4hN^2+2qhN.
\]
Using $R_{\{e,f\}}<2$, \eqref{ext-taisho-remainder}, and
$q/N\le1/(100M^2)$, we obtain
\begin{align*}
 K(\widehat G_M+\{e,f\})
 &<\left(4h+\frac{h}{50M^2}+\frac{52}{25}\right)N^2\\
 &<7hN^2.
\end{align*}
The last inequality holds for $M\ge10$ and $h\ge1$.
Since the optimal pair has no larger final index than $\{e,f\}$,
\[
 \frac{K(\widehat G_M+S_g)}{K(\widehat G_M+S_*)}
 >\frac{3M}{28h}\ge\frac{M}{10h}.
\]

It remains to verify the structural restrictions.  New internal tree vertices
have degree at most three.  The core degrees, after every possible addition,
are at most five at $s,t$, four at $u$, and two at $v$; all internal path vertices
have degree two.  Thus the maximum degree is at most five, even if all three
candidates are added.  Attaching trees preserves the single cycle of the
baseline core.

For completeness, a tree decomposition has vertex bags of size at most three,
with each edge covered by a bag and the bags containing any vertex forming a
connected subtree.  Before subdividing its four long paths, the core is covered
by the two adjacent bags $\{s,u,t\}$ and $\{s,t,v\}$, for every candidate subset.
For a path $a=x_0,x_1,\ldots,x_\ell=b$ replacing an edge of this small core,
attach to a bag containing $a,b$ the bag chain
\[
 \{a,b,x_{\ell-1}\},\
 \{a,x_{\ell-1},x_{\ell-2}\},\ \ldots,\
 \{a,x_2,x_1\}.
\]
All four paths have length at least two.  These bags cover their edges and
preserve the connected-bag condition.  Any retained direct candidate edge is
already covered by a core bag.  Each pendant-tree edge is covered by a new
two-vertex bag attached recursively at a bag containing its parent.  This is a
tree decomposition of width at most two.

Finally $q=\Theta(M^2)$, $N=\Theta(M^4)$, and $h=\Theta(\log M)$.  Therefore
$n=\Theta(M^4)$, and the asymptotic claims follow.  In particular
$M/h\longrightarrow\infty$.
\end{proof}

\subsection{What this extension closes}

The construction proves that the unbounded degrees of the pendant stars are
unnecessary for the candidate-constrained obstruction: maximum degree five and
treewidth two still permit vanishing greedy gain ratios and unbounded final
Kirchhoff ratios.  The exact attachment identity separates this structural
improvement from the electrical core, whose estimates were already established
in the tournament's three-candidate construction.

Achterberg and Kooij distinguish the disappearance of submodularity-based
certificates from actual poor greedy performance \cite{ref052}.
Their default candidate set consists of all nonedges.  Our prescribed
three-candidate theorem, including its bounded-degree strengthening here,
does not resolve their unchanged unrestricted formulation.  In particular,
allowing shortcuts involving the new pendant vertices would change the
algorithm's available actions and invalidate the present first-choice proof.
The result is an additional proved theorem in the constrained model; no claim
of comprehensive historical priority is made.


\clearpage
\markboth{Further deductions}{Reiwa extension}
\section{Exact universal entropy contraction for arbitrary recombination laws}
\label{ext-reiwa-section}

This extension starts from the weighted total-correlation inequality in
Reiwa R05 and determines both a sharp entropy-production constant and a
sharp finite-time contraction profile.  The normal-minorant lemma used
below is established prior work \cite[Lemma~3.10]{ref067}.
Caputo and Sinclair \cite{ext-reiwa-CS} define the universal
entropy-production constant studied here and compute the associated
entropy-subadditivity constants for four recombination laws.  The
fragmentation representation is a version of the classical branching
representation of nonlinear recombination; see also
\cite{ext-reiwa-CP}.  We give complete proofs of the extensions, without
making a claim of historical priority.  In particular, the universal
optimization over product reference measures below must not be confused
with the optimization for one fixed reference measure.

\subsection{A sharp inequality for random partitions}

Fix $V=[n]$, $n\geq2$, and finite coordinate spaces.  For a probability
law $P$ on their product, write
\[
 \pi_P=\bigotimes_{i\in V}P_i,
 \qquad T(P)=D(P\Vert\pi_P)=\sum_{i\in V}H(P_i)-H(P).
\]
Entropies and relative entropies use natural logarithms.  For a set
partition $\mathcal B$ of $V$, define
\[
 R_{\mathcal B}(P)=\bigotimes_{B\in\mathcal B}P_B,
 \qquad
 k_{\mathcal B}(S)=\#\{B\in\mathcal B:B\cap S\ne\varnothing\}.
\]
If $M$ is a probability law on set partitions, put
\[
 R_M(P)=\sum_{\mathcal B}M(\mathcal B)R_{\mathcal B}(P),
 \qquad b_M(S)=\mathbb E_M k_{\mathcal B}(S),
 \qquad
 \rho_M=\max_{\substack{S\subseteq V\\|S|\geq2}}
          \frac{|S|-b_M(S)}{|S|-1}.
\]
Thus $0\leq\rho_M\leq1$.

\begin{lemma}[Known normal-minorant reduction]\label{ext-reiwa-minorant}
Every normalized, monotone, submodular function $h$ on $2^V$ has a
minorant of the form
\[
 h'(A)=\sum_{\varnothing\ne S\subseteq V}c_S
                 \mathbf1\{A\cap S\ne\varnothing\},\qquad c_S\geq0,
\]
such that $h'(V)=h(V)$ and $h'(\{i\})=h(\{i\})$ for every $i$.
Consequently, with $T_h(A)=\sum_{i\in A}h(\{i\})-h(A)$,
\begin{equation}\label{ext-reiwa-partition-entropy}
 \mathbb E_M\sum_{B\in\mathcal B}T_h(B)\leq\rho_M T_h(V).
\end{equation}
\end{lemma}
\begin{proof}
For completeness, the minorant can be constructed by induction.  Choose
$v\in V$ and write $U=V\setminus\{v\}$.  The contraction
$k_0(A)=h(A\cup\{v\})-h(\{v\})$ is a polymatroid on $U$.
Let $k\leq k_0$ be an inductively constructed normal minorant preserving
the full-set and singleton values.  Set
\[
 a_v=h(\{v\}),\qquad
 a_i=h(\{i\})+h(\{v\})-h(\{i,v\})\quad(i\in U),
 \qquad m(A)=\max_{i\in A}a_i,
\]
where $m(\varnothing)=0$.  Submodularity and monotonicity give
$0\leq a_i\leq a_v$.  The function $m$ is a nonnegative combination of
coverage functions, by decomposing its finitely many positive levels.
Define $h'(A)=k(A\cap U)+m(A)$.  If $v\in A$, the minorant inequality
follows from contraction.  If $v\notin A$, submodularity gives
\[
 m(A)\leq h(A)+h(\{v\})-h(A\cup\{v\}),
\]
so again $h'(A)\leq h(A)$.  The singleton and full-set equalities follow
directly from the construction.  The one-element base case is immediate.

For this minorant,
\[
 T_{h'}(A)=\sum_S c_S(|A\cap S|-1)_+,
 \qquad T_{h'}(V)=T_h(V),\qquad T_h(A)\leq T_{h'}(A).
\]
Since every $\mathcal B$ partitions $V$,
\[
 \sum_{B\in\mathcal B}(|B\cap S|-1)_+=|S|-k_{\mathcal B}(S).
\]
Taking expectations and using the definition of $\rho_M$ proves
\eqref{ext-reiwa-partition-entropy}.  This proof also makes explicit
where the known normal-minorant lemma enters the argument.
\end{proof}

\begin{theorem}[Exact contraction under a random partition]
\label{ext-reiwa-random-partition}
For every finite law $P$,
\begin{equation}\label{ext-reiwa-partition-contraction}
 T(R_M(P))\leq\rho_M T(P).
\end{equation}
The constant is optimal over all finite coordinate spaces and laws with
$T(P)>0$.  The same optimal constant is obtained by restricting to
strictly positive laws on $\{0,1\}^n$, while allowing their one-coordinate
marginals to vary.
\end{theorem}
\begin{proof}
Every recombined law has the same one-coordinate marginals as $P$.
Convexity of relative entropy and
Lemma~\ref{ext-reiwa-minorant}, applied to $h(A)=H(P_A)$, give
\[
 T(R_M(P))\leq\mathbb E_M D(R_{\mathcal B}(P)\Vert\pi_P)
 =\mathbb E_M\sum_{B\in\mathcal B}T(P_B)\leq\rho_M T(P).
\]

For sharpness fix $S$ of size $s\geq2$.  Let every coordinate in $S$
equal one common Bernoulli$(\varepsilon)$ variable, and let the other
coordinates be independent fair bits, independent of that variable.
Call this law $P^\varepsilon$.  Then
$T(P^\varepsilon)=(s-1)h_2(\varepsilon)$.
Conditional on $\mathcal B$, the restriction of $R_{\mathcal B}(P^\varepsilon)$
to $S$ consists of $k_{\mathcal B}(S)$ independent rare bits copied on
the occupied blocks.  If $Q^\varepsilon=R_M(P^\varepsilon)$, then
\[
 Q^\varepsilon_S(0^s)
 =\mathbb E_M(1-\varepsilon)^{k_{\mathcal B}(S)}
 =1-b_M(S)\varepsilon+O(\varepsilon^2).
\]
For each nonzero configuration $x$, its probability is
$c_x\varepsilon+O(\varepsilon^2)$ with $c_x\geq0$, and
$\sum_{x\ne0}c_x=b_M(S)$.  There are finitely many configurations;
therefore
\[
 H(Q^\varepsilon_S)
 =b_M(S)\varepsilon\log(1/\varepsilon)+O(\varepsilon).
\]
All one-coordinate marginals in $S$ are still Bernoulli$(\varepsilon)$.
Consequently,
\[
 \lim_{\varepsilon\downarrow0}
 \frac{T(R_M(P^\varepsilon))}{T(P^\varepsilon)}
 =\frac{s-b_M(S)}{s-1}.
\]
Maximizing over $S$ proves sharpness.  For each fixed $\varepsilon>0$,
replace $P^\varepsilon$ by
$(1-\eta)P^\varepsilon+\eta\pi_{P^\varepsilon}$ and then let
$\eta\downarrow0$.  These laws are strictly positive, keep the same
one-coordinate marginals, and their ratios converge by continuity.
This establishes the final assertion without interchanging limits.
\end{proof}

\subsection{The universal entropy-production constant}

Let $\nu$ be any probability distribution on subsets $A\subseteq V$;
empty and full subsets are allowed.  Define
\[
 Q_A(P)=P_A\otimes P_{A^c},\qquad
 \Psi_\nu(P)=\sum_A\nu(A)Q_A(P),\qquad
 q_\nu(S)=\sum_A\nu(A)\mathbf1\{0<|A\cap S|<|S|\},
\]
and set
\begin{equation}\label{ext-reiwa-kappa}
 \kappa_\nu=
 \min_{\substack{S\subseteq V\\|S|\geq2}}
       \frac{q_\nu(S)}{|S|-1}.
\end{equation}
For a strictly positive law $P$, its entropy production is
\begin{equation}\label{ext-reiwa-production}
 \mathcal I_\nu(P)=
 \sum_A\nu(A)\bigl[D(P\Vert Q_A(P))+D(Q_A(P)\Vert P)\bigr].
\end{equation}
The universal constant is
\[
 \delta_\nu=\inf\frac{\mathcal I_\nu(P)}{T(P)},
\]
where the infimum runs over all finite coordinate spaces and strictly
positive laws with $T(P)>0$.  Restricting to positive laws does not
change the usual extended-value definition: finite zero coordinates
can be removed, and mixing with the product of the marginals
approximates any remaining finite-production boundary law.

\begin{theorem}[Exact universal production and one-step constants]
\label{ext-reiwa-production-exact}
For every subset law $\nu$,
\begin{equation}\label{ext-reiwa-delta-exact}
 \boxed{\displaystyle
 \delta_\nu=\kappa_\nu
 =\min_{|S|\geq2}\frac{q_\nu(S)}{|S|-1}.}
\end{equation}
Furthermore,
\[
 \sup_{P:T(P)>0}\frac{T(\Psi_\nu(P))}{T(P)}=1-\kappa_\nu.
\]
Both universal optimizations can be restricted to strictly positive
binary laws.  The constant $\kappa_\nu$ is positive exactly when $\nu$
separates every pair of distinct coordinates with positive probability.
\end{theorem}
\begin{proof}
Discard empty blocks in the partition $\{A,A^c\}$.  For its induced law
$M$, one has $b_M(S)=1+q_\nu(S)$ for nonempty $S$.
Theorem~\ref{ext-reiwa-random-partition} proves the one-step assertion.
Lemma~\ref{ext-reiwa-minorant} also gives
\[
 \sum_A\nu(A)\bigl[T(P_A)+T(P_{A^c})\bigr]
       \leq(1-\kappa_\nu)T(P).
\]
Since
\[
 D(P\Vert Q_A(P))
 =T(P)-T(P_A)-T(P_{A^c}),
\]
the nonnegativity of the reverse divergences proves
$\mathcal I_\nu(P)\geq\kappa_\nu T(P)$.

Here is a strictly positive sharpness construction; in particular, it
does not differentiate a law with zeros or exchange small-time and
small-probability limits.  Fix $S$ of size $s\geq2$ and put
$N=2^s-2$.  For $0<\varepsilon<e^{-4}$ choose
$\theta=(\log(1/\varepsilon))^{-1/2}$, and on $\{0,1\}^S$ let
\begin{equation}\label{ext-reiwa-positive-witness}
 P(0^s)=1-\varepsilon,\qquad
 P(1^s)=\varepsilon(1-\theta),\qquad
 P(x)=\frac{\varepsilon\theta}{N}\quad(x\ne0^s,1^s).
\end{equation}
Outside $S$, use independent fair bits independent of $P_S$.
Each marginal in $S$ has success probability
$\varepsilon(1-\theta/2)$.  For each nonempty $C\subseteq S$, write
$P_C(X_C\ne0)=a_C\varepsilon$.  Then
$1-\theta\leq a_C\leq1$, and the binary decomposition according to
$\{X_C=0\}$ gives, uniformly in $\theta$,
\[
 H(P_C)=a_C\varepsilon\log(1/\varepsilon)+O_s(\varepsilon).
\]
For $C=S$, $a_S=1$.  It follows that
\begin{equation}\label{ext-reiwa-witness-T}
 T(P)=(s-1+o(1))\varepsilon\log(1/\varepsilon).
\end{equation}

If $A$ does not split $S$, then $Q_A(P)=P$.  If it splits $S$ into
nonempty $C$ and $D$, the forward divergence satisfies
\[
 D(P\Vert Q_A(P))
 =H(P_C)+H(P_D)-H(P_S)
 =(1+o(1))\varepsilon\log(1/\varepsilon).
\]
To bound the reverse divergence, it suffices to work on $S$ and put
$Q=P_C\otimes P_D$.  For every nonzero $x$,
$P(x)\geq\varepsilon\theta/N$ and $Q(x)\leq\varepsilon$;
also $Q(X_S\ne0)\leq2\varepsilon$.  At zero,
$Q(0)/P(0)\leq1/(1-\varepsilon)$, so its contribution is at most
$2\varepsilon$.  Hence
\[
 D(Q\Vert P)\leq
 2\varepsilon+2\varepsilon\log(N/\theta)
 =o\bigl(\varepsilon\log(1/\varepsilon)\bigr).
\]
There are finitely many cuts.  Combining these bounds with
\eqref{ext-reiwa-witness-T} proves
\[
 \lim_{\varepsilon\downarrow0}
 \frac{\mathcal I_\nu(P)}{T(P)}
       =\frac{q_\nu(S)}{s-1}.
\]
Choose a minimizing $S$ to obtain \eqref{ext-reiwa-delta-exact}.
Finally, a pair that is never separated makes the minimum zero.
Conversely, any $S$ of size at least two contains a pair $i,j$, and
$q_\nu(S)\geq q_\nu(\{i,j\})>0$ under pair separation.
\end{proof}

To identify precisely the convention in \cite{ext-reiwa-CS}, let
$\pi$ be a strictly positive product measure, let $P=f\pi$, and
assume $P_i=\pi_i$.  With
$f_A=\mathbb E_\pi[f\mid X_A]$, the density of $Q_A(P)$ is
$f_Af_{A^c}$.  Thus \eqref{ext-reiwa-production} is exactly
\[
 \sum_A\nu(A)\,
 \pi\!\left[(f-f_Af_{A^c})
                \log\frac{f}{f_Af_{A^c}}\right].
\]
Consequently \eqref{ext-reiwa-delta-exact} concerns the infimum of their
$\delta(\pi,\nu)$ over all product $\pi$ and all finite alphabets,
with no extra factor of two.  It does not assert that every fixed
$\delta(\pi,\nu)$ equals $\kappa_\nu$.

\subsection{Exact contraction at every finite time}

Consider the continuous-time normalization
\begin{equation}\label{ext-reiwa-ode}
 \frac{d}{dt}P_t=\Psi_\nu(P_t)-P_t,\qquad P_0=P,
\end{equation}
and the discrete iterates $P^{(k)}=\Psi_\nu^k(P)$.  This polynomial
ordinary differential equation has a unique global solution on the
finite simplex.  Set $q_\nu(\{i\})=0$ and
$r_\nu(U)=1-q_\nu(U)$ for every nonempty $U$.

\begin{theorem}[Sharp finite-time profiles]\label{ext-reiwa-profiles}
For $t\geq0$ and integers $k\geq1$, define
\begin{align}
 \rho_{\nu}^{\mathrm{ct}}(t)
 &=\max_{|S|\geq2}\frac{1}{|S|-1}
   \sum_{\substack{U\subseteq S\\|U|\geq2}}
        (-1)^{|U|}e^{-q_\nu(U)t},
       \label{ext-reiwa-continuous-profile}\\
 \rho_{\nu}^{\mathrm{dt}}(k)
 &=\max_{|S|\geq2}\frac{1}{|S|-1}
   \sum_{\substack{U\subseteq S\\|U|\geq2}}
        (-1)^{|U|}r_\nu(U)^k.
       \label{ext-reiwa-discrete-profile}
\end{align}
Then
\[
 \sup_{P:T(P)>0}\frac{T(P_t)}{T(P)}
       =\rho_{\nu}^{\mathrm{ct}}(t),\qquad
 \sup_{P:T(P)>0}\frac{T(P^{(k)})}{T(P)}
       =\rho_{\nu}^{\mathrm{dt}}(k).
\]
Both suprema are unchanged if only positive binary laws are allowed,
with their marginals unrestricted.  Despite their alternating-sum
form, both profiles lie in $[0,1]$.
\end{theorem}
\begin{proof}
Start a partition process from $\{V\}$.  In discrete time every current
block $B$ independently samples $A\sim\nu$ and is replaced by the
nonempty sets $B\cap A$ and $B\cap A^c$.  In continuous time each
current block has an independent rate-one clock, and at each ring the
same update is made.  Updates producing no split are allowed.  There
are at most $n$ blocks, so the continuous process is nonexplosive.
Write its partition at generation $k$ or time $t$ as $\Pi_k$ or
$\Pi_t$.

We first justify the ancestry identity, including its time normalization:
\begin{equation}\label{ext-reiwa-ancestry}
 P^{(k)}=\mathbb E R_{\Pi_k}(P),\qquad
 P_t=\mathbb E R_{\Pi_t}(P).
\end{equation}
For any nonempty $B\subseteq V$, run this process starting from $\{B\}$
with cuts obtained by intersecting independent $\nu$-cuts with $B$.
Let $F_t^B(P_B)$ be the expected product of its terminal block
marginals, and set $F_t^{\varnothing}=1$.
Restriction of a partition process on $V$ to $B$ has this law: distinct
blocks meeting $B$ carry independent rate-one clocks, and each
restricted block is split by exactly the indicated restricted cut.
Thus $(F_t^V(P))_B=F_t^B(P_B)$.

Condition on the first rate-one clock ring of the initial block $V$.
If it occurs at time $s\leq t$ with cut $A$, its nonempty children
evolve independently for the remaining time $t-s$.  Empty or full
cuts simply restart the same block with an independent future clock.
Consequently
\[
 F_t^V(P)=e^{-t}P+
 \int_0^t e^{-s}\sum_A\nu(A)
       F_{t-s}^{A}(P_A)\otimes F_{t-s}^{A^c}(P_{A^c})\,ds
 =e^{-t}P+\int_0^t e^{-s}\Psi_\nu(F_{t-s}^V(P))\,ds.
\]
After changing variables this is the variation-of-constants formula
for \eqref{ext-reiwa-ode}; uniqueness proves its ancestry identity.
The same root-split argument with $k$ remaining generations proves
the discrete identity by induction.  Independence here is between
different ancestral blocks; one must not reuse one common cut for
all blocks at a generation.

For a nonempty $U$, once its coordinates lie in different ancestral
blocks they can never rejoin.  While they are still in one block, a
clock ring separates them with probability $q_\nu(U)$, independently
of the containing block and of the past.  It follows that
\[
 \Pr(U\text{ lies in one block of }\Pi_t)=e^{-q_\nu(U)t},
 \qquad
 \Pr(U\text{ lies in one block of }\Pi_k)=r_\nu(U)^k.
\]
For every deterministic partition $\mathcal B$, inclusion--exclusion
inside each of its occupied blocks gives the exact identity
\[
 k_{\mathcal B}(S)=
 \sum_{\varnothing\ne U\subseteq S}(-1)^{|U|+1}
       \mathbf1\{U\text{ lies in one block of }\mathcal B\}.
\]
Taking expectations evaluates $b_M(S)$ for the partition laws in
\eqref{ext-reiwa-ancestry}.  The singleton terms contribute $|S|$.
Theorem~\ref{ext-reiwa-random-partition} now yields exactly
\eqref{ext-reiwa-continuous-profile} and
\eqref{ext-reiwa-discrete-profile}, including sharpness and positivity.
\end{proof}

\begin{corollary}[Best universal exponential rate]
\label{ext-reiwa-exponential}
Under the normalization \eqref{ext-reiwa-ode},
\[
 T(P_t)\leq e^{-\kappa_\nu t}T(P)
\]
for all finite laws and all $t\geq0$.  No larger exponent works
uniformly with prefactor one over all finite laws and alphabets.
Moreover,
\[
 (\rho_{\nu}^{\mathrm{ct}})'(0+)=-\kappa_\nu.
\]
\end{corollary}
\begin{proof}
For positive $P$, the flow remains positive since $P_t\geq e^{-t}P$.
The single-coordinate marginals, and therefore $\pi_P$, are conserved.
The identity
\[
 \sum_x(P(x)-Q_A(P)(x))\log\frac{Q_A(P)(x)}{\pi_P(x)}=0
\]
holds because the logarithm is a sum of one function on $A$ and one
on $A^c$, whose expectations agree under $P$ and $Q_A(P)$.
It follows that
$\frac{d}{dt}T(P_t)=-\mathcal I_\nu(P_t)$.
Theorem~\ref{ext-reiwa-production-exact} and integration give the
bound.  Approximation gives it on the boundary of the simplex.
If a larger uniform exponent existed, differentiation at zero for
every strictly positive initial law would contradict
\eqref{ext-reiwa-delta-exact}.

For the profile derivative, initially $\Pi_0=\{V\}$, and the first
clock ring increases the number of blocks meeting $S$ by one precisely
when its cut splits $S$.  Hence
$\mathbb E k_{\Pi_t}(S)=1+q_\nu(S)t+O(t^2)$.
The maximum of the finitely many corresponding ratios is
$1-\kappa_\nu t+O(t^2)$, proving the derivative assertion.
\end{proof}

\subsection{Examples and scope}

The formula for $\kappa_\nu$ recovers the four subadditivity constants
in \cite{ext-reiwa-CS}, and
Theorem~\ref{ext-reiwa-production-exact} gives the same values for the
universal production constant.  The cut laws and values are
\[
\begin{array}{c|c}
 \text{cut law}&\delta_\nu=\kappa_\nu\\ \hline
 \nu(\{i\})=1/n&1/(n-1)\\
 \nu(\{1,\ldots,j\})=1/(n+1),\ 0\leq j\leq n&1/(n+1)\\
 \nu(A)=2^{-n}&(1-2^{1-n})/(n-1)\\
 \nu(A)=a^{|A|}(1-a)^{n-|A|},\ 0\leq a\leq1
       &[1-a^n-(1-a)^n]/(n-1)
\end{array}
\]
For single-site cuts, $q_\nu(S)=|S|/n$ for proper $S$ of size at
least two, whereas $q_\nu(V)=1$; the minimum follows immediately.
For prefix cuts, a set $S$ is split by exactly
$\max S-\min S$ of the $n+1$ possible cuts, which is at least
$|S|-1$, with equality for intervals.  For Bernoulli cuts,
$q_\nu(S)=1-a^{|S|}-(1-a)^{|S|}$.  The sequence
$[1-a^s-(1-a)^s]/(s-1)$ decreases in $s\geq2$: its numerator is zero
at $s=1$ and has nonincreasing increments
$(1-a)a^s+a(1-a)^s$.  Thus the minimum is attained at $s=n$;
$a=1/2$ gives uniform crossover.

For a concrete finite-time example, take $n=3$ and single-site cuts.
Pairs have separation probability $2/3$ and the triple has probability
$1$.  The triple dominates the pair terms in
\eqref{ext-reiwa-continuous-profile}, giving
\[
 \rho_{\nu}^{\mathrm{ct}}(t)
   =\frac{3e^{-2t/3}-e^{-t}}{2}\qquad(t\geq0).
\]
Its initial derivative is $-1/2$, whereas its large-time exponential
rate is $2/3$.  This distinguishes the best universal exponent with
prefactor one from the eventual decay rate.

All sharpness statements allow the product reference measure to vary;
the positive witnesses approach a rare-event boundary as
$\varepsilon\downarrow0$.  The results therefore do not determine a
sharp constant for a prescribed product measure bounded away from
zero.  They also do not address the sharp entropy-production constant
of the finite-particle generalized-transposition system in
\cite{ext-reiwa-CP}, or the fixed-equilibrium total-variation cutoff
problem studied in \cite{ext-reiwa-KS}.  These are different
optimizations.


\clearpage
\markboth{Further deductions}{Heisei extension}
\section{A maximal rational spectral space from local pole budgets}
\label{ext-heisei-section}

This section extends the positive-shift resonance-chain argument from
polynomials to rational functions.  A prescribed denominator can first be
handled by conjugation.  More substantially, a local pole-raising hypothesis
produces one denominator that works simultaneously for every finite rational
invariant space, every rational eigenfunction, and every rational Jordan
chain.  The hypothesis holds for every positive-order monic rational
differential operator.  The local indicial method is classical in rational
solution algorithms \cite{ref089}; the result below combines
it with a uniform spectral dimension bound and a maximal invariant-space
certificate.  No historical-priority claim is made.

Write $\partial=d/dx$ and let
\[
 T=\sum_{j=0}^q a_j(x)\partial^j,
 \qquad a_j\in\C(x),\quad q\ge1,\quad a_q\ne0.
\]
For a rational function, its degree at infinity is numerator degree minus
denominator degree.  Assume
\[
 s=\max_{a_j\ne0}\bigl(\deg_\infty a_j-j\bigr)>0.
\]
Thus $s$ is a positive integer.  For each index attaining this maximum, write
$a_j(x)=\alpha_jx^{j+s}+O(x^{j+s-1})$, and define
\begin{equation}
 h(z)=\sum_{\deg_\infty a_j-j=s}\alpha_j(z)_j,
 \qquad (z)_j=z(z-1)\cdots(z-j+1),\quad (z)_0=1.
 \label{ext-heisei-h}
\end{equation}
The polynomial $h$ is nonzero.  If a nonzero rational function $f$ has
$f(x)=c x^d+O(x^{d-1})$ at infinity, with $d\in\mathbb Z$, then
\begin{equation}
 Tf=c h(d)x^{d+s}+O(x^{d+s-1}).
 \label{ext-heisei-infinity}
\end{equation}
The formula remains valid for negative $d$.

For a finite set $H\subseteq\mathbb Z_{\ge0}$, set
\[
 \mathcal A_s(H)=\{d\ge0:d+ks\in H\text{ for some integer }k\ge0\},
 \qquad N_s(H)=|\mathcal A_s(H)|.
\]
Equivalently, if $M_r$ is the largest member of $H$ congruent to $r$ modulo
$s$, omit residues with no such member and write
\begin{equation}
 N_s(H)=\sum_r\left(\frac{M_r-r}{s}+1\right).
 \label{ext-heisei-chain-count}
\end{equation}
In particular, $N_s(\varnothing)=0$.

\subsection{A prescribed denominator}

Let $\mathfrak d\in\C[x]$ be monic and nonzero, put $K=\deg\mathfrak d$, and
write
\[
 \mathcal R_{\mathfrak d}=\mathfrak d^{-1}\C[x],
 \qquad
 H_{\mathfrak d}=\{n\in\mathbb Z_{\ge0}:h(n-K)=0\}.
\]
Membership in $\mathcal R_{\mathfrak d}$ means that the reduced denominator
divides $\mathfrak d$; no restriction on numerator degree is imposed initially.

\begin{theorem}[Fixed-denominator resonance bound]
\label{ext-heisei-fixed}
Every finite-dimensional $T$-invariant space
$W\subseteq\mathcal R_{\mathfrak d}$ satisfies
\[
 \dim W\le N_s(H_{\mathfrak d}).
\]
If $H_{\mathfrak d}=\varnothing$, no such nonzero space exists.  Otherwise,
with $M=\max H_{\mathfrak d}$, every such space is contained in
$\mathfrak d^{-1}\Pcal_M$.  There is a unique largest such invariant space,
computable by a finite block-kernel calculation.
\end{theorem}

\begin{proof}
Conjugate by multiplication by $\mathfrak d$:
\[
 \widehat T=\mathfrak d T\mathfrak d^{-1}.
\]
This is again a rational differential operator.  If $p$ has degree $n$ and
leading coefficient $c$, then $p/\mathfrak d=cx^{n-K}+O(x^{n-K-1})$, so
\[
 \widehat Tp=c h(n-K)x^{n+s}+O(x^{n+s-1}).
\]
Conjugation preserves the positive shift $s$ and replaces the leading
multiplier by $h(z-K)$.  Indeed its coefficient shifts are at most $s$ by
the product rule, and the displayed nonzero multiplier precludes a smaller
shift.

The space $\widehat W=\mathfrak d W$ is a finite invariant polynomial space.
Choose a coefficient-echelon basis of $\widehat W$ with distinct leading
degrees and let $S$ be its degree set.  This is also the set of degrees of
all nonzero elements of $\widehat W$, and $|S|=\dim W$.  If
$n\in S\setminus H_{\mathfrak d}$, invariance and the displayed leading
term force $n+s\in S$.  Iteration in a finite set must reach
$H_{\mathfrak d}$.  Hence $S\subseteq\mathcal A_s(H_{\mathfrak d})$, proving
the bound.  The maximum degree in any nonzero invariant space belongs to
$H_{\mathfrak d}$, and is therefore at most $M$.

For the finite construction, let $\operatorname{pp}(f)$ denote the
polynomial part of a rational function at infinity and let $\pi_M$ truncate
polynomial coefficients above degree $M$.  On $\Pcal_M$ set
\[
 Lp=\pi_M\operatorname{pp}(\widehat Tp),
 \qquad Cp=Q(\widehat Tp-Lp),
\]
where a fixed nonzero polynomial $Q$ clears all coefficient denominators
of $\widehat T$.  The map $C$ takes values in a fixed finite-dimensional
polynomial space and records both poles and high-degree polynomial terms.
Define
\begin{equation}
 U_* =\bigcap_{j=0}^{M}\ker(CL^j),
 \qquad W_{\mathfrak d,*}=\mathfrak d^{-1}U_*.
 \label{ext-heisei-kernel}
\end{equation}
Any invariant polynomial space in $\Pcal_M$ has zero leakage at every
iterate, so it lies in $U_*$.  Conversely Cayley--Hamilton for the
$(M+1)$-dimensional map $L$ implies that $CL^{M+1}$ vanishes on $U_*$.
It follows that $L(U_*)\subseteq U_*$.  Since $C(U_*)=0$, the operators
$\widehat T$ and $L$ agree on $U_*$, which is consequently invariant and
maximal.  Conjugation proves the assertion about $W_{\mathfrak d,*}$.
\end{proof}

\subsection{A denominator derived uniformly from the operator}

Let $\mathcal S$ be the finite set consisting of the zeros and poles of
$a_q$, together with the poles of the other nonzero coefficients.  At a
finite point $c$, write
\[
 a_j(x)=b_{j,c}(x-c)^{\nu_{j,c}}
       +O((x-c)^{\nu_{j,c}+1}),\qquad b_{j,c}\ne0,
\]
for each nonzero $a_j$, where $\nu_{j,c}=\operatorname{ord}_c a_j$.  Define
\begin{equation}
 \sigma_c=\max_{a_j\ne0}(j-\nu_{j,c}),
 \qquad
 g_c(z)=\sum_{j-\nu_{j,c}=\sigma_c}b_{j,c}(z)_j.
 \label{ext-heisei-local-data}
\end{equation}
The distinct falling factorials show that $g_c$ is nonzero.  For
$c\notin\mathcal S$, one has $\sigma_c=q$ and
$g_c(z)=a_q(c)(z)_q$, which has no negative integer root.

\begin{lemma}[A pole bound for an invariant space]
\label{ext-heisei-local-lemma}
Suppose $\sigma_c>0$ and $W\subseteq\C(x)$ is a finite-dimensional
$T$-invariant space.  If $k>0$ is the largest pole order at $c$ of an
element of $W$, then $g_c(-k)=0$.
\end{lemma}

\begin{proof}
Pole order in a finite-dimensional rational space is bounded by the maximum
of the pole orders of any basis.  Choose $f\in W$ with largest pole order
$k$, and write $f=u(x-c)^{-k}+O((x-c)^{-k+1})$, $u\ne0$.  The product rule
and the local expansions give
\[
 Tf=u g_c(-k)(x-c)^{-k-\sigma_c}
       +O((x-c)^{-k-\sigma_c+1}).
\]
If $g_c(-k)\ne0$, then $Tf\in W$ has pole order $k+\sigma_c>k$, a
contradiction.  Therefore the multiplier must vanish.
\end{proof}

\begin{theorem}[Complete rational spectral space under pole raising]
\label{ext-heisei-global}
Assume $\sigma_c>0$ for every $c\in\mathcal S$.  Set
\begin{equation}
 K_c=\max\bigl(\{k\in\mathbb Z_{>0}:g_c(-k)=0\}\cup\{0\}\bigr),
 \qquad
 \mathfrak d(x)=\prod_{c\in\mathcal S}(x-c)^{K_c}.
 \label{ext-heisei-denominator}
\end{equation}
Then every finite-dimensional rational invariant space is contained in
$\mathcal R_{\mathfrak d}$.  The space $W_{\mathfrak d,*}$ of
Theorem~\ref{ext-heisei-fixed} is therefore the unique largest
finite-dimensional rational invariant space of $T$, and
\[
 \dim W_{\mathfrak d,*}\le N_s(H_{\mathfrak d}).
\]
It contains every rational generalized eigenvector: if
$f\in\C(x)$ and $(T-\lambda)^k f=0$, then $f\in W_{\mathfrak d,*}$.
If $T|_{W_{\mathfrak d,*}}$ has matrix $K_*$ in any basis, the entire
rational eigenvalue set is
\[
 \{\lambda\in\C:Tf=\lambda f\text{ for some }0\ne f\in\C(x)\}
   =\{\lambda:\det(\lambda I-K_*)=0\}.
\]
For a zero-dimensional space, the determinant is interpreted as $1$.
\end{theorem}

\begin{proof}
At a point outside $\mathcal S$, the preceding lemma and the absence of
negative integer roots of $(z)_q$ exclude every pole of an invariant
space.  At $c\in\mathcal S$, the same lemma bounds its largest pole order
by $K_c$.  Thus multiplication by $\mathfrak d$ removes every finite pole
of every member, leaving a polynomial.  All finite rational invariant
spaces are therefore included in the fixed-denominator class of
Theorem~\ref{ext-heisei-fixed}, which proves global maximality and the
dimension bound.

If $(T-\lambda)^k f=0$, all its intermediate iterates belong to $\C(x)$
because $T$ has rational coefficients.  Their span is finite-dimensional
and $T$-invariant, hence lies in the maximal space.  Finally every rational
eigenfunction belongs to that space, and every eigenvector of its finite
matrix gives a nonzero rational eigenfunction.  These two inclusions prove
the spectral equality.
\end{proof}

\begin{corollary}[Monic rational operators]
\label{ext-heisei-monic}
If $a_q$ is a nonzero constant and $s>0$, all conclusions of
Theorem~\ref{ext-heisei-global} hold.  In particular, there are finitely
many rational eigenvalues and one finite-dimensional space contains all
rational Jordan chains of all lengths.
\end{corollary}

\begin{proof}
At every finite point the leading derivative contributes
$q-\operatorname{ord}_c a_q=q>0$, so $\sigma_c\ge q>0$.
\end{proof}

For algebraic coefficient data the construction is effective: determine the
finitely many exceptional points, compute the local multipliers and their
negative integer roots, form the denominator, and perform the finite
division and kernel operations in~\eqref{ext-heisei-kernel}.  The assertion
is one of termination, not polynomial-time complexity.  In contrast to
solving one prescribed equation $(T-\lambda)f=0$, the denominator and
dimension bound are uniform over all energies and their rational Jordan
chains.

\begin{corollary}[Arithmetic restriction on rational eigenvalues]
\label{ext-heisei-arithmetic}
Under Theorem~\ref{ext-heisei-global}, suppose the coefficients of $T$
belong to $F(x)$ for a number field $F$.  Every rational eigenvalue
$\lambda$ is algebraic over $F$ and satisfies
\[
 [F(\lambda):F]\le \dim W_{\mathfrak d,*}
       \le N_s(H_{\mathfrak d}).
\]
An associated nonzero rational eigenfunction can be chosen in
$F(\lambda)(x)$.
\end{corollary}

\begin{proof}
The exceptional set is algebraic over $F$.  Field embeddings permute its
points and preserve the integers $K_c$, so the monic denominator
$\mathfrak d$ in~\eqref{ext-heisei-denominator} belongs to $F[x]$.
The conjugated operator, leakage equations, and a kernel basis can all be
formed over $F$.  Consequently $K_*$ can be chosen over $F$, and its
characteristic polynomial has degree $\dim W_{\mathfrak d,*}$.  This
proves the degree bound.  A nonzero vector in the kernel of
$K_*-\lambda I$ can be chosen over $F(\lambda)$, proving the assertion
about the eigenfunction.
\end{proof}

\subsection{The previously excluded Jordan chain is recovered exactly}

Consider the rational operator
\begin{equation}
 T=x^4\partial^2+(2x^3+1)\partial+\frac1x.
 \label{ext-heisei-jordan-example}
\end{equation}
Its infinity data are $s=2$ and $h(z)=z(z+1)$.  Its only exceptional finite
point is $0$, where $\sigma_0=1$, $g_0(z)=z+1$, and $K_0=1$.
Thus $\mathfrak d=x$, $H_{\mathfrak d}=\{0,1\}$, and
$N_2(H_{\mathfrak d})=2$.  Direct calculation gives
\[
 xTx^{-1}=x^4\partial^2+\partial,
 \qquad
 W_{\mathfrak d,*}=\operatorname{Span}_{\C}\{x^{-1},1\},
\]
with restriction matrix
\[
 K_*=\begin{pmatrix}0&1\\0&0\end{pmatrix}
 \quad\text{in the ordered basis }(x^{-1},1).
\]
Indeed $T1=x^{-1}$ and $T(x^{-1})=0$.  The two-dimensional chain attains
the rational dimension bound.  The maximal \emph{polynomial} invariant
space is nevertheless zero: the polynomial resonance set is $\{0\}$,
but $T1$ is not polynomial.  Hence this extension recovers the chain in
its correct rational ambient space; it does not reinstate the false claim
that a polynomial initial vector alone guarantees a polynomial Jordan
chain.

\subsection{Why the local hypothesis cannot simply be discarded}

For
\[
 T=x(x^2+1)\partial
\]
the infinity shift is $s=2$ and $h(z)=z$.  The polynomial invariant-space
bound is one.  However, for every integer $k\ge0$,
\[
 f_k(x)=(1+x^{-2})^k,
 \qquad Tf_k=-2k f_k.
\]
Thus there are infinitely many distinct rational eigenvalues and rational
invariant spaces of arbitrarily large dimension.  At the zeros
$0,i,-i$ of the leading coefficient, $\sigma_c=0$, so the pole-raising
hypothesis fails.  This example proves that positive degree shift at
infinity alone cannot extend the polynomial spectral finiteness theorem
to rational functions.  Conversely, the failure of our sufficient local
hypothesis is not claimed to force an infinite rational spectrum.


\clearpage\part*{Source Concordance and Corrections}
\phantomsection\addcontentsline{toc}{part}{Source Concordance and Corrections}
\markboth{Source concordance}{Corrections}
\section*{Complete source concordance}
\phantomsection\addcontentsline{toc}{section}{Complete source concordance}
All sixty contributions below are reproduced mathematically in the indicated sections. Titles, abstracts, and local reference labels have been normalized for a single manuscript. The four corrections listed after the table are incorporated in the text.
\begin{longtable}{@{}p{0.075\textwidth}p{0.76\textwidth}r@{}}
\toprule Code & Contribution & Page\\\midrule\endhead

S01 & Sharp finite-moment reconstruction on zero-forcing graphs & \pageref{contrib-S01}\\

S02 & Optimal adversarial correction of triangle-minor tomography & \pageref{contrib-S02}\\

S03 & A third-order logarithmic test for partition matroids and cactus graphs & \pageref{contrib-S03}\\

S04 & The exact third-order inverse-power threshold for rank-two basis polynomials & \pageref{contrib-S04}\\

S05 & A sharp universal cubic inequality for real-stable polynomials & \pageref{contrib-S05}\\

S06 & Sharp fourth-order inverse-power thresholds on Lorentz cones & \pageref{contrib-S06}\\

S07 & Explicit rank-two alpha-permanent witnesses near the half-integer threshold & \pageref{contrib-S07}\\

S08 & All finite-order thresholds on the three-dimensional Lorentz cone & \pageref{contrib-S08}\\

S09 & Fifth-order Lorentz thresholds and five-point spherical energy & \pageref{contrib-S09}\\

S10 & Quantitative rigidity and exact certificates for critical rank-two permanents & \pageref{contrib-S10}\\

M01 & What a Stopped Coin Reveals: Exact Moment Rank of Sequential Mixture Experiments & \pageref{contrib-M01}\\

M02 & Logarithmic Recovery from Finite-Mean Coin Trees: A Sharp Minimax Theorem & \pageref{contrib-M02}\\

M03 & Finite Cost Does Not Control Inverse Difficulty: Random-Cap Mixtures with Prescribed Slow Rates & \pageref{contrib-M03}\\

M04 & Exact Robust Recovery of a Distributional Functional: Oscillation Duality and Small Atomic Certificates & \pageref{contrib-M04}\\

M05 & Sign Variation Controls Moment Inversion: Sharp Transport Exponents and Robust Crossing Certificates & \pageref{contrib-M05}\\

M06 & The Full Stretched-Exponential Cap Curve: Sharp Wasserstein Recovery Across a Tail Phase Transition & \pageref{contrib-M06}\\

M07 & Polynomial Observation Tails: Multilevel Recovery and Sharp Noncritical Minimax Rates & \pageref{contrib-M07}\\

M08 & The Information Spectrum of Random-Cap Binomial Mixtures: Exact Diagonalization and Sharp Tail Asymptotics & \pageref{contrib-M08}\\

M09 & The Critical Random-Cap Mixture Rate: A Common-Prior Proof of the Wasserstein Logarithm & \pageref{contrib-M09}\\

M10 & When Is Random-Cap Mixture Recovery Root-\texorpdfstring{$n$}{n}? An Exact Information Criterion and Critical Tail Hierarchies & \pageref{contrib-M10}\\

K01 & Integral Compression of Block Constant Perturbations & \pageref{contrib-K01}\\

K02 & Critical Groups of Uniform Blowups of Matched Trees & \pageref{contrib-K02}\\

K03 & Matching Numbers Control the Prime Parts of Tree Blowup Groups & \pageref{contrib-K03}\\

K04 & The Critical Group of Every Uniform Tree Blowup & \pageref{contrib-K04}\\

K05 & Smith Groups of Surjective Coupled Presentations & \pageref{contrib-K05}\\

K06 & Potentials and Exact Exponents of Nonuniform Tree Blowups & \pageref{contrib-K06}\\

K07 & Two Independent Scales in the Critical Groups of Tree Blowups & \pageref{contrib-K07}\\

K08 & Pendant Twin Attachments Transfer Critical Groups to Weighted Bases & \pageref{contrib-K08}\\

K09 & The Exceptional Prime in Mixed Clique--Independent Tree Blowups & \pageref{contrib-K09}\\

K10 & What All Uniform Blowup Critical Groups Can and Cannot Recover & \pageref{contrib-K10}\\

T01 & Weighted Star Defects Extremize Resistance, Heat Traces, and Tree Sums & \pageref{contrib-T01}\\

T02 & Exact Minimax Conductance Design under Fractional Edge Failures & \pageref{contrib-T02}\\

T03 & An Exponent-Independent Support Transition in Optimal Split-Graph Design & \pageref{contrib-T03}\\

T04 & Fractional-Power Optimal Designs and a Sharp Universal Laplacian Threshold & \pageref{contrib-T04}\\

T05 & Exact Resistance-Diameter Design on Trees with Heterogeneous Costs & \pageref{contrib-T05}\\

T06 & Sharp Random-Tree Energy Inflation:\ Exact Bond Ranges and Theta-Core Rigidity & \pageref{contrib-T06}\\

T07 & Exact Sign-Robustness Certificates for Resistance Reconstruction & \pageref{contrib-T07}\\

T08 & Universal Submodularity of Rank-One Inverse Gains:\ A Commutation Classification and a Three-Vertex Obstruction & \pageref{contrib-T08}\\

T09 & Higher Spectral Volumes of Random Trees:\ Sharp Quotient-Partition Ranges & \pageref{contrib-T09}\\

T10 & Arbitrarily Bad Greedy Link Addition on Unweighted Networks\ with Three Candidate Edges & \pageref{contrib-T10}\\

R01 & Overlap-Affinity Certificates for Hypergraph Product Deficits & \pageref{contrib-R01}\\

R02 & Global Marginal Inconsistency: Attained Dual Certificates and Quantitative Recovery & \pageref{contrib-R02}\\

R03 & Exact Weighted Entropic Projection onto Parity Polytopes & \pageref{contrib-R03}\\

R04 & Sharp Graph-Supported Entropy Stability and Positive Equality Cases & \pageref{contrib-R04}\\

R05 & The Exact Total-Correlation Remainder for Every Fractional Entropy Partition & \pageref{contrib-R05}\\

R06 & Local Entropy Stability Has a Different Hypergraph Spectrum & \pageref{contrib-R06}\\

R07 & Sharp Same-Family Comparisons of Symmetric Fractional Entropy Gaps & \pageref{contrib-R07}\\

R08 & Entropic Frustration on Signed Graphs: Exact Cactus Solutions and a Complexity Barrier & \pageref{contrib-R08}\\

R09 & Common Blocks and Quantitative Intersection in Pairwise Entropy Extremizers & \pageref{contrib-R09}\\

R10 & Constructive Rounding of Entropy Near-Extremizers to Disjoint Product Blocks & \pageref{contrib-R10}\\

H01 & A Degree-Filtration Bound for Algebraic Schr\"odinger Energies & \pageref{contrib-H01}\\

H02 & Pole Budgets and Finite Certificates for Rational Schr\"odinger Spectra & \pageref{contrib-H02}\\

H03 & The Monomial-Module Dichotomy for Polynomial Schr\"odinger Gauges & \pageref{contrib-H03}\\

H04 & Maximal Polynomial Modules and Their Hidden Jacobi Matrices & \pageref{contrib-H04}\\

H05 & Two-State Polynomial Modules: Divisibility, Decatic Classification, and Exact Nodal Indices & \pageref{contrib-H05}\\

H06 & A Spectral-Index Gap in Degree-Fourteen Two-State Polynomial Potentials & \pageref{contrib-H06}\\

H07 & Symmetry-Breaking Branches in the Complete Quartic-Generator Classification & \pageref{contrib-H07}\\

H08 & A Complete Three-State Decatic Module with Exact Levels Zero, Four, and Eight & \pageref{contrib-H08}\\

H09 & A Complex Nonexistence Certificate for Four-State Even Decatic Modules & \pageref{contrib-H09}\\

H10 & Sharp Resonance-Chain Bounds for Arbitrary-Order Rational Differential Spectra & \pageref{contrib-H10}\\

\bottomrule\end{longtable}
\section*{Incorporated corrections}
\phantomsection\addcontentsline{toc}{section}{Incorporated corrections}
\begin{enumerate}[leftmargin=*]
\item \textbf{K02, star counterexample.} Both the actual group $(\Z/2)^2\oplus\Z/8$ and the naively predicted group $\Z/4\oplus\Z/8$ have order $32$. Their generator ranks differ. The example is corrected in place; the main theorem is unchanged.
\item \textbf{K08, verification count.} The finalized source reports 120 completed compressed-core checks, with base size at most six. The earlier count of 200 with base size at most ten included an unfinished batch. Only the completed count is retained; the proof does not depend on those checks.
\item \textbf{H08, polynomial congruence.} The printed letters ``pmod'' in the discriminant remainder are replaced by the mathematical congruence notation. No coefficient or conclusion is changed.
\item \textbf{H10, Jordan-chain scope.} The maximal polynomial invariant space contains a generalized eigenvector only when its entire Jordan chain remains polynomial. The unrestricted rational-ambient reading is false. The corrected theorem and an explicit counterexample are included in place.
\end{enumerate}

\clearpage
\markboth{References}{References}
\phantomsection\addcontentsline{toc}{section}{References}
\begingroup\small
\begin{thebibliography}{999}
\bibitem{ref001} D. Burgarth and K. Maruyama, \emph{Indirect Hamiltonian identification through a small gateway}, New Journal of Physics 11 (2009), 103019. \url{https://arxiv.org/abs/0903.0612}.
\bibitem{ref002} D. Burgarth, \emph{Identifying combinatorially symmetric hidden Markov models}, preprint (2017). Primary preprint record: \url{https://arxiv.org/abs/1709.02932}. The citation is used for the established zero-forcing identifiability principle, not for a claim that its precise hypotheses coincide with ours.
\bibitem{ref003} J. Urschel, V.-E. Brunel, A. Moitra, and P. Rigollet, \emph{Learning determinantal point processes with moments and cycles}, Proceedings of Machine Learning Research 70 (2017), 3511--3520. \url{https://proceedings.mlr.press/v70/urschel17a.html}.
\bibitem{ref004} P. Sol\'e and T. Zaslavsky, \emph{A coding approach to signed graphs}, SIAM Journal on Discrete Mathematics (1994), DOI 10.1137/S0895480189174374. \url{https://doi.org/10.1137/S0895480189174374}. This reference is background for switching and coding, not a claim about the decoder proved here.
\bibitem{ref005} A. D. Scott and A. D. Sokal, \emph{Complete monotonicity for inverse powers of some combinatorially defined polynomials}, Acta Mathematica 213 (2014), 323--392. \url{https://arxiv.org/abs/1301.2449}. Expanded author version: \url{https://people.maths.ox.ac.uk/scott/Papers/compmono.pdf}. See in particular the small-exponent classifications in Section 7.
\bibitem{ref006} P. Br\"and\'en, \emph{Polynomials with the half-plane property and matroid theory}, Advances in Mathematics 216 (2007), 302--320. \url{https://arxiv.org/abs/math/0605678}.
\bibitem{ref007} J. Borcea, P. Br\"and\'en, and T. M. Liggett, \emph{Negative dependence and the geometry of polynomials}, Journal of the American Mathematical Society 22 (2009), 521--567. \url{https://arxiv.org/abs/0707.2340}.
\bibitem{ref008} D. Wei, Riesz kernels of hyperbolic polynomials: positivity, admissible exponents and Jordan rigidity, preprint (2026), version 2. \url{https://arxiv.org/abs/2609.31592}. Cited for current context only; its claims are not dependencies of this proof.
\bibitem{ref009} P. Br\"and\'en, Solutions to two problems on permanents, \emph{Linear Algebra Appl.} 436 (2012), 53--58. \url{https://arxiv.org/abs/1104.3531}.
\bibitem{ref010} D. Vere-Jones, A generalization of permanents and determinants, \emph{Linear Algebra Appl.} 111 (1988), 119--124. The generalized MacMahon identity is also stated in \cite[Theorem 2.4]{ref009}; only its self-contained squarefree proof is needed here.
\bibitem{ref011} A. Tumanov, Minimal biquadratic energy of five particles on a 2-sphere, \emph{Indiana Univ. Math. J.} 62 (2013), 1717--1731. \url{https://arxiv.org/abs/1208.5044}.
\bibitem{ref012} H. Teicher, \emph{Identifiability of mixtures}, Annals of Mathematical Statistics 32 (1961), 244--248. \url{https://doi.org/10.1214/aoms/1177705155}.
\bibitem{ref013} H. Teicher, \emph{Identifiability of finite mixtures}, Annals of Mathematical Statistics 34 (1963), 1265--1269. \url{https://doi.org/10.1214/aoms/1177708944}.
\bibitem{ref014} K. Tian, W. Kong and G. Valiant, \emph{Learning populations of parameters}, NeurIPS 2017. \url{https://arxiv.org/abs/1709.02707}.
\bibitem{ref015} R. K. Vinayak, W. Kong, G. Valiant and S. M. Kakade, \emph{Maximum likelihood estimation for learning populations of parameters}, ICML 2019. \url{https://proceedings.mlr.press/v97/vinayak19a.html}.
\bibitem{ref016} C. Scott, \emph{A rate of convergence for mixture proportion estimation, with application to learning from noisy labels}, AISTATS 2015. \url{https://web.eecs.umich.edu/~cscott/pubs/aistats15.pdf}.
\bibitem{ref017} D. L. Donoho, \emph{Statistical estimation and optimal recovery}, Annals of Statistics 22 (1994), 238--270. \url{https://doi.org/10.1214/aos/1176325367}. Author-hosted text: \url{https://web.stanford.edu/dept/statistics/cgi-bin/donoho/wp-content/uploads/2018/08/SEOR.pdf}.
\bibitem{ref018} Y. Wu and P. Yang, \emph{Optimal estimation of Gaussian mixtures via denoised method of moments}, Annals of Statistics 48 (2020), 1981--2007. \url{https://arxiv.org/abs/1807.07237}. Author-hosted manuscript: \url{https://www.stat.yale.edu/~yw562/preprints/GM.pdf}.
\bibitem{ref019} M. B. Giles, \emph{Multilevel Monte Carlo methods}, Acta Numerica 24 (2015), 259--328. \url{https://doi.org/10.1017/S096249291500001X}.
\bibitem{ref020} M. B. Giles, T. Nagapetyan and K. Ritter, \emph{Multilevel Monte Carlo approximation of distribution functions and densities}, SIAM/ASA Journal on Uncertainty Quantification 3 (2015), 267--295. \url{https://doi.org/10.1137/140960086}.
\bibitem{ref021} M.-M. Derriennic, \emph{Sur l'approximation de fonctions int\'egrables sur $[0,1]$ par des polyn\^omes de Bernstein modifi\'es}, Journal of Approximation Theory 31 (1981), 325--343. \url{https://doi.org/10.1016/0021-9045(81)90101-5}.
\bibitem{ref022} P. Diaconis, K. Khare and L. Saloff-Coste, \emph{Gibbs sampling, exponential families and orthogonal polynomials}, Statistical Science 23 (2008), 151--178. \url{https://arxiv.org/abs/0808.3852}.
\bibitem{ref023} R. D. Gill and B. Y. Levit, \emph{Applications of the van Trees inequality: a Bayesian Cram\'er--Rao bound}, Bernoulli 1 (1995), 59--79. Author manuscript: \url{https://dspace.library.uu.nl/bitstream/handle/1874/1190/van_trees.pdf}.
\bibitem{ref024} M. Tsang, \emph{Ziv--Zakai-type error bounds for general statistical models} (2023). \url{https://arxiv.org/abs/2306.08660}.
\bibitem{ref025} J. Rousseau and C. Scricciolo, \emph{Wasserstein convergence in Bayesian and frequentist deconvolution models}, Ann. Statist. 52 (2024). \url{https://arxiv.org/abs/2309.15300}.
\bibitem{ref026} M. Trabs, \emph{Information bounds for inverse problems with application to deconvolution and L\'evy models}, Ann. Inst. H. Poincar\'e Probab. Statist. 51 (2015), 1620--1650. \url{https://doi.org/10.1214/14-AIHP627}.
\bibitem{ref027} J. Dedecker, A. Fischer and B. Michel, \emph{Improved rates for Wasserstein deconvolution with ordinary smooth error in dimension one}, Electron. J. Statist. 9 (2015), 234--265. \url{https://doi.org/10.1214/15-EJS997}; author version \url{https://arxiv.org/abs/1404.0646}.
\bibitem{ref028} X. Dong, W. Guo and G. Jiang, \emph{On the critical group of the $k$-partite graph}, 2024. \url{https://arxiv.org/abs/2409.02654}.
\bibitem{ref029} C. A. Alfaro, H. Corrales and C. E. Valencia, \emph{Critical ideals of signed graphs with twin vertices}, Advances in Applied Mathematics 86 (2017), 99--131. \url{https://arxiv.org/abs/1504.06257}.
\bibitem{ref030} R. P. Stanley, \emph{Smith normal form in combinatorics}, 2016. \url{https://arxiv.org/abs/1602.00166}.
\bibitem{ref031} Keio, \emph{Integral compression of block constant perturbations}, tournament Round 1, 2026. All needed arguments are reproduced here.
\bibitem{ref032} B. Jacobson, A. Niedermaier and V. Reiner, \emph{Critical groups for complete multipartite graphs and Cartesian products of complete graphs}, Journal of Graph Theory 44 (2003), 231--250. \url{https://doi.org/10.1002/jgt.10139}.
\bibitem{ref033} F. Shokrieh, \emph{The monodromy pairing and discrete logarithm on the Jacobian of finite graphs}, \url{https://arxiv.org/abs/0907.4764}.
\bibitem{ref034} X.-h. Tan, S.-y. Wang, H.-y. Qian, \emph{The Critical Group of the Corona of Trees and Cycles}, 2018, \url{https://doi.org/10.12783/dtcse/icmsa2018/23275}.
\bibitem{ref035} H. Corrales, C. E. Valencia, \emph{Critical ideals of trees}, 2015, revised 2017, \url{https://arxiv.org/abs/1504.06239}.
\bibitem{ref036} A. Ghosh, S. Boyd and A. Saberi, \emph{Minimizing Effective Resistance of a Graph}, SIAM Review 50 (2008), 37--66. \url{https://web.stanford.edu/~boyd/papers/eff_res.html}.
\bibitem{ref037} K. Xu, K. C. Das and X.-D. Zhang, \emph{Ordering connected graphs by their Kirchhoff indices}. \url{https://arxiv.org/abs/1602.07039}.
\bibitem{ref038} S. Tamura, \emph{Edge-defect matrices and stability of the Kirchhoff index for complete graphs with deleted edges} (2026). \url{https://arxiv.org/abs/2608.05924}.
\bibitem{ref039} S. Boyd, \emph{Convex Optimization of Graph Laplacian Eigenvalues}, Proceedings ICM (2006). \url{https://web.stanford.edu/~boyd/papers/cvx_opt_graph_lapl_eigs.html}.
\bibitem{ref040} Taisho, \emph{Weighted Star Defects Extremize Resistance, Heat Traces, and Tree Sums}, research round 1 (2026). The needed result is reproved above.
\bibitem{ref041} \emph{Resistance distances and the Moon-type formula of a vertex-weighted complete split graph}, Discrete Applied Mathematics (2024). Publisher record: \url{https://www.sciencedirect.com/science/article/pii/S0166218X2400341X}.
\bibitem{ref042} A. P. Riascos and J. L. Mateos, \emph{Fractional dynamics on networks: Emergence of anomalous diffusion and L\'evy flights}, Physical Review E 90 (2014), 032809. \url{https://doi.org/10.1103/PhysRevE.90.032809}.
\bibitem{ref043} E. Estrada, \emph{Path Laplacians versus fractional Laplacians as nonlocal operators on networks}, New Journal of Physics 23 (2021), 073049. \url{https://doi.org/10.1088/1367-2630/ac14ac}.
\bibitem{ref044} B. Ben-Moshe, M. Elkin, E. Omri and L.-A. Gottlieb, \emph{Optimizing budget allocation for center and median points}, Theoretical Computer Science (2016). \url{https://doi.org/10.1016/j.tcs.2016.02.013}. Preprint: \url{https://arxiv.org/abs/1406.2107}.
\bibitem{ref045} A. Kassel, R. Kenyon and W. Wu, Random two-component spanning forests, \emph{Ann. Inst. H. Poincar\'e Probab. Stat.} 51 (2015), 1457--1464. \url{https://doi.org/10.1214/14-AIHP625}.
\bibitem{ref046} H. Richman, F. Shokrieh and C. Wu, Counting two-forests and random cut size via potential theory, arXiv:2308.03859. In particular Section 3.2 and the classical resistance formula in Section 4.1. \url{https://arxiv.org/abs/2308.03859}.
\bibitem{ref047} R. Lyons and Y. Peres, \emph{Probability on Trees and Networks}, Cambridge University Press, 2016, Chapter 4. Author-hosted text: \url{https://rdlyons.pages.iu.edu/prbtree/book.pdf}.
\bibitem{ref048} M. Hlad\'ik, Tolerances, robustness and parametrization of matrix properties related to optimization problems, \emph{Optimization} 68 (2019), 667--690. \url{https://arxiv.org/abs/1709.07629}.
\bibitem{ref049} R. B. Bapat, Resistance matrix of a weighted graph, \emph{MATCH Commun. Math. Comput. Chem.} 50 (2004), 73--82. Journal archive: \url{https://match.pmf.kg.ac.rs/electronic_versions/Match50/match50_73-82.pdf}.
\bibitem{ref050} S. J. Kirkland and M. Neumann, The M-matrix group generalized inverse problem for weighted trees, \emph{SIAM J. Matrix Anal. Appl.} 19 (1998), 226--234. \url{https://doi.org/10.1137/S0895479896304927}. This prior classification rules out claiming the weighted-tree inverse-sign problem as new.
\bibitem{ref051} L. F. O. Chamon and A. Ribeiro, Approximate Supermodularity Bounds for Experimental Design, \emph{NeurIPS} 2017. \url{https://arxiv.org/abs/1711.01501}.
\bibitem{ref052} M. A. Achterberg and R. E. Kooij, On the Non-Submodularity of the Problem of Adding Links to Minimize the Effective Graph Resistance, \emph{J. Graph Algorithms Appl.} 30 (2026), 109--132. \url{https://doi.org/10.7155/jgaa.v30i1.3192}; preprint \url{https://arxiv.org/abs/2501.03363}.
\bibitem{ref053} P. Chebotarev and E. Shamis, Matrix-Forest Theorems. \url{https://arxiv.org/abs/math/0602575}.
\bibitem{ref054} M. Derezi\'nski and M. K. Warmuth, Reverse Iterative Volume Sampling for Linear Regression, \emph{J. Mach. Learn. Res.} 19 (2018), 1--39. \url{https://jmlr.org/papers/v19/17-781.html}.
\bibitem{ref055} Taisho, Sharp Random-Tree Energy Inflation: Exact Bond Ranges and Theta-Core Rigidity, tournament research round 6, 24 October 2026. The present paper is self-contained.
\bibitem{ref057} P. Yalamanchili and A. Bhaskara, Clustering-Based Edge Augmentation for Minimizing the Kirchhoff Index, \emph{AISTATS}, PMLR 300 (2026), 3745--3753. \url{https://proceedings.mlr.press/v300/yalamanchili26a.html}.
\bibitem{ref058} M. Predari, L. Berner, R. Kooij and H. Meyerhenke, Greedy optimization of resistance-based graph robustness with global and local edge insertions, \emph{Social Network Analysis and Mining} 13 (2023), article 130. \url{https://doi.org/10.1007/s13278-023-01137-1}.
\bibitem{ref059} H. Finner, A generalization of H\"older's inequality and some probability inequalities, \emph{Ann. Probab.} 20 (1992), 1893--1901. \url{https://doi.org/10.1214/aop/1176989534}.
\bibitem{ref060} E. A. Carlen and D. Cordero-Erausquin, Subadditivity of the entropy and its relation to Brascamp--Lieb type inequalities. \url{https://arxiv.org/abs/0710.0870}.
\bibitem{ref061} J. Liu, T. A. Courtade, P. Cuff and S. Verd\'u, Brascamp--Lieb inequality and its reverse: an information theoretic view. \url{https://arxiv.org/abs/1605.02818}.
\bibitem{ref062} J. Bennett and S. Oh, Stability for the adjoint Finner inequality and improved lower bounds on the $k$-plane transform (2026). \url{https://arxiv.org/abs/2610.29041}.
\bibitem{ref063} M. J. Wainwright and M. I. Jordan, Graphical models, exponential families, and variational inference, \emph{Foundations and Trends in Machine Learning} 1 (2008), 1--305. Author-hosted text: \url{https://people.eecs.berkeley.edu/~jordan/sail/readings/wainwright-jordan-fnt.pdf}.
\bibitem{ref064} J. Feldman, M. J. Wainwright and D. R. Karger, Using linear programming to decode binary linear codes, \emph{IEEE Trans. Inform. Theory} 51 (2005), 954--972. \url{https://doi.org/10.1109/TIT.2004.842696}. Author publication record: \url{https://people.csail.mit.edu/jonfeld/pubs.html}.
\bibitem{ref065} M. J. Wainwright, T. S. Jaakkola and A. S. Willsky, A new class of upper bounds on the log partition function, \emph{IEEE Trans. Inform. Theory} 51 (2005), 2313--2335. \url{https://people.csail.mit.edu/tommi/papers/WaiJaaWil05_Upper.pdf}.
\bibitem{ref066} G. Jakhar, G. R. Kurri, S. Chillara and V. M. Prabhakaran, Fractional subadditivity of submodular functions: equality conditions and their applications (2025). \url{https://arxiv.org/abs/2501.12058}.
\bibitem{ref067} M. Abo Khamis, P. G. Kolaitis, H. Q. Ngo and D. Suciu, Bag query containment and information theory, PODS 2020, Lemma 3.10. \url{https://homes.cs.washington.edu/~suciu/bag-containment-pods-2020.pdf}.
\bibitem{ref068} M. Madiman and P. Tetali, Information inequalities for joint distributions, with interpretations and applications, \emph{IEEE Trans. Inform. Theory} 56 (2010), 2699--2713. \url{https://www.stat.yale.edu/~mm888/Pubs/2010/MT10-joint.pdf}.
\bibitem{ref069} M. Rzeszut and M. Wojciechowski, Hoeffding decomposition in $H^1$ spaces. \url{https://arxiv.org/abs/1906.01405}.
\bibitem{ref070} A. Makur and L. Zheng, Linear bounds between contraction coefficients for $f$-divergences. \url{https://arxiv.org/abs/1510.01844}.
\bibitem{ref071} T. Nishiyama and I. Sason, On relations between the relative entropy and $\chi^2$-divergence, generalizations and applications. \url{https://arxiv.org/abs/2004.11197}.
\bibitem{ref072} M. Madiman and P. Tetali, \emph{Information inequalities for joint distributions, with interpretations and applications}, IEEE Trans. Inform. Theory 56 (2010), 2699--2713. \url{https://arxiv.org/abs/0901.0044}.
\bibitem{ref073} Q. Chen and R. W. Yeung, \emph{Partition-symmetrical entropy functions}, IEEE Trans. Inform. Theory 62 (2016), 5385--5402. \url{https://arxiv.org/abs/1407.7405}.
\bibitem{ref074} M. J. Wainwright and M. I. Jordan, \emph{Graphical models, exponential families, and variational inference}, Foundations and Trends in Machine Learning 1 (2008), 1--305. \url{https://people.eecs.berkeley.edu/~wainwrig/Papers/WaiJor08_FTML.pdf}.
\bibitem{ref075} M. R. Garey, D. S. Johnson, and L. Stockmeyer, \emph{Some simplified NP-complete graph problems}, Theoret. Comput. Sci. 1 (1976), 237--267. \url{https://doi.org/10.1016/0304-3975(76)90059-1}.
\bibitem{ref076} J. Weed, \emph{An explicit analysis of the entropic penalty in linear programming}, Proc. COLT, PMLR 75 (2018), 1841--1855. \url{https://proceedings.mlr.press/v75/weed18a.html}.
\bibitem{ref077} P. G\'acs and J. K\"orner, \emph{Common information is far less than mutual information}, Problems of Control and Information Theory 2 (1973), 149--162. Author-hosted copy: \url{https://cs-web.bu.edu/faculty/gacs/papers/commoninf.pdf}.
\bibitem{ref078} P. Caputo, G. Menz, and P. Tetali, \emph{Approximate tensorization of entropy at high temperature}, Ann. Fac. Sci. Toulouse Math. 24 (2015), 691--716. \url{https://arxiv.org/abs/1405.0608}.
\bibitem{ref079} I. Bardet, A. Capel, and C. Rouz\'e, \emph{Approximate tensorization of the relative entropy for noncommuting conditional expectations}, 2020. \url{https://arxiv.org/abs/2001.07981}. Cited for the broader framework, not used as a black-box theorem.
\bibitem{ref080} P. B. Acosta-Hum\'anez, D. Bl\'azquez-Sanz and H. Venegas-G\'omez, \emph{Liouvillian solutions for second order linear differential equations with polynomial coefficients}, arXiv:1908.07666, especially the polynomial-potential criterion and Theorem 4.2. \url{https://arxiv.org/abs/1908.07666}.
\bibitem{ref081} C. M. Bender and G. V. Dunne, \emph{Quasi-exactly solvable systems and orthogonal polynomials}, J. Math. Phys. 37 (1996), 6--11. The sextic family is classical, not introduced here. \url{https://doi.org/10.1063/1.531373}.
\bibitem{ref082} J. J. Kovacic, \emph{An algorithm for solving second order linear homogeneous differential equations}, J. Symbolic Computation 2 (1986), 3--43. \url{https://doi.org/10.1016/S0747-7171(86)80010-4}.
\bibitem{ref083} A. V. Turbiner, \emph{Quasi-exactly-solvable differential equations}. \url{https://arxiv.org/abs/hep-th/9409068}.
\bibitem{ref084} G. Post and A. V. Turbiner, \emph{Classification of linear differential operators with an invariant subspace of monomials}, Russian J. Math. Phys. 3 (1995), 113--122. \url{https://arxiv.org/abs/funct-an/9307001}.
\bibitem{ref085} V. M. Tkachuk, \emph{Quasi-exactly solvable potentials with two known eigenstates}. \url{https://arxiv.org/abs/quant-ph/9801021}.
\bibitem{ref086} W. Schweiger and W. H. Klink, \emph{Polynomial potentials and nilpotent groups}, Eur. Phys. J. Plus 140 (2025), 559. \url{https://doi.org/10.1140/epjp/s13360-025-06485-9}.
\bibitem{ref087} V. M. Tkachuk, \emph{Supersymmetric approach for generating quasi-exactly solvable potentials with arbitrary two known eigenstates}. \url{https://arxiv.org/abs/quant-ph/0107115}.
\bibitem{ref088} D. G\'omez-Ullate, N. Kamran and R. Milson, \emph{Quasi-exact solvability beyond the sl(2) algebraization}. \url{https://arxiv.org/abs/nlin/0601053}.
\bibitem{ref089} S. A. Abramov, \emph{Rational solutions of linear differential and difference equations with polynomial coefficients}, USSR Comput. Math. Math. Phys. 29 (1989), no. 6, 7--12. \url{https://doi.org/10.1016/S0041-5553(89)80002-3}.
\bibitem{integration-lee-kennedy} J. Lee and E. H. Kennedy, \emph{Minimax Rates for Learning Smooth Populations of Parameters} (2026). \url{https://arxiv.org/abs/2610.04135}.
\bibitem{ext-showa-Pate} T. H. Pate, \emph{The best lower bound for the permanent of a correlation
matrix of rank two}, Linear and Multilinear Algebra \textbf{51} (2003),
263--278. \url{https://doi.org/10.1080/0308108031000081304}.
\bibitem{ext-reiwa-CS} P. Caputo and A. Sinclair,
\emph{Entropy production in nonlinear recombination models},
Bernoulli \textbf{24} (2018), no.~4B, 3246--3282.
arXiv:1609.06897.
\bibitem{ext-reiwa-CP} P. Caputo and D. Parisi,
\emph{Nonlinear recombinations and generalized random transpositions},
Annales Henri Lebesgue \textbf{7} (2024), 1245--1299.
doi:10.5802/ahl.219; arXiv:2207.04775.
\bibitem{ext-reiwa-KS} J. Kim and I. Seo,
\emph{Cutoff phenomenon for inhomogeneous nonlinear recombination in
arbitrary finite product spaces},
Journal of Statistical Physics \textbf{193} (2026), article~69.
doi:10.1007/s10955-026-03640-z; arXiv:2510.02136.
\end{thebibliography}
\endgroup
\end{document}
